\pdfoutput=1
\documentclass[11pt]{amsart}
\usepackage[margin=1.1in]{geometry}
\usepackage{amsmath,amssymb,amsthm,mathrsfs,mathtools}
\usepackage{enumitem}
\usepackage{microtype}
\usepackage{hyperref}
\usepackage{aliascnt}
\usepackage[nameinlink,capitalize]{cleveref}
\hypersetup{colorlinks=true,linkcolor=blue,citecolor=blue,urlcolor=blue}
\allowdisplaybreaks
\pdfstringdefDisableCommands{%
  \def\({}%
  \def\){}%
  \def\tilde#1{#1}%
  \def\mathfrak#1{#1}%
  \def\mathrm#1{#1}%
  \def\mathbb#1{#1}%
  \def\tau{tau}%
  \def\lambda{lambda}%
  \def\infty{infty}%
}

\theoremstyle{plain}
\newtheorem{theorem}{Theorem}[section]
\newaliascnt{lemma}{theorem}
\newtheorem{lemma}[lemma]{Lemma}
\aliascntresetthe{lemma}
\newaliascnt{proposition}{theorem}
\newtheorem{proposition}[proposition]{Proposition}
\aliascntresetthe{proposition}
\newaliascnt{corollary}{theorem}
\newtheorem{corollary}[corollary]{Corollary}
\aliascntresetthe{corollary}
\theoremstyle{definition}
\newaliascnt{definition}{theorem}
\newtheorem{definition}[definition]{Definition}
\aliascntresetthe{definition}
\theoremstyle{remark}
\newaliascnt{remark}{theorem}
\newtheorem{remark}[remark]{Remark}
\aliascntresetthe{remark}

\crefname{theorem}{theorem}{theorems}
\Crefname{theorem}{Theorem}{Theorems}
\crefname{lemma}{lemma}{lemmas}
\Crefname{lemma}{Lemma}{Lemmas}
\crefname{proposition}{proposition}{propositions}
\Crefname{proposition}{Proposition}{Propositions}
\crefname{corollary}{corollary}{corollaries}
\Crefname{corollary}{Corollary}{Corollaries}
\crefname{definition}{definition}{definitions}
\Crefname{definition}{Definition}{Definitions}
\crefname{remark}{remark}{remarks}
\Crefname{remark}{Remark}{Remarks}

\newcommand{\dd}{\mathrm d}
\newcommand{\RR}{\mathbb R}
\newcommand{\C}{\mathbb C}
\newcommand{\supp}{\mathrm{supp}}
\newcommand{\nn}{\nabla}
\newcommand{\tr}{\mathrm{tr}}
\newcommand{\fint}{\mathop{\ooalign{\(\int\)\cr\hidewidth\(-\)\hidewidth\cr}}}

\title[Type II Regularity]{Type II Regularity for Navier--Stokes}
\author{Guillem Duran Ballester}
\email{guillem@fragile.tech}
\date{}

\begin{document}

\begin{abstract}
We establish a local retained-branch exclusion criterion for the
three-dimensional Navier--Stokes equations near a possible singular point and
record the named exit rows together with their local closures.  From a positive
Caffarelli--Kohn--Nirenberg concentration sequence we select renormalized local
windows around a retained core and analyze the resulting alternatives by
compactness, pressure reconstruction, and local energy estimates.  The argument
combines a compact single-core exclusion criterion, a repaired-gauge
representation for nondegenerate concentration cores, local Calder\'on--Zygmund
pressure control, a Caccioppoli estimate on compact cylinders, multibubble and
cascade reductions, and scale-collapse cost estimates.  The retained compact
scale-collapse branch is closed by local state-space routing into the
scale-rigid discharge, so no additional stationary or self-similar
classification theorem is used implicitly.  Thin negative drift is closed as an
independent terminal row by diffuse occupation-state routing.  Branches that
leave the retained state through selected compactness-package failure are now
routed to more specific local rows.  Modulation failure, exterior
concentration, failure of the retained spacetime activity floor, and
noncanonical carrier/cost failure
are not treated as hidden inputs; they are listed as named exits and closed
by their own local closure theorems.
\end{abstract}

\maketitle

\setcounter{tocdepth}{2}
\tableofcontents


\clearpage
\section{Introduction}

The three-dimensional Navier--Stokes regularity problem is governed by the
interaction between local energy control, scale-critical concentration, and the
possible formation of singular structures at arbitrarily small parabolic
scales.  Since Leray's construction of global weak solutions
\cite{Leray1934}, the theory of suitable weak solutions and the partial
regularity theorem of Scheffer and Caffarelli--Kohn--Nirenberg have provided the
basic local language for this interaction
\cite{Scheffer1976,CaffarelliKohnNirenberg1982,Lin1998}.  In that theory,
singularity formation is detected by scale-invariant velocity-pressure
quantities, and the critical \(L^3\) scale plays a central role alongside the
local energy inequality, pressure reconstruction, and compactness.  The
analytic background also includes the classical local and critical-space
theory for Navier--Stokes
\cite{Ladyzhenskaya1969,Temam1977,ConstantinFoias1988,Kato1984,LemarieRieusset2002},
the \(L_{3,\infty}\) regularity theorem of Escauriaza--Seregin--\v Sver\'ak
\cite{EscauriazaSereginSverak2003}, and the rigidity theorem of
Ne\v cas--R\r u\v zi\v cka--\v Sver\'ak for Leray self-similar profiles
\cite{NecasRuzickaSverak1996}.  These classical rigidity results are cited as
background; the retained Type II branch is closed by the local state-space
reductions below, not by importing an external stationary or self-similar
classification theorem.

This paper studies the local Type II alternative near a possible singular point.
Starting from a positive Caffarelli--Kohn--Nirenberg concentration sequence, we
select renormalized local windows around a retained core and analyze all
branches that can arise from the scale-center dynamics.  The result in the
present version is a retained-state exclusion theorem together with an explicit
ledger of named exits and a closure theorem for each remaining ledger row.
Pressure reconstruction, repaired-gauge compactness, terminal profile decoupling, cost
comparison, and local state-space closure are proved on retained branches or
used to route a first failure to one of those named exits.  No additional
stationary, ancient, or self-similar classification theorem is invoked
implicitly.

The proof begins with local concentration and the Type I/Type II separation.
For a nondegenerate retained core, the solution is rewritten in local
scale-translation variables.  The repaired-gauge representation produces a
renormalized Navier--Stokes equation with controlled modulation terms, local
pressure reconstruction modulo functions of time, and compact-window suitable
estimates.  Within this representation, a compact single-core branch with a
retained positive velocity core and vanishing localized dissipation is
impossible: local compactness gives a spatially constant limit, the zero
constant contradicts retained velocity by strong \(L^3\) convergence, and a
nonzero constant gives a bounded selected-window limit, which leaves the strict
selected-window retained branch and is later physically routed by the
bounded-window closure theorem.  Pressure-only harmonic concentration is routed to
the pressure/exterior rows rather than closed by the single-core velocity
argument.

The remaining sections remove the ways in which a Type II branch could avoid
that compact single-core mechanism.  Caccioppoli estimates convert finite outer
velocity-pressure control into inner windowed \(H^1\) control, so rough-core
loss is returned to an explicit compact-cylinder alternative.  Multibubble and
cascade arguments separate active profiles, same-point cascades, and
gauge-degenerate branches from a single retained core.  Scale-collapse branches
are treated by localized cost identities: finite scale-collapsing cost, and the
corresponding finite absolute cost, are incompatible with a nonvanishing
retained core, while thick autonomous windows lead only to the stated
modulation, compactness, or retained compact scale-collapse state; the latter
is routed into the scale-rigid discharge rather than into a new Liouville
assumption.  The diffuse thin-drift boundary face is handled by normalized
negative-drift occupation blocks: persistent windows return to the thick-window
ledger, while genuinely diffuse blocks make the negative scale term a
vanishing local perturbation and route the branch to another local row or to a
closed retained row.  The terminal profile and corrected monotonicity sections
then turn these local reductions into a branchwise classification: retained states are
eliminated, while first failures outside the retained state are recorded as
named exits.

The conclusion should therefore be read in this local state-space sense.  In the
retained branch, a positive local Type II concentration sequence cannot survive
the repaired-gauge, compactness, multibubble, rough-core, compact-cost, and
retained scale-collapse reductions.  A failure of one of the local analytic rows
is not left as an unstated input; it is recorded as a named state-space
exit.  In the final assembly, the remaining named exits are then closed by
their own local routing arguments, giving the stated no-local-Type-II
conclusion without importing global critical estimates.

\clearpage

\section{Local Single-Core Criterion}\label{sec:single-core-criterion}

\subsection{Setup and definitions}

For \(R>0\), write
\[
  Q_R=B_R\times(-R^2,0).
\]
For a suitable pair \((v,q)\) on \(Q_R\), define
\[
  A_v(R)=R^{-1}\operatorname*{ess\,sup}_{-R^2<s<0}
  \int_{B_R}|v(y,s)|^2\,dy,
\]
\[
  E_v(R)=R^{-1}\int_{Q_R}|\nabla v|^2\,dy\,ds,
\]
\[
  C_v(R)=R^{-2}\int_{Q_R}|v|^3\,dy\,ds,
\]
\[
  D_q(R)=R^{-2}\int_{-R^2}^0\int_{B_R}
  |q(y,s)-(q)_{B_R}(s)|^{3/2}\,dy\,ds.
\]
All pressure quantities are understood modulo functions of time.

\begin{definition}[Bounded selected-window limit]\label{paper1:def:bounded-window-limit}
Let \((V_n,P_n)\) be a sequence of represented local rescalings on a compact
cylinder \(Q_r\).  A selected-window subsequence has a bounded selected-window limit
on \(Q_r\) if, after passing to a subsequence,
\[
  V_n\to V_*
  \qquad\text{strongly in }L^3(Q_r),
\]
where \(V_*\in L^\infty(Q_r)\).  The limit is nonzero if
\(V_*\not\equiv0\) on \(Q_r\).
\end{definition}

\begin{definition}[Strict selected-window Type II branch]\label{paper1:def:typeII}
A represented selected-window branch is called a strict selected-window Type II
branch if it has positive local Caffarelli--Kohn--Nirenberg concentration on
some fixed compact cylinder and no selected-window subsequence has a nonzero
bounded selected-window limit on any compact cylinder in the sense of
\cref{paper1:def:bounded-window-limit}.
\end{definition}

\begin{remark}[Strict selected-window Type II versus physical Type II]
\label{paper1:rem:strict-selected-window-vs-physical-typeII}
\Cref{paper1:def:typeII} is a local branch condition inside the
selected-window state-space decomposition.  It is not the definition of a
physical Type II singular germ; physical local Type II germs are defined later
in \cref{paper6:def:physical-local-typeII-germ} by failure of the local
Type I scale-invariant bound at a singular point.  A bounded selected-window
limit is therefore not discarded as ``not physically Type II''; it leaves the
strict selected-window branch and is routed through the bounded-window local
closure row by
\cref{paper3:lem:scale-rigid-bounded-limit-exit,paper6:thm:bounded-selected-window-physical-routing}.
Only after that row is closed, or reselected into one of the named local
alternatives, may the proof continue in the strict selected-window branch.
\end{remark}

\begin{definition}[Compact single-core vanishing-dissipation branch]\label{paper1:def:single-core-branch}
Fix radii \(0<r_*<R_*\), constants \(M<\infty\), \(\eta_0>0\), and a sequence
of represented local rescalings \((V_n,P_n,a_n,b_n)\) on \(Q_{R_*}\).  The
sequence is a compact single-core vanishing-dissipation branch if:
\begin{enumerate}[label=\textup{(\roman*)}]
\item each \((V_n,P_n)\) is suitable on \(Q_{R_*}\) and solves the represented
      Navier--Stokes equation
      \begin{equation}\label{paper1:eq:represented-ns}
      \partial_s V_n+(V_n\cdot\nabla)V_n+\nabla P_n
      =\nu\Delta V_n+a_n(s)(V_n+y\cdot\nabla V_n)+b_n(s)\cdot\nabla V_n,
      \qquad \nabla\cdot V_n=0;
      \end{equation}
\item the compact-cylinder suitable estimates are uniformly bounded:
      \[
      A_{V_n}(R_*)+E_{V_n}(R_*)+C_{V_n}(R_*)+D_{P_n}(R_*)\le M;
      \]
\item the selected core has retained positive velocity concentration on \(Q_{r_*}\):
      \[
      C_{V_n}(r_*)\ge\eta_0
      \qquad\text{for all }n;
      \]
      Pressure-only concentration is not included in this compact single-core
      row; it is routed to the harmonic/exterior-pressure alternative.
\item the localized scale and center are nondegenerate for the selected core,
      so the representation theorem gives the variables in
      \eqref{paper1:eq:represented-ns};
\item the modulation coefficients are uniformly bounded on the selected cylinder:
      \[
      \|a_n\|_{L^\infty(-R_*^2,0)}+\|b_n\|_{L^\infty(-R_*^2,0)}\le M;
      \]
\item the branch has vanishing localized dissipation in the sense of
      \cref{paper1:def:localized-dissipation} below.
\end{enumerate}
\end{definition}

\begin{definition}[Vanishing localized dissipation]\label{paper1:def:localized-dissipation}
Let \(r_*<\rho<R_*\).  Selected windows are first translated and rescaled to the
fixed cylinders \(Q_{r_*}\Subset Q_\rho\Subset Q_{R_*}\).  The localized
dissipation on \(Q_\rho\) is
\[
  \mathcal D_n(\rho)=\int_{Q_\rho}|\nabla V_n|^2\,dy\,ds.
\]
The branch has vanishing localized dissipation if, after passing to a
selected-window subsequence, there is \(\rho\in(r_*,R_*)\) such that
\[
  \mathcal D_n(\rho)\to0.
\]
\end{definition}

\subsection{Local analytic inputs}

The compact single-core argument is based on the following local analytic statements.

\begin{theorem}[Local concentration dichotomy]
\label{paper1:thm:paper0-dichotomy}
Let \((u,p)\) be a suitable weak solution on
\(\mathbb R^3\times(T-\delta,T)\), and suppose \(\Sigma(T)\ne\emptyset\).
Then, for every \(x_0\in\Sigma(T)\), using the pressure-normalized quantities
\(C,D\), there are \(\eta>0\) and \(r_n\downarrow0\) such that
\[
  C((x_0,T),r_n)+D((x_0,T),r_n)\ge\eta
  \qquad\text{for all }n.
\]
Moreover, at each singular point exactly one of the following alternatives
holds: a local Type I scale-invariant bound holds, or the point lies in the
local Type II alternative, meaning that no such local Type I bound holds.
\end{theorem}

\begin{proof}
By the Caffarelli--Kohn--Nirenberg epsilon-regularity criterion, there is a
universal \(\varepsilon_{\rm CKN}>0\) such that if
\[
  C((x_0,T),r)+D((x_0,T),r)<\varepsilon_{\rm CKN}
\]
for all sufficiently small \(r\), then \(u\) is locally bounded in a smaller
backward cylinder ending at \((x_0,T)\).  Since \(x_0\in\Sigma(T)\), this
regularity conclusion is impossible.  Hence for every \(k\) there exists
\(0<r_k<1/k\) with
\[
  C((x_0,T),r_k)+D((x_0,T),r_k)\ge\varepsilon_{\rm CKN}.
\]
After passing to a decreasing subsequence and setting
\(\eta=\varepsilon_{\rm CKN}\), the displayed concentration sequence follows.
The Type I/Type II split is the dichotomy between existence and failure of the
local Type I scale-invariant bound at the same point.
\end{proof}

\begin{proposition}[Passage to selected-window branches]
\label{paper1:prop:paper0-typeII-reduction}
Let \((x_0,T)\) be a concentration point in the physical local Type II
alternative of \cref{paper1:thm:paper0-dichotomy}.  After choosing the positive
concentration scales given there and applying the local scale-translation
selection used in the present section to a nondegenerate compact single core,
the represented selected windows have positive local CKN concentration on a
fixed compact cylinder.  Exactly one of the following selected-window
alternatives is then recorded:
\begin{enumerate}[label=\textup{(\roman*)}]
\item the selected branch is a strict selected-window Type II branch in the
sense of \cref{paper1:def:typeII};
\item a selected-window subsequence has a nonzero bounded selected-window limit
in the sense of \cref{paper1:def:bounded-window-limit}.  This is the
bounded-window row and is routed by
\cref{paper6:thm:bounded-selected-window-physical-routing}.
\end{enumerate}
\end{proposition}

\begin{proof}
The local concentration dichotomy supplies the positive CKN concentration
sequence and the failure of the local Type I scale-invariant bound at the
chosen physical point.  The scale-translation selection rewrites those
shrinking cylinders as fixed selected windows.  On the selected windows there
are then two mutually exclusive possibilities: either no selected-window
subsequence has a nonzero bounded selected-window limit, which is precisely
\cref{paper1:def:typeII}, or such a subsequence exists.  In the latter case the
branch is not discarded as physically impossible; it is sent to the
bounded-window physical routing theorem.
\end{proof}

\begin{proposition}[Local representation input]\label{paper1:prop:representation-input}
For a nondegenerate selected single core, the local scale-translation
construction gives represented variables \((V_n,P_n,a_n,b_n)\) satisfying
\eqref{paper1:eq:represented-ns} on compact cylinders.  Failure of the localized
nondegeneracy condition is a multibubble, cascade, or gauge-degenerate
alternative, not a compact single-core branch.
\end{proposition}
\begin{proof}
On the nondegenerate selected-core branch the scale-center construction provides
absolutely continuous scale and center functions
\(\lambda_n(t)>0\), \(x_n(t)\), with \(\lambda_n\) bounded above and below on
each selected compact window after normalization.  Define
\[
  y=\frac{x-x_n(t)}{\lambda_n(t)},\qquad
  \frac{ds}{dt}=\lambda_n(t)^{-2},
\]
and set
\[
  u(x,t)=\lambda_n(t)^{-1}V_n(y,s),
  \qquad
  p(x,t)=\lambda_n(t)^{-2}P_n(y,s).
\]
The distributional chain rule is justified by the absolute continuity of the
chart and the local suitable-energy bounds on the physical solution.  Testing
the physical Navier--Stokes equations against pullbacks of compactly supported
test fields gives, on every compact selected cylinder,
\[
  \partial_s V_n+(V_n\cdot\nabla)V_n+\nabla P_n
  =
  \nu\Delta V_n+a_n(s)(V_n+y\cdot\nabla V_n)+b_n(s)\cdot\nabla V_n,
  \qquad \nabla\cdot V_n=0,
\]
where
\[
  a_n(s)=\lambda_n(t)\dot\lambda_n(t),\qquad
  b_n(s)=\lambda_n(t)\dot x_n(t).
\]
Suitability and the local CKN quantities are invariant under this parabolic
change of variables up to the fixed selected-window constants.
\par
If the localized nondegeneracy needed to choose a single compact scale-center
chart fails, then the ordered selection has detected either separated active
cores, a same-point scale cascade, or failure of the gauge chart.  These are
precisely the multibubble, cascade, and gauge-degenerate rows, so the failure
cannot be part of the compact single-core branch.
\end{proof}

\begin{proposition}[Local pressure decomposition]
\label{paper1:prop:pressure-reconstruction}
For every \(B_r\Subset B_R\), the represented pressure decomposes as
\[
  P_n=P_{{\rm loc},n}+H_n,
  \qquad
  -\Delta P_{{\rm loc},n}=\partial_i\partial_j(\zeta V_{n,i}V_{n,j}),
\]
where \(\zeta\in C_c^\infty(B_R)\) equals one on \(B_r\), and \(H_n\) is
harmonic in \(B_r\).  The local part satisfies the Calder\'on--Zygmund bound
\cite{CalderonZygmund1952,Stein1970}
\[
  \|P_{{\rm loc},n}\|_{L^{3/2}(B_R)}
  \le C\|V_n\|_{L^3(B_R)}^2
\]
for a.e. time. The harmonic part is controlled only by the local pressure norm,
modulo spatial means; no vanishing of the harmonic part follows from
\(V_n\to0\) without an additional pressure-row hypothesis.
\end{proposition}
\begin{proof}
Fix a time for which \(V_n(\cdot,s)\in L^3(B_R)\) and
\(P_n(\cdot,s)\in L^{3/2}(B_R)\).  Let
\[
  P_{{\rm loc},n}
  =
  (-\Delta)^{-1}\partial_i\partial_j(\zeta V_{n,i}V_{n,j})
\]
in \(\mathbb R^3\), with the Newtonian potential normalized modulo constants.
Since \(\zeta V_{n,i}V_{n,j}\in L^{3/2}(\mathbb R^3)\), the Calderon--Zygmund
inequality gives
\[
  \|P_{{\rm loc},n}\|_{L^{3/2}(\mathbb R^3)}
  \le C\|V_n\|_{L^3(B_R)}^2,
\]
and hence the stated bound on \(B_R\).
\par
Taking the divergence of the represented equation and using
\(\nabla\cdot V_n=0\) gives
\[
  -\Delta P_n=\partial_i\partial_j(V_{n,i}V_{n,j})
\]
in distributions, modulo functions of time.  Since \(\zeta\equiv1\) on \(B_r\),
\[
  -\Delta(P_n-P_{{\rm loc},n})=0
  \qquad\text{in }B_r .
\]
Thus \(H_n:=P_n-P_{{\rm loc},n}\) is harmonic in \(B_r\).  The pressure is only
defined up to functions of time, so this decomposition is understood with the
same spatial-mean normalization used in the CKN pressure quantity.  The harmonic
term is therefore controlled through the local pressure oscillation, not through
the velocity alone.
\end{proof}

\begin{lemma}[Localization of pressure oscillation bounds]\label{paper1:lem:D-localization}
Let \(0<r<R\), and let \(q\in L^{3/2}(Q_R)\).  Then
\[
  D_q(r)\le C\left(\frac{R}{r}\right)^2D_q(R),
\]
where \(C\) is universal.
\end{lemma}

\begin{proof}
For a.e. \(s\in(-r^2,0)\), Jensen's inequality gives
\[
  |(q)_{B_r}(s)-(q)_{B_R}(s)|^{3/2}
  \le \frac1{|B_r|}\int_{B_r}|q(y,s)-(q)_{B_R}(s)|^{3/2}\,dy .
\]
Hence
\[
  \int_{B_r}|q-(q)_{B_r}(s)|^{3/2}\,dy
  \le
  C\int_{B_r}|q-(q)_{B_R}(s)|^{3/2}\,dy .
\]
Integrating over \((-r^2,0)\) and using \(Q_r\subset Q_R\),
\[
\begin{aligned}
  D_q(r)
  &\le
  Cr^{-2}\int_{-r^2}^0\int_{B_r}
  |q-(q)_{B_R}(s)|^{3/2}\,dy\,ds \\
  &\le
  Cr^{-2}\int_{-R^2}^0\int_{B_R}
  |q-(q)_{B_R}(s)|^{3/2}\,dy\,ds
  =
  C\left(\frac{R}{r}\right)^2D_q(R).
\end{aligned}
\]
\end{proof}

\subsection{Compactness and localized pressure compactness}

\begin{proposition}[Compactness on selected windows]\label{paper1:prop:compactness}
Let \((V_n,P_n)\) be suitable on \(Q_{R_*}\), solve \eqref{paper1:eq:represented-ns}, and satisfy
\[
  A_{V_n}(R_*)+E_{V_n}(R_*)+C_{V_n}(R_*)+D_{P_n}(R_*)\le M,
\]
with
\[
  \|a_n\|_{L^\infty(-R_*^2,0)}+\|b_n\|_{L^\infty(-R_*^2,0)}\le M.
\]
Then for every \(r<R_*\), after passing to a subsequence there are \((V,P)\)
such that
\[
  V_n\stackrel{*}{\rightharpoonup}V
  \quad\text{in }L^\infty(-r^2,0;L^2(B_r)),
\]
\[
  V_n\rightharpoonup V
  \quad\text{in }L^2(-r^2,0;H^1(B_r)),
\]
\[
  V_n\to V
  \quad\text{strongly in }L^3(Q_r),
\]
and
\[
  P_n-(P_n)_{B_r}\rightharpoonup P
  \quad\text{in }L^{3/2}(Q_r).
\]
After passing to a further subsequence, \(a_n\) and \(b_n\) converge weak-* in
\(L^\infty\) to bounded coefficients, and the limit satisfies the corresponding
represented equation distributionally on smaller cylinders.
\end{proposition}

\begin{proof}
The bounds on \(A\) and \(E\) give weak-* compactness in \(L^\infty L^2\) and
weak compactness in \(L^2H^1\) on smaller cylinders.  Fix
\(r<r_1<R_*\), and choose \(m\) so large that
\[
  H^m_0(B_{r_1})\hookrightarrow W^{1,\infty}_0(B_{r_1}).
\]
We claim that \(\partial_sV_n\) is uniformly bounded in
\[
  L^1(-r_1^2,0;H^{-m}(B_{r_1})).
\]
Let \(\varphi\in H^m_0(B_{r_1})\).  Pairing \eqref{paper1:eq:represented-ns} with
\(\varphi\), integrating by parts in space where needed, and using the
embedding above gives
\[
  |\langle (V_n\cdot\nabla)V_n,\varphi\rangle|
  =
  \left|\int_{B_{r_1}}V_{n,i}V_{n,j}\partial_j\varphi_i\,dy\right|
  \le
  C\|V_n\|_{L^3(B_{r_1})}^2\|\varphi\|_{H^m},
\]
\[
  |\langle \nu\Delta V_n,\varphi\rangle|
  \le
  C\|\nabla V_n\|_{L^2(B_{r_1})}\|\varphi\|_{H^m},
\]
\[
  |\langle \nabla P_n,\varphi\rangle|
  =
  \left|\int_{B_{r_1}}(P_n-(P_n)_{B_{r_1}})\nabla\cdot\varphi\,dy\right|
  \le
  C\|P_n-(P_n)_{B_{r_1}}\|_{L^{3/2}(B_{r_1})}\|\varphi\|_{H^m}.
\]
The modulation terms obey
\[
  |\langle a_nV_n,\varphi\rangle|
  \le
  C\|V_n\|_{L^2(B_{r_1})}\|\varphi\|_{H^m},
\]
\[
  |\langle a_n\,y\cdot\nabla V_n,\varphi\rangle|
  =
  \left|\int_{B_{r_1}}a_n V_{n,i}\partial_j(y_j\varphi_i)\,dy\right|
  \le
  C\|V_n\|_{L^2(B_{r_1})}\|\varphi\|_{H^m},
\]
and
\[
  |\langle b_n\cdot\nabla V_n,\varphi\rangle|
  =
  \left|\int_{B_{r_1}}V_{n,i}b_{n,j}\partial_j\varphi_i\,dy\right|
  \le
  C\|V_n\|_{L^2(B_{r_1})}\|\varphi\|_{H^m},
\]
where the constant uses \(r_1\) and the uniform \(L^\infty\) bounds for
\(a_n,b_n\).  Integrating these estimates in time gives a uniform bound because
\[
  \int_{-r_1^2}^0\|V_n(s)\|_{L^3(B_{r_1})}^2\,ds
  \le
  r_1^{2/3}\|V_n\|_{L^3(Q_{r_1})}^2,
\]
\[
  \int_{-r_1^2}^0\|\nabla V_n(s)\|_{L^2(B_{r_1})}\,ds
  \le
  r_1\|\nabla V_n\|_{L^2(Q_{r_1})},
\]
\[
  \int_{-r_1^2}^0
  \|P_n(s)-(P_n)_{B_{r_1}}(s)\|_{L^{3/2}(B_{r_1})}\,ds
  \le
  r_1^{2/3}
  \|P_n-(P_n)_{B_{r_1}}\|_{L^{3/2}(Q_{r_1})},
\]
and
\[
  \int_{-r_1^2}^0\|V_n(s)\|_{L^2(B_{r_1})}\,ds
  \le
  r_1^2
  \|V_n\|_{L^\infty(-r_1^2,0;L^2(B_{r_1}))}.
\]
The uniform \(C\), \(E\), \(D\), and \(A\) bounds therefore prove the claimed
time-derivative bound; the \(D\)-bound on \(Q_{r_1}\) follows from the bound on
\(Q_{R_*}\) by \cref{paper1:lem:D-localization}.

The compact embedding \(H^1(B_{r_1})\Subset L^2(B_r)\), together with
Aubin--Lions compactness \cite{Aubin1963,LionsMagenes1972,Simon1987}, gives
\[
  V_n\to V\qquad\text{strongly in }L^2(Q_r)
\]
after passing to a subsequence.  The same \(L^\infty L^2\cap L^2H^1\) bounds
give the parabolic interpolation estimate
\[
  \|V_n\|_{L^{10/3}(Q_r)}\le C(r,M).
\]
Indeed, for a.e. \(s\), the Sobolev and Gagliardo--Nirenberg inequalities give
\[
  \|V_n(s)\|_{L^{10/3}(B_r)}
  \le
  C\|V_n(s)\|_{L^2(B_r)}^{2/5}
  \|\nabla V_n(s)\|_{L^2(B_r)}^{3/5}.
\]
Raising to the power \(10/3\), integrating in time, and using the \(A\) and
\(E\) bounds yields
\[
  \int_{-r^2}^0\|V_n(s)\|_{L^{10/3}(B_r)}^{10/3}\,ds
  \le
  C
  \left(\operatorname*{ess\,sup}_{-r^2<s<0}
  \|V_n(s)\|_{L^2(B_r)}^2\right)^{2/3}
  \int_{Q_r}|\nabla V_n|^2\,dy\,ds
  \le C(r,M).
\]
The same estimate applies to \(V\) by weak lower semicontinuity, and therefore
to \(V_n-V\) by the triangle inequality.  Interpolating between \(L^2\) and
\(L^{10/3}\), one obtains
\[
  \|V_n-V\|_{L^3(Q_r)}
  \le
  \|V_n-V\|_{L^2(Q_r)}^{1/6}
  \|V_n-V\|_{L^{10/3}(Q_r)}^{5/6}.
\]
The second factor is uniformly bounded and the first tends to zero, so
\[
  V_n\to V\qquad\text{strongly in }L^3(Q_r).
\]

The pressure compactness follows from the uniform \(D\)-bound:
after subtracting spatial means,
\[
  P_n-(P_n)_{B_r}
  \rightharpoonup P
  \qquad\text{in }L^{3/2}(Q_r).
\]
We pass to the limit distributionally using the preceding convergences.  It
remains to justify the modulation terms.  For every compactly supported
test field \(\varphi\),
\[
  \int a_n(s)\, y\cdot\nabla V_n\cdot\varphi
  =
  -\int a_n(s)\,V_{n,i}\,\partial_j(y_j\varphi_i),
\]
where summation over repeated spatial indices is understood.  Equivalently,
\[
  \partial_j(y_j\varphi_i)=3\varphi_i+y_j\partial_j\varphi_i .
\]
Hence this term passes to the limit by strong local \(L^2\) convergence of
\(V_n\) and weak-* convergence of \(a_n\): if \(F_n\to F\) in \(L^1\) and
\(a_n\rightharpoonup^* a\) in \(L^\infty\), then
\[
  \int a_nF_n-\int aF
  =
  \int a_n(F_n-F)+\int (a_n-a)F\to0.
\]
The term \(a_nV_n\) is immediate from the same argument, while
\[
  \int b_n(s)\cdot\nabla V_n\cdot\varphi
  =
  -\int V_{n,i}\,b_{n,j}(s)\,\partial_j\varphi_i
\]
passes to the limit by strong local \(L^2\) convergence of \(V_n\) and weak-*
convergence of \(b_n\).  These passages give the limiting represented equation
in the sense of distributions.
\end{proof}

\begin{remark}[Use of bounded modulation]\label{paper1:rem:bounded-modulation}
The uniform \(L^\infty\) bounds on \(a_n\) and \(b_n\) enter the compactness
argument in exactly two places.  First, they give the negative Sobolev
time-derivative bound for the three lower-order modulation terms in
\eqref{paper1:eq:represented-ns}.  Second, after subsequential weak-* convergence of
\(a_n,b_n\), they allow the modulation terms to pass to the distributional
limit once \(V_n\to V\) strongly in \(L^2_{\rm loc}\).
\end{remark}

\begin{lemma}[Stability of the localized pressure part]\label{paper1:lem:pressure-stability}
Let \(V_n\to V\) strongly in \(L^3(Q_R)\), and use the pressure decomposition
of \cref{paper1:prop:pressure-reconstruction}.  Then for every \(r<R\),
\[
  P_{{\rm loc},n}\to P_{\rm loc}
  \quad\text{strongly in }L^{3/2}(Q_r),
\]
where \(-\Delta P_{\rm loc}=\partial_i\partial_j(\zeta V_iV_j)\).  If
\(V\equiv0\) on \(Q_R\), then
\[
  P_{{\rm loc},n}\to0
  \quad\text{strongly in }L^{3/2}(Q_r).
\]
This lemma does not assert convergence of the full mean-normalized pressure:
the harmonic component may contain local parasitic pressure modes independent
of \(V_n\).
\end{lemma}

\begin{proof}
Strong convergence in \(L^3\) gives
\[
  V_{n,i}V_{n,j}\to V_iV_j
  \quad\text{strongly in }L^{3/2}(Q_R).
\]
Indeed,
\[
  \|V_{n,i}V_{n,j}-V_iV_j\|_{L^{3/2}(Q_R)}
  \le
  \bigl(\|V_n\|_{L^3(Q_R)}+\|V\|_{L^3(Q_R)}\bigr)
  \|V_n-V\|_{L^3(Q_R)}.
\]
Calder\'on--Zygmund estimates for the compactly supported localized source yield
strong convergence of the localized pressures in \(L^{3/2}\) on smaller balls:
\[
  \|P_{{\rm loc},n}-P_{\rm loc}\|_{L^{3/2}(Q_r)}
  \le
  C_{r,R}
  \|V_{n,i}V_{n,j}-V_iV_j\|_{L^{3/2}(Q_R)}.
\]

If \(V\equiv0\), then \(P_{\rm loc}=0\), and the first part gives the displayed
convergence of \(P_{{\rm loc},n}\). The harmonic part \(H_n\) is not controlled
by Calder\'on--Zygmund from the velocity source and may persist even when
\(V_n\to0\); it is therefore handled only in the pressure/exterior rows, not in
the compact single-core zero-velocity contradiction.
\end{proof}

\begin{lemma}[Zero-dissipation limit]\label{paper1:lem:zero-dissipation}
If
\[
  V_n\rightharpoonup V\quad\text{in }L^2(I;H^1(B_\rho))
\]
and
\[
  \int_I\int_{B_\rho}|\nabla V_n|^2\,dy\,ds\to0,
\]
then \(\nabla V=0\) a.e. on \(B_\rho\times I\).  Hence for a.e. time,
\(V(\cdot,s)\) is spatially constant on \(B_\rho\).
\end{lemma}

\begin{proof}
Weak lower semicontinuity gives
\[
  \int_I\int_{B_\rho}|\nabla V|^2\,dy\,ds
  \le \liminf_{n\to\infty}\int_I\int_{B_\rho}|\nabla V_n|^2\,dy\,ds=0.
\]
Thus \(\nabla V=0\) a.e. on \(B_\rho\times I\).  By Fubini, for a.e.
\(s\in I\) one has \(\nabla V(\cdot,s)=0\) in \(L^2(B_\rho)\).  Since
\(B_\rho\) is connected, every \(H^1(B_\rho)\) function with zero weak
gradient is a.e. equal to a constant.  Hence \(V(\cdot,s)\) is spatially
constant for a.e. \(s\).
\end{proof}

\begin{lemma}[Spatially constant limits]\label{paper1:lem:constant-limit}
Let \(V\) be a limit on \(Q_\rho\) and suppose \(\nabla V=0\) a.e. there.
Then \(V(y,s)=c(s)\) on \(Q_\rho\), with
\(c\in L^\infty(-\rho^2,0)\).  If \(c\not\equiv0\) on \(Q_\rho\), then
the approximating selected-window subsequence has a nonzero bounded
selected-window limit.
\end{lemma}

\begin{proof}
The spatial constancy follows from \cref{paper1:lem:zero-dissipation}.  Since
\(V(\cdot,s)=c(s)\) on \(B_\rho\), the local \(L^\infty_sL^2_y\) bound gives
\[
  |c(s)|^2|B_\rho|\le \operatorname*{ess\,sup}_{-\rho^2<s<0}
  \int_{B_\rho}|V(y,s)|^2\,dy,
\]
so \(c\in L^\infty(-\rho^2,0)\).  If \(c\not\equiv0\) on \(Q_\rho\), then the
strong \(L^3\) convergence of \cref{paper1:prop:compactness} gives a subsequential
limit to a nonzero bounded function on \(Q_\rho\), which is precisely a nonzero
bounded selected-window limit in the sense of \cref{paper1:def:bounded-window-limit}.
\end{proof}

\subsection{Contradiction lemmas}

\begin{lemma}[Persistence of retained velocity core]\label{paper1:lem:persistence}
For a compact single-core branch in the sense of \cref{paper1:def:single-core-branch},
there are fixed \(r_*>0\) and \(\eta_0>0\) such that along the selected-window
subsequence,
\[
  C_{V_n}(r_*)\ge\eta_0
  \qquad\text{for all }n.
\]
\end{lemma}

\begin{proof}
The retained velocity lower bound is part of the definition of the selected
compact core and is preserved by passing to subsequences of the selected
windows.
\end{proof}

\begin{lemma}[Zero constant contradicts retained velocity]\label{paper1:lem:zero-contradiction}
Let \(r_*<\rho\).  Suppose the compactness limit in \cref{paper1:prop:compactness} is
spatially constant on \(Q_\rho\) and is identically zero on \(Q_{r_*}\).  Then
\[
  C_{V_n}(r_*)\to0.
\]
Consequently this case contradicts \cref{paper1:lem:persistence}.
\end{lemma}

\begin{proof}
Strong \(L^3(Q_{r_*})\) convergence to zero gives
\[
  C_{V_n}(r_*)
  =
  r_*^{-2}\int_{Q_{r_*}}|V_n|^3\,dy\,ds
  =
  r_*^{-2}\|V_n\|_{L^3(Q_{r_*})}^3
  \to0.
\]
This contradicts the retained velocity lower bound
\[
  C_{V_n}(r_*)\ge\eta_0.
\]
No assertion is made about \(D_{P_n}(r_*)\): a pressure-only harmonic mode may
remain and is not part of the retained compact single-core velocity row.
\end{proof}

\begin{lemma}[Nonzero constant exits the strict selected-window branch]\label{paper1:lem:nonzero-contradiction}
If the compactness limit is spatially constant and nonzero on a compact
subcylinder, then the original sequence leaves the strict selected-window
branch.
\end{lemma}

\begin{proof}
By \cref{paper1:lem:constant-limit}, the selected-window subsequence has a nonzero
bounded selected-window limit.  By \cref{paper1:def:typeII}, such a sequence is
assigned to the bounded-window alternative rather than to the strict
selected-window branch.  Physical local Type II germs are then routed by
\cref{paper6:thm:bounded-selected-window-physical-routing}.
\end{proof}

\subsection{Main theorem}

\begin{theorem}[Local single-core Type II criterion]\label{paper1:thm:single-core-criterion}
Under the analytic inputs
\cref{paper1:prop:paper0-typeII-reduction,paper1:prop:representation-input,paper1:prop:pressure-reconstruction},
a compact single-core vanishing-dissipation branch in the sense of
\cref{paper1:def:single-core-branch} cannot be a local Type II sequence.
\end{theorem}

\begin{proof}
By the definition of vanishing localized dissipation, after passing to a
selected-window subsequence there is a fixed radius \(\rho\in(r_*,R_*)\) such
that
\[
  \int_{Q_\rho}|\nabla V_n|^2\,dy\,ds\to0
\]
and the positive core bound on \(Q_{r_*}\) remains valid along this subsequence.
Apply compactness on the fixed compact
cylinder \(Q_\rho\) and pass to a subsequence.  The compactness proposition
gives
\[
  V_n\rightharpoonup V\quad\text{in }L^2(-\rho^2,0;H^1(B_\rho))
\]
and strong convergence in \(L^3\) on every smaller cylinder.  Since the
localized dissipation tends to zero, the zero-dissipation lemma gives
\[
  \nabla V=0\quad\text{a.e. on }Q_\rho.
\]
Thus \(V(y,s)=c(s)\) on \(Q_\rho\), and in particular on the selected core
cylinder \(Q_{r_*}\).

If this spatially constant limit is zero on the core cylinder, the zero-limit
contradiction lemma gives
\[
  C_{V_n}(r_*)\to0,
\]
contradicting persistence of the retained velocity core.

If the spatially constant limit is not identically zero on \(Q_{r_*}\), the
nonzero-limit lemma produces a nonzero bounded selected-window limit, so the
strict selected-window branch stops at the bounded-window alternative.  In the
physical assembly that bounded-window event is routed by
\cref{paper6:thm:bounded-selected-window-physical-routing}.  The two cases exhaust all
spatially constant limits, so the branch cannot occur.
\end{proof}

\begin{remark}[Minimal local inputs]
The proof uses two local analytic inputs from the Type II argument: the
localized representation for a nondegenerate selected core, the compact-ball
pressure reconstruction with stability on smaller balls.  The
compact-cylinder \(A,E,C,D\) bounds and vanishing localized dissipation are
part of the branch hypothesis, not separate analytic inputs.  Multibubble,
cascade, gauge-degenerate, and compactness-failure alternatives are outside
this single-core theorem; they are treated by the local multibubble and
cascade exclusion theorem.
\end{remark}

\clearpage
\section{Repaired-Gauge Construction}\label{sec:repaired-gauge-construction}
\subsection{Setup and notation}\label{paper2:sec:setup-notation}

Let \(u:\RR^{3}\times (T-\delta,T)\to \RR^{3}\), \(p:\RR^{3}\times (T-\delta,T)\to\RR\)
be a suitable weak Navier--Stokes solution with viscosity \(\nu>0\).
For \(z=(x_0,t_0)\) and \(R>0\),
\[
Q_R(z):=B_R(x_0)\times (t_0-R^2,t_0),\qquad
Q_R^*:=B_R\times I_*.
\]
Fix a reference physical center-time \((x_\bullet,t_\star)\) around which \(B_{R_*}(x_\bullet)\) is taken.
For a fixed spatial center \(x_0\) we also use
\[
Q_R(t):=Q_R(x_0,t).
\]
Fix a compact window \(I_*:=[\tau_0-\ell,\tau_0+\ell]\), \(\ell>0\), and write
\(I_{\rm ph}:=[t_\star-\Lambda_{\max}^2\ell,\; t_\star+\Lambda_{\max}^2\ell]\).

For a cylinder \(Q_R(x_0,t_0)=B_R(x_0)\times (t_0-R^2,t_0)\), set
\[
\begin{aligned}
A(u;Q_R(x_0,t_0))&:=R^{-1}\sup_{s\in(t_0-R^2,t_0)}
\int_{B_R(x_0)}|u(x,s)|^2\,dx,\\
E(u;Q_R(x_0,t_0))&:=R^{-1}\int_{Q_R(x_0,t_0)}|\nabla u|^2\,dx\,ds,\\
C(u;Q_R(x_0,t_0))&:=R^{-2}\int_{Q_R(x_0,t_0)}|u|^3\,dx\,ds,\\
D(p;Q_R(x_0,t_0))&:=R^{-2}\int_{t_0-R^2}^{t_0}\int_{B_R(x_0)}
\left|p(x,s)-\langle p(s)\rangle_{B_R(x_0)}\right|^{3/2}\,dx\,ds.
\end{aligned}
\]
Here \(\langle p(s)\rangle_{B_R(x_0)}:=|B_R|^{-1}\int_{B_R(x_0)}p(x,s)\,dx\);
pressure is always understood modulo time-dependent constants.

\begin{definition}[Single-core hypotheses]\label{paper2:def:SC}
Let \((R_*,r_*,M,\eta,\kappa,\Lambda_{\min},\Lambda_{\max})\) satisfy
\(\Lambda_{\min}\le \Lambda_{\max}<\infty\).
We say \((u,p)\) satisfies \(\mathrm{SC}(R_*,r_*,M,\eta,\kappa)\) on \(Q_{R_*}(x_\bullet,t_\star)\)
if the following hold:
\begin{enumerate}
\item[(\textup{SC1})] \(A(u;Q_{R_*}(x_\bullet,t)) + E(u;Q_{R_*}(x_\bullet,t))
  + C(u;Q_{R_*}(x_\bullet,t)) + D(p;Q_{R_*}(x_\bullet,t))\le M\) for every
  \(t\in[t_\star-R_*^2,t_\star)\).
\item[(\textup{SC2})] There exists a physically selected core map
  \((x_{\rm c}(t),\rho_{\rm c}(t))\) on \(I_{\rm ph}\) such that
  \(\rho_{\rm c}(t)\in(0,R_*]\), \(x_{\rm c}(t)\in B_{R_*/2}(x_\bullet)\), and for some
  \(\Theta_0>0\),
  \[
  \int_{B_{r_*}(0)}|\rho_{\rm c}(t)u(x_{\rm c}(t)+\rho_{\rm c}(t)y,t)|^3\dd y\ge \eta
  \quad \text{for a.e. }t\in I_{\rm ph}.
  \]
\item[(\textup{SC3})] There exists \(0<r_{\rm sep}\le r_*/2\) such that for every
  \(0<r\le r_{\rm sep}\), every competitor pair
  \[
  \mathcal C_{\rm core}(t):=\left\{(\bar x,\bar\rho):|\bar x-x_\bullet|\le \frac{R_*}{2},\
  0<\bar\rho\le R_*\right\},
  \]
  \((\bar x,\bar\rho)\in \mathcal C_{\rm core}(t)\setminus\{(x_{\rm c}(t),\rho_{\rm c}(t))\}\)
  with
  \[
  \rho_{\rm c}(t)^{-1}|\bar x-x_{\rm c}(t)|
  +\left|\frac{\bar\rho}{\rho_{\rm c}(t)}-1\right|\ge r_{\rm sep}
  \]
  has separated local weighted mass in the same core-scale variables:
  \[
  \left|\int_{B_r}|\rho_{\rm c}(t)u(x_{\rm c}(t)+\rho_{\rm c}(t)y,t)|^3\,\dd y
  - \int_{B_r}|\bar\rho\,u(\bar x+\bar\rho y,t)|^3\,\dd y\right|
  \ge \eta.
  \]
  This quantitative non-competing core condition is imposed on the
  selected physical core directly; it is not a consequence of the repaired-gauge
  construction.
\item[(\textup{SC4})] For the selected core,
  \[
  \sup_{t\in I_{\rm ph}}\int_{B_{2R_*}}|y|^{-p}|\rho_{\rm c}(t)u(x_{\rm c}(t)+\rho_{\rm c}(t) y,t)|^3\,dy\le M,
  \]
  for some \(0<p<3\).
\item[(\textup{SC5})] The preliminary frame associated with the selected core is
  the absolutely continuous representative
  \[
  0<\Lambda_{\min}\le \Lambda_{\rm pre}(\tau)\le \Lambda_{\max},\quad
  \Lambda_{\rm pre}\in AC(I_*),\quad
  X_{\rm pre}\in AC(I_*;\RR^3).
  \]
  It is tied to the selected physical core by the time change
  \(dt_{\rm c}/d\tau=\Lambda_{\rm pre}(\tau)^2\), \(t_{\rm c}(\tau_0)=t_\star\):
  one has
  \[
  X_{\rm pre}(\tau)=x_{\rm c}(t_{\rm c}(\tau)),\qquad
  \Lambda_{\rm pre}(\tau)=\rho_{\rm c}(t_{\rm c}(\tau))
  \]
  for a.e. \(\tau\in I_*\).
\end{enumerate}
\end{definition}
\noindent
The parameter \(\kappa\) denotes the quantitative constants obtained in the local gauge estimates and is not used in items \textup{(SC1)--(SC5)}.

\begin{remark}
The selected core \((x_{\rm c},\rho_{\rm c})\) and its separation/nondegeneracy
conditions are specified in physical variables; no repaired-gauge constraint is used in
their definition.  This avoids the circularity of defining the core via the chart
to be constructed.
\end{remark}

\subsection{Preliminary frame map from the selected core}\label{paper2:sec:preliminary-frame}

The transport map is constructed from the selected core; the core selection
itself is part of the hypotheses.

Define
\[
\frac{dt_{\rm c}}{d\tau}=\Lambda_{\rm pre}(\tau)^2,\qquad t_{\rm c}(\tau_0)=t_\star,
\]
and use the compatibility in \textup{(SC5)} to identify
\(\Lambda_{\rm pre}(\tau)=\rho_{\rm c}(t_{\rm c}(\tau))\) and
\(X_{\rm pre}(\tau)=x_{\rm c}(t_{\rm c}(\tau))\) a.e.

\begin{proposition}[Preliminary frame map regularity]\label{paper2:prop:prelim-frame}
Under \(\mathrm{SC}(R_*,r_*,M,\eta,\kappa)\), the map
\[
\Phi_{\rm pre}(\tau,Y):=(x,t)=(X_{\rm pre}(\tau)+\Lambda_{\rm pre}(\tau)Y,\ t_{\rm c}(\tau))
\]
is well-defined and absolutely continuous in \(\tau\).
Moreover, \(1/\Lambda_{\rm pre}\), \(\Lambda_{\rm pre}\in L^\infty(I_*)\),
\(X_{\rm pre}\in W^{1,1}_{\rm loc}(I_*;\RR^3)\), and for every \(R\le R_*\),
\[
\Phi_{\rm pre}(Q_R^*)\subset B_{R_*/2+\Lambda_{\max}R}(x_\bullet)\times I_{\rm ph}.
\]
For each fixed time \(\tau\), the spatial map
\(Y\mapsto X_{\rm pre}(\tau)+\Lambda_{\rm pre}(\tau)Y\) is bi-Lipschitz on
\(B_R\), with \(\det \nabla_Y\Phi_{\rm pre}=\Lambda_{\rm pre}^3\), hence
Jacobian and inverse-Jacobian bounds depend only on \(\Lambda_{\min},\Lambda_{\max}\).
\end{proposition}
\begin{proof}
From \textup{(SC5)}, \(\Lambda_{\rm pre}\in AC(I_*)\),
\(0<\Lambda_{\min}\le \Lambda_{\rm pre}\le\Lambda_{\max}<\infty\), and
\(X_{\rm pre}\in AC(I_*;\mathbb R^3)\). The scalar ODE
\[
\frac{dt_{\rm c}}{d\tau}=\Lambda_{\rm pre}(\tau)^2,\qquad t_{\rm c}(\tau_0)=t_\star
\]
therefore has an absolutely continuous strictly increasing solution on \(I_*\), and
\[
\frac{dt_{\rm c}}{d\tau}=\Lambda_{\rm pre}(\tau)^2
\quad\text{a.e. on }I_*.
\]
The compatibility condition \textup{(SC5)} identifies this absolutely continuous frame
with the selected physical core a.e.; no differentiability of \(x_{\rm c}\) or
\(\rho_{\rm c}\) beyond the frame regularity in \textup{(SC5)} is used.
For fixed \(\tau\), \(\nabla_Y\Phi_{\rm pre}(\tau,\cdot)=\Lambda_{\rm pre}(\tau)I\), so
\[
\det\nabla_Y\Phi_{\rm pre}(\tau,\cdot)=\Lambda_{\rm pre}(\tau)^3,
\qquad
\Lambda_{\min}^3\le \det\nabla_Y\Phi_{\rm pre}\le \Lambda_{\max}^3.
\]
Hence, if \(Y_1,Y_2\in B_R\),
\[
\Lambda_{\min}|Y_1-Y_2|
\le |\Phi_{\rm pre}(\tau,Y_1)-\Phi_{\rm pre}(\tau,Y_2)|
\le \Lambda_{\max}|Y_1-Y_2|,
\]
and therefore \(\Phi_{\rm pre}(\tau,\cdot)\) is bi-Lipschitz on \(B_R\). Its inverse has
Lipschitz constant \(\Lambda_{\min}^{-1}\), so Jacobian and inverse-Jacobian bounds
are controlled by \(\Lambda_{\min},\Lambda_{\max}\).

For image inclusions, if \(x=X_{\rm pre}(\tau)+\Lambda_{\rm pre}(\tau)Y\) with
\(|Y|\le R\), then
\[
|x-X_{\rm pre}(\tau)|\le \Lambda_{\max}R,\qquad
X_{\rm pre}(\tau)\in B_{R_*/2}(x_\bullet),
\]
so \(x\in B_{R_*/2+\Lambda_{\max}R}(x_\bullet)\). As \(\tau\in I_*\), we also have
\(t=t_{\rm c}(\tau)\in I_{\rm ph}\). Therefore
\[
\Phi_{\rm pre}(Q_R^*)\subset B_{R_*/2+\Lambda_{\max}R}(x_\bullet)\times I_{\rm ph}.
\]
\end{proof}
\begin{definition}[Preliminary profile]\label{paper2:def:prelim-profile}
For \((x,t)=\Phi_{\rm pre}(\tau,Y)\), define
\[
V_{\rm pre}(Y,\tau):=\Lambda_{\rm pre}(\tau)\,u(x,t),\qquad
P_{\rm pre}(Y,\tau):=\Lambda_{\rm pre}(\tau)^2\,p(x,t).
\]
Hence \(u(x,t)=\Lambda_{\rm pre}(\tau)^{-1}V_{\rm pre}(Y,\tau)\) and
\(p(x,t)=\Lambda_{\rm pre}(\tau)^{-2}P_{\rm pre}(Y,\tau)\) under \(\Phi_{\rm pre}\).
\end{definition}

\subsection{Exact renormalized change of variables}\label{paper2:sec:change}

Define renormalized fields on \(Q_{2R}^*\) by
\[
Y=\frac{x-X(\tau)}{\Lambda(\tau)},\qquad
u(x,t)=\Lambda(\tau)^{-1}V(Y,\tau),\qquad
p(x,t)=\Lambda(\tau)^{-2}P(Y,\tau),
\]
where \(dt/d\tau=\Lambda(\tau)^2\), \(X:\,I_*\to\RR^3\), \(\Lambda:\,I_*\to(0,\infty)\),
and \(\Lambda,X\in AC(I_*)\).
Write
\[
a(\tau):=\frac{\Lambda'(\tau)}{\Lambda(\tau)},\qquad
b(\tau):=\frac{X'(\tau)}{\Lambda(\tau)}.
\]

\begin{lemma}[Exact differential identities]\label{paper2:lem:change-id}
For smooth data and fixed \((x,t)\), the chain rule gives
\[
\partial_t u
=\Lambda^{-3}\left(\partial_\tau V-(b+aY)\cdot\nabla_Y V-aV\right),
\]
\[
\nn_x u=\Lambda^{-2}\nn_Y V,\qquad \Delta_x u=\Lambda^{-3}\Delta_Y V,\qquad
\nn_x p=\Lambda^{-3}\nn_Y P.
\]
In particular
\[
\nn_x\cdot u=\Lambda^{-2}\nn_Y\cdot V.
\]
\end{lemma}
\begin{proof}
For fixed \((x,t)\), \(\tau=\tau(t)\) and \(Y=(x-X(\tau))/\Lambda(\tau)\).
Differentiate
\[
u(x,t)=\Lambda(\tau)^{-1}V\bigl(Y,\tau\bigr)
\]
in \(t\):
\[
\partial_t u
=-\frac{\Lambda'}{\Lambda^2}V
+\frac1\Lambda\partial_\tau V\,\frac{d\tau}{dt}
+\frac1\Lambda\nabla_YV\cdot\frac{dY}{dt}.
\]
Since \(d\tau/dt=\Lambda^{-2}\) and
\[
Y_\tau=-\Lambda^{-1}X'-\frac{\Lambda'}{\Lambda}Y,\qquad
\frac{dY}{dt}=Y_\tau\,\frac{d\tau}{dt},
\]
therefore
\[
\partial_tu
=\Lambda^{-3}\left(\partial_\tau V-
\left(\frac{X'}{\Lambda}+\frac{\Lambda'}{\Lambda}Y\right)\cdot\nabla_YV
-\frac{\Lambda'}{\Lambda}V\right)
=\Lambda^{-3}\left(\partial_\tau V-(b+aY)\cdot\nabla_YV-aV\right).
\]
For spatial derivatives, \(Y=(x-X)/\Lambda\) gives \(\nabla_xY=\Lambda^{-1}I\), hence
\[
\nabla_xu=\Lambda^{-2}\nabla_YV,\qquad
\Delta_xu=\Lambda^{-3}\Delta_YV.
\]
Similarly \(p(x,t)=\Lambda^{-2}P(Y,\tau)\) yields
\[
\nabla_x p=\Lambda^{-3}\nabla_YP.
\]
Taking divergence of \(u=\Lambda^{-1}V\) gives
\[
\nabla_x\cdot u=\Lambda^{-2}\nabla_Y\cdot V.
\]
\end{proof}

\begin{theorem}[Renormalized Navier--Stokes equation]\label{paper2:thm:ren-NS}
Assume \(u,p\) solve
\[
\partial_t u+(u\cdot\nn_x)u+\nn_x p-\nu\Delta_xu=0,\qquad \nn_x\cdot u=0
\]
in distributions. Under the transformation above, \(V,P\) satisfy, on the image
renormalized cylinder, the exact equation
\[
\partial_\tau V+(V\cdot\nn)V+\nn P
=\nu\Delta V
+a(\tau)\bigl(V+Y\cdot\nn V\bigr)+b(\tau)\cdot\nn V,\qquad \nn\cdot V=0,
\]
with the convention fixed by
\[
a=\frac{\Lambda'}{\Lambda},\qquad b=\frac{X'}{\Lambda}.
\]
Conversely, any \(V,P\) satisfying this equation pull back to a distributional
solution \(u,p\) of Navier--Stokes on the physical image.
\end{theorem}
\begin{proof}
To avoid hidden regularity loss, first write the physical equation tested against
\(\Phi\in C_c^\infty\) in the physical cylinder and then use the smooth pull-back formula
(an approximation argument gives distributions). By Lemma~\ref{paper2:lem:change-id},
\[
\partial_tu+\bigl(u\cdot\nabla_x\bigr)u+\nabla_xp-\nu\Delta_xu=0
\]
becomes
\[
\Lambda^{-3}\!\left(\partial_\tau V-(b+aY)\cdot\nabla V-aV\right)
+\Lambda^{-3}(V\cdot\nabla)V
+\Lambda^{-3}\nabla P
-\Lambda^{-3}\nu\Delta V=0.
\]
Multiplying by \(\Lambda^3\) gives
\[
\partial_\tau V+(V\cdot\nabla)V+\nabla P
=
\nu\Delta V
+a\bigl(V+Y\cdot\nabla V\bigr)+b\cdot\nabla V.
\]
The divergence condition is
\[
\nabla_x\cdot u=\Lambda^{-2}\nabla\cdot V=0
\quad\Longrightarrow\quad
\nabla\cdot V=0.
\]
Conversely, assume \((V,P)\) solve this equation on the renormalized cylinder.
Set \(u,p\) by the same change-of-variables map.
Then the preceding identities run in reverse and give
 \[
\partial_t u+(u\cdot\nabla_x)u+\nabla_xp-\nu\Delta_xu=0,\qquad
\nabla_x\cdot u=0
\]
in \(\mathcal D'\) on the physical image.
\end{proof}

\subsection{Domain localization and support}\label{paper2:sec:domains}

Let \(R>0\) be fixed and \(0<\ell\le \ell_0\). Set \(Q_R^*=B_R\times I_*\).
With \(\Phi\) as above define physical support maps
\[
\mathcal Q^{\rm ph}_{R}:=\Phi(Q_R^*),\qquad
\mathcal Q^{\rm ph}_{2R}:=\Phi(Q_{2R}^*).
\]

\begin{proposition}[Compact support transport]\label{paper2:prop:domain}
Assume \(\Lambda_{\min}\le \Lambda\le \Lambda_{\max}\) on \(I_*\).
Then
\[
\begin{gathered}
X(\tau)\in B_{R_*/2}(x_\bullet),\qquad |x-X(\tau)|\le \Lambda_{\max}R,\\
\mathcal Q^{\rm ph}_{R}\subset B_{R_*/2+ \Lambda_{\max}R}(x_\bullet)\times I_{\rm ph},\\
\mathcal Q^{\rm ph}_{2R}\subset B_{R_*/2+2\Lambda_{\max}R}(x_\bullet)\times I_{\rm ph},
\end{gathered}
\]
for every \(\tau\in I_*\) and \(x=X(\tau)+\Lambda(\tau)Y\) with \(Y\in B_{2R}\).
In particular, if \(|c(\tau)|\le \delta\) from Proposition~\ref{paper2:prop:preliminary-correction},
then \(X(\tau)\in B_{R_*/2+\Lambda_{\max}\delta}(x_\bullet)\) on \(I_*\).
Hence, \(\chi_R\in C_c^\infty(B_{2R})\) pulls back to a compactly supported
physical cutoff that is smooth in the spatial variables and absolutely
continuous in time under the AC chart, with support depending only on
\(R,\Lambda_{\min},\Lambda_{\max}\).
\end{proposition}
\begin{proof}
Fix \(\tau\in I_*\) and \(Y\in B_R\). Since \(x=X(\tau)+\Lambda(\tau)Y\),
\[
|x-X(\tau)|\le |Y|\,\Lambda_{\max}\le \Lambda_{\max}R,\qquad
\Lambda_{\max}^{-1}|x-X(\tau)|\le |Y|.
\]
Because \(X(\tau)\in B_{R_*/2}(x_\bullet)\), this gives
\[
x\in B_{R_*/2+\Lambda_{\max}R}(x_\bullet).
\]
Similarly, if \(Y\in B_{2R}\), then \(x\in B_{R_*/2+2\Lambda_{\max}R}(x_\bullet)\).
Together with \(t=t_{\rm c}(\tau)\in I_{\rm ph}\), these give the claimed
inclusions for \(\mathcal Q_{R}^{\rm ph}\) and \(\mathcal Q_{2R}^{\rm ph}\).

For pull-backs of cutoffs, \(\Phi\) is bi-Lipschitz at every \(\tau\), so composition
\[
\chi_R^{\rm ph}(x,t):=\chi_R\!\left(\frac{x-X(\tau(t))}{\Lambda(\tau(t))}\right)
\]
is smooth in \(x\) for each time and is AC in time when the chart is AC. It is
supported where \(Y\in B_{2R}\), hence supported where
\((x,t)\in\mathcal Q^{\rm ph}_{2R}\). The support size depends
only on \(R,\Lambda_{\min},\Lambda_{\max}\) and the fixed \(C_c^\infty\) norm of
\(\chi_R\), not on the particular branch.
\end{proof}

\subsection{Transport of local quantities under the chart}\label{paper2:sec:local-transfer}

Let \(t=t(\tau)=t_{\rm c}(\tau)\) and, for \(R>0\), define
\[
J_\tau^{-}(R):=\bigl(t(\tau)-\Lambda_{\max}^{2}R^2,\ t(\tau)\bigr].
\]
Define renormalized local quantities on \(Q_R^*(\tau)\) by
\[
\begin{aligned}
A_V(\tau,R)&:=\sup_{s\in(\tau-R^2,\tau)}\int_{B_R}|V(\cdot,s)|^2\,dY,\qquad
E_V(\tau,R):=\int_{\tau-R^2}^\tau\!\!\int_{B_R}|\nabla V|^2\,dY\,d s,\\
C_V(\tau,R)&:=\int_{\tau-R^2}^\tau\!\!\int_{B_R}|V|^3\,dY\,d s,\qquad
D_V(\tau,R):=\int_{\tau-R^2}^\tau\!\!\int_{B_R}\Bigl|P-\langle P\rangle_{B_R}\Bigr|^{3/2}\,dY\,d s.
\end{aligned}
\]
For interval notation and a center \(X\), define
\[
A_u(J,R;X):=\sup_{s\in J}\int_{B_R(X)}|u(\cdot,s)|^2\,dx,\quad
E_u(J,R;X):=\int_J\int_{B_R(X)}|\nabla u|^2\,dx\,ds,
\]
\[
C_u(J,R;X):=\int_J\int_{B_R(X)}|u|^3\,dx\,ds,\quad
D_u(J,R;X):=\int_J\int_{B_R(X)}|p-\langle p\rangle_{B_R(X)}|^{3/2}\,dx\,ds.
\]
In the inequalities below we abbreviate
\[
\begin{aligned}
A_u(J,R)&:=A_u(J,R;X(\tau)),& E_u(J,R)&:=E_u(J,R;X(\tau)),\\
C_u(J,R)&:=C_u(J,R;X(\tau)),& D_u(J,R)&:=D_u(J,R;X(\tau)).
\end{aligned}
\]
and set
\[
\Delta_X(\tau,R):=\sup_{s\in(\tau-R^2,\tau)}|X(s)-X(\tau)|,
\qquad
R_{\rm ph}(\tau,R):=\Lambda_{\max}R+\Delta_X(\tau,R).
\]

\begin{proposition}[Cylinder comparison on compact windows]\label{paper2:prop:local-transfer}
Assume \(\Lambda_{\min}\le \Lambda\le \Lambda_{\max}\) and formulas of
Lemma~\ref{paper2:lem:change-id}.
Then for every \(\tau\in I_*\), \(R>0\),
\[
A_V(\tau,R)\le \Lambda_{\min}^{-1}A_u(J_\tau^{-}(R),R_{\rm ph}(\tau,R)),
\]
\[
E_V(\tau,R)\le \Lambda_{\min}^{-1}E_u(J_\tau^{-}(R),R_{\rm ph}(\tau,R)),
\]
\[
C_V(\tau,R)\le \Lambda_{\min}^{-2}C_u(J_\tau^{-}(R),R_{\rm ph}(\tau,R)),
\]
\[
D_V(\tau,R)\le \Lambda_{\min}^{-2}D_u(J_\tau^{-}(R),R_{\rm ph}(\tau,R)).
\]
In particular, all four renormalized quantities are bounded whenever the physical
quantities are bounded on the corresponding compact physical cylinders.
\end{proposition}

\begin{proof}
These are direct changes of variables under \(\Phi_{\rm pre}\):
\[
V(\cdot,\tau)=\Lambda(\tau)u(\cdot,t(\tau)),\qquad
P(\cdot,\tau)=\Lambda(\tau)^2 p(\cdot,t(\tau)),
\]
\[
\nabla_YV=\Lambda^2\nabla_xu,\qquad dY\,d\tau=\Lambda^{-5}\,dx\,dt,\qquad
d\tau=\Lambda^{-2}\,dt.
\]
Insert these formulas into each definition and use
\(\Lambda_{\min}\le\Lambda\le\Lambda_{\max}\), plus
\[
\tau-s\in(0,R^2)\implies t(\tau)-\Lambda_{\max}^{2}R^2\le t(s)\le t(\tau).
\]
Moreover, if \(s\in(\tau-R^2,\tau)\) and \(x=X(s)+\Lambda(s)Y\) with \(|Y|<R\), then
\[
|x-X(\tau)|\le |X(s)-X(\tau)|+\Lambda(s)|Y|
\le \Delta_X(\tau,R)+\Lambda_{\max}R
=R_{\rm ph}(\tau,R).
\]
So the physical image of \(B_R\times(\tau-R^2,\tau)\) is contained in
\(B_{R_{\rm ph}(\tau,R)}(X(\tau))\times J_\tau^{-}(R)\), which yields the
stated comparisons.
\end{proof}

\subsection{Local repaired gauge map}\label{paper2:sec:gauge-map}

Fix \(R>0\), \(0<p<3\), and choose a radial cutoff
\(\chi_R\in C_c^\infty(B_{2R})\) with \(0\le\chi_R\le1\),
\(\chi_R\equiv1\) on \(B_R\), and \(Y\cdot\nabla\chi_R\le0\). Let the local
energy profile space be
\[
\mathcal V_R:=\left\{V\in L^3(B_{2R})\cap W^{1,3/2}(B_{2R};\RR^3):
|Y|^{-p/3}V\in L^3(B_{2R})\right\},
\]
with norm
\[
\|V\|_{\mathcal V_R}:=\|V\|_{L^3(B_{2R})}
+\||Y|^{-p/3}V\|_{L^3(B_{2R})}
+\|\nn V\|_{L^{3/2}(B_{2R})}.
\]
For differentiating the repaired scale and centering gauges we use the
admissible gauge class
\[
\mathcal A_{p,R}:=\left\{V:
\begin{array}{l}
V\in L^3(B_{4R})\cap W^{1,3}_{\rm loc}(B_{4R};\RR^3),\\[2mm]
\displaystyle \int_{B_{4R}}|Y|^{-p}|V|^3\,dY<\infty,\quad
\displaystyle \int_{B_{4R}}|Y|^{-p}|V|\,|Y\cdot\nabla V|\,dY<\infty,\\[2mm]
\displaystyle \int_{B_{4R}}|Y|^{-p-1}|V|^3\,dY<\infty
\end{array}\right\}.
\]
The larger ball is only a domain buffer: all gauge integrals are supported in
\(B_{2R}\), while \((\lambda,c)\) is restricted near \((1,0)\) so that
\(c+\lambda B_{2R}\subset B_{4R}\).

For the actual repaired-gauge implicit-function step we use the following
strong local chart class.

\begin{definition}[Strong translation-admissible gauge class]
\label{paper2:def:strong-translation-admissible-gauge-class}
Fix \(0<p<3\) and \(R>0\).  Let \(\chi_R\in C_c^\infty(B_{2R})\).
A profile \(V\) belongs to \(\mathcal A^{\rm tr,C^1}_{p,R}\) if
\[
   V\in \mathcal A_{p,R}\cap\mathcal V_R
\]
and there exists \(\delta_0>0\) such that the scale functional
\((c,\lambda)\mapsto\mathcal J(c,\lambda;V)\) and the three centering
moments \((c,\lambda)\mapsto\mathcal M_j(c,\lambda;V)\), \(j=1,2,3\),
are finite and actually \(C^1\) for
\[
   |c|+|\lambda-1|<\delta_0,\qquad \lambda>0:
\]
\[
   \mathcal J(c,\lambda;V)
   :=
   \int_{B_{2R}}\chi_R(Y)|Y|^{-p}|V(c+\lambda Y)|^3\,dY ,
\]
\[
\begin{aligned}
   \mathcal J_{c_k}(c,\lambda;V)
   &:=
   \lambda^{-1}
   \left[
   p\int_{B_{2R}}\chi_R(Y)Y_k|Y|^{-p-2}|V(c+\lambda Y)|^3\,dY \right.\\
   &\hspace{2.7cm}\left.
   -
   \int_{B_{2R}}\partial_k\chi_R(Y)|Y|^{-p}|V(c+\lambda Y)|^3\,dY
   \right],
\end{aligned}
\]
and
\[
\begin{aligned}
   \mathcal J_{\lambda}(c,\lambda;V)
   &:=
   -\lambda^{-1}
   \int_{B_{2R}}
   \left[(3-p)\chi_R(Y)+Y\cdot\nabla\chi_R(Y)\right]
   |Y|^{-p}|V(c+\lambda Y)|^3\,dY .
\end{aligned}
\]
For the centering moments set
\[
   \mathcal M_j(c,\lambda;V)
   :=
   \int_{B_{2R}}Y_j\chi_R(Y)|V(c+\lambda Y)|^2\,dY,
   \qquad j=1,2,3,
\]
and require that their actual derivatives are the displayed maps
\[
   \mathcal M_{j,c_k}(c,\lambda;V)
   :=
   -\lambda^{-1}
   \int_{B_{2R}}
   \bigl(\delta_{jk}\chi_R(Y)+Y_j\partial_k\chi_R(Y)\bigr)
   |V(c+\lambda Y)|^2\,dY,
\]
\[
   \mathcal M_{j,\lambda}(c,\lambda;V)
   :=
   -\lambda^{-1}
   \int_{B_{2R}}
   Y_j\bigl(4\chi_R(Y)+Y\cdot\nabla\chi_R(Y)\bigr)
   |V(c+\lambda Y)|^2\,dY,
\]
for \(j,k\in\{1,2,3\}\).  In other words,
\[
  \partial_{c_k}\mathcal J=\mathcal J_{c_k},
  \quad
  \partial_\lambda\mathcal J=\mathcal J_\lambda,
  \quad
  \partial_{c_k}\mathcal M_j=\mathcal M_{j,c_k},
  \quad
  \partial_\lambda\mathcal M_j=\mathcal M_{j,\lambda}
\]
on the chart.  Thus these are not merely candidate derivative maps: the strong
chart requires actual \(C^1\) differentiability of the gauge functionals.  If
this property fails on a selected branch, the branch is routed to the
repaired-gauge degeneracy row before any implicit-function theorem is used.
\end{definition}

\begin{lemma}[Gauge chart entry test]
\label{paper2:lem:strong-gauge-chart-entry-test}
\label{paper2:lem:gauge-differentiability-entry}
At every point where the repaired gauge normalization is used, exactly one of
the following occurs:
\begin{enumerate}[label=\textup{(\roman*)}]
\item the profile belongs to the strong translation-admissible class
\(\mathcal A^{\rm tr,C^1}_{p,R}\), and the gauge Jacobian is nondegenerate;
\item the profile fails to belong to
\(\mathcal A^{\rm tr,C^1}_{p,R}\), or the gauge Jacobian degenerates.  This is
the repaired-gauge degeneracy row.
\end{enumerate}
\end{lemma}

\begin{proof}
This is the ordered gauge test in the local state-space algorithm.  The test is
performed before applying the implicit-function theorem.  If the strong chart
conditions and Jacobian nondegeneracy hold, the branch proceeds through the
regular gauge chart.  If either condition fails, the first failed test is the
repaired-gauge degeneracy row.  No closure theorem is used in this routing
step; closure of the degeneracy row is supplied later by the gauge-degeneracy
closure theorem.
\end{proof}

\begin{definition}[Gauge actions]\label{paper2:def:gauge-action}
For \((\lambda,c)\in (0,\infty)\times\RR^3\), define
\[
V_{\lambda,c}(Y):=\lambda V(c+\lambda Y).
\]
For pressure profiles use
\[
P_{\lambda,c}(Y):=\lambda^2P(c+\lambda Y).
\]
We also write \(V^{\lambda,c}:=V_{\lambda,c}\) and \(P^{\lambda,c}:=P_{\lambda,c}\).
\end{definition}

\begin{definition}[Finite-dimensional gauge map]\label{paper2:def:gauge-map}
For \(V\in\mathcal V_R\) and \( (\lambda,c) \in (0,\infty)\times\RR^3\), set
\[
\mathcal F(\lambda,c;V)
:=\bigl(\mathcal F_0,\mathcal F_1,\mathcal F_2,\mathcal F_3\bigr),
\]
with
\[
\mathcal F_0(\lambda,c;V):=
\int_{\RR^3}\chi_R(Y)\,|Y|^{-p}|V_{\lambda,c}(Y)|^3\dd Y-\Theta_0,\qquad
\mathcal F_j(\lambda,c;V):=\int_{\RR^3}Y_j\,\chi_R(Y)\,|V_{\lambda,c}(Y)|^2\dd Y
\]
for \(j=1,2,3\). The gauge conditions are \(\mathcal F(\lambda,c;V)=0\).
Codomain is \(\RR^4\), domain \((0,\infty)\times\RR^3\times\mathcal V_R\).
\end{definition}

\begin{definition}[Localized modulation matrix]\label{paper2:def:gauge-nondeg}
For \(J\subset I_*\) and a profile path \(W:J\to\mathcal V_R\), define
\[
J_W(\tau):=
D_{(\lambda,c)}\mathcal F\bigl(1,0;W(\tau)\bigr)
\in\mathbb R^{4\times4}.
\]
The rows are ordered as the scale constraint followed by the three centering
constraints.
\end{definition}

\begin{lemma}[Weighted integrability and differentiation under integrals]\label{paper2:lem:weighted-finite}
Assume \(V,W\in\mathcal A_{p,R}\). Then \(G_0(V):=\int\chi_R|Y|^{-p}|V|^3\)
is finite and
\[
\frac{\dd}{\dd\epsilon}\Big|_{\epsilon=0}
\int \chi_R|Y|^{-p}|V+\epsilon W|^3\,\dd Y
=\int 3\chi_R|Y|^{-p}|V|\,V\cdot W\,\dd Y.
\]
\end{lemma}
\begin{proof}
For \(W\in\mathcal A_{p,R}\), define
\[
F_\varepsilon(Y):=\chi_R(Y)|Y|^{-p}|V(Y)+\varepsilon W(Y)|^3.
\]
For each \(Y\), \(\varepsilon\mapsto F_\varepsilon(Y)\) is \(C^1\) and
\[
\partial_\varepsilon F_\varepsilon
=3\chi_R|Y|^{-p}|V+\varepsilon W|(V+\varepsilon W)\cdot W.
\]
For \(|\varepsilon|\le1\),
\[
|F_\varepsilon(Y)|+|\partial_\varepsilon F_\varepsilon|
\le C\chi_R|Y|^{-p}\bigl(|V|^3+|W|^3\bigr)\in L^1(B_{2R})
\]
by the defining weighted \(L^3\) condition of \(\mathcal A_{p,R}\).
Hence dominated convergence allows differentiation under the integral, and
\[
\frac{\dd}{\dd\varepsilon}\Big|_{\varepsilon=0}\int_{B_{2R}}\chi_R|Y|^{-p}|V+\epsilon W|^3\,\dd Y
=\int_{B_{2R}}3\chi_R|Y|^{-p}|V|\,V\cdot W\,\dd Y.
\]
\end{proof}

\begin{lemma}[Gauge covariance]\label{paper2:lem:gauge-covariance}
For \(V\in\mathcal V_R\), \((\lambda,c)\in(0,\infty)\times\RR^3\), and
\(\alpha=0,\dots,3\),
\[
\mathcal F_\alpha(1,0;V_{\lambda,c})=\mathcal F_\alpha(\lambda,c;V).
\]
\end{lemma}

\begin{proof}
For each component this is a direct substitution:
\[
\mathcal F_0(1,0;V_{\lambda,c})
=\int\chi_R(Y)|Y|^{-p}|V_{\lambda,c}(Y)|^3-\Theta_0\,\dd Y
=\mathcal F_0(\lambda,c;V),
\]
\[
\mathcal F_j(1,0;V_{\lambda,c})
=\int Y_j\chi_R(Y)|V_{\lambda,c}(Y)|^2\,\dd Y
=\mathcal F_j(\lambda,c;V),\quad j=1,2,3.
\]
\end{proof}

\begin{lemma}[Weighted gauge functionals are \(C^1\) on the strong chart]
\label{paper2:lem:weighted-gauge-C1-strong-chart}
\label{paper2:lem:weighted-translation-differentiability}
Let \(0<p<3\), let \(\chi_R\in C_c^\infty(B_{2R})\), and let
\(V\in\mathcal A^{\rm tr,C^1}_{p,R}\).  Then the weighted scale and centering
gauge functionals are \(C^1\) on the strong chart and have the displayed
derivatives of \cref{paper2:def:strong-translation-admissible-gauge-class}.
\end{lemma}

\begin{proof}
This is the defining property of the strong chart
\(\mathcal A^{\rm tr,C^1}_{p,R}\).  If a selected branch does not satisfy this
actual \(C^1\) property, the ordered gauge-chart entry test
\cref{paper2:lem:strong-gauge-chart-entry-test} routes it to the
repaired-gauge degeneracy row before the implicit-function theorem is invoked.
\end{proof}

\begin{proposition}[Differentiability of \(\mathcal F\)]\label{paper2:prop:gauge-C1}
For each fixed \(V\in\mathcal A^{\rm tr,C^1}_{p,R}\), the map
\((\lambda,c)\mapsto \mathcal F(\lambda,c;V)\) is \(C^1\) in a neighborhood of
\((1,0)\). In each component \( \mathcal F_\alpha\) all derivatives under the
integral sign are legitimate on the retained gauge chart.  By
\cref{paper2:lem:strong-gauge-chart-entry-test}, every retained
repaired-gauge branch either lies in \(\mathcal A^{\rm tr,C^1}_{p,R}\) on the
selected chart, or has already exited through the repaired-gauge degeneracy
row before this proposition is invoked.
\end{proposition}

\begin{proof}
The \(C^1\) regularity of the gauge functionals follows from
\cref{paper2:lem:weighted-gauge-C1-strong-chart}.  The nondegeneracy of the
Jacobian is part of the regular gauge-chart alternative in
\cref{paper2:lem:strong-gauge-chart-entry-test}.  The implicit-function theorem
therefore gives the claimed local \(C^1\) gauge map.  If either the strong
translation-admissible chart condition or the Jacobian nondegeneracy fails, the
branch is routed to the repaired-gauge degeneracy row and is not treated by
this proposition.
Thus each component admits pointwise \(C^1\) formulas and these are continuous in
\((\lambda,c)\) near \((1,0)\). Therefore \((\lambda,c)\mapsto\mathcal F(\lambda,c;V)\in C^1\)
near \((1,0)\).
\end{proof}

\begin{proposition}[Full local Jacobian at \((1,0)\)]\label{paper2:prop:full-jacobian}
Let \(V\in\mathcal A^{\rm tr,C^1}_{p,R}\).
Set \(V_0(Y)=V(Y)\).
At \((\lambda,c)=(1,0)\), the matrix
\[
J(V):=D_{(\lambda,c)}\mathcal F(1,0;V)\in \RR^{4\times4}
\]
has rows
\[
J_\alpha=\bigl(\partial_\lambda \mathcal F_\alpha,\ \partial_{c_1}\mathcal F_\alpha,\
\partial_{c_2}\mathcal F_\alpha,\ \partial_{c_3}\mathcal F_\alpha\bigr),\quad \alpha=0,\dots,3.
\]
Writing
\[
\mathcal J_*:=\mathcal J(0,1;V),\qquad
\mathcal J_{\lambda,*}:=\mathcal J_{\lambda}(0,1;V),\qquad
\mathcal J_{c_k,*}:=\mathcal J_{c_k}(0,1;V),
\]
and
\[
\mathcal M_{j,*}:=\mathcal M_j(0,1;V),\qquad
\mathcal M_{j,\lambda,*}:=\mathcal M_{j,\lambda}(0,1;V),\qquad
\mathcal M_{j,c_k,*}:=\mathcal M_{j,c_k}(0,1;V),
\]
one has
\begin{align}
\partial_\lambda \mathcal F_0(1,0;V)
&=
3\mathcal J_*+\mathcal J_{\lambda,*}
=p\int\chi_R|Y|^{-p}|V_0|^3\,\dd Y
-\int (Y\cdot\nabla\chi_R)|Y|^{-p}|V_0|^3\,\dd Y,\label{paper2:eq:jac00}\\
\partial_{c_k}\mathcal F_0(1,0;V)
&=\mathcal J_{c_k,*}
=p\int \chi_R(Y)Y_k|Y|^{-p-2}|V_0(Y)|^3\,dY
-\int (\partial_k\chi_R)(Y)|Y|^{-p}|V_0(Y)|^3\,dY,\label{paper2:eq:jac0k}\\
\partial_\lambda \mathcal F_j(1,0;V)
&=2\mathcal M_{j,*}+\mathcal M_{j,\lambda,*}
=-\int_{\RR^3}Y_j\bigl(2\chi_R+Y\cdot\nabla\chi_R\bigr)|V_0|^2\,\dd Y,\label{paper2:eq:jacjlambda}\\
\partial_{c_k}\mathcal F_j(1,0;V)
&=\mathcal M_{j,c_k,*}
=-\int_{\RR^3}\bigl(\delta_{jk}\chi_R+Y_j\partial_k\chi_R\bigr)|V_0|^2\,\dd Y,\label{paper2:eq:jacjk}
\end{align}
for \(j,k=1,2,3\).
\end{proposition}

\begin{proof}
By gauge covariance,
\[
\mathcal F_0(\lambda,c;V)=\lambda^3\mathcal J(c,\lambda;V)-\Theta_0,
\qquad
\mathcal F_j(\lambda,c;V)=\lambda^2\mathcal M_j(c,\lambda;V).
\]
Since \(V\in\mathcal A^{\rm tr,C^1}_{p,R}\),
\cref{paper2:lem:weighted-gauge-C1-strong-chart} gives the \(C^1\) chart
derivatives for \(\mathcal J\), and the same strong chart package supplies the
\(C^1\) derivatives for \(\mathcal M_j\).  Applying the product rule at
\((c,\lambda)=(0,1)\) yields
\[
\partial_\lambda\mathcal F_0(1,0;V)=3\mathcal J(0,1;V)+\mathcal J_\lambda(0,1;V),
\qquad
\partial_{c_k}\mathcal F_0(1,0;V)=\mathcal J_{c_k}(0,1;V),
\]
and
\[
\partial_\lambda\mathcal F_j(1,0;V)=2\mathcal M_j(0,1;V)+\mathcal M_{j,\lambda}(0,1;V),
\qquad
\partial_{c_k}\mathcal F_j(1,0;V)=\mathcal M_{j,c_k}(0,1;V).
\]
The displayed integral formulas are exactly the explicit \((0,1)\)-formulas for
these chart derivatives.  In particular, the singular-weight first row is read
through \(\mathcal J_\lambda\) and \(\mathcal J_{c_k}\), not through the raw
term \(\chi_R|Y|^{-p}|V|V\cdot\partial_kV\).
\end{proof}

\begin{proposition}[Translation block and Schur-complement invertibility]\label{paper2:thm:quant-jac}
Let \(V:J\to\mathcal A^{\rm tr,C^1}_{p,R}\) satisfy the repaired gauge constraints
\[
\mathcal F(1,0;V(\tau))=0\qquad\text{for a.e. }\tau\in J.
\]
Write the modulation matrix in block form
\[
J(V(\tau))=
\begin{pmatrix}
\alpha(\tau)&u(\tau)^T\\
v(\tau)&A(\tau)
\end{pmatrix},
\]
where \(\alpha=\partial_\lambda\mathcal F_0\), \(u_\ell=\partial_{c_\ell}\mathcal F_0\),
\(v_j=\partial_\lambda\mathcal F_j\), and \(A_{j\ell}=\partial_{c_\ell}\mathcal F_j\).
Assume that for a.e. \(\tau\in J\)
\[
m_\chi(\tau):=\int |V(Y,\tau)|^2\chi_R(Y)\,dY\ge m_0>0,
\]
\[
\int_{A_R}|V(Y,\tau)|^2\,dY\le\varepsilon_0,\qquad
3C_\chi\varepsilon_0\le \frac{m_0}{2},
\]
where \(A_R=\{R\le |Y|\le2R\}\) and
\[
C_\chi:=\max_{j,\ell}\|Y_j\partial_\ell\chi_R\|_{L^\infty}.
\]
Assume also
\[
C_u:=\operatorname*{ess\,sup}_{\tau\in J}\int |Y|^{-p-1}|V(Y,\tau)|^3\,dY<\infty
\]
and, with \(C_\nabla:=\||Y|\nabla\chi_R\|_{L^\infty}\),
\[
6(p+C_\nabla)\left(\frac{2}{m_0}\right)C_u\,R C_\nabla\varepsilon_0
\le \frac{p\Theta_0}{2}.
\]
Then \(J(V(\tau))\) is invertible for a.e. \(\tau\in J\), and
\[
\operatorname*{ess\,sup}_{\tau\in J}\|J(V(\tau))^{-1}\|_{\infty\to\infty}
\le K(R,p,\Theta_0,m_0,C_u,\varepsilon_0,\chi_R).
\]
\end{proposition}
\begin{proof}
By \cref{paper2:prop:full-jacobian}, the first row is given by the strong-chart
derivatives
\[
\alpha=3\mathcal J(0,1;V(\tau))+\mathcal J_\lambda(0,1;V(\tau)),
\qquad
u_\ell=\mathcal J_{c_\ell}(0,1;V(\tau)),
\]
and the remaining rows by
\[
v_j=2\mathcal M_j(0,1;V(\tau))+\mathcal M_{j,\lambda}(0,1;V(\tau)),
\qquad
A_{j\ell}=\mathcal M_{j,c_\ell}(0,1;V(\tau)).
\]
Thus the Jacobian is evaluated entirely in the strong chart already built into
\(\mathcal A^{\rm tr,C^1}_{p,R}\), and no singular derivative of \(V\) under the
weight is used.  The entries are finite by \(V(\tau)\in\mathcal A_{p,R}\) and the
compact support of \(\chi_R\).  On the gauge surface,
\(\mathcal J(0,1;V(\tau))=\Theta_0\).  Hence the scale entry satisfies the
signed transversality bound
\[
\alpha(\tau)=
p\int\chi_R|Y|^{-p}|V(Y,\tau)|^3\,dY
-\int (Y\cdot\nabla\chi_R)|Y|^{-p}|V(Y,\tau)|^3\,dY
\ge p\Theta_0
\]
on the gauge surface, because \(Y\cdot\nabla\chi_R\le0\). Thus
\(|\alpha|\ge p\Theta_0\).

For the centering block, integration by parts gives
\[
A_{j\ell}
=-\delta_{j\ell}\int |V|^2\chi_R
-\int |V|^2Y_j\partial_\ell\chi_R.
\]
Writing \(A=-m_\chi I_3+E\), the annular support of \(\nabla\chi_R\) gives
\[
|E_{j\ell}|\le C_\chi\int_{A_R}|V|^2\le C_\chi\varepsilon_0,\qquad
\|E\|_{\infty\to\infty}\le3C_\chi\varepsilon_0\le m_0/2.
\]
Thus \(A=-m_\chi(I_3-m_\chi^{-1}E)\) and
\(\|m_\chi^{-1}E\|_{\infty\to\infty}\le1/2\), so \(A^{-1}\) exists and
\(\|A^{-1}\|_{\infty\to\infty}\le2/m_0\).

For the scale-to-translation block, the scale-constraint formula gives
\[
\left|u_\ell\right|
\le p\int |Y|^{-p-1}|V|^3\,dY
+\int |\partial_\ell\chi_R|\,|Y|^{-p}|V|^3\,dY
\le (p+C_\nabla)C_u,
\]
so \(\|u\|_\infty\le (p+C_\nabla)C_u\). For the translation-to-scale block,
\[
v_j=-\int |V|^2Y_j(Y\cdot\nabla\chi_R)\,dY
\]
because the centering constraint cancels the interior term. Therefore
\[
\|v\|_\infty\le2R\,C_\nabla\varepsilon_0.
\]
The Schur complement \(S=\alpha-u^TA^{-1}v\) satisfies
\[
|u^TA^{-1}v|
\le 3\|u\|_\infty\|A^{-1}\|_{\infty\to\infty}\|v\|_\infty
\le 6(p+C_\nabla)\left(\frac{2}{m_0}\right)C_uRC_\nabla\varepsilon_0
\le \frac{p\Theta_0}{2}.
\]
Since \(|\alpha|\ge p\Theta_0\), this gives \(|S|\ge p\Theta_0/2\).
The block inverse formula
\[
J^{-1}=
\begin{pmatrix}
S^{-1} & -S^{-1}u^TA^{-1}\\
-A^{-1}vS^{-1}&A^{-1}+A^{-1}vS^{-1}u^TA^{-1}
\end{pmatrix}
\]
then gives the displayed bound.
\end{proof}

\subsection{Retained repaired-gauge selection in time}\label{paper2:sec:gauge-solve}

\begin{definition}[Retained nondegenerate repaired-gauge regime]\label{paper2:def:admissible-gauge-path}
Let \(J\subset I_*\) be compact. The retained nondegenerate repaired-gauge
regime over the preliminary profile \(V_{\rm pre}\) consists of branches for
which there are
\[
\lambda\in AC(J;(0,\infty)),\qquad c\in AC(J;\mathbb R^3),
\]
such that, for a.e. \(\tau\in J\),
\[
V(\tau):=V_{\rm pre}^{\lambda(\tau),c(\tau)}
\in\mathcal A^{\rm tr,C^1}_{p,R},\qquad
\mathcal F(1,0;V(\tau))=0.
\]
The path is quantitatively nondegenerate if the translation-block, annular,
weighted first-moment, and Schur-complement bounds in
Proposition~\ref{paper2:thm:quant-jac} hold on \(J\).
If no such local AC selection exists, if the strong chart condition
\(\mathcal A^{\rm tr,C^1}_{p,R}\) fails, or if the quantitative
nondegeneracy hypotheses fail persistently, the branch lies in the
gauge-degenerate,
multibubble, or cascade alternatives rather than in the retained regime.
\end{definition}

\begin{proposition}[Gauge regularity on the retained regime]\label{paper2:prop:preliminary-correction}
Let \((\lambda,c)\) be the repaired-gauge selection of the retained
nondegenerate regime on \(J\). Then
\[
\mathcal F\bigl(\lambda(\tau),c(\tau);V_{\rm pre}(\tau)\bigr)=0
\quad\text{for a.e. }\tau\in J,
\]
and \((\lambda,c)\in AC(J;\mathbb R^4)\) by definition. If the path is
quantitatively nondegenerate, then
\[
\operatorname*{ess\,sup}_{\tau\in J}\|J(V(\tau))^{-1}\|<\infty.
\]
\end{proposition}
\begin{proof}
The constraint follows from Lemma~\ref{paper2:lem:gauge-covariance} applied to
\(V(\tau)=V_{\rm pre}^{\lambda(\tau),c(\tau)}\). The absolute continuity is part
of Definition~\ref{paper2:def:admissible-gauge-path}. The hypotheses of
\Cref{paper2:thm:quant-jac} are part of the retained nondegenerate regime
through the strong chart entry test
\Cref{paper2:lem:strong-gauge-chart-entry-test}. The uniform inverse bound is
Proposition~\ref{paper2:thm:quant-jac}.
\end{proof}

\begin{definition}[Repaired chart]\label{paper2:def:repaired-chart}
Define
\[
\Lambda(\tau):=\Lambda_{\rm pre}(\tau)\lambda(\tau),\qquad
X(\tau):=X_{\rm pre}(\tau)+\Lambda_{\rm pre}(\tau)c(\tau),
\]
\[
V(\tau):=V_{\rm pre}^{\lambda(\tau),c(\tau)},\qquad
P(\tau):=P_{\rm pre}^{\lambda(\tau),c(\tau)}\ \text{(same physical pressure pull-back under }\Phi).
\]
\end{definition}

\begin{theorem}[AC repaired parameters]\label{paper2:thm:repaired-existence}
Assume \(\mathrm{SC}(R_*,r_*,M,\eta,\kappa)\), and suppose the branch lies in the
retained nondegenerate repaired-gauge regime on \(J\), with selection
\((\lambda,c)\). Then the corrected
chart \((\Lambda,X)\) is absolutely continuous on \(J\), satisfies
\[
a(\tau)=\frac{\Lambda'(\tau)}{\Lambda(\tau)}\in L^1_{\rm loc}(J),\qquad
b(\tau)=\frac{X'(\tau)}{\Lambda(\tau)}\in L^1_{\rm loc}(J),
\]
and the gauge constraints \(\mathcal F(1,0;V(\tau))=0\) hold a.e.
Equivalently,
\[
\mathcal F(\lambda(\tau),c(\tau);V_{\rm pre}(\tau))=0
\]
for the repaired-gauge selection \((\lambda,c)\).
\end{theorem}
\begin{proof}
By Definition~\ref{paper2:def:admissible-gauge-path}, \((\lambda,c)\in AC(J)\) and the
repaired gauge constraints hold for \(V(\tau)\) a.e.
By definition \(V(\tau)=V_{\rm pre}^{\lambda(\tau),c(\tau)}\), so Lemma~\ref{paper2:lem:gauge-covariance}
applied at \((\lambda,c)=(\lambda(\tau),c(\tau))\) gives
\[
\mathcal F(1,0;V_{\rm pre}^{\lambda,c})=\mathcal F(\lambda,c;V_{\rm pre})
\]
implies \(\mathcal F(1,0;V(\tau))=0\) a.e.

Since \(\Lambda=\Lambda_{\rm pre}\lambda\) and \(X=X_{\rm pre}+\Lambda_{\rm pre}c\),
and \(\Lambda_{\rm pre},X_{\rm pre}\) are AC by \textup{(SC5)}, the products and sums
preserve AC, hence \((\Lambda,X)\in AC(I_*)\). Differentiating,
\[
a=\frac{\Lambda'}{\Lambda}
=\frac{\Lambda_{\rm pre}'}{\Lambda_{\rm pre}}+\frac{\lambda'}{\lambda}\in L^1_{\rm loc},\qquad
b=\frac{X'}{\Lambda}
=\frac{X_{\rm pre}'+\Lambda_{\rm pre}'c+\Lambda_{\rm pre}c'}
{\Lambda_{\rm pre}\lambda}\in L^1_{\rm loc}.
\]
This gives the claimed regularity.
\end{proof}

\begin{remark}[Gauge degeneracy]
Failure of invertibility of \(J\) is the gauge-degenerate alternative to the
retained single-core regime for this fixed renormalization mechanism.  In the
Type II decomposition of alternatives this case belongs to the
multibubble/cascade and scale-rigidity alternatives, rather than to the
representation theorem.
\end{remark}

\subsection{Pressure transport and decomposition}\label{paper2:sec:pressure}
\begin{lemma}[Renormalized pressure equation]\label{paper2:lem:pressure-eq}
If \(V,P\) satisfy the renormalized Navier--Stokes equation in
Theorem~\ref{paper2:thm:ren-NS}, then
\[
-\Delta P=\partial_i\partial_j(V_iV_j)\qquad\text{in }\mathcal D'(Q_{2R}^*).
\]
\end{lemma}
\begin{proof}
Take test \(\varphi\in C_c^\infty(Q_{2R}^*)\). Since \(V,P\in L^2_{\rm loc}\times
L^{3/2}_{\rm loc}\) on \(Q_{2R}^*\), all pairings below are legitimate.
Expand
\[
\nn\cdot\bigl((V\cdot\nabla)V\bigr)
=\sum_{i,j}\partial_i\partial_j(V_iV_j)-\sum_{i,j}\partial_iV_i\,\partial_jV_j.
\]
Because \(\nn\cdot V=0\), this reduces to \(\partial_i\partial_j(V_iV_j)\).
Taking divergence of
\[
\partial_\tau V+(V\cdot\nabla)V+\nabla P
=\nu\Delta V+a(V+Y\cdot\nabla V)+b\cdot\nabla V
\]
and using \(\nn\cdot V=0\), \(\nn\cdot(V+Y\cdot\nabla V)=0\), \(\nn\cdot(b\cdot\nabla V)=b\cdot\nn(\nn\cdot V)=0\),
and \(\nn\cdot\Delta V=\Delta(\nn\cdot V)=0\), the following identity holds in \(\mathcal D'\):
\[
\Delta P+\partial_i\partial_j(V_iV_j)=0.
\]
Hence \(-\Delta P=\partial_i\partial_j(V_iV_j)\) on \(Q_{2R}^*\).
\end{proof}

\begin{theorem}[Localized pressure decomposition]\label{paper2:thm:pressure-decomp}
Let \(\zeta\in C_c^\infty(B_{2R})\), \(\zeta\equiv1\) on \(B_R\), and assume
\(V\in L^1_{\mathrm{loc}}(I_*;L^3(B_{2R}))\), \(P\in L^1_{\mathrm{loc}}(I_*;L^{3/2}(B_{2R}))\).
For a.e. \(\tau\in I_*\) there exists
\[
P(\cdot,\tau)=P_{\rm loc}(\cdot,\tau)+H(\cdot,\tau)
\]
on \(B_R\) such that
\[
\begin{cases}
-\Delta P_{\rm loc}(\cdot,\tau)=\partial_i\partial_j\bigl(\zeta V_iV_j\bigr),\\
-\Delta H(\cdot,\tau)=0\ \text{on }B_R,\quad
\fint_{B_R}H(\cdot,\tau)\,dY=0.
\end{cases}
\]
Moreover
\[
\|P_{\rm loc}(\tau)\|_{L^{3/2}(B_R)}
\lesssim \|V(\tau)\|_{L^3(B_{2R})}^2,\qquad
\|H(\tau)\|_{L^q(B_r)}\lesssim_{q,R,r}
\|P(\tau)\|_{L^{3/2}(B_{2R})}+\|V(\tau)\|_{L^3(B_{2R})}^2
\]
for each \(1\le q<\infty\) and each \(0<r<R\), with constants depending only on \(R\), \(r\), and \(q\).
\end{theorem}

\begin{proof}
For fixed \(\tau\), set
\[
P_{\rm loc}= \mathcal R_i\mathcal R_j(\zeta V_iV_j),\qquad
H:=P-P_{\rm loc}.
\]
\(\zeta V_iV_j\in L^{1}_{\mathrm{loc}}(I_*;L^{3/2}(B_{2R}))\), so the Calderón--Zygmund
operator \(f\mapsto \mathcal R_i\mathcal R_j f\) gives
\[
P_{\rm loc}(\cdot,\tau)=\mathcal R_i\mathcal R_j(\zeta V_iV_j)(\cdot,\tau),
\qquad
\|P_{\rm loc}(\tau)\|_{L^{3/2}(B_R)}
\lesssim \|V(\tau)\|_{L^3(B_{2R})}^2.
\]
Because \(\mathcal R_i\mathcal R_j\) is linear bounded on \(L^{3/2}\),
\[
P_{\rm loc}\in L^{1}_{\mathrm{loc}}(I_*;L^{3/2}(B_R)).
\]
Since \(-\Delta P_{\rm loc}=\partial_i\partial_j(\zeta V_iV_j)\) and \(\zeta\equiv1\) on \(B_R\),
 \[
-\Delta H(\cdot,\tau)
=\partial_i\partial_j\bigl((1-\zeta)V_iV_j\bigr)=0 \quad\text{in }B_R,
\]
so \(H(\cdot,\tau)\) is harmonic on \(B_R\). Imposing \(\fint_{B_R}H=0\) fixes the
time-dependent additive constant uniquely in each time slice.
For \(q\in[1,\infty)\) and \(0<r<R\), harmonicity and local
\(L^p\)-to-\(L^q\) interior estimates imply
\[
\|H(\tau)\|_{L^q(B_r)}\lesssim_{q,R,r}\|H(\tau)\|_{L^{3/2}(B_R)}.
\]
Finally, \(\|H(\tau)\|_{L^{3/2}(B_R)}\le \|P(\tau)\|_{L^{3/2}(B_R)}+\|P_{\rm loc}(\tau)\|_{L^{3/2}(B_R)}\),
giving the announced bound in terms of \(P\) and \(V\).
\end{proof}
\begin{remark}[Measurability and normalization]\label{paper2:rem:pressure-time}
If \(V,P\) are measurable in \(\tau\), then \(\tau\mapsto P_{\rm loc}(\cdot,\tau)\)
is measurable in \(L^{3/2}(B_R)\) by boundedness of \(\mathcal R_i\mathcal R_j\), hence
\(\tau\mapsto H(\cdot,\tau)\) is measurable in each \(L^q(B_r)\), \(0<r<R\).
The normalization \(\fint_{B_R}H(\cdot,\tau)\,dY=0\) fixes the time-dependent additive constant
in \(H\) uniquely for each \(\tau\), and therefore the decomposition is a legitimate
renormalized pressure splitting.
\end{remark}

\subsection{Modulation from differentiated constraints}\label{paper2:sec:modulation}
\begin{definition}[Admissible modulation-differentiability row]\label{paper2:def:mod-diff-row}
On a compact window \(J\subset I_*\), a retained repaired-gauge branch satisfies
the modulation-differentiability row if, for each \(\alpha=0,1,2,3\),
\[
G_\alpha(\tau):=\mathcal F_\alpha(1,0;V(\tau))
\]
belongs to \(AC(J)\), and there are functions
\(\mathcal R_\alpha\in L^1(J)\) such that
\begin{equation}\label{paper2:eq:mod-row}
G_\alpha'(\tau)
=-\mathcal R_\alpha(\tau)
+a(\tau)\,\partial_\lambda\mathcal F_\alpha(1,0;V(\tau))
+\sum_{k=1}^3 b_k(\tau)\,\partial_{c_k}\mathcal F_\alpha(1,0;V(\tau))
\end{equation}
for a.e. \(\tau\in J\).  The terms \(\mathcal R_\alpha\) are defined as the
limits obtained by testing the represented renormalized equation against smooth
compactly supported truncations of the Gateaux gradients of
\(\mathcal F_\alpha\), then passing to the limit in the local energy and pressure
topologies.  Thus \(\mathcal R_\alpha\) is a distributional truncated-test limit,
not the unsupported pairing of an \(L^1_x\) coefficient with a raw
\(\partial_\tau V\).

If this row fails on a compact window, the branch is routed to the
modulation-failure exit rather than to the retained modulation theorem.
\end{definition}

\begin{lemma}[Differentiated repaired constraints under the modulation row]\label{paper2:lem:diff-f}
If the repaired branch satisfies \cref{paper2:def:mod-diff-row} on \(J\), then
the repaired constraints have the a.e. derivative identity
\eqref{paper2:eq:mod-row} on \(J\), with
\(\mathcal R=(\mathcal R_0,\ldots,\mathcal R_3)\in L^1(J)^4\).
\end{lemma}

\begin{proof}
This is exactly the retained-row assertion in
\cref{paper2:def:mod-diff-row}.  The definition records the admissible limiting
procedure: one first pairs the equation with smooth compactly supported
truncations of the Gateaux gradients and only then passes to the limit.  No
step uses an \(L^1_x\)-against-\(L^2_x\) pairing with \(\partial_\tau V\).
\end{proof}

\begin{lemma}[Profile and chart direction identities]\label{paper2:lem:profile-direction}
For \(\mathcal A_\tau=V+Y\cdot\nn V\) and \(\alpha=0,1,2,3\),
\[
D_V\mathcal F_\alpha(1,0;V)[\mathcal A_\tau]
=\partial_\lambda\mathcal F_\alpha(1,0;V).
\]
For the translation directions, the Jacobian entries are read from the strong
chart of \cref{paper2:prop:full-jacobian}:
\[
\partial_{c_k}\mathcal F_0(1,0;V)=\mathcal J_{c_k}(0,1;V),
\qquad
\partial_{c_k}\mathcal F_j(1,0;V)=\mathcal M_{j,c_k}(0,1;V)
\]
for \(j,k=1,2,3\).  In particular, the singular weighted translation entries in
the first row are not identified through the raw profile variation
\(D_V\mathcal F_0(1,0;V)[\partial_kV]\), but through the strong-chart
derivative \(\mathcal J_{c_k}\).
\end{lemma}

\begin{proof}
From \(V_{\lambda,c}(Y)=\lambda V(c+\lambda Y)\),
\[
\partial_\lambda V_{\lambda,c}\big|_{(1,0)}=V+Y\cdot\nabla V=\mathcal A_\tau.
\]
Substituting this into the profile derivative formula for \(\mathcal F_\alpha\)
gives the scale identity.  The translation identities are exactly the chart
derivatives already computed in \cref{paper2:prop:full-jacobian}.
\end{proof}

\begin{theorem}[Modulation system]\label{paper2:thm:mod-sys}
Let \((V,P)\) be the repaired pair on the compact window under consideration,
and assume the branch satisfies the modulation-differentiability row
\cref{paper2:def:mod-diff-row}. Writing \(\theta:=(a,b)\), let
\(\mathcal R_\alpha\) be the truncated-test limit functions from
\cref{paper2:def:mod-diff-row}.
Then for a.e. \(\tau\in I_*\),
\[
\partial_\lambda\mathcal F_\alpha(1,0;V)\,a(\tau)
+\sum_{k=1}^3\partial_{c_k}\mathcal F_\alpha(1,0;V)\,b_k(\tau)
=\mathcal R_\alpha(\tau),
\]
Equivalently,
\[
J(\tau)\theta(\tau)=\mathcal R(\tau).
\]
Hence
\[
\theta(\tau)=J(\tau)^{-1}\mathcal R(\tau)\in L^1_{\rm loc}(I_*)^4,
\qquad |\theta(\tau)|\le K\,|\mathcal R(\tau)|.
\]
\end{theorem}

\begin{proof}
Set \(G_\alpha(\tau):=\mathcal F_\alpha(1,0;V(\tau))\). By Theorem~\ref{paper2:thm:repaired-existence},
\(G_\alpha(\tau)=0\) a.e.; by \cref{paper2:lem:diff-f}, \(G_\alpha\in AC\) and
\eqref{paper2:eq:mod-row} holds. Hence \(G_\alpha'=0\) a.e., and
\eqref{paper2:eq:mod-row} gives
\[
\partial_\lambda\mathcal F_\alpha(1,0;V)\,a(\tau)
+\sum_{k=1}^3\partial_{c_k}\mathcal F_\alpha(1,0;V)\,b_k(\tau)
=\mathcal R_\alpha(\tau),
\]
which is equivalent to \(J(\tau)\theta(\tau)=\mathcal R(\tau)\).
The row gives \(\mathcal R\in L^1_{\rm loc}(I_*)^4\), and
Proposition~\ref{paper2:thm:quant-jac} gives \(|J^{-1}(\tau)|\le K\) a.e.; hence
\[
\theta(\tau)=J^{-1}(\tau)\mathcal R(\tau),\quad
\theta\in L^1_{\rm loc}(I_*),\quad |\theta|\le K|\mathcal R|\text{ a.e.}
\]
as claimed.
\end{proof}

\begin{corollary}[Regularity of \(a,b\)]
For the repaired pair of Theorem~\ref{paper2:thm:mod-sys}, \(a,b\in L^1_{\rm loc}(I_*)\).
If in addition \(\mathcal R\in L^\infty_{\rm loc}(I_*)\), then \(a,b\in
L^\infty_{\rm loc}(I_*)\).
\end{corollary}

\begin{proof}
By Theorem~\ref{paper2:thm:mod-sys}, \(\theta=(a,b)=J^{-1}\mathcal R\) a.e., and
\(J^{-1}\in L^\infty_{\rm loc}(I_*)\) by Theorem~\ref{paper2:thm:quant-jac}. The first claim is therefore
immediate from \(\mathcal R\in L^1_{\rm loc}\), while the second follows if \(\mathcal R\in L^\infty_{\rm loc}\).
\end{proof}

\subsection{Repaired suitable structure in renormalized variables}\label{paper2:sec:suitable}

\begin{lemma}[AC chart-generated tests are admissible]\label{paper2:lem:ac-chart-tests}
Let \((u,p)\) be suitable on a compact physical cylinder, and let
\(\Lambda,X\in AC(J)\) on a compact interval \(J\), with
\(0<\lambda_-\le \Lambda\le\lambda_+<\infty\) and
\[
a=\frac{\Lambda'}{\Lambda},\qquad b=\frac{X'}{\Lambda}
\quad\text{in }L^1(J).
\]
Let \(t=t(\tau)\) be defined by \(dt/d\tau=\Lambda(\tau)^2\).  If
\(\phi\in C_c^\infty(B_R\times J)\) is nonnegative and its chart pullback has a
positive support buffer inside the physical cylinder, then
\[
\psi(x,t):=\Lambda(\tau(t))^{-1}
\phi\!\left(\frac{x-X(\tau(t))}{\Lambda(\tau(t))},\tau(t)\right)^2
\]
is an admissible nonnegative test for the physical local energy inequality.
The chain-rule identities for \(\nabla_x\psi\), \(\Delta_x\psi\), and
\(\partial_t\psi\) hold for a.e. \(t\).
\end{lemma}

\begin{proof}
For each fixed time, \(x\mapsto\psi(x,t)\) is smooth and compactly supported;
the AC chart gives
\[
\psi,\ \nabla_x\psi,\ \Delta_x\psi\in L^\infty,
\qquad
\partial_t\psi\in L^1_tL^\infty_x
\]
on the compact cylinder, with the time derivative computed by the a.e. chain
rule.  Extend \(\psi\) by zero across a small time buffer and set
\(\psi^\varepsilon:=\rho_\varepsilon*_t\psi\), where \(\rho_\varepsilon\) is a
nonnegative time mollifier and \(\varepsilon\) is smaller than the support
buffer. Then \(\psi^\varepsilon\ge0\), \(\psi^\varepsilon\in C_c^\infty\) in
spacetime with compact support in the same physical cylinder, and
\[
\psi^\varepsilon\to\psi,
\quad \nabla_x\psi^\varepsilon\to\nabla_x\psi,
\quad \Delta_x\psi^\varepsilon\to\Delta_x\psi
\quad\text{in }L^\infty,
\]
while
\[
\partial_t\psi^\varepsilon\to\partial_t\psi
\quad\text{in }L^1_tL^\infty_x.
\]
Apply the physical local energy inequality to \(\psi^\varepsilon\).  The suitable
bounds \(|u|^2\in L_t^\infty L_x^1\), \(|\nabla u|^2\in L^1\), and
\(|u|^3+|p||u|\in L^1\) allow passage to the limit in the right-hand side by
dominated convergence and in the nonnegative dissipation term by lower
semicontinuity. This gives the physical local energy inequality with \(\psi\).
\end{proof}

\begin{proposition}[Transport of suitable local energy inequality]\label{paper2:prop:ren-suitable}
Suppose \((u,p)\) is suitable on \(\mathcal Q^{\rm ph}_{2R}\), and
\((V,P)\) is obtained by the repaired map \(\Phi\).
Set
\[
\tilde a(\tau):=\Lambda(\tau)^{-1}\Lambda'(\tau),\qquad
\tilde b(\tau):=\Lambda(\tau)^{-1}X'(\tau).
\]
Then for every nonnegative \(\phi\in C_c^\infty(Q_R^*)\) and \(\tau_1<\tau_2\),
\[
\begin{aligned}
\mathscr E_\phi(V,P)
&:=\frac12|V|^2\partial_\tau(\phi^2)
+\frac{\nu}{2}|V|^2\Delta_Y\phi^2
+\frac12|V|^2\,V\cdot\nn(\phi^2)\\
&\quad +(P_{\rm loc}+H)\,V\cdot\nn(\phi^2)
-\frac{\tilde a}{2}|V|^2\phi^2
-\frac{\tilde a}{2}|V|^2\,Y\cdot\nn(\phi^2)
-\frac12|V|^2\,\tilde b\cdot\nn(\phi^2).
\end{aligned}
\]
\begin{multline*}
\frac12\int_{B_R}|V(\tau_2)|^2\phi(\cdot,\tau_2)^2
+\nu\int_{\tau_1}^{\tau_2}\!\!\int_{B_R}|\nn V|^2\phi^2\,dY\,d\tau
\le \frac12\int_{B_R}|V(\tau_1)|^2\phi(\cdot,\tau_1)^2\\
 +\int_{\tau_1}^{\tau_2}\!\!\int_{B_R}\mathscr E_\phi(V,P)\,dY\,d\tau.
\end{multline*}
The right-hand side is controlled on compact windows by
\(\|\tilde a\|_{L^1},\|\tilde b\|_{L^1},M\), and the local bounds from
Theorem~\ref{paper2:thm:pressure-decomp}, and support bounds from Section~\ref{paper2:sec:domains}.
Hence \((V,P)\) is locally suitable on \(Q_R^*\).
\end{proposition}

\begin{proof}
Set
\[
\psi(x,t):=\Lambda(\tau(t))^{-1}\phi\!\left(\frac{x-X(\tau(t))}{\Lambda(\tau(t))},\tau(t)\right)^2,
\qquad
\tau=\tau(t),\qquad d\tau/dt=\Lambda^{-2},
\]
and let \(t_i=t(\tau_i)\), \(i=1,2\).  By
Proposition~\ref{paper2:prop:prelim-frame}, Section~\ref{paper2:sec:domains},
and \cref{paper2:lem:ac-chart-tests}, the physical local energy inequality is
valid with this nonnegative AC chart-generated test \(\psi\).
\[
\nabla_x\psi=\Lambda^{-2}\nabla_Y\phi^2,\quad
\Delta_x\psi=\Lambda^{-3}\Delta_Y\phi^2,\quad
\partial_t\psi
=-\Lambda^{-3}\tilde a\phi^2+\Lambda^{-3}\partial_\tau(\phi^2)-\Lambda^{-3}(\tilde b+\tilde aY)\cdot\nabla_Y\phi^2.
\]

Using \(u=\Lambda^{-1}V\), \(p=\Lambda^{-2}P\), \(\nabla_xu=\Lambda^{-2}\nabla_YV\),
\(dx\,dt=\Lambda^{5}\,dY\,d\tau\), each term in the physical inequality becomes
\[
\frac12\int |u|^2\psi\big|_{t_1}^{t_2}
=\frac12\int |V|^2\phi^2\big|_{\tau_1}^{\tau_2},
\]
\[
\nu\int |\nabla_xu|^2\psi\,dxdt
=\nu\int |\nabla V|^2\phi^2\,dY\,d\tau,
\]
\[
\frac12\int |u|^2\partial_t\psi\,dxdt
=\frac12\int |V|^2\left(\partial_\tau(\phi^2)-\tilde a\phi^2-\tilde a\,Y\cdot\nabla_Y\phi^2-\tilde b\cdot\nabla_Y\phi^2\right)\,dY\,d\tau,
\]
\[
\frac{\nu}{2}\int |u|^2\Delta_x\psi\,dxdt
=\frac{\nu}{2}\int |V|^2\Delta_Y\phi^2\,dY\,d\tau,
\]
\[
\int\left(\frac{|u|^2}{2}u+p\,u\right)\cdot\nabla_x\psi\,dxdt
=\int\left(\frac12|V|^2V+(P_{\rm loc}+H)V\right)\cdot\nabla_Y\phi^2\,dY\,d\tau.
\]
Substituting these identities and grouping terms gives the proposition's right-hand
side and coefficients \(\tilde a,\tilde b\).
\end{proof}
\subsection{Stability under subsequences}\label{paper2:sec:stability}
\begin{proposition}\label{paper2:prop:stability}
Let \((u^{(n)},p^{(n)})\) satisfy \(\mathrm{SC}(R_*,r_*,M,\eta,\kappa)\) with the same
data parameters and let \((\Lambda^{(n)},X^{(n)},V^{(n)},P^{(n)},a^{(n)},b^{(n)})\)
be obtained from retained quantitatively nondegenerate repaired-gauge selections
on the same \(I_*\) satisfying the modulation-differentiability row,
with the constants in Proposition~\ref{paper2:thm:quant-jac} uniform in \(n\). Assume also that
\((a^{(n)},b^{(n)})\) is uniformly integrable on every compact time subinterval
of \(I_*\) (in particular this holds under the \(L^\infty\) modulation bounds
used in the scale-rigid applications). After extraction \(n_j\to\infty\), there exist
subsequences and limits \((\Lambda,X,V,P,a,b)\) such that
\[
\Lambda^{(n_j)}\to\Lambda,\quad X^{(n_j)}\to X
\quad\text{uniformly on compact time subintervals},
\]
\[
V^{(n_j)}\to V\ \text{in }L^2_{\rm loc}(I_*;L^2_{\rm loc}),\quad
V^{(n_j)}\rightharpoonup V\ \text{in }L^2_{\rm loc}(I_*;H^1_{\rm loc}),
\]
\[
P^{(n_j)}\rightharpoonup P\ \text{in }L^{3/2}_{\rm loc}(Q_R^*),\quad
(a^{(n_j)},b^{(n_j)})\rightharpoonup(a,b)\ \text{in }L^1_{\rm loc}(I_*)^4.
\]
Moreover the limit chart satisfies
\[
\Lambda'=a\Lambda,\qquad X'=b\Lambda
\]
in distributions, equivalently
\[
a=\frac{\Lambda'}{\Lambda},\qquad b=\frac{X'}{\Lambda}
\]
a.e.\ on \(I_*\), and the limit obeys all conclusions below.
\end{proposition}

\begin{proof}
From \(\mathrm{SC}\) and the retained repaired-gauge selections, \((V^{(n)},P^{(n)})\) are uniformly bounded in
\(L^\infty_\tau L^2(B_{2R})\cap L^2_\tau H^1(B_{2R})\) and \(L^{3/2}_{\rm loc}(Q_R^*)\).
Theorem~\ref{paper2:thm:quant-jac} gives uniform invertibility bounds for \(J^{(n)}\), and
Theorem~\ref{paper2:thm:mod-sys} gives \((a^{(n)},b^{(n)})\) uniformly in \(L^1_{\rm loc}\) (since
\(\mathcal R^{(n)}\in L^1_{\rm loc}\)). Since on compact windows
\[
(\Lambda^{(n)})'=a^{(n)}\Lambda^{(n)},\qquad
(X^{(n)})'=b^{(n)}\Lambda^{(n)},
\]
and the charts have uniform upper and lower scale bounds, \(\Lambda^{(n)}\) and
\(X^{(n)}\) are uniformly bounded in \(W^{1,1}_{\rm loc}\). The assumed local
uniform integrability and Dunford--Pettis provide weak subsequential limits for
\((a^{(n)},b^{(n)})\). Thus, by BV compactness and one-dimensional subsequences,
\[
\Lambda^{(n_j)}\to\Lambda,\qquad X^{(n_j)}\to X
\quad\text{uniformly on compact time subintervals},
\]
and, after extraction,
\[
  a^{(n_j)}\rightharpoonup a,\qquad b^{(n_j)}\rightharpoonup b,\qquad
  a^{(n_j)}\Lambda^{(n_j)}\rightharpoonup a\Lambda,
  \qquad b^{(n_j)}\Lambda^{(n_j)}\rightharpoonup b\Lambda
\]
in \(L^1_{\rm loc}\). Since \(\Lambda^{(n_j)}\to\Lambda\) uniformly on compact
time intervals after extraction, the distributional derivative identities pass
to the limit and give
\[
\Lambda'=a\Lambda,
\qquad X'=b\Lambda.
\]
The chart bounds keep \(\Lambda\) locally bounded away from zero, so equivalently
\(a=\Lambda'/\Lambda\) and \(b=X'/\Lambda\) a.e.
From the renormalized equation,
\[
\partial_\tau V^{(n)}\in
L^1_\tau(H^{-1}_{\rm loc})+L^{4/3}_\tau(L^{12/11}_{\rm loc}),
\]
uniformly on compact windows. Aubin--Lions compactness on bounded balls gives
\[
V^{(n_j)}\to V\ \text{in }L^2_{\rm loc}(I_*;L^2_{\rm loc}),\quad
V^{(n_j)}\rightharpoonup V\ \text{in }L^2_{\rm loc}(I_*;H^1_{\rm loc}).
\]
By Calderón--Zygmund and Theorem~\ref{paper2:thm:pressure-decomp}, \(P^{(n_j)}\rightharpoonup P\) in
\(L^{3/2}_{\rm loc}(Q_R^*)\), and with normalized pressure split \(P=P_{\rm loc}+H\),
the same for \(P_{\rm loc}\) and \(H\). Passing to limits in
\(\mathcal F(1,0;V^{(n)}(\tau))=0\), Theorem~\ref{paper2:thm:ren-NS}, and the linear modulation system
yields the corresponding limit constraints and equations.

\end{proof}

\subsection{Final representation theorem}\label{paper2:sec:final}

\begin{theorem}[Chart and repaired-gauge representation]\label{paper2:thm:A}
Assume \(\mathrm{SC}(R_*,r_*,M,\eta,\kappa)\), and suppose the branch lies in the
retained repaired-gauge regime with selection \((\lambda,c)\) on
\(J=[\tau_0-\ell,\tau_0+\ell]\subset I_*\). Then there are fields \((V,P)\) on \(Q_{2R}^J\), with
\((\Lambda,X)\in AC(J)\),
\[
a=\frac{\Lambda'}{\Lambda},\qquad b=\frac{X'}{\Lambda},
\qquad a,b\in L^1_{\rm loc}(J),
\]
such that
\[
u(x,t)=\Lambda(\tau)^{-1}V\!\left(\frac{x-X(\tau)}{\Lambda(\tau)},\tau\right),\quad
\frac{dt}{d\tau}=\Lambda(\tau)^2,
\]
and \(\mathcal F(1,0;V(\tau))=0\) for \(\tau\in J\).
\end{theorem}

\begin{proof}
From \textup{(SC5)} and Proposition~\ref{paper2:prop:prelim-frame}, \((\Lambda_{\rm pre},X_{\rm pre})\)
are AC on \(I_*\), so \(V_{\rm pre}\) is defined on \(Q_{2R}^*\).
The retained repaired-gauge regime supplies
\((\lambda,c)\in AC(J;(0,\infty)\times\mathbb R^3)\) and
\(\mathcal F(1,0;V_{\rm pre}^{\lambda,c})=0\) a.e.
By Definition~\ref{paper2:def:repaired-chart} and Theorem~\ref{paper2:thm:repaired-existence}, setting
\[
\Lambda(\tau)=\Lambda_{\rm pre}(\tau)\lambda(\tau),\qquad
X(\tau)=X_{\rm pre}(\tau)+\Lambda_{\rm pre}(\tau)c(\tau),\qquad
V(\tau)=V_{\rm pre}^{\lambda(\tau),c(\tau)},
\]
and \(P(\tau)=P_{\rm pre}^{\lambda(\tau),c(\tau)}\),
gives \( \Lambda,X\in AC(J)\) and
\[
a=\frac{\Lambda'}{\Lambda},\qquad b=\frac{X'}{\Lambda},
\qquad a,b\in L^1_{\rm loc}(J),
\]
Moreover,
\[
\mathcal F(1,0;V(\tau))=0\ \text{for a.e. }\tau\in J.
\]
Finally \(u=\Lambda^{-1}V(\cdot,\tau)\circ\Phi^{-1}\) by definition of the repaired map, hence the
physical identity with \(\frac{dt}{d\tau}=\Lambda^2\) follows.
\end{proof}

\begin{theorem}[Pressure and modulation]\label{paper2:thm:B}
Assume, in addition to \autoref{paper2:thm:A}, that the retained repaired-gauge
selection is quantitatively nondegenerate in the sense of
Proposition~\ref{paper2:thm:quant-jac} and satisfies the
modulation-differentiability row of \cref{paper2:def:mod-diff-row}.
Then \((V,P)\) satisfies the renormalized Navier--Stokes
equation with the logarithmic coefficient convention of Theorem~\ref{paper2:thm:ren-NS}; there is a
decomposition \(P=P_{\rm loc}+H\) on \(Q_R^*\) as in
Theorem~\ref{paper2:thm:pressure-decomp}; and \((a,b)\) is the unique solution of
\[
J(\tau)\binom{a(\tau)}{b(\tau)}=\mathcal R(\tau),\qquad
\binom{a}{b}\in L^1_{\rm loc}(J)^4,
\]
with bounds \(|(a,b)|\le K|\mathcal R|\), where \(K\) depends on
\(R,p,\Theta_0,m_0,C_u,\varepsilon_0,\chi_R\).
\end{theorem}

\begin{proof}
Equation and coefficient convention are inherited directly from Theorem~\ref{paper2:thm:ren-NS} applied to
\((\Lambda,X)\) from Theorem~\ref{paper2:thm:A} and the definition of \(V\).
By Theorem~\ref{paper2:thm:pressure-decomp} applied at each time slice, the decomposition on \(Q_R^J\)
is defined and unique once \(\fint_{B_R}H=0\) is imposed.
The modulation identity is Theorem~\ref{paper2:thm:mod-sys}; uniqueness follows from invertibility of
\(J\) (Theorem~\ref{paper2:thm:quant-jac}) and the pointwise identity \((a,b)=J^{-1}\mathcal R\).
Uniform bounds on \((a,b)\) are obtained by the bound on \(J^{-1}\) there and the stated bound on
\(\mathcal R\).
\end{proof}

\begin{theorem}[Compact-window representation theorem]\label{paper2:thm:C}
Under the assumptions of \cref{paper2:thm:A,paper2:thm:B}, for every compact
window \(J\subset I_*\) and each fixed \(R>0\), setting
\(Q_R^J:=B_R\times J\),
one has:
\begin{enumerate}
\item the exact renormalized equation on \(Q_{2R}^J\),
\item the four repaired-gauge constraints \(\mathcal F(1,0;V(\tau))=0\),
\item bounds on \(a,b\in L^1_{\rm loc}(J)\) (and \(L^\infty_{\rm loc}\) under
  \(\mathcal R\in L^\infty_{\rm loc}\)),
\item local pressure decomposition \(P=P_{\rm loc}+H\) with \(L^{3/2}\) local bounds and
  harmonic remainder controls,
	\item renormalized local suitable energy inequality on \(Q_R^J\),
	\item stability under subsequences in the topologies above for families whose
        modulation coefficients are locally uniformly integrable, in particular
        in the \(L^\infty\)-modulated retained rows used below.
\end{enumerate}
The later reductions use only these conclusions.
\end{theorem}

\begin{proof}
From Theorem~\ref{paper2:thm:A} and the choice of \(J\subset I_*\), all quantities are defined on \(Q_{2R}^J\).
Items (1) and (2) are Theorems~\ref{paper2:thm:ren-NS} and~\ref{paper2:thm:repaired-existence};
item (3) is Theorem~\ref{paper2:thm:mod-sys} together with the inverse bound from
Proposition~\ref{paper2:thm:quant-jac}; item (4) is Theorem~\ref{paper2:thm:pressure-decomp};
item (5) is Proposition~\ref{paper2:prop:ren-suitable}; item (6) is
Proposition~\ref{paper2:prop:stability} under the stated local uniform
integrability condition. Combining these results gives the
stated compact-window representation.
\end{proof}

\begin{remark}
The preceding results separate two alternatives:
\begin{enumerate}
\item the nondegenerate single-core repaired-gauge regime, where the representation theorems apply;
\item gauge degeneracy (\(\det J=0\)) or compact-structure failure, which belongs
to the multibubble, cascade, or scale-rigidity alternatives.
\end{enumerate}
\end{remark}

\clearpage
\section{Multibubble and Cascade Exclusion}\label{sec:multibubble-cascade-exclusion}
\medskip
\noindent\textbf{Analytic inputs.}

The local reduction theorem uses the following analytic inputs from the
preceding sections, with profile decompositions and decoupling estimates in
their standard critical-space sense \cite{Gallagher2001,BrezisLieb1983}.  The
remaining work is the multibubble, cascade, and terminal-frame decoupling.
\begin{enumerate}[label=\textup{(\roman*)}]
\item From \cref{sec:single-core-criterion}, the compact single-core criterion \cref{paper1:thm:single-core-criterion}: a
compact single-core finite-cost Type II branch with positive local CKN
concentration, nondegenerate scale-center selection, pressure reconstruction,
bounded compact-cylinder suitable estimates, bounded modulation, and vanishing
localized dissipation is excluded.
\item From \cref{sec:repaired-gauge-construction}, the renormalized equation \cref{paper2:thm:ren-NS}, pressure decomposition \cref{paper2:thm:pressure-decomp}, modulation system \cref{paper2:thm:mod-sys}, and repaired-gauge representation \cref{paper2:thm:A,paper2:thm:B,paper2:thm:C}: the renormalized equation,
pressure decomposition modulo functions of time, bounded modulation
coefficients, renormalized suitable structure, and stability of these structures
under subsequences.
\item From \cref{sec:rough-core-exclusion}, the compact-cylinder Caccioppoli estimate \cref{paper4:thm:caccioppoli} and rough-core exclusion \cref{paper4:cor:rough-core}.
\item From \cref{sec:scale-collapse-cost}, the collapse cost exclusion \cref{paper5:thm:collapse-core-floor,paper5:thm:cost-exclusion} and autonomous reduced-limit theorem \cref{paper5:thm:autonomous-limit} for scale-rigid or scale-collapsing branches.
\end{enumerate}

\subsection{Local setting and active frames}

\medskip
\noindent\textit{Normalized local Navier--Stokes setting.}
All quantities are written in rescaled local variables \((x,t)\in Q_R=B_R\times(-R^2,0)\)
with viscosity \(\nu=1\) (critical scaling has absorbed \(\nu\)).  A local pair
\((u_n,p_n)\) on such a cylinder satisfies
\[
  \partial_tu_n-\Delta u_n+(u_n\cdot\nabla)u_n+\nabla p_n=0,\qquad \nabla\cdot u_n=0,
\]
and pressure is normalised through spatial averaging
\[
  (p_n)_{B_\rho}(t):=|B_\rho|^{-1}\int_{B_\rho}p_n(x,t)\,dx
\]
to factor out time-dependent constants.

For $r>0$, $x_0\in\mathbb R^3$, and $t_0\in\mathbb R$, let
\[
  Q_r(x_0,t_0):=B_r(x_0)\times(t_0-r^2,t_0),
  \qquad
  Q_r:=Q_r(0,0)=B_r\times(-r^2,0).
\]

Let $(u_n,p_n)$ be parabolic rescalings around a potential singular point with
positive local Caffarelli--Kohn--Nirenberg concentration on compact cylinders.
Concretely, for a base triple \((x_n,t_n,\lambda_n)\),
\[
  u_n(y,s)=\lambda_n u(x_n+\lambda_n y,\ t_n+\lambda_n^2 s),\qquad
  p_n(y,s)=\lambda_n^2 p(x_n+\lambda_n y,\ t_n+\lambda_n^2 s).
\]
For each $n$ and $R>0$, define the scaling-invariant local quantities
\[
  A_n(R):=R^{-1}\operatorname*{ess\,sup}_{-R^2<t<0}\int_{B_R}|u_n(x,t)|^2\,dx,
\]
\[
  E_n(R):=R^{-1}\int_{Q_R}|\nabla u_n(x,t)|^2\,dx\,dt,
\]
\[
  C_n(R):=R^{-2}\int_{Q_R}|u_n(x,t)|^3\,dx\,dt,
  \qquad
  D_n(R):=R^{-2}\int_{-R^2}^0\int_{B_R}|p_n-(p_n)_{B_R}(t)|^{3/2}\,dx\,dt,
\]
where
\[
  (p_n)_{B_R}(t):=|B_R|^{-1}\int_{B_R}p_n(x,t)\,dx.
\]

For an active frame \((x_n,\lambda_n)\), \(\lambda_n>0\), define the critical rescaling
\[
  \mathcal T_{x_n,\lambda_n}f(y):=\lambda_n f(x_n+\lambda_n y).
\]
If \(\phi\in L^3_\sigma(\mathbb R^3)\) is observed in this frame, its contribution is
\[
  \phi_n(y)=\mathcal T_{x_n,\lambda_n}\phi(y)=\lambda_n\phi(x_n+\lambda_n y).
\]
The critical Jacobian identity is exact:
\[
  \|\mathcal T_{x_n,\lambda_n}f\|_{L^3(E)}^3
  =\int_{x_n+\lambda_nE}|f(x)|^3\,dx
\]
for every measurable \(E\subset\mathbb R^3\).  In particular
\[
  \|\mathcal T_{x_n,\lambda_n}f\|_{L^3(\mathbb R^3)}^3
  =\|f\|_{L^3(\mathbb R^3)}^3,
\]
so \(\mathcal T_{x_n,\lambda_n}\) is an isometry on \(L^3(\mathbb R^3)\).

\begin{definition}[Bounded selected-window limit]
A selected-window subsequence has a bounded selected-window limit on a compact
cylinder $Q_R$ if there exists $u_\infty\in L^\infty(Q_R)$ such that, after
extracting a subsequence,
\[
  \forall\varepsilon>0\ \exists N:\ \forall n\ge N,\ \|u_n-u_\infty\|_{L^3(Q_R)}<\varepsilon.
\]
The limit is nonzero if $u_\infty\not\equiv0$ on $Q_R$.
\end{definition}

\begin{definition}[Local Type II sequence]
A positive local concentration sequence is called a local Type II sequence if it
satisfies \cref{paper1:def:typeII}; equivalently, no selected-window subsequence
converges strongly in \(L^3\) on any compact cylinder \(Q_R\) to a nonzero bounded
selected-window limit.
\end{definition}

\begin{definition}[Active frame]\label{paper3:def:active-frame}
Fix a concentrating time sequence \(t_n\uparrow T\) and a compact time interval
\(I\Subset(-\infty,0]\).  An active frame is a tuple
\[
  \mathfrak F^j=\bigl(x^j_n,\lambda^j_n,I,U^j_n,P^j_n\bigr),
  \qquad \lambda^j_n\downarrow0,
\]
where \(x^j_n\in\mathbb R^3\) and
\[
  U^j_n(y,s):=\lambda^j_n u(x^j_n+\lambda^j_n y,t_n+(\lambda^j_n)^2s),
  \qquad
  P^j_n(y,s):=(\lambda^j_n)^2p(x^j_n+\lambda^j_n y,t_n+(\lambda^j_n)^2s)
\]
on \(B_{2R}\times I\) for every fixed compact cylinder under consideration.
The frame is active at threshold \(\eta_*>0\) if for some fixed \(R_0<\infty\)
\[
  \liminf_{n\to\infty}\int_{Q_{R_0}}|U^j_n|^3\,dy\,ds\ge \eta_*^3.
\]
\end{definition}

\begin{definition}[Relations between active frames]\label{paper3:def:frame-relations}
Let \(\mathfrak F^i\) and \(\mathfrak F^j\) be two active frames.  After passing
to a subsequence exactly one of the following parameter regimes holds.
\begin{enumerate}[label=\textup{(\roman*)}]
\item They are \emph{same-point comparable-scale} if
\[
  0<c\le\frac{\lambda^i_n}{\lambda^j_n}\le C<\infty,\qquad
  \frac{|x^i_n-x^j_n|}{\max(\lambda^i_n,\lambda^j_n)}\le C
\]
for all large \(n\).
\item They form a \emph{same-point cascade} if, after relabeling,
\[
  \frac{\lambda^i_n}{\lambda^j_n}\to0,\qquad
  \frac{|x^i_n-x^j_n|}{\lambda^j_n}=O(1).
\]
The frame \(i\) is then the inner frame and \(j\) is an outer frame.
\item They are \emph{separated-point} if
\[
  \frac{|x^i_n-x^j_n|}{\lambda^i_n}\to\infty
  \quad\text{and}\quad
  \frac{|x^i_n-x^j_n|}{\lambda^j_n}\to\infty .
\]
\end{enumerate}
\end{definition}

\begin{definition}[Terminal frame and perturbative remainder]\label{paper3:def:terminal-frame}
Given a finite active family, a terminal frame is chosen by selecting an active
scale with minimal order:
\[
  \lambda^{j_*}_n\le C_j\lambda^j_n
  \qquad\text{for every active }j
\]
after passing to a subsequence.  Comparable frames at the same physical point
are grouped before this choice; if several grouped frames remain tied, one is
chosen by a fixed lexicographic ordering of their labels.  A remainder \(r_n\)
in a terminal frame is perturbative on \(B_{2R}\times I\) if
\[
  \|r_n\|_{L^\infty_I L^3(B_{2R})}\to0
\]
and its equation residual and pressure satisfy
\[
  \|f_n\|_{L^1_IH^{-1}(B_R)}+\|\nabla\pi^r_n\|_{L^1_IH^{-1}(B_R)}\to0
\]
for every compact \(B_R\times I\).
\end{definition}

\begin{definition}[Local multibubble/cascade branch]\label{paper3:def:local-multibubble}
A local multibubble/cascade branch is a positive-concentration local Type II
branch for which at least one of the following alternatives occurs:
\begin{enumerate}[label=\textup{(\roman*)}]
\item the compact-cylinder suitable estimates $A_n+E_n+C_n+D_n$ are unbounded on some
  fixed rescaled cylinder;
\item positive CKN mass splits into two or more comparable active cores;
\item a strict same-point scale cascade appears inside bounded rescaled cylinders;
\item separated physical points remain active in different local rescaled frames;
\item the localized scale or centering selection degenerates (no unique core
  selected).
\end{enumerate}
\end{definition}

\begin{definition}[Active profile family]\label{paper3:def:active-family}
An active profile family is a family
\[
  \{(\phi^j,x_n^j,\lambda_n^j)\}_{j\in J},
  \qquad \phi^j\in L^3_\sigma(\mathbb R^3),
\]
with finite or countable index set $J$, extracted along a positive local
concentration sequence.  In physical variables write
\[
  \Lambda_{\lambda,x_0}\phi(x):=\lambda^{-1}\phi\left(\frac{x-x_0}{\lambda}\right).
\]
For every finite \(J_0\subset J\) the family gives an expansion
\[
  u(\cdot,t_n)=\sum_{j\in J_0}\Lambda_{\lambda^j_n,x^j_n}\phi^j+r^{J_0}_n
\]
on the compact terminal windows under consideration.  After the finite active
subfamily has been selected, the corresponding remainder satisfies the
perturbative bounds of \cref{paper3:def:terminal-frame}.  The parameters are
considered only after grouping same-point comparable-scale profiles into
compound profiles as in \cref{paper3:def:compound-profile}.  A profile is active if
\[
  \|\phi^j\|_{L^3(\mathbb R^3)}\ge\eta_*>0.
\]
The family has finite critical mass when
\[
  \sum_{j\in J}\|\phi^j\|_{L^3(\mathbb R^3)}^3\le M_*<\infty .
\]
It captures positive local concentration if every compact terminal-window
sequence with
\[
  \liminf_{n\to\infty}\int_{Q_R}|u_n|^3>0
\]
either contains an active profile observed in that window or is absorbed into
the selected terminal frame; equivalently, no positive terminal-window
concentration remains in the perturbative remainder.
\end{definition}

\begin{definition}[Compound profile]\label{paper3:def:compound-profile}
Let \(J_k\subset J\) be a same-point comparable-scale class.  Choose one
representative \((x_n^{j_k},\lambda_n^{j_k})\).  After passing to a subsequence,
\[
  \rho_j:=\lim_{n\to\infty}\frac{\lambda^j_n}{\lambda^{j_k}_n}\in(0,\infty),
  \qquad
  z_j:=\lim_{n\to\infty}\frac{x^j_n-x^{j_k}_n}{\lambda^{j_k}_n}\in\mathbb R^3 .
\]
The compound profile in the representative frame is the \(L^3_\sigma\)-sum
\[
  \Phi^k(y):=\sum_{j\in J_k}\rho_j^{-1}
  \phi^j\left(\frac{y-z_j}{\rho_j}\right),
\]
whenever \(J_k\) is finite; finiteness follows for active classes from
\cref{paper3:lem:finite-active-count}.  If \(\|\Phi^k\|_{L^3}<\varepsilon_\star\), where
\(\varepsilon_\star\) is the critical small-data threshold, the class is
perturbative.  Otherwise \(\Phi^k\) is treated as one active object in its
representative frame.
\end{definition}

\begin{lemma}[Finite active profile count]\label{paper3:lem:finite-active-count}
Every active profile family with finite critical mass contains only finitely many
active profiles.  More precisely,
\[
  \#J\le M_*\eta_*^{-3}.
\]
\end{lemma}

\begin{proof}
For each active profile $j\in J$,
\[
  \|\phi^j\|_{L^3(\mathbb R^3)}^3\ge\eta_*^3.
\]
Let $J_0\subset J$ be finite. Summing over $J_0$ gives
\[
  \#J_0\,\eta_*^3
  \le
  \sum_{j\in J_0}\|\phi^j\|_{L^3(\mathbb R^3)}^3
  \le
  \sum_{j\in J}\|\phi^j\|_{L^3(\mathbb R^3)}^3
  \le M_*.
\]
Hence $\#J_0\le M_*\eta_*^{-3}$.  Taking the supremum over all finite
subsets $J_0$ yields the same bound for the cardinality of $J$, so $J$ is finite
and
\[
  \#J\le M_*\eta_*^{-3}.
\]
\end{proof}

\begin{lemma}[Exhaustive parameter partition]\label{paper3:lem:parameter-partition}
Let \(\mathfrak F^i,\mathfrak F^j\) be two active frames.  After passing to a
subsequence, either the pair is same-point comparable-scale, same-point cascade,
or separated-point in the sense of \cref{paper3:def:frame-relations}; if none of these
relations holds in a selected compact frame, then one of the two contributions
is locally invisible there in \(L^3\).
\end{lemma}

\begin{proof}
Pass to a subsequence so that
\[
  \frac{\lambda^i_n}{\lambda^j_n}\to \ell\in[0,\infty]
\]
in the extended half-line.  If \(0<\ell<\infty\), then either
\[
  \frac{|x^i_n-x^j_n|}{\max(\lambda^i_n,\lambda^j_n)}
\]
is bounded along a further subsequence, giving the same-point comparable-scale
case, or it tends to infinity, giving separated points because the two scales
are comparable.
\par
If \(\ell=0\), then either
\[
  \frac{|x^i_n-x^j_n|}{\lambda^j_n}
\]
is bounded, giving a same-point cascade with \(i\) inner and \(j\) outer, or it
tends to infinity.  In the latter case the \(j\)-profile is locally invisible in
the \(i\)-frame by \cref{paper3:lem:separated-vanish}.  The case \(\ell=\infty\) is the
same after interchanging \(i\) and \(j\).  These alternatives exhaust all
subsequential limits of the scale ratio and normalized center distance.
\end{proof}

\begin{theorem}[Minimal bad-configuration extraction]\label{paper3:thm:minimal-bad-extraction}
Let a positive local Type II branch admit an active profile family with finite
critical mass, perturbative remainder in the sense of
\cref{paper3:def:terminal-frame}, and concentration capture in the sense of
\cref{paper3:def:active-family}.  If the branch is not already a single active
terminal frame satisfying the single-core hypotheses, then after passing to a
subsequence one of the following occurs:
\begin{enumerate}[label=\textup{(\roman*)}]
\item compact-cylinder suitable control fails, and
\cref{paper3:lem:compactness-failure} produces a bounded-cylinder concentration core;
\item there are at least two active frames in a same-point comparable-scale
class, hence a compound profile in the sense of \cref{paper3:def:compound-profile};
\item there is a same-point cascade;
\item there are separated active physical points;
\item the scale-center selection of the selected compound or terminal frame is
degenerate.
\end{enumerate}
In every case the positive local concentration is carried by an active frame or
compound profile, not by the perturbative remainder.
\end{theorem}

\begin{proof}
If compact-cylinder suitable control fails, item \((i)\) follows from
\cref{paper3:lem:compactness-failure}.  Assume therefore that the compact-cylinder
family is bounded on the windows under consideration.
\par
By concentration capture, positive local concentration cannot remain solely in
the perturbative remainder.  Hence either one selected active frame carries the
terminal-window concentration, or another active frame remains non-negligible.
If there is one nondegenerate selected active frame and no competing active
object, the branch is in the single-core regime.  Since this case is
excluded by hypothesis, either the selected scale-center map is degenerate,
giving item \((v)\), or at least two active frames remain.
\par
For two active frames, apply \cref{paper3:lem:parameter-partition}.  Same-point
comparable-scale frames give item \((ii)\), same-point scale separation gives
item \((iii)\), and separated points give item \((iv)\).  The remaining
alternative in \cref{paper3:lem:parameter-partition} is local invisibility of one
contribution in the selected compact frame; such a contribution is absorbed
into the perturbative remainder and cannot carry the retained positive local
concentration by the capture property.  This gives the extraction statement.
\end{proof}

\subsection{Elementary local decoupling estimates}

\begin{lemma}[Critical rescalings that escape are locally invisible]\label{paper3:lem:escaping-l3}
Let \(\phi\in L^3(\mathbb R^3)\), let \(R<\infty\), and define
\[
  W_n(y)=\rho_n\phi(z_n+\rho_n y),
\]
where \(\rho_n>0\).  Assume that the sets
\[
  z_n+\rho_nB_R
\]
either have measure tending to zero, or escape to spatial infinity in the sense
that for every compact set \(K\subset\mathbb R^3\),
\[
  (z_n+\rho_nB_R)\cap K=\varnothing
\]
for all sufficiently large \(n\).  Then
\[
  \|W_n\|_{L^3(B_R)}\to0.
\]
\end{lemma}

\begin{proof}
By the critical change of variables,
\[
  \|W_n\|_{L^3(B_R)}^3
  =\int_{B_R}\rho_n^3|\phi(z_n+\rho_n y)|^3\,dy
  =\int_{z_n+\rho_n B_R}|\phi(x)|^3\,dx.
\]
Set
\[
  m_n:=\int_{z_n+\rho_n B_R}|\phi(x)|^3\,dx.
\]
Since \(\|W_n\|_{L^3(B_R)}^3=m_n\), it is enough to show \(m_n\to0\).
\par
Fix \(\varepsilon>0\).  We split into the two structural alternatives.

\smallskip
\noindent\textit{Case 1: shrinking measure.}
If \(|z_n+\rho_nB_R|\to0\), then by absolute continuity of the integral
\(|\phi|^3\in L^1(\mathbb R^3)\) there exists \(\delta>0\) such that
\[
  |E|<\delta\Longrightarrow \int_E|\phi|^3\,dx<\varepsilon.
\]
Choose \(N\) with \(|z_n+\rho_nB_R|<\delta\) for all \(n\ge N\).  Then
\[
  m_n=\int_{z_n+\rho_nB_R}|\phi|^3\,dx<\varepsilon,\qquad n\ge N.
\]

\smallskip
\noindent\textit{Case 2: escape to infinity.}
Assume now that for every compact set \(K\subset\mathbb R^3\),
\((z_n+\rho_nB_R)\cap K=\varnothing\) for all sufficiently large \(n\).  By
absolute continuity at infinity, choose compact \(K\subset\mathbb R^3\) such that
\[
  \int_{\mathbb R^3\setminus K}|\phi|^3<\varepsilon.
\]
For every such compact \(K\), there is \(N(K)\) with
\((z_n+\rho_nB_R)\cap K=\varnothing\) for all \(n\ge N(K)\).  For these
indices,
\[
  z_n+\rho_nB_R \subset \mathbb R^3\setminus K,
\]
and therefore, for \(n\ge N(K)\),
\[
  m_n=\int_{z_n+\rho_nB_R}|\phi|^3\,dx
  \le
  \int_{\mathbb R^3\setminus K}|\phi|^3
  <\varepsilon.
\]
In either case, for every \(\varepsilon>0\), \(m_n<\varepsilon\) for \(n\) sufficiently
large.  Hence \(m_n\to0\), and therefore
\[
  \|W_n\|_{L^3(B_R)}\to0.
\]
\end{proof}

\begin{lemma}[Separated active profiles vanish in the selected frame]\label{paper3:lem:separated-vanish}
Let \((x^p_n,\lambda^p_n)\) and \((x^q_n,\lambda^q_n)\) be two active frames
with separated physical centers in the sense that
\[
  \frac{|x^q_n-x^p_n|}{\lambda^p_n}\to\infty.
\]
If \(\phi^q\in L^3(\mathbb R^3)\), then the \(q\)-profile observed in the
\(p\)-frame converges to zero locally in \(L^3\): for every \(R<\infty\),
\[
  \left\|
  \frac{\lambda^p_n}{\lambda^q_n}
  \phi^q\left(\frac{x^p_n-x^q_n}{\lambda^q_n}
       +\frac{\lambda^p_n}{\lambda^q_n}y\right)
  \right\|_{L^3(B_R)}\to0.
\]
\end{lemma}

\begin{proof}
The displayed function has the form in \cref{paper3:lem:escaping-l3} with
\[
  z_n=\frac{x^p_n-x^q_n}{\lambda^q_n},
  \qquad
  \rho_n=\frac{\lambda^p_n}{\lambda^q_n}.
\]
The corresponding physical set in the \(q\)-profile variables is
\[
  z_n+\rho_nB_R
  =\frac{x^p_n-x^q_n}{\lambda^q_n}
   +\frac{\lambda^p_n}{\lambda^q_n}B_R.
\]
For every \(y\in B_R\),
\[
  |z_n+\rho_n y|
  \ge |z_n|-\rho_nR.
\]
Set \(E_n:=z_n+\rho_nB_R\).  Then
\[
  \operatorname{dist}(E_n,0)
  =|z_n|-\rho_nR
  =\frac{\lambda^p_n}{\lambda^q_n}\left(\frac{|x^q_n-x^p_n|}{\lambda^p_n}-R\right).
\]
Since \(\frac{|x^q_n-x^p_n|}{\lambda^p_n}\to\infty\), there exists \(N_0(R)\) such that
\[
  \frac{|x^q_n-x^p_n|}{\lambda^p_n}-R\ge1,\quad n\ge N_0(R),
\]
and hence for such \(n\), \(\operatorname{dist}(E_n,0)\ge\rho_n\).
\par
There are two cases.
\begin{enumerate}[label=\textup{(\alph*)}]
\item If \(\rho_n\to0\), then
 \[
   |z_n+\rho_nB_R|=\rho_n^3|B_R|\to0.
 \]
 Then the first alternative in \cref{paper3:lem:escaping-l3} applies.
\item Otherwise, after passing to a subsequence, \(\rho_n\ge\rho_\ast>0\).  Then
 \[
   \operatorname{dist}(z_n+\rho_nB_R,0)\ge \rho_\ast\left(\frac{|x^q_n-x^p_n|}{\lambda^p_n}-R\right)\to\infty,
 \]
 so the sets escape every fixed compact subset of \(\mathbb R^3\).  The
 second alternative in \cref{paper3:lem:escaping-l3} applies.
\end{enumerate}
In both subcases the \(L^3(B_R)\)-norm of the displayed expression converges to
zero.
\end{proof}

\begin{lemma}[Strict outer scales vanish in the innermost frame]\label{paper3:lem:outer-vanish}
Let several active profiles have the same physical center and scales
\(\lambda^1_n,\ldots,\lambda^J_n\).  Suppose \(\lambda^1_n\) is an innermost
scale and
\[
  \frac{\lambda^1_n}{\lambda^j_n}\to0
  \qquad (j\ne1).
\]
Then every outer profile \(\phi^j\in L^3(\mathbb R^3)\), viewed in the
\(\lambda^1_n\)-frame, converges to zero in \(L^3(B_R)\) for each fixed
\(R<\infty\).  Any constant ambient velocity arising from the regular part of
the outer flow is absorbed into the translation modulation.
\end{lemma}

\begin{proof}
In the innermost variables, the \(j\)-th outer contribution has the form
\[
  W^j_n(y)=\frac{\lambda^1_n}{\lambda^j_n}
  \phi^j\left(\frac{x^1_n-x^j_n}{\lambda^j_n}
  +\frac{\lambda^1_n}{\lambda^j_n}y\right).
\]
Set
\[
  \rho_n:=\frac{\lambda^1_n}{\lambda^j_n}\to0,
  \qquad
  z_n^j:=\frac{x^1_n-x^j_n}{\lambda^j_n}.
\]
Then
\[
  W^j_n(y)=\rho_n\phi^j\left(z_n^j+\rho_n y\right),
\]
and by the critical change of variables,
\[
  \|W^j_n\|_{L^3(B_R)}^3
  =\int_{z_n^j+\rho_nB_R}|\phi^j(x)|^3\,dx.
\]
Since \(\rho_n\to0\), the translated domains have
\[
  |z_n^j+\rho_nB_R|=\rho_n^3|B_R|\to0.
\]
By absolute continuity of the integral for \(\phi^j\in L^3(\mathbb R^3)\),
\[
  \|W^j_n\|_{L^3(B_R)}\to0.
\]
To treat a general \(\phi^j\), let \(\phi^j_\varepsilon\in C_c^\infty(\mathbb R^3)\) with
\(\|\phi^j-\phi^j_\varepsilon\|_{L^3}<\varepsilon\).  By critical isometry,
\[
  \|W^j_n-(W^j_n)_\varepsilon\|_{L^3(B_R)}
  \le
  \|\phi^j-\phi^j_\varepsilon\|_{L^3} < \varepsilon,
\]
where \((W^j_n)_\varepsilon\) is the same scaling of \(\phi^j_\varepsilon\).
The compactly supported term has vanishing local norm by the shrinking-domain
argument above, and then letting first \(n\to\infty\), then \(\varepsilon\downarrow0\)
implies the claimed convergence.
A constant vector field from the regular part of an outer flow contributes only
a Galilean drift in the selected frame; it is therefore incorporated into the
translation modulation and does not remain as an exterior perturbation.
\end{proof}

\begin{lemma}[Compactness failure produces another core]\label{paper3:lem:compactness-failure}
If the compact-cylinder suitable estimates are unbounded on a fixed rescaled
cylinder, then after passing to a subsequence the local \(C_n+D_n\) quantity is
unbounded on a comparable cylinder.  Consequently the branch contains an
additional local CKN concentration core.
\end{lemma}

\begin{proof}
Fix \(r>0\).  Apply the local Caccioppoli inequality with outer radius \(2r\).
For each index \(n\), suitable weak solutions satisfy, with \(C_r\) independent
of \(n\),
\[
  A_n(r)+E_n(r)
  \le C_r\bigl(1+C_n(2r)+C_n(2r)^{2/3}+D_n(2r)\bigr),
\]
where the pressure is normalized by subtracting spatial means
\((p_n)_{B_{2r}}(t)\).
\par
Suppose that \(A_n(r)+E_n(r)+C_n(r)+D_n(r)\) is unbounded along a subsequence.
If \(C_n(2r)+D_n(2r)\) were bounded on this subsequence, the preceding estimate
would give uniform boundedness of \(A_n(r)+E_n(r)\).  Monotonicity of the local
integrals also gives \(C_n(r)+D_n(r)\le C_n(2r)+D_n(2r)\), contradicting
unboundedness of the full family on \(Q_r\).  Therefore, after extracting a
subsequence,
\[
  C_n(2r)+D_n(2r)\to\infty .
\]
\par
Define the localized CKN concentration at variable center and radius
\[
  G_n(x,t,\rho):=\rho^{-2}\int_{Q_\rho(x,t)}|u_n|^3\,dx\,ds
  +\rho^{-2}\int_{Q_\rho(x,t)}|p_n-(p_n)_{B_\rho}(s)|^{3/2}\,dx\,ds.
\]
Set
\[
  \mathbf M_n:=\sup\{G_n(x,t,\rho):0<\rho\le 2r,\ Q_\rho(x,t)\subset Q_{4r}\}.
\]
The choice \((x,t,\rho)=(0,0,2r)\) is admissible because \(Q_{2r}(0,0)\subset Q_{4r}\), so
\[
  \mathbf M_n\ge C_n(2r)+D_n(2r)\to\infty.
\]
After extracting a further subsequence, we may assume \(\mathbf M_n\ge1\) for all \(n\).
Choose \((x_n,t_n,\rho_n)\) with \(0<\rho_n\le2r\), \(Q_{\rho_n}(x_n,t_n)\subset Q_{4r}\),
such that
\[
  G_n(x_n,t_n,\rho_n)\ge\tfrac12\,\mathbf M_n\ge\tfrac12.
\]
Rescale around \((x_n,t_n,\rho_n)\):
\[
  v_n(y,s):=\rho_nu_n(x_n+\rho_n y,t_n+\rho_n^2 s),\qquad
  q_n(y,s):=\rho_n^2p_n(x_n+\rho_n y,t_n+\rho_n^2 s).
\]
Using the critical Jacobian identity,
\[
  \int_{Q_1}|v_n|^3\,dy\,ds
  +\int_{Q_1}|q_n-(q_n)_{B_1}(s)|^{3/2}\,dy\,ds
  =G_n(x_n,t_n,\rho_n)\ge\tfrac12.
\]
Hence this new frame carries positive local critical mass, i.e. a new local CKN
concentration core, as required.
\end{proof}

\subsection{Terminal active-frame decoupling}

\begin{definition}[Admissible terminal active family]\label{paper3:def:admissible-terminal}
Fix a compact cylinder \(B_{2R}\times I\) and a selected active frame.  A
terminal active family is admissible on this cylinder if the following hold.
\begin{enumerate}[label=\textup{(\roman*)}]
\item The active profiles satisfy the finite-mass and active-mass conditions of
\cref{paper3:def:active-family}.
\item Every nonselected active profile is either separated from the selected
physical point in the sense of \cref{paper3:lem:separated-vanish}, or has the same
physical point and a strictly outer scale in the sense of
\cref{paper3:lem:outer-vanish}, after comparable same-point scales have been grouped.
\item The rescaled remainder \(r_n\) satisfies
\[
  \|r_n\|_{L^\infty_I L^3(B_{2R})}\to0,
\]
and there exist fields \((\pi_n^r,f_n)\) such that on \(B_{2R}\times I\),
\[
  \partial_t r_n-\Delta r_n+a_n\bigl(r_n+y\cdot\nabla r_n\bigr)+b_n\cdot\nabla r_n
  +\nabla\pi_n^r=f_n,\qquad
  \nabla\cdot r_n=0,
\]
with
\[
  \|f_n\|_{L^1_IH^{-1}(B_R)}\to0,
\]
\[
  \|\nabla\pi_n^r\|_{L^1_IH^{-1}(B_R)}\to0.
\]
\item The selected velocity satisfies
\[
  \sup_n\|U_n\|_{L^\infty_I L^3(B_{2R})}<\infty .
\]
\item The selected modulation coefficients satisfy
\[
  \sup_n\bigl(\|a_n\|_{L^\infty(I)}+\|b_n\|_{L^\infty(I)}\bigr)<\infty .
\]
\item Local pressure decomposition is stable under the splitting: if tensor
coefficients tend to zero in \(L^1_IL^{3/2}(B_R)\), then the corresponding
pressure gradients tend to zero in \(L^1_IH^{-1}(B_R)\), after subtracting
spatial means.
\item Cutoff commutators generated by localizing to \(B_R\) are controlled by
the \(L^\infty_I L^3\), \(L^1_I L^{3/2}\), and pressure-gradient norms appearing
above.  In particular, for any compact \(\chi_R\in C_c^\infty(B_{2R})\) with
\(\chi_R\equiv1\) on \(B_R\), the commutators satisfy
\[
  \|[\Delta,\chi_R]u\|_{H^{-1}(B_R)}
  +\|[\nabla\chi_R,\,\nabla\cdot](u\otimes u)\|_{H^{-1}(B_R)}
  \le C_R\|u\|_{L^3(B_{2R})}^2.
\]
uniformly for \(u\in L^3(B_{2R})\) and the same type of estimate for
\(S_n\)-quadratic commutator terms.
\end{enumerate}
\end{definition}

\begin{remark}[Terminal \(L^\infty_I L^3\) inputs are perturbative hypotheses]
\label{paper3:rem:terminal-perturbative-interface}
The \(L^\infty_I L^3\)-vanishing of \(r_n\) in item \textup{(iii)} and the
\(L^\infty_I L^3\) bound on \(U_n\) in item \textup{(iv)} are defining
inputs of an admissible perturbative terminal active family.  They are not
derived from the later spacetime compactness package of
\cref{paper5:def:selected-scale-collapse-compactness-package}.  If a physical
ordered branch does not provide these pointwise critical inputs on the chosen
compact cylinder, then the branch is routed first by the critical local upper
row of \cref{paper6a:lem:terminal-local-critical-mass-routing} or by an
earlier active-family failure, rather than being treated by
\cref{paper3:thm:terminal-decoupling}.
\end{remark}

\begin{theorem}[Terminal profile-evolution decoupling]\label{paper3:thm:terminal-decoupling}
Consider an admissible terminal active family on \(B_{2R}\times I\).  Let
\(U_n\) be the selected-frame velocity and let \(S_n\) be the sum, in that
frame, of all nonselected active profiles together with the rescaled remainder.
Then on \(B_R\times I\):
\begin{enumerate}[label=\textup{(\roman*)}]
\item \(\nabla\cdot S_n=0\) distributionally;
\item \(S_n\to0\) in \(L^\infty_I L^3(B_{2R})\);
\item \(U_n\otimes S_n+S_n\otimes U_n+S_n\otimes S_n\to0\) in
\(L^1_I L^{3/2}(B_R)\);
\item
\[
  \begin{aligned}
  &\exists\,\pi_n,\ F_n:\quad
  \partial_t S_n-\Delta S_n
  +a_n\bigl(y\cdot\nabla S_n + S_n\bigr)
  +b_n\cdot\nabla S_n+\nabla\pi_n =F_n,\\
  &\|F_n\|_{L^1_IH^{-1}(B_R)}\to0,\qquad
  \nabla\cdot S_n=0\quad\text{in }\mathcal D'(B_R\times I).
  \end{aligned}
\]
\[
  \|\nabla\pi_n\|_{L^1_IH^{-1}(B_R)}\to0.
\]
\item every pressure source containing at least one factor of \(S_n\) has
vanishing pressure gradient in \(L^1_IH^{-1}(B_R)\);
\item localization errors introduced by compact cutoffs vanish in
\(L^1_IH^{-1}(B_R)\);
\item the effective modulation parameters of the selected frame remain bounded;
\item after subtracting \(S_n\), the selected frame remains a positive finite-mass
local Type II branch with pressure reconstruction and the compact-cylinder
bounds required by the single-core analysis.
\end{enumerate}
\end{theorem}

\begin{proof}
Each profile in the family is divergence-free, and the rescaled remainder is
divergence-free.  Since divergence commutes with the Navier--Stokes critical
rescalings, finite sums remain divergence-free in distributions.  This establishes
\((i)\).

We prove \((ii)\).  By admissibility, every nonselected active profile is either
separated from the selected physical point or is a same-point strictly outer
scale.  The separated profiles vanish locally by \cref{paper3:lem:separated-vanish},
and the same-point outer profiles vanish locally by \cref{paper3:lem:outer-vanish}.
The remainder tends to zero in \(L^\infty_I L^3(B_{2R})\) by
\cref{paper3:def:admissible-terminal}.  The number of active profiles is finite by
\cref{paper3:lem:finite-active-count}.  Since every summand in \(S_n\)
tends to zero in \(L^\infty_I L^3(B_{2R})\), the triangle inequality gives
\[
  \|S_n\|_{L^\infty_I L^3(B_{2R})}
  \le
  \sum_{j\ne j_*}\|S^j_n\|_{L^\infty_I L^3(B_{2R})}
  +\|r_n\|_{L^\infty_I L^3(B_{2R})}\to0.
\]
This establishes \((ii)\).

For \((iii)\), admissibility gives
\[
  \sup_n\|U_n\|_{L^\infty_I L^3(B_{2R})}<\infty.
\]
Hence, by H\"older's inequality,
\[
  \|U_n\otimes S_n\|_{L^1_IL^{3/2}(B_R)}
  \le |I|\|U_n\|_{L^\infty_IL^3(B_R)}
       \|S_n\|_{L^\infty_IL^3(B_R)}\to0,
\]
and the same estimate applies to \(S_n\otimes U_n\).  Finally,
\[
  \|S_n\otimes S_n\|_{L^1_IL^{3/2}(B_R)}
  \le |I|\|S_n\|_{L^\infty_IL^3(B_R)}^2\to0.
\]

For \((iv)\), write the represented Navier--Stokes equation for the full
profile sum \(U_n+S_n\) and subtract the represented equation for the selected
component \(U_n\).  The remaining equation for \(S_n\) has the displayed linear
left-hand side
\[
  \partial_t S_n-\Delta S_n
  +a_n\bigl(y\cdot\nabla S_n+S_n\bigr)+b_n\cdot\nabla S_n.
\]
Collect all right-hand terms into
\[
  F_n:= -\nabla\cdot(U_n\otimes S_n+S_n\otimes U_n+S_n\otimes S_n)
       +f_n+E_n^{\mathrm{cut}}+E_n^{\mathrm{mod}},
\]
where \(f_n\) is the remainder forcing from item \((iii)\) of admissibility, and
\(E_n^{\mathrm{cut}},E_n^{\mathrm{mod}}\) are localization and modulation errors.
The nonlinear stresses tend to zero in \(L^1_IL^{3/2}(B_R)\), hence in
\(L^1_IH^{-1}(B_R)\), because \(L^{3/2}(B_R)\hookrightarrow H^{-1}(B_R)\) on a
bounded three-dimensional ball.  The residual term vanishes by
\cref{paper3:def:admissible-terminal}.  The cutoff terms vanish by the localization
part of admissibility, recorded explicitly in \((vi)\).  This gives the
displayed convergence in \(L^1_IH^{-1}(B_R)\).

For \((v)\), decompose the pressure source of \(U_n+S_n\) into the pure selected
source, the pure exterior source, and the two mixed sources.  The mixed sources
are
\[
  \partial_i\partial_j(U_{n,i}S_{n,j}+S_{n,i}U_{n,j})
  \quad\text{and}\quad
  \partial_i\partial_j(S_{n,i}S_{n,j}).
\]
By \((iii)\), their tensor coefficients tend to zero in
\(L^1_IL^{3/2}(B_R)\).  The pressure-stability part of admissibility,
equivalently local Calder\'on--Zygmund estimates plus harmonic pressure control
modulo spatial means, gives convergence of the corresponding pressure gradients
to zero in \(L^1_IH^{-1}(B_R)\).  Pure exterior sources are generated by
profiles that are locally invisible in the selected frame; applying the same
pressure stability after \cref{paper3:lem:separated-vanish,paper3:lem:outer-vanish} gives
their vanishing pressure gradients.

For \((vi)\), let \(\chi_R\) be a cutoff equal to one on \(B_R\) and supported
in \(B_{2R}\).  The commutators are finite sums of terms of the following
types:
\[
  (\nabla\chi_R)\cdot\nabla S_n,
  \qquad
  (\Delta\chi_R)S_n,
  \qquad
  (\nabla\chi_R)\cdot (U_n\otimes S_n+S_n\otimes U_n+S_n\otimes S_n),
\]
plus analogous pressure cutoff terms.  Their supports lie in the fixed annulus
\(B_{2R}\setminus B_R\).  By the localization part of admissibility, these
commutators are controlled by the norms already shown to vanish in \((ii)\),
\((iii)\), and \((v)\).  Hence all localization errors vanish in
\(L^1_IH^{-1}(B_R)\).

For \((vii)\), the only exterior contribution that can remain nonzero at first
order in a compact selected frame is a constant ambient velocity.  This has
already been absorbed into the translation modulation.  The remaining exterior
profiles vanish locally, so their contribution to the modulation equations is
small in the same local \(L^3\) topology used above.  The selected modulation
coefficients are bounded by admissibility; hence the effective parameters remain
bounded.

For \((viii)\), the positive CKN mass of the selected profile is unchanged by
subtracting a term that tends to zero in local \(L^3\) and whose
mixed pressure sources vanish by \((v)\).  The pressure reconstruction is stable
under the same decomposition by the pressure-stability condition in
\cref{paper3:def:admissible-terminal}.  The compact-cylinder bounds pass to the
selected component because the full family is bounded and the nonretained terms
vanish on compact cylinders.  Thus the selected frame remains a positive
finite-mass local Type II branch with the analytic inputs needed for the
single-core analysis.
\end{proof}

\subsection{Same-point and separated-point reductions}

\begin{theorem}[Same-point cascade reduction]\label{paper3:thm:same-point-reduction}
Every same-point multibubble or strict same-point cascade reduces, in the
selected innermost or compound frame, to one of the following: the single-core
criterion, scale-rigidity, or a perturbative terminal branch covered by
\cref{paper3:thm:terminal-decoupling}, provided the terminal family generated in the
selected innermost frame is admissible in the sense of
\cref{paper3:def:admissible-terminal}.
\end{theorem}

\begin{proof}
\noindent\textit{Step 1: scale-comparability classes.}
Split active same-point profiles into equivalence classes by
\[
  i\sim j \iff \exists c,C\in(0,\infty):
  c\le \frac{\lambda^i_n}{\lambda^j_n}\le C \quad \text{for all }n
\]
after passing to a subsequence.
Thus within each class there are constants \(c_{ij},C_{ij}\) with
\(c_{ij}\le\lambda^i_n/\lambda^j_n\le C_{ij}\) for all \(n\), so one may define
a single common rescaling sequence (choose \(\lambda_n^{(k)}=\lambda_n^{j_k}\)) and write
all class-\(k\) profiles in the same coordinates.
Within each class, finite sums of divergence-free profiles remain divergence-free.

\smallskip
\noindent\textit{Step 2: compound profile reduction inside each class.}
For each class set
\[
  \Phi^k_n:=\sum_{j\in J_k}\lambda^j_n\phi^j(x_n^j+\lambda^j_n y).
\]
If the \(L^3\)-size of this rescaled sum satisfies
\[
  \left\|\sum_{j\in J_k}\phi^j\right\|_{L^3(\mathbb R^3)}<\varepsilon_\star,
\]
with \(\varepsilon_\star\) the perturbative small-data threshold, then standard local
small-data theory for Navier--Stokes in the critical space eliminates this class
from the active alternatives.

If the class is above threshold, it is retained as a single compound active core;
at most one active core from that class is needed for local selection.

\smallskip
\noindent\textit{Step 3: local alternatives.}
If the retained compound core is unique and its scale-centering Jacobian is
nondegenerate, we obtain the single selected core alternative.
If this Jacobian is degenerate, the localized choice of active center or scale is
not unique, which is the gauge-degenerate alternative in
\cref{paper3:def:local-multibubble}.

\smallskip
\noindent\textit{Step 4: strict same-point scale separation.}
Assume now the branch is not reduced to comparable scales and choose the innermost
active scale \(\lambda^1_n\) in that physical point.  By
\cref{paper3:lem:outer-vanish}, every profile with larger scale vanishes in the
\(\lambda^1_n\)-frame in local \(L^3\), up to constant drift terms.
These constant drifts are absorbed into translation modulation by admissibility.
Let \(S_n\) be the sum of all such outer profiles plus the rescaled remainder
\(r_n\).  Then, by Step 2 and \cref{paper3:lem:outer-vanish},
\[
  S_n\to0\quad\text{in }L^\infty_I L^3(B_{2R}).
\]

Subtracting the equation for the innermost selected profile from the full-frame
equation gives a perturbation equation for the selected profile whose forcing
belongs to \(L^1_I H^{-1}(B_R)\) and vanishes.  Thus
\cref{paper3:thm:terminal-decoupling} applies, and one of three outcomes holds:
\begin{enumerate}[label=\textup{(\alph*)}]
\item \textbf{Single-core branch:}
  after subtracting \(S_n\), the branch is governed by one active component, so it
  reduces to the compact single-core branch of \cref{paper3:prop:local-single-core-exclusion};
\item \textbf{Scale-rigid branch:}
  the relative Jacobian between comparable scales becomes rigid in the sense of
  the later scale-rigidity closure package, so the branch is routed to the
  scale-rigid local row eventually closed by
  \cref{paper3:prop:scale-rigidity-discharged};
\item \textbf{Purely perturbative branch:}
  \(S_n\to0\) and the selected profile satisfies a vanishing forcing perturbation
  in \(L^1_IH^{-1}(B_R)\).  No independent same-point cascade mechanism remains;
  after subtracting \(S_n\), the selected component is reduced to the single-core
  or scale-rigid local row.
\end{enumerate}

These are the strict same-point scale-separation alternatives, so the
same-point reduction is exhausted.
\end{proof}


\begin{theorem}[Separated-point decoupling]\label{paper3:thm:separated-point-reduction}
Separated active physical points are locally invisible in each other's active
rescaled frames.  Consequently every separated-point multibubble reduces in each
selected active frame to a single-core, scale-rigid, or perturbative terminal
branch, provided the terminal family in that selected active frame is
admissible in the sense of \cref{paper3:def:admissible-terminal}.
\end{theorem}

\begin{proof}
Fix an active point \(p\) with frame \((x^p_n,\lambda^p_n)\).  Let \(q\ne p\)
be another active point.  By separation,
\[
  \frac{|x^q_n-x^p_n|}{\lambda^p_n}\to\infty.
\]
Thus \cref{paper3:lem:separated-vanish} gives local \(L^3\) convergence to zero of the
\(q\)-profile in the \(p\)-frame on every compact ball.  The same conclusion
holds for each profile centered at any physical point different from \(p\).
Since only finitely many active profiles remain after the critical mass
threshold is imposed, the sum of all profiles outside the \(p\)-frame tends to
zero locally in \(L^3\).

The pure pressure generated by a separated profile is locally invisible by the
same change of variables whenever its profile pressure belongs to
\(L^{3/2}_{\mathrm{loc}}\) with the standard Calder\'on--Zygmund control.  Mixed
pressure terms contain at least one exterior velocity factor, hence vanish in
\(L^1H^{-1}_{\mathrm{loc}}\) by the pressure part of
\cref{paper3:thm:terminal-decoupling}.  Localization and modulation commutators are
localized to \(B_{2R}\setminus B_R\), so their \(H^{-1}\)-norm is controlled by
the \(L^3\)-smallness already proved in \cref{paper3:lem:separated-vanish,paper3:lem:outer-vanish}
and by admissibility item \((\text{vii})\).

Therefore, in the selected \(p\)-frame, all other physical points contribute
only a perturbative forcing tending to zero in \(L^1H^{-1}\) on compact
cylinders.  The selected frame remains a positive finite-mass branch with the
same compact-cylinder and pressure reconstruction properties.  Hence the branch
reduces to the single-core case, the scale-rigid case, or the perturbative
terminal branch.  Since \(p\) was arbitrary, the same reduction holds in every
active frame.
\end{proof}

\subsection{Multibubble and cascade exclusion}

The following statements isolate the analytic inputs needed for the
multibubble reduction: the (already proved) single-core exclusion, the
scale-rigidity reduction (proved below), and terminal active-frame
admissibility (deduced here from the terminal nonlinear decoupling theorem
\cref{paper8:thm:terminal-nonlinear-decoupling}).  They are stated
explicitly so that the following theorem uses no unstated compactness,
pressure, or profile completeness conclusion.

\begin{proposition}[Local single-core exclusion]\label{paper3:prop:local-single-core-exclusion}
Every compact single-core branch with retained positive local velocity concentration,
nondegenerate scale-center selection, pressure reconstruction, bounded
compact-cylinder suitable estimates, bounded modulation, and vanishing localized
dissipation is excluded from the local Type II regime.
\end{proposition}

\begin{proof}
This is precisely the conclusion of the local single-core Type II criterion,
\cref{paper1:thm:single-core-criterion}, applied to the represented compact
single-core branch in the renormalized chart.  The hypotheses listed here are
the input data of that theorem: retained positive local velocity concentration provides the
nondegenerate selected core, the nondegenerate scale-center selection together
with the pressure reconstruction and bounded compact-cylinder suitable
estimates are the inputs of the compactness statement
\cref{paper1:prop:compactness}, bounded modulation is
\cref{paper1:rem:bounded-modulation}, and vanishing localized dissipation is
\cref{paper1:def:localized-dissipation}.  The conclusion of
\cref{paper1:thm:single-core-criterion} is exactly the stated exclusion.
\end{proof}

\begin{proposition}[Local scale-rigidity exclusion]\label{paper3:prop:local-scale-rigidity-exclusion}
Every scale-rigid local branch produced by the same-point or separated-point
reductions, including a branch with degenerate relative scale-center Jacobian,
either has a nonzero bounded selected-window limit or reduces to the single-core
branch in \cref{paper3:prop:local-single-core-exclusion}.  Consequently no such branch is compatible
with the local Type II condition.
\end{proposition}

\begin{proof}
This is proved at the end of the subsection as
\cref{paper3:prop:scale-rigidity-discharged}.  The present proposition is a
forward reference used by the multibubble reduction.
\end{proof}

\begin{remark}[Status of \cref{paper3:prop:local-scale-rigidity-exclusion}]
\label{paper3:rem:scale-rigidity-discharge}
The remainder of this subsection proves
\cref{paper3:prop:local-scale-rigidity-exclusion} by an explicit local compactness
stratification.  The proof consists of the local lemmas
\cref{paper3:lem:scale-rigid-bounded-limit-exit,paper3:lem:scale-rigid-pressure-modulation-stability,paper3:lem:scale-rigid-chart-comparability,paper3:lem:scale-rigid-two-sided-ckn-transfer,paper3:lem:scale-rigid-ckn-in-gauge,paper3:lem:scale-rigid-scaled-pressure-decay,paper3:lem:scale-rigid-dyadic-ckn-persistence,paper3:lem:scale-rigid-exterior-pressure-routing,paper3:lem:scale-rigid-pressure-core-activation,paper3:lem:scale-rigid-active-core-capture,paper3:lem:scale-rigid-no-subconcentration-bounded,paper3:lem:scale-rigid-degenerate-jacobian,paper3:thm:scale-rigid-single-core-compatibility}
and the main decomposition
\cref{paper3:thm:scale-rigid-stratum-decomposition}.  The final
\cref{paper3:prop:scale-rigidity-discharged} then proves the conclusion of
\cref{paper3:prop:local-scale-rigidity-exclusion} from these results.  All proofs use only the
local Type II definition \cref{paper1:def:typeII}, the represented variables of
\cref{paper2:thm:C}, the local pressure decomposition of
\cref{paper2:thm:pressure-decomp}, the modulation \(L^\infty\) control of
\cref{paper2:thm:B}, the same-point and separated-point reductions
\cref{paper3:thm:same-point-reduction,paper3:thm:separated-point-reduction},
the minimal bad-configuration extraction
\cref{paper3:thm:minimal-bad-extraction}, the single-core exclusion
\cref{paper1:thm:single-core-criterion}, the terminal CKN regularity criterion
\cref{paper8:thm:terminal-ckn-regularity}, and Banach--Alaoglu plus
Calder\'on--Zygmund on fixed cylinders \cite{CalderonZygmund1952,Stein1970}.
No global PDE bound, no Liouville theorem, and no profile-decomposition
classification are invoked.
\end{remark}

\subsubsection{Local stratification of scale-rigid branches}

We now identify the named scale-rigid local stratum and prove that every
branch placed there by
\cref{paper3:thm:same-point-reduction,paper3:thm:separated-point-reduction}
either exits the local Type II class through
\cref{paper1:def:typeII} or relaxes back to one of the already excluded local
strata.

\begin{definition}[Scale-rigid local branch]
\label{paper3:def:scale-rigid-branch}
A scale-rigid local branch is a represented branch
\((V_n,P_n,a_n,b_n)\) on a fixed compact selected cylinder \(Q_{R_*}\),
produced by Step~3 or Step~4 of \cref{paper3:thm:same-point-reduction} or by
\cref{paper3:thm:separated-point-reduction}, such that:
\begin{enumerate}[label=\textup{(R\arabic*)}]
\item the compact-cylinder suitable estimates are uniformly bounded,
      \[
        A_{V_n}(R_*)+E_{V_n}(R_*)+C_{V_n}(R_*)+D_{P_n}(R_*)\le M;
      \]
\item the branch carries retained positive local velocity concentration on a fixed inner
      cylinder \(Q_{r_*}\Subset Q_{R_*}\),
      \[
        C_{V_n}(r_*)\ge\eta_0;
      \]
      If the retained concentration is pressure-only, the branch is assigned
      instead to the harmonic pressure row or to the exterior-pressure row.
\item the modulation coefficients are bounded on the selected cylinder,
      \[
        \|a_n\|_{L^\infty(-R_*^2,0)}+\|b_n\|_{L^\infty(-R_*^2,0)}\le M;
      \]
\item the relative scale-center selection between the retained core and any
      previously selected core is rigid in one of the senses produced by
      \cref{paper3:thm:same-point-reduction}: either the relative
      scale-center Jacobian degenerates along the subsequence (the
      gauge-degenerate alternative of
      \cref{paper3:def:local-multibubble}\,(v)), or the inner core is the
      innermost member of a same-point cascade in which all outer profiles
      have been absorbed into the perturbative remainder by
      \cref{paper3:lem:outer-vanish}.
\end{enumerate}
\end{definition}

\begin{remark}[Inputs (R1)--(R3) are automatic on the selected cylinder]
\label{paper3:rem:scale-rigid-inputs-automatic}
The bounds (R1)--(R3) are not extra hypotheses but are inherited from the
preceding reductions.  Bound (R1) is the compact-cylinder hypothesis carried
through \cref{paper3:thm:minimal-bad-extraction}; if it failed,
\cref{paper3:lem:compactness-failure} would assign the branch to the
multicore alternative, not to the scale-rigid stratum.  Bound (R2) is the
velocity-carrier part of the positive concentration capture in
\cref{paper3:def:active-family}: after pressure-only carriers have been routed
to the harmonic/exterior-pressure row, the retained scale-rigid row is entered
only by a fixed positive \(L^3\) velocity carrier on \(Q_{r_*}\).  Bound (R3)
is the selected-cylinder modulation control of
\cref{paper2:thm:B}, which is part of the represented-variable output of
\cref{paper2:thm:C}.  No additional analytic input is needed to enter the
scale-rigid stratum.
\end{remark}

\begin{lemma}[Bounded selected-window limits have VMO pressure]
\label{paper3:lem:bounded-window-vmo-pressure}
Let \((V_n,P_n)\) be a represented suitable sequence on a parabolic cylinder
\(Q_R\), with the local pressure atlas and compact-cylinder suitable bounds
available on every \(Q_{R'}\Subset Q_R\).  Assume that, after subtracting
functions of time from \(P_n\),
\[
   V_n\to V_*
   \qquad\text{strongly in }L^3(Q_R),
\]
and that
\[
   V_*\in L^\infty(Q_R).
\]
Then, after passing to a subsequence, there exists a pressure \(P_*\) such that
\[
   P_n\rightharpoonup P_*
   \qquad\text{weakly in }L^{3/2}_{\mathrm{loc}}(Q_R),
\]
and \((V_*,P_*)\) is a suitable weak solution on every compact subcylinder of
\(Q_R\).  Moreover, for every \(Q_{R''}\Subset Q_R\),
\[
   \lim_{\rho\downarrow0}
   \sup_{Q_\rho(z)\Subset Q_{R''}}
   \left[
      \rho^{-2}\iint_{Q_\rho(z)} |V_*|^3
      +
      \rho^{-2}\iint_{Q_\rho(z)}
      |P_*-(P_*)_{B_\rho(x_z)}(t)|^{3/2}
   \right]
   =0.
\]
Finally, on every smaller subcylinder, the local Calder\'on--Zygmund pressure
pieces associated to \(V_n\otimes V_n\) converge strongly in \(L^{3/2}\) to the
corresponding local pressure piece associated to \(V_*\otimes V_*\).
\end{lemma}

\begin{proof}
Fix \(Q_{4r}(z_0)\Subset Q_R\), and let \(\zeta\in C_c^\infty(B_{4r}(x_0))\)
satisfy \(\zeta\equiv1\) on \(B_{3r}(x_0)\).  Decompose, after subtracting a
function of time,
\[
   P_n=P_n^{\mathrm{loc}}+H_n,
\]
where
\[
   -\Delta P_n^{\mathrm{loc}}
   =
   \partial_i\partial_j\bigl(\zeta (V_n)_i(V_n)_j\bigr)
\]
in \(B_{4r}(x_0)\), and \(H_n(\cdot,t)\) is harmonic in \(B_{3r}(x_0)\) for
a.e. \(t\).
Since
\[
   V_n\to V_*
   \qquad\text{strongly in }L^3(Q_{4r}(z_0)),
\]
we have
\[
   V_n\otimes V_n\to V_*\otimes V_*
   \qquad\text{strongly in }L^{3/2}(Q_{4r}(z_0)).
\]
Thus the Calder\'on--Zygmund estimate gives
\[
   P_n^{\mathrm{loc}}\to P_*^{\mathrm{loc}}
   \qquad\text{strongly in }L^{3/2}(Q_{4r}(z_0)),
\]
where
\[
   -\Delta P_*^{\mathrm{loc}}
   =
   \partial_i\partial_j\bigl(\zeta (V_*)_i(V_*)_j\bigr).
\]
The pressure-atlas bounds give \(H_n\) bounded in
\(L^{3/2}_tL^{3/2}_x(Q_{3r}(z_0))\).  Since \(H_n(\cdot,t)\) is harmonic in
space, the interior harmonic estimates imply, for each fixed finite \(m\),
\[
   \|\nabla_x^m H_n(t)\|_{L^{3/2}(B_{2r}(x_0))}
   \le C_{m,r}\|H_n(t)\|_{L^{3/2}(B_{3r}(x_0))}
\]
for a.e. \(t\).  Hence, after passing to a subsequence,
\[
   H_n\rightharpoonup H_*
   \qquad\text{weakly in }L^{3/2}_tW^{m,3/2}_x(B_{2r}(x_0))
\]
for every fixed finite \(m\), and \(H_*(\cdot,t)\) is harmonic in
\(B_{2r}(x_0)\) for a.e. \(t\).  Set \(P_*:=P_*^{\mathrm{loc}}+H_*\).  Passing
to the limit in the represented equation
uses the strong \(L^3\) convergence of \(V_n\) and the weak \(L^{3/2}\)
compactness of \(P_n\); the local energy inequality passes to the limit by
lower semicontinuity of the dissipation and strong convergence of the cubic
velocity terms.  Thus \((V_*,P_*)\) is suitable on compact subcylinders.

Since \(V_*\in L^\infty(Q_R)\), for every compact \(Q''\Subset Q_R\),
\[
   C_{V_*}(2\ell;z)
   \lesssim \|V_*\|_{L^\infty(Q_R)}^3\ell^3
\]
uniformly for \(Q_{2\ell}(z)\Subset Q''\).  Fix \(Q'\Subset Q_R\), choose
\(Q''\) with \(Q'\Subset Q''\Subset Q_R\), and choose a compact outer scale
\(\ell_0>0\) so that \(Q_{2\ell_0}(z)\Subset Q''\) for all centres \(z\in Q'\).
The local
\(L^{3/2}\)-bound for \(P_*\) gives a finite uniform bound for
\(D_{P_*}(\ell_0;z)\) over these centres.  Iterating
\cref{paper3:lem:scale-rigid-scaled-pressure-decay} from the fixed outer scale
\(\ell_0\), with a fixed \(0<\theta\le1/4\), gives
\[
   D_{P_*}(\theta^k\ell_0;z)
   \le
   (C_{\rm dec}\theta^{5/2})^kD_{P_*}(\ell_0;z)
   +C\|V_*\|_{L^\infty(Q_R)}^3
     \sum_{j=0}^{k-1}(C_{\rm dec}\theta^{5/2})^{k-1-j}\theta^{3j}\ell_0^3 .
\]
Taking \(\theta\) small enough that \(C_{\rm dec}\theta^{5/2}<1\), the right
side tends to zero as \(k\to\infty\).  The pressure-localization estimate
\cref{paper1:lem:D-localization} between adjacent geometric scales then yields
\[
   \lim_{\rho\downarrow0}
   \sup_{Q_\rho(z)\Subset Q'}
   D_{P_*}(z,\rho)=0 .
\]
Together with the trivial estimate
\[
   C_{V_*}(z,\rho)\lesssim \|V_*\|_{L^\infty(Q_R)}^3\rho^3,
\]
this gives the desired velocity-pressure CKN smallness for the limit on compact
subcylinders.

To transfer this smallness to the approximating sequence, decompose
\[
   P_n-P_*
   =
   (P_n^{\mathrm{loc}}-P_*^{\mathrm{loc}})
   +(H_n-H_*)
\]
modulo functions of time on \(B_{2r}(x_0)\).  The first term converges
strongly in \(L^{3/2}\) by the Calder\'on--Zygmund estimate already proved.
For the harmonic term, interior harmonic estimates give
\[
   \|\nabla_x(H_n-H_*)(t)\|_{L^\infty(B_{3r/2}(x_0))}
   \le C_r\|H_n(t)-H_*(t)\|_{L^{3/2}(B_{2r}(x_0))}
\]
for a.e. \(t\).  Hence for every \(Q_\rho(z)\Subset Q_{R''}\),
\[
   \rho^{-2}\iint_{Q_\rho(z)}
   |(H_n-H_*)-((H_n-H_*))_{B_\rho(x_z)}(t)|^{3/2}
   \le C_r \rho^{5/2}
   \int_{t_z-\rho^2}^{t_z}
   \|H_n(t)-H_*(t)\|_{L^{3/2}(B_{2r})}^{3/2}\,dt .
\]
The right-hand side is bounded uniformly in \(n\) by the
\(L^{3/2}_tL^{3/2}_x\)-bound on the harmonic pieces and therefore tends to zero
uniformly as \(\rho\downarrow0\).  This yields the required small-scale
pressure decay for \(P_n\) after first taking \(\rho\) small and then \(n\)
large.  The strong convergence of the local pressure pieces has already been
proved above.
\end{proof}

\begin{lemma}[Bounded selected-window limits are CKN-small away from pressure failure]
\label{paper3:lem:bounded-window-ckn-smallness}
Let \((V_n,P_n)\) be represented suitable windows on a compact cylinder
\(Q_R(z_0)\).  Suppose that, after passing to a subsequence,
\[
   V_n\to V_*
   \qquad\text{strongly in }L^3(Q_R(z_0)),
   \qquad V_*\in L^\infty(Q_R(z_0)).
\]
Assume also that the local pressure atlas / pressure reconstruction row has
not failed on \(Q_R(z_0)\).  Then for every compact subcylinder
\(Q\Subset Q_R(z_0)\) and every \(\varepsilon>0\) there exists
\(0<\rho_Q<R\), depending only on \(Q\), \(\varepsilon\), the local pressure
atlas, and \(\|V_*\|_{L^\infty(Q_R)}\), such that
\[
   C_{V_n}(z,\rho_Q)+D_{P_n}(z,\rho_Q)<\varepsilon
\]
for all sufficiently large \(n\) and every admissible parabolic cylinder
\(Q_{\rho_Q}(z)\Subset Q\).  If this conclusion fails, then the first failed
local test is the pressure atlas, pressure gauge, or pressure compactness row.
\end{lemma}

\begin{proof}
Fix \(Q\Subset Q_R(z_0)\).  Because \(V_*\in L^\infty(Q_R)\), for every
\(z\in Q\) one has
\[
   C_{V_*}(z,\rho)
   =\rho^{-2}\int_{Q_\rho(z)}|V_*|^3\,dx\,dt
   \le \|V_*\|_{L^\infty(Q_R)}^3\,\rho^3\to0
\]
as \(\rho\downarrow0\), uniformly for \(Q_\rho(z)\Subset Q\).  Choose
\(\rho_Q\) so that this velocity contribution is below \(\varepsilon/8\) on a
finite cover of \(Q\).  Strong convergence in \(L^3(Q_R)\) then gives
\[
   C_{V_n-V_*}(z,\rho_Q)<\varepsilon/8
\]
for all sufficiently large \(n\) and every admissible center in the cover.
Hence the velocity part is below \(\varepsilon/2\).

By \cref{paper3:lem:bounded-window-vmo-pressure}, the bounded selected-window
limit has uniformly vanishing velocity--pressure CKN quantity at sufficiently
small scales on compact subcylinders.
The strong convergence of \(V_n\) in \(L^3\), together with strong convergence
of the local Calder\'on--Zygmund pressure pieces and harmonic pressure decay,
therefore transfers this smallness to \((V_n,P_n)\) for all large \(n\),
unless the pressure-atlas or harmonic-pressure bounds fail.  Such a failure is
precisely one of the named pressure rows.

Combining the velocity and pressure estimates gives the displayed CKN
smallness.  If any pressure step is unavailable, the first unavailable item is
exactly the pressure atlas, pressure gauge, or pressure compactness row, so
the bounded window is not retained as an independent terminal branch.
\end{proof}

\begin{theorem}[Bounded selected-window row exits]
\label{paper3:lem:scale-rigid-bounded-limit-exit}
Let \((V_n,P_n,a_n,b_n)\) be a represented sequence on a compact selected
cylinder \(Q_r\), and assume that, after passing to a selected-window
subsequence, there is a nonzero \(V_*\in L^\infty(Q_r)\) such that
\[
  V_n\to V_*
  \qquad\text{strongly in }L^3(Q_r).
\]
Then this sequence cannot be a terminal branch of the strict selected-window
Type II row.  More precisely, after running the ordered local tests, exactly
one of the following occurs:
\begin{enumerate}[label=\textup{(\roman*)}]
\item the local pressure atlas, pressure gauge, or pressure compactness row
      fails;
\item the bounded selected window is CKN-small on a smaller compact core and
      is locally regular/removable by the Caffarelli--Kohn--Nirenberg
      regularity criterion;
\item the retained concentration is not carried by the bounded compact core
      and is reselected as a smaller active core, an exterior core, a
      multibubble/cascade label, a gauge-degenerate label, a pressure label, or
      another named local row in the ordered state-space decomposition.
\end{enumerate}
Consequently a bounded selected-window limit is a closed local routing row, not
an independent retained terminal Type II state.  In particular, by
\cref{paper1:def:typeII}, the branch is not a strict selected-window Type II
branch; the physical bounded-window event is handled by
\cref{paper6:thm:bounded-selected-window-physical-routing}.
\end{theorem}

\begin{proof}
The existence of a nonzero bounded selected-window limit is exactly the
negation of the strict selected-window condition in \cref{paper1:def:typeII},
so the branch is not a strict selected-window Type II branch.  This
observation alone is only branch routing, not physical closure.

To close the bounded row locally, apply
\cref{paper3:lem:bounded-window-ckn-smallness}.  If the pressure-atlas
hypotheses required there fail, item \textup{(i)} occurs.  Otherwise, on each
compact subcylinder of the bounded window, all sufficiently large selected
windows are CKN-small after shrinking once.  By the
Caffarelli--Kohn--Nirenberg regularity criterion, the bounded compact core is
regular and removable, giving item \textup{(ii)}.

If the original positive CKN concentration of the physical Type II germ still
survives once this regular compact core is removed, then it cannot be carried
by the bounded compact core.  The ordered local selection is then rerun on
the remaining concentration.  Since the state-space tests are ordered before
terminal retention, the surviving concentration is recorded as an inner
active core, exterior core, multibubble/cascade, gauge-degenerate, pressure,
or other named local row, giving item \textup{(iii)}.  No global Type I bound,
global \(L^3\) bound, or global profile decomposition is used.
\end{proof}

\begin{lemma}[Pressure and modulation stability under rigid-frame limits]
\label{paper3:lem:scale-rigid-pressure-modulation-stability}
Let \((V_n,P_n,a_n,b_n)\) be a scale-rigid local branch on \(Q_{R_*}\) in the
sense of \cref{paper3:def:scale-rigid-branch}.  Fix \(\rho\in(r_*,R_*)\).
After passing to a subsequence:
\begin{enumerate}[label=\textup{(\roman*)}]
\item \(V_n\rightharpoonup V_*\) weakly in
      \(L^3(Q_\rho)\cap L^2_t H^1_x(Q_\rho)\), and
      \(V_n\to V_*\) strongly in \(L^3(Q_\rho')\) for every
      \(\rho'\in(r_*,\rho)\);
\item the mean-normalized pressure
      \(\widetilde P_n:=P_n-\overline{P_n}_{B_\rho}\) converges weakly in
      \(L^{3/2}(Q_\rho)\) to a limit \(P_*\) such that
      \(-\Delta P_*=\partial_i\partial_j(V_{*,i}V_{*,j})\) on \(Q_\rho\) in
      the sense of distributions modulo a function of time;
\item \(a_n\rightharpoonup^* a_*\) and \(b_n\rightharpoonup^* b_*\)
      weakly-\(*\) in \(L^\infty(-R_*^2,0)\), with
      \(\|a_*\|_\infty+\|b_*\|_\infty\le M\);
\item the limit \((V_*,P_*,a_*,b_*)\) is a suitable weak solution of the
      represented Navier--Stokes equation \eqref{paper1:eq:represented-ns} on
      \(Q_\rho\), with \(V_*\in L^3(Q_\rho)\cap L^2_tH^1_x(Q_\rho)\).
\end{enumerate}
This lemma is only a compactness and stability statement.  It does not assert
boundedness of \(V_*\).  Bounded selected-window limits are obtained below from
the separate no-subconcentration/CKN argument.
\end{lemma}

\begin{proof}
\textit{(i) Compactness of \(V_n\).}  By (R1), the family \((V_n)\) is bounded
in \(L^3(Q_\rho)\) (from \(C_{V_n}(R_*)\)), in \(L^2_t H^1_x(Q_\rho)\) (from
\(A_{V_n}(R_*)+E_{V_n}(R_*)\)), and in \(L^\infty_t L^2_x(Q_\rho)\) (from
\(A_{V_n}(R_*)\)).  Banach--Alaoglu yields the weak limits in
\(L^3(Q_\rho)\) and \(L^2_t H^1_x(Q_\rho)\) along a subsequence.  The
represented Navier--Stokes equation \eqref{paper1:eq:represented-ns} bounds
\(\partial_s V_n\) in \(L^{4/3}_t H^{-1}_x(Q_\rho)\) by Sobolev product
estimates on the bounded cylinder, since each term on the right-hand side is
controlled by \(A_{V_n}(R_*)\), \(E_{V_n}(R_*)\), \(C_{V_n}(R_*)\),
\(D_{P_n}(R_*)\), and \(\|a_n\|_\infty+\|b_n\|_\infty\le M\); the modulation
correction \(a_n(s)(V_n+y\cdot\nabla V_n)+b_n(s)\cdot\nabla V_n\) is bounded
in \(L^2_t L^{6/5}_x(Q_\rho)\subset L^{4/3}_t H^{-1}_x(Q_\rho)\) on the
bounded cylinder.  Hence Aubin--Lions \cite{Aubin1963,Simon1987} applied on
the fixed cylinder \(Q_\rho\) yields strong convergence
\(V_n\to V_*\) in \(L^3(Q_{\rho'})\) for every \(\rho'\in(r_*,\rho)\).

\textit{(ii) Pressure stability.}  By (R1) the local pressure quantity
\(D_{P_n}(R_*)\) is bounded, equivalently \(\widetilde P_n\) is bounded in
\(L^{3/2}(Q_\rho)\) by the local Calder\'on--Zygmund mean-normalization
estimate \cite{CalderonZygmund1952,Stein1970} applied to the local pressure
identity in \cref{paper2:thm:pressure-decomp}.  Banach--Alaoglu yields
\(\widetilde P_n\rightharpoonup P_*\) weakly in \(L^{3/2}(Q_\rho)\).
Pass to the limit in the local pressure identity
\[
  -\Delta\widetilde P_n
  =\partial_i\partial_j(V_{n,i}V_{n,j})
  -\partial_i\partial_j(\overline{V_{n,i}V_{n,j}}_{B_\rho})
\]
on \(Q_\rho\).  The right-hand side converges in
\(\mathcal D'(Q_\rho)\) by the strong \(L^3\) convergence of (i) (which
implies weak \(L^{3/2}\) convergence of \(V_{n,i}V_{n,j}\) by H\"older).
This proves (ii).

\textit{(iii) Modulation limits.}  By (R3),
\(a_n,b_n\) are bounded in \(L^\infty(-R_*^2,0)\); Banach--Alaoglu applied to
the dual of \(L^1\) yields the weak-\(*\) limits with the same bound \(M\).

\textit{(iv) Suitable weak solution status.}  Each
\((V_n,P_n)\) is suitable on \(Q_\rho\) (by (R1) and the represented-equation
output of \cref{paper2:thm:C}).  The local energy inequality and the
distributional momentum equation pass to the limit using (i)--(iii) and
H\"older on the fixed cylinder \(Q_\rho\): the dissipation term is lower
semicontinuous under weak \(L^2_t H^1_x\) convergence, the convection term
converges by strong \(L^3(Q_{\rho'})\) convergence, and the pressure-velocity
pairing
\(\int \widetilde P_n V_n\cdot\nabla\chi\) converges by weak \(L^{3/2}\)
times strong \(L^3\) on the fixed cylinder.  The modulation correction
converges by weak-\(*\) \(L^\infty\) for \((a_n,b_n)\) against the strong
\(L^1_t\) convergence of \((V_n,y\cdot\nabla V_n,\nabla V_n)\) tested against
fixed test functions on \(Q_\rho\).  This proves the represented suitable
limit in the displayed local energy class.  No \(L^\infty\) bound is inferred
from the compact-cylinder estimates.
\end{proof}

\begin{lemma}[Compact chart comparability in scale-rigid gauges]
\label{paper3:lem:scale-rigid-chart-comparability}
Let \((V,P,a,b)\) be a represented branch on \(Q_{2\rho}(z_0)\) obtained from
the repaired-gauge chart of \cref{paper2:thm:C}.  Assume that the associated
chart \((\Lambda,X)\) satisfies
\[
  a=\frac{\Lambda'}{\Lambda},\qquad b=\frac{X'}{\Lambda},\qquad
  \|a\|_{L^\infty(I_{2\rho})}+\|b\|_{L^\infty(I_{2\rho})}\le M
\]
on the time interval \(I_{2\rho}\) of \(Q_{2\rho}(z_0)\), and set
\(\Lambda_0:=\Lambda(s_0)>0\) for \(z_0=(y_0,s_0)\).  There exists
\(\rho_c=\rho_c(M,\Lambda_0,\rho)>0\) such that, for every
\(0<r\le\rho_c\), the chart on \(Q_{2r}(z_0)\) obeys
\[
  \frac{\Lambda_0}{2}\le\Lambda(s)\le\frac{3\Lambda_0}{2},
  \qquad |X(s)-X(s_0)|\le 6\Lambda_0Mr^2,
  \qquad s\in(s_0-4r^2,s_0].
\]
After a harmless concentric shrink, if needed, the physical image remains
inside the support buffer of \cref{paper2:prop:domain}.  Consequently
\(\cref{paper2:prop:domain,paper2:prop:local-transfer}\) apply on the retained
subcylinder with
\[
  \Lambda_{\min}=\Lambda_0/2,\qquad \Lambda_{\max}=3\Lambda_0/2,
\]
and represented and physical CKN quantities on comparable cylinders are
equivalent with constants depending only on
\(M,\Lambda_0,r,\rho\) and the fixed buffer.
\end{lemma}

\begin{proof}
The identities
\[
a=\frac{\Lambda'}{\Lambda},\qquad b=\frac{X'}{\Lambda}
\]
are part of the exact repaired change of variables in
\cref{paper2:lem:change-id,paper2:thm:A,paper2:thm:C}. Hence, for
\(|s-s_0|\le4r^2\),
\[
  \left|\log\frac{\Lambda(s)}{\Lambda_0}\right|
  \le \int_s^{s_0}|a(\sigma)|\,d\sigma
  \le 4Mr^2,\qquad
  |X(s)-X(s_0)|
  \le \int_s^{s_0}\Lambda(\sigma)|b(\sigma)|\,d\sigma .
\]
Choose
\[
  \rho_c\le \min\{\rho/4,(\log(3/2)/(4M))^{1/2}\}
\]
when \(M>0\), and choose \(\rho_c=\rho/4\) when \(M=0\).  Then
\(e^{-4Mr^2}\Lambda_0\le\Lambda(s)\le e^{4Mr^2}\Lambda_0\), which implies the displayed
lower-upper scale bounds.  Substituting the upper scale bound into the center
integral gives
\[
  |X(s)-X(s_0)|\le (3\Lambda_0/2)\,M\,4r^2=6\Lambda_0Mr^2.
\]
A final concentric shrink, depending only on the fixed distance to the boundary
of the selected compact window, keeps the image inside the domain buffer.
\Cref{paper2:prop:domain,paper2:prop:local-transfer} then give the stated
comparison of the represented and physical quantities.  No solution
boundedness, and no vanishing dissipation, is used.
\end{proof}

\begin{lemma}[Two-sided local CKN transfer in comparable charts]
\label{paper3:lem:scale-rigid-two-sided-ckn-transfer}
Let \((V,P)\) be related to a physical suitable pair \((u,p)\) by the repaired
scale-translation chart on \(Q_{2r}(z_0)\), and assume the chart comparability
conclusion of \cref{paper3:lem:scale-rigid-chart-comparability} there:
\[
  0<\Lambda_{\min}\le \Lambda(s)\le \Lambda_{\max}<\infty,
  \qquad |X(s)-X(s_0)|\le M r^2 .
\]
Set
\[
  \mathcal C_V(z_0,r):=
  r^{-2}\int_{Q_r(z_0)}|V|^3
  +r^{-2}\int_{Q_r(z_0)}|P-(P)_{B_r}|^{3/2},
\]
and define the physical CKN quantity \(\mathcal C_u(x,t,\rho)\) analogously,
with pressure normalized by spatial averages.  There are constants
\[
  N_{\rm tr}<\infty,\qquad c_{\rm tr}\in(0,1),\qquad C_{\rm tr}<\infty,
\]
depending only on \(M,\Lambda_{\min},\Lambda_{\max}\) and the fixed support
buffer, and physical cylinders
\[
  Q^{\rm ph}_{\ell}:=
  B_{\rho_\ell}(x_\ell)\times(t_\ell-\rho_\ell^2,t_\ell),
  \qquad
  c_{\rm tr}\Lambda_{\min}r\le \rho_\ell\le C_{\rm tr}\Lambda_{\max}r,
  \qquad 1\le \ell\le N_{\rm tr},
\]
such that the physical image of \(Q_{c_{\rm tr}r}(z_0)\) is contained in
\(\cup_\ell Q^{\rm ph}_{\ell}\), each \(Q^{\rm ph}_{\ell}\) is contained in the
physical image of \(Q_r(z_0)\), and
\[
  \mathcal C_u(x_\ell,t_\ell,\rho_\ell)
  \le C_{\rm tr}\mathcal C_V(z_0,r)
  \qquad (1\le \ell\le N_{\rm tr}).
\]
Conversely, if a finite physical cylinder cover of the image of \(Q_r(z_0)\)
with comparable radii satisfies
\(\mathcal C_u(x_\ell,t_\ell,\rho_\ell)\le\varepsilon\) for every member of the
cover, then
\[
  \mathcal C_V(z_0,c_{\rm tr}r)\le C_{\rm tr}N_{\rm tr}\varepsilon .
\]
\end{lemma}

\begin{proof}
The chart is
\[
  x=X(s)+\Lambda(s)y,\qquad dt=\Lambda(s)^2\,ds,
  \qquad V(y,s)=\Lambda(s)u(x,t),\qquad P(y,s)=\Lambda(s)^2p(x,t).
\]
The bounds on \(\Lambda\) and \(X\) imply that the image of a represented
parabolic cylinder is trapped between parabolic cylinders whose radii are
comparable to \(\Lambda(s_0)r\).  A standard Besicovitch covering with the
bounded eccentricity supplied by
\(\Lambda_{\min},\Lambda_{\max}\) gives the finite cover and the bounded
overlap constant \(N_{\rm tr}\).

For the velocity term the change of variables gives, on every comparison
cylinder,
\[
  \rho_\ell^{-2}\int_{Q^{\rm ph}_{\ell}}|u|^3\,dx\,dt
  \le C_{\rm tr} r^{-2}\int_{Q_r(z_0)}|V|^3\,dy\,ds .
\]
For the pressure term we use the invariant seminorm
\(\inf_{c(s)}\int |P-c(s)|^{3/2}\).  Spatial averages are equivalent to this
seminorm up to the universal factor \(2^{3/2}\), because for every ball \(B\)
and \(q\ge1\),
\[
  \int_B |F-F_B|^q\le 2^q\inf_c\int_B|F-c|^q .
\]
The same Jacobian identity therefore gives the displayed physical estimate for
the pressure oscillation after subtracting spatial means.  Applying the same
argument to the inverse chart, whose scale is bounded between
\(\Lambda_{\max}^{-1}\) and \(\Lambda_{\min}^{-1}\) on the comparison cover,
gives the converse estimate.  No direction of
\cref{paper2:prop:local-transfer} is used as a black box here; both estimates
come from the exact change of variables and the pressure oscillation seminorm.
\end{proof}

\begin{lemma}[CKN regularity in comparable repaired gauges]
\label{paper3:lem:scale-rigid-ckn-in-gauge}
Let \((V,P,a,b)\) be a represented suitable weak solution on \(Q_{2r}(z_0)\)
obtained from the repaired-gauge chart of \cref{paper2:thm:C}.  Assume that the
chart satisfies the comparability conclusion of
\cref{paper3:lem:scale-rigid-chart-comparability} on this cylinder, with
constants \(\Lambda_{\min},\Lambda_{\max}\), and that
\[
  \|a\|_{L^\infty}+\|b\|_{L^\infty}\le M
\]
there.  There are constants
\[
  \varepsilon_{\rm rg}
  =\varepsilon_{\rm rg}(M,\Lambda_{\min},\Lambda_{\max})>0,
  \qquad
  c_{\rm rg}=c_{\rm rg}(M,\Lambda_{\min},\Lambda_{\max})\in(0,1),
\]
such that if
\[
  r^{-2}\int_{Q_r(z_0)} |V|^3
  +r^{-2}\int_{Q_r(z_0)} |P-(P)_{B_r}|^{3/2}
  \le\varepsilon_{\rm rg},
\]
then \(V\) is bounded on \(Q_{c_{\rm rg}r}(z_0)\), with a bound depending only
on \(M,\Lambda_{\min},\Lambda_{\max},r\), and the preceding displayed
quantity.
\end{lemma}

\begin{proof}
Apply the forward direction of
\cref{paper3:lem:scale-rigid-two-sided-ckn-transfer} to \(Q_r(z_0)\).  Choose
\(\varepsilon_{\rm rg}\) so small that every physical cylinder in the resulting
finite comparison cover has CKN quantity below the threshold
\(\varepsilon_{\rm CKN}\) from
\cref{paper8:thm:terminal-ckn-regularity}.  That theorem gives boundedness of
the physical velocity on smaller physical cylinders.  Pulling those estimates
back through the same comparable repaired chart gives the asserted
\(L^\infty\) bound for \(V\) on \(Q_{c_{\rm rg}r}(z_0)\).  Thus boundedness is
obtained only from epsilon regularity after two-sided chart transfer and small
CKN quantity, not from compactness or gauge degeneracy.
\end{proof}

\begin{lemma}[Scaled local pressure decay]
\label{paper3:lem:scale-rigid-scaled-pressure-decay}
Let \((V,P)\) be a represented suitable pair on \(Q_{2r}(z_0)\), and assume the
localized pressure decomposition of \cref{paper2:thm:pressure-decomp} is
available with a cutoff supported in \(B_{2r}(y_0)\) and equal to \(1\) on
\(B_r(y_0)\).  For
\[
  C_V(\rho;z_0):=\rho^{-2}\int_{Q_\rho(z_0)}|V|^3,\qquad
  D_P(\rho;z_0):=\rho^{-2}\int_{Q_\rho(z_0)}
       |P-(P)_{B_\rho}|^{3/2},
\]
there is a universal constant \(C_{\rm dec}\) such that, whenever
\(0<\theta\le1/4\),
\[
  D_P(\theta r;z_0)
  \le C_{\rm dec}\theta^{5/2}D_P(r;z_0)
     +C_{\rm dec}\theta^{-2}C_V(2r;z_0).
\]
\end{lemma}

\begin{proof}
Write \(z_0=(y_0,s_0)\).  For a.e. time
\(s\in(s_0-r^2,s_0]\), apply \cref{paper2:thm:pressure-decomp} on
\(B_r(y_0)\) with cutoff supported in \(B_{2r}(y_0)\):
\[
  P(\cdot,s)=P_{\rm loc}(\cdot,s)+H(\cdot,s),
  \qquad -\Delta H(\cdot,s)=0\quad\text{in }B_r(y_0).
\]
The Calder\'on--Zygmund part satisfies
\[
  \int_{B_{\theta r}(y_0)}|P_{\rm loc}|^{3/2}
  \le C\int_{B_{2r}(y_0)}|V|^3 .
\]
After integrating over the shorter time interval and multiplying by
\((\theta r)^{-2}\), this contributes at most
\(C\theta^{-2}C_V(2r;z_0)\).

For the harmonic part, the interior oscillation estimate for harmonic
functions gives, for \(q=3/2\),
\[
  \int_{B_{\theta r}(y_0)}
  |H-(H)_{B_{\theta r}}|^q
  \le C\theta^{3+q}
  \int_{B_r(y_0)}|H-(H)_{B_r}|^q .
\]
Since \(3+q=9/2\), scaling the cylinder by \((\theta r)^{-2}\) leaves the
factor \(\theta^{5/2}\).  Moreover
\[
  \int_{B_r}|H-(H)_{B_r}|^{3/2}
  \le C\int_{B_r}|P-(P)_{B_r}|^{3/2}
     +C\int_{B_r}|P_{\rm loc}|^{3/2},
\]
and the last term is again bounded by \(C\int_{B_{2r}}|V|^3\).  Absorbing the
resulting \(C\theta^{5/2}C_V(2r;z_0)\) into
\(C\theta^{-2}C_V(2r;z_0)\) gives the displayed estimate.  Spatial averages may
be replaced by the minimizing constants at the cost of a universal factor, as
in \cref{paper3:lem:scale-rigid-two-sided-ckn-transfer}.
\end{proof}

\begin{lemma}[Dyadic persistence of CKN concentration]
\label{paper3:lem:scale-rigid-dyadic-ckn-persistence}
Let \((V,P)\) be suitable on a compact cylinder \(Q_\rho\), and let
\[
  \mathcal C_{V,P}(z,r):=
  r^{-2}\int_{Q_r(z)}|V|^3
  +r^{-2}\int_{Q_r(z)}|P-(P)_{B_r}|^{3/2}.
\]
If \(z_*\in Q_\rho\) is a Caffarelli--Kohn--Nirenberg concentration point in
the sense that
\[
  \limsup_{r\downarrow0}\mathcal C_{V,P}(z_*,r)\ge\kappa_*>0,
\]
then for every \(0<\theta<1\) and every sufficiently small outer radius
\(r_0\) with \(Q_{2r_0}(z_*)\Subset Q_\rho\), there exist
\(z_m\to z_*\), integers \(L_m\to\infty\), and dyadic radii
\(\rho_m=\theta^{L_m}r_0\) such that
\[
  Q_{\rho_m}(z_m)\Subset Q_{r_0}(z_*),\qquad
  \mathcal C_{V,P}(z_m,\rho_m)\ge c_\theta\kappa_* ,
\]
where \(c_\theta>0\) depends only on \(\theta\).
\end{lemma}

\begin{proof}
Choose radii \(\sigma_m\downarrow0\) and points \(z_m\to z_*\) with
\(Q_{\sigma_m}(z_m)\Subset Q_{r_0}(z_*)\) and
\(\mathcal C_{V,P}(z_m,\sigma_m)\ge\kappa_*/2\).  Pick \(L_m\) so that
\[
  \theta^{L_m+1}r_0<\sigma_m\le \theta^{L_m}r_0=:\rho_m .
\]
Then \(L_m\to\infty\) and \(\rho_m\le\theta^{-1}\sigma_m\).  The velocity term
on \(Q_{\rho_m}(z_m)\) controls the velocity term on \(Q_{\sigma_m}(z_m)\) with
loss at most \(\theta^2\), because the normalization changes by
\((\sigma_m/\rho_m)^2\).  For the pressure term, spatial averages are compared
through the minimizing constant:
\[
  \int_{B_{\sigma_m}}|P-(P)_{B_{\sigma_m}}|^{3/2}
  \le 2^{3/2}\int_{B_{\sigma_m}}|P-c|^{3/2}
  \le 2^{3/2}\int_{B_{\rho_m}}|P-c|^{3/2}.
\]
Taking \(c=(P)_{B_{\rho_m}}\), integrating in time, and accounting for the same
normalization loss gives
\[
  \mathcal C_{V,P}(z_m,\rho_m)\ge 2^{-5/2}\theta^2
  \mathcal C_{V,P}(z_m,\sigma_m).
\]
Thus the conclusion holds with \(c_\theta=2^{-7/2}\theta^2\).
\end{proof}

\begin{lemma}[Exterior pressure persistence is a structural exit]
\label{paper3:lem:scale-rigid-exterior-pressure-routing}
Let a scale-rigid branch have a pressure lower bound on shrinking selected
cylinders that is not accounted for by the local source
\(\partial_i\partial_j(V_iV_j)\) in any compact selected frame after the
decomposition of \cref{paper2:thm:pressure-decomp}.  Then the branch enters the
separated/exterior structural alternative of
\cref{paper3:def:local-multibubble}\,(iv), and it does not remain in the
residual scale-rigid stratum.
\end{lemma}

\begin{proof}
The unaccounted part of the pressure is harmonic in the selected compact frame
and is generated by sources at positive distance in the corresponding physical
variables.  If the exterior velocity source has a terminal CKN concentration on
one of the annuli separating it from the selected core, then
\cref{paper8:thm:exterior-stratification} gives an exterior
concentration branch.  In the local variables this is a separated active
physical point, hence the separated-point alternative
\cref{paper3:def:local-multibubble}\,(iv) and the reduction
\cref{paper3:thm:separated-point-reduction} apply.

If no such exterior concentration exists on the separating annuli, then
\cref{paper8:thm:exterior-regularity-no-concentration} gives regular exterior
control and
\cref{paper8:lem:exterior-pressure-vanishes} gives convergence of the exterior
pressure contribution modulo functions of time on the selected compact frame.
This contradicts the assumed persistent pressure lower bound there.  Thus the
only non-vanishing exterior-pressure case is already a separated/exterior
structural exit.
\end{proof}

\begin{lemma}[Persistent CKN mass activates an \(L^3\) core or exits structurally]
\label{paper3:lem:scale-rigid-pressure-core-activation}
Let \((V_n,P_n)\) be a compact-cylinder scale-rigid sequence on \(Q_{2R}\) with
\[
  \sup_n\bigl(C_{V_n}(2R)+D_{P_n}(2R)\bigr)\le M_0,
\]
and assume the pressure decomposition of
\cref{paper2:thm:pressure-decomp} is available on \(Q_{2R}\).  Fix
\(\kappa_0>0\) and choose
\(\theta\in(0,1/8)\) so that
\[
  C_{\rm dec}\theta^{5/2}\le \frac14 .
\]
Then there are constants \(\eta_{\rm act}>0\) and \(L_0\in\mathbb N\),
depending only on \(M_0,R,\kappa_0,\theta\), the fixed-scale pressure bound
below, and the pressure-decomposition constants, with the following property.
Suppose \(0<R_0<R\), \(z_n\to z\),
\(Q_{2R_0}(z_n)\Subset Q_{2R}\) for all large \(n\), and suppose the fixed
outer scale used for the iteration has a uniform local pressure bound
\[
  D_{P_n}(R_0;z_n)\le D_0<\infty .
\]
Let \(L_n\to\infty\), \(\rho_n:=\theta^{L_n}R_0\), and assume
\[
  \mathcal C_n(\rho_n;z_n):=
  \rho_n^{-2}\int_{Q_{\rho_n}(z_n)} |V_n|^3
  +\rho_n^{-2}\int_{Q_{\rho_n}(z_n)}
       |P_n-(P_n)_{B_{\rho_n}}|^{3/2}
  \ge \kappa_0
\]
for all large \(n\).  Then, after passing to a subsequence, either:
\begin{enumerate}[label=\textup{(\alph*)}]
\item for some integers \(0\le j_n<L_n\), with
      \(s_n:=2\theta^{j_n}R_0\),
      \[
        s_n^{-2}\int_{Q_{s_n}(z_n)}|V_n|^3\ge\eta_{\rm act}^3 ,
      \]
      so the rescaled branch at scale \(s_n\) contains an active \(L^3\) core;
\item the retained pressure lower bound is exterior pressure persistence, hence
      the branch is assigned to the separated/exterior structural alternative by
      \cref{paper3:lem:scale-rigid-exterior-pressure-routing}.
\end{enumerate}
\end{lemma}

\begin{proof}
Choose \(L_0\) so large that
\[
  (1/4)^{L_0}D_0\le \kappa_0/16,
\]
and choose \(\eta_{\rm act}>0\) so small that
\[
  C_{\rm dec}\theta^{-2}\sum_{m=0}^{\infty}(1/4)^m\eta_{\rm act}^3
  \le \kappa_0/16,
  \qquad
  4\theta^{-2}\eta_{\rm act}^3\le \kappa_0/4 .
\]
For large \(n\), \(L_n\ge L_0\).  Suppose item \textup{(a)} fails.  Then for
each \(0\le j<L_n\),
\[
  C_{V_n}(2\theta^jR_0;z_n)<\eta_{\rm act}^3 .
\]
Iterating \cref{paper3:lem:scale-rigid-scaled-pressure-decay} from the fixed
outer scale \(R_0\) down to \(\theta^{L_n}R_0=\rho_n\) gives
\[
  D_{P_n}(\rho_n;z_n)
  \le (1/4)^{L_n}D_{P_n}(R_0;z_n)
     +C_{\rm dec}\theta^{-2}
      \sum_{j=0}^{L_n-1}(1/4)^{L_n-1-j}
      C_{V_n}(2\theta^jR_0;z_n)
  \le \kappa_0/8 .
\]
Here the initial pressure is bounded by \(D_0\) because the iteration starts
from this fixed compact scale, not from the shrinking target scale.
Moreover \(Q_{\rho_n}(z_n)\subset Q_{2\theta^{L_n-1}R_0}(z_n)\), so failure of
item \textup{(a)} at \(j=L_n-1\) gives
\[
  \rho_n^{-2}\int_{Q_{\rho_n}(z_n)}|V_n|^3
  \le 4\theta^{-2}\eta_{\rm act}^3\le \kappa_0/4 .
\]
Hence \(\mathcal C_n(\rho_n;z_n)<\kappa_0\), contradicting the assumed lower
bound, unless the pressure term used in the lower bound is not the local
pressure governed by the compact-frame decomposition.  The latter possibility
is exactly exterior pressure persistence and is item \textup{(b)} by
\cref{paper3:lem:scale-rigid-exterior-pressure-routing}.
\end{proof}

\begin{lemma}[Activated cores are captured by the finite active ledger]
\label{paper3:lem:scale-rigid-active-core-capture}
Let a scale-rigid branch inherit an active profile family in the sense of
\cref{paper3:def:active-family}, with finite critical mass \(M_*\) and
concentration capture.  Suppose
\cref{paper3:lem:scale-rigid-pressure-core-activation} produces represented
scales \(s_n>0\) and centers \(z_n\) such that
\[
  s_n^{-2}\int_{Q_{s_n}(z_n)}|V_n|^3\ge \eta_{\rm act}^3 .
\]
After lowering the active threshold once, if necessary, to
\(\eta_*^{\rm new}:=\min\{\eta_*,\eta_{\rm act}/2\}\), exactly one of the
following occurs after passing to a subsequence:
\begin{enumerate}[label=\textup{(\alph*)}]
\item the activated core is already represented by an existing active profile
      or compound profile in the inherited family;
\item it is absorbed into the selected terminal frame, so it does not create an
      independent residual obstruction;
\item it is assigned to the profile-completeness obligation rather than to the
      scale-rigid residual branch.
\end{enumerate}
In particular, under the stated concentration-capture hypothesis, pressure-core
activation cannot create an unregistered active object inside the residual
scale-rigid branch.
\end{lemma}

\begin{proof}
Rescale the branch around \((z_n,s_n)\) inside the represented window; in
physical variables the corresponding scale is the selected physical scale
multiplied by \(s_n\), hence still tends to zero along the Type II sequence.
The displayed lower bound is exactly the active-frame condition in
\cref{paper3:def:active-frame} on \(Q_1\), with threshold \(\eta_{\rm act}\).
The concentration-capture clause in
\cref{paper3:def:active-family} says that a compact terminal-window sequence
with positive \(L^3\) mass is either seen by the active family or is absorbed
into the selected terminal frame; it cannot remain in the perturbative
remainder.  Hence alternatives \textup{(a)} and \textup{(b)} hold unless the
frame is genuinely absent from the captured profile family.  That last
possibility is not a scale-rigid residual mechanism; it is precisely failure of
the terminal active-family completeness/capture input and is assigned to that
separate profile-completeness obligation.  Therefore the scale-rigid refinement
below only uses active objects that already belong to the finite captured
family, together with objects absorbed into the selected terminal frame.
\end{proof}

\begin{lemma}[No sub-concentration gives a bounded selected-window limit]
\label{paper3:lem:scale-rigid-no-subconcentration-bounded}
Let \((V_n,P_n,a_n,b_n)\) be a scale-rigid local branch and let
\((V_*,P_*)\) be a compact limit supplied by
\cref{paper3:lem:scale-rigid-pressure-modulation-stability} on
\(Q_\rho\).  Then, after shrinking once to a compact subcylinder
\(Q_{\rho'}\Subset Q_\rho\), exactly one of the following alternatives holds:
\begin{enumerate}[label=\textup{(\alph*)}]
\item \((V_*,P_*)\) has no Caffarelli--Kohn--Nirenberg concentration point in
      \(Q_{\rho'}\).  Then \(V_*\in L^\infty(Q_{\rho''})\) for every
      \(Q_{\rho''}\Subset Q_{\rho'}\), and the selected subsequence has a
      nonzero bounded selected-window limit on some such \(Q_{\rho''}\).
\item \((V_*,P_*)\) has a Caffarelli--Kohn--Nirenberg concentration point in
      \(Q_{\rho'}\).  Then the original branch either contains a further
      captured active \(L^3\) concentration core after rescaling around that
      point, the pressure mass is exterior and the branch is already in the
      separated/exterior structural alternative, or the core is assigned to the
      separate profile-completeness obligation.  According to its relative
      scale-center position, every captured core enters one of the already named
      multicore, same-point cascade, or separated/exterior structural
      alternatives, rather than remaining a residual scale-rigid branch.
\end{enumerate}
\end{lemma}

\begin{proof}
Let \(\varepsilon_{\rm rg}\) be the represented-gauge regularity threshold from
\cref{paper3:lem:scale-rigid-ckn-in-gauge}.  Before using that threshold,
apply \cref{paper3:lem:scale-rigid-chart-comparability} on each retained
compact subcylinder.  The chart stability in \cref{paper2:prop:stability},
together with the \(L^\infty\) modulation bound (R3), gives a limiting chart
\((\Lambda_*,X_*)\) with \(a_*=\Lambda_*'/\Lambda_*\) and
\(b_*=X_*'/\Lambda_*\) on the retained
subwindow.  Thus the compactness of the subcylinder and (R3) give a uniform
lower-upper scale comparison after one shrink, so the constants in
\cref{paper3:lem:scale-rigid-ckn-in-gauge} depend only on the retained chart
geometry and not on any unproved boundedness of \(V_n\) or \(V_*\).

Assume first that there is no CKN concentration point of \((V_*,P_*)\) in
\(Q_{\rho'}\).  For every point \(z\in Q_{\rho'}\) there is a parabolic
cylinder \(Q_{r_z}(z)\Subset Q_\rho\) such that the scale-invariant
velocity-pressure quantity of \((V_*,P_*)\) on \(Q_{r_z}(z)\) is below
\(\varepsilon_{\rm rg}\).  A finite subcover of each
\(Q_{\rho''}\Subset Q_{\rho'}\), followed by
\cref{paper3:lem:scale-rigid-ckn-in-gauge} on the cover, gives
\(V_*\in L^\infty(Q_{\rho''})\).

It remains to check that the bounded limit is nonzero.  If \(V_*\equiv0\) on
the retained inner cylinder, then the strong \(L^3\) convergence from
\cref{paper3:lem:scale-rigid-pressure-modulation-stability} and the
zero-velocity argument of \cref{paper1:lem:zero-contradiction} imply
\[
  C_{V_n}(r_*)\to0,
\]
contradicting the retained lower bound (R2) in
\cref{paper3:def:scale-rigid-branch}.  Hence \(V_*\not\equiv0\) on some
smaller compact cylinder.  Since \(V_n\to V_*\) strongly in \(L^3\) there,
the selected subsequence has a nonzero bounded selected-window limit.

Now suppose a CKN concentration point \(z_*\) remains in \(Q_{\rho'}\).  Then
there is \(\kappa_*>0\) such that
\[
  \limsup_{r\downarrow0}\mathcal C_{V_*,P_*}(z_*,r)\ge \kappa_* .
\]
Choose \(R_0>0\) with \(Q_{2R_0}(z_*)\Subset Q_\rho\), and choose
\(\theta\in(0,1/8)\) satisfying the pressure-decay smallness condition in
\cref{paper3:lem:scale-rigid-pressure-core-activation}.  By
\cref{paper3:lem:scale-rigid-dyadic-ckn-persistence}, there are
\(z_m\to z_*\), \(L_m\to\infty\), and
\(\rho_m=\theta^{L_m}R_0\) such that
\[
  \mathcal C_{V_*,P_*}(z_m,\rho_m)\ge c_\theta\kappa_* .
\]
For each fixed \(m\), the strong \(L^3\) convergence of \(V_n\) and weak lower
semicontinuity of the pressure oscillation seminorm in
\(L^{3/2}/\{\text{spatial constants}\}\) imply, along a diagonal subsequence,
\[
  \mathcal C_{V_n,P_n}(z_m,\rho_m)\ge \frac12 c_\theta\kappa_* .
\]
Relabeling the diagonal sequence, apply
\cref{paper3:lem:scale-rigid-pressure-core-activation} with
\(\kappa_0=\frac12c_\theta\kappa_*\) on the compact cylinder where (R1) holds.
The additional fixed-scale pressure bound required by that lemma holds because
\(R_0\) is fixed, \(Q_{2R_0}(z_m)\Subset Q_\rho\) for all large \(m\), and the
mean-normalized pressures are uniformly bounded in \(L^{3/2}(Q_\rho)\) by
\cref{paper3:lem:scale-rigid-pressure-modulation-stability}; comparison of
spatial averages on the fixed nested balls gives
\(D_{P_n}(R_0;z_m)\le D_0(R_0,\rho,M)\).
If the persistent CKN mass is produced by an interior velocity core, then
\cref{paper3:lem:scale-rigid-active-core-capture} either identifies that core
with the finite captured active ledger or assigns it to the separate
profile-completeness obligation, after the one-time lowering of the active
threshold.  If
the mass is carried by harmonic pressure generated outside the selected compact
frame, \cref{paper3:lem:scale-rigid-exterior-pressure-routing} assigns the
branch to the separated/exterior structural alternative, so it exits the
residual scale-rigid stratum.

The relative position of this new core with respect to the retained scale is
exhausted by \cref{paper3:lem:parameter-partition}: it is comparable at the
same point, strictly separated in scale at the same point, separated in
physical center, or locally invisible in the retained frame.  The invisible
case cannot carry the retained positive concentration by the capture clause in
\cref{paper3:def:active-family}.  The remaining cases are precisely the
multicore, same-point cascade, and separated/exterior structural alternatives
already named in the manuscript.  Thus the branch does not remain in the
residual scale-rigid stratum in this case.
\end{proof}

\begin{theorem}[Single-core reduction compatibility]
\label{paper3:thm:scale-rigid-single-core-compatibility}
Let \((V_n,P_n,a_n,b_n)\) be a scale-rigid local branch in the sense of
\cref{paper3:def:scale-rigid-branch}, and assume the branch reduces, after
the rigid-frame limit of
\cref{paper3:lem:scale-rigid-pressure-modulation-stability}, to a single
retained nondegenerate core on a fixed inner cylinder
\(Q_{r_*}\Subset Q_{\rho}\Subset Q_{R_*}\).  Then, after passing to a
subsequence, one of the following alternatives holds:
\begin{enumerate}[label=\textup{(\roman*)}]
\item the branch has a nonzero bounded selected-window limit;
\item the original selected sequence has vanishing localized dissipation on a
      retained compact cylinder, hence is a compact single-core
      vanishing-dissipation branch in the sense of
      \cref{paper1:def:single-core-branch} and is excluded by
      \cref{paper1:thm:single-core-criterion};
\item a further active concentration core is produced, so the branch leaves the
      single-core reduction and enters a multicore, cascade, or
      separated/exterior structural alternative.
\end{enumerate}
\end{theorem}

\begin{proof}
If the original selected sequence has vanishing localized dissipation on some
\(Q_{\rho'}\), \(Q_{r_*}\Subset Q_{\rho'}\Subset Q_{R_*}\), then we first verify
\cref{paper1:def:single-core-branch}.  Suitability and the represented
equation are supplied by \cref{paper2:thm:C}.  Compact-cylinder bounds,
positive selected-core CKN concentration, and bounded modulation are exactly
(R1), (R2), and (R3) of \cref{paper3:def:scale-rigid-branch}.  The
nondegenerate scale-center selection is the hypothesis of this single-core
case.  The vanishing localized dissipation is the present hypothesis on the
original selected sequence, not on a residual after subtracting a limit.  Thus
the branch satisfies every clause of
\cref{paper1:def:single-core-branch}, and
\cref{paper1:thm:single-core-criterion} excludes it.  This is item
\textup{(ii)}.

The compact limit on a slightly smaller cylinder is covered by
\cref{paper3:lem:scale-rigid-no-subconcentration-bounded}.  If the
no-subconcentration alternative holds, item \textup{(i)} follows.  If a further
active concentration core is produced, item \textup{(iii)} follows.
\end{proof}

\begin{lemma}[Degenerate Jacobian resolution]
\label{paper3:lem:scale-rigid-degenerate-jacobian}
Let \((V_n,P_n,a_n,b_n)\) be a represented sequence reaching the
gauge-degenerate alternative \cref{paper3:def:local-multibubble}\,(v),
i.e.\ along the subsequence the relative scale-center Jacobian of the
selected compound or terminal frame degenerates.  Then, after passing to a
further subsequence and reselecting the active center in the sense of
\cref{paper3:def:active-family} on the same physical point, exactly one of
the following four resolutions occurs:
\begin{enumerate}[label=\textup{(D\arabic*)}]
\item the reselection produces an additional active core or active compound
      carrying positive \(L^3\) mass at the fixed active threshold; hence the
      branch enters the multicore, same-point cascade, or separated/exterior
      alternative
      \cref{paper3:def:local-multibubble}\,(ii)--(iv) and is treated by
      \cref{paper3:thm:minimal-bad-extraction};
\item the reselection collapses the compound to a single nondegenerate core
      and the branch enters the single-core compatibility dichotomy governed by
      \cref{paper3:thm:scale-rigid-single-core-compatibility};
\item the reselection produces a perturbative compound that is below the
      critical small-data threshold of Step~2 of
      \cref{paper3:thm:same-point-reduction}, hence is absorbed into the
      perturbative remainder and treated by
      \cref{paper3:thm:terminal-decoupling};
\item the reselection still has degenerate scale-center Jacobian on every
      further subsequence and has no further active subconcentration; in this
      case the no-subconcentration alternative of
      \cref{paper3:lem:scale-rigid-no-subconcentration-bounded} gives a
      nontrivial bounded selected-window limit on a compact cylinder.
\end{enumerate}
\end{lemma}

\begin{proof}
By the gauge-degeneracy alternative
\cref{paper3:def:local-multibubble}\,(v), the localized scale or centering
selection of the selected compound or terminal frame is non-unique along the
subsequence.  Apply the parameter partition lemma
\cref{paper3:lem:parameter-partition} to the selected compound together with
any candidate alternate selection produced by the gauge degeneracy.  Three
mutually exclusive structural outcomes are possible.

\noindent\textit{Case A: alternate selection is comparable to the original at
the same physical point.}  Then both selections lie in the same scale
comparability class of Step~1 of \cref{paper3:thm:same-point-reduction}.
Step~2 of that theorem groups them into a single compound profile.  If the
compound is below the small-data threshold, we are in case (D3); otherwise
the compound is above threshold and Step~3 of the theorem applies.  If the
compound is then nondegenerate, we are in case (D2); if it is still
degenerate, the gauge degeneracy persists on the reselected compound, and we
proceed to Case~D below.

\noindent\textit{Case B: alternate selection is at a strictly different
scale at the same physical point.}  By
\cref{paper3:lem:outer-vanish}, the strictly outer profile vanishes locally
in the innermost frame; absorbing it into the perturbative remainder lands
us in case (D3).  If instead the reselection produces a strictly inner core
not previously detected, then either compact-cylinder suitable control fails,
in which case \cref{paper3:lem:compactness-failure} produces the new CKN core,
or compact control persists and
\cref{paper3:lem:scale-rigid-pressure-core-activation} converts the retained
CKN mass into an active \(L^3\) core unless the branch has already exited to
the separated/exterior pressure alternative.  These outcomes are case (D1).

\noindent\textit{Case C: alternate selection is at a separated physical
point.}  By \cref{paper3:thm:separated-point-reduction}, the new selection
is locally invisible in the original frame, hence absorbed into the
perturbative remainder; this is case (D3).  If the new selection retains
positive CKN concentration in its own frame, then
\cref{paper3:lem:scale-rigid-pressure-core-activation} makes it an additional
separated active core, or else assigns it to the separated/exterior pressure
alternative.  Both outcomes are case (D1) via
\cref{paper3:def:local-multibubble}\,(iv).

\noindent\textit{Case D: persistent degeneracy.}  Suppose every reselection
on every further subsequence still has degenerate scale-center Jacobian.
Apply \cref{paper3:lem:scale-rigid-pressure-modulation-stability} and then
\cref{paper3:lem:scale-rigid-no-subconcentration-bounded}.  If the compact
limit has a CKN subconcentration point, that lemma produces an additional
active core; this is case (D1).  If no such subconcentration remains, the same
lemma gives a nonzero bounded selected-window limit; this is case (D4).  No
\(L^\infty\) bound is inferred from the Jacobian degeneracy itself.
\end{proof}

\begin{theorem}[Scale-rigid stratum decomposition]
\label{paper3:thm:scale-rigid-stratum-decomposition}
Every scale-rigid local branch in the sense of
\cref{paper3:def:scale-rigid-branch} enters one of the following named local
strata after passing to a subsequence:
\begin{enumerate}[label=\textup{(\roman*)}]
\item the multicore alternative \cref{paper3:def:local-multibubble}\,(ii)--(iv),
      governed by \cref{paper3:thm:minimal-bad-extraction};
\item the compact single-core stratum, governed by
      \cref{paper3:thm:scale-rigid-single-core-compatibility};
\item the perturbative remainder of \cref{paper3:def:active-family}, governed
      by \cref{paper3:thm:terminal-decoupling};
\item the bounded selected-window limit alternative
      \cref{paper1:def:bounded-window-limit}, which by
      \cref{paper3:lem:scale-rigid-bounded-limit-exit} is incompatible with
      \cref{paper1:def:typeII}.
\end{enumerate}
\end{theorem}

\begin{proof}
By \cref{paper3:def:scale-rigid-branch}\,(R4), the branch enters the
scale-rigid stratum either via a degenerate relative scale-center Jacobian or
as an innermost cascade core after outer profiles have been absorbed
perturbatively.

\noindent\textit{Degenerate Jacobian case.}  Apply
\cref{paper3:lem:scale-rigid-degenerate-jacobian}.  Cases (D1)--(D4) of that
lemma give items (i)--(iv) respectively.

\noindent\textit{Innermost cascade case.}  By Step~4 of
\cref{paper3:thm:same-point-reduction}, after subtracting the outer-profile
sum \(S_n\to0\) in \(L^\infty_I L^3(B_{2R})\) and absorbing constant drifts
into the translation modulation, the residual is governed by
\cref{paper3:thm:terminal-decoupling}.  That theorem produces three
outcomes already enumerated as (a) single-core, (b) scale-rigid (degenerate
Jacobian), or (c) purely perturbative.  Case (a) is item (ii) here via
\cref{paper3:thm:scale-rigid-single-core-compatibility}; case (c) is item
(iii); case (b) loops back into the degenerate Jacobian case treated above
and is therefore reduced by
\cref{paper3:lem:scale-rigid-degenerate-jacobian} to one of (i)--(iv).
\end{proof}

\begin{lemma}[Finite refinement of scale-rigid structural exits]
\label{paper3:lem:scale-rigid-refinement-terminates}
Let \(\omega\) be a scale-rigid local branch produced by
\cref{paper3:thm:same-point-reduction} or
\cref{paper3:thm:separated-point-reduction}, with the inherited active family
of finite critical mass and with concentration capture in the sense of
\cref{paper3:def:active-family}.  Resolve every structural exit of
\cref{paper3:thm:scale-rigid-stratum-decomposition} by applying only
\cref{paper3:thm:minimal-bad-extraction},
\cref{paper3:thm:same-point-reduction},
\cref{paper3:thm:separated-point-reduction}, and
\cref{paper3:thm:terminal-decoupling}, never
\cref{paper3:thm:multi}.  Then the refinement terminates after finitely many
ledger-decreasing refinements.  At the terminal residual branch one has either a nonzero
bounded selected-window limit or a compact single-core
vanishing-dissipation branch excluded by
\cref{paper1:thm:single-core-criterion}.
\end{lemma}

\begin{proof}
After the active threshold has been lowered once, every core produced by
\cref{paper3:lem:scale-rigid-pressure-core-activation} is handled by
\cref{paper3:lem:scale-rigid-active-core-capture}: it is already in the
captured inherited family, is absorbed into the terminal frame, or exits to the
separate profile-completeness obligation.  The residual scale-rigid ledger
therefore uses only the captured active family.  Write
\(\eta_*^{\rm new}\) for the lowered active threshold.  By
\cref{paper3:lem:finite-active-count}, the resulting family contains at most
\[
  N_*:=\lfloor M_*(\eta_*^{\rm new})^{-3}\rfloor
\]
active profiles.  Group same-point comparable-scale profiles into compound
profiles as in \cref{paper3:def:compound-profile}.  There are therefore only
finitely many active objects and only finitely many unordered pair relations
among them.

Define the refinement ledger to consist of:
\begin{enumerate}[label=\textup{(L\arabic*)}]
\item every active object not yet selected, grouped, absorbed, or declared
      terminal;
\item every pair relation not yet classified as comparable same-point, strict
      same-point cascade, separated point, or locally invisible;
\item every degenerate scale-center selection not yet resolved by
      \cref{paper3:lem:scale-rigid-degenerate-jacobian};
\item every perturbative remainder not yet removed by
      \cref{paper3:thm:terminal-decoupling}.
\end{enumerate}
The ledger is finite because it is built from a finite active family and
finitely many compounds.  Order ledgers lexicographically by the number of
unresolved entries in \textup{(L1)}, \textup{(L2)}, \textup{(L3)}, and
\textup{(L4)}.

At each nonterminal step, no multibubble exclusion theorem is used.  If
\cref{paper3:thm:scale-rigid-stratum-decomposition} gives a multicore,
same-point cascade, or separated/exterior structural exit, apply
\cref{paper3:thm:minimal-bad-extraction} to select the corresponding active
object.  Comparable same-point objects are grouped into a compound profile;
below-threshold compounds are removed by the small-data step in
\cref{paper3:thm:same-point-reduction}, while above-threshold compounds count
as active objects.  Strict same-point cascades are reduced by
\cref{paper3:thm:same-point-reduction}; separated physical centers are reduced
by \cref{paper3:thm:separated-point-reduction}.  Perturbative exterior and
remainder terms are discarded only through
\cref{paper3:thm:terminal-decoupling}, whose output preserves the selected
positive concentration and produces no new active object.

Each move strictly decreases the ledger.  Adding a previously unseen active
core consumes one entry of \textup{(L1)}, and this can happen at most \(N_*\)
times.  Grouping comparable same-point objects, ordering a same-point cascade,
separating two physical centers, or declaring one object locally invisible
resolves a previously unresolved pair and consumes one entry of \textup{(L2)}.
Applying \cref{paper3:lem:scale-rigid-degenerate-jacobian} resolves the
selected degenerate frame or produces an active structural exit already counted
in \textup{(L1)}--\textup{(L2)}, so it consumes one entry of \textup{(L3)}.  Applying
\cref{paper3:thm:terminal-decoupling} to a perturbative remainder consumes one
entry of \textup{(L4)}.  Once an entry is resolved it is not reintroduced: later
passes use the grouped compound, the selected innermost representative, the
separated-frame assignment, or the discarded perturbative remainder.  Hence no
cycle is possible.

After the ledger is exhausted, item \textup{(i)} of
\cref{paper3:thm:scale-rigid-stratum-decomposition} cannot occur again, because
it would create either a new active object or an unresolved active relation,
contradicting exhaustion of \textup{(L1)}--\textup{(L2)}.  Item \textup{(iii)}
cannot carry the retained positive concentration by the capture clause in
\cref{paper3:def:active-family}; a purely perturbative remainder has already
been absorbed through \textup{(L4)}.
Item \textup{(iv)} is exactly the nonzero bounded selected-window-limit
alternative.

It remains only to handle item \textup{(ii)}.  Apply
\cref{paper3:thm:scale-rigid-single-core-compatibility}.  Its bounded-limit
alternative gives the first terminal outcome.  Its further-active-core
alternative would add a new ledger entry, impossible after ledger exhaustion.
Its vanishing-localized-dissipation alternative verifies
\cref{paper1:def:single-core-branch} for the original selected sequence and is
excluded by \cref{paper1:thm:single-core-criterion}.  These are precisely the
two terminal residual outcomes claimed.
\end{proof}

\begin{proposition}[Scale-rigidity exclusion]
\label{paper3:prop:scale-rigidity-discharged}
Every scale-rigid local branch produced by
\cref{paper3:thm:same-point-reduction} or
\cref{paper3:thm:separated-point-reduction}, including a branch with
degenerate relative scale-center Jacobian, either has a nonzero bounded
selected-window limit on a compact cylinder (and so fails
\cref{paper1:def:typeII} by
\cref{paper3:lem:scale-rigid-bounded-limit-exit}) or reduces to the compact
single-core branch governed by
\cref{paper1:thm:single-core-criterion} (equivalently by the local
proposition \cref{paper3:prop:local-single-core-exclusion}).  Consequently no such branch is
compatible with the local Type II condition.  This proves the earlier
scale-rigidity exclusion statement recorded as
\cref{paper3:prop:local-scale-rigidity-exclusion} in the stated class.
\end{proposition}

\begin{proof}
Let \(\omega\) be a scale-rigid local branch.
\Cref{paper3:thm:scale-rigid-stratum-decomposition} places \(\omega\) into
one of items (i)--(iv).

If item (i) occurs, the branch enters a named multicore, same-point cascade, or
separated/exterior structural exit.  Resolve that exit by the finite
refinement procedure of
\cref{paper3:lem:scale-rigid-refinement-terminates}.  This procedure uses only
\cref{paper3:thm:minimal-bad-extraction},
\cref{paper3:thm:same-point-reduction},
\cref{paper3:thm:separated-point-reduction}, and
\cref{paper3:thm:terminal-decoupling}; it does not invoke
\cref{paper3:thm:multi}, so there is no circular use of
\cref{paper3:prop:local-scale-rigidity-exclusion}.  The refinement terminates at a bounded
selected-window limit or at the compact single-core
vanishing-dissipation branch.

In item (ii), apply
\cref{paper3:thm:scale-rigid-single-core-compatibility}.  Its bounded-limit
case is the first alternative in the proposition.  Its
vanishing-localized-dissipation case verifies
\cref{paper1:def:single-core-branch} for the original selected sequence and is
excluded by \cref{paper1:thm:single-core-criterion}; this is the
``reduces to the single-core branch'' alternative.  Its further-active-core
case is resolved by the same finite refinement lemma used for item (i).

In item (iii), the branch is the perturbative remainder; by
\cref{paper3:thm:terminal-decoupling} it produces no independent
multibubble obstruction.  The retained positive concentration is carried by
the selected component, not by the perturbative remainder, so the residual
selected component is reduced to item (ii) or item (iv) by
\cref{paper3:lem:scale-rigid-refinement-terminates}.

In item (iv), by
\cref{paper3:lem:scale-rigid-bounded-limit-exit} the branch is not a local
Type II sequence in the sense of \cref{paper1:def:typeII}.  This is the
``nonzero bounded selected-window limit'' alternative.

In every case the branch either has a nonzero bounded selected-window limit
or reduces to the single-core branch.  This is the conclusion of
\cref{paper3:prop:local-scale-rigidity-exclusion}.
\end{proof}

\begin{proposition}[Terminal admissibility in active frames]\label{paper3:prop:terminal-admissibility-active-frames}
Assume finite local critical norm, terminal windowed state-space completeness in
the sense of \cref{paper8:thm:terminal-windowed-profile-completeness}, local
perturbative residual stability, the exterior annulus alternative
\cref{paper8:cor:exterior-annulus-alternative}, the repaired-gauge representation
of \cref{paper2:thm:C}, and local Caccioppoli regularity.  Then every terminal
family produced by the same-point and separated-point reductions is admissible
on each compact selected cylinder in the sense of
\cref{paper3:def:admissible-terminal}.
\end{proposition}

\begin{proof}
The terminal nonlinear local-state decoupling theorem,
\cref{paper8:thm:terminal-nonlinear-decoupling}, applies under the stated
hypotheses and produces, on every compact terminal cylinder, the conclusions
of \cref{paper8:def:terminal-nonlinear-decoupling}.  We verify each clause of
\cref{paper3:def:admissible-terminal}.

Items (i) and (ii) are the active-family enumeration of
\cref{paper3:def:active-family} together with the same-point and
separated-point grouping of \cref{paper3:lem:separated-vanish,paper3:lem:outer-vanish};
they are properties of the active family that survive the terminal-frame
construction by definition of the terminal active class.

Item (iii) is the perturbative remainder hypothesis of
\cref{paper3:def:admissible-terminal}.  For the admissible active families
constructed here, the velocity decoupling
\cref{paper8:lem:terminal-velocity-decoupling}, the linear residual decoupling
\cref{paper8:lem:terminal-linear-residual}, and the pressure decoupling
\cref{paper8:lem:terminal-pressure-decoupling} supply the required
\(L^\infty_I L^3\) vanishing of the rescaled remainder \(r_n\) together with
the existence of \((\pi_n^r,f_n)\) having the stated
\(L^1_IH^{-1}(B_R)\) decay.  This pointwise critical hypothesis is not
derived from the later spacetime compactness package.

Item (iv) is the companion \(L^\infty_I L^3(B_{2R})\) hypothesis on the
selected component in \cref{paper3:def:admissible-terminal}.  On the active
families produced in
\cref{paper8:lem:terminal-scale-collapse-admissibility} it is read from the
retained local-state bounds; if a physical branch fails this input on the
chosen compact cylinder, then the branch is routed by
\cref{paper6a:lem:terminal-local-critical-mass-routing} rather than treated by
the perturbative terminal-frame theorem.

Item (v) is the bounded-modulation hypothesis, supplied by the repaired-gauge
representation \cref{paper2:thm:C} and \cref{paper2:thm:B}.

Item (vi) is the local pressure-decomposition stability of
\cref{paper8:lem:terminal-pressure-decoupling}, whose proof relies only on the
Calder\'on--Zygmund pressure-tail decoupling
\cref{paper8:lem:pressure-kernel-tail} together with the kernel boundedness on
\(L^{3/2}\) provided by the classical Calder\'on--Zygmund theorem
\cite{CalderonZygmund1952,Stein1970}.

Item (vii) follows from the Leibniz formula
\[
  [\Delta,\chi_R]u
  =
  (\Delta\chi_R)u+2\nabla\chi_R\cdot\nabla u
\]
and the identity
\[
  [\nabla\chi_R,\nabla\cdot](u\otimes u)
  =
  (\nabla\chi_R)(u\otimes u).
\]
These are combined with the dual estimate
\[
  \|fu\|_{H^{-1}(B_R)}
  \le C_R\|f\|_{C^1(\overline B_R)}\|u\|_{L^{6/5}(B_R)}
\]
and the embedding
\(L^3(B_R)\hookrightarrow L^{6/5}(B_R)\) on bounded balls.  Indeed, for
\(\varphi\in H^1_0(B_R)\),
\[
  \langle\nabla\chi_R\cdot\nabla u,\varphi\rangle
  =-\langle u,(\Delta\chi_R)\varphi+\nabla\chi_R\cdot\nabla\varphi\rangle,
\]
giving \(\|[\Delta,\chi_R]u\|_{H^{-1}(B_R)}\le C_R\|u\|_{L^3(B_{2R})}\).
The pointwise bound
\(|(\nabla\chi_R)(u\otimes u)|\le\|\nabla\chi_R\|_\infty|u|^2\) gives
\(\|(\nabla\chi_R)(u\otimes u)\|_{L^{3/2}(B_R)}\le C_R\|u\|_{L^3(B_{2R})}^2\),
and \(L^{3/2}\hookrightarrow H^{-1}\) on \(B_R\) supplies the displayed bound.
The combined estimate
\[
  \|[\Delta,\chi_R]u\|_{H^{-1}(B_R)}
  +\|[\nabla\chi_R,\,\nabla\cdot](u\otimes u)\|_{H^{-1}(B_R)}
  \le C_R\bigl(\|u\|_{L^3(B_{2R})}+\|u\|_{L^3(B_{2R})}^2\bigr)
  \le C_R'\|u\|_{L^3(B_{2R})}^2
\]
holds whenever \(\|u\|_{L^3(B_{2R})}\) is bounded below by a fixed constant,
which is the case in the active selected frame.  The same estimate applies to
the \(S_n\)-quadratic commutator terms thanks to the uniform
\(L^\infty_I L^3(B_{2R})\) control of \(U_n\) and the
\(L^1_I L^{3/2}(B_R)\) control of the cross stresses produced by
\cref{paper8:lem:terminal-stress-decoupling}.

All items in \cref{paper3:def:admissible-terminal} are therefore satisfied,
which is the asserted admissibility.
\end{proof}

\begin{theorem}[Local multibubble and cascade exclusion]\label{paper3:thm:multi}
Use the local single-core exclusion
\cref{paper3:prop:local-single-core-exclusion}, the discharged scale-rigidity
closure \cref{paper3:prop:scale-rigidity-discharged}, and the terminal admissibility in
active frames \cref{paper3:prop:terminal-admissibility-active-frames}.  Let a positive local
Type II branch admit an active family satisfying the hypotheses of
\cref{paper3:thm:minimal-bad-extraction}.  Then no local multibubble, cascade, or
gauge-degenerate Type II case can occur in this class.
\end{theorem}

\begin{proof}
Let a local Type II branch be in the multibubble/cascade class of
\cref{paper3:def:local-multibubble}.  Apply \cref{paper3:thm:minimal-bad-extraction}.  The
positive local concentration is carried by an active frame or compound local state,
and one of the alternatives listed there occurs.
\par
\textit{Case 1: compactness-failure reduction.}
Compactness failure yields, after changing radius by a fixed factor and
rescaling as in \cref{paper3:lem:compactness-failure}, a bounded-cylinder concentration
branch.  Relabel this branch as an active core.  After passing to a subsequence,
the new core and every previously retained active core have one of three
relations: comparable scale at the same physical point, strictly separated scale
at the same physical point, or separated physical center.  Comparable same-point
cores are grouped into compound local states.  A compound state with no
detecting compact cylinder is removed by local perturbative residual stability;
an active compound state is treated as one
selected core in its active frame.  If the scale-center selection inside such a
compound core is degenerate, the branch is gauge-degenerate and belongs to the
non-uniqueness alternative in \cref{paper3:def:local-multibubble}, rather than to
a separate core dynamics.
\par
\textit{Case 1a: gauge degeneracy.}
If the scale-center selection of the selected compound or terminal frame is
degenerate, then the branch is in the gauge-degenerate alternative of
\cref{paper3:def:local-multibubble}.  By the repaired-gauge representation
from \cref{sec:repaired-gauge-construction}, nondegeneracy is the hypothesis needed to enter the
single-core branch.  Persistent degeneracy therefore belongs to the
scale-rigid local row eventually discharged by \cref{paper3:prop:scale-rigidity-discharged}; if the degeneracy
disappears after passing to a subsequence, the branch enters the single-core or
terminal-frame alternatives below.
\par
\textit{Case 1b: same-point comparable scales.}
Same-point comparable-scale active frames are grouped into a compound local
state by the local state-space selection rule.  If the compound state has no
detecting compact cylinder, it is part of the perturbative residual.  If it has
the local detection threshold and its scale-center selection is nondegenerate,
it is one selected active core and enters the single-core branch.  If the
corresponding scale-center selection is degenerate, it is covered by Case 1a.
Thus comparable same-point states do not produce an additional multibubble
mechanism.
\par
\textit{Case 2: same-point cascading.}
Strict same-point scale cascades are treated by
\cref{paper3:thm:same-point-reduction}.  In the innermost frame, all outer scales are
locally invisible after constant drifts are absorbed.  The branch therefore
reduces to one of the three terminal alternatives
\[
  \text{(A) single-core},\quad \text{(B) scale-rigid},\quad \text{(C) terminal perturbative}.
\]
\par
\textit{Case 3: separated-point interactions.}
For separated physical points, \cref{paper3:thm:separated-point-reduction} gives the same
three terminal alternatives as in Case 2:
\[
  \text{(A) single-core},\quad \text{(B) scale-rigid},\quad \text{(C) terminal perturbative}.
\]

The single-core alternative is excluded by \cref{paper3:prop:local-single-core-exclusion}.  The
scale-rigid alternative is excluded by \cref{paper3:prop:scale-rigidity-discharged}, which includes
the scale-collapse and autonomous-limit exclusions in \cref{sec:scale-collapse-cost} in this setting.  The
terminal perturbative alternative gives no independent multibubble obstruction:
\cref{paper3:thm:terminal-decoupling} removes all exterior profiles and remainders from
the selected active equation in \(L^1H^{-1}_{\mathrm{loc}}\), preserves the positive
finite-mass selected branch, and leaves only the single-core or scale-rigid
possibilities already excluded.  Consequently every alternative in
\cref{paper3:def:local-multibubble} is incompatible with the local Type II regime.
\end{proof}

\begin{corollary}[Removal of the multibubble alternative]\label{paper3:cor:multi-removal}
Under the hypotheses of \cref{paper3:thm:multi}, every failure of the single
nondegenerate compact core alternative is excluded.
\end{corollary}

\begin{proof}
By \cref{paper3:def:local-multibubble}, failure of the single nondegenerate compact
core alternative means that at least one of the following occurs: the
compact-cylinder suitable estimates are unbounded; positive local CKN mass splits
between active cores; a strict same-point scale cascade occurs; separated
physical points remain active in distinct local frames; or the localized scale
or center selection is degenerate.  The compactness-failure case produces an
additional core by \cref{paper3:lem:compactness-failure}.  Same-point and separated
cases are reduced by \cref{paper3:thm:same-point-reduction,paper3:thm:separated-point-reduction}.
\Cref{paper3:thm:multi} excludes the remaining cases.  Hence no failure of the
single nondegenerate compact core alternative remains.
\end{proof}

\clearpage
\section{Local Caccioppoli and Rough-Core Exclusion}\label{sec:rough-core-exclusion}
\subsection{Analytic inputs and notation}

The renormalized local energy inequality from the preceding renormalized-chart
analysis is used as an analytic input.  The renormalized suitable structure and
pressure reconstruction are not rederived in this section.  The estimates below
are local implications under the displayed compact-cylinder conditions; the
terminal assembly verifies or routes these conditions by the ordered local
state-space analysis.

\begin{proposition}[Renormalized local energy inequality]\label{paper4:prop:renormalized-local-energy-inequality}
Let $(V,P,a,b)$ be defined on a compact renormalized cylinder
$Q_{2R,J}:=B_{2R}\times J$, where $J\subset\mathbb R$ is a bounded interval.
Assume
\begin{equation}\label{paper4:eq:renorm-ns}
\partial_\tau V+(V\cdot\nabla)V+\nabla P-\nu\Delta V
= a(\tau)V+a(\tau)(y\cdot\nabla V)+b(\tau)\cdot\nabla V,
\qquad
\nabla\cdot V=0,
\end{equation}
with $\nu>0$, $\nabla\cdot V=0$, and $a,b\in L^\infty(J)$. For every
nonnegative $\phi\in C_c^\infty(Q_{2R,J})$ and every $\tau_1<\tau_2$ in $J$,
\begin{equation}\label{paper4:eq:local-energy}
\begin{aligned}
&\int_{B_{2R}}\frac{|V(\cdot,\tau_2)|^2}{2}\phi(\cdot,\tau_2)
-\int_{B_{2R}}\frac{|V(\cdot,\tau_1)|^2}{2}\phi(\cdot,\tau_1)
\\
&\quad+\nu\int_{\tau_1}^{\tau_2}\!\int_{B_{2R}}\phi\,|\nabla V|^2
\le
\int_{\tau_1}^{\tau_2}\!\int_{B_{2R}}\frac{|V|^2}{2}(\partial_\tau\phi+\nu\Delta\phi)
\\
&\quad+\int_{\tau_1}^{\tau_2}\!\int_{B_{2R}}\frac{|V|^2}{2}V\cdot\nabla\phi
+\int_{\tau_1}^{\tau_2}\!\int_{B_{2R}}PV\cdot\nabla\phi
\\
&\quad+\int_{\tau_1}^{\tau_2}\!\int_{B_{2R}}
a\frac{|V|^2}{2}(-\phi-y\cdot\nabla\phi)
\\
&\quad-\int_{\tau_1}^{\tau_2}\!\int_{B_{2R}}\frac{|V|^2}{2}b\cdot\nabla\phi.
\end{aligned}
\end{equation}
The pressure is understood modulo functions of time, and the mean-normalized
pressure on compact balls belongs to the local class used below.
\end{proposition}

\begin{proof}
\Cref{paper2:prop:ren-suitable} transports the physical local energy
inequality of Caffarelli--Kohn--Nirenberg
\cite{CaffarelliKohnNirenberg1982} into the renormalized variables
\((y,\tau)\) and yields the inequality stated above with test function of the
form \(\phi^2\), where \(\phi\in C_c^\infty(Q_R^*)\) is nonnegative.  Setting
\(\Phi:=\phi^2\) and noting that every nonnegative
\(\Phi\in C_c^\infty(Q_{2R,J})\) is the limit, in
\(C_c^k(Q_{2R,J})\) for every \(k\), of squares of nonnegative smooth
compactly supported functions \(\phi_\varepsilon:=(\Phi+\varepsilon^2)^{1/2}-\varepsilon\),
the inequality extends by approximation to all such \(\Phi\); this is the
standard reformulation used in the theory of suitable weak solutions
\cite{CaffarelliKohnNirenberg1982,Lin1998}.  The compact-window
representation theorem \cref{paper2:thm:C} provides the cylinder
\(Q_{2R,J}=B_{2R}\times J\) and the divergence-free renormalized equation
\eqref{paper4:eq:renorm-ns}, while
\cref{paper2:thm:pressure-decomp} supplies the localized pressure
decomposition with the mean-normalized pressure on compact balls.  Identifying
the modulation coefficients \(a=\tilde a\) and \(b=\tilde b\) of
\cref{paper2:prop:ren-suitable} with those of
\eqref{paper4:eq:renorm-ns} and grouping the renormalized error terms exactly
as in \eqref{paper4:eq:local-energy} yields the stated inequality.
\end{proof}

Throughout, for a space--time set $E\subset\mathbb R^3\times\mathbb R$, write
\[
\int_E f:=\int_E f(y,\tau)\,dy\,d\tau.
\]
For $r>0$ and a bounded interval $J\subset\mathbb R$, define
\[
Q_{r,J}:=B_r\times J,\qquad |Q_{r,J}|:=|B_r|\,|J|.
\]
For a ball $B\subset\mathbb R^3$, denote the spatial mean of pressure by
\[
(P)_B(\tau):=\frac1{|B|}\int_B P(y,\tau)\,dy.
\]
The local quantities are
\[
\mathcal A_J(r):=\operatorname*{ess\,sup}_{\tau\in J}\int_{B_r}|V(y,\tau)|^2\,dy,
\]
\[
\mathcal H_J(r):=\int_J\!\int_{B_r}|\nabla V|^2\,dy\,d\tau,
\]
\[
\mathcal C_J(r):=\int_J\!\int_{B_r}|V|^3\,dy\,d\tau,
\qquad
\mathcal D_J(r):=\int_J\!\int_{B_r}|P-(P)_{B_r}(\tau)|^{3/2}\,dy\,d\tau.
\]

\begin{lemma}[Finite-measure conversion]\label{paper4:lem:l2-from-l3}
For every bounded interval $J$ and every $r>0$,
\[
\int_J\!\int_{B_r}|V|^2
\le |Q_{r,J}|^{1/3}\mathcal C_J(r)^{2/3}
\le |Q_{r,J}|^{1/3}\bigl(1+\mathcal C_J(r)\bigr).
\]
\end{lemma}

\begin{proof}
H\"older's inequality on $Q_{r,J}$ gives
\[
\int_J\!\int_{B_r}|V|^2
\le |Q_{r,J}|^{1/3}\left(\int_J\!\int_{B_r}|V|^3\right)^{2/3}.
\]
The second inequality follows from $s^{2/3}\le 1+s$ for $s\ge0$.
\end{proof}

\begin{lemma}[Comparison of pressure normalizations]\label{paper4:lem:pressure-means}
Let $1\le p<\infty$ and $0<r\le R$. For almost every $\tau$ such that
$P(\tau)\in L^p(B_R)$,
\[
\|P(\tau)-(P)_{B_r}(\tau)\|_{L^p(B_r)}
\le 2\|P(\tau)-(P)_{B_R}(\tau)\|_{L^p(B_R)}.
\]
Consequently,
\[
\mathcal D_J(r)\le 2^{3/2}\mathcal D_J(R)
\]
whenever the right-hand side is finite.
\end{lemma}

\begin{proof}
Set $c=(P)_{B_R}(\tau)$. Then
\[
\|P-(P)_{B_r}\|_{L^p(B_r)}
\le \|P-c\|_{L^p(B_r)}+|B_r|^{1/p}|(P)_{B_r}-c|.
\]
Since
\[
|(P)_{B_r}-c|\le |B_r|^{-1}\int_{B_r}|P-c|
\le |B_r|^{-1/p}\|P-c\|_{L^p(B_r)},
\]
the displayed pointwise estimate follows. Raising to the power $3/2$ and
integrating in time gives the last assertion.
\end{proof}

\subsection{Compact-cylinder Caccioppoli estimate}

\begin{theorem}[Compact-cylinder Caccioppoli estimate]\label{paper4:thm:caccioppoli}
Let $R>0$, let $J\subset\mathbb R$ be a bounded open interval, and let
$J'\Subset J$ be a compact subinterval. Assume
\cref{paper4:prop:renormalized-local-energy-inequality} on $Q_{2R,J}$ and assume
\[
\mathcal C_J(2R)+\mathcal D_J(2R)<\infty.
\]
Then there is a constant
\[
C=C\bigl(R,|J|,\nu,\|a\|_{L^\infty(J)},\|b\|_{L^\infty(J)},
\operatorname{dist}(J',\partial J)\bigr)
\]
such that
\[
\mathcal A_{J'}(R)+\mathcal H_{J'}(R)
\le C\Bigl(1+
\mathcal C_J(2R)+\mathcal D_J(2R)
\Bigr).
\]
The estimate is outer-to-inner: all right-hand quantities are measured on
$B_{2R}\times J$, while the conclusion is on $B_R\times J'$.
\end{theorem}

\begin{proof}
\textbf{Step 1: localized cut-off and the two energy inequalities.}
Set
\[
\delta:=\operatorname{dist}(J',\partial J)>0.
\]
Choose $\eta\in C_c^\infty(J)$ and $\zeta\in C_c^\infty(B_{2R})$ with
\[
0\le\eta\le1,\qquad
0\le\zeta\le1,\qquad
\eta\equiv1\ \text{on}\ J',\qquad
\zeta\equiv1\ \text{on}\ B_R,
\]
and
\[
\|\eta'\|_{L^\infty(J)}\le C_\eta\,\delta^{-1},
\qquad
\|\nabla\zeta\|_{L^\infty(B_{2R})}\le C_\zeta R^{-1},
\qquad
\|\Delta\zeta\|_{L^\infty(B_{2R})}\le C_\zeta R^{-2}.
\]
Set
\[
\phi(y,\tau):=\eta(\tau)\,\zeta(y)^2.
\]
Then $\phi\in C_c^\infty(Q_{2R,J})$, $0\le\phi\le1$, and
\[
\phi\equiv1\text{ on }B_R\times J'.
\]
Choose $\tau_-<\inf J'$ and $\tau_+>\sup J'$ in $J$ with
$\eta(\tau_-)=\eta(\tau_+)=0$. Applying \eqref{paper4:eq:local-energy} on
$(\tau_-,\tau_+)$ gives
\[
\nu\int_{\tau_-}^{\tau_+}\!\int_{B_{2R}}\phi\,|\nabla V|^2
\le |R_\tau|+|R_\Delta|+|R_{\mathrm{conv}}|+|R_{\mathrm{pres}}|+|R_{\mathrm{mod}}|,
\]
where
\[
R_\tau:=\int_J\!\int_{B_{2R}}\frac{|V|^2}{2}\,\partial_\tau\phi,\qquad
R_\Delta:=\nu\int_J\!\int_{B_{2R}}\frac{|V|^2}{2}\,\Delta\phi,
\]
\[
R_{\mathrm{conv}}:=\int_J\!\int_{B_{2R}}\bigl(\tfrac{|V|^2}{2}V\bigr)\cdot\nabla\phi,\qquad
R_{\mathrm{pres}}:=\int_J\!\int_{B_{2R}}\bigl(P-(P)_{B_{2R}}(\tau)\bigr)V\cdot\nabla\phi,
\]
\[
R_{\mathrm{mod}}:=\int_J\!\int_{B_{2R}}
a\frac{|V|^2}{2}(-\phi-y\cdot\nabla\phi)
-\int_J\!\int_{B_{2R}}\frac{|V|^2}{2}b\cdot\nabla\phi.
\]
For almost every $\sigma\in J'$, applying \eqref{paper4:eq:local-energy} on
$(\tau_-,\sigma)$ gives the same right-hand estimates and leaves the terminal
term $\frac12\int_{B_R}|V(\sigma)|^2$ on the left. Therefore the bounds below
control both $\mathcal H_{J'}(R)$ and $\mathcal A_{J'}(R)$.
Define
\[
L:=\int_J\!\int_{B_{2R}}|V|^2,\qquad
C_3:=\mathcal C_J(2R),\qquad D_{3/2}:=\mathcal D_J(2R).
\]
By \cref{paper4:lem:l2-from-l3},
\[
L\le |Q_{2R,J}|^{1/3}C_3^{2/3}\le C(R,|J|)(1+C_3).
\]

Since $\phi(y,\tau)=\eta(\tau)\zeta(y)^2$,
\[
\partial_\tau\phi=\eta'(\tau)\zeta(y)^2,
\qquad
\nabla\phi=2\eta(\tau)\zeta(y)\nabla\zeta(y),
\qquad
\Delta\phi=\eta(\tau)\Delta(\zeta(y)^2),
\]
hence
\[
\|\partial_\tau\phi\|_{L^\infty(Q_{2R,J})}\le C_\eta\delta^{-1},
\qquad
\|\nabla\phi\|_{L^\infty(Q_{2R,J})}\le 2C_\zeta R^{-1},
\qquad
\|\Delta\phi\|_{L^\infty(Q_{2R,J})}\le C_\zeta R^{-2}.
\]

\textbf{Step 2: derivative terms from the cut-off.}
\[
|R_\tau|
\le\frac12\|\partial_\tau\phi\|_{L^\infty(Q_{2R,J})}L
\le C\delta^{-1}(1+C_3),
\]
\[
|R_\Delta|
\le\frac{\nu}{2}\|\Delta\phi\|_{L^\infty(Q_{2R,J})}L
\le C\nu R^{-2}(1+C_3).
\]

\textbf{Step 3: convection term.}
Using $|\nabla\phi|\le CR^{-1}$,
\[
|R_{\mathrm{conv}}|
\le \frac12\|\nabla\phi\|_{L^\infty(Q_{2R,J})}\int_J\!\int_{B_{2R}}|V|^3
\le CR^{-1}C_3.
\]

\textbf{Step 4: pressure term.}
For each time $\tau$,
\[
\int_{B_{2R}}(P-(P)_{B_{2R}}(\tau))V\cdot \nabla\phi
=\int_{B_{2R}}PV\cdot \nabla\phi,
\]
with
\[
\int_{B_{2R}}(P)_{B_{2R}}(\tau)V\cdot \nabla\phi
=(P)_{B_{2R}}(\tau)\int_{B_{2R}}V\cdot \nabla\phi
=-(P)_{B_{2R}}(\tau)\int_{B_{2R}}\phi\,\nabla\cdot V=0
\]
(using $\phi\in C_c^\infty(B_{2R})$ and $\nabla\cdot V=0$).
Thus subtracting any time-dependent pressure constant is legitimate in the
pressure flux term.
Hence
\[
|R_{\mathrm{pres}}|
\le\|\nabla\phi\|_{L^\infty(Q_{2R,J})}
\int_J
\|P-(P)_{B_{2R}}(\tau)\|_{L_x^{3/2}(B_{2R})}
\|V(\tau,\cdot)\|_{L_x^3(B_{2R})}
\,d\tau
\le CR^{-1}D_{3/2}^{2/3}C_3^{1/3}.
\]
Young's inequality gives
\[
|R_{\mathrm{pres}}|\le C(R)(C_3+D_{3/2}).
\]

\textbf{Step 5: modulation terms.}
The local energy inequality already contains the modulation terms
after the spatial integrations by parts from the repaired-gauge equation:
\[
\int_{B_{2R}}(y\cdot\nabla V)\cdot V\,\phi
=\frac12\int_{B_{2R}}y\cdot\nabla(|V|^2)\phi
=-\frac12\int_{B_{2R}}|V|^2\nabla\cdot(y\phi),
\]
and
\[
\int_{B_{2R}}(b(\tau)\cdot\nabla V)\cdot V\,\phi
=\frac12\int_{B_{2R}}b(\tau)\cdot\nabla(|V|^2)\phi
=-\frac12\int_{B_{2R}}|V|^2b(\tau)\cdot\nabla\phi.
\]
Thus we estimate the additional terms directly. Split
\[
R_{\mathrm{mod}}=R_a+R_{ay}+R_b,
\]
where
\[
R_a:=-\int_J\!\int_{B_{2R}}a\,\frac{|V|^2}{2}\phi,
\quad
R_{ay}:=-\int_J\!\int_{B_{2R}}a\,\frac{|V|^2}{2}\,y\cdot\nabla\phi,
\quad
R_b:=-\int_J\!\int_{B_{2R}}\frac{|V|^2}{2}b\cdot\nabla\phi.
\]
\[
M_a:=\|a\|_{L^\infty(J)},\qquad M_b:=\|b\|_{L^\infty(J)}.
\]
For $R_a$,
\[
|R_a|
\le \frac12 M_aL
\le C M_a(1+C_3).
\]
For the scale-gradient part, $|y|\le2R$ on $B_{2R}$ and
$\|\nabla\phi\|_{L^\infty(Q_{2R,J})}\le CR^{-1}$, hence
\[
|y\cdot\nabla\phi|\le C.
\]
Therefore
\[
|R_{ay}|\le C M_aL\le C M_a(1+C_3).
\]
For $R_b$,
\[
|R_b|\le C M_bR^{-1}L\le C M_bR^{-1}(1+C_3).
\]

\textbf{Step 6: absorption and conclusion.}
Collecting estimates from Steps~2--5 yields
\[
\nu\int_{\tau_-}^{\tau_+}\!\int_{B_{2R}}\phi|\nabla V|^2
\le C\bigl(1+C_3+D_{3/2}\bigr),
\]
where $C$ depends only on
\[
R,|J|,\nu,\|a\|_{L^\infty(J)},\|b\|_{L^\infty(J)},\delta.
\]
Since $\phi\equiv1$ on $B_R\times J'$,
\[
\mathcal H_{J'}(R)
\le C\bigl(1+\mathcal C_J(2R)+\mathcal D_J(2R)\bigr).
\]
Repeating the same estimates with terminal time $\sigma\in J'$ gives
\[
\frac12\int_{B_R}|V(y,\sigma)|^2\,dy
\le C\bigl(1+\mathcal C_J(2R)+\mathcal D_J(2R)\bigr)
\]
for almost every $\sigma\in J'$. Taking the essential supremum over
$\sigma\in J'$ and renaming $C$ proves the theorem.
\end{proof}

\begin{definition}[Local compact-cylinder control]
Local Caffarelli--Kohn--Nirenberg control on $Q_{r,J}$ means
\[
\mathcal C_J(r)+\mathcal D_J(r)<\infty.
\]
Local windowed $H^1$ control on $Q_{r,J}$ means
\[
\mathcal H_J(r)<\infty.
\]
\end{definition}

\begin{corollary}[Outer control implies inner suitable control]\label{paper4:cor:outer-inner}
Let $R>0$, let $J\subset\mathbb R$ be a bounded open interval, and let
$J'\Subset J$ be a compact subinterval. Assume \cref{paper4:prop:renormalized-local-energy-inequality} on $Q_{2R,J}$ and
$\mathcal C_J(2R)+\mathcal D_J(2R)<\infty$. Then
\[
\mathcal A_{J'}(R)+\mathcal H_{J'}(R)<\infty.
\]
More quantitatively, the bound is the one in \cref{paper4:thm:caccioppoli}. If an
argument uses pressure means on a smaller ball, \cref{paper4:lem:pressure-means}
converts those normalizations to the outer-ball normalization above.
\end{corollary}

\begin{proof}
The estimate follows from \cref{paper4:thm:caccioppoli}. The comparison of pressure normalizations follows from
\cref{paper4:lem:pressure-means}.
\end{proof}

\begin{corollary}[Rough-core exclusion on compact windows]\label{paper4:cor:rough-core}
Let $R>0$, let $J\subset\mathbb R$ be a bounded open interval, and let
$J'\Subset J$ be a compact subinterval. Assume \cref{paper4:prop:renormalized-local-energy-inequality} on $Q_{2R,J}$. If
\[
\mathcal C_J(2R)+\mathcal D_J(2R)<\infty,
\]
then
\[
\mathcal H_{J'}(R)<\infty.
\]
Equivalently, if $\mathcal H_{J'}(R)=+\infty$ for an inner cylinder
$Q_{R,J'}$, then for every enclosing cylinder $Q_{2R,J}$ with
$J'\Subset J$ on which the local energy inequality and bounded
modulation hypotheses hold,
\[
\mathcal C_J(2R)=+\infty
\qquad\text{or}\qquad
\mathcal D_J(2R)=+\infty.
\]
\end{corollary}

\begin{proof}
The direct implication is the gradient part of \cref{paper4:thm:caccioppoli}. The
last assertion is its contrapositive with the same outer-to-inner geometry:
finite outer $\mathcal C+\mathcal D$ on $Q_{2R,J}$ forces finite inner
$\mathcal H$ on $Q_{R,J'}$.
\end{proof}

\begin{remark}[Meaning of rough core]
Here a rough core means failure of local windowed $H^1$
control on compact inner renormalized cylinders while the modulation
coefficients are bounded. \Cref{paper4:cor:rough-core} shows that such a failure is
not independent of compact-cylinder Caffarelli--Kohn--Nirenberg control: it can
occur only if the velocity-pressure bounds $\mathcal C+\mathcal D$ fail on
every relevant larger enclosing compact cylinder.
\end{remark}

\subsection{Consequences for the assembly}

We shall use the following consequences in the final assembly.
\begin{enumerate}[label=\textup{(\arabic*)}]
\item On every compact renormalized window satisfying the local energy
inequality and bounded modulation, finite $\mathcal C+\mathcal D$ on
$Q_{2R,J}$ implies finite $\mathcal A+\mathcal H$ on $Q_{R,J'}$ whenever
$J'\Subset J$.
\item The pressure normalization may be taken on any comparable ball, after
using \cref{paper4:lem:pressure-means}.
\item Failure of local windowed $H^1$ control on an inner compact cylinder
forces failure of the compact-cylinder $\mathcal C+\mathcal D$ bounds on each
larger enclosing cylinder to which \cref{paper4:thm:caccioppoli} applies.
\end{enumerate}

\clearpage
\section{Scale-Collapse Cost and Autonomous Limits}\label{sec:scale-collapse-cost}
\medskip
\noindent\textit{Notation and functional conventions.}
For a bounded spatial region \(K\subset\mathbb R^3\), we use
\[
  (f,g)_{L^2(K)}:=\int_K f\cdot g\,dx,\qquad
  \langle F,\phi\rangle_{H^{-1}(K),H^1_0(K)}
  :=\int_0^T\langle F(\cdot,s),\phi(\cdot,s)\rangle_{H^{-1}(K),H^1_0(K)}\,ds
\]
whenever \(F\in L^1(0,T;H^{-1}(K))\) and \(\phi\in L^\infty(0,T;H^1_0(K))\). Weak limits and weak convergences are taken in the distributional sense on \(C_c^\infty\)-test functions.
\[
  \|f\|_{L^p_tL^q_x}:=\left(\int_0^T\|f(\cdot,s)\|_{L^q(\mathbb R^3)}^p\,ds\right)^{1/p},
\]
for any fixed \(T>0\), with the usual modification \(p=\infty\).

\subsection{Renormalized scale-collapse setting}

\medskip
\noindent\textit{Position of the result.}
The scale-collapse analysis has two components:
\begin{enumerate}[label=(\arabic*)]
  \item Sections~\ref{paper5:sec:cost}--\ref{paper5:sec:finite-cost} contain the cost
    exclusion arguments, which do not use compactness.
  \item Section~\ref{paper5:sec:autonomous} contains the autonomous compactness
    statement under explicitly stated compactness hypotheses.
\end{enumerate}

Throughout the scale-collapse analysis, the renormalized representation established earlier is assumed: on the
time interval under consideration, the fields \((V,P)\) satisfy \eqref{paper5:eq:renorm-ns} and its
distributional weak formulation.  In particular, the solenoidal formulation is obtained by
restricting the test fields below to divergence-free fields.

Let \((u,p)\) be a divergence-free suitable weak solution of
\[
  \partial_tu+(u\cdot\nabla)u+\nabla p=\nu\Delta u,
  \qquad \nabla\cdot u=0,
\]
in \(\mathbb R^3\), with viscosity \(\nu>0\).  Let
\(\lambda(\tau)>0\), \(X(\tau)\in\mathbb R^3\), and a measurable gauge function
\(c(t)\) be given, and define the renormalized fields by
\[
  u(x,t)=\lambda(\tau)^{-1}V\left(\frac{x-X(\tau)}{\lambda(\tau)},\tau\right),
  \qquad
  p(x,t)=\lambda(\tau)^{-2}P\left(\frac{x-X(\tau)}{\lambda(\tau)},\tau\right)+c(t).
\]
Pressure is defined only up to an arbitrary function of time; \(c(t)\) has no effect on
\(\nabla p\), so it is irrelevant for all subsequent estimates.
Assume that \((V,P)\) satisfies on \((0,\infty)\times\mathbb R^3\)
\begin{equation}\label{paper5:eq:renorm-ns}
  \partial_\tau V+(V\cdot\nabla)V+\nabla P
  =\nu\Delta V+a(\tau)(V+y\cdot\nabla V)+b(\tau)\cdot\nabla V,
  \qquad \nabla\cdot V=0.
\end{equation}
For every \(\psi\in C_c^\infty((0,\infty)\times\mathbb R^3;\mathbb R^3)\) with \(\nabla\cdot\psi=0\),
the pair \((V,P)\) satisfies the weak formulation
\[
\begin{aligned}
0={}&\int_0^\infty\!\!\int_{\mathbb R^3}
\bigl(
-V\cdot\partial_\tau\psi
-(V\otimes V):\nabla\psi
+\nu\nabla V:\nabla\psi
\bigr)\,dy\,d\tau \\
&+\int_0^\infty\!\!\int_{\mathbb R^3}
\bigl(
-a(\tau)(V+y\cdot\nabla V)\cdot\psi
-(b(\tau)\cdot\nabla V)\cdot\psi
-P\,\nabla\cdot\psi
\bigr)\,dy\,d\tau .
\end{aligned}
\]
For every \(0<T<\infty\), restricting the \(\tau\)-integration to \((0,T)\) gives the same identity on \((0,T)\times\mathbb R^3\).
Equivalently, for every
\(\psi\in C_c^\infty((0,\infty)\times\mathbb R^3;\mathbb R^3)\), not necessarily
divergence free,
\[
\begin{aligned}
0={}&\int_0^\infty\!\!\int_{\mathbb R^3}
\bigl(
-V\cdot\partial_\tau\psi
-(V\otimes V):\nabla\psi
+\nu\nabla V:\nabla\psi
-P\,\nabla\cdot\psi
\bigr)\,dy\,d\tau \\
&+\int_0^\infty\!\!\int_{\mathbb R^3}
\bigl(
-a(\tau)(V+y\cdot\nabla V)\cdot\psi
-(b(\tau)\cdot\nabla V)\cdot\psi
\bigr)\,dy\,d\tau .
\end{aligned}
\]

Define positive and negative parts of the scalar field \(a\) by
\[
  a_+(\tau):=\max\{a(\tau),0\},
  \qquad
  a_-(\tau):=\max\{-a(\tau),0\}.
\]
Then for a.e. \(\tau\), \(a=a_+-a_-\), \(a_+,a_-\ge0\), and \(a_+a_-=0\).  Assume
\[
  a\in L^1_{\mathrm{loc}}([\tau_0,\infty)),
  \qquad
  b\in L^1_{\mathrm{loc}}([\tau_0,\infty);\mathbb R^3),
\]
and the scale convention
\begin{equation}\label{paper5:eq:log-scale-convention}
  \frac{d}{d\tau}\log\lambda(\tau)=a(\tau)
\end{equation}
for a.e. \(\tau\ge\tau_0\), with \(\log\lambda\) absolutely continuous on compact intervals.

All integrals below are interpreted in \([0,\infty]\). When we write
\(\int_{\tau_0}^{\infty}\), partial integrals are taken over \([\tau_0,T]\) and
monotone convergence is used as \(T\to\infty\).

Fix \(R>0\).  Choose
\(
  \Phi\in C_c^\infty(\mathbb R^3;[0,1])
\)
satisfying
\begin{equation}\label{paper5:eq:Phi-properties}
  \Phi(y)=1\ \text{if }|y|\le R,
  \quad
  \operatorname{supp}_y\Phi\subset B_{2R}.
\end{equation}
Define the localized kinetic mass
\[
  M_R(\tau):=\int_{\mathbb R^3}|V(y,\tau)|^2\Phi(y)
  \,dy.
\]
Since \(\Phi\) is fixed in \(\tau\), \(M_R\) is measurable in \(\tau\), and only this
measurability is used below.

\begin{definition}[Nonvanishing localized core]
A renormalized branch has nonvanishing localized core if there exist
constants \(m_0>0\) and \(\tau_1\ge\tau_0\) such that
\[
  M_R(\tau)\ge m_0\quad\text{for a.e. }\tau\in[\tau_1,\infty).
\]
In applications, this lower bound is supplied by the selected-core construction; here it is stated as an explicit hypothesis.
\end{definition}

\begin{definition}[Scale-collapse and absolute scale-drift densities]
Define
\[
  \mathcal C_{\mathrm{sc}}(\tau):=a_-(\tau)M_R(\tau),
  \qquad
  \mathcal C_{\mathrm{abs}}(\tau):=\nu\int_{\mathbb R^3}|\nabla V(y,\tau)|^2\Phi(y)
  \,dy+|a(\tau)|M_R(\tau).
\]
\(\mathcal C_{\mathrm{sc}}\) is the negative drift weighted by localized core mass, while
\(\mathcal C_{\mathrm{abs}}\) controls both collapse drift and local dissipation.
Their total costs are extended real-valued by
\[
  \int_{\tau_0}^{\infty}\mathcal C_{\mathrm{sc}}(\tau)
  \,d\tau,
  \qquad
  \int_{\tau_0}^{\infty}\mathcal C_{\mathrm{abs}}(\tau)
  \,d\tau.
\]
\end{definition}

\begin{definition}[Finite-cost branch]\label{paper5:def:finite-cost-branch}
A branch has finite scale-collapsing cost if
\[
  \int_{\tau_0}^{\infty}\mathcal C_{\mathrm{sc}}(\tau)
  \,d\tau<\infty.
\]
It has finite absolute cost if
\[
  \int_{\tau_0}^{\infty}\mathcal C_{\mathrm{abs}}(\tau)
  \,d\tau<\infty.
\]
\end{definition}

\subsection{Elementary alternative}
\label{paper5:sec:cost}

\begin{lemma}[Finite-or-infinite alternative]\label{paper5:lem:finite-infinite}
Let \(f\ge0\) be measurable on \([\tau_0,\infty)\). Then exactly one of the following holds:
\[
  \int_{\tau_0}^{\infty}f(\tau)
  \,d\tau<\infty
  \quad\text{or}\quad
  \int_{\tau_0}^{\infty}f(\tau)
  \,d\tau=\infty.
\]
\end{lemma}

\begin{proof}
\begin{enumerate}[label=(\arabic*)]
  \item For \(T>\tau_0\), define partial costs
  \[
    F(T):=\int_{\tau_0}^{T}f(\tau)
    \,d\tau\in[0,\infty].
  \]

  \item If \(T_2>T_1\), then by nonnegativity of \(f\),
  \[
    F(T_2)-F(T_1)=\int_{T_1}^{T_2}f(\tau)
    \,d\tau\ge0,
  \]
  so \(F\) is monotone nondecreasing.

  \item A monotone function on \([\tau_0,\infty)\) has an extended limit
  \(L\in[0,\infty]\),
  \[
    L:=\lim_{T\to\infty}F(T).
  \]
  By definition of improper integral, this limit equals
  \(\int_{\tau_0}^{\infty}f(\tau)\,d\tau\).

  \item Therefore \(L\) is either finite or equal to \(+\infty\), and these
  cases are mutually exclusive.
\end{enumerate}
Hence exactly one item in the lemma holds.
\end{proof}

\begin{lemma}[Logarithmic scale identity]\label{paper5:lem:log-scale}
Under \eqref{paper5:eq:log-scale-convention}, for all \(\tau_1,\tau_2\) with
\(\tau_0\le\tau_1<\tau_2<\infty\),
\[
  \int_{\tau_1}^{\tau_2}a(\tau)
  \,d\tau
  =\log\lambda(\tau_2)-\log\lambda(\tau_1).
\]
\end{lemma}

\begin{proof}
\begin{enumerate}[label=(\arabic*)]
  \item Since \(a\in L^1_{\mathrm{loc}}\), the map
  \(\log\lambda\) is absolutely continuous on \([\tau_1,\tau_2]\), hence has a
  representative with a.e. derivative equal to \(a\).

  \item By the fundamental theorem of calculus for absolutely continuous functions,
  for every \(\tau\in[\tau_1,\tau_2]\),
  \[
    \log\lambda(\tau)=\log\lambda(\tau_1)+\int_{\tau_1}^{\tau}a(s)
    \,ds.
  \]

  \item Setting \(\tau=\tau_2\) yields the identity.
\end{enumerate}
\end{proof}

\subsection{Scale collapse versus localized core lower bounds}

\begin{theorem}[Collapse with nonvanishing localized core forces infinite
scale-collapsing cost]\label{paper5:thm:collapse-core-floor}
Assume:
\begin{enumerate}[label=(\arabic*)]
  \item \(\lambda(\tau)\to0\) as \(\tau\to\infty\);
  \item the branch has nonvanishing localized core.
\end{enumerate}
Then
\[
  \int_{\tau_0}^{\infty}\mathcal C_{\mathrm{sc}}(\tau)
  \,d\tau
  =\int_{\tau_0}^{\infty}a_-(\tau)M_R(\tau)
  \,d\tau=\infty.
\]
\end{theorem}

\begin{proof}
\begin{enumerate}[label=(\arabic*)]
  \item Choose \(m_0>0\) and \(\tau_1\ge\tau_0\) from the nonvanishing
  core assumption so that
  \[
    M_R(\tau)\ge m_0\quad\text{for a.e. }\tau\in[\tau_1,\infty).
  \]

  \item By Lemma~\ref{paper5:lem:log-scale}, for any \(\tau_2>\tau_1\),
  \[
    \int_{\tau_1}^{\tau_2}a(\tau)
    \,d\tau
    =\log\lambda(\tau_2)-\log\lambda(\tau_1).
  \]
  Taking \(\tau_2\to\infty\) and using \(\lambda(\tau_2)\to0\),
  \[
    \lim_{\tau_2\to\infty}\int_{\tau_1}^{\tau_2}a(\tau)
    \,d\tau=-\infty.
  \]

  \item Decompose \(a=a_+-a_-\).  Then for each \(\tau_2>\tau_1\),
  \[
    \int_{\tau_1}^{\tau_2}a(\tau)
    \,d\tau
    =\int_{\tau_1}^{\tau_2}a_+(\tau)
    \,d\tau-\int_{\tau_1}^{\tau_2}a_-(\tau)
    \,d\tau
    \ge-\int_{\tau_1}^{\tau_2}a_-(\tau)
    \,d\tau.
    \]
    Hence
    \[
    \int_{\tau_1}^{\tau_2}a_-(\tau)
    \,d\tau\ge
    -\int_{\tau_1}^{\tau_2}a(\tau)
    \,d\tau.
    \]

  \item Taking limits and using Step~2 yields
  \[
    \int_{\tau_1}^{\infty}a_-(\tau)
    \,d\tau=\infty.
  \]
  Indeed, by Step~2,
  \[
    -\int_{\tau_1}^{\tau_2}a(\tau)\,d\tau\to+\infty.
  \]
  Since the partial integrals of \(a_-\) dominate this quantity, they are unbounded.
  Because \(a_-\ge0\), monotone convergence gives
  \(\int_{\tau_1}^{\infty}a_-(\tau)\,d\tau=\infty\).

  \item Split the cost into early and late times:
  \[
    \int_{\tau_0}^{\infty}a_-M_R
    \,d\tau
    =\int_{\tau_0}^{\tau_1}a_-M_R
    \,d\tau+
    \int_{\tau_1}^{\infty}a_-M_R
    \,d\tau.
  \]
  The first term is finite or infinite independently of the argument; the second term
  satisfies
  \[
    \int_{\tau_1}^{\infty}a_-(\tau)M_R(\tau)
    \,d\tau
    \ge m_0\int_{\tau_1}^{\infty}a_-(\tau)
    \,d\tau
    =\infty
  \]
  by the bound in (1) and Step~4.  Hence the total integral is infinite.
\end{enumerate}
\end{proof}

\subsection{Finite-cost exclusions}
\label{paper5:sec:finite-cost}

\begin{theorem}[Scale-collapse cost exclusion]\label{paper5:thm:cost-exclusion}
No branch with finite scale-collapsing cost (Definition~\ref{paper5:def:finite-cost-branch}) can satisfy simultaneously
\(\lambda(\tau)\to0\) and nonvanishing localized core.
\end{theorem}

\begin{proof}
Assume, for contradiction, that such a branch exists.  The two hypotheses are exactly
those of Theorem~\ref{paper5:thm:collapse-core-floor}; therefore
\[
  \int_{\tau_0}^{\infty}\mathcal C_{\mathrm{sc}}(\tau)
  \,d\tau=\infty.
\]
This contradicts finite scale-collapsing cost.  Therefore no such branch exists.
\end{proof}

\begin{lemma}[Absolute-cost domination]\label{paper5:lem:absolute-cost}
For a.e. \(\tau\in[\tau_0,\infty)\),
\[
  \mathcal C_{\mathrm{abs}}(\tau)\ge\mathcal C_{\mathrm{sc}}(\tau).
\]
Hence, pointwise and after integrating in \(\tau\),
\[
  \int_{\tau_0}^{\infty}\mathcal C_{\mathrm{sc}}(\tau)\,d\tau
  \le
  \int_{\tau_0}^{\infty}\mathcal C_{\mathrm{abs}}(\tau)\,d\tau,
\]
as extended real numbers.
In particular, finite absolute cost implies finite scale-collapsing cost, and
equivalently
\[
  \int_{\tau_0}^{\infty}\mathcal C_{\mathrm{sc}}(\tau)\,d\tau=\infty
  \quad\Longrightarrow\quad
  \int_{\tau_0}^{\infty}\mathcal C_{\mathrm{abs}}(\tau)\,d\tau=\infty.
\]
\end{lemma}

\begin{proof}
At each time \(\tau\),
\[
  \mathcal C_{\mathrm{abs}}(\tau)
  =\nu\int|\nabla V|^2\Phi
  +(a_++a_-)M_R
  \ge(a_++a_-)M_R\ge a_-M_R
  =\mathcal C_{\mathrm{sc}}(\tau),
\]
since \(a_+,a_-\ge0\) and \(\nu\int|\nabla V|^2\Phi\ge0\).  Integrating over
\([\tau_0,\infty)\) yields the claimed implication.
\end{proof}

\begin{corollary}[Absolute scale-cost exclusion]\label{paper5:cor:absolute-cost-exclusion}
No branch with finite absolute cost in the sense of Definition~\ref{paper5:def:finite-cost-branch}
can satisfy simultaneously \(\lambda(\tau)\to0\) and nonvanishing localized core.
\end{corollary}

\begin{proof}
Assume instead such a branch exists.  By Theorem~\ref{paper5:thm:collapse-core-floor},
\(\int\mathcal C_{\mathrm{sc}}=\infty\).  By Lemma~\ref{paper5:lem:absolute-cost}, the
absolute cost is pointwise above \(\mathcal C_{\mathrm{sc}}\), so its integral must also be
infinite, contradicting finite absolute cost.
\end{proof}

\begin{remark}
The exclusions above are purely integral and use neither compactness nor any
autonomous-limit argument.
\end{remark}

\subsection{Autonomous reduced-limit theorem from the selected-window estimates}
\label{paper5:sec:autonomous}

The preceding exclusions are integral statements and do not rely on compactness.  On a fixed finite autonomous window, compactness requires not only the cost bound but also the selected-window estimates supplied by the preceding renormalized representation, pressure reconstruction, and local Caccioppoli estimates.

\medskip
\noindent\textit{Analytic inputs.}
The preceding sections supply the following ingredients on selected compact
renormalized windows:
\begin{enumerate}[label=(\arabic*)]
  \item the repaired-gauge renormalized equation and weak formulation for \((V,P)\);
  \item pressure reconstruction modulo spatial means, with local \(L^{3/2}\) bounds on compact
    cylinders;
  \item local windowed bounds
  \(V\in L^\infty_tL^2_x\cap L^2_tH^1_x\) on compact cylinders;
  \item a retained compact core lower bound preventing the selected profile from vanishing.
\end{enumerate}
Together with the thick-window convergence of the modulation coefficients, these estimates imply
the compactness and limit passage below.

\begin{definition}[Thick autonomous windows]\label{paper5:def:thick-autonomous-windows}
A branch has \emph{thick autonomous windows} if there exist constants
\(T>0\), \(c>0\), \(M_1\in(0,\infty)\), a sequence \((\tau_n)\) with
\(\tau_n\to\infty\), and limits \(a_\infty\in\mathbb R\),
\(b_\infty\in\mathbb R^3\), such that for every \(n\):
\begin{enumerate}[label=(\roman*)]
  \item
  \[
    \frac1T\int_{\tau_n}^{\tau_n+T}a_-(\tau)M_R(\tau)
    \,d\tau\ge c;
  \]
  \item
  \[
    0\le M_R(\tau)\le M_1\quad\text{for a.e. }\tau\in[\tau_n,\tau_n+T];
  \]
  \item
  \[
    a(\tau_n+\cdot)\to a_\infty\ \text{in }L^1(0,T),
    \qquad
    b(\tau_n+\cdot)\to b_\infty\ \text{in }L^1(0,T;\mathbb R^3).
  \]
\end{enumerate}
\end{definition}
\begin{remark}
This notion isolates a limiting drift coefficient \(a_\infty<0\).  The
compactness to extract a profile limit is obtained from the selected-window compactness estimates,
not from the scalar cost inequalities of Sections~\ref{paper5:sec:cost}--\ref{paper5:sec:finite-cost}.
\end{remark}

\begin{lemma}[Thick drift implies negative autonomous limit parameter]\label{paper5:lem:negative-parameter}
On thick autonomous windows, the limiting autonomous coefficient is negative:
\[
  a_\infty\le-\frac{c}{M_1}<0.
\]
\end{lemma}

\begin{proof}
\begin{enumerate}[label=(\arabic*)]
  \item From (ii), for a.e. \(\tau\in[\tau_n,\tau_n+T]\),
  \(0\le M_R(\tau)\le M_1\).  Therefore
  \[
    \frac{a_-(\tau)M_R(\tau)}{M_1}
    \le a_-(\tau).
  \]
  Therefore
  \[
    \frac1T\int_{\tau_n}^{\tau_n+T}a_-(\tau)
    \,d\tau
    \ge \frac1{M_1T}\int_{\tau_n}^{\tau_n+T}a_-(\tau)M_R(\tau)
    \,d\tau
    \ge \frac{c}{M_1}.
  \]

  \item Set \(a_n(s):=a(\tau_n+s)\).  By (iii), \(a_n\to a_\infty\) in
  \(L^1(0,T)\). The map \(r\mapsto r^-\) is 1-Lipschitz, so
  \(a_n^-\to a_\infty^-\) in \(L^1(0,T)\):
  \[
    \|a_n^--a_\infty^-\|_{L^1(0,T)}\le
    \|a_n-a_\infty\|_{L^1(0,T)}\to0.
  \]
  Thus
  \[
  \lim_{n\to\infty}\frac1T\int_{\tau_n}^{\tau_n+T}a_-(\tau)
    \,d\tau
    =\frac1T\int_0^T(a_\infty)^-
    \,ds=(a_\infty)^-.
  \]

  \item Combine Step~1 and Step~2:
  \((a_\infty)^-\ge c/M_1\), i.e.\ \(a_\infty\le-c/M_1<0\).
\end{enumerate}
\end{proof}

\begin{proposition}[Selected-window compact upper estimates]\label{paper5:prop:selected-window-suitable-estimates}
Let \((\tau_n)\) be a thick autonomous-window sequence and set
\[
  V_n(y,s):=V(y,\tau_n+s),\qquad
  P_n(y,s):=P(y,\tau_n+s),
\]
\[
  \alpha_n(s):=a(\tau_n+s),\qquad
  \beta_n(s):=b(\tau_n+s),
  s\in(0,T).
\]
Assume that the repaired-gauge representation of \cref{paper2:thm:C}, the
localized pressure decomposition of \cref{paper2:thm:pressure-decomp}, and the
compact-cylinder Caccioppoli estimate \cref{paper4:thm:caccioppoli} are
available on the selected windows.  Assume also that the ordered local
validation has not entered a local-energy, pressure-reconstruction, or
rough-core failure row on these windows.  Then, after passing to a subsequence,
for every bounded Lipschitz domain \(K\Subset\mathbb R^3\), there is
\(C_K<\infty\), independent of \(n\), such that
\[
  \|V_n\|_{L^\infty(0,T;L^2(K))}
  +\|V_n\|_{L^2(0,T;H^1(K))}
  +\|P_n-(P_n)_K\|_{L^{3/2}((0,T)\times K)}
  \le C_K,
\]
where
\[
  (P_n)_K(s):=\frac1{|K|}\int_KP_n(y,s)\,dy.
\]
Each \(V_n\) is divergence free.
\end{proposition}

\begin{proof}
The repaired-gauge representation supplies the shifted equation and
divergence-free condition.  The localized pressure decomposition gives the
\(L^{3/2}\) pressure oscillation bound after subtracting spatial means.  The
compact-cylinder Caccioppoli estimate on a slightly larger compact cylinder
gives the local \(L^\infty_tL^2_x\) and \(L^2_tH^1_x\) velocity bounds.  If any
of these inputs is missing, the ordered validation records the first missing
input as a local-energy, pressure-reconstruction, or rough-core failure row,
which is outside the retained compact-estimate branch.  Hence, on the retained
branch, the displayed estimates hold.
\end{proof}

\begin{corollary}[Derivation of selected-window compact upper estimates]
\label{paper5:prop:selected-window-derived}
Under the hypotheses of
\cref{paper5:prop:selected-window-suitable-estimates}, the local compact upper
estimate package of that proposition holds on every compact
\(K\Subset\mathbb R^3\).
\end{corollary}

\begin{proof}
This is exactly \cref{paper5:prop:selected-window-suitable-estimates}.  The
corollary is kept only as a reference point for later sections that cite the
derived estimate package.
\end{proof}

\begin{lemma}[Retained selected-window mass floor]
\label{paper5:lem:selected-window-mass-floor}
\label{paper5:lem:retained-selected-window-mass-floor}
Let \((\tau_n)\) be a selected fixed-length terminal window sequence of length
\(T>0\) for a represented Type II branch produced by the selected-core
construction of the preceding sections, and write
\[
  V_n(y,s):=V(y,\tau_n+s),\qquad s\in(0,T).
\]
If the loss-of-retained-active-core row does not occur on these selected
windows, then after passing to a subsequence there exist
\[
  0<\theta<T/2,\qquad R>0,\qquad \eta>0,
\]
such that for all sufficiently large \(n\),
\[
  \int_{T/2-\theta}^{T/2+\theta}\int_{B_R}|V_n(y,s)|^3\,dy\,ds\ge\eta .
\]
\end{lemma}

\begin{proof}
The selected-core construction retains a positive local active window whenever
the loss-of-retained-active-core row is excluded.  If no fixed compact
space-time subwindow carried a positive \(L^3\)-mass floor after passing to
subsequences, then every candidate retained core would lose its active
space-time mass under further window selection, which is precisely the
loss-of-retained-active-core row.  Therefore a subsequence admits one compact
ball and one compact time interval with the displayed positive \(L^3\) floor.
\end{proof}

\begin{lemma}[Local time regularity on autonomous windows]\label{paper5:lem:window-time-regularity}
Assume thick autonomous windows and Proposition~\ref{paper5:prop:selected-window-suitable-estimates}.  Let
\(K_0\Subset K_1\Subset\mathbb R^3\) be bounded Lipschitz domains, and choose an integer
\(m\) so large that \(H^m_0(K_1)\hookrightarrow W^{1,\infty}_0(K_1)\).  Then
\[
  \{\partial_sV_n\}_{n\ge1}
  \quad\text{is bounded in }L^1(0,T;H^{-m}(K_1)).
\]
\end{lemma}

\begin{proof}
Fix \(\zeta\in C_c^\infty(K_1)\) with \(\zeta\equiv1\) on a neighborhood of \(K_0\).
Let \(\phi\in H^m_0(K_1;\mathbb R^3)\).  Since \(H^m_0(K_1)\hookrightarrow W^{1,\infty}_0(K_1)\),
there is \(C_m\) such that
\[
  \|\phi\|_{L^\infty(K_1)}+\|\nabla\phi\|_{L^\infty(K_1)}
  \le C_m\|\phi\|_{H^m(K_1)}.
\]
The shifted equation is
\[
  \partial_sV_n
  =-(V_n\cdot\nabla)V_n-\nabla P_n+\nu\Delta V_n
  +\alpha_n(V_n+y\cdot\nabla V_n)+\beta_n\cdot\nabla V_n
\]
in distributions.  We estimate each term against \(\phi\).

For the convection term, integration by parts gives
\[
  \left|\langle (V_n\cdot\nabla)V_n,\phi\rangle\right|
  =
  \left|\int_{K_1}V_{n,i}V_{n,j}\partial_j\phi_i\,dy\right|
  \le C_m\|V_n(s)\|_{L^2(K_1)}^2\|\phi\|_{H^m(K_1)}.
\]
For the viscous term,
\[
  \left|\langle \nu\Delta V_n,\phi\rangle\right|
  =
  \nu\left|\int_{K_1}\nabla V_n:\nabla\phi\,dy\right|
  \le \nu C_m\|\nabla V_n(s)\|_{L^2(K_1)}|K_1|^{1/2}\|\phi\|_{H^m(K_1)}.
\]
For the pressure term, subtracting the spatial mean does not change the gradient:
\[
  \left|\langle\nabla P_n,\phi\rangle\right|
  =
  \biggl|\int_{K_1}(P_n-(P_n)_{K_1})\,\nabla\cdot\phi\,dy\biggr|
  \le C_m |K_1|^{1/3}
  \|P_n-(P_n)_{K_1}\|_{L^{3/2}(K_1)}
  \|\phi\|_{H^m(K_1)}.
\]
For the \(V_n\)-part of the scaling drift,
\[
  |\alpha_n(s)|\,\left|\int_{K_1}V_n\cdot\phi\,dy\right|
  \le C_m|\alpha_n(s)|\,|K_1|^{1/2}
  \|V_n(s)\|_{L^2(K_1)}\|\phi\|_{H^m(K_1)}.
\]
For the \(y\cdot\nabla V_n\)-part, integrate by parts:
\[
  \int_{K_1}(y\cdot\nabla V_n)\cdot\phi\,dy
  =
  -\int_{K_1}V_n\cdot(y\cdot\nabla\phi+3\phi)\,dy,
\]
because \(\phi\) has zero trace.  Hence
\[
  |\alpha_n(s)|\,\left|\int_{K_1}(y\cdot\nabla V_n)\cdot\phi\,dy\right|
  \le C_{K_1,m}|\alpha_n(s)|\|V_n(s)\|_{L^2(K_1)}\|\phi\|_{H^m(K_1)}.
\]
Similarly,
\[
  \left|\int_{K_1}(\beta_n\cdot\nabla V_n)\cdot\phi\,dy\right|
  =
  \left|\int_{K_1}V_n\cdot(\beta_n\cdot\nabla\phi)\,dy\right|
  \le C_m|\beta_n(s)|\,|K_1|^{1/2}
  \|V_n(s)\|_{L^2(K_1)}\|\phi\|_{H^m(K_1)}.
\]
Combining the estimates and taking the supremum over
\(\|\phi\|_{H^m(K_1)}\le1\), we obtain
\[
\begin{aligned}
  \|\partial_sV_n(s)\|_{H^{-m}(K_1)}
  \le C_{K_1,m}\big(&
  \|V_n(s)\|_{L^2(K_1)}^2
  +\|\nabla V_n(s)\|_{L^2(K_1)}\\
  &+\|P_n(s)-(P_n)_{K_1}(s)\|_{L^{3/2}(K_1)}
  +( |\alpha_n(s)|+|\beta_n(s)| )\|V_n(s)\|_{L^2(K_1)}
  \big).
\end{aligned}
\]
Integrating in \(s\in(0,T)\), using Proposition~\ref{paper5:prop:selected-window-suitable-estimates}, and using
\(\alpha_n\to a_\infty\), \(\beta_n\to b_\infty\) in \(L^1\), hence uniform \(L^1\)-bounds
for \(\alpha_n,\beta_n\), proves the claimed \(L^1(0,T;H^{-m})\) bound.
\end{proof}

\begin{lemma}[Local compactness on autonomous windows]\label{paper5:lem:autonomous-window-compactness}
Assume thick autonomous windows and Proposition~\ref{paper5:prop:selected-window-suitable-estimates}.  Then after
passing to a subsequence there exist
\[
  U\in L^\infty(0,T;L^2_{\mathrm{loc}}(\mathbb R^3))
  \cap L^2(0,T;H^1_{\mathrm{loc}}(\mathbb R^3)),
  \qquad
  \Pi\in L^{3/2}_{\mathrm{loc}}((0,T)\times\mathbb R^3),
\]
such that, for every compact \(K\Subset\mathbb R^3\),
\begin{enumerate}[label=(\alph*)]
  \item
  \[
    V_n\to U
    \quad\text{strongly in }L^2(0,T;L^2(K));
  \]
  \item
  \[
    V_n\to U
    \quad\text{strongly in }L^3((0,T)\times K);
  \]
  \item
  \[
    V_n\rightharpoonup U
    \quad\text{weakly in }L^2(0,T;H^1(K));
  \]
  \item
  \[
    V_n\rightharpoonup^\ast U
    \quad\text{weakly-* in }L^\infty(0,T;L^2(K));
  \]
  \item
  \[
    P_n-(P_n)_K\rightharpoonup\Pi_K
    \quad\text{weakly in }L^{3/2}((0,T)\times K),
  \]
  where the local distributions \(\nabla\Pi_K\) agree on overlaps and define
  a pressure gradient \(\nabla\Pi\) modulo functions of time.
\end{enumerate}
\end{lemma}

\begin{proof}
Fix \(K_0\Subset K_1\Subset\mathbb R^3\).  Proposition~\ref{paper5:prop:selected-window-suitable-estimates}
gives boundedness of \(V_n\) in \(L^2(0,T;H^1(K_1))\), and
Lemma~\ref{paper5:lem:window-time-regularity} gives boundedness of \(\partial_sV_n\) in
\(L^1(0,T;H^{-m}(K_1))\) for \(m\) large.  Since
\[
  H^1(K_1)\Subset L^2(K_0)\hookrightarrow H^{-m}(K_1),
\]
the Aubin--Lions--Simon compactness theorem implies, after extracting a subsequence,
\[
  V_n\to U
  \quad\text{strongly in }L^2(0,T;L^2(K_0)).
\]
The local \(L^\infty_tL^2_x\) and \(L^2_tH^1_x\) bounds imply a uniform
\(L^{10/3}((0,T)\times K_0)\) bound by standard interpolation.  Interpolating
between the strong \(L^2\) convergence and this uniform \(L^{10/3}\) bound
gives
\[
  V_n\to U
  \qquad\text{strongly in }L^q((0,T)\times K_0)
\]
for every \(2\le q<10/3\), in particular for \(q=3\).
The local \(L^\infty L^2\) and \(L^2H^1\) bounds also give, after extraction,
\[
  V_n\rightharpoonup^\ast U
  \quad\text{in }L^\infty(0,T;L^2(K_0)),
  \qquad
  V_n\rightharpoonup U
  \quad\text{in }L^2(0,T;H^1(K_0)).
\]
The pressure bound in Proposition~\ref{paper5:prop:selected-window-suitable-estimates} gives weak compactness of
\(P_n-(P_n)_{K_0}\) in \(L^{3/2}((0,T)\times K_0)\).  A diagonal extraction over balls
\(B_j\) gives a single subsequence for every compact \(K\).  On overlaps, two pressure limits
can differ only by a function of \(s\), because both represent the same distributional pressure
gradient in the limit equation.  Thus the gradients assemble into \(\nabla\Pi\), with \(\Pi\)
defined modulo functions of time.
\end{proof}

\begin{lemma}[Compatibility of local pressure limits]\label{paper5:lem:pressure-compatibility}
Let \(K_1,K_2\Subset\mathbb R^3\) be bounded Lipschitz domains with
nonempty connected overlap \(K_{12}:=K_1\cap K_2\).  Suppose that, after passing to a common
subsequence,
\[
  P_n-(P_n)_{K_i}\rightharpoonup \Pi_i
  \quad\text{in }L^{3/2}((0,T)\times K_i),
  \qquad i=1,2.
\]
Then there exists \(c_{12}\in L^{3/2}(0,T)\) such that
\[
  \Pi_1(y,s)-\Pi_2(y,s)=c_{12}(s)
  \quad\text{for a.e. }(y,s)\in K_{12}\times(0,T).
\]
Consequently the pressure gradient \(\nabla\Pi\) obtained from the local limits is globally
well-defined on \((0,T)\times\mathbb R^3\), with the pressure itself defined modulo functions
of time.
\end{lemma}

\begin{proof}
For each \(n\),
\[
  \bigl(P_n-(P_n)_{K_1}\bigr)-\bigl(P_n-(P_n)_{K_2}\bigr)
  =(P_n)_{K_2}(s)-(P_n)_{K_1}(s)
\]
on \(K_{12}\).  Hence its spatial gradient is zero in
\(\mathcal D'((0,T)\times K_{12})\).  Passing to the weak limit gives
\[
  \nabla_y(\Pi_1-\Pi_2)=0
  \quad\text{in }\mathcal D'((0,T)\times K_{12}).
\]
Since \(K_{12}\) is connected, the standard distributional characterization of functions with
zero spatial gradient implies that \(\Pi_1-\Pi_2=c_{12}(s)\) for some
\(c_{12}\in L^{3/2}(0,T)\).  The last assertion follows because adding a time-dependent scalar
to pressure does not change \(\nabla\Pi\).
\end{proof}

\begin{theorem}[Autonomous reduced-limit theorem]\label{paper5:thm:autonomous-limit}
Assume that a branch has thick autonomous windows and satisfies
Proposition~\ref{paper5:prop:selected-window-suitable-estimates} on the same window sequence.  Let
\((U,\Pi)\) be the limit extracted in Lemma~\ref{paper5:lem:autonomous-window-compactness}.  Then
\[
  \nabla\cdot U=0,
\]
\[
  U\not\equiv0
\]
and \((U,\Pi)\) satisfies
\begin{equation}\label{paper5:eq:autonomous-limit}
  \partial_sU+(U\cdot\nabla)U+\nabla\Pi
  =\nu\Delta U+a_\infty(U+y\cdot\nabla U)+b_\infty\cdot\nabla U,
  \qquad
  \nabla\cdot U=0,
\end{equation}
in \(\mathcal D'((0,T)\times\mathbb R^3)\).  Equivalently, for every
\(\varphi\in C_c^\infty((0,T)\times\mathbb R^3;\mathbb R^3)\),
\[
  \int_0^T\!\int
  \bigl(
  -U\cdot\partial_s\varphi
  -(U\otimes U):\nabla\varphi
  +\nu\nabla U:\nabla\varphi
  -\Pi\,\nabla\cdot\varphi
  -a_\infty(U+y\cdot\nabla U)\cdot\varphi
  -(b_\infty\cdot\nabla U)\cdot\varphi
  \bigr)\,dy\,ds=0.
\]
\end{theorem}

\begin{proof}
\emph{Step 1: divergence-free property of the limit.}

Each \(V_n\) is divergence free by Proposition~\ref{paper5:prop:selected-window-suitable-estimates}, so for any compactly supported scalar
\(\psi\in C_c^\infty((0,T)\times\mathbb R^3)\),
\[
  \int_{(0,T)\times\mathbb R^3}V_n\cdot\nabla\psi\,dy\,ds=0.
\]
Hence, for every such \(\psi\), since \(\nabla\psi\in C_c^\infty((0,T)\times\mathbb R^3)\) and \(V_n\to U\) in
\(L^2(0,T;L^2_{\mathrm{loc}})\),
\[
  \int_{(0,T)\times\mathbb R^3}(V_n-U)\cdot\nabla\psi\,dy\,ds\to0,
\]
we deduce
\[
  \int_{(0,T)\times\mathbb R^3}U\cdot\nabla\psi\,dy\,ds=0.
\]
Hence \(\nabla\cdot U=0\) in \(\mathcal D'((0,T)\times\mathbb R^3)\).


\emph{Step 2: nontriviality of the limit profile.}
By the retained-window branch of
\cref{paper5:thm:core-retention}, after passing to the same subsequence there
exist \(R>0\), \(\chi\in C_c^\infty(B_R)\), \(\chi\ge0\), and \(\eta>0\) such
that
\[
  \int_0^T\int_{\mathbb R^3}|V_n(y,s)|^3\chi(y)\,dy\,ds\ge\eta
\]
for all sufficiently large \(n\).  By
\cref{paper5:lem:autonomous-window-compactness}, \(V_n\to U\) strongly in
\(L^3((0,T)\times B_R)\).  Hence
\[
  \int_0^T\int_{\mathbb R^3}|U(y,s)|^3\chi(y)\,dy\,ds
  =
  \lim_{n\to\infty}
  \int_0^T\int_{\mathbb R^3}|V_n(y,s)|^3\chi(y)\,dy\,ds
  \ge\eta.
\]
Thus \(U\not\equiv0\).

\emph{Step 3: shifted weak formulation for \(V_n\).}
Take any test field
\(\varphi\in C_c^\infty((0,T)\times\mathbb R^3;\mathbb R^3)\), and set
\[
K_\varphi:=\overline{\{x\in\mathbb R^3:\exists s\in(0,T),\ \varphi(x,s)\neq0\}}
\subset\mathbb R^3,\qquad K_\varphi\ \text{bounded}.
\]
Since \((V,P)\) satisfies~\eqref{paper5:eq:renorm-ns} in distributions, the change of variables
\(\tau=\tau_n+s\) gives the shifted weak formulation
\begin{equation}\label{paper5:eq:weak-n}
\begin{aligned}
0={}&\int_0^T\!\int_{\mathbb R^3}
\Bigl(
-V_n\cdot\partial_s\varphi
-(V_n\otimes V_n):\nabla\varphi
+\nu\nabla V_n:\nabla\varphi
-P_n\nabla\cdot\varphi
\Bigr)\,dy\,ds\\
&+\int_0^T\!\int_{\mathbb R^3}
\Bigl(
-a(\tau_n+s)(V_n+y\cdot\nabla V_n)\cdot\varphi
-(b(\tau_n+s)\cdot\nabla V_n)\cdot\varphi
\Bigr)\,dy\,ds.
\end{aligned}
\end{equation}

\emph{Step 4: term-by-term passage to the limit.}
Write \(\alpha_n(s):=a(\tau_n+s)\), \(\beta_n(s):=b(\tau_n+s)\).

\medskip
\noindent\textit{(i) Time derivative.}

By strong convergence in \(L^2(0,T;L^2(K_\varphi))\),
\[
  \left|\int_0^T\!\int (V_n-U)\cdot\partial_s\varphi
  \,dy\,ds\right|
  \le
  \|V_n-U\|_{L^2(0,T;L^2(K_\varphi))}
  \|\partial_s\varphi\|_{L^2(0,T;L^2(K_\varphi))}
  \to0.
\]
Hence
\[\int V_n\cdot\partial_s\varphi\to\int U\cdot\partial_s\varphi.
\]

\medskip
\noindent\textit{(ii) Convection term.}
Decompose
\[
  V_n\otimes V_n-U\otimes U=(V_n-U)\otimes V_n+U\otimes(V_n-U).
\]
Then
\[
\begin{aligned}
&\left|\int_0^T\!\int((V_n\otimes V_n-U\otimes U):\nabla\varphi)
  \,dy\,ds\right|\\
&\quad\le
  \|\nabla\varphi\|_{L^\infty_{x,t}}
  \|V_n-U\|_{L^2(0,T;L^2(K_\varphi))}
  \bigl(\|V_n\|_{L^2(0,T;L^2(K_\varphi))}
  +\|U\|_{L^2(0,T;L^2(K_\varphi))}\bigr)
  \to0,
\end{aligned}
\]
using boundedness of \(V_n\) in \(L^2(0,T;L^2(K_\varphi))\). Therefore
\(
  \int(V_n\otimes V_n):\nabla\varphi
  \to\int(U\otimes U):\nabla\varphi.
\)

\medskip
\noindent\textit{(iii) Viscous term.}
From weak convergence of \(V_n\) in \(L^2(0,T;H^1_{\mathrm{loc}})\),
\[
  \int_0^T\!\int\nabla V_n:\nabla\varphi
  \,dy\,ds
  \to
  \int_0^T\!\int\nabla U:\nabla\varphi
  \,dy\,ds.
\]

\medskip
\noindent\textit{(iv) Pressure term.}
Let \(K\) be a bounded Lipschitz domain with \(K_\varphi\Subset K\).  Since
\(\varphi\) is compactly supported in \(K\),
\[
  \int_K (P_n)_K(s)\,\nabla\cdot\varphi(y,s)\,dy
  =
  (P_n)_K(s)\int_K\nabla\cdot\varphi(y,s)\,dy=0.
\]
Therefore
\[
  \int_0^T\!\int_{\mathbb R^3}P_n\nabla\cdot\varphi\,dy\,ds
  =
  \int_0^T\!\int_K(P_n-(P_n)_K)\nabla\cdot\varphi\,dy\,ds.
\]
By Lemma~\ref{paper5:lem:autonomous-window-compactness},
\[
  P_n-(P_n)_K\rightharpoonup\Pi_K
  \quad\text{in }L^{3/2}((0,T)\times K),
\]
and \(\nabla\cdot\varphi\in L^3((0,T)\times K)\).  The compatibility
Lemma~\ref{paper5:lem:pressure-compatibility} permits us to write the resulting pressure distribution
as \(\Pi\), modulo functions of time.  Hence
\[
  \int_0^T\!\int_{\mathbb R^3}P_n\nabla\cdot\varphi\,dy\,ds
  \to
  \int_0^T\!\int_K\Pi_K\nabla\cdot\varphi\,dy\,ds.
\]
The value of the last integral is independent of the chosen \(K\), because replacing
\(\Pi_K\) by \(\Pi_K+c(s)\) changes the integral by
\(\int_0^Tc(s)\int_K\nabla\cdot\varphi\,dy\,ds=0\).  The limit is therefore the pressure term
\(\int\Pi\nabla\cdot\varphi\).

\medskip
\noindent\textit{(v) Drift term with \(a\): part involving \(V_n\).}
Define
\[
  Y_n^0(s):=\int_{\mathbb R^3}V_n(y,s)\cdot\varphi(y,s)
  \,dy,
  \qquad
  Y^0(s):=\int_{\mathbb R^3}U(y,s)\cdot\varphi(y,s)
  \,dy.
\]
For each \(s\), Cauchy--Schwarz gives
\[
  |Y_n^0(s)|
  \le
  \|V_n(s)\|_{L^2(K_\varphi)}
  \|\varphi(s)\|_{L^2(K_\varphi)}.
\]
By Proposition~\ref{paper5:prop:selected-window-suitable-estimates}, there is \(C_0<\infty\)
such that \(\sup_n\|Y_n^0\|_{L^\infty(0,T)}\le C_0\). Moreover,
\[
  Y_n^0\to Y^0\quad\text{in }L^2(0,T),
\]
because \(V_n\to U\) in \(L^2(0,T;L^2(K_\varphi))\).
Now
\[
  \int_0^T\alpha_n(s)Y_n^0(s)
  \,ds
  =\int_0^T(\alpha_n(s)-a_\infty)Y_n^0(s)
  \,ds+a_\infty\int_0^TY_n^0(s)
  \,ds.
\]
For the first term,
\[
  \left|\int_0^T(\alpha_n-a_\infty)Y_n^0\,ds\right|\le
\sup_{s\in(0,T)}|Y_n^0(s)|\,\|\alpha_n-a_\infty\|_{L^1(0,T)}\to0.
\]
For the second term, strong convergence in \(L^2(0,T)\) implies
\(\int_0^TY_n^0\to\int_0^TY^0\), so
\[
  \int_0^T\alpha_n(s)\int_{\mathbb R^3}V_n\cdot\varphi
  \,dy\,ds
  \to a_\infty\int_0^T\int_{\mathbb R^3}U\cdot\varphi
  \,dy\,ds.
\]
Thus the signed contribution in \eqref{paper5:eq:weak-n} satisfies
\[
  -\int_0^T\alpha_n(s)\int_{\mathbb R^3}V_n\cdot\varphi
  \,dy\,ds
  \to -a_\infty\int_0^T\int_{\mathbb R^3}U\cdot\varphi
  \,dy\,ds.
\]

\medskip
\noindent\textit{(vi) Drift term with \(a\): part involving \(y\cdot\nabla V_n\).}
Since \(\varphi\in C_c^\infty\), integrating by parts in \(y\) gives
\[
  \int_{\mathbb R^3}(y\cdot\nabla V_n)\cdot\varphi
  =-\int_{\mathbb R^3}V_n\cdot(y\cdot\nabla\varphi+3\varphi)
  \,dy.
\]
Define
\[
  Y_n^1(s):=\int_{\mathbb R^3}V_n\cdot(y\cdot\nabla\varphi+3\varphi)
  \,dy,
  \qquad
  Y^1(s):=\int_{\mathbb R^3}U\cdot(y\cdot\nabla\varphi+3\varphi)
  \,dy.
\]
For each \(s\), Cauchy--Schwarz gives
\[
  |Y_n^1(s)|
  \le
  \|V_n(s)\|_{L^2(K_\varphi)}
  \|y\cdot\nabla\varphi(s)+3\varphi(s)\|_{L^2(K_\varphi)}.
\]
The \(L^\infty(0,T;L^2(K_\varphi))\) bound in
Proposition~\ref{paper5:prop:selected-window-suitable-estimates} gives
\[
  \sup_n\|Y_n^1\|_{L^\infty(0,T)}\le C_1<\infty.
\]
Also,
\[
  \|Y_n^1-Y^1\|_{L^2(0,T)}
  \le
  \|V_n-U\|_{L^2(0,T;L^2(K_\varphi))}
  \|y\cdot\nabla\varphi+3\varphi\|_{L^\infty(0,T;L^2(K_\varphi))}
  \to0.
\]
Therefore
\[
  -\int_0^T\alpha_n(s)
  \int_{\mathbb R^3}(y\cdot\nabla V_n)\cdot\varphi
  \,dy\,ds
  \to
  a_\infty\int_0^T\!\int_{\mathbb R^3}U\cdot(y\cdot\nabla\varphi+3\varphi)
  \,dy\,ds.
\]
Integrating by parts once more in the limit gives
\[
  a_\infty\int_0^T\!\int_{\mathbb R^3}U\cdot(y\cdot\nabla\varphi+3\varphi)
  \,dy\,ds
  =
  -a_\infty\int_0^T\!\int_{\mathbb R^3}(y\cdot\nabla U)\cdot\varphi
  \,dy\,ds.
\]

\medskip
\noindent\textit{(vii) Term with \(b\).}
Since \(\beta_n(s)\) is independent of \(y\),
\[
  \int_{\mathbb R^3}(\beta_n\cdot\nabla V_n)\cdot\varphi
  =-\int_{\mathbb R^3}V_n\cdot(\beta_n\cdot\nabla\varphi)
  \,dy,
\]
so
\[
  -\int_0^T\!\int(\beta_n\cdot\nabla V_n)\cdot\varphi
  \,dy\,ds
  =
  \int_0^T\!\int V_n\cdot(\beta_n\cdot\nabla\varphi)
  \,dy\,ds.
\]
Let
\[
\begin{aligned}
Z_n(s)&:=\left(\int_{\mathbb R^3}V_n\cdot\partial_1\varphi\,dy,\,
\int_{\mathbb R^3}V_n\cdot\partial_2\varphi\,dy,\,
\int_{\mathbb R^3}V_n\cdot\partial_3\varphi\,dy\right),\\
Z(s)&:=\left(\int_{\mathbb R^3}U\cdot\partial_1\varphi\,dy,\,
\int_{\mathbb R^3}U\cdot\partial_2\varphi\,dy,\,
\int_{\mathbb R^3}U\cdot\partial_3\varphi\,dy\right).
\end{aligned}
\]
Then
\[
  \int_{\mathbb R^3}V_n\cdot(\beta_n\cdot\nabla\varphi)\,dy
  =\beta_n(s)\cdot Z_n(s),\qquad
  \int_{\mathbb R^3}U\cdot(b_\infty\cdot\nabla\varphi)\,dy
  =b_\infty\cdot Z(s).
\]
Consequently
\[
  \left|\int_0^T\!(\beta_n-b_\infty)\cdot Z_n\,ds\right|
  \le
  \|\beta_n-b_\infty\|_{L^1(0,T)}\sup_{s\in(0,T)}|Z_n(s)|\to0,
\]
because
\[
  |Z_n(s)|\le
  \|V_n(s)\|_{L^2(K_\varphi)}
  \left(\sum_{j=1}^3\|\partial_j\varphi(s)\|_{L^2(K_\varphi)}^2\right)^{1/2}
\]
and Proposition~\ref{paper5:prop:selected-window-suitable-estimates} bounds the right-hand side
uniformly in \(n\) and \(s\). Furthermore,
\[
  \|Z_n-Z\|_{L^2(0,T;\mathbb R^3)}
  \le
  \|V_n-U\|_{L^2(0,T;L^2(K_\varphi))}
  \left(\sum_{j=1}^3\|\partial_j\varphi\|_{L^\infty(0,T;L^2(K_\varphi))}^2\right)^{1/2}
  \to0.
\]
Therefore
\[
  \int_0^T b_\infty\cdot Z_n(s)\,ds
  \to
  \int_0^T b_\infty\cdot Z(s)\,ds.
\]
Hence
\[
  -\int_0^T\!\int (\beta_n\cdot\nabla V_n)\cdot\varphi
  \,dy\,ds
  \to
  \int_0^T\!\int U\cdot(b_\infty\cdot\nabla\varphi)
  \,dy\,ds.
\]
Since \(b_\infty\) is constant in \(y\), integration by parts yields
\[
  \int_0^T\!\int U\cdot(b_\infty\cdot\nabla\varphi)
  \,dy\,ds
  =
  -\int_0^T\!\int (b_\infty\cdot\nabla U)\cdot\varphi
  \,dy\,ds.
\]

\emph{Step 5: assemble the limit.}
Passing to the limit in \eqref{paper5:eq:weak-n} in each term (Steps 4(i)--4(vii)) gives
\[
  \int_0^T\!\int
  \bigl(
  -U\cdot\partial_s\varphi
  -(U\otimes U):\nabla\varphi
  +\nu\nabla U:\nabla\varphi
  -\Pi\nabla\cdot\varphi
  -a_\infty(U+y\cdot\nabla U)\cdot\varphi
  -(b_\infty\cdot\nabla U)\cdot\varphi
  \bigr)
  \,dy\,ds=0.
\]
Since this holds for every compactly supported vector field \(\varphi\), equation~\eqref{paper5:eq:autonomous-limit} holds in distributions on
\((0,T)\times\mathbb R^3\).
\end{proof}

\begin{remark}
The theorem asserts only the compact finite-window autonomous limit from the
selected-window estimates.  A rigidity contradiction for an arbitrary
autonomous limit would require a separate Liouville theorem for the class
containing \((U,\Pi)\).  The retained compact scale-collapse branch used in the
terminal proof is handled differently: failures of retained inputs are named
local exits, and the remaining retained autonomous case is routed by
\cref{paper5:thm:compact-scale-collapse-rigidity-discharged} into the
scale-rigid discharge.
\end{remark}

\subsection{Decomposition of alternatives of the scale-collapse branch}
\label{paper5:sec:branch-closure}

The cost argument in \cref{paper5:sec:cost,paper5:sec:finite-cost} excludes the finite-cost scale-collapse branch.  The remaining
scale-collapse obstruction is treated by a decomposition of alternatives, following the local
estimate and autonomous-reduction argument of the preceding sections.  The conclusion separates the excluded branches from the residual alternatives; it does not assert a Liouville theorem for every autonomous negative-drift equation.

\begin{theorem}[Windowwise retained-core spacetime alternative]\label{paper5:thm:core-retention}
Let \((\tau_n)\) be a selected fixed-length terminal window sequence of length \(T>0\) for a represented
Type II branch produced by the selected-core construction of the preceding sections,
and write
\[
  V_n(y,s):=V(y,\tau_n+s),\qquad s\in(0,T).
\]
Then, after passing to a subsequence, exactly one of the following alternatives
holds:
\begin{enumerate}[label=\textup{(\roman*)}]
  \item the retained spacetime activity floor fails on these selected windows
  (equivalently, the loss-of-retained-active-core row occurs);
  \item there exist \(R>0\), \(\chi\in C_c^\infty(B_R)\), \(\chi\ge0\), and
  \(\eta>0\) such that
  \[
    \int_0^T\int_{\mathbb R^3} |V_n(y,s)|^3\chi(y)\,dy\,ds
    \ge \eta
  \]
  for all sufficiently large \(n\).
\end{enumerate}
\end{theorem}

\begin{proof}
The selected-core construction marks, on each admissible selected window, an
inner compact renormalized core carrying positive localized CKN velocity mass
on the time set used in the window analysis.  If no subsequence admitted a
uniform spacetime \(L^3\) lower bound with one fixed inner cutoff, then after
shrinking once inside the candidate core every further selected-window
subsequence would lose the retained spacetime activity floor.  This is exactly the
loss-of-retained-active-core row in the ordered local first-failure ledger.
Therefore, once that row is excluded, a subsequence exists on which one compact
inner cutoff and one positive threshold work uniformly on all remaining windows,
yielding alternative \textup{(ii)}.
\end{proof}

\begin{theorem}[Autonomous modulation selection]\label{paper5:thm:autonomous-modulation-selection}
Let a represented scale-collapse branch satisfy the modulation bounds obtained in the repaired
gauge construction.  Suppose that a sequence of windows \([\tau_n,\tau_n+T]\) is selected in the
autonomous regime.  Then, after passing to a subsequence, there exist constants
\(a_\infty\in\mathbb R\) and \(b_\infty\in\mathbb R^3\) such that
\[
  a(\tau_n+\cdot)\to a_\infty
  \quad\text{in }L^1(0,T),
  \qquad
  b(\tau_n+\cdot)\to b_\infty
  \quad\text{in }L^1(0,T;\mathbb R^3).
\]
\end{theorem}

\begin{proof}
The autonomous-window selection in the repaired-gauge modulation analysis gives convergence
of the modulation functions to constant autonomous parameters on the fixed window length
\(T\).  Passing to the corresponding subsequence gives the displayed
\(L^1\)-convergences.
\end{proof}

\begin{theorem}[Scale-collapse decomposition of alternatives]\label{paper5:thm:scale-collapse-stratification}
Let a represented Type II branch have genuine scale collapse and nonvanishing selected core.
Then exactly one of the following alternatives occurs after passing to selected terminal
windows:
\begin{enumerate}[label=(\roman*)]
  \item the branch has finite scale-collapsing cost, which is excluded by
  Theorem~\ref{paper5:thm:cost-exclusion};
  \item the branch has finite absolute cost, which is excluded by
  Corollary~\ref{paper5:cor:absolute-cost-exclusion};
  \item the branch admits thick autonomous windows in the sense of
  Definition~\ref{paper5:def:thick-autonomous-windows}, satisfies the selected-window estimates of
  Proposition~\ref{paper5:prop:selected-window-suitable-estimates}, and has autonomous modulation convergence;
  \item the branch has a thin-drift alternative: no selected fixed-length window carries a uniform
  positive lower bound for the averaged weighted negative drift;
  \item the branch has an autonomous-modulation alternative: thick windows exist, but no subsequence
  has a constant-coefficient \(L^1\) modulation limit in the sense needed for
  Theorem~\ref{paper5:thm:autonomous-limit};
  \item the branch has a compactness alternative: the local estimate package or the
  retained-core input needed for
  Proposition~\ref{paper5:prop:selected-window-suitable-estimates} fails on the selected windows.
\end{enumerate}
\end{theorem}

\begin{proof}
The scale-collapse selection theorem decomposes terminal scale-collapse windows according to the
averaged negative-drift mass
\[
  \frac1T\int_{\tau_n}^{\tau_n+T}a_-(\tau)M_R(\tau)\,d\tau.
\]
If the accumulated cost is finite, alternatives (i) or (ii) apply according to the cost
functional under consideration.  If the averaged negative-drift mass stays bounded below on a
subsequence, one obtains thick windows; otherwise the branch lies in the thin-drift alternative.
On thick windows, either the local estimate package holds, in which case
\cref{paper5:prop:selected-window-suitable-estimates} applies directly, or the
required representation, pressure, local-energy, or rough-core input fails on
the selected windows, in which case the branch lies in the compactness
alternative.  Separately, if the retained-core alternative fails on the
selected windows, that is the loss-of-retained-active-core row rather than an
input to the upper-estimate package.  If
the estimates hold but no subsequence has the required constant-coefficient modulation limit,
the branch lies in the autonomous-modulation alternative.  The remaining case is the thick
autonomous alternative.
\end{proof}

\begin{definition}[Retained compact scale-collapse state]
\label{paper5:def:retained-compact-scale-collapse-state}
A represented scale-collapse branch is in the retained compact
scale-collapse state if, on the selected terminal windows, it has:
\begin{enumerate}[label=\textup{(\roman*)}]
\item the repaired-gauge representation and local pressure reconstruction of
\cref{paper2:thm:C,paper2:thm:pressure-decomp};
\item the compact-cylinder suitable estimates used in
\cref{paper5:prop:selected-window-suitable-estimates};
\item thick autonomous selected windows in the sense of
\cref{paper5:def:thick-autonomous-windows}, together with a subsequential
constant-coefficient \(L^1\) modulation limit in the sense used by
\cref{paper5:thm:autonomous-limit};
\item bounded modulation on the selected compact cylinders;
\item a retained active Caffarelli--Kohn--Nirenberg mass floor on some fixed
inner selected cylinder;
\item the finite scale-collapsing-cost and finite absolute-cost alternatives
have already been removed by
\cref{paper5:thm:cost-exclusion,paper5:cor:absolute-cost-exclusion};
\item no unresolved multicore, same-point cascade,
separated/exterior, rough-core, autonomous-modulation-failure, or
compactness-package first-failure label after the ordered local state-space
tests have been applied.
\end{enumerate}
\end{definition}

\begin{lemma}[Failure of retained compact scale-collapse is a named exit]
\label{paper5:lem:retained-scale-collapse-first-failure}
Let a represented genuine scale-collapse branch with nonvanishing selected core
be tested against
\cref{paper5:def:retained-compact-scale-collapse-state} after the ordered
scale-collapse stratification has been applied.  If the branch is not in the
retained compact scale-collapse state, then its first failed retained-state
clause is one of the named local exits: representation/pressure failure,
thin-drift, finite scale-cost or finite absolute-cost exclusion, compactness
package first failure routed by
\cref{paper5:thm:selected-compactness-failure-routed},
autonomous-modulation failure, no retained active core,
multicore/same-point cascade, separated/exterior concentration, or rough-core
loss.
\end{lemma}

\begin{proof}
The clauses of
\cref{paper5:def:retained-compact-scale-collapse-state} are ordered in the same
way as the local tests.  Failure of clause \textup{(i)} is exactly the
representation or pressure-reconstruction exit.  Failure of
\textup{(ii)} is the compactness alternative in
\cref{paper5:thm:scale-collapse-stratification}, which is routed by
\cref{paper5:thm:selected-compactness-failure-routed}.  Failure of
\textup{(iii)} is either the thin-drift alternative or the
autonomous-modulation alternative in that theorem.  Failure of
\textup{(iv)} is a modulation-control exit in the repaired-gauge state-space
tests.  Failure of \textup{(v)} means the selected scale-collapse branch carries
no retained active Caffarelli--Kohn--Nirenberg core and is therefore not the
retained terminal branch.  Failure of \textup{(vi)} is precisely one of the two
cost alternatives already excluded by
\cref{paper5:thm:cost-exclusion,paper5:cor:absolute-cost-exclusion}.  Failure
of \textup{(vii)} is, by definition, one of the listed ordered local
state-space exits.  Thus no failed retained-state input remains unlabelled.
\end{proof}

\begin{definition}[Selected scale-collapse compactness package]
\label{paper5:def:selected-scale-collapse-compactness-package}
Let \((V_n,P_n,a_n,b_n)\) be a represented selected scale-collapse
window sequence on \(B_{2R}\times I\), where
\[
  V_n(y,s)=V(y,\tau_n+s),\qquad
  P_n(y,s)=P(y,\tau_n+s),\qquad s\in I,
\]
and \(I\Subset\mathbb R\) is a fixed model interval.  The sequence satisfies
the selected scale-collapse compactness package on \(B_{2R}\times I\) if the
following rows all hold with constants independent of \(n\):
\begin{enumerate}[label=\textup{(\roman*)}]
\item the repaired-gauge equation, divergence-free condition, local energy
inequality, and local pressure reconstruction modulo functions of time hold on
the window;
\item there is \(C_{R,I}<\infty\) such that
\[
  \sup_n\Big(
  \|V_n\|_{L^\infty(I;L^2(B_{2R}))}
  +\|V_n\|_{L^2(I;H^1(B_{2R}))}
  +\|P_n-c_{n,2R}\|_{L^{3/2}(B_{2R}\times I)}
  \Big)\le C_{R,I};
\]
\item the modulation has a local selected-window bound
\[
  \sup_n\int_I\bigl(|a_n(s)|+|b_n(s)|\bigr)\,ds\le C_{R,I};
\]
\item the retained selected core has a positive inner spacetime \(L^3\) floor:
for some \(0<r<R\), a compact interval \(J_*\Subset I\), and
\(\eta_*>0\),
\[
  \int_{J_*}\int_{B_r}|V_n(y,s)|^3\,dy\,ds\ge\eta_*
\]
after passing to a tail of the sequence.  This is the compact-window retained
mass floor supplied by
\cref{paper5:lem:selected-window-mass-floor,paper8:thm:active-profile-mass-floor};
\item the critical local upper rows have passed uniformly: for every compact
subcylinder \(B_\rho\times J\Subset B_{2R}\times I\) used by the terminal
local state-space test, there is \(C_{\rho,J}<\infty\), independent of \(n\),
such that
\[
  \sup_n\int_J\int_{B_\rho}|V_n(y,s)|^3\,dy\,ds\le C_{\rho,J}.
\]
Failure of this uniform row is not part of the retained package; it is routed
by \cref{paper6a:lem:terminal-local-critical-mass-routing}.
\end{enumerate}
This package is a local compact-cylinder statement.  It contains no whole-space
\(L^3\) bound, no global critical profile decomposition, and no assumption that
localized dissipation vanishes.
\end{definition}

\begin{lemma}[Spacetime active-core floor is the correct retained nonvanishing hypothesis]
\label{paper5:lem:spacetime-active-core-floor-suffices}
In the selected compactness package, every downstream use of the retained
active-core lower bound is implied by the spacetime \(L^3\) floor in
\cref{paper5:def:selected-scale-collapse-compactness-package}.  More precisely:
\begin{enumerate}[label=\textup{(\alph*)}]
\item if \(V_n\to V_\infty\) strongly in \(L^3_{\rm loc}\) on a compact
subcylinder \(Q':=B_r\times J\Subset B_{2R}\times I\), then
\[
\int_{Q'} |V_\infty|^3\,dy\,ds\ge \eta_*,
\]
provided \(Q'\) is the inner cylinder and \(\eta_*\) is the retained floor from
item \textup{(iv)} of
\cref{paper5:def:selected-scale-collapse-compactness-package}; in particular
\(V_\infty\not\equiv0\) on \(Q'\);
\item if one passes instead to an occupation-state or Young-state limit, then
the same retained spacetime \(L^3\) floor yields a positive visible local
occupation coordinate on \(Q'\);
\item no later use of the active-core hypothesis requires an a.e.-in-time
\(L^2\) lower floor once the local routing and compactness theorems are
formulated with the spacetime \(L^3\) activity coordinates.
\end{enumerate}
\end{lemma}

\begin{proof}
Part \textup{(a)} is immediate from strong \(L^3\)-convergence on \(Q'\):
\[
\int_{Q'}|V_\infty|^3\,dy\,ds
=
\lim_{n\to\infty}\int_{Q'}|V_n|^3\,dy\,ds
\ge \eta_*.
\]
For part \textup{(b)}, choose a nonnegative cutoff
\(\phi\in C_c^\infty(B_{2R}\times I)\) such that
\[
  0\le \phi\le1,
  \qquad
  \phi\equiv1
  \quad\text{on the retained floor cylinder }Q'.
\]
This is possible because \(Q'\Subset B_{2R}\times I\).  Then
\[
\int \phi|V_n|^3\,dy\,ds
\ge
\int_{Q'}|V_n|^3\,dy\,ds
\ge \eta_* ,
\]
so every occupation-state or Young-state limit records a positive local
occupation coordinate.  No shrinking of the retained floor cylinder is used.
Part \textup{(c)} is exactly how the later compactness and routing arguments use
the retained floor: they need a nonvanishing spacetime activity signal on the
retained compact window, not an a.e.-time \(L^2\) persistence statement.
\end{proof}

\begin{lemma}[Selected compactness package gives local strong compactness]
\label{paper5:lem:selected-compactness-aubin-lions}
Assume the selected scale-collapse compactness package of
\cref{paper5:def:selected-scale-collapse-compactness-package} on
\(B_{2R}\times I\).  Then, after passing to a subsequence,
\[
  V_n\to V_\infty
  \quad\text{strongly in }L^3(B_R\times I).
\]
For the pressure, for every
\(B_\rho\Subset B_{R_1}\Subset B_R\) and every cutoff
\(\zeta\in C_c^\infty(B_R)\) with \(\zeta\equiv1\) on \(B_{R_1}\), define
\[
  P_\infty^{\rm loc}
  :=
  \mathcal R_i\mathcal R_j
  \bigl(\zeta(V_\infty\otimes V_\infty)_{ij}\bigr).
\]
Then either
\[
  P_n-c_{n,R_1}(s)-P_\infty^{\rm loc}\to0
  \quad\text{strongly in }L^{3/2}(B_\rho\times I),
\]
and hence the pressure gradients converge locally modulo functions of time on
the retained branch,
or the branch has already entered the pressure/exterior row routed by
\cref{paper8:lem:harmonic-cross-pressure-routing}.
\end{lemma}

\begin{proof}
Choose \(m\) so large that
\[
  H^1(B_R)\Subset L^2(B_R)\hookrightarrow H^{-m}(B_{2R}).
\]
The local \(L^\infty_sL^2_y\), \(L^2_sH^1_y\), pressure, and modulation bounds
in the compactness package allow the repaired-gauge equation to be tested
against \(H^m_0(B_{2R})\).  The same estimates as in
\cref{paper5:lem:window-time-regularity} give
\[
  \sup_n\|\partial_sV_n\|_{L^1(I;H^{-m}(B_R))}<\infty .
\]
Aubin--Lions--Simon applied to
\[
  H^1(B_R)\Subset L^2(B_R)\hookrightarrow H^{-m}(B_R)
\]
therefore gives strong convergence in \(L^2(B_R\times I)\), after passing to a
subsequence.  The uniform \(L^\infty_sL^2_y\cap L^2_sH^1_y\) bound gives the
standard parabolic interpolation estimate
\[
  \sup_n\|V_n\|_{L^{10/3}(B_R\times I)}<\infty .
\]
Interpolating the strong \(L^2\) convergence with this uniform \(L^{10/3}\)
bound gives strong convergence in \(L^p(B_R\times I)\) for every
\(2\le p<10/3\), in particular for \(p=3\).

For the pressure, fix \(B_\rho\Subset B_{R_1}\Subset B_R\) and a cutoff
\(\zeta\equiv1\) on \(B_{R_1}\), \(\operatorname{supp}\zeta\Subset B_R\).  With
the local pressure \(P_\infty^{\rm loc}\) defined in the statement, the pressure
atlas gives, modulo functions of time on \(B_{R_1}\),
\[
  P_n-c_{n,R_1}(s)-P_\infty^{\rm loc}
  =
  \mathcal R_i\mathcal R_j\bigl(\zeta(V_n\otimes V_n
  -V_\infty\otimes V_\infty)_{ij}\bigr)+H_{n,R_1}.
\]
The strong \(L^3(B_R\times I)\) convergence gives
\[
  V_n\otimes V_n\to V_\infty\otimes V_\infty
  \quad\text{strongly in }L^{3/2}(B_R\times I),
\]
so Calder\'on--Zygmund continuity gives strong convergence of the local part in
\(L^{3/2}(B_{R_1}\times I)\).  If the mean-zero harmonic remainder
\(H_{n,R_1}\) does not converge to zero in
\(L^{3/2}(B_\rho\times I)\), then
\cref{paper8:lem:harmonic-cross-pressure-routing} with
\((q,p)=(3/2,3/2)\) assigns the branch to the pressure/exterior row.  Otherwise
the pressure convergence is strong locally modulo functions of time.  In all
cases pressure compactness has been obtained by a local decomposition or routed
to a named row; no weak lower-semicontinuity of pressure mass is used.
\end{proof}

\begin{lemma}[Pressure-only compactness failure is a pressure/exterior row]
\label{paper5:lem:selected-compactness-pressure-only-routing}
Let a selected scale-collapse window sequence satisfy the repaired-gauge,
local pressure-reconstruction, and velocity rows of
\cref{paper5:def:selected-scale-collapse-compactness-package} on
\(B_{2R}\times I\), and suppose that, after passing to a subsequence,
\[
  V_n\to V_\infty
  \quad\text{strongly in }L^3(B_R\times I).
\]
Fix \(B_\rho\Subset B_{R_1}\Subset B_R\) and
\(\zeta\in C_c^\infty(B_R)\) with \(\zeta\equiv1\) on \(B_{R_1}\), and set
\[
  P_\infty^{\rm loc}
  :=
  \mathcal R_i\mathcal R_j
  \bigl(\zeta(V_\infty\otimes V_\infty)_{ij}\bigr).
\]
Write the normalized pressure defect on \(B_{R_1}\), modulo functions of time,
as
\[
  Q_n
  :=
  P_n-c_{n,R_1}(s)-P_\infty^{\rm loc}
  =
  Q_n^{\rm loc}+H_{n,R_1},
\]
where
\[
  Q_n^{\rm loc}
  =
  \mathcal R_i\mathcal R_j\bigl(\zeta(V_n\otimes V_n
  -V_\infty\otimes V_\infty)_{ij}\bigr)
\]
and \(H_{n,R_1}\) is mean-zero harmonic on \(B_{R_1}\).  If the normalized
pressure defect has a positive lower bound on the smaller ball,
\[
  \limsup_{n\to\infty}
  \|Q_n\|_{L^{3/2}(B_\rho\times I)}>0,
\]
then the branch is in the ordered pressure/exterior alternative.
\end{lemma}

\begin{proof}
The strong \(L^3(B_R\times I)\) convergence gives
\[
  V_n\otimes V_n-V_\infty\otimes V_\infty
  \to0
  \quad\text{strongly in }L^{3/2}(B_R\times I).
\]
Calder\'on--Zygmund continuity gives
\[
  Q_n^{\rm loc}\to0
  \quad\text{strongly in }L^{3/2}(B_{R_1}\times I).
\]
Thus the positive lower bound for \(Q_n\) on \(B_\rho\times I\) forces
\[
  \limsup_{n\to\infty}
  \|H_{n,R_1}\|_{L^{3/2}(B_\rho\times I)}>0 .
\]
By
\cref{paper8:lem:harmonic-cross-pressure-routing}, together with the
quantitative annular-tail converse
\cref{paper8:lem:nested-ball-harmonic-tail}, such a nonvanishing harmonic part
can only come from pressure-reconstruction failure or from a far-field
pressure carrier.  These are exactly the pressure/exterior rows of the ordered
local state space.
\end{proof}

\begin{lemma}[Failure of selected compactness produces a local defect label]
\label{paper5:lem:selected-compactness-defect-label}
Let a represented genuine scale-collapse branch have a nonvanishing selected
core and be tested on canonical selected windows \(B_{2R}\times I\).  If the
selected scale-collapse compactness package of
\cref{paper5:def:selected-scale-collapse-compactness-package} fails on those
windows, then after passing to a subsequence the first failed row is one of the
following local labels:
\begin{enumerate}[label=\textup{(\roman*)}]
\item representation or pressure-reconstruction failure;
\item autonomous-modulation failure or modulation-driven recentering;
\item loss of retained active core, including pressure-only critical-mass
detection;
\item rough-core loss of local windowed \(H^1\) control;
\item exterior or separated-core concentration;
\item multicore, comparable-core, or same-point cascade;
\item a critical local \(L^3\) upper-mass failure routed by
\cref{paper6a:lem:terminal-local-critical-mass-routing}.
\end{enumerate}
If none of these labels occurs, the compactness package holds and
\cref{paper5:lem:selected-compactness-aubin-lions} gives strong local
compactness on the selected windows.
\end{lemma}

\begin{proof}
The rows in
\cref{paper5:def:selected-scale-collapse-compactness-package} are tested in
the displayed order.  Failure of the repaired-gauge equation, local energy
inequality, or pressure reconstruction is item \textup{(i)}.  Failure of the
local selected-window modulation bound is exactly the
autonomous-modulation or recentering row, giving item \textup{(ii)}.

If the retained spacetime activity floor fails, the selected candidate is not the
retained terminal component.  The terminal local critical-mass routing lemma
\cref{paper6a:lem:terminal-local-critical-mass-routing} records this as lower
active-core loss or, when the CKN signal is carried only by pressure, as the
pressure-only critical-mass row.  This gives item \textup{(iii)}.

Assume the preceding rows hold and the local velocity/pressure estimate row
fails.  If the failure is the \(L^2_sH^1_y\) part on a retained compact
cylinder, then it is the rough-core row by
\cref{paper6:thm:rough-core}, giving item \textup{(iv)}.  If the failure is
that the compact-cylinder \(L^\infty_sL^2_y\) or pressure mass cannot be
bounded on any fixed enclosing ball, apply a finite cover of the selected
retained core and the local pressure atlas.  Either a positive velocity
concentration remains in one member of the finite cover, or the mass leaves
every fixed member.  In the first case the ordered local state-space
decomposition reads a compact active label: a comparable retained core,
multicore label, or same-point cascade.  In the second case the label is
exterior or separated-core concentration.  If the positive diagnostic is
pressure-only, \cref{paper5:lem:selected-compactness-pressure-only-routing}
routes it to the pressure/exterior row.  This gives items
\textup{(iii)}, \textup{(v)}, and \textup{(vi)}.

It remains to consider failure of the critical local upper row.  By
\cref{paper6a:lem:terminal-local-critical-mass-routing}, such a failure
produces an additional compact active label at the same retained frame, a
comparable finite-mass retained class, a multibubble/multicore label, a
cascade label, a rough-core label, or an exterior/separated label.  These are
exactly items \textup{(iv)}--\textup{(vii)}.  Comparable same-point classes are
not left as a separate compactness defect: they are grouped by
\cref{paper6a:lem:retained-frame-maximality} and
\cref{paper8:lem:compound-same-point}.

If no listed first failure occurs, all rows in the selected scale-collapse
compactness package hold.  The local strong compactness conclusion then follows
from \cref{paper5:lem:selected-compactness-aubin-lions}.  Thus selected-window
compactness failure has no unlabelled terminal form.
\end{proof}

\begin{lemma}[Comparable compact defects are grouped before retention]
\label{paper5:lem:selected-compactness-comparable-defect}
Suppose a selected compactness defect produced by
\cref{paper5:lem:selected-compactness-defect-label} is based at the same
physical point as the retained selected scale-collapse core and has scale
comparable to the retained frame.  Then either it is absorbed into the retained
compound core before the retained scale-collapse state is declared, or the
branch has already entered the multicore/cascade/exterior local alternative.
\end{lemma}

\begin{proof}
If the relative scale-center parameters stay in a compact subset of the
same-point comparable chart, \cref{paper8:lem:compound-same-point} assembles
the corresponding positive local states into one compound state, and
\cref{paper6a:lem:retained-frame-maximality} says that any surviving comparable
retained class is already absorbed into the selected retained component.  If
the relative parameters leave every compact subset, the ordered state-space
compactification records separated scale, separated center, exterior leakage,
or same-point cascade before the retained compact scale-collapse state is
reached.  Hence a comparable compactness defect cannot remain as a separate
unresolved retained compactness failure.
\end{proof}

\begin{theorem}[Selected compactness failure is routed]
\label{paper5:thm:selected-compactness-failure-routed}
Let a represented genuine scale-collapse branch have nonvanishing selected
core and be tested on the canonical selected windows supplied by the local
state-space construction.  Then the compactness alternative in
\cref{paper5:thm:scale-collapse-stratification} is not an independent terminal
exit.  More precisely, either the selected scale-collapse compactness package
holds on every retained compact subcylinder after passing to a subsequence, or
the first failed row is routed to one of the named local alternatives:
representation/pressure failure, autonomous-modulation failure, loss of
retained active core or pressure-only critical mass, rough-core loss,
exterior/separated-core concentration, multicore/same-point cascade, or a
critical local upper-mass label routed by
\cref{paper6a:lem:terminal-local-critical-mass-routing}.
\end{theorem}

\begin{proof}
Choose a countable nested exhaustion of the retained compact selected core by
compact cylinders
\[
  B_{R_j}\times I_j\Subset B_{2R_j}\times I_j,
  \qquad j=1,2,\dots .
\]
Apply \cref{paper5:lem:selected-compactness-defect-label} to each member of
this exhaustion in the ordered local test.  If a row fails on some member, take
the least \(j\) for which a failure occurs and then the first failed row in the
ordered list on that cylinder.  The lemma assigns it to one of the named local
alternatives.  If comparable compact defects occur at the retained physical
point and comparable scale, then
\cref{paper5:lem:selected-compactness-comparable-defect} absorbs them into the
retained compound core or routes them to the multicore/cascade/exterior rows.

If no row fails on the exhaustion, the package holds on each compact
subcylinder.  Apply
\cref{paper5:lem:selected-compactness-aubin-lions} on the first cylinder, then
on the second along the extracted subsequence, and continue diagonally.  The
diagonal subsequence gives strong local \(L^3\)
compactness and pressure compactness modulo functions of time, with any
nonvanishing harmonic pressure remainder routed by
\cref{paper8:lem:harmonic-cross-pressure-routing}.  Therefore compactness
failure is represented by a local first-failure
classification.  It either returns a compact retained sequence or enters a
more specific named local row; it is not a separate terminal obstruction.
\end{proof}

\begin{lemma}[Retained autonomous scale collapse enters the scale-rigid stratum]
\label{paper5:lem:autonomous-scale-collapse-is-scale-rigid}
Let a represented Type II branch be in the retained compact scale-collapse
state of \cref{paper5:def:retained-compact-scale-collapse-state}.  Suppose it
has thick autonomous windows and satisfies
\cref{paper5:prop:selected-window-suitable-estimates} on the same window sequence.  Then
the shifted selected-window sequence
\[
  V_n(y,s)=V(y,\tau_n+s),\qquad
  P_n(y,s)=P(y,\tau_n+s),
\]
restricted to a fixed compact cylinder carrying the retained active mass, is a
scale-rigid local branch in the sense of
\cref{paper3:def:scale-rigid-branch}, or it is the compact single-core branch
excluded by \cref{paper1:thm:single-core-criterion}, unless the branch has
already entered one of the named local alternatives excluded from the retained
compact state.
\end{lemma}

\begin{proof}
The compact-cylinder suitable bounds in
\cref{paper5:def:retained-compact-scale-collapse-state}\textup{(ii)} give
\cref{paper3:def:scale-rigid-branch}\textup{(R1)} on a fixed compact selected
cylinder after shrinking once.  The retained active CKN mass floor in
\textup{(v)} gives \textup{(R2)} on an inner cylinder.  Bounded modulation in
\textup{(iv)} gives \textup{(R3)}.

It remains to verify the scale-rigid structural clause \textup{(R4)}.  Apply
the minimal active-family extraction
\cref{paper3:thm:minimal-bad-extraction} to the selected retained core and to
the nearest previously selected active scale-center, if such a scale-center is
present.  If no such active relation is present, the branch is the compact
single-core state rather than a genuine scale-collapse residual; this is the
single-core branch excluded by \cref{paper1:thm:single-core-criterion}, not a
remaining case in this lemma.

Otherwise the pair is an input to the same-point or separated-point reduction,
\cref{paper3:thm:same-point-reduction,paper3:thm:separated-point-reduction}.
Those reductions are an ordered partition.  Their non-rigid outputs are exactly
the named exits: another active core at comparable scale, a strict same-point
cascade not yet absorbed, a separated physical point, an exterior
pressure/concentration label, a rough-core label, or failure of the compact
estimates/modulation convergence needed for the autonomous scale-collapse
passage.  The comparable, same-point, separated, and exterior outcomes are the
multicore, same-point cascade, or separated/exterior alternatives of
\cref{paper3:def:local-multibubble}; the rough-core outcome is the local
windowed-\(H^1\) failure alternative; and the last two outcomes are the
compactness or autonomous-modulation alternatives of
\cref{paper5:thm:scale-collapse-stratification}.  All are excluded from the
retained compact state by
\cref{paper5:def:retained-compact-scale-collapse-state}\textup{(vii)}.

After those exits are removed, the same ordered reductions leave only the
structural outputs named in
\cref{paper3:def:scale-rigid-branch}\textup{(R4)}: either the relative
scale-center Jacobian degenerates, or the selected core is the innermost member
after all strict outer same-point states have been absorbed into the
perturbative remainder by \cref{paper3:lem:outer-vanish}.  Thus the shifted
sequence is produced by the scale-rigid reduction and satisfies \textup{(R4)}
on the selected compact cylinder.
\end{proof}

\begin{theorem}[Scale-collapse first failures are routed local alternatives]
\label{paper5:thm:scale-collapse-first-failure-closure}
Let a represented branch have genuine scale collapse and a nonvanishing
selected core.  If one of the retained compact scale-collapse inputs fails
before the branch reaches
\cref{paper5:def:retained-compact-scale-collapse-state}, then the first failure
is one of the named local alternatives: finite scale cost, finite absolute
scale cost, autonomous-modulation failure, compactness-package first failure
routed to a named local row,
exterior/separated-core loss, rough-core loss, multicore/same-point cascade, or
scale-rigid residual behavior; the thin-drift first failure is routed by
\cref{paper6a:thm:thin-negative-drift-closed} to one of these rows or to a
closed retained row.  None of these first failures remains inside the retained
compact scale-collapse branch.
\end{theorem}

\begin{proof}
Apply the ordered scale-collapse stratification
\cref{paper5:thm:scale-collapse-stratification} and the first-failure ledger of
\cref{paper5:lem:retained-scale-collapse-first-failure}.  Finite scale cost and
finite absolute scale cost are the cost alternatives excluded by
\cref{paper5:thm:cost-exclusion,paper5:cor:absolute-cost-exclusion}.  Thin
drift is first recorded as the scale-drift boundary face of the scale-collapse
stratification, but
\cref{paper6a:thm:thin-negative-drift-closed} proves that it is not terminal:
persistent windows return to the thick-window ledger, while diffuse occupation
blocks route to another named row or to a closed retained row.  Failure of
autonomous modulation is a
repaired-gauge/modulation boundary face.  Failure
of compactness on the selected windows is routed by
\cref{paper5:thm:selected-compactness-failure-routed}: either the local
compactness package holds after passing to a subsequence, or the first failed
row is a representation/pressure row, active-core or pressure-only row,
rough-core row, exterior/separated-core row, multicore/same-point cascade row,
or critical local upper-mass row.  Comparable same-point defects are grouped by
\cref{paper5:lem:selected-compactness-comparable-defect}.  Thus compactness
failure is not a separate retained terminal state.  The
remaining rigid residual is the scale-rigid local branch discharged by
\cref{paper3:prop:scale-rigidity-discharged}.  These are exactly the named
local alternatives used in the state-space decomposition, so none is a
retained compact scale-collapse state.
\end{proof}

\begin{theorem}[Compact scale-collapse rigidity from local routing]
\label{paper5:thm:compact-scale-collapse-rigidity-discharged}
Let a represented Type II branch have genuine scale collapse, nonvanishing
selected core, and belong to the retained compact scale-collapse state of
\cref{paper5:def:retained-compact-scale-collapse-state}.  Then no compact
scale-collapse branch remains.  In particular, compact scale-collapse rigidity
is supplied by the retained local state-space branch itself, not by an
independent stationary input.
\end{theorem}

\begin{proof}
Apply the scale-collapse stratification
\cref{paper5:thm:scale-collapse-stratification}.  The finite scale-collapsing
cost and finite absolute cost cases are excluded by
\cref{paper5:thm:cost-exclusion,paper5:cor:absolute-cost-exclusion}.

The thin-drift, autonomous-modulation, and compactness alternatives are not
discarded by definition.  A branch satisfying the retained compact
scale-collapse-state hypotheses has thick autonomous selected windows,
selected-window estimates, and autonomous modulation convergence by
\cref{paper5:def:retained-compact-scale-collapse-state}\textup{(ii),(iii)}.
Thus the thin-drift, autonomous-modulation, and compactness alternatives of
\cref{paper5:thm:scale-collapse-stratification} are incompatible with the
present retained state.

It remains to consider the thick autonomous case.  By
\cref{paper5:lem:autonomous-scale-collapse-is-scale-rigid}, the shifted
selected-window sequence is either a scale-rigid local branch or the compact
single-core branch, unless the branch has already left the retained compact
state.  Since the present branch is retained, the exit case is unavailable.  If
the compact single-core branch occurs, it is excluded by
\cref{paper1:thm:single-core-criterion}.  If the scale-rigid branch occurs,
\cref{paper3:prop:scale-rigidity-discharged} applies: the branch either has a
nonzero bounded selected-window limit, which leaves the strict
selected-window retained branch by
\cref{paper3:lem:scale-rigid-bounded-limit-exit}, or it reduces
to the compact single-core branch, again excluded by
\cref{paper1:thm:single-core-criterion}.  Hence the thick autonomous retained
case is impossible.

\end{proof}

\begin{theorem}[Scale-collapse alternatives]\label{paper5:thm:scale-collapse-alternatives}
Let a represented Type II branch have genuine scale collapse and a nonvanishing selected core.
Then one of the following mutually exclusive outcomes occurs:
\begin{enumerate}[label=(\roman*)]
  \item the finite scale-collapsing cost exclusion applies;
  \item the finite absolute cost exclusion applies;
  \item the branch has a thin-drift alternative, which is routed by
  \cref{paper6a:thm:thin-negative-drift-closed};
  \item the branch has an autonomous-modulation alternative;
  \item the branch has a compactness-package first failure, which is routed by
  \cref{paper5:thm:selected-compactness-failure-routed};
  \item the branch enters the retained compact scale-collapse state, in which
  case \cref{paper5:thm:compact-scale-collapse-rigidity-discharged} excludes it.
\end{enumerate}
\end{theorem}

\begin{proof}
Apply \cref{paper5:thm:scale-collapse-stratification}.  The finite-cost alternatives are
items \textup{(i)} and \textup{(ii)}.  The thin-drift, modulation, and
compactness alternatives are items \textup{(iv)}--\textup{(vi)} of that
theorem; the thin-drift item is routed by
\cref{paper6a:thm:thin-negative-drift-closed}, and the compactness item is
routed by \cref{paper5:thm:selected-compactness-failure-routed}.  It remains
only to treat item \textup{(iii)}, the thick autonomous branch with
selected-window estimates and autonomous modulation convergence.  If one of the
retained compact-state inputs in
\cref{paper5:def:retained-compact-scale-collapse-state} fails, then
\cref{paper5:thm:scale-collapse-first-failure-closure} assigns the branch to a
named local alternative, hence to one of the preceding nonretained outcomes.
Otherwise the branch is in the retained compact scale-collapse state and
\cref{paper5:thm:compact-scale-collapse-rigidity-discharged} excludes it.
\end{proof}

\clearpage
\section{Critical Tightness and Windowed-Control Estimates}
\label{sec:c6-c7-cost-exclusion}

This section records the local estimates used to exclude failure of uniform
critical tightness and failure of local windowed \(H^1\) control.  The
statements are phrased as analytic alternatives for a renormalized
Navier--Stokes solution \(V(\tau)\).

\subsection{Uniform critical tightness}

\begin{definition}[Uniform critical tightness]\label{paper7:def:critical-tightness}
A renormalized orbit \(V(\tau)\), \(\tau\ge\tau_0\), is uniformly critically
tight if
\[
  \forall\varepsilon>0\ \exists R_\varepsilon<\infty:
  \sup_{\tau\ge\tau_0}
  \int_{|y|>R_\varepsilon}|V(y,\tau)|^3\,dy<\varepsilon.
\]
The non-tight alternative is the failure of this property.
\end{definition}

\begin{lemma}[Critical precompactness implies uniform tightness]
\label{paper7:lem:global-l3-compactness-tightness}
If
\[
  \mathcal O:=\{V(\tau):\tau\ge\tau_0\}
\]
is contained in a compact subset of \(L^3(\mathbb R^3)\), then \(V\) is
uniformly critically tight.
\end{lemma}

\begin{proof}
Compact subsets of \(L^3(\mathbb R^3)\) are uniformly tight.  Let
\(K\subset L^3\) be compact and fix \(\varepsilon>0\).  Set
\(\delta=\varepsilon^{1/3}/2\).  Choose a finite \(\delta\)-net
\(f_1,\dots,f_N\) in \(L^3\).  For each \(i\), choose \(R_i\) so that
\[
  \|f_i\|_{L^3(|y|>R_i)}<\delta.
\]
Let \(R=\max_iR_i\).  For any \(f\in K\), choose \(i\) with
\(\|f-f_i\|_3<\delta\).  Then
\[
  \|f\|_{L^3(|y|>R)}
  \le
  \|f-f_i\|_3+\|f_i\|_{L^3(|y|>R)}
  <\varepsilon^{1/3}.
\]
Taking the cube gives
\[
  \int_{|y|>R}|f(y)|^3\,dy<\varepsilon.
\]
Applying this to the compact set containing \(\mathcal O\) gives uniform
critical tightness.
\end{proof}

\begin{definition}[Compact core with vanishing critical remainder]
\label{paper7:def:compact-core-vanishing-remainder}
The orbit has a compact critical core with vanishing remainder if there exist
a compact set \(\mathcal K^{\rm c}\subset L^3(\mathbb R^3)\), a map
\(Q:[\tau_0,\infty)\to \mathcal K^{\rm c}\), and a remainder
\(r(\tau)\in L^3(\mathbb R^3)\), such that
\[
  V(\tau)=Q(\tau)+r(\tau),\qquad
  \lim_{\tau\to\infty}\|r(\tau)\|_{L^3}=0,
\]
and, for every \(T>\tau_0\), the set
\[
  \{V(\tau):\tau\in[\tau_0,T]\}
\]
is uniformly \(L^3\)-tight.
\end{definition}

\begin{lemma}[Compact core plus vanishing remainder implies tightness]
\label{paper7:lem:core-remainder-tightness}
If \cref{paper7:def:compact-core-vanishing-remainder} holds, then \(V\) is
uniformly critically tight.
\end{lemma}

\begin{proof}
Fix \(\varepsilon>0\) and set \(\delta=\varepsilon^{1/3}\).  By compactness of
\(\mathcal K^{\rm c}\), apply
\cref{paper7:lem:global-l3-compactness-tightness} to \(\mathcal K^{\rm c}\)
and choose \(R_1\) such that
\[
  \sup_{q\in \mathcal K^{\rm c}}\|q\|_{L^3(|y|>R_1)}<\delta/3.
\]
Since \(\|r(\tau)\|_3\to0\), choose \(T\) so that
\[
  \|r(\tau)\|_3<\delta/3
  \qquad\text{for all }\tau\ge T.
\]
Then for \(\tau\ge T\),
\[
  \|V(\tau)\|_{L^3(|y|>R_1)}
  \le
  \|Q(\tau)\|_{L^3(|y|>R_1)}+\|r(\tau)\|_3
  <2\delta/3.
\]
By finite-initial-window tightness on \([\tau_0,T]\), choose \(R_2\) such that
\[
  \sup_{\tau\in[\tau_0,T]}
  \int_{|y|>R_2}|V(y,\tau)|^3\,dy<\varepsilon.
\]
With \(R=\max(R_1,R_2)\), the \(L^3\)-norm of the tail of \(V(\tau)\) is
\(<\delta\) for all \(\tau\ge T\), while the tail integral is
\(<\varepsilon\) on \([\tau_0,T]\).  Cubing on the tail \(\tau\ge T\) gives
the required \(L^3\)-tail integral \(<\varepsilon\) for all
\(\tau\ge\tau_0\).
\end{proof}

\begin{definition}[No-splitting profile condition]\label{paper7:def:no-splitting-condition}
A critical profile decomposition has no critical-mass splitting if it has:
\begin{enumerate}[label=\textup{(\roman*)}]
\item a primary compact core family \(\mathcal K^{\rm c}\subset L^3\);
\item asymptotic \(L^3\)-decoupling of core, secondary profiles, and
remainder;
\item single-profile saturation, so every non-primary profile has zero
critical \(L^3\)-mass;
\item a remainder whose \(L^3\)-norm vanishes on the renormalized tail;
\item finite-initial-window tightness.
\end{enumerate}
\end{definition}

\begin{lemma}[No splitting gives compact core with vanishing remainder]
\label{paper7:lem:no-splitting-gives-core-rem}
If \cref{paper7:def:no-splitting-condition} holds, then
\cref{paper7:def:compact-core-vanishing-remainder} holds.
\end{lemma}

\begin{proof}
By single-profile saturation and nonnegative \(L^3\)-mass decoupling, every
secondary profile has zero critical \(L^3\)-mass and therefore vanishes in
\(L^3\).  The profile decomposition then reduces to a primary compact core
plus a remainder whose \(L^3\)-norm tends to zero on the renormalized tail.
These provide the core \(Q(\tau)\), compact set \(\mathcal K^{\rm c}\), and
remainder \(r(\tau)\) required by
\cref{paper7:def:compact-core-vanishing-remainder}, with
finite-initial-window tightness supplied as part of
\cref{paper7:def:no-splitting-condition}.
\end{proof}

\begin{corollary}[No splitting gives uniform critical tightness]
\label{paper7:cor:no-splitting-critical-tightness}
If \cref{paper7:def:no-splitting-condition} holds, then \(V\) is uniformly
critically tight.
\end{corollary}

\begin{proof}
\Cref{paper7:lem:no-splitting-gives-core-rem} gives a compact core with
vanishing critical remainder, and
\cref{paper7:lem:core-remainder-tightness} gives uniform critical tightness.
\end{proof}

\begin{definition}[Finite critical profile approximation]
\label{paper7:def:finite-profile-tight}
The orbit admits a finite critical profile approximation if there are finitely
many profiles
\[
  \mathcal B^{II}_{NS}:=\{B_1,\dots,B_N\}\subset L^3(\mathbb R^3)
\]
such that, for every \(\delta>0\), there is \(\tau_\delta\) with
\[
  \forall\tau\ge\tau_\delta\ \exists i(\tau)\in\{1,\dots,N\}:
  \|V(\tau)-B_{i(\tau)}\|_{L^3}<\delta,
\]
and the finite initial window
\[
  \{V(\tau):\tau\in[\tau_0,\tau_\delta]\}
\]
is uniformly \(L^3\)-tight for each \(\delta>0\).
\end{definition}

\begin{lemma}[Finite critical profile approximation implies tightness]
\label{paper7:lem:finite-profile-tight}
If \cref{paper7:def:finite-profile-tight} holds, then \(V\) is uniformly
critically tight.
\end{lemma}

\begin{proof}
Fix \(\varepsilon>0\) and set \(\delta=\varepsilon^{1/3}\).  Since the profile family is finite and each
\(B_i\in L^3(\mathbb R^3)\), there is \(R_1\) such that
\[
  \max_{1\le i\le N}\|B_i\|_{L^3(|y|>R_1)}<\delta/3.
\]
Choose the approximation tolerance \(\delta/3\).  For
\(\tau\ge\tau_{\delta/3}\), pick \(i(\tau)\) with
\(\|V(\tau)-B_{i(\tau)}\|_3<\delta/3\).  Then
\[
  \|V(\tau)\|_{L^3(|y|>R_1)}
  \le
  \|V(\tau)-B_{i(\tau)}\|_3
  +\|B_{i(\tau)}\|_{L^3(|y|>R_1)}
  <2\delta/3.
\]
By finite-initial-window tightness, choose \(R_2\) so that the \(L^3\) tail is
\(<\varepsilon\) on \([\tau_0,\tau_{\delta/3}]\).  With
\(R=\max(R_1,R_2)\), the tail integral is \(<\varepsilon\) for every
\(\tau\ge\tau_0\), so the orbit is uniformly \(L^3\)-tight.
\end{proof}

\begin{definition}[Compact parameter realization]\label{paper7:def:compact-param-realization}
The orbit has a compact parameter realization if there exist a compact
parameter set \(\Theta_{II}\), a continuous map
\[
  \Theta_{II}\ni\theta\mapsto V_\theta\in L^3(\mathbb R^3),
\]
a curve \(\theta(\tau)\in\Theta_{II}\), and a remainder \(r(\tau)\to0\) in
\(L^3\), such that
\[
  V(\tau)=V_{\theta(\tau)}+r(\tau),
\]
and, for every \(T>\tau_0\), the set
\[
  \{V(\tau):\tau\in[\tau_0,T]\}
\]
is uniformly \(L^3\)-tight.
\end{definition}

\begin{lemma}[Compact realization implies tightness]
\label{paper7:lem:compact-realization-tight}
If \cref{paper7:def:compact-param-realization} holds, then \(V\) is uniformly
critically tight.
\end{lemma}

\begin{proof}
The continuous image of compact \(\Theta_{II}\) in \(L^3\) is compact.  Thus
\[
  \mathcal K^{\rm c}:=\{V_\theta:\theta\in\Theta_{II}\}
\]
is compact in \(L^3\).  The representation
\(V(\tau)=V_{\theta(\tau)}+r(\tau)\) with \(r(\tau)\to0\) in \(L^3\) gives
the compact-core vanishing-remainder situation, because finite-initial-window
tightness is included in \cref{paper7:def:compact-param-realization}.
\Cref{paper7:lem:core-remainder-tightness} gives uniform critical tightness.
\end{proof}

\begin{theorem}[Exclusion of the non-tight alternative]\label{paper7:thm:critical-tightness-criteria}
For a renormalized Type II Navier--Stokes solution, any one of the following
conditions gives uniform critical tightness:
\begin{enumerate}[label=\textup{(\roman*)}]
\item \(\mathcal O\Subset L^3(\mathbb R^3)\);
\item compact core with vanishing critical remainder;
\item no-splitting profile condition;
\item finite critical profile approximation;
\item compact parameter realization.
\end{enumerate}
\end{theorem}

\begin{proof}
The five alternatives are handled respectively by
\cref{paper7:lem:global-l3-compactness-tightness,paper7:lem:core-remainder-tightness,paper7:cor:no-splitting-critical-tightness,paper7:lem:finite-profile-tight,paper7:lem:compact-realization-tight}.
In every case \(V\) is uniformly critically tight, so the failure of uniform
critical tightness is unavailable.
\end{proof}

\begin{corollary}[After exclusion of the non-tight alternative]\label{paper7:cor:critical-tightness-reduces-alternatives}
Assume a finite-cost renormalized Type II solution not excluded by the compact
cost-divergence criterion has a two-alternative classification whose only
remaining alternatives are failure of uniform critical tightness and failure of
tail windowed \(H^1\) control.  If
\cref{paper7:def:critical-tightness} holds, then the only remaining alternative is
failure of tail windowed \(H^1\) control.
\end{corollary}

\begin{proof}
The two-alternative classification gives exactly one of the two alternatives.  Uniform
critical tightness rules out failure of tightness, so only the windowed
\(H^1\)-failure alternative remains.
\end{proof}

\subsection{\texorpdfstring{Tail windowed \(H^1\) control}{Tail windowed H1 control}}

\begin{lemma}[Uniform local pressure control]\label{paper7:lem:uniform-pressure-from-c3}
Assume
\[
  M_{3,R}:=\sup_{\tau\ge\tau_0}\|V(\tau)\|_{L^3(B_{4R})}<\infty
\]
and
\[
  -\Delta P=\partial_i\partial_j(V_iV_j)
  \quad\text{in }\mathbb R^3
\]
for each \(\tau\ge\tau_0\).  Assume also the local tail estimate
\[
  \mathcal T_R(\tau):=
  \sup_{x\in B_R}\int_{|z|>2R}
  \frac{|V(z,\tau)|^2}{|x-z|^5}\,dz,
  \qquad
  \sup_{\tau\ge\tau_0}\mathcal T_R(\tau)\le CR^{-4}M_{3,R}^2 .
\]
Then there is a measurable choice of constants \(c_R(\tau)\in\mathbb R\) such
that
\[
  \sup_{\tau\ge\tau_0}
  \|P(\tau)-c_R(\tau)\|_{L^{3/2}(B_R)}
  \le C_RM_{3,R}^2 .
\]
\end{lemma}

\begin{proof}
The local pressure decomposition gives
\[
  \|P(\tau)-c_R(\tau)\|_{L^{3/2}(B_R)}
  \le
  C\left(
  \|V(\tau)\|_{L^3(B_{4R})}^2+R^2\mathcal T_R(\tau)
  \right).
\]
The tail hypothesis gives
\[
  \sup_{\tau\ge\tau_0}\mathcal T_R(\tau)\le CR^{-4}M_{3,R}^2,
\]
and the local critical bound gives
\[
  \|V(\tau)\|_{L^3(B_{4R})}\le M_{3,R}.
\]
Combining the two estimates yields
\[
  \sup_{\tau\ge\tau_0}
  \|P(\tau)-c_R(\tau)\|_{L^{3/2}(B_R)}
  \le C_RM_{3,R}^2 .
\]
The constants \(c_R(\tau)\) can be chosen by fixing the displayed pressure
decomposition, hence measurably.
\end{proof}

\begin{proposition}[Renormalized Caccioppoli gives uniform windowed gradients]
\label{paper7:prop:uniform-windowed-gradient}
Assume \(V,P\) are smooth on compact renormalized cylinders and solve
\[
  \partial_\tau V +(V\cdot\nabla)V+\nabla P
  =
  \nu\Delta V+a(\tau)(V+y\cdot\nabla V)+b(\tau)\cdot\nabla V,
  \qquad \nabla\cdot V=0.
\]
Assume local compact-cylinder critical control on the outer cylinders used in
the estimate, bounded modulation
\[
  M_{ab}:=\sup_{\tau\ge\tau_0}(|a(\tau)|+|b(\tau)|)<\infty,
\]
and pressure normalized as in \cref{paper7:lem:uniform-pressure-from-c3}.  Then
for every \(R>0\) there exists
\[
  C_R=C(R,\nu,M_{3,R},M_{ab})
\]
such that
\[
  \sup_{T\ge\tau_0+1}
  \int_T^{T+1}\int_{B_R}|\nabla V(y,\tau)|^2\,dy\,d\tau
  \le C_R.
\]
\end{proposition}

\begin{proof}
Fix \(R>0\) and \(T\ge\tau_0+1\).  Use the renormalized Caccioppoli estimate
with outer radius \(2R\), inner radius \(R\), and time intervals
\[
  I_{2R}:=(T-1,T+2),\qquad I_R:=[T,T+1].
\]
Thus \(R_{\rm out}-R_{\rm in}=R\) and the inner time length is \(1\).  Keeping
the pressure term before taking absolute values gives
\[
\begin{aligned}
\nu\int_T^{T+1}\int_{B_R}|\nabla V|^2
\le{}&
C_R\iint_{(T-1,T+2)\times B_{2R}}|V|^2
+C_R\iint_{(T-1,T+2)\times B_{2R}}|V|^3\\
&+C_R\mathcal P_{T,R}
+C_RM_{ab}\iint_{(T-1,T+2)\times B_{2R}}|V|^2,
\end{aligned}
\]
where \(C_R\) absorbs the powers of \(R^{-1}\), \(R^{-2}\), and \(\nu\), and
\[
  \mathcal P_{T,R}:=
  \left|
  \iint_{(T-1,T+2)\times B_{2R}}P\,V\cdot\nabla\zeta
  \right|.
\]
The \(L^2\)-term is controlled by the local critical bound and boundedness of
\(B_{2R}\):
\[
  \int_{B_{2R}}|V(\tau)|^2\,dy
  \le |B_{2R}|^{1/3}\|V(\tau)\|_{L^3(B_{2R})}^2
  \le C_RM_{3,R}^2.
\]
Since the outer time interval has length \(3\),
\[
  \iint_{(T-1,T+2)\times B_{2R}}|V|^2\le C_RM_{3,R}^2,
\]
and similarly
\[
  \iint_{(T-1,T+2)\times B_{2R}}|V|^3\le 3M_{3,R}^3.
\]
For any time-dependent constant \(c(\tau)\),
\[
  \int_{B_{2R}}c(\tau)V\cdot\nabla\zeta\,dy
  =
  -\int_{B_{2R}}c(\tau)\zeta\,\nabla\cdot V\,dy=0.
\]
Thus \(P\) may be replaced by \(P-c_{2R}(\tau)\).  Hölder's inequality gives
\[
  \mathcal P_{T,R}
  \le
  C_R
  \iint_{(T-1,T+2)\times B_{2R}}
  |P-c_{2R}(\tau)|\,|V|\,dy\,d\tau .
\]
For almost every \(\tau\),
\[
  \int_{B_{2R}}|P-c_{2R}(\tau)|\,|V(\tau)|\,dy
  \le
  \|P(\tau)-c_{2R}(\tau)\|_{L^{3/2}(B_{2R})}
  \|V(\tau)\|_{L^3(B_{2R})}.
\]
The pressure estimate with radius \(2R\) gives
\[
  \|P(\tau)-c_{2R}(\tau)\|_{L^{3/2}(B_{2R})}
  \le C_RM_{3,R}^2,
\]
and the local critical bound gives
\[
  \|V(\tau)\|_{L^3(B_{2R})}\le M_{3,R}.
\]
Hence \(\mathcal P_{T,R}\le C_RM_{3,R}^3\).  Substituting these estimates into
Caccioppoli yields
\[
  \nu\int_T^{T+1}\int_{B_R}|\nabla V|^2
  \le
  C_R\left(M_{3,R}^2+M_{3,R}^3+M_{ab}M_{3,R}^2\right),
\]
uniformly in \(T\ge\tau_0+1\).  Dividing by \(\nu\) and renaming the constant
proves the claim.
\end{proof}

\begin{corollary}[Tail windowed local \(H^1\) bound]\label{paper7:cor:windowed-h1}
Under the hypotheses of \cref{paper7:prop:uniform-windowed-gradient}, for
every \(R>0\),
\[
  \sup_{T\ge\tau_0+1}
  \int_T^{T+1}\|V(\tau)\|_{H^1(B_R)}^2\,d\tau<\infty.
\]
\end{corollary}

\begin{proof}
The gradient part is \cref{paper7:prop:uniform-windowed-gradient}.  The
\(L^2\)-part follows from the same bounded-set H\"older estimate:
\[
  \int_T^{T+1}\|V(\tau)\|_{L^2(B_R)}^2\,d\tau
  \le
  C_R\int_T^{T+1}\|V(\tau)\|_{L^3(B_{2R})}^2\,d\tau
  \le C_RM_{3,R}^2.
\]
Adding the two bounds gives the asserted \(H^1(B_R)\) estimate.
\end{proof}

\begin{corollary}[Exclusion of failure of local windowed \(H^1\) control]
\label{paper7:cor:windowed-h1-exclusion-tail}
Assume the renormalized Type II solution satisfies the repaired-gauge
renormalized Navier--Stokes equation on compact cylinders, local
compact-cylinder critical control
\[
  \sup_{\tau\ge\tau_0}\|V(\tau)\|_{L^3(B_{2R})}<\infty
\]
for each relevant \(R\), bounded modulation
\[
  \sup_{\tau\ge\tau_0}(|a(\tau)|+|b(\tau)|)<\infty,
\]
and pressure reconstruction
\[
  -\Delta P=\partial_i\partial_j(V_iV_j).
\]
Then failure of local windowed \(H^1\) control,
\[
  \exists m\ge1:
  \sup_{T\ge\tau_0+1}
  \int_T^{T+1}\|V(\tau)\|_{H^1(B_m)}^2\,d\tau=\infty
\]
is impossible.
\end{corollary}

\begin{proof}
Apply \cref{paper7:cor:windowed-h1} with \(R=m\).  The resulting finite bound
contradicts the displayed failure of local windowed \(H^1\) control.
\end{proof}

\subsection{Modulation and weighted forcing inputs}

\begin{definition}[Modulation matrix Schur data]\label{paper7:def:modmatrix-schur}
Write the repaired modulation matrix in block form
\[
  M(V)=
  \begin{pmatrix}
  \alpha(V)&u(V)^T\\
  v(V)&A(V)
  \end{pmatrix}.
\]
Assume
\[
  \alpha(V(\tau))\ge \alpha_0>0,\qquad
  \|A(V(\tau))^{-1}\|_{\infty\to\infty}\le C_A,
\]
\[
  \|u(V(\tau))\|_\infty\le C_u',
  \qquad
  \|v(V(\tau))\|_\infty\le C_v',
\]
and
\[
  S(V(\tau)):=
  \alpha(V(\tau))-u(V(\tau))^TA(V(\tau))^{-1}v(V(\tau))
  \ge\sigma_0>0.
\]
\end{definition}

\begin{lemma}[Uniform inverse from Schur data]\label{paper7:lem:modmatrix-inverse}
Under \cref{paper7:def:modmatrix-schur},
\[
  \sup_{\tau\ge\tau_0}\|M(V(\tau))^{-1}\|_{\infty\to\infty}
  \le
  \sigma_0^{-1}
  +3\sigma_0^{-1}C_u'C_A
  +\sigma_0^{-1}C_A C_v'
  +C_A
  +3C_A^2C_v'C_u'\sigma_0^{-1}.
\]
\end{lemma}

\begin{proof}
For each \(\tau\), write
\[
  M(V(\tau))=
  \begin{pmatrix}
  \alpha&u^T\\
  v&A
  \end{pmatrix},
  \qquad
  S=\alpha-u^TA^{-1}v.
\]
By the translation-block inverse bound, \(A^{-1}\) exists and is uniformly
bounded.  By the Schur gap, \(S^{-1}\) exists and
\(|S^{-1}|\le\sigma_0^{-1}\).  The block inverse formula gives
\[
  M^{-1}=
  \begin{pmatrix}
  S^{-1} & -S^{-1}u^TA^{-1}\\
  -A^{-1}vS^{-1} & A^{-1}+A^{-1}vS^{-1}u^TA^{-1}
  \end{pmatrix}.
\]
Using the \(\ell^\infty\) matrix norm and the finite dimension of the mixed
blocks,
\[
  \|S^{-1}u^TA^{-1}\|_{\infty\to\infty}
  \le 3\sigma_0^{-1}C_u'C_A,
\]
\[
  \|A^{-1}vS^{-1}\|_{\infty\to\infty}
  \le \sigma_0^{-1}C_A C_v',
\]
and
\[
  \|A^{-1}vS^{-1}u^TA^{-1}\|_{\infty\to\infty}
  \le 3C_A^2C_v'C_u'\sigma_0^{-1}.
\]
Adding the four block estimates gives the displayed uniform bound.
\end{proof}

\begin{lemma}[Scale-normalization nondegeneracy gives Schur data]
\label{paper7:lem:gauge-nondeg-schur}
Assume the scale normalization constraint is satisfied, the scale derivative obeys
\[
  \alpha(V)=DG_{\rm sc}(V)[Z_{\rm sc}(V)]=p\Theta_0>0,
\]
the translation block has a uniform inverse bound, and the mixed blocks obey
the weighted first-moment bound and the Schur-compatible estimate
\[
  \alpha(V(\tau))-u(V(\tau))^TA(V(\tau))^{-1}v(V(\tau))
  \ge \frac{\alpha_0}{2}.
\]
Then the hypotheses of \cref{paper7:def:modmatrix-schur} hold with
\(\sigma_0=\alpha_0/2\).
\end{lemma}

\begin{proof}
The scale-normalization identity gives
\(\alpha(V)=DG_{\rm sc}(V)[Z_{\rm sc}(V)]=p\Theta_0\), hence the scale
transversality bound.  The translation-block inverse hypothesis gives the
uniform bound on \(A^{-1}\).  The mixed-block hypotheses give the bounds on
\(u\) and \(v\), and the displayed estimate is the Schur gap with
\(\sigma_0=\alpha_0/2\).
\end{proof}

\begin{definition}[Weighted scale-force components]\label{paper7:def:scale-force-pieces}
Let \(0<p<3\) and set
\[
  H(V,P)=\nu\Delta V-(V\cdot\nabla)V-\nabla P.
\]
Define
\[
  I_\Delta(\tau):=\int_{\mathbb R^3}|y|^{-p}|V|V\cdot\Delta V\,dy,
\]
\[
  I_{\rm nl}(\tau):=\int_{\mathbb R^3}|y|^{-p}|V|V\cdot((V\cdot\nabla)V)\,dy,
\]
and
\[
  I_P(\tau):=\int_{\mathbb R^3}|y|^{-p}|V|V\cdot\nabla P\,dy,
\]
whenever these integrals are absolutely convergent or are defined by the
distributional pairing used for the renormalized solution.
\end{definition}

\begin{lemma}[Scale-force decomposition]\label{paper7:lem:scale-force-decomp}
If
\[
  \sup_\tau|I_\Delta(\tau)|\le C_\Delta,\qquad
  \sup_\tau|I_{\rm nl}(\tau)|\le C_{\rm nl},\qquad
  \sup_\tau|I_P(\tau)|\le C_P,
\]
then
\[
  \sup_{\tau\ge\tau_0}
  \left|
  \int |y|^{-p}|V|V\cdot H(V,P)\,dy
  \right|
  \le \nu C_\Delta+C_{\rm nl}+C_P.
\]
\end{lemma}

\begin{proof}
By the definition of \(H\),
\[
  \int |y|^{-p}|V|V\cdot H(V,P)\,dy
  =
  \nu I_\Delta-I_{\rm nl}-I_P.
\]
The triangle inequality gives the displayed uniform bound.
\end{proof}

\begin{lemma}[Weighted fourth moment controls the nonlinear scale component]
\label{paper7:lem:scale-nonlinear-from-l4}
Assume \(V(\tau)\) is divergence-free and belongs to the approximation class
needed to justify the weighted integration by parts.  If
\[
  \sup_{\tau\ge\tau_0}
  \int_{\mathbb R^3}|y|^{-p-1}|V(y,\tau)|^4\,dy<\infty,
\]
then
\[
  |I_{\rm nl}(\tau)|
  \le
  \frac{p}{3}\int_{\mathbb R^3}|y|^{-p-1}|V(y,\tau)|^4\,dy.
\]
\end{lemma}

\begin{proof}
For smooth compactly supported approximants,
\[
  |V|V\cdot((V\cdot\nabla)V)=\frac13 V\cdot\nabla(|V|^3).
\]
Since \(\nabla\cdot V=0\), integration by parts gives
\[
  I_{\rm nl}
  =\frac13\int |y|^{-p}V\cdot\nabla(|V|^3)\,dy
  =-\frac13\int |V|^3V\cdot\nabla(|y|^{-p})\,dy.
\]
Because
\[
  \nabla(|y|^{-p})=-p|y|^{-p-2}y,
\]
we obtain
\[
  I_{\rm nl}
  =\frac{p}{3}\int |y|^{-p-2}(V\cdot y)|V|^3\,dy.
\]
Therefore
\[
  |I_{\rm nl}|
  \le
  \frac{p}{3}\int |y|^{-p-1}|V|^4\,dy.
\]
Taking the supremum in \(\tau\) proves the assertion.  The general admissible
case follows by the same approximation convention used for the weighted scale
normalization.
\end{proof}

\begin{corollary}[Pointwise scale-force bound]\label{paper7:cor:scale-force-bound}
If
\[
  \sup_\tau |I_\Delta(\tau)|<\infty,\qquad
  \sup_{\tau\ge\tau_0}\int |y|^{-p-1}|V(y,\tau)|^4\,dy<\infty,\qquad
  \sup_\tau |I_P(\tau)|<\infty,
\]
then
\[
  \sup_{\tau\ge\tau_0}
  \left|
  \int |y|^{-p}|V|V\cdot H(V,P)\,dy
  \right|<\infty.
\]
\end{corollary}

\begin{proof}
\Cref{paper7:lem:scale-nonlinear-from-l4} bounds \(I_{\rm nl}\).  Together
with the Laplacian and pressure bounds, this gives the hypotheses of
\cref{paper7:lem:scale-force-decomp}.
\end{proof}

\begin{lemma}[Bounded modulation from the modulation system]
\label{paper7:lem:bounded-mod-from-system}
Assume
\[
  M(V(\tau))
  \begin{pmatrix}
  a(\tau)\\ b(\tau)
  \end{pmatrix}
  =-F(V(\tau)),
\]
\[
  \sup_{\tau\ge\tau_0}\|M(V(\tau))^{-1}\|_{\infty\to\infty}<\infty,
  \qquad
  \sup_{\tau\ge\tau_0}\|F(V(\tau))\|_\infty<\infty.
\]
Then
\[
  \sup_{\tau\ge\tau_0}(|a(\tau)|+\|b(\tau)\|_\infty)<\infty.
\]
\end{lemma}

\begin{proof}
The system gives
\[
  \begin{pmatrix}
  a(\tau)\\ b(\tau)
  \end{pmatrix}
  =-M(V(\tau))^{-1}F(V(\tau)),
\]
so
\[
  \max\{|a(\tau)|,\|b(\tau)\|_\infty\}
  \le
  \|M(V(\tau))^{-1}\|_{\infty\to\infty}\|F(V(\tau))\|_\infty.
\]
Taking the supremum in \(\tau\) gives bounded modulation.
\end{proof}

\begin{proposition}[Windowed \(H^1\) exclusion from modulation bounds]\label{paper7:prop:windowed-h1-from-modulation}
Assume a renormalized Type II solution has the repaired-gauge orbit and pressure
reconstruction, bounded critical \(L^3\) norm, uniformly invertible modulation
matrix, uniformly bounded modulation forcing vector, and the compact-cylinder
regularity needed for Caccioppoli.  Then
\[
  \forall R>0:\qquad
  \sup_{T\ge\tau_0+1}
  \int_T^{T+1}\|V(\tau)\|_{H^1(B_R)}^2\,d\tau<\infty,
\]
and failure of local windowed \(H^1\) control is excluded for this solution.
\end{proposition}

\begin{proof}
The inverse matrix and forcing bounds imply bounded modulation by
\cref{paper7:lem:bounded-mod-from-system}.  The renormalized solution supplies the
renormalized equation and pressure reconstruction.  Bounded critical norm
gives the local compact-cylinder critical bounds, and the regularity hypothesis
allows the Caccioppoli estimate.  Therefore
\cref{paper7:cor:windowed-h1} applies.  Failure of local windowed \(H^1\)
control is the negation of the displayed tail windowed \(H^1\) estimate on
some bounded ball.
\end{proof}

\begin{corollary}[Two-alternative exclusion after windowed control]\label{paper7:cor:windowed-control-update}
Assume a finite-cost renormalized Type II solution not excluded by the compact
cost-divergence criterion has reached the two-alternative classification.  If
the hypotheses of \cref{paper7:prop:windowed-h1-from-modulation} are supplied for that
solution, then failure of local windowed \(H^1\) control is excluded.  If, in
addition, uniform critical tightness is available, then no finite-cost
alternative remains outside the compact cost-divergence criterion.
\end{corollary}

\begin{proof}
\Cref{paper7:prop:windowed-h1-from-modulation} gives the positive tail windowed
\(H^1\) bound, contradicting failure of local windowed \(H^1\) control.  If
uniform critical tightness is also available, then failure of tightness is
excluded as well,
and the two-alternative classification has no remaining finite-cost alternative
outside the compact cost-divergence criterion.
\end{proof}

\begin{theorem}[Exclusion of local windowed \(H^1\) failure]\label{paper7:thm:windowed-h1-failure-exclusion}
If every renormalized Type II solution in the class satisfies the hypotheses of
\cref{paper7:prop:windowed-h1-from-modulation}, then no such solution can satisfy
\[
  \exists m\ge1:
  \sup_{T\ge\tau_0+1}
  \int_T^{T+1}\|V(\tau)\|_{H^1(B_m)}^2\,d\tau=\infty.
\]
\end{theorem}

\begin{proof}
Apply \cref{paper7:prop:windowed-h1-from-modulation} to each renormalized solution.  The
displayed condition is the negation of the resulting windowed
\(H^1\)-bound for \(R=m\).
\end{proof}

\subsection{Window-Integrated Modulation Forcing and Weighted Scale Components}

This subsection records sufficient analytic estimates for the modulation
forcing.  The estimates are stated directly for the differentiated modulation
conditions and for the weighted scale component; no pointwise modulation bound
is used unless it is explicitly assumed.

Let the modulation forcing operator be
\[
  H(V,P):=\nu\Delta V-(V\cdot\nabla)V-\nabla P.
\]
For gauge functionals \(G_k\), \(0\le k\le3\), define
\[
  F_k(V):=DG_k(V)[H(V,P)].
\]
Thus the modulation system has the form
\[
  M(V(\tau))
  \begin{pmatrix}
  a(\tau)\\ b(\tau)
  \end{pmatrix}
  =
  -F(V(\tau)).
\]

\begin{lemma}[Forcing bound for the translation constraints]
\label{paper7:lem:translation-constraint-forcing}
Let
\[
  G_j(V)=\int_{\mathbb R^3}y_j|V(y)|^2\psi_R(y)\,dy,
  \qquad 1\le j\le3,
\]
where \(\psi_R\in C_c^\infty(B_{2R})\), and set
\[
  \Phi_j(y):=y_j\psi_R(y).
\]
For
\[
  F_j(V):=DG_j(V)[H(V,P)]
  =
  2\int_{\mathbb R^3}\Phi_j\,V\cdot H(V,P)\,dy,
\]
assume, uniformly for \(\tau\ge\tau_0\),
\[
  \|V(\tau)\|_{L^3(B_{2R})}\le M_3,\qquad
  \|\nabla V(\tau)\|_{L^2(B_{2R})}\le M_\nabla,
\]
and
\[
  \inf_{c\in\mathbb R}
  \|P(\tau)-c\|_{L^{3/2}(B_{2R})}\le M_P.
\]
Then
\[
  \sup_{\tau\ge\tau_0}\max_{1\le j\le3}|F_j(V(\tau))|
  \le
  C(R,\nu,\psi_R)
  \left(M_\nabla^2+M_3^2+M_3^3+M_PM_3\right).
\]
\end{lemma}

\begin{proof}
Fix \(j\).  Since \(\Phi_j\) is supported in \(B_{2R}\),
\[
  F_j=2\nu I_\Delta-2I_{\rm nl}-2I_P,
\]
where
\[
  I_\Delta:=\int \Phi_jV\cdot\Delta V,\qquad
  I_{\rm nl}:=\int \Phi_jV\cdot(V\cdot\nabla V),
  \qquad
  I_P:=\int \Phi_jV\cdot\nabla P.
\]
For the Laplacian term, using compact support of \(\Phi_j\) and summing over
the component index \(\alpha\),
\[
\begin{aligned}
  I_\Delta
  &=\sum_{\alpha=1}^3\int \Phi_j V_\alpha\Delta V_\alpha\,dy        \\
  &=-\sum_{\alpha,k}\int \partial_k(\Phi_jV_\alpha)\partial_kV_\alpha\,dy\\
  &=-\int \Phi_j|\nabla V|^2\,dy
    -\sum_{\alpha,k}\int(\partial_k\Phi_j)V_\alpha\partial_kV_\alpha\,dy .
\end{aligned}
\]
Equivalently,
\[
  I_\Delta
  =
  -\int \Phi_j|\nabla V|^2
  -\int(\nabla\Phi_j\cdot\nabla V)\cdot V .
\]
Consequently, by the Cauchy--Schwarz inequality,
\[
  |I_\Delta|
  \le
  \|\Phi_j\|_{L^\infty}\|\nabla V\|_{L^2(B_{2R})}^2
  +
  \|\nabla\Phi_j\|_{L^\infty}
  \|V\|_{L^2(B_{2R})}\|\nabla V\|_{L^2(B_{2R})}.
\]
H\"older's inequality on the bounded set \(B_{2R}\) gives
\[
  \|V\|_{L^2(B_{2R})}
  \le |B_{2R}|^{1/6}\|V\|_{L^3(B_{2R})}
  \le C_RM_3.
\]
Therefore
\[
  |I_\Delta|
  \le C_R(M_\nabla^2+M_3M_\nabla)
  \le C_R(M_\nabla^2+M_3^2).
\]
The last inequality is Young's inequality, \(ab\le(a^2+b^2)/2\).

For the nonlinear term,
\[
  V\cdot(V\cdot\nabla V)
  =
  V_iV_\alpha\partial_iV_\alpha
  =
  \frac12V_i\partial_i(|V|^2)
  =
  \frac12V\cdot\nabla(|V|^2).
\]
Since \(\nabla\cdot V=0\) and \(\Phi_j\) is compactly supported,
\[
  I_{\rm nl}
  =
  \frac12\int \Phi_jV\cdot\nabla(|V|^2)
  =
  -\frac12\int |V|^2V\cdot\nabla\Phi_j.
\]
Thus
\[
  |I_{\rm nl}|
  \le
  \frac12\|\nabla\Phi_j\|_{L^\infty}
  \|V\|_{L^3(B_{2R})}^3
  \le C_RM_3^3.
\]

For the pressure term, replace \(P\) by \(P-c(\tau)\).  Since
\(\nabla\cdot V=0\), the identity
\[
  \nabla\cdot(\Phi_j V)=V\cdot\nabla\Phi_j
\]
holds distributionally.  Therefore
\[
  I_P
  =
  \int \Phi_jV\cdot\nabla(P-c)
  =
  -\int(P-c)V\cdot\nabla\Phi_j.
\]
Consequently
\[
  |I_P|
  \le
  \|\nabla\Phi_j\|_{L^\infty}
  \|P-c\|_{L^{3/2}(B_{2R})}
  \|V\|_{L^3(B_{2R})}
  \le C_RM_PM_3.
\]
Taking the infimum over \(c\), summing the three estimates in \(F_j\), and
then taking the supremum in \(\tau\) proves the bound.
\end{proof}

\begin{remark}[Local source of the gradient bound]
\label{paper7:rem:translation-gradient-input}
The pointwise local gradient bound in
\cref{paper7:lem:translation-constraint-forcing} is not a free terminal
assumption in the retained proof.  Whenever the pointwise lemma is invoked on a
retained compact window, the bound is supplied at admissible times by the local
Caccioppoli/regularity package, or else the first failed row is the local
\(H^1\)-loss or rough-core row before the translation estimate is used.  Thus the
translation forcing estimate never obtains windowed \(H^1\)-control by feeding
that same conclusion back into its hypotheses.
\end{remark}

\begin{definition}[Weighted scale forcing]
\label{paper7:def:weighted-scale-forcing}
Let
\[
  G_{\rm sc}(V)=\int_{\mathbb R^3}|y|^{-p}|V|^3\,dy-\Theta_0,
  \qquad 0<p<3.
\]
The weighted scale component is
\[
  F_0(V):=DG_{\rm sc}(V)[H(V,P)]
  =
  3\int_{\mathbb R^3}|y|^{-p}|V|V\cdot H(V,P)\,dy.
\]
All weighted pairings in this subsection are assumed to be defined either as
absolutely convergent integrals or as specified distributional pairings
represented by \(L^1_{\rm loc}(d\tau)\) functions.
\end{definition}

\begin{lemma}[Weighted scale forcing bound]
\label{paper7:lem:scale-component-weighted-forcing}
Assume
\[
  \sup_{\tau\ge\tau_0}
  \left|
  \int_{\mathbb R^3}|y|^{-p}|V|V\cdot H(V,P)\,dy
  \right|<\infty.
\]
Then
\[
  \sup_{\tau\ge\tau_0}|F_0(V(\tau))|<\infty.
\]
\end{lemma}

\begin{proof}
By the derivative formula in \cref{paper7:def:weighted-scale-forcing},
\[
  |F_0(V(\tau))|
  \le
  3\left|
  \int_{\mathbb R^3}|y|^{-p}|V|V\cdot H(V,P)\,dy
  \right|.
\]
Taking the supremum in \(\tau\) gives the result.
\end{proof}

\begin{corollary}[Pointwise modulation forcing bound]
\label{paper7:cor:pointwise-modulation-forcing}
Assume the hypotheses of \cref{paper7:lem:translation-constraint-forcing} for
the three translation components and the weighted scale forcing hypothesis of
\cref{paper7:lem:scale-component-weighted-forcing}.  Then
\[
  \sup_{\tau\ge\tau_0}\|F(V(\tau))\|_\infty<\infty .
\]
\end{corollary}

\begin{proof}
The vector \(F(V)\in\mathbb R^4\) has the scale component \(F_0\) and the
three translation components \(F_j\), \(1\le j\le3\).  The scale component is
bounded by \cref{paper7:lem:scale-component-weighted-forcing}, and the
translation components are bounded by
\cref{paper7:lem:translation-constraint-forcing}.  Taking the maximum over the
four components proves the \(\ell^\infty\)-bound.
\end{proof}

\begin{corollary}[Analytic inputs for pointwise modulation forcing]
\label{paper7:cor:pointwise-modulation-inputs}
Pointwise boundedness of the modulation forcing follows from the following
analytic estimates:
\begin{enumerate}[label=\textup{(\roman*)}]
\item local \(L^3\), local \(H^1\), and local pressure-oscillation control on
the support of the compactly supported centering constraints;
\item weighted control of the Laplacian, pressure, and fourth-moment terms in
the scale component, sufficient to bound the weighted scale forcing.
\end{enumerate}
The scale-component estimate is independent of the algebraic scale
transversality
identity
\[
  DG_{\rm sc}(V)[Z_{\rm sc}(V)]=p\Theta_0.
\]
\end{corollary}

\begin{proof}
The first item is the set of hypotheses used in
\cref{paper7:lem:translation-constraint-forcing}.  The second item bounds the
weighted scale component through \cref{paper7:lem:scale-component-weighted-forcing},
or through the explicit decomposition in
\cref{paper7:cor:reduced-window-scale-force} below in the corresponding
integrated setting.  \Cref{paper7:cor:pointwise-modulation-forcing} then gives
the modulation forcing bound.  The displayed transversality identity controls
the modulation matrix, not the forcing vector, so it is a separate algebraic
input.
\end{proof}

\begin{lemma}[Window-integrated forcing gives window-integrated modulation]
\label{paper7:lem:window-l1-forcing-modulation}
Assume
\[
  \sup_{\tau\ge\tau_0}
  \|M(V(\tau))^{-1}\|_{\infty\to\infty}\le C_M<\infty
\]
and
\[
  \sup_{T\ge\tau_0}
  \int_T^{T+1}\|F(V(\tau))\|_\infty\,d\tau
  \le C_F<\infty .
\]
Then
\[
  \sup_{T\ge\tau_0}
  \int_T^{T+1}\bigl(|a(\tau)|+\|b(\tau)\|_\infty\bigr)\,d\tau
  \le 2C_MC_F.
\]
\end{lemma}

\begin{proof}
For almost every \(\tau\),
\[
  \begin{pmatrix}
  a(\tau)\\ b(\tau)
  \end{pmatrix}
  =
  -M(V(\tau))^{-1}F(V(\tau)).
\]
Hence
\[
  \max\{|a(\tau)|,\|b(\tau)\|_\infty\}
  \le C_M\|F(V(\tau))\|_\infty.
\]
Since
\[
  |a(\tau)|+\|b(\tau)\|_\infty
  \le
  2\max\{|a(\tau)|,\|b(\tau)\|_\infty\},
\]
we have, for every \(T\ge\tau_0\),
\[
\begin{aligned}
  \int_T^{T+1}\bigl(|a(\tau)|+\|b(\tau)\|_\infty\bigr)\,d\tau
  &\le
  2C_M\int_T^{T+1}\|F(V(\tau))\|_\infty\,d\tau  \\
  &\le 2C_MC_F .
\end{aligned}
\]
Taking the supremum over \(T\) proves the claim.
\end{proof}

\begin{lemma}[Selection from an \(L^1\) modulation bound]
\label{paper7:lem:unit-window-selection}
Let \(I_n=[T_n,T_n+1]\), where \(T_n\to\infty\).  Let
\[
  d_n:I_n\to[0,\infty),\qquad q_n:I_n\to[0,\infty)
\]
be measurable functions satisfying
\[
  \int_{I_n}d_n(\tau)\,d\tau\to0,
  \qquad
  \sup_n\int_{I_n}q_n(\tau)\,d\tau\le Q<\infty .
\]
Then, after choosing measurable representatives, there are times
\(\tau_n\in I_n\) such that
\[
  d_n(\tau_n)\to0
\]
and, if \(Q>0\),
\[
  q_n(\tau_n)\le2Q
  \qquad\text{for every }n.
\]
If \(Q=0\), the times may be chosen so that
\[
  q_n(\tau_n)=0,\qquad d_n(\tau_n)\to0.
\]
\end{lemma}

\begin{proof}
Assume first that \(Q>0\), and set
\[
  E_n:=\{\tau\in I_n:q_n(\tau)\le2Q\}.
\]
By Chebyshev's inequality,
\[
  |I_n\setminus E_n|
  \le
  \frac{1}{2Q}\int_{I_n}q_n(\tau)\,d\tau
  \le \frac12.
\]
Thus \(|E_n|\ge1/2\).  Since \(d_n\ge0\),
\[
  \inf_{\tau\in E_n}d_n(\tau)
  \le
  \frac{1}{|E_n|}\int_{E_n}d_n(\tau)\,d\tau
  \le
  2\int_{I_n}d_n(\tau)\,d\tau.
\]
Choose \(\tau_n\in E_n\), outside the null set on which the representatives
are undefined, such that
\[
  d_n(\tau_n)
  \le
  2\int_{I_n}d_n(\tau)\,d\tau+\frac1n.
\]
The right-hand side tends to zero, hence \(d_n(\tau_n)\to0\), and the
definition of \(E_n\) gives \(q_n(\tau_n)\le2Q\).

If \(Q=0\), then \(q_n=0\) almost everywhere on \(I_n\).  Applying the same
averaging argument to \(d_n\) on the full-measure set where \(q_n=0\) gives
the asserted times.
\end{proof}

\begin{corollary}[Selection on unit windows]
\label{paper7:cor:unit-window-selection}
Assume
\[
  q(\tau):=|a(\tau)|+\|b(\tau)\|_\infty
\]
satisfies
\[
  \sup_{T\ge\tau_0}\int_T^{T+1}q(\tau)\,d\tau<\infty .
\]
For every sequence of unit windows \(I_n=[T_n,T_n+1]\) and every nonnegative
measurable density \(d_n\) satisfying
\[
  \int_{I_n}d_n(\tau)\,d\tau\to0,
\]
there are times \(\tau_n\in I_n\) such that
\[
  d_n(\tau_n)\to0,\qquad
  \sup_n\bigl(|a(\tau_n)|+\|b(\tau_n)\|_\infty\bigr)<\infty .
\]
\end{corollary}

\begin{proof}
Apply \cref{paper7:lem:unit-window-selection} with \(q_n=q|_{I_n}\) and
\[
  Q=\sup_n\int_{I_n}q(\tau)\,d\tau .
\]
This \(Q\) is finite because of the uniform window \(L^1\)-bound.  If
\(Q>0\), the selected times satisfy
\[
  |a(\tau_n)|+\|b(\tau_n)\|_\infty=q(\tau_n)\le2Q.
\]
If \(Q=0\), the same lemma gives \(q(\tau_n)=0\).  In both cases
\(d_n(\tau_n)\to0\) and the modulation is uniformly bounded along the
selected times.
\end{proof}

\begin{lemma}[Product-form windowed modulation estimate]
\label{paper7:lem:product-form-window-modulation}
Assume
\[
  \sup_{T\ge\tau_0}
  \int_T^{T+1}
  \|M(V(\tau))^{-1}\|_{\infty\to\infty}
  \|F(V(\tau))\|_\infty\,d\tau<\infty.
\]
Then
\[
  \sup_{T\ge\tau_0}
  \int_T^{T+1}\bigl(|a(\tau)|+\|b(\tau)\|_\infty\bigr)\,d\tau<\infty .
\]
\end{lemma}

\begin{proof}
The modulation system gives
\[
  \begin{pmatrix}a(\tau)\\ b(\tau)\end{pmatrix}
  =
  -M(V(\tau))^{-1}F(V(\tau)).
\]
Therefore
\[
  |a(\tau)|+\|b(\tau)\|_\infty
  \le
  2\|M(V(\tau))^{-1}\|_{\infty\to\infty}
  \|F(V(\tau))\|_\infty.
\]
Integrating this pointwise estimate on \([T,T+1]\) and taking the supremum in
\(T\) gives the result.
\end{proof}

\begin{lemma}[Integrated translation forcing bound]
\label{paper7:lem:window-translation-forcing}
Let \(G_j\), \(\Phi_j\), \(H(V,P)\), and \(F_j(V)\) be as in
\cref{paper7:lem:translation-constraint-forcing}.  Assume that, for some
\(R>0\),
\[
  \sup_{\tau\ge\tau_0}\|V(\tau)\|_{L^3(B_{2R})}\le M_3,
\]
\[
  \sup_{T\ge\tau_0}
  \int_T^{T+1}\|\nabla V(\tau)\|_{L^2(B_{2R})}^2\,d\tau
  \le A_R,
\]
and
\[
  \sup_{\tau\ge\tau_0}\inf_{c\in\mathbb R}
  \|P(\tau)-c\|_{L^{3/2}(B_{2R})}\le M_P.
\]
Then
\[
  \sup_{T\ge\tau_0}
  \int_T^{T+1}\max_{1\le j\le3}|F_j(V(\tau))|\,d\tau
  \le
  C(R,\nu,\psi_R)(A_R+M_3^2+M_3^3+M_PM_3).
\]
\end{lemma}

\begin{proof}
Fix \(j\).  The proof of
\cref{paper7:lem:translation-constraint-forcing} gives, for almost every
\(\tau\),
\[
  |F_j(V(\tau))|
  \le
  C(R,\nu,\psi_R)
  \left(
  \|\nabla V(\tau)\|_{L^2(B_{2R})}^2
  +M_3\|\nabla V(\tau)\|_{L^2(B_{2R})}
  +M_3^3+M_PM_3
  \right).
\]
Integrate this estimate over \(I_T=[T,T+1]\).  The first term contributes at
most \(A_R\).  For the mixed term, Cauchy's inequality on the unit interval
gives
\[
  \int_{I_T}M_3\|\nabla V(\tau)\|_{L^2(B_{2R})}\,d\tau
  \le
  M_3
  \left(
  \int_{I_T}\|\nabla V(\tau)\|_{L^2(B_{2R})}^2\,d\tau
  \right)^{1/2}
  \le M_3A_R^{1/2}.
\]
Since
\[
  M_3A_R^{1/2}\le\frac12(M_3^2+A_R),
\]
the mixed term is bounded by \(C(A_R+M_3^2)\).  The terms \(M_3^3\) and
\(M_PM_3\) are constant on the unit window.  Hence
\[
  \int_T^{T+1}|F_j(V(\tau))|\,d\tau
  \le
  C(R,\nu,\psi_R)(A_R+M_3^2+M_3^3+M_PM_3).
\]
Since \(\max_{1\le j\le3}|F_j|\le\sum_{j=1}^3|F_j|\), taking the maximum
over the three centering constraints changes only the constant.
\end{proof}

\begin{remark}[Independence of the windowed gradient hypothesis]
\label{paper7:rem:window-translation-input}
The windowed gradient input in
\cref{paper7:lem:window-translation-forcing} is an independent
hypothesis whenever the integrated translation estimate is used to prove
windowed \(H^1\)-control.  It may come, for example, from a preliminary
Caccioppoli estimate depending only on already available data, from a direct
approximation or regularity estimate on the support of the centering
functions, or from an independent verification obtained without using the
conclusion of the windowed \(H^1\) estimate under consideration.
\end{remark}

\begin{lemma}[Window-integrated scale forcing]
\label{paper7:lem:window-scale-forcing}
Assume the window-integrated weighted scale-forcing bound
\[
  \sup_{T\ge\tau_0}
  \int_T^{T+1}
  \left|
  \int_{\mathbb R^3}|y|^{-p}|V|V\cdot H(V,P)\,dy
  \right|\,d\tau
  <\infty,
\]
where the weighted pairing is interpreted as in
\cref{paper7:def:weighted-scale-forcing}.  Then
\[
  \sup_{T\ge\tau_0}\int_T^{T+1}|F_0(V(\tau))|\,d\tau<\infty .
\]
\end{lemma}

\begin{proof}
By the identity
\[
  F_0(V)
  =
  3\int_{\mathbb R^3}|y|^{-p}|V|V\cdot H(V,P)\,dy
\]
we have, for each \(T\ge\tau_0\),
\[
  \int_T^{T+1}|F_0(V(\tau))|\,d\tau
  =
  3\int_T^{T+1}
  \left|
  \int_{\mathbb R^3}|y|^{-p}|V|V\cdot H(V,P)\,dy
  \right|\,d\tau .
\]
Taking the supremum over \(T\) gives the asserted \(L^1\)-bound.
\end{proof}

\begin{lemma}[Full integrated modulation forcing]
\label{paper7:lem:window-modulation-forcing}
Assume the integrated translation forcing estimate
\[
  \sup_{T\ge\tau_0}
  \int_T^{T+1}\max_{1\le j\le3}|F_j(V(\tau))|\,d\tau<\infty
\]
and the window-integrated weighted scale-forcing bound in
\cref{paper7:lem:window-scale-forcing}.  Then
\[
  \sup_{T\ge\tau_0}
  \int_T^{T+1}\|F(V(\tau))\|_\infty\,d\tau<\infty .
\]
\end{lemma}

\begin{proof}
The vector \(F(V)\in\mathbb R^4\) consists of \(F_0\) and the three
translation components \(F_j\), \(1\le j\le3\).  Since
\[
  \|F(V(\tau))\|_\infty
  \le
  |F_0(V(\tau))|+\max_{1\le j\le3}|F_j(V(\tau))|,
\]
for almost every \(\tau\).  Integrating this inequality on \([T,T+1]\) gives
\[
  \int_T^{T+1}\|F(V(\tau))\|_\infty\,d\tau
  \le
  \int_T^{T+1}|F_0(V(\tau))|\,d\tau
  +
  \int_T^{T+1}\max_{1\le j\le3}|F_j(V(\tau))|\,d\tau .
\]
The first term is uniformly bounded in \(T\) by
\cref{paper7:lem:window-scale-forcing}, and the second term is uniformly
bounded by the assumed integrated translation forcing estimate.  Taking the
supremum over \(T\) proves the claim.
\end{proof}

\begin{corollary}[Integrated modulation criterion on selected windows]
\label{paper7:cor:selected-window-modulation}
Assume the uniform inverse bound for \(M(V(\tau))\), the integrated
translation forcing estimate, and the window-integrated weighted scale-forcing
bound.  Then
\[
  \sup_{T\ge\tau_0}
  \int_T^{T+1}\bigl(|a(\tau)|+\|b(\tau)\|_\infty\bigr)\,d\tau<\infty.
\]
Consequently every sequence of unit windows with vanishing nonnegative
measurable density contains selected times where this density vanishes
asymptotically and the modulation remains uniformly bounded.
\end{corollary}

\begin{proof}
\Cref{paper7:lem:window-modulation-forcing} gives the
window-integrated bound on \(\|F(V)\|_\infty\).  Combining this with the
uniform inverse bound in
\cref{paper7:lem:window-l1-forcing-modulation} gives the
window-integrated modulation bound.  The selection conclusion is
\cref{paper7:cor:unit-window-selection}.
\end{proof}

\begin{definition}[Weighted scale-force components]
\label{paper7:def:scale-force-components}
Let \(w(y)=|y|^{-p}\), \(0<p<3\).  Define
\[
  F_0^\Delta(V)
  :=
  3\nu\int_{\mathbb R^3}w|V|V\cdot\Delta V\,dy,
\]
\[
  F_0^{\rm nl}(V)
  :=
  -3\int_{\mathbb R^3}w|V|V\cdot(V\cdot\nabla V)\,dy,
\]
and
\[
  F_0^P(V,P)
  :=
  -3\int_{\mathbb R^3}w|V|V\cdot\nabla P\,dy.
\]
The three pairings are required to be compatible with the decomposition of
\(H(V,P)\) into diffusion, transport, and pressure.
\end{definition}

\begin{lemma}[Integrated scale-force decomposition]
\label{paper7:lem:scale-force-decomposition-window}
For every time at which the three pairings in
\cref{paper7:def:scale-force-components} are defined,
\[
  F_0(V(\tau))
  =
  F_0^\Delta(V(\tau))
  +F_0^{\rm nl}(V(\tau))
  +F_0^P(V(\tau),P(\tau)).
\]
\end{lemma}

\begin{proof}
Substitute
\[
  H(V,P)=\nu\Delta V-(V\cdot\nabla)V-\nabla P
\]
into
\[
  F_0(V)=3\int w|V|V\cdot H(V,P)\,dy.
\]
Then
\[
\begin{aligned}
  F_0(V)
  &=
  3\nu\int w|V|V\cdot\Delta V\,dy
  -3\int w|V|V\cdot(V\cdot\nabla V)\,dy  \\
  &\qquad
  -3\int w|V|V\cdot\nabla P\,dy .
\end{aligned}
\]
The three terms coincide with \(F_0^\Delta(V)\), \(F_0^{\rm nl}(V)\), and
\(F_0^P(V,P)\), respectively.  Thus linearity of the weighted pairing gives
the stated decomposition.
\end{proof}

\begin{lemma}[Component bounds imply integrated scale-force control]
\label{paper7:lem:scale-components-control}
Assume
\[
  \sup_{T\ge\tau_0}
  \int_T^{T+1}|F_0^\Delta(V(\tau))|\,d\tau<\infty,
\]
\[
  \sup_{T\ge\tau_0}
  \int_T^{T+1}|F_0^{\rm nl}(V(\tau))|\,d\tau<\infty,
\]
and
\[
  \sup_{T\ge\tau_0}
  \int_T^{T+1}|F_0^P(V(\tau),P(\tau))|\,d\tau<\infty.
\]
Then the window-integrated weighted scale-forcing bound in
\cref{paper7:lem:window-scale-forcing} holds.
\end{lemma}

\begin{proof}
By \cref{paper7:lem:scale-force-decomposition-window},
\[
  |F_0(V(\tau))|
  \le
  |F_0^\Delta(V(\tau))|
  +|F_0^{\rm nl}(V(\tau))|
  +|F_0^P(V(\tau),P(\tau))|
\]
for almost every \(\tau\).  Integrating on \([T,T+1]\), taking the supremum
in \(T\), and applying the three component bounds gives
\[
  \sup_{T\ge\tau_0}\int_T^{T+1}|F_0(V(\tau))|\,d\tau<\infty.
\]
Since \(F_0\) is three times the weighted pairing in
\cref{paper7:def:weighted-scale-forcing}, this is equivalent to the
window-integrated weighted scale-forcing bound.
\end{proof}

\begin{definition}[Singular-weight convective integration by parts]
\label{paper7:def:singular-weight-convective-ibp}
The singular-weight convective integration-by-parts condition means that, for
almost every \(\tau\),
\[
  -\int_{\mathbb R^3}w\,V\cdot\nabla(|V|^3)\,dy
  =
  \int_{\mathbb R^3}|V|^3V\cdot\nabla w\,dy,
\]
with no boundary contribution from \(y=0\) or from spatial infinity, either
classically or by convergence of smooth compactly supported divergence-free
approximants.
\end{definition}

\begin{lemma}[Convective scale-force bound]
\label{paper7:lem:convective-scale-force}
Assume the condition of \cref{paper7:def:singular-weight-convective-ibp} and
\[
  \sup_{T\ge\tau_0}
  \int_T^{T+1}\int_{\mathbb R^3}
  |y|^{-p-1}|V(y,\tau)|^4\,dy\,d\tau<\infty .
\]
Then
\[
  \sup_{T\ge\tau_0}
  \int_T^{T+1}|F_0^{\rm nl}(V(\tau))|\,d\tau<\infty.
\]
More precisely,
\[
  |F_0^{\rm nl}(V(\tau))|
  \le
  p\int_{\mathbb R^3}|y|^{-p-1}|V(y,\tau)|^4\,dy
\]
for almost every \(\tau\).
\end{lemma}

\begin{proof}
Fix such a time \(\tau\) and suppress it from the notation.  Since
\[
  \partial_i(|V|^3)=3|V|V_j\partial_iV_j
\]
for each \(i\), and therefore
\[
  V\cdot\nabla(|V|^3)
  =
  3|V|V_iV_j\partial_iV_j
  =
  3|V|\,V\cdot((V\cdot\nabla)V),
\]
we have
\[
  w|V|V\cdot(V\cdot\nabla V)
  =
  \frac13 w\,V\cdot\nabla(|V|^3).
\]
Therefore
\[
  F_0^{\rm nl}(V)
  =
  -3\int w|V|V\cdot(V\cdot\nabla V)
  =
  -\int w\,V\cdot\nabla(|V|^3).
\]
Using \(\nabla\cdot V=0\) and
\cref{paper7:def:singular-weight-convective-ibp},
\[
  -\int w\,V\cdot\nabla(|V|^3)
  =
  \int |V|^3V\cdot\nabla w.
\]
Since
\[
  \nabla w(y)=-p|y|^{-p-2}y,\qquad y\ne0,
\]
we have
\[
  |V\cdot\nabla w|
  =
  p|y|^{-p-2}|V\cdot y|
  \le p|y|^{-p-1}|V|.
\]
Thus
\[
  |F_0^{\rm nl}(V)|
  \le
  p\int |y|^{-p-1}|V|^4\,dy.
\]
Integrating over \([T,T+1]\) and taking the supremum in \(T\) proves the
window-integrated convective bound.
\end{proof}

\begin{definition}[Annular convective regularity]
\label{paper7:def:annular-convective-regularity}
Annular convective regularity means that, for almost every \(\tau\),
\(V(\tau)\) is divergence free and, on every compact annulus
\(0<r<|y|<R<\infty\), the chain-rule identities
\[
  \partial_i(|V|^3)=3|V|V_j\partial_iV_j,
  \qquad
  V\cdot\nabla(|V|^3)=3|V|V_iV_j\partial_iV_j
\]
hold in the sense of distributions, either classically or by passage to the
limit from smooth divergence-free approximants on annuli.
\end{definition}

\begin{lemma}[Annular regularity gives singular-weight integration by parts]
\label{paper7:lem:annular-regularity-ibp}
Assume annular convective regularity and the weighted \(L^4\)-window bound
\[
  \sup_{T\ge\tau_0}
  \int_T^{T+1}\int_{\mathbb R^3}
  |y|^{-p-1}|V(y,\tau)|^4\,dy\,d\tau<\infty.
\]
Then the singular-weight convective integration-by-parts condition of
\cref{paper7:def:singular-weight-convective-ibp} holds for almost every \(\tau\).
\end{lemma}

\begin{proof}
Fix a time \(\tau\) for which
\[
  \int_{\mathbb R^3}|y|^{-p-1}|V(y,\tau)|^4\,dy<\infty
\]
and annular convective regularity holds.  Such times form a full-measure set
on every unit window.

Let \(\eta\in C^\infty([0,\infty))\) satisfy \(0\le\eta\le1\),
\(\eta(s)=0\) for \(s\le1\), and \(\eta(s)=1\) for \(s\ge2\).  Define
\[
  \chi_{\varepsilon,R}(y)
  :=
  \eta(|y|/\varepsilon)\eta(R/|y|).
\]
Then \(\chi_{\varepsilon,R}\) is supported in
\(\{\varepsilon<|y|<R\}\), equals one on
\(\{2\varepsilon<|y|<R/2\}\), and satisfies
\[
  |\nabla\chi_{\varepsilon,R}(y)|
  \le
  C\varepsilon^{-1}\mathbf 1_{\{\varepsilon<|y|<2\varepsilon\}}
  +CR^{-1}\mathbf 1_{\{R/2<|y|<R\}}.
\]
By annular convective regularity, the following integration by parts is valid
on the annulus \(\{\varepsilon<|y|<R\}\).  Since
\(\nabla\cdot V=0\) and \(w\chi_{\varepsilon,R}|V|^3V\) is compactly
supported in this annulus,
\[
  0
  =
  \int \nabla\cdot\bigl(w\chi_{\varepsilon,R}|V|^3V\bigr)\,dy .
\]
Expanding the divergence gives
\[
  0
  =
  \int |V|^3V\cdot\nabla(w\chi_{\varepsilon,R})\,dy
  +
  \int w\chi_{\varepsilon,R}V\cdot\nabla(|V|^3)\,dy ,
\]
and hence
\[
  -\int w\chi_{\varepsilon,R}V\cdot\nabla(|V|^3)
  =
  \int |V|^3V\cdot\nabla(w\chi_{\varepsilon,R}).
\]
Expanding the derivative,
\[
  -\int w\chi_{\varepsilon,R}V\cdot\nabla(|V|^3)
  =
  \int \chi_{\varepsilon,R}|V|^3V\cdot\nabla w
  +\int w|V|^3V\cdot\nabla\chi_{\varepsilon,R}.
\]
On the inner cutoff region \(\{\varepsilon<|y|<2\varepsilon\}\), one has
\(|y|\simeq\varepsilon\), and on the outer cutoff region
\(\{R/2<|y|<R\}\), one has \(|y|\simeq R\).  Therefore the derivative bound
for \(\chi_{\varepsilon,R}\) implies
\[
\begin{aligned}
  \left|
  \int w|V|^3V\cdot\nabla\chi_{\varepsilon,R}
  \right|
  &\le
  C\int_{\{\varepsilon<|y|<2\varepsilon\}}
  |y|^{-p-1}|V|^4\,dy \\
  &\quad{}
  +C\int_{\{R/2<|y|<R\}}
  |y|^{-p-1}|V|^4\,dy .
\end{aligned}
\]
The integrand \(|y|^{-p-1}|V|^4\) belongs to \(L^1(\mathbb R^3)\) at the
fixed time.  Hence absolute continuity of the Lebesgue integral gives
\[
  \int_{\{\varepsilon<|y|<2\varepsilon\}}
  |y|^{-p-1}|V|^4\,dy\to0
  \qquad\text{as }\varepsilon\downarrow0,
\]
and integrability at infinity gives
\[
  \int_{\{R/2<|y|<R\}}
  |y|^{-p-1}|V|^4\,dy\to0
  \qquad\text{as }R\uparrow\infty.
\]
Moreover
\[
  |\chi_{\varepsilon,R}|V|^3V\cdot\nabla w|
  \le
  p|y|^{-p-1}|V|^4\in L^1(\mathbb R^3),
\]
so dominated convergence gives
\[
  \int \chi_{\varepsilon,R}|V|^3V\cdot\nabla w
  \to
  \int |V|^3V\cdot\nabla w.
\]
Passing first \(\varepsilon\downarrow0\) and then \(R\uparrow\infty\) yields
\[
  -\int_{\mathbb R^3}w\,V\cdot\nabla(|V|^3)\,dy
  =
  \int_{\mathbb R^3}|V|^3V\cdot\nabla w\,dy,
\]
with no boundary contribution at \(y=0\) or at infinity.
\end{proof}

\begin{corollary}[Reduced integrated scale-force criterion]
\label{paper7:cor:reduced-window-scale-force}
Assume annular convective regularity, the weighted \(L^4\)-window bound,
and the two component estimates
\[
  \sup_{T\ge\tau_0}
  \int_T^{T+1}|F_0^\Delta(V(\tau))|\,d\tau<\infty,
\]
\[
  \sup_{T\ge\tau_0}
  \int_T^{T+1}|F_0^P(V(\tau),P(\tau))|\,d\tau<\infty.
\]
Then the window-integrated weighted scale-forcing bound holds.
\end{corollary}

\begin{proof}
\Cref{paper7:lem:annular-regularity-ibp} gives the singular-weight
integration-by-parts condition.  \Cref{paper7:lem:convective-scale-force}
then gives the integrated convective scale-force bound.  Combining this with
the diffusive and pressure component estimates, and applying
\cref{paper7:lem:scale-components-control}, gives the window-integrated
weighted scale-forcing bound.
\end{proof}

\begin{corollary}[Remaining analytic alternatives for the scale component]
\label{paper7:cor:scale-component-alternatives}
Before the singular-weight integration by parts is justified, failure of the
window-integrated weighted scale-forcing bound is reduced to failure of at
least one of the following four analytic statements:
\begin{enumerate}[label=\textup{(\roman*)}]
\item the integrated diffusive weighted scale-force bound;
\item the singular-weight convective integration-by-parts identity;
\item the weighted \(L^4\)-window bound;
\item the integrated weighted pressure scale-force bound.
\end{enumerate}
If annular convective regularity is available, then the second item is replaced
by failure of annular convective regularity itself, together with the same
weighted \(L^4\)-window input.
\end{corollary}

\begin{proof}
\Cref{paper7:lem:scale-components-control} shows that the diffusive,
convective, and pressure component bounds imply the integrated scale-force
bound.  \Cref{paper7:lem:convective-scale-force} reduces the
convective component to the singular-weight integration-by-parts identity and
the weighted \(L^4\)-window bound.  Finally,
\cref{paper7:lem:annular-regularity-ibp} shows that annular convective
regularity and the weighted \(L^4\)-window bound imply the singular-weight
integration-by-parts identity.  Therefore the listed failures exhaust the
remaining ways in which this scale-component estimate can fail.
\end{proof}

\subsection{Navier--Stokes Cost-Divergence Exclusion}

\begin{definition}[Type II exclusion conclusion]\label{paper7:def:cost-exclusion-conclusion}
For a renormalized Type II solution, the exclusion conclusion means that after
the active supercritical scaling condition and the compact cost-divergence
criterion have both been established, the solution is excluded from the Type II
class under consideration.  This is stronger than cost divergence alone, which
only states that the solution satisfies the compact Type II cost-divergence
criterion.
\end{definition}

\begin{definition}[Local cost-to-exclusion data]
\label{paper7:def:local-cost-to-exclusion-data}
Let \(\omega\) be a repaired-gauge Navier--Stokes Type II branch on a terminal
renormalized tail \([\tau_0,\infty)\).  Local cost-to-exclusion data for the
validated compact cost consist of a nonnegative locally integrable cost
\(\mathcal D_\omega\), a locally absolutely continuous corrected localized
energy \(\mathcal H_\omega\), a constant \(c_\omega>0\), and a nonnegative
locally integrable remainder \(\mathcal R_\omega\) such that:
\begin{enumerate}[label=\textup{(\roman*)}]
\item the validated compact cost-divergence conclusion is
\[
  \int_{\tau_0}^{\infty}\mathcal D_\omega(\tau)\,d\tau=\infty;
\]
\item \(\mathcal H_\omega(\tau_0)<\infty\) and
\[
  \inf_{\tau\ge\tau_0}\mathcal H_\omega(\tau)>-\infty;
\]
\item the corrected localized energy inequality
\[
  \frac{d}{d\tau}\mathcal H_\omega(\tau)
  +c_\omega\mathcal D_\omega(\tau)
  +\mathcal R_\omega(\tau)\le0
\]
holds in distributions on \((\tau_0,\infty)\).
\end{enumerate}
For the identity cost one takes
\(\mathcal D_\omega=\tilde{\mathfrak D}_{R_0}\).  For an explicitly comparable
localized cost, \(\mathcal D_\omega\) is the localized cost whose divergence is
validated by the compact criterion, provided that the same cost is controlled
by the corrected localized energy inequality on the terminal tail.
\end{definition}

\begin{definition}[Energy-controlled validated cost]
\label{paper7:def:energy-controlled-validated-cost}
Let \((V,P,a,b)\) be a repaired-gauge branch and let
\[
  \tilde{\mathfrak D}_{R_0}(\tau)
  =
  \nu\int |\nabla V|^2\phi_{R_0}
  +a_+(\tau)\int |V|^2\phi_{R_0}
\]
be the identity localized Navier--Stokes cost.  A validated compact cost
\(\mathfrak C_{\rm II}^{NS}\) is energy-controlled by the identity cost on a
terminal tail if there are constants \(C_1<\infty\) and
\(\tau_1\ge\tau_0\) such that
\[
  0\le \mathfrak C_{\rm II}^{NS}(\tau)
  \le C_1\tilde{\mathfrak D}_{R_0}(\tau)
  \qquad\text{for a.e. }\tau\ge\tau_1 .
\]
The identity cost is energy-controlled with \(C_1=1\) and
\(\tau_1=\tau_0\).  For a non-identity cost, this is the upper comparison
needed for a corrected localized energy inequality controlling
\(\tilde{\mathfrak D}_{R_0}\) to also control the validated cost.
\end{definition}

\begin{definition}[Local monotonicity carrier for a validated cost]
\label{paper7:def:local-monotonicity-carrier}
Let \(\omega\) be a repaired-gauge Navier--Stokes Type II branch on a terminal
renormalized tail \([\tau_0,\infty)\), and let
\(\mathfrak C_{\rm II}^{NS}\) be the validated compact cost for that branch.
A local monotonicity carrier for \(\mathfrak C_{\rm II}^{NS}\) consists of a
nonnegative locally integrable carrier cost \(\mathcal A_\omega\), a locally
absolutely continuous corrected localized energy \(\mathcal H_\omega\), a
nonnegative locally integrable remainder \(\mathcal R_\omega\), and constants
\(c_A>0\), \(0<C_A<\infty\) such that:
\begin{enumerate}[label=\textup{(\roman*)}]
\item \(\mathcal H_\omega(\tau_0)<\infty\) and
\[
  \inf_{\tau\ge\tau_0}\mathcal H_\omega(\tau)>-\infty;
\]
\item the carrier monotonicity inequality
\[
  \frac{d}{d\tau}\mathcal H_\omega(\tau)
  +c_A\mathcal A_\omega(\tau)
  +\mathcal R_\omega(\tau)\le0
\]
holds in distributions on \((\tau_0,\infty)\);
\item the validated compact cost is controlled by the carrier on the terminal
tail:
\[
  0\le \mathfrak C_{\rm II}^{NS}(\tau)
  \le C_A\mathcal A_\omega(\tau)
  \qquad\text{for a.e. }\tau\ge\tau_0 .
\]
\end{enumerate}
The carrier is local: \(\mathcal A_\omega\) may be the fixed cutoff identity
cost, a moving-cutoff identity cost, or another localized cost with an explicit
terminal comparison to the validated compact cost.
\end{definition}

\begin{lemma}[Monotonicity carrier gives local cost-to-exclusion data]
\label{paper7:lem:monotonicity-carrier-gives-cost-to-exclusion-data}
Let \(\omega\) be a repaired-gauge Navier--Stokes Type II branch.  Assume that
the validated compact cost \(\mathfrak C_{\rm II}^{NS}\) is nonnegative,
locally integrable on finite renormalized intervals, and satisfies
\[
  \int_{\tau_0}^{\infty}
  \mathfrak C_{\rm II}^{NS}(\tau)\,d\tau=\infty .
\]
If \(\mathfrak C_{\rm II}^{NS}\) admits a local monotonicity carrier in the
sense of \cref{paper7:def:local-monotonicity-carrier}, then \(\omega\) carries
the local cost-to-exclusion data of
\cref{paper7:def:local-cost-to-exclusion-data} with
\(\mathcal D_\omega=\mathfrak C_{\rm II}^{NS}\).
\end{lemma}

\begin{proof}
Let
\((\mathcal A_\omega,\mathcal H_\omega,\mathcal R_\omega^{\rm car},c_A,C_A)\)
be the carrier data.  The lower bound and finite initial value of
\(\mathcal H_\omega\) are exactly the energy bounds required in
\cref{paper7:def:local-cost-to-exclusion-data}.  The terminal comparison gives
\[
  \mathcal A_\omega(\tau)
  \ge C_A^{-1}\mathfrak C_{\rm II}^{NS}(\tau)
  \qquad\text{for a.e. }\tau\ge\tau_0 .
\]
Substituting this into the carrier monotonicity inequality yields
\[
  \frac{d}{d\tau}\mathcal H_\omega(\tau)
  +\frac{c_A}{C_A}\mathfrak C_{\rm II}^{NS}(\tau)
  +\mathcal R_\omega^{\rm car}(\tau)\le0
\]
in distributions.  Thus
\cref{paper7:def:local-cost-to-exclusion-data} holds with
\[
  \mathcal D_\omega=\mathfrak C_{\rm II}^{NS},
  \qquad c_\omega=c_A/C_A,
  \qquad \mathcal R_\omega=\mathcal R_\omega^{\rm car}.
\]
The divergent-cost item is the validated compact cost-divergence conclusion.
\end{proof}

\begin{lemma}[Fixed-cutoff monotonicity gives local cost-to-exclusion data]
\label{paper7:lem:fixed-monotonicity-gives-cost-to-exclusion-data}
Let \(\omega\) be represented on a terminal tail by a repaired-gauge
Navier--Stokes branch \((V,P,a,b)\).  Let \(\phi_{R_0}\) be the cutoff used in
the validated compact cost, let \(\mathfrak C_{\rm II}^{NS}\) be a
nonnegative locally integrable validated compact cost, and set
\[
  \mathcal E_\phi(\tau)
  =
  \frac12\int |V(y,\tau)|^2\phi_{R_0}(y)\,dy .
\]
Assume, after increasing \(\tau_0\) if necessary, that:
\begin{enumerate}[label=\textup{(\roman*)}]
\item \(\mathcal E_\phi(\tau_0)<\infty\);
\item the fixed-cutoff localized energy identity
\eqref{paper6a:eq:fixed-local-energy} holds on the terminal tail;
\item the fixed-cutoff error density \(B_{R_0}\) of
\cref{paper6a:def:fixed-error-density} satisfies
\[
  \int_{\tau_0}^{\infty}B_{R_0}(\tau)\,d\tau<\infty;
\]
\item the validated compact cost \(\mathfrak C_{\rm II}^{NS}\) is
energy-controlled by \(\tilde{\mathfrak D}_{R_0}\) in the sense of
\cref{paper7:def:energy-controlled-validated-cost};
\item the validated compact cost-divergence conclusion is
\[
  \int_{\tau_0}^{\infty}
  \mathfrak C_{\rm II}^{NS}(\tau)\,d\tau=\infty .
\]
\end{enumerate}
Then \(\omega\) carries the local cost-to-exclusion data of
\cref{paper7:def:local-cost-to-exclusion-data} with
\(\mathcal D_\omega=\mathfrak C_{\rm II}^{NS}\).
\end{lemma}

\begin{proof}
Let \(C_1,\tau_1\) be the constants in the energy-control comparison.  All
statements are terminal.  Since \(\mathfrak C_{\rm II}^{NS}\) is locally
integrable on finite intervals, removing any finite initial subinterval does
not change divergence of its tail integral.  After the allowed increase and
relabeling of \(\tau_0\), we may therefore assume that
\(\tau_0\ge\tau_1\), that \(\mathcal E_\phi(\tau_0)<\infty\), and that the
upper comparison holds on the whole tail.
\par
Define the corrected localized energy
\[
  \mathcal H_\omega(\tau)
  =
  \mathcal E_\phi(\tau)+\int_\tau^\infty B_{R_0}(s)\,ds .
\]
The finite-error hypothesis makes the correction tail finite for every
\(\tau\ge\tau_0\) and locally absolutely continuous.  Since
\(\mathcal E_\phi\ge0\) and \(B_{R_0}\ge0\), one has
\[
  \inf_{\tau\ge\tau_0}\mathcal H_\omega(\tau)\ge0,
  \qquad
  \mathcal H_\omega(\tau_0)<\infty .
\]
\Cref{paper6a:thm:fixed-corrected-monotonicity} gives, in distributions,
\[
  \frac{d}{d\tau}\mathcal H_\omega(\tau)
  +\frac12\tilde{\mathfrak D}_{R_0}(\tau)\le0 .
\]
By energy-control of the validated cost,
\[
  \tilde{\mathfrak D}_{R_0}(\tau)
  \ge C_1^{-1}\mathfrak C_{\rm II}^{NS}(\tau)
  \quad\text{for a.e. }\tau\ge\tau_0 .
\]
Therefore
\[
  \frac{d}{d\tau}\mathcal H_\omega(\tau)
  +\frac{1}{2C_1}\mathfrak C_{\rm II}^{NS}(\tau)\le0 .
\]
This is the corrected localized energy inequality in
\cref{paper7:def:local-cost-to-exclusion-data}, with
\[
  \mathcal D_\omega=\mathfrak C_{\rm II}^{NS},\qquad
  c_\omega=(2C_1)^{-1},\qquad
  \mathcal R_\omega=0 .
\]
The divergent-cost item is exactly the validated compact cost-divergence
conclusion.  Hence all local cost-to-exclusion data hold.
\end{proof}

\begin{lemma}[Moving-cutoff monotonicity gives local cost-to-exclusion data]
\label{paper7:lem:moving-monotonicity-gives-cost-to-exclusion-data}
Let \(\omega\) be represented on a terminal tail by a repaired-gauge
Navier--Stokes branch \((V,P,a,b)\), and let
\(\mathfrak C_{\rm II}^{NS}\) be a nonnegative locally integrable validated
compact cost.  Let
\[
  \Phi(\tau,y)=\phi(y/R(\tau))
\]
be a moving cutoff as in \cref{paper6a:def:moving-error-density}, with
\(R(\tau)\ge R_0\) locally absolutely continuous and nondecreasing, and set
\[
  \mathcal E_R(\tau)=\frac12\int |V(y,\tau)|^2\Phi(\tau,y)\,dy,
  \qquad
  \tilde{\mathfrak D}_{R(\tau)}(\tau)
  =
  \nu\int|\nabla V|^2\Phi
  +a_+(\tau)\int|V|^2\Phi .
\]
Assume, after increasing \(\tau_0\) if necessary, that:
\begin{enumerate}[label=\textup{(\roman*)}]
\item \(\mathcal E_R(\tau_0)<\infty\);
\item the localized energy identity is valid with the time-dependent cutoff
\(\Phi\);
\item the moving finite-error condition holds, for instance through summable
moving-annulus errors:
\[
  \int_{\tau_0}^{\infty}B_{\mathrm{mov}}(\tau)\,d\tau<\infty;
\]
\item the validated compact cost is controlled by the moving identity cost on
the terminal tail: for some \(C_m<\infty\),
\[
  0\le \mathfrak C_{\rm II}^{NS}(\tau)
  \le C_m\tilde{\mathfrak D}_{R(\tau)}(\tau)
  \qquad\text{for a.e. }\tau\ge\tau_0;
\]
\item the validated compact cost-divergence conclusion is
\[
  \int_{\tau_0}^{\infty}
  \mathfrak C_{\rm II}^{NS}(\tau)\,d\tau=\infty .
\]
\end{enumerate}
Then \(\omega\) carries the local cost-to-exclusion data of
\cref{paper7:def:local-cost-to-exclusion-data} with
\(\mathcal D_\omega=\mathfrak C_{\rm II}^{NS}\).
\end{lemma}

\begin{proof}
The moving finite-error condition gives
\(B_{\mathrm{mov}}\in L^1(\tau_0,\infty)\); if it is verified through
\cref{paper6a:def:summable-moving-errors}, this is precisely
\cref{paper6a:lem:summable-moving-errors}.  Define
\[
  \mathcal H_\omega(\tau)
  =
  \mathcal E_R(\tau)+\int_\tau^\infty B_{\mathrm{mov}}(s)\,ds .
\]
The finite moving-error condition makes the correction tail finite and locally
absolutely continuous.  Since \(\mathcal E_R\ge0\) and
\(B_{\mathrm{mov}}\ge0\),
\[
  \inf_{\tau\ge\tau_0}\mathcal H_\omega(\tau)\ge0,
  \qquad
  \mathcal H_\omega(\tau_0)<\infty .
\]
\Cref{paper6a:thm:moving-corrected-monotonicity} gives
\[
  \frac{d}{d\tau}\mathcal H_\omega(\tau)
  +\frac12\tilde{\mathfrak D}_{R(\tau)}(\tau)\le0
\]
in distributions.  Thus
\(\tilde{\mathfrak D}_{R(\tau)}\) is a local monotonicity carrier for the
validated compact cost, with \(c_A=1/2\), \(C_A=C_m\), and
\(\mathcal R_\omega=0\).  Applying
\cref{paper7:lem:monotonicity-carrier-gives-cost-to-exclusion-data} gives the
claimed local cost-to-exclusion data.
\end{proof}

\begin{lemma}[Corrected monotonicity bounds the validated cost]
\label{paper7:lem:corrected-monotonicity-bounds-cost}
Let \(\omega\) carry local cost-to-exclusion data in the sense of
\cref{paper7:def:local-cost-to-exclusion-data}, except do not assume the
divergence condition in item \textup{(i)}.  Then
\[
  \int_{\tau_0}^{\infty}\mathcal D_\omega(\tau)\,d\tau<\infty .
\]
\end{lemma}

\begin{proof}
Since \(\mathcal H_\omega\) is locally absolutely continuous and the corrected
localized energy inequality holds in distributions, integration over
\([\tau_0,T]\) gives, for every \(T>\tau_0\),
\[
  \mathcal H_\omega(T)-\mathcal H_\omega(\tau_0)
  +c_\omega\int_{\tau_0}^{T}\mathcal D_\omega(\tau)\,d\tau
  +\int_{\tau_0}^{T}\mathcal R_\omega(\tau)\,d\tau
  \le0 .
\]
The remainder is nonnegative, hence
\[
  c_\omega\int_{\tau_0}^{T}\mathcal D_\omega(\tau)\,d\tau
  \le
  \mathcal H_\omega(\tau_0)-\mathcal H_\omega(T)
  \le
  \mathcal H_\omega(\tau_0)
  -\inf_{\sigma\ge\tau_0}\mathcal H_\omega(\sigma).
\]
The right-hand side is finite and independent of \(T\).  Letting
\(T\to\infty\) and using monotone convergence for the nonnegative cost proves
the claim.
\end{proof}

\begin{theorem}[Local Navier--Stokes cost-to-exclusion theorem]
\label{paper7:thm:local-cost-to-exclusion}
Let \(\omega\) be a renormalized Navier--Stokes Type II branch on a retained
terminal compact stratum.  If \(\omega\) has active supercritical scaling, the
validated compact Type II cost-divergence conclusion, and local
cost-to-exclusion data of
\cref{paper7:def:local-cost-to-exclusion-data}, then \(\omega\) is excluded from
the Type II class under consideration.  Equivalently, on this local stratum,
\[
  \Pi_{\rm NS,II}(\omega):
  \mathcal S_\omega\wedge\mathcal B_\omega
  \Longrightarrow
  \mathcal E_\omega
\]
is a theorem.
\end{theorem}

\begin{proof}
Assume, for contradiction, that the branch remains in the Type II class on the
retained compact stratum.  The local cost-to-exclusion data give a corrected
localized energy inequality for the same nonnegative cost
\(\mathcal D_\omega\) whose divergent tail integral is the validated compact
cost-divergence conclusion.  By
\cref{paper7:lem:corrected-monotonicity-bounds-cost},
\[
  \int_{\tau_0}^{\infty}\mathcal D_\omega(\tau)\,d\tau<\infty .
\]
This contradicts the validated compact cost-divergence conclusion.  Hence the
branch cannot remain in the Type II class under consideration.
\end{proof}

\begin{definition}[Pointwise cost-divergence exclusion data]\label{paper7:def:cost-exclusion-data}
For a renormalized Type II solution \(\omega\), the pointwise
cost-divergence exclusion data are available when the following local analytic
statements hold:
\begin{enumerate}[label=\textup{(\roman*)}]
\item class membership: \(\omega\) lies in the renormalized Navier--Stokes
Type II class and satisfies the Type II analytic data before the Type II cost
criterion is evaluated;
\item cost compatibility: the active compact Type II cost criterion uses either
the Navier--Stokes identity-cost criterion
\[
  \mathfrak C_{\rm II}^{NS}=\tilde{\mathfrak D}_{R_0},
\]
or a localized cost explicitly comparable to
\(\tilde{\mathfrak D}_{R_0}\) on a terminal tail in the sense of
\cref{def:c4-cost-comparison}; in either case its conclusion is the validated
compact Type II cost-divergence conclusion for \(\omega\);
\item stability of the active scaling condition: the active supercritical
scaling condition remains in force after the repaired-gauge representation,
cost comparison, and criterion evaluation are attached;
\item stability of the cost-divergence conclusion: every renormalized
continuation preserving the active scaling condition and the validated compact
Type II cost-divergence conclusion remains in the same compact Type II
cost-divergence regime, unless it first enters one of the local alternatives
already isolated above;
\item local monotonicity-carrier data: on the retained compact stratum, the
validated compact cost admits a carrier in the sense of
\cref{paper7:def:local-monotonicity-carrier}.  This may be supplied by the
fixed or moving carrier subidentity of
\cref{paper7:def:carrier-subidentities}, the corresponding finite-error
condition, and the corresponding cost comparison.  Exact fixed or moving
localized energy identities are stronger sufficient inputs through
\cref{paper7:lem:fixed-monotonicity-gives-cost-to-exclusion-data,paper7:lem:moving-monotonicity-gives-cost-to-exclusion-data}.
\end{enumerate}
\end{definition}

\begin{lemma}[Canonical cost compatibility]
\label{paper7:lem:canonical-cost-compatibility}
Let \(\omega\) be represented on a terminal tail by a repaired-gauge
Navier--Stokes branch \((V,P,a,b)\).  Assume that the localized identity cost
\[
  \tilde{\mathfrak D}_{R_0}(\tau)
  =
  \nu\int_{\mathbb R^3}|\nabla V(y,\tau)|^2\phi_{R_0}(y)\,dy
  +
  a_+(\tau)\int_{\mathbb R^3}|V(y,\tau)|^2\phi_{R_0}(y)\,dy
\]
is measurable, nonnegative, and locally integrable on finite
\(\tau\)-intervals.  If the compact Type II criterion is evaluated with the
canonical Navier--Stokes identity cost, then
\[
  \mathfrak C_{\rm II}^{NS}=\tilde{\mathfrak D}_{R_0}
\]
is a valid Type II scale-exclusion cost, the comparison holds with
\[
  c_0=1,\qquad \tau_1=\tau_0,
\]
and the validated compact cost-divergence conclusion for \(\omega\) is the
identity-cost divergence statement
\[
  \int_{\tau_0}^{\infty}
  \tilde{\mathfrak D}_{R_0}(\tau)\,d\tau=\infty .
\]
If the compact criterion is not evaluated with this identity cost, the branch
is not being treated by the canonical compatibility statement and must instead
come with an explicit comparison as in
\cref{paper7:lem:comparable-cost-compatibility}.
\end{lemma}

\begin{proof}
The first assertion is the identity-cost case of the localized cost
definitions.  Since
\(\tilde{\mathfrak D}_{R_0}\) is assumed measurable, nonnegative, and
locally integrable on finite intervals, the canonical choice
\[
  \mathfrak C_{\rm II}^{NS}:=\tilde{\mathfrak D}_{R_0}
\]
satisfies the defining properties of a Type II scale-exclusion cost.  The
comparison inequality is the identity
\[
  \mathfrak C_{\rm II}^{NS}(\tau)
  =
  \tilde{\mathfrak D}_{R_0}(\tau)
  \ge
  1\cdot\tilde{\mathfrak D}_{R_0}(\tau)
\]
for a.e. \(\tau\ge\tau_0\), so \(c_0=1\) and \(\tau_1=\tau_0\).  Therefore
divergence of the active compact cost is precisely divergence of
\(\tilde{\mathfrak D}_{R_0}\) on the terminal tail.  This is the
canonical case of \cref{lem:c4-identity-cost,thm:c4-cost-comparison}.
\end{proof}

\begin{lemma}[Comparable cost compatibility]
\label{paper7:lem:comparable-cost-compatibility}
Let \(\omega\) be represented on a terminal tail by a repaired-gauge
Navier--Stokes branch \((V,P,a,b)\).  Let
\(\mathfrak C_{\rm II}^{NS}\) be a nonnegative, measurable, locally integrable
localized cost on finite \(\tau\)-intervals.  Assume that, on a later terminal
tail \([\tau_1,\infty)\), there is a constant \(c_0>0\) such that
\[
  \mathfrak C_{\rm II}^{NS}(\tau)
  \ge
  c_0\,\tilde{\mathfrak D}_{R_0}(\tau)
  \quad\text{for a.e. }\tau\ge\tau_1 .
\]
Then every identity-cost divergence conclusion
\[
  \int_{\tau_1}^{\infty}
  \tilde{\mathfrak D}_{R_0}(\tau)\,d\tau=\infty
\]
implies the corresponding compact Type II cost-divergence conclusion
\[
  \int_{\tau_1}^{\infty}
  \mathfrak C_{\rm II}^{NS}(\tau)\,d\tau=\infty .
\]
Consequently the use of \(\mathfrak C_{\rm II}^{NS}\) introduces no additional
exit from the compact-cost stratum beyond the explicit failure of this
comparison.
\end{lemma}

\begin{proof}
For every \(T>\tau_1\), the assumed comparison gives
\[
  \int_{\tau_1}^{T}
  \mathfrak C_{\rm II}^{NS}(\tau)\,d\tau
  \ge
  c_0\int_{\tau_1}^{T}
  \tilde{\mathfrak D}_{R_0}(\tau)\,d\tau .
\]
Letting \(T\to\infty\) proves the implication from identity-cost divergence to
divergence of the selected compact Type II cost.  This is precisely the
one-sided comparison needed by \cref{def:c4-cost-comparison} and used in
\cref{thm:c4-cost-comparison}.  If the inequality is unavailable, the selected
diagnostic cost is not connected to the Navier--Stokes identity dissipation on
this terminal tail; this is a cost-comparison failure, not a new unlabelled
Type II branch.
\end{proof}

\begin{lemma}[Local cost well-posedness]
\label{paper7:lem:local-cost-well-posed}
Let \((V,P,a,b)\) be a repaired-gauge branch on a finite renormalized interval
\(I=[\alpha,\beta]\), and let \(\phi_{R_0}\in C_c^\infty(B_{R_0})\),
\(\phi_{R_0}\ge0\).  Assume:
\begin{enumerate}[label=\textup{(\roman*)}]
\item \(V\in L^2(I;H^1(B_{R_0}))\);
\item \(a\in L^1(I)\);
\item either \(V\in L^\infty(I;L^3(\mathbb R^3))\), or more locally
\[
  \operatorname*{ess\,sup}_{\tau\in I}
  \|V(\tau)\|_{L^2(B_{R_0})}<\infty .
\]
\end{enumerate}
Then the localized identity cost
\[
  \tilde{\mathfrak D}_{R_0}(\tau)
  =
  \nu\int |\nabla V(y,\tau)|^2\phi_{R_0}(y)\,dy
  +
  a_+(\tau)\int |V(y,\tau)|^2\phi_{R_0}(y)\,dy
\]
is measurable, nonnegative, and integrable on \(I\).  Consequently, on every
compact renormalized window satisfying these local bounds,
\(\tilde{\mathfrak D}_{R_0}\) is a well-defined finite-interval cost.
\end{lemma}

\begin{proof}
The two summands are measurable because \(V\) is strongly measurable in the
local Sobolev spaces on the finite interval and \(a_+\) is measurable.  They are
nonnegative since \(\nu>0\), \(\phi_{R_0}\ge0\), and \(a_+\ge0\).
\par
For the gradient term, boundedness of \(\phi_{R_0}\) and the support condition
give
\[
  \int_I\int |\nabla V|^2\phi_{R_0}
  \le
  \|\phi_{R_0}\|_{L^\infty}
  \int_I\|\nabla V(\tau)\|_{L^2(B_{R_0})}^2\,d\tau
  <\infty .
\]
For the drift term, first suppose
\[
  \operatorname*{ess\,sup}_{\tau\in I}
  \|V(\tau)\|_{L^2(B_{R_0})}<\infty .
\]
Then
\[
  \int_I a_+(\tau)\int |V|^2\phi_{R_0}\,dy\,d\tau
  \le
  \|\phi_{R_0}\|_{L^\infty}
  \left(\operatorname*{ess\,sup}_{\tau\in I}
  \|V(\tau)\|_{L^2(B_{R_0})}^2\right)
  \|a_+\|_{L^1(I)}
  <\infty .
\]
If instead \(V\in L^\infty(I;L^3(\mathbb R^3))\), then Hölder on \(B_{R_0}\)
gives the local \(L^2\) bound
\[
  \|V(\tau)\|_{L^2(B_{R_0})}^2
  \le
  |B_{R_0}|^{1/3}\|V(\tau)\|_{L^3(\mathbb R^3)}^2
\]
for a.e. \(\tau\in I\), reducing to the previous case.  Thus both summands are
integrable on \(I\).
\end{proof}

\begin{lemma}[Local modulation integrability from absolutely continuous repaired gauges]
\label{paper7:lem:local-modulation-integrability}
Let \(u,p\) be a Navier--Stokes solution on a terminal physical time interval,
and let \((V,P)\) be its repaired-gauge renormalization on a terminal
renormalized tail.  Assume that, on every compact renormalized subinterval, the
physical concentration chart and the repaired-gauge path satisfy
\cref{paper6a:def:ac-initial-concentration-chart,paper6a:def:ac-repaired-gauge-path}.
Let \(a,b\) be the final modulation coefficients in the repaired-gauge
renormalized equation, as in
\cref{paper6a:lem:repaired-modulation-coefficients}.  Then
\[
  a,b\in L^1_{\rm loc}([\tau_0,\infty)).
\]
In particular, for every unit terminal window
\[
  I_n=[\tau_0+n,\tau_0+n+1],
\]
one has \(a\in L^1(I_n)\).  If the absolutely continuous chart, the
absolutely continuous repaired gauge, or the resulting local integrability of
the modulation coefficients is unavailable on such a compact window, then the
branch does not lie in the retained compact-cost stratum and falls into the
corresponding representation, repaired-gauge, or modulation-integrability
alternative.
\end{lemma}

\begin{proof}
On each compact renormalized window, the assumed absolutely continuous chart
and repaired-gauge path place the branch under the hypotheses of
\cref{paper6a:thm:ac-repaired-gauge-consequences}.  In particular,
\cref{paper6a:lem:ac-chart-regularity} gives
\[
  a,b\in L^1_{\rm loc},
\]
and \cref{paper6a:lem:suitable-gauge-pressure-modulation} identifies these
locally integrable functions as the final coefficients in the repaired-gauge
equation.  Restricting to \(I_n\) gives \(a\in L^1(I_n)\).  If the absolutely
continuous local chart, repaired-gauge path, or resulting coefficient
integrability is missing, then the ordered stratification records the first
missing input as a representation, repaired-gauge, or modulation-integrability
failure.
\end{proof}

\begin{lemma}[Local state data give local \(L^2\) windows]
\label{paper7:lem:local-state-data-local-l2}
Let \((V,P,a,b)\) be a repaired-gauge Type II branch on a terminal tail.
Fix \(R_0<\infty\) and a unit terminal window \(I_n\).  Suppose the retained
local suitable-energy data on the compact repaired-gauge cylinder supply
\[
  \operatorname*{ess\,sup}_{\tau\in I_n}
  \|V(\tau)\|_{L^2(B_{R_0})}<\infty,
\]
Then
\[
  \operatorname*{ess\,sup}_{\tau\in I_n}
  \|V(\tau)\|_{L^2(B_{R_0})}
\]
is finite.  Thus the local \(L^2\)-input required for the finite-window
identity cost is available on every retained compact cylinder from local
suitable-energy data alone.  Failure of local \(L^2\) control on the selected
compact cylinder is recorded by the ordered local validation as a local-energy
or rough-core exit.
\end{lemma}

\begin{proof}
The displayed local \(L^\infty_\tau L^2_y(B_{R_0})\) bound is precisely the
compact-cylinder local-energy component of the retained state data.  Hence the
claimed finiteness is immediate.  If that local \(L^2\) control were absent,
the ordered validation would stop at the local-energy or rough-core row before
any compact cost is evaluated.  Therefore the compact-cost channel never uses
the later compact-cost validation layer to manufacture the needed local \(L^2\)
input; the input is already part of the retained local state data.
\end{proof}

\begin{corollary}[Compact-window data give cost well-posedness]
\label{paper7:cor:compact-window-data-cost-well-posed}
Let \((V,P,a,b)\) be a repaired-gauge branch on a terminal tail.  Fix
\(R_0<\infty\).  Suppose that, for every unit terminal window
\[
  I_n=[\tau_0+n,\tau_0+n+1],
  \qquad n\in\mathbb N,
\]
the branch has local windowed \(H^1(B_{R_0})\) control and satisfies one of
the local \(L^2\)-window hypotheses of
\cref{paper7:lem:local-state-data-local-l2}.  Suppose also that the local
modulation-integrability conclusion of
\cref{paper7:lem:local-modulation-integrability} holds.  Then
\(\tilde{\mathfrak D}_{R_0}\) is a measurable nonnegative cost, integrable
on each \(I_n\).  If the branch also has the identity cost or an explicitly
comparable localized cost, then the compact Type II cost criterion can be
evaluated on every retained terminal window.
\par
Failure of the conclusion is exactly one of the local failures in the
hypotheses: compact-cylinder \(L^2\) control, windowed \(H^1\) control,
modulation-integrability, or cost-comparison failure.
\end{corollary}

\begin{proof}
\Cref{paper7:lem:local-state-data-local-l2} supplies the
\(L^\infty(I_n;L^2(B_{R_0}))\) input on every unit window, and
\cref{paper7:lem:local-modulation-integrability} supplies \(a\in L^1(I_n)\).
The assumed windowed \(H^1(B_{R_0})\) control gives
\[
  V\in L^2(I_n;H^1(B_{R_0})).
\]
Thus the hypotheses of \cref{paper7:lem:local-cost-well-posed} hold on each
\(I_n\).  The identity-cost case is then covered by
\cref{paper7:lem:canonical-cost-compatibility}; the explicitly comparable case
is covered by \cref{paper7:lem:comparable-cost-compatibility}.  Hence the
criterion is locally meaningful on every retained terminal window.  Conversely,
if the conclusion fails before criterion evaluation, at least one of the listed
local inputs is absent; by the ordered stratification, that absence is recorded
as the corresponding local alternative.
\end{proof}

\begin{proposition}[Active scaling stability criterion]
\label{paper7:prop:active-scaling-stability-package}
Let \(\omega\) be an admissible local Type II branch with active supercritical
scaling on a terminal tail.  A passage from a raw Type II chart to a
repaired-gauge chart is called scale-admissible on a later tail if its raw and
repaired physical scales satisfy
\[
  0<m\le
  \frac{\lambda_{\rm rep}(t)}{\lambda_{\rm raw}(t)}
  \le M<\infty
\]
on that tail.  On the compact-cost stratum, the active supercritical scaling
condition is unchanged by the following operations:
\begin{enumerate}[label=\textup{(\roman*)}]
\item restriction to a later terminal tail;
\item passage from an admissible raw Type II chart to a scale-admissible
repaired-gauge representation;
\item replacement of the canonical cost by an explicitly comparable localized
cost;
\item evaluation of the compact Type II cost criterion.
\end{enumerate}
If the repaired chart is not scale-admissible on any later tail, the branch is
not retained in the compact-cost stratum; it is assigned instead to the
corresponding representation, repaired-gauge, or scale-degeneracy alternative.
\end{proposition}

\begin{proof}
The active supercritical scaling condition is a terminal property of the
scale-center chart and is invariant under removal of a finite initial
renormalized interval.  Thus restriction from \([\tau_0,\infty)\) to
\([\tau_1,\infty)\), \(\tau_1\ge\tau_0\), cannot change the terminal scale
alternative.
\par
Suppose next that a scale-admissible repaired-gauge representation is
available.  Two-sided comparability of \(\lambda_{\rm rep}\) and
\(\lambda_{\rm raw}\) implies that any terminal statement depending only on the
scale class, such as supercriticality relative to the Type I rate, holds for
one scale if and only if it holds for the other after passing to the same
tail.  The repaired gauge changes coordinates and modulation variables inside
this comparable scale class; it does not replace the branch by a different
terminal scaling alternative.  If the representation cannot be constructed, or
if every repaired chart loses scale comparability, then the first failure is a
representation, repaired-gauge, or scale-degeneracy alternative.
\par
Cost comparison is scale-passive.  It changes only the nonnegative diagnostic
functional used to record cost divergence and does not alter the center,
physical scale, renormalized time, or active scaling hypothesis.  If the
comparison is unavailable, the failure is cost-comparison failure.  Finally,
criterion evaluation is a logical operation applied to the already constructed
branch and cost.  It does not modify the branch or its scale.  Hence active
scaling is stable through the listed operations exactly on the local stratum on
which those operations are admissible.
\end{proof}

\begin{lemma}[No unrecorded exit from compact cost divergence]
\label{paper7:lem:no-unrecorded-cost-escape}
Let \(\omega\) be an admissible local Type II branch with active supercritical
scaling on a terminal tail.  Before declaring that a continuation lies in the
compact cost-divergence stratum, run the ordered local validation tests:
\begin{enumerate}[label=\textup{(\roman*)}]
\item Type II analytic data and exhaustion data;
\item raw-chart selection, repaired-gauge representation, and scale
admissibility;
\item local compactness decomposition into the alternatives of
\cref{paper6:thm:state-decomposition};
\item retained finite critical-mass window, read as the compact-window
finite critical-mass input of
\cref{paper6a:lem:terminal-local-critical-mass-routing};
\item the critical-tightness/exterior-leakage test at the retained scale;
\item local windowed \(H^1\) control;
\item local modulation integrability, cost well-posedness, and compatibility of
the chosen compact Type II cost;
\item local monotonicity-carrier data for the validated compact cost, supplied
by the fixed-cutoff criterion, the moving-cutoff criterion, or another explicit
localized carrier;
\item the compact Type II cost-divergence conclusion for the validated cost.
\end{enumerate}
Then exactly one of the following occurs:
\begin{enumerate}[label=\textup{(\alph*)}]
\item all tests pass, and \(\omega\) remains in the compact Type II
cost-divergence regime on the retained terminal branch;
\item the first failed test gives one of the already named local alternatives:
ordered exhaustion failure, representation failure, repaired-gauge or
scale-degeneracy failure, multibubble/multicore splitting, cascade, exterior
tightness failure, loss of local windowed \(H^1\) control, finite-cost
scale-collapse alternative, scale-rigid residual alternative, critical-mass
alternative, modulation-integrability failure, cost-compatibility failure,
carrier-subidentity failure, finite fixed-cutoff error failure, summable
moving-annulus error failure, carrier-comparison failure for the validated
cost, or failure of compact cost divergence for the validated cost.
\end{enumerate}
Thus a renormalized continuation preserving active scaling and validated compact
cost divergence cannot leave the compact cost-divergence regime without being
recorded by the local compactness stratification.
\end{lemma}

\begin{proof}
The proof is the ordered-alternative argument.  First test whether the Type II
analytic and exhaustion data hold.  If not, the first failure is one of the
ordered exhaustion failures of \cref{def:c1-exhaustion-failures}.  If they
hold, test for raw-chart selection, repaired-gauge representation, and scale
admissibility.  Failure at this stage is one of the ordered representation
failures of \cref{def:c2-representation-failures}, possibly with the specific
scale-degeneracy recorded by
\cref{paper7:prop:active-scaling-stability-package}.
\par
Once a scale-admissible repaired branch is available, apply the local
compactness decomposition \cref{paper6:thm:state-decomposition}.  If the branch
falls into the multibubble, multicore, cascade, gauge-degenerate, finite-cost
scale-collapse, local \(H^1\)-loss, or scale-rigid residual alternatives, it has
left the compact single-core cost-divergence stratum through a named local
case.  In particular, an additional exterior concentration mechanism is not
discarded; if it retains critical mass it appears as a multibubble or multicore
branch, and if it escapes the retained compact region it appears as failure of
critical tightness.
\par
The first alternative in \cref{paper6:thm:state-decomposition} is retained only
as a compact single-core branch; it is not yet identified with compact cost
divergence.  On that retained branch the remaining tests below decide whether
the branch belongs to the compact cost-divergence stratum, has finite validated
cost, or exits through one of the remaining compact tests.
\par
On the retained compact single-core branch, test the finite critical-mass row
using the compact-window version of
\cref{paper6a:lem:terminal-local-critical-mass-routing}: failure of the
retained spacetime activity floor means non-retention, and failure of a finite
local upper bound produces
a comparable retained label or a named adjacent active label.  Thus failure of
the finite critical-mass row is recorded as a critical-mass or adjacent local
alternative, not as a hidden global estimate.  Once the row passes, the
remaining compact tests are critical tightness and local windowed \(H^1\)
control.  Failure of either is precisely one of the alternatives in
\cref{def:c5-compact-tests}.
\par
If all compact tests pass, local cost well-posedness and compatibility are
evaluated.  Local modulation integrability is supplied by
\cref{paper7:lem:local-modulation-integrability}; the local cost
well-posedness is supplied by
\cref{paper7:cor:compact-window-data-cost-well-posed}.  The compatibility is
the identity case of \cref{paper7:lem:canonical-cost-compatibility} or the
explicit comparison case of
\cref{paper7:lem:comparable-cost-compatibility}; failure of both is a
cost-compatibility failure.  Next test whether the validated compact cost has a
local monotonicity carrier.  In the fixed-cutoff case this means the fixed
carrier subidentity, finite fixed-cutoff error, and energy-control comparison.
In the moving-cutoff case this means the moving carrier subidentity, summable
moving-annulus errors, and moving-cost comparison.  Exact localized energy
identities are sufficient for these subidentities, but the one-sided suitable
local energy inequality is the actual input needed.  Failure of these inputs is
recorded as a carrier-subidentity failure, a finite fixed-cutoff error failure,
a summable moving-annulus error failure, or a carrier-comparison failure for
the validated cost.
Finally, evaluate the compact cost-divergence conclusion for the validated cost.
If it fails, the branch is a finite-cost or non-divergent compact branch and is
not counted as a compact cost-divergence branch.  The tests are ordered, so the
first failure is unique.  If no failure occurs, the branch has the analytic data,
scale-admissible repaired-gauge representation, compact single-core
local compactness position, the retained finite critical-mass window, critical
tightness, local windowed \(H^1\) control, local modulation integrability, a
compatible compact cost, local monotonicity-carrier data, and the validated
compact cost-divergence conclusion.  Together with the assumed active scaling,
these hypotheses place the branch in the compact Type II cost-divergence
regime.
Therefore there is no unrecorded exit.
\end{proof}

\begin{theorem}[Carrier validation alternative on the retained compact stratum]
\label{paper7:thm:carrier-validation-retained-stratum}
Let \(\omega\) be an admissible local Type II branch with active supercritical
scaling.  Suppose that the ordered validation tests in
\cref{paper7:lem:no-unrecorded-cost-escape} have reached the retained compact
single-core branch, that cost well-posedness and cost compatibility have
passed, and that the validated compact cost-divergence conclusion holds.  Then
exactly one of the following alternatives occurs:
\begin{enumerate}[label=\textup{(\roman*)}]
\item the validated compact cost admits local monotonicity-carrier data in the
sense of \cref{paper7:def:local-monotonicity-carrier};
\item the branch leaves the retained compact cost-divergence stratum through
the first failed carrier-validation test: carrier-subidentity failure, finite
fixed-cutoff error failure, summable moving-annulus error failure, or
carrier-comparison failure for the validated cost.
\end{enumerate}
Consequently, on the retained compact cost-divergence stratum, where the
adjacent carrier-failure alternatives are not part of the stratum, the
validated compact cost admits local monotonicity-carrier data.
\end{theorem}

\begin{proof}
After the branch has passed the Type II analytic, representation,
compactness, retained finite critical mass, critical-tightness, local windowed \(H^1\),
modulation, cost well-posedness, and cost-compatibility tests, the next ordered
test in \cref{paper7:lem:no-unrecorded-cost-escape} is precisely the
local monotonicity-carrier test for the validated compact cost.  If this test
passes, alternative \textup{(i)} holds.  If it fails, the first failed input is
recorded as one of the named adjacent carrier failures listed in
\cref{paper7:lem:no-unrecorded-cost-escape}: failure of the carrier
subidentity, failure of the fixed finite-error condition, failure of the moving
finite-error condition through the summable moving-annulus test, or failure of
the carrier comparison for the validated cost.  The ordered first-failure
convention makes the two alternatives disjoint and exhaustive.
\par
By definition, the retained compact cost-divergence stratum keeps only the
branch for which all preceding compact tests and the carrier-validation test
have passed before the validated compact cost-divergence conclusion is
attached.  Therefore a branch that remains in this stratum cannot realize one
of the carrier-failure alternatives, and the local monotonicity carrier exists.
\end{proof}

\begin{proposition}[Local verification of the cost-exclusion data]
\label{paper7:prop:local-cost-exclusion-data-verification}
On a retained terminal branch in the compact Type II cost-divergence stratum,
the pointwise data of \cref{paper7:def:cost-exclusion-data} are local
analytic statements, not additional global assumptions.  More precisely, assume
that the branch:
\begin{enumerate}[label=\textup{(\roman*)}]
\item has the local Type II analytic data and has not entered an exhaustion,
representation, modulation-integrability, multibubble/multicore, cascade,
exterior-tightness, scale-collapse, scale-rigid, critical-mass, or local
\(H^1\)-loss alternative;
\item admits a scale-admissible repaired-gauge representation on retained
windows from absolutely continuous local scale-center and repaired-gauge data;
\item has the retained finite critical-mass window, critical tightness at the
retained scale, local windowed \(H^1\) control, active supercritical scaling,
and the compact cost-divergence conclusion for the validated cost;
\item uses either the canonical identity cost
\(\tilde{\mathfrak D}_{R_0}\) or an explicitly comparable localized cost;
\item remains in the retained compact cost-divergence stratum after the
carrier-validation alternatives of
\cref{paper7:thm:carrier-validation-retained-stratum} have been evaluated.
\end{enumerate}
Then class membership, cost compatibility, stability of the active scaling
condition, stability of the validated cost-divergence conclusion, and the
cost-to-exclusion implication all hold for this branch.  Consequently, on this
local stratum, all clauses of \cref{paper7:def:cost-exclusion-data} are verified;
no estimate outside the retained local stratum and no exclusion of secondary
exterior concentration are being assumed.
\end{proposition}

\begin{proof}
Class membership is the retained local Type II analytic-data alternative after
the ordered exhaustion and representation tests have passed.
The local inputs needed to evaluate the compact cost are available by
\cref{paper7:lem:local-modulation-integrability,paper7:lem:local-state-data-local-l2}
and \cref{paper7:cor:compact-window-data-cost-well-posed}.  Cost compatibility
then follows from
\cref{paper7:lem:canonical-cost-compatibility} in the identity-cost case and
from \cref{paper7:lem:comparable-cost-compatibility} in the explicitly
comparable-cost case.  Stability of the active scaling condition is
\cref{paper7:prop:active-scaling-stability-package}.  Stability of the
validated cost-divergence conclusion, with all exits assigned to named local
alternatives, is \cref{paper7:lem:no-unrecorded-cost-escape}.  These prove
clauses \textup{(i)}--\textup{(iv)} of
\cref{paper7:def:cost-exclusion-data}.  Since the branch remains in the
retained compact cost-divergence stratum after carrier validation,
\cref{paper7:thm:carrier-validation-retained-stratum} gives local
monotonicity-carrier data for the validated compact cost.  These data provide
the local cost-to-exclusion data through
\cref{paper7:lem:monotonicity-carrier-gives-cost-to-exclusion-data}.  Then
\cref{paper7:thm:local-cost-to-exclusion} gives
\(\Pi_{\rm NS,II}(\omega)\), which is the analytic content of clause
\textup{(v)}.
\end{proof}

\begin{definition}[Classwise retained compact coverage]\label{paper7:def:classwise-retained-compact-coverage}
The renormalized Type II class has classwise retained compact coverage for the
validated compact cost if every branch with active supercritical scaling and
validated compact cost divergence either remains in the retained compact
cost-divergence stratum after the ordered validation tests of
\cref{paper7:lem:no-unrecorded-cost-escape}, or has already entered one of the
named adjacent alternatives in that ordered stratification.
\end{definition}

\begin{theorem}[Classwise retained compact coverage from ordered validation]
\label{paper7:thm:classwise-retained-compact-coverage}
Every admissible local Type II branch with active supercritical scaling and
validated compact Type II cost divergence has classwise retained compact
coverage in the sense of
\cref{paper7:def:classwise-retained-compact-coverage}.
\end{theorem}

\begin{proof}
Apply the ordered validation tests of
\cref{paper7:lem:no-unrecorded-cost-escape} to the branch.  The lemma gives two
possibilities.  If all tests pass, the branch remains in the compact Type II
cost-divergence regime on the retained terminal branch, which is precisely the
retained compact cost-divergence stratum.
\par
If a test fails, the first failed test is one of the named local alternatives in
\cref{paper7:lem:no-unrecorded-cost-escape}.  Since the present theorem assumes
the validated compact cost-divergence conclusion, the last listed failure,
failure of compact cost divergence for the validated cost, cannot occur.  Every
other first failure is an adjacent compactness stratum: exhaustion,
representation, repaired-gauge or scale-degeneracy failure,
multibubble/multicore splitting, cascade, exterior tightness failure, local
\(H^1\)-loss, finite-cost scale-collapse, scale-rigid residual behavior,
critical-mass failure, modulation-integrability failure, cost-compatibility
failure, carrier-subidentity failure, fixed-cutoff finite-error failure,
moving-annulus summability failure, or carrier-comparison failure.  Therefore
every branch satisfying the hypotheses either remains retained or has already
entered a named adjacent alternative.  This is exactly classwise retained
compact coverage.
\end{proof}

\begin{remark}[Scope of classwise retained compact coverage]
\label{paper7:rem:uniform-cost-scope}
\Cref{paper7:prop:local-cost-exclusion-data-verification} is deliberately local.
It proves the interface hypotheses for a branch only after the branch has been
placed in the validated compact cost-divergence stratum by the compactness
stratification.  Branches with multibubble, cascade, exterior leakage,
scale-collapse, scale-rigid, or rough-core behavior are not forced into this
stratum; they are assigned to their own local alternatives.  Thus the
cost-divergence exclusion step is unconditional on the retained compact
stratum, while a classwise statement is obtained only after the adjacent
strata have been treated by their own local arguments.  The coverage part itself
is no longer an assumption; it is supplied by
\cref{paper7:thm:classwise-retained-compact-coverage}.
\end{remark}

\begin{definition}[Abstract cost-divergence exclusion theorem applicability]
\label{paper7:def:abstract-cost-exclusion-applicability}
An abstract Type II cost-divergence exclusion theorem is applicable to a
renormalized Type II solution if the following analytic identifications are
available:
\begin{enumerate}[label=\textup{(\roman*)}]
\item an abstract Type II exclusion theorem based on localized
renormalization-cost divergence;
\item hypotheses supplying supercritical scaling, a Type II concentration
scenario with bounded physical energy, and a localized energy monotonicity
formula at the active scale \(\lambda(t)\);
\item an identification of the theorem's renormalization-cost integral with a
validated compact Navier--Stokes cost \(\mathfrak C_{\rm II}^{NS}\), where
either
\[
  \mathfrak C_{\rm II}^{NS}=\tilde{\mathfrak D}_{R_0}
\]
or \(\mathfrak C_{\rm II}^{NS}\) is a localized cost satisfying the terminal
one-sided comparison with \(\tilde{\mathfrak D}_{R_0}\) in
\cref{def:c4-cost-comparison};
\item a domain embedding placing the repaired-gauge Navier--Stokes orbit and
its scale variables into the parabolic domain expected by the theorem without
changing the active scaling condition, validated cost-divergence conclusion, or
exclusion conclusion;
\item identification of the theorem's conclusion with
\cref{paper7:def:cost-exclusion-conclusion}.
\end{enumerate}
\end{definition}

\begin{lemma}[Applicability of the abstract cost-divergence theorem]
\label{paper7:lem:abstract-cost-exclusion}
Assume a renormalized Type II solution has the active supercritical scaling
condition, the validated compact Type II cost-divergence conclusion, and the
abstract applicability data of
\cref{paper7:def:abstract-cost-exclusion-applicability}.
Then the solution is excluded from the Type II class under consideration.
\end{lemma}

\begin{proof}
The abstract Type II theorem is applied with supercritical scaling,
bounded-energy Type II concentration, and the localized energy monotonicity
formula at the active scale.  The cost identification places the validated
compact Navier--Stokes cost in the theorem's renormalization-cost integral.  In
the identity case this cost is \(\tilde{\mathfrak D}_{R_0}\).  In the
comparable-cost case, the terminal one-sided comparison in
\cref{def:c4-cost-comparison} and
\cref{paper7:lem:comparable-cost-compatibility} show that the validated
Navier--Stokes cost-divergence conclusion implies divergence of the localized
cost used by the abstract theorem.  The domain embedding places the
repaired-gauge Navier--Stokes solution in the theorem's parabolic domain, and
the conclusion is identified with \cref{paper7:def:cost-exclusion-conclusion}.
Therefore the theorem gives exclusion from the Type II class.
\end{proof}

\begin{theorem}[Pointwise Navier--Stokes cost-divergence exclusion]
\label{paper7:thm:pointwise-cost-exclusion}
Let \(\omega\) be a renormalized Type II solution.  If \(\omega\) has the
active supercritical scaling condition, the validated compact Type II
cost-divergence conclusion, and the pointwise data of
\cref{paper7:def:cost-exclusion-data}, then \(\omega\) is excluded from the
Type II class under consideration.
\end{theorem}

\begin{proof}
Class membership and stability of the active scaling condition ensure that
\(\omega\) is evaluated in the renormalized Type II class and that the active
supercritical scaling condition is still the relevant one.  Cost compatibility
ensures that the validated compact cost-divergence conclusion is the
Navier--Stokes compact cost conclusion for \(\omega\), not an external or
incomparable cost statement.  Stability of the cost-divergence conclusion rules
out an exit from the compact Type II cost-divergence regime to a Type II case
outside the compact cost-divergence criterion inside the same class.
\par
It remains only to justify the transition from cost divergence to exclusion.
The local monotonicity-carrier data in \cref{paper7:def:cost-exclusion-data}
give the local cost-to-exclusion data by
\cref{paper7:lem:monotonicity-carrier-gives-cost-to-exclusion-data}.
Therefore
\cref{paper7:thm:local-cost-to-exclusion} applies to the active supercritical
scaling condition and the validated compact cost-divergence conclusion, giving
the claimed exclusion.  In an abstract theorem setting, the same implication is
also licensed by \cref{paper7:lem:abstract-cost-exclusion}.
\end{proof}

\begin{corollary}[Cost-divergence exclusion on the retained compact stratum]
\label{paper7:cor:retained-stratum-cost-exclusion}
Let \(\omega\) be a renormalized Type II solution with active supercritical
scaling and validated compact Type II cost divergence.  If \(\omega\) remains
in the retained compact cost-divergence stratum after the ordered validation
tests of \cref{paper7:lem:no-unrecorded-cost-escape}, then \(\omega\) is
excluded from the Type II class under consideration.
\end{corollary}

\begin{proof}
By \cref{paper7:prop:local-cost-exclusion-data-verification}, the pointwise
cost-divergence exclusion data of \cref{paper7:def:cost-exclusion-data} are
verified on this retained stratum.  Applying
\cref{paper7:thm:pointwise-cost-exclusion} gives the exclusion conclusion.
\end{proof}

\begin{definition}[Adjacent-stratum closure for the compact-cost regime]
\label{paper7:def:compact-cost-adjacent-closure}
The compact-cost regime has adjacent-stratum closure if every named adjacent
alternative produced by
\cref{paper7:lem:no-unrecorded-cost-escape} for a branch with active
supercritical scaling and validated compact Type II cost divergence has its own
stratum-local argument giving the same Type II exclusion conclusion, either
directly or by placing the branch in a named non-compact-cost
alternative whose closure theorem already gives the Type II exclusion
conclusion.
\end{definition}

\begin{corollary}[Uniform Navier--Stokes cost-divergence exclusion]
\label{paper7:cor:uniform-cost-exclusion}
Assume adjacent-stratum closure for the compact-cost regime in the sense of
\cref{paper7:def:compact-cost-adjacent-closure}.  Then every
renormalized Type II solution with active supercritical scaling condition and
validated compact Type II cost-divergence conclusion is either excluded in its
adjacent stratum or excluded by
\cref{paper7:cor:retained-stratum-cost-exclusion}.
\end{corollary}

\begin{proof}
By \cref{paper7:thm:classwise-retained-compact-coverage}, a branch with active
supercritical scaling and validated compact cost divergence either lies in one
of the named adjacent strata or remains in the retained compact
cost-divergence stratum.  The adjacent cases are handled by the
adjacent-stratum closure hypothesis.  In the retained compact case,
\cref{paper7:cor:retained-stratum-cost-exclusion} applies.
\end{proof}

\begin{definition}[Adjacent exits from cost-divergence exclusion]
\label{paper7:def:cost-exclusion-missing-hypotheses}
If the active supercritical scaling condition and validated compact
cost-divergence conclusion are present but the branch is not retained in the
compact cost-divergence stratum, then the first exit from the cost-exclusion
argument lies in the ordered list:
\[
  \begin{gathered}
  \text{class membership},\quad
  \text{cost compatibility},\quad
  \text{stability of the active scaling condition},\\
  \text{stability of the validated cost-divergence conclusion},\quad
  \text{local monotonicity-carrier validation}.
  \end{gathered}
\]
On the retained compact stratum none of these exits occurs.  The final item is
expanded by \cref{paper7:thm:carrier-validation-retained-stratum} into the
carrier-subidentity, fixed-finite-error, moving-annulus, and
carrier-comparison adjacent strata.
\end{definition}

\begin{lemma}[Carrier exits are adjacent strata, not retained compact branches]
\label{paper7:lem:missing-hypotheses-not-alternatives}
Assume a renormalized Type II solution already has the active supercritical
scaling condition and the validated compact Type II cost-divergence conclusion.
If one of the exits in \cref{paper7:def:cost-exclusion-missing-hypotheses}
occurs, then the branch is not in the retained compact cost-divergence stratum.
If no such exit occurs, the local cost-to-exclusion data are available.
\end{lemma}

\begin{proof}
The retained compact cost-divergence stratum is reached only after the ordered
tests of \cref{paper7:lem:no-unrecorded-cost-escape} have passed.  Hence a
failure of class membership, cost compatibility, active-scaling stability,
cost-divergence stability, or carrier validation is by definition an exit to an
adjacent stratum.  Conversely, if no such exit occurs, the carrier-validation
test has passed.  By \cref{paper7:thm:carrier-validation-retained-stratum}, the
validated compact cost admits local monotonicity-carrier data, and then
\cref{paper7:lem:monotonicity-carrier-gives-cost-to-exclusion-data} gives the
local cost-to-exclusion data.
\end{proof}

\begin{definition}[Pre-carrier validated compact-cost regime]
\label{paper7:def:precarrier-validated-compact-channel}
A renormalized Type II branch lies in the pre-carrier validated compact-cost
regime if, on a terminal tail, it has active supercritical scaling and the
validated compact Type II cost-divergence conclusion, and the ordered validation
tests of \cref{paper7:lem:no-unrecorded-cost-escape} have passed through the
following stages:
\begin{enumerate}[label=\textup{(\roman*)}]
\item the Type II analytic data and ordered exhaustion tests;
\item raw-chart selection, repaired-gauge representation, pressure
reconstruction, modulation data, and scale-admissibility of the repaired chart;
\item the local compactness decomposition without entering a multibubble,
multicore, cascade, exterior-tightness, local \(H^1\)-loss,
finite-cost scale-collapse, or scale-rigid residual alternative;
\item the retained finite critical-mass window, namely the compact-window
finite critical-mass input supplied by
\cref{paper6a:lem:terminal-local-critical-mass-routing};
\item critical tightness at the retained scale and local windowed \(H^1\)
control;
\item local modulation integrability, local cost well-posedness, and
compatibility of the chosen compact Type II cost with the validated
Navier--Stokes cost.
\end{enumerate}
This regime deliberately stops immediately before local monotonicity-carrier
validation.  Thus it contains all compact-cost interface data needed to ask for
the fixed-cutoff or moving-cutoff carrier estimates, but it does not assert
that those estimates have already been proved.
The finite critical-mass row is read locally in the final no-Type-II assembly:
failure of compact-window lower or upper critical mass is routed by the local
terminal critical-mass and exterior/state-space alternatives.
\end{definition}

\begin{definition}[Fixed and moving carrier subidentities]
\label{paper7:def:carrier-subidentities}
For the fixed cutoff \(\phi_{R_0}\), the fixed carrier subidentity is the
distributional estimate
\[
  \frac{d}{d\tau}\mathcal E_\phi(\tau)
  +\frac12\tilde{\mathfrak D}_{R_0}(\tau)
  \le B_{R_0}(\tau),
  \qquad
  \mathcal E_\phi(\tau)=\frac12\int |V|^2\phi_{R_0}.
\]
For a moving cutoff \(\Phi(\tau,y)=\phi(y/R(\tau))\), the moving carrier
subidentity is the distributional estimate
\[
  \frac{d}{d\tau}\mathcal E_R(\tau)
  +\frac12\tilde{\mathfrak D}_{R(\tau)}(\tau)
  \le B_{\mathrm{mov}}(\tau),
  \qquad
  \mathcal E_R(\tau)=\frac12\int |V|^2\Phi .
\]
Here \(B_{R_0}\) and \(B_{\mathrm{mov}}\) are the absolute error densities of
\cref{paper6a:def:fixed-error-density,paper6a:def:moving-error-density}.  The
subidentity is the one-sided estimate needed for corrected monotonicity; an
exact localized energy identity is sufficient but not necessary.
\end{definition}

\begin{lemma}[Suitable local energy gives carrier subidentities]
\label{paper7:lem:suitable-energy-gives-carrier-subidentities}
Assume the repaired-gauge representation of a pre-carrier branch includes the
local suitable energy inequality on compact renormalized cylinders.  Then the
fixed carrier subidentity holds for every nonnegative compactly supported fixed
cutoff.  The moving carrier subidentity holds for every nonnegative
finite-valued moving cutoff \(\Phi(\tau,y)=\phi(y/R(\tau))\) with
\(R\in AC_{\mathrm{loc}}\), after the usual approximation of \(R\) by smooth
positive radii on compact \(\tau\)-intervals.
\end{lemma}

\begin{proof}
Apply the local suitable energy inequality in the physical variables to the
pullback of the nonnegative renormalized test function.  The repaired-gauge
change of variables is locally absolutely continuous in time and smooth in
space on compact cylinders, so the distributional chain rule gives the
renormalized local energy inequality with the drift terms
\(a(V+y\cdot\nabla V)\) and \(b\cdot\nabla V\).  For a fixed cutoff this is
the fixed localized energy identity with equality replaced by \(\le\).  Adding
\(\frac12\tilde{\mathfrak D}_{R_0}\) to both sides and bounding the cutoff,
pressure, translation, and scale terms by their absolute values gives the fixed
carrier subidentity.
\par
For a moving cutoff, use smooth positive approximations \(R_k\to R\) in
\(AC_{\mathrm{loc}}\) on each compact time interval.  The local energy
inequality applies to \(\Phi_k(\tau,y)=\phi(y/R_k(\tau))\), and the preceding
argument gives the moving subidentity with \(B_{\mathrm{mov},k}\).  Passing to
the limit in distributions uses the local finite-energy and local \(H^1\)
bounds already present on the pre-carrier regime.  The term
\(\partial_\tau\Phi\) is interpreted as the \(L^1_{\mathrm{loc}}\) derivative
coming from \(R\in AC_{\mathrm{loc}}\).  The limit is exactly the moving
carrier subidentity.
\end{proof}

\begin{definition}[Fixed-cutoff localized-energy closure criterion]
\label{paper7:def:fixed-carrier-closure-package}
A branch in the pre-carrier validated compact-cost regime satisfies the
fixed-cutoff localized-energy closure criterion if, for the cutoff \(\phi_{R_0}\) used to
evaluate the validated compact cost and for
\[
  \mathcal E_\phi(\tau)=\frac12\int |V(y,\tau)|^2\phi_{R_0}(y)\,dy,
\]
the following local analytic statements hold on a terminal tail:
\begin{enumerate}[label=\textup{(\roman*)}]
\item \(\mathcal E_\phi(\tau_0)<\infty\);
\item the fixed carrier subidentity of
\cref{paper7:def:carrier-subidentities} is valid for the repaired-gauge branch;
\item the fixed-cutoff error density \(B_{R_0}\) of
\cref{paper6a:def:fixed-error-density} has finite tail integral,
\[
  \int_{\tau_0}^{\infty}B_{R_0}(\tau)\,d\tau<\infty;
\]
\item the validated compact cost is energy-controlled by
\(\tilde{\mathfrak D}_{R_0}\) in the sense of
\cref{paper7:def:energy-controlled-validated-cost}.
\end{enumerate}
\end{definition}

\begin{definition}[Moving-cutoff localized-energy closure criterion]
\label{paper7:def:moving-carrier-closure-package}
A branch in the pre-carrier validated compact-cost regime satisfies the
moving-cutoff localized-energy closure criterion if there is a finite-valued
nondecreasing locally absolutely continuous radius \(R(\tau)\ge R_0\), with
\(R(\tau)\to\infty\) and moving cutoff
\(\Phi(\tau,y)=\phi(y/R(\tau))\), such that, on a terminal tail:
\begin{enumerate}[label=\textup{(\roman*)}]
\item
\[
  \mathcal E_R(\tau_0)
  =
  \frac12\int |V(y,\tau_0)|^2\Phi(\tau_0,y)\,dy
  <\infty;
\]
\item the moving carrier subidentity of
\cref{paper7:def:carrier-subidentities} is valid with the time-dependent cutoff
\(\Phi\);
\item the moving-annulus errors are summable in the sense of
\cref{paper6a:def:summable-moving-errors}, which implies
\[
  \int_{\tau_0}^{\infty}B_{\mathrm{mov}}(\tau)\,d\tau<\infty
\]
by \cref{paper6a:lem:summable-moving-errors};
\item the validated compact cost is controlled by the moving identity cost on
the terminal tail: for some \(C_m<\infty\),
\[
  0\le \mathfrak C_{\rm II}^{NS}(\tau)
  \le C_m\tilde{\mathfrak D}_{R(\tau)}(\tau)
  \qquad\text{for a.e. }\tau\ge\tau_0 .
\]
\end{enumerate}
\end{definition}

\begin{theorem}[Non-carrier adjacent exits are identified before carrier validation]
\label{paper7:thm:noncarrier-exits-route}
Let \(\omega\) be an admissible local Type II branch with active supercritical
scaling and validated compact Type II cost divergence.  Apply the ordered tests
of \cref{paper7:lem:no-unrecorded-cost-escape}.  If a first failed test exists
and is not one of the carrier-validation failures
\[
  \text{carrier subidentity},\quad
  \text{fixed finite error},\quad
  \text{moving-annulus summability},\quad
  \text{carrier comparison},
\]
then the branch has already failed to enter the pre-carrier validated
compact-cost regime of
\cref{paper7:def:precarrier-validated-compact-channel}.  More precisely, its
first failure is one of the following named non-carrier strata:
\[
\begin{gathered}
\text{ordered exhaustion failure, representation failure, repaired-gauge or
scale-degeneracy failure},\\
\text{multibubble or multicore splitting, cascade, exterior-tightness failure,
local \(H^1\)-loss},\\
\text{finite-cost scale-collapse, scale-rigid residual behavior,
critical-mass failure},\\
\text{modulation-integrability failure, cost-well-posedness failure, or
cost-compatibility failure}.
\end{gathered}
\]
In any ambient class in which the corresponding local closure theorems for
these named strata are imposed, such a branch is not a remaining compact-cost
branch.
\end{theorem}

\begin{proof}
The validation procedure in
\cref{paper7:lem:no-unrecorded-cost-escape} is ordered.  The
pre-carrier regime is the set of branches for which all tests before
local monotonicity-carrier validation have passed.  Therefore any first failure
before the carrier-validation row prevents membership in
\cref{paper7:def:precarrier-validated-compact-channel}.
\par
The same lemma lists every possible first failure before carrier validation:
the exhaustion failures are those of
\cref{def:c1-exhaustion-failures}; representation and repaired-gauge failures
are those of
\cref{def:c2-representation-failures,def:c2r-minimal-failures} together with
scale-degeneracy from
\cref{paper7:prop:active-scaling-stability-package}; the compactness failures
are the alternatives in
\cref{paper6:thm:state-decomposition}; critical-mass failures are those routed
by \cref{paper6a:lem:terminal-local-critical-mass-routing}; compact tightness
and local \(H^1\)-loss are the tests of \cref{def:c5-compact-tests}; and the
modulation, cost-well-posedness, and compatibility failures are exactly the
failures recorded in
\cref{paper7:lem:local-modulation-integrability,paper7:cor:compact-window-data-cost-well-posed,paper7:lem:canonical-cost-compatibility,paper7:lem:comparable-cost-compatibility}.
Thus the displayed list is exhaustive for non-carrier first failures.
\par
If the surrounding Type II class includes the local closure theorem assigned to
the named stratum, then the branch is handled there.  It is therefore not a
branch remaining to be treated by the retained compact-cost localized-energy argument.
\end{proof}

\begin{theorem}[Pre-carrier reduction to carrier validation]
\label{paper7:thm:precarrier-reduces-to-carrier}
Let \(\omega\) be in the pre-carrier validated compact-cost regime.  Then
exactly one of the following alternatives occurs:
\begin{enumerate}[label=\textup{(\roman*)}]
\item the validated compact cost admits local monotonicity-carrier data in the
sense of \cref{paper7:def:local-monotonicity-carrier};
\item the branch exits the retained compact cost-divergence stratum through the
first failed carrier-validation test: carrier-subidentity failure,
fixed-cutoff finite-error failure, moving-annulus summability failure, or
carrier-comparison failure for the validated cost.
\end{enumerate}
If the first alternative occurs, \(\omega\) is excluded by the retained
compact-cost argument.
\end{theorem}

\begin{proof}
By definition of the pre-carrier regime, every ordered validation test before
carrier validation has passed, and the validated compact cost-divergence
conclusion is already present.  Thus the next unresolved row in
\cref{paper7:lem:no-unrecorded-cost-escape} is precisely the local
monotonicity-carrier row.  Applying
\cref{paper7:thm:carrier-validation-retained-stratum} gives the two
alternatives listed in the statement, and the first-failure convention makes
them disjoint.
\par
If the carrier exists, then
\cref{paper7:lem:monotonicity-carrier-gives-cost-to-exclusion-data} gives the
local cost-to-exclusion data, and
\cref{paper7:thm:local-cost-to-exclusion} excludes the branch.  Equivalently,
the same conclusion is
\cref{paper7:cor:retained-stratum-cost-exclusion}.
\end{proof}

\begin{theorem}[Fixed-cutoff carrier closure]
\label{paper7:thm:fixed-carrier-closure}
Let \(\omega\) be in the pre-carrier validated compact-cost regime.  If
\(\omega\) satisfies the fixed-cutoff localized-energy closure criterion of
\cref{paper7:def:fixed-carrier-closure-package}, then \(\omega\) is excluded
from the Type II class under consideration.
\end{theorem}

\begin{proof}
The pre-carrier regime supplies active supercritical scaling, local
well-posedness and compatibility of the validated compact cost, and the
validated compact cost-divergence conclusion.  The fixed-cutoff localized-energy criterion
supplies finite initial localized energy, the fixed carrier subidentity, finite
fixed-cutoff error, and energy-control comparison of the validated cost with
\(\tilde{\mathfrak D}_{R_0}\).  Let \(C_1,\tau_1\) be the constants in that
comparison and increase the tail initial time so that the comparison holds from
\(\tau_0\) onward.  Define
\[
  \mathcal H_\omega(\tau)
  =
  \mathcal E_\phi(\tau)+\int_\tau^\infty B_{R_0}(s)\,ds .
\]
The finite-error condition makes \(\mathcal H_\omega\) locally absolutely
continuous and bounded below.  The fixed carrier subidentity gives
\[
  \frac{d}{d\tau}\mathcal H_\omega(\tau)
  +\frac12\tilde{\mathfrak D}_{R_0}(\tau)\le0
\]
in distributions.  Since
\(\mathfrak C_{\rm II}^{NS}\le C_1\tilde{\mathfrak D}_{R_0}\), this implies
\[
  \frac{d}{d\tau}\mathcal H_\omega(\tau)
  +\frac{1}{2C_1}\mathfrak C_{\rm II}^{NS}(\tau)\le0 .
\]
Together with the validated compact cost-divergence conclusion, these are the
local cost-to-exclusion data with
\(\mathcal D_\omega=\mathfrak C_{\rm II}^{NS}\).  Equivalently, the fixed
criterion is a successful carrier-validation row.  Since all pre-carrier tests
and the validated compact cost-divergence conclusion are already present,
\(\omega\) is on the retained compact cost-divergence stratum after carrier
validation.  Applying \cref{paper7:thm:local-cost-to-exclusion} gives the Type
II exclusion conclusion.
\end{proof}

\begin{theorem}[Moving-cutoff carrier closure]
\label{paper7:thm:moving-carrier-closure}
Let \(\omega\) be in the pre-carrier validated compact-cost regime.  If
\(\omega\) satisfies the moving-cutoff localized-energy closure criterion of
\cref{paper7:def:moving-carrier-closure-package}, then \(\omega\) is excluded
from the Type II class under consideration.
\end{theorem}

\begin{proof}
The pre-carrier regime supplies active supercritical scaling, local
well-posedness and compatibility of the validated compact cost, and the
validated compact cost-divergence conclusion.  The moving-cutoff carrier
criterion supplies the finite initial moving localized energy, the
moving carrier subidentity, summable moving-annulus errors, and
the terminal comparison
\[
  0\le \mathfrak C_{\rm II}^{NS}(\tau)
  \le C_m\tilde{\mathfrak D}_{R(\tau)}(\tau).
\]
By \cref{paper6a:lem:summable-moving-errors}, the summability condition gives
\(B_{\mathrm{mov}}\in L^1(\tau_0,\infty)\).  Define
\[
  \mathcal H_\omega(\tau)
  =
  \mathcal E_R(\tau)+\int_\tau^\infty B_{\mathrm{mov}}(s)\,ds .
\]
The finite moving-error condition makes \(\mathcal H_\omega\) locally
absolutely continuous and bounded below.  The moving carrier subidentity gives
\[
  \frac{d}{d\tau}\mathcal H_\omega(\tau)
  +\frac12\tilde{\mathfrak D}_{R(\tau)}(\tau)\le0
\]
in distributions.  The moving comparison gives
\[
  \frac{d}{d\tau}\mathcal H_\omega(\tau)
  +\frac{1}{2C_m}\mathfrak C_{\rm II}^{NS}(\tau)\le0 .
\]
Together with the validated compact cost-divergence conclusion, these are the
local cost-to-exclusion data with
\(\mathcal D_\omega=\mathfrak C_{\rm II}^{NS}\).  Equivalently, the moving
criterion is a successful carrier-validation row.  Since all pre-carrier tests
and the validated compact cost-divergence conclusion are already present,
\(\omega\) is on the retained compact cost-divergence stratum after carrier
validation.  Applying \cref{paper7:thm:local-cost-to-exclusion} excludes the
branch.
\end{proof}

\begin{lemma}[Pre-carrier local energy gives finite initial localized energy]
\label{paper7:lem:precarrier-finite-local-energy}
Let \(\omega\) be in the pre-carrier validated compact-cost regime, represented
by \(V\) on a terminal tail.  Fix any later time \(T\ge\tau_0\).  After
possibly replacing the tail by one whose initial time
\(\tau_1\in[T,T+1]\) is a compact-cylinder local-energy Lebesgue time, every
finite \(R<\infty\) satisfies
\[
  \int_{B_R}|V(y,\tau_1)|^2\,dy<\infty .
\]
Consequently the initial localized-energy clauses in the fixed and moving
localized-energy closure criteria are available after a harmless terminal
restriction for every compactly supported fixed cutoff and every moving cutoff
with finite initial radius.  This is a local compact-cylinder statement; it does
not require \(V(\tau)\in L^3(\mathbb R^3)\) or a whole-space critical bound.
\end{lemma}

\begin{proof}
The physical branch is suitable and the repaired-gauge chart is locally
absolutely continuous on compact renormalized cylinders in the pre-carrier
regime.  Pulling back the local energy class gives
\[
  V\in L^\infty_{\rm loc}([\tau_0,\infty);L^2(B_R))
  \quad\text{for every }R<\infty
\]
after choosing the usual representative, or at minimum
\[
  V\in L^2_{\rm loc}([\tau_0,\infty);L^2(B_R))
\]
with local-energy Lebesgue times on every compact interval.  If the stronger
representative is unavailable at a prescribed endpoint, choose
\(\tau_1\in[T,T+1]\) which is a Lebesgue time simultaneously for the finitely
many compact supports used by the fixed or moving cutoff under consideration.
This is possible because the local \(L^2\)-in-time norms are finite on
\([T,T+1]\).
\par
At such a \(\tau_1\), every compactly supported fixed cutoff has finite
localized energy.  If \(\Phi(\tau,y)=\phi(y/R(\tau))\) and
\(R(\tau_1)<\infty\), then \(\Phi(\tau_1,\cdot)\) is compactly supported, so
the same local \(L^2\) finiteness gives \(\mathcal E_R(\tau_1)<\infty\).
Passing from \(\tau_0\) to \(\tau_1\) removes only a finite initial
renormalized interval and therefore does not change any terminal divergence or
tail-integrability alternative.
\end{proof}

\begin{lemma}[Nested moving cutoffs give carrier comparison]
\label{paper7:lem:nested-moving-carrier-comparison}
Let \(\omega\) be in the pre-carrier validated compact-cost regime, and
assume the validated compact cost is energy-controlled by the fixed identity
cost:
\[
  0\le \mathfrak C_{\rm II}^{NS}(\tau)
  \le C_1\tilde{\mathfrak D}_{R_0}(\tau)
  \qquad\text{for a.e. }\tau\ge\tau_1 .
\]
Let \(\Phi(\tau,y)\) be a moving cutoff such that, after increasing the initial
time if necessary,
\[
  \Phi(\tau,y)\ge \phi_{R_0}(y)
  \qquad\text{for all }y\in\mathbb R^3
  \text{ and a.e. }\tau\ge\tau_1 .
\]
Then the validated compact cost is controlled by the moving identity cost:
\[
  0\le \mathfrak C_{\rm II}^{NS}(\tau)
  \le C_1\tilde{\mathfrak D}_{R(\tau)}(\tau)
  \qquad\text{for a.e. }\tau\ge\tau_1 .
\]
In particular, for the canonical identity cost the carrier comparison is
automatic with \(C_1=1\).
\end{lemma}

\begin{proof}
Because \(\Phi(\tau,\cdot)\ge\phi_{R_0}\) and both cutoffs are nonnegative,
\[
  \nu\int |\nabla V|^2\phi_{R_0}
  \le
  \nu\int |\nabla V|^2\Phi(\tau,\cdot),
\]
and, since \(a_+(\tau)\ge0\),
\[
  a_+(\tau)\int |V|^2\phi_{R_0}
  \le
  a_+(\tau)\int |V|^2\Phi(\tau,\cdot).
\]
Adding these two inequalities gives
\[
  \tilde{\mathfrak D}_{R_0}(\tau)
  \le
  \tilde{\mathfrak D}_{R(\tau)}(\tau)
  \qquad\text{for a.e. }\tau\ge\tau_1 .
\]
Combining this with the fixed energy-control comparison yields the displayed
moving comparison.  The canonical identity cost is energy-controlled by
\cref{paper7:def:energy-controlled-validated-cost} with \(C_1=1\).
\end{proof}

\begin{definition}[Finite-error carrier selection]
\label{paper7:def:finite-error-carrier-selection}
The pre-carrier validated compact-cost regime has finite-error carrier
selection if every branch in that regime whose validated compact cost is
energy-controlled by the fixed identity cost satisfies, after passing to a
later terminal tail if necessary, at least one of the following two local
alternatives:
\begin{enumerate}[label=\textup{(\roman*)}]
\item the fixed carrier subidentity of
\cref{paper7:def:carrier-subidentities} is valid and
\[
  \int_{\tau_0}^{\infty}B_{R_0}(\tau)\,d\tau<\infty ;
\]
\item there exists a finite-valued nondecreasing locally absolutely continuous
radius \(R(\tau)\ge R_0\), with \(R(\tau)\to\infty\) and moving cutoff
\(\Phi(\tau,y)=\phi(y/R(\tau))\), such that
\[
  \Phi(\tau,y)\ge\phi_{R_0}(y)
  \qquad\text{for all }y\text{ and a.e. }\tau,
\]
the moving carrier subidentity of
\cref{paper7:def:carrier-subidentities} is valid with the time-dependent cutoff
\(\Phi\), and the moving-annulus errors are summable in the sense of
\cref{paper6a:def:summable-moving-errors}.
\end{enumerate}
This condition contains only the local one-sided energy estimate and finite-error
selection statements.  It does not include finite initial localized energy or
carrier comparison, because those are supplied by
\cref{paper7:lem:precarrier-finite-local-energy,paper7:lem:nested-moving-carrier-comparison}.
\end{definition}

\begin{definition}[Annular finite-error selection]
\label{paper7:def:annular-finite-error-selection}
The pre-carrier validated compact-cost regime has annular finite-error
selection if every branch in that regime whose validated compact cost is
energy-controlled by the fixed identity cost satisfies, after passing to a
later terminal tail if necessary, at least one of the following two alternatives:
\begin{enumerate}[label=\textup{(\roman*)}]
\item the fixed error density has finite tail integral,
\[
  \int_{\tau_0}^{\infty}B_{R_0}(\tau)\,d\tau<\infty;
\]
\item there exists a finite-valued nondecreasing locally absolutely continuous
nested radius \(R(\tau)\ge R_0\), with \(R(\tau)\to\infty\) and
\(\Phi(\tau,y)=\phi(y/R(\tau))\ge\phi_{R_0}(y)\), such that the moving-annulus
errors are summable in the sense of
\cref{paper6a:def:summable-moving-errors}.
\end{enumerate}
This is the purely finite-error part of the carrier problem.  It does not
include local energy subidentities, finite initial energy, or carrier comparison.
\end{definition}

\begin{definition}[Componentwise fixed-or-moving finite-error estimates]
\label{paper7:def:componentwise-annular-finite-error}
For a branch in the pre-carrier validated compact-cost regime, the
componentwise annular finite-error estimates hold if, after passing to a later
terminal tail if necessary, either the six fixed-cutoff components
\[
\begin{gathered}
  I_1=\left|\int |V|^2\Delta\phi_{R_0}\right|,\qquad
  I_2=\left|\int |V|^2V\cdot\nabla\phi_{R_0}\right|,\qquad
  I_3=\left|\int P\,V\cdot\nabla\phi_{R_0}\right|,\\
  I_4=|b|\left|\int |V|^2\nabla\phi_{R_0}\right|,\qquad
  I_5=a_-M_\phi,\qquad
  I_6=a_+\int |V|^2\bigl(-y\cdot\nabla\phi_{R_0}\bigr)
\end{gathered}
\]
belong to \(L^1(\tau_0,\infty)\), or there is a finite-valued nondecreasing
locally absolutely continuous nested
radius \(R(\tau)\ge R_0\), with \(R(\tau)\to\infty\) and
\(\Phi(\tau,y)=\phi(y/R(\tau))\ge\phi_{R_0}(y)\), such that each of the seven
moving annular components
\[
\begin{gathered}
  J_1=\left|\int |V|^2\partial_\tau\Phi\right|,\qquad
  J_2=\left|\int |V|^2\Delta\Phi\right|,\qquad
  J_3=\left|\int |V|^2V\cdot\nabla\Phi\right|,\\
  J_4=\left|\int P\,V\cdot\nabla\Phi\right|,\qquad
  J_5=|b|\left|\int |V|^2\nabla\Phi\right|,\\
  J_6=a_-\int |V|^2\Phi,\qquad
  J_7=a_+\int |V|^2\bigl(-y\cdot\nabla\Phi\bigr)
\end{gathered}
\]
belongs to \(L^1(\tau_0,\infty)\).
\end{definition}

\begin{remark}[What the componentwise estimates isolate]
\label{paper7:rem:componentwise-annular-residual}
The quantities in
\cref{paper7:def:componentwise-annular-finite-error} are exactly the fixed and
moving carrier-error components: in the fixed-cutoff case, viscous cutoff,
convection, pressure-tail, center drift, negative scale drift, and the
positive scale-shell drift; in the moving-cutoff case, cutoff transition, viscous
cutoff, convection, pressure-tail, center drift, negative scale drift, and the
positive scale-shell drift.  These diagnostics need not be closed by a single
global tail estimate.  In the stratified use of Section~7, persistent failure
of the cutoff, viscous, convection, or pressure diagnostics is read as shell
mass that cannot be evacuated from selected windows and is therefore assigned to
exterior leakage, exterior concentration, loss of critical tightness, or local
windowed \(H^1\)-loss.  Persistent failure of the translation diagnostic is
read as a modulation-driven adjacent branch, while persistent failure of the
positive scale-shell term is read as a scale-dynamics branch.  Thus the
component list records the local diagnostics seen by the carrier test; it is
not intended to force the proof into a single classical PDE closure argument.
For radial nonincreasing cutoffs, the scale-shell term
\[
  -\frac{a}{2}\int |V|^2 y\cdot\nabla\Phi
\]
has only an \(a_+\)-contribution in the residual carrier error.  The \(a_-\)
part is favorable and is already absorbed into the negative-drift alternative.
The negative scale-drift component \(J_6\) is not to be replaced by an
unproved equivalence with compact-cost divergence: if \(J_6\) is not summable,
it must be handled through the explicit scale-drift alternatives of
\cref{paper6a:cor:scale-drift-alternatives,paper6a:thm:moving-cutoff-scale-collapse-alternative}
or treated by a direct summability estimate.
\end{remark}

\begin{lemma}[Componentwise estimates give annular finite-error selection]
\label{paper7:lem:componentwise-gives-annular-finite-error}
If every branch in the pre-carrier validated compact-cost regime whose
validated compact cost is energy-controlled by the fixed identity cost
satisfies the componentwise annular finite-error estimates of
\cref{paper7:def:componentwise-annular-finite-error}, then the regime has
annular finite-error selection in the sense of
\cref{paper7:def:annular-finite-error-selection}.
\end{lemma}

\begin{proof}
Let \(\omega\) be such a branch.  If the fixed alternative in
\cref{paper7:def:componentwise-annular-finite-error} holds, with
\(I_1,\ldots,I_6\) denoting the fixed components appearing in that definition,
then
\cref{paper6a:def:fixed-error-density} gives
\[
  B_{R_0}
  =
  \frac{\nu}{2}I_1+\frac12 I_2+I_3+\frac12 I_4+\frac12 I_5+\frac12 I_6 .
\]
Since \(B_{R_0}\) is a finite positive linear combination of integrable fixed
components, \(B_{R_0}\in L^1(\tau_0,\infty)\).  This is the first alternative
in \cref{paper7:def:annular-finite-error-selection}.
Otherwise the definition supplies a finite-valued nested moving radius
\(R(\tau)\ge R_0\), \(R(\tau)\to\infty\), for which the seven moving annular
components \(J_1,\ldots,J_7\) are integrable.  These seven components are
exactly the summability requirements in
\cref{paper6a:def:summable-moving-errors}, with the same moving cutoff
\(\Phi\).  Hence the moving-annulus errors are summable, which is the second
alternative in \cref{paper7:def:annular-finite-error-selection}.
\end{proof}

\begin{theorem}[Finite-error selection reduced to annular finite-error]
\label{paper7:thm:finite-error-selection-reduced-to-annular}
Assume that repaired-gauge suitable local energy gives the fixed and moving
carrier subidentities in the sense of
\cref{paper7:lem:suitable-energy-gives-carrier-subidentities}.  If the
pre-carrier regime has annular finite-error selection in the sense of
\cref{paper7:def:annular-finite-error-selection}, then it has finite-error
carrier selection in the sense of
\cref{paper7:def:finite-error-carrier-selection}.
\end{theorem}

\begin{proof}
Let \(\omega\) be a branch in the pre-carrier regime whose validated compact
cost is energy-controlled by the fixed identity cost.  By annular finite-error
selection, after restricting to a later tail either the fixed error density is
integrable or there is a nested moving radius for which the moving-annulus
errors are summable.  The repaired-gauge suitable local energy lemma supplies
the fixed carrier subidentity in the first case and the moving carrier
subidentity in the second.  These are exactly the two alternatives in
\cref{paper7:def:finite-error-carrier-selection}.
\end{proof}

\begin{theorem}[Carrier completion reduced to finite-error selection]
\label{paper7:thm:carrier-completion-reduced-to-finite-error}
Assume the pre-carrier validated compact-cost regime has finite-error carrier
selection in the sense of
\cref{paper7:def:finite-error-carrier-selection}.  Then every branch in the
pre-carrier regime whose validated compact cost is energy-controlled by the
fixed identity cost satisfies the fixed-cutoff or moving-cutoff carrier closure
criterion.  Consequently that branch is excluded from the Type II class under
consideration.
\end{theorem}

\begin{proof}
Let \(\omega\) be such a branch.  The finite initial localized energy required
by either localized-energy criterion is supplied locally, after a harmless
terminal restriction, by
\cref{paper7:lem:precarrier-finite-local-energy}.  The assumed energy-control
comparison gives the fixed carrier-comparison clause.
\par
If the fixed alternative in
\cref{paper7:def:finite-error-carrier-selection} holds, then the fixed
carrier subidentity and the finite \(B_{R_0}\)-tail are precisely the two
remaining clauses of the fixed-cutoff localized-energy closure criterion.  Hence
\cref{paper7:thm:fixed-carrier-closure} excludes the branch.
\par
If the moving alternative holds, the summable moving-annulus condition gives
\[
  \int_{\tau_0}^{\infty}B_{\mathrm{mov}}(\tau)\,d\tau<\infty
\]
by \cref{paper6a:lem:summable-moving-errors}.  The moving carrier subidentity
is part of the moving alternative.  The nested cutoff condition and fixed
energy-control comparison give the moving carrier comparison by
\cref{paper7:lem:nested-moving-carrier-comparison}.  Together with finite
initial moving energy from
\cref{paper7:lem:precarrier-finite-local-energy}, these are exactly the moving
localized-energy closure criterion.  Therefore
\cref{paper7:thm:moving-carrier-closure} excludes the branch.
\par
The two alternatives in finite-error carrier selection exhaust the branch, so
the claimed fixed-or-moving carrier completion and exclusion follow.
\end{proof}

\begin{theorem}[Localized-energy criterion closes the pre-carrier regime]
\label{paper7:thm:carrier-package-closes-precarrier}
Let \(\omega\) be in the pre-carrier validated compact-cost regime.  If
\(\omega\) satisfies either the fixed-cutoff localized-energy closure criterion of
\cref{paper7:def:fixed-carrier-closure-package} or the moving-cutoff
localized-energy closure criterion of \cref{paper7:def:moving-carrier-closure-package}, then
\(\omega\) is excluded from the Type II class under consideration.
\end{theorem}

\begin{proof}
In the fixed-cutoff case apply
\cref{paper7:thm:fixed-carrier-closure}.  In the moving-cutoff case apply
\cref{paper7:thm:moving-carrier-closure}.  These are the only cases in the
stated disjunction.
\end{proof}

\begin{definition}[Carrier-estimate completion for the compact-cost regime]
\label{paper7:def:carrier-estimate-completion}
The compact-cost regime has carrier-estimate completion if every branch in the
pre-carrier validated compact-cost regime satisfies, after passing to a later
terminal tail if necessary, either the fixed-cutoff localized-energy closure criterion of
\cref{paper7:def:fixed-carrier-closure-package} or the moving-cutoff
localized-energy closure criterion of \cref{paper7:def:moving-carrier-closure-package}.
\end{definition}

\begin{corollary}[Finite-error selection gives carrier-estimate completion]
\label{paper7:cor:finite-error-selection-gives-carrier-completion}
Assume every branch in the pre-carrier validated compact-cost regime uses a
validated compact cost that is energy-controlled by the fixed identity cost; in
particular this holds for the canonical identity cost.  If finite-error carrier
selection holds in the sense of
\cref{paper7:def:finite-error-carrier-selection}, then the compact-cost regime
has carrier-estimate completion in the sense of
\cref{paper7:def:carrier-estimate-completion}.
\end{corollary}

\begin{proof}
Let \(\omega\) be any branch in the pre-carrier validated compact-cost regime.
By the energy-control hypothesis, \(\omega\) is one of the branches covered by
\cref{paper7:thm:carrier-completion-reduced-to-finite-error}.  That theorem
gives either the fixed-cutoff localized-energy closure criterion or the moving-cutoff
localized-energy closure criterion on a terminal tail.  This is exactly
\cref{paper7:def:carrier-estimate-completion}.  The parenthetical canonical
case follows from
\cref{paper7:def:energy-controlled-validated-cost}.
\end{proof}

\begin{corollary}[Annular finite-error selection gives carrier-estimate completion]
\label{paper7:cor:annular-finite-error-selection-gives-carrier-completion}
Assume every branch in the pre-carrier validated compact-cost regime uses a
validated compact cost that is energy-controlled by the fixed identity cost.
Assume also that repaired-gauge suitable local energy gives the fixed and
moving carrier subidentities in the sense of
\cref{paper7:lem:suitable-energy-gives-carrier-subidentities}.  If annular
finite-error selection holds in the sense of
\cref{paper7:def:annular-finite-error-selection}, then the compact-cost regime
has carrier-estimate completion.
\end{corollary}

\begin{proof}
\Cref{paper7:thm:finite-error-selection-reduced-to-annular} converts annular
finite-error selection into finite-error carrier selection.  Then
\cref{paper7:cor:finite-error-selection-gives-carrier-completion} gives
carrier-estimate completion.
\end{proof}

\begin{theorem}[Adjacent closure reduced to carrier-estimate completion]
\label{paper7:thm:adjacent-closure-reduced-to-carrier-completion}
Assume the non-carrier adjacent strata listed in
\cref{paper7:thm:noncarrier-exits-route} are closed by their corresponding
local compactness or interface arguments in the ambient Type II class.  If the
compact-cost regime has carrier-estimate completion in the sense of
\cref{paper7:def:carrier-estimate-completion}, then the compact-cost regime
has adjacent-stratum closure in the sense of
\cref{paper7:def:compact-cost-adjacent-closure}.
\end{theorem}

\begin{proof}
Let \(\omega\) be a branch with active supercritical scaling and validated
compact Type II cost divergence that does not remain in the retained compact
cost-divergence stratum.  By
\cref{paper7:thm:classwise-retained-compact-coverage}, \(\omega\) enters one
of the named adjacent strata produced by
\cref{paper7:lem:no-unrecorded-cost-escape}.
\par
If the adjacent stratum is non-carrier, then
\cref{paper7:thm:noncarrier-exits-route} identifies it with one of the
previously named exhaustion, representation, compactness, critical-mass,
compact-test, modulation, cost-well-posedness, or cost-compatibility strata.
By the first hypothesis of the present theorem, that stratum is closed by its
own local argument, so the Type II exclusion conclusion holds there, or the
branch is assigned to a previously closed non-compact-cost alternative.
\par
It remains to consider a branch that reaches the pre-carrier regime.  By
carrier-estimate completion, after restriction to a later tail if needed, it
satisfies either the fixed-cutoff or moving-cutoff localized-energy closure criterion.  In
the fixed or moving case,
\cref{paper7:thm:carrier-package-closes-precarrier} gives the Type II
exclusion conclusion.  Thus every named adjacent stratum arising from the
compact-cost regime has a local exclusion or an assignment to a previously closed
alternative, which is exactly
\cref{paper7:def:compact-cost-adjacent-closure}.
\end{proof}

\begin{definition}[Stratified moving-carrier alternative]
\label{paper7:def:stratified-moving-carrier-routing}
Assume the validated compact cost is energy-controlled by the fixed identity
cost.  The pre-carrier validated compact-cost regime has stratified
moving-cutoff alternative if every branch \(\omega\) in that regime, after
passing to a later terminal tail if necessary, satisfies at least one of the
following local alternatives:
\begin{enumerate}[label=\textup{(\roman*)}]
\item \(\omega\) satisfies the fixed-cutoff localized-energy closure criterion of
\cref{paper7:def:fixed-carrier-closure-package};
\item \(\omega\) satisfies the moving-cutoff localized-energy closure criterion of
\cref{paper7:def:moving-carrier-closure-package};
\item \(\omega\) satisfies the exterior/leakage alternative: either
there are unit windows \(I_n=[T_n,T_n+1]\), selected times \(\tau_n\in I_n\),
radii \(R_n\to\infty\), and \(\eta>0\) such that
\[
  \sup_n\bigl(|a(\tau_n)|+\|b(\tau_n)\|_\infty\bigr)<\infty,
\]
\[
  \frac{1}{R_n^2}\int_{A_{R_n}}|V(y,\tau_n)|^2\,dy
  +
  \frac{1}{R_n}\int_{A_{R_n}}|V(y,\tau_n)|^3\,dy
  +
  \frac{1}{R_n}\int_{\widetilde A_{R_n}}|V(y,\tau_n)|^3\,dy
  \ge \eta
\]
for every \(n\), and this selected-window shell persistence is assigned to one
of the already named critical-tightness failure, exterior concentration, or
local windowed \(H^1\)-loss branches; or the normalized transition-carrier
blocks of the moving cutoff have a local occupation-state limit with the
transition-leakage carrier label of
\cref{paper7:lem:transition-carrier-label-is-exterior}, which is identified
there with the same exterior/leakage cluster.
\item \(\omega\) satisfies the translation alternative: there are unit
windows \(I_n=[T_n,T_n+1]\), selected times \(\tau_n\in I_n\), radii
\(R_n\to\infty\), cutoffs \(\Phi_n(y)=\phi(y/R_n)\), and \(\eta>0\) such that
\[
  \sup_n\bigl(|a(\tau_n)|+\|b(\tau_n)\|_\infty\bigr)<\infty,
\]
\[
  \frac{|b(\tau_n)|}{R_n}\int_{A_{R_n}}|V(y,\tau_n)|^2\,dy\ge\eta
\]
for every \(n\), and this persistent translation transport is assigned to the
already named modulation-integrability failure or to a multicore/exterior
branch produced after selected-window recentering;
\item \(\omega\) satisfies the positive-scale alternative: either
there are unit windows \(I_n=[T_n,T_n+1]\), selected times
\(\tau_n\in I_n\), radii \(R_n\to\infty\), cutoffs
\(\Phi_n(y)=\phi(y/R_n)\), and \(\eta>0\) such that
\[
  \sup_n\bigl(|a(\tau_n)|+\|b(\tau_n)\|_\infty\bigr)<\infty,
\]
\[
  a_+(\tau_n)\int |V(y,\tau_n)|^2
  \bigl(-y\cdot\nabla\Phi_n(y)\bigr)\,dy\ge\eta
\]
for every \(n\), and this persistent positive scale forcing is assigned to the
already named scale-rigid residual alternative or to a bounded selected-window
limit refinement of that alternative; or the normalized positive-scale-work
carrier blocks have a retained compact occupation-state limit whose compact-core
realization satisfies the scale-rigid branch clauses of
\cref{paper3:def:scale-rigid-branch}.  If the occupation state instead has an
earlier exterior/leakage, modulation-driven recentering, or negative-drift
label, the branch is read in that earlier alternative rather than in the
positive-scale alternative.
\item \(\omega\) enters the negative scale-drift obstruction of
\cref{paper6a:def:negative-drift}.
\end{enumerate}
The fixed-cutoff case is auxiliary here.  The canonical local alternative is the
moving one: either a nested moving cutoff closes the carrier, or a persistent
failure produces one of the displayed selected-window diagnostics and is assigned
to the corresponding adjacent stratum.
\end{definition}

\begin{remark}[How the moving diagnostics match the carrier components]
\label{paper7:rem:stratified-moving-routing}
The exterior/leakage alternative in
\cref{paper7:def:stratified-moving-carrier-routing} groups the moving cutoff,
viscous, convection, and pressure diagnostics.  The point is that the canonical
windowwise shell package
\cref{paper7:lem:windowwise-shell-pressure-translation}, together with the
local harmonic-pressure routing
\cref{paper8:lem:harmonic-cross-pressure-routing,paper8:lem:nested-ball-harmonic-tail},
routes pressure failure to enlarged-shell \(L^3\) persistence, pressure-only
critical mass, or exterior concentration.  Thus the pressure obstruction is
not an independent global theorem; it is another way of detecting local shell
mass that refuses to leave the selected windows.  The
translation alternative contains the \(J_5\) term and is meant to be combined
with the bounded selected-window modulation provided by
\cref{paper7:cor:unit-window-selection,paper7:cor:selected-window-modulation}.
The positive-scale alternative contains the \(J_7\) term and isolates the only
genuinely positive scale-dynamics obstruction left after the favorable
\(a_-\)-part has been separated.  The negative-drift term \(J_6\) is already
assigned to Section~6 and is not to be replaced by an unproved compact-cost
equivalence.  The occupation-state clauses are included so that a diffuse
nonintegrable error is not incorrectly converted into a pointwise lower bound;
it is instead read as an occupation state in the local ordered stratification.
\end{remark}

\begin{definition}[Moving test family]
\label{paper7:def:moving-test-family}
For a branch in the pre-carrier validated compact-cost regime, a moving test
family is a finite-valued nondecreasing locally absolutely continuous radius
\(R(\tau)\ge R_0\), with \(R(\tau)\to\infty\), together with the nested moving
cutoff
\[
  \Phi(\tau,y)=\phi(y/R(\tau))\ge\phi_{R_0}(y).
\]
For such a family, let \(J_1,\ldots,J_7\) denote the seven moving carrier
components of \cref{paper7:def:componentwise-annular-finite-error}.
\end{definition}

\begin{lemma}[A successful moving test family gives moving carrier closure]
\label{paper7:lem:successful-moving-test-family}
Assume repaired-gauge suitable local energy gives the moving carrier
subidentity in the sense of
\cref{paper7:lem:suitable-energy-gives-carrier-subidentities}.  Let \(\omega\)
be a branch in the pre-carrier validated compact-cost regime whose validated
compact cost is energy-controlled by the fixed identity cost.  If there exists
a moving test family in the sense of \cref{paper7:def:moving-test-family} such
that, for some \(\tau_1\ge\tau_0\),
\[
  J_k\in L^1(\tau_1,\infty)
  \qquad\text{for each }1\le k\le7,
\]
then \(\omega\) satisfies the moving-cutoff localized-energy closure criterion of
\cref{paper7:def:moving-carrier-closure-package} after restriction to the later
tail \([\tau_1,\infty)\).  In particular, \(\omega\) is excluded by
\cref{paper7:thm:moving-carrier-closure}.
\end{lemma}

\begin{proof}
Restrict the branch and the moving test family to the tail \([\tau_1,\infty)\).
Let \(R(\tau)\) and \(\Phi(\tau,y)\) be the restricted moving test family, with
components \(J_1,\ldots,J_7\).  By
\cref{paper6a:def:summable-moving-errors}, the integrability of these seven
components is exactly the summable moving-annulus condition for \(\Phi\).
Hence \(\int_{\tau_1}^\infty B_{\mathrm{mov}}(\tau)\,d\tau<\infty\) by
\cref{paper6a:lem:summable-moving-errors}.  The suitable local energy lemma
gives the moving carrier subidentity for the same cutoff.  Finite initial
localized energy on the pre-carrier regime is automatic by
\cref{paper7:lem:precarrier-finite-local-energy}, and the moving carrier
comparison is automatic by
\cref{paper7:lem:nested-moving-carrier-comparison}.  These are exactly the
clauses of \cref{paper7:def:moving-carrier-closure-package}.  Therefore
\cref{paper7:thm:moving-carrier-closure} excludes \(\omega\).
\end{proof}

\begin{lemma}[Retained pressure tails vanish on windowwise annuli]
\label{paper7:lem:retained-pressure-tail-windowwise}
Let \(\omega\) be a branch in the pre-carrier validated compact-cost regime,
represented by \((V,P,a,b)\) on a terminal tail, and let \(I=[A,B]\) be a
compact renormalized time interval.  Assume that no pressure-reconstruction,
exterior-leakage, separated-core, or critical-tightness failure label occurs on
the retained branch over \(I\).  Then, for every radial cutoff
\(\phi\in C_c^\infty(B_2)\) with \(\phi=1\) on \(B_1\),
\[
  \lim_{S\to\infty}
  \sup_{\substack{R\in AC(I),\ R(\tau)\ge S\\ R(\tau)<\infty}}
  \int_I
  \left|
  \int P(y,\tau)V(y,\tau)\cdot\nabla\phi(y/R(\tau))\,dy
  \right|d\tau
  =0 .
\]
Consequently, on the retained branch one may choose any prescribed summable
sequence \(\varepsilon_n>0\) and radii \(S_n\to\infty\) so that every
finite-valued moving radius \(R(\tau)\ge S_n\) on the \(n\)-th unit window has
pressure flux at most \(\varepsilon_n\) on that window.
\end{lemma}

\begin{proof}
Argue by contradiction on a compact interval \(I\).  If the displayed limit
fails, then there are \(\eta>0\), radii \(S_m\to\infty\), and finite-valued
moving radii \(R_m\in AC(I)\), \(R_m\ge S_m\), such that the pressure flux
through the annulus of \(\phi(\cdot/R_m(\tau))\) has \(L^1_\tau\)-mass at least
\(\eta\).
\par
Let
\[
  A_m(\tau)=\{R_m(\tau)<|y|<2R_m(\tau)\},\qquad
  \widetilde A_m(\tau)=\{R_m(\tau)/2<|y|<4R_m(\tau)\}.
\]
The cutoff gradient is supported in \(A_m(\tau)\) and satisfies
\(|\nabla\phi(y/R_m(\tau))|\le C_\phi/R_m(\tau)\).  On a finite partition of
\(\widetilde A_m(\tau)\), rescaled to a fixed reference annulus, the localized
Calder\'on--Zygmund pressure \(P_m^{\rm loc}\) generated by
\({\bf 1}_{\widetilde A_m(\tau)}V\otimes V\) obeys the scale-invariant bound
\[
  \|P_m^{\rm loc}(\tau)\|_{L^{3/2}(A_m(\tau))}
  \le
  C_\phi\|V(\tau)\|_{L^3(\widetilde A_m(\tau))}^2 .
\]
Therefore Hölder's inequality gives
\[
\begin{aligned}
  \int_I\left|
  \int P_m^{\rm loc}V\cdot\nabla\phi(y/R_m(\tau))\,dy
  \right|d\tau
  &\le
  C_\phi\int_I
  R_m(\tau)^{-1}
  \|V(\tau)\|_{L^3(\widetilde A_m(\tau))}^3\,d\tau  \\
  &\le
  C_{\phi,I}S_m^{-1}
  \sup_{\tau\in I}\int_{|y|>S_m/2}|V(y,\tau)|^3\,dy .
\end{aligned}
\]
This tends to zero because the retained branch has passed the critical
tightness row on \(I\).  Thus the local Calder\'on--Zygmund pressure part
cannot carry the lower bound \(\eta\).
\par
The remaining pressure contribution is harmonic on each retained rescaled
subannulus after subtracting functions of time.  If its pressure work did not
vanish, then on at least one element of the finite annular partition the
normalized harmonic pressure would have a positive \(L^{3/2}\)-mass after
rescaling.  The quantitative converse in
\cref{paper8:lem:nested-ball-harmonic-tail} would then force a positive
\(L^{3/2}\)-mass of the recorded annular source \(V\otimes V\), equivalently a
positive \(L^3\)-mass of \(V\), on a retained exterior annular cell.  This is
exactly an exterior-leakage or separated-core label under
\cref{paper8:thm:exterior-stratification}.  If the harmonic piece is not
generated by a recorded annular source, the ordered pressure atlas reads it as
a pressure-reconstruction failure or as a far-field pressure tail; the latter
is routed by
\cref{paper8:lem:harmonic-cross-pressure-routing} and
\cref{paper8:lem:exterior-pressure-vanishes}.  Every resulting label is
excluded by the retained-branch hypotheses on \(I\).  Hence the harmonic
pressure cannot carry the lower bound \(\eta\) either.
\par
Therefore the displayed uniform tail limit holds.  Applying it on each unit
window with \(\varepsilon_n\) gives the final scheduling statement.
\end{proof}

\begin{lemma}[Critical tightness gives a tailored windowwise radius schedule]
\label{paper7:lem:windowwise-tight-radius-schedule}
Let \(\omega\) be a branch in the pre-carrier validated compact-cost regime,
represented by \(V\) on a terminal tail.  For the unit windows
\[
  I_n=[\tau_0+n,\tau_0+n+1],\qquad n\in\mathbb N,
\]
set
\[
  A_n:=1+\int_{I_n}\bigl(|a(\tau)|+\|b(\tau)\|_\infty\bigr)\,d\tau .
\]
Then there exists a nondecreasing sequence \(R_n\to\infty\), with
\[
  R_n\ge 4\cdot 2^n,
\]
such that
\[
  \sup_{\tau\in I_n}\int_{|y|>R_n/4}|V(y,\tau)|^3\,dy
  \le
  2^{-3n}A_n^{-3/2}
  \qquad\text{for every }n.
\]
The sequence may also be chosen so that every finite-valued moving radius
\(\mathcal R(\tau)\ge R_n\) on \(I_n\) satisfies
\[
  \int_{I_n}
  \left|
  \int P(y,\tau)V(y,\tau)\cdot\nabla\phi(y/\mathcal R(\tau))\,dy
  \right|d\tau
  \le 2^{-n}.
\]
\end{lemma}

\begin{proof}
Since \(\omega\) lies in the pre-carrier validated compact-cost regime, the
critical-tightness test and the modulation-integrability test in
\cref{paper7:lem:no-unrecorded-cost-escape} have already passed.  Thus
\cref{paper7:def:critical-tightness} holds and
\cref{paper7:lem:local-modulation-integrability} gives
\[
  a,b\in L^1_{\rm loc}([\tau_0,\infty)).
\]
In particular each \(A_n\) is finite.
\par
For every \(n\), set
\[
  \varepsilon_n:=2^{-3n}A_n^{-3/2}>0.
\]
By uniform critical tightness, for each \(\varepsilon_n\) there exists
\(S_n<\infty\) such that
\[
  \sup_{\tau\ge\tau_0}\int_{|y|>S_n}|V(y,\tau)|^3\,dy\le \varepsilon_n .
\]
Since the branch is retained in the pre-carrier regime, pressure reconstruction,
exterior leakage, and critical tightness have not failed on \(I_n\).  Applying
\cref{paper7:lem:retained-pressure-tail-windowwise} with tolerance \(2^{-n}\),
choose \(P_n<\infty\) such that every finite-valued moving radius
\(R(\tau)\ge P_n\) on \(I_n\) has pressure flux at most \(2^{-n}\).
Choose \(R_n\) recursively so that
\[
  R_n\ge \max\{4S_n,\;P_n,\;4\cdot 2^n,\;R_{n-1}\}
\]
for \(n\ge1\), and \(R_0\ge \max\{4S_0,P_0,4\}\).  Then \(R_n\) is nondecreasing,
\(R_n\to\infty\), \(R_n\ge 4\cdot2^n\), and
\[
  \sup_{\tau\in I_n}\int_{|y|>R_n/4}|V(y,\tau)|^3\,dy
  \le
  \sup_{\tau\ge\tau_0}\int_{|y|>S_n}|V(y,\tau)|^3\,dy
  \le \varepsilon_n ,
\]
which is the desired estimate.
\end{proof}

\begin{lemma}[Canonical windowwise moving family controls shell, pressure, and translation terms]
\label{paper7:lem:windowwise-shell-pressure-translation}
Let \(\omega\) be a branch in the pre-carrier validated compact-cost regime.
Let \(R_n\) be the sequence supplied by
\cref{paper7:lem:windowwise-tight-radius-schedule}.  On the unit windows
\(I_n=[\tau_0+n,\tau_0+n+1]\), define the canonical interpolated radius
\(R(\tau)\) by
\[
  R(\tau)
  :=
  R_n+(\tau-(\tau_0+n))(R_{n+1}-R_n),
  \qquad \tau\in I_n,
\]
and the associated moving cutoff
\[
  \Phi(\tau,y):=\phi(y/R(\tau)).
\]
Write
\[
\begin{gathered}
  J_2(\tau)=\left|\int |V|^2\Delta\Phi(\tau,\cdot)\right|,\qquad
  J_3(\tau)=\left|\int |V|^2V\cdot\nabla\Phi(\tau,\cdot)\right|,\\
  J_4(\tau)=\left|\int P\,V\cdot\nabla\Phi(\tau,\cdot)\right|,\qquad
  J_5(\tau)=|b(\tau)|\left|\int |V|^2\nabla\Phi(\tau,\cdot)\right|.
\end{gathered}
\]
Then
\[
  \int_{\tau_0}^{\infty}
  \bigl(J_2(\tau)+J_3(\tau)+J_4(\tau)+J_5(\tau)\bigr)\,d\tau
  <\infty.
\]
\end{lemma}

\begin{proof}
The sequence \(R_n\) is nondecreasing and tends to \(+\infty\).  Hence the
piecewise affine interpolation \(R(\tau)\) is finite-valued, nondecreasing,
locally absolutely continuous, and satisfies \(R(\tau)\to\infty\).  Thus it is
an admissible moving test family in the sense of
\cref{paper7:def:moving-test-family}.
\par
By \cref{paper7:lem:windowwise-tight-radius-schedule},
\[
  \sup_{\tau\in I_n}\int_{|y|>R_n/4}|V(y,\tau)|^3\,dy
  \le \varepsilon_n:=2^{-3n}A_n^{-3/2},
\]
where
\[
  A_n=1+\int_{I_n}\bigl(|a(\tau)|+\|b(\tau)\|_\infty\bigr)\,d\tau.
\]
For \(\tau\in I_n\), one has \(R(\tau)\ge R_n\), hence
\[
  A_{R(\tau)}\cup\widetilde A_{R(\tau)}
  \subset \{|y|>R(\tau)/4\}\subset \{|y|>R_n/4\}.
\]
Therefore, for every \(\tau\in I_n\),
\[
  \int_{A_{R(\tau)}}|V(y,\tau)|^3\,dy
  +
  \int_{\widetilde A_{R(\tau)}}|V(y,\tau)|^3\,dy
  \le \varepsilon_n .
\]
Using Hölder on the annulus \(A_{R(\tau)}\),
\[
  \int_{A_{R(\tau)}}|V(y,\tau)|^2\,dy
  \le
  |A_{R(\tau)}|^{1/3}
  \left(\int_{A_{R(\tau)}}|V(y,\tau)|^3\,dy\right)^{2/3}
  \le
  C_\phi R(\tau)\varepsilon_n^{2/3}.
\]
\par
Since \(|\Delta\Phi(\tau,\cdot)|\le C_\phi R(\tau)^{-2}\mathbf 1_{A_{R(\tau)}}\),
\[
  J_2(\tau)
  \le
  C_\phi R(\tau)^{-2}\int_{A_{R(\tau)}}|V|^2
  \le
  C_\phi R(\tau)^{-1}\varepsilon_n^{2/3}
  \le
  C_\phi R_n^{-1}\varepsilon_n^{2/3}.
\]
Integrating on \(I_n\) and using \(R_n\ge 4\cdot 2^n\),
\[
  \int_{I_n}J_2(\tau)\,d\tau
  \le
  C_\phi R_n^{-1}\varepsilon_n^{2/3}
  \le
  C_\phi 2^{-3n}2^{-2n}A_n^{-1}.
\]
\par
Since \(|\nabla\Phi(\tau,\cdot)|\le C_\phi R(\tau)^{-1}\mathbf 1_{A_{R(\tau)}}\),
\[
  J_3(\tau)
  \le
  C_\phi R(\tau)^{-1}\int_{A_{R(\tau)}}|V|^3
  \le
  C_\phi R_n^{-1}\varepsilon_n,
\]
and therefore
\[
  \int_{I_n}J_3(\tau)\,d\tau
  \le
  C_\phi R_n^{-1}\varepsilon_n
  =
  C_\phi R_n^{-1}2^{-3n}A_n^{-3/2}.
\]
\par
The pressure component is controlled by the retained pressure-tail scheduling
included in \cref{paper7:lem:windowwise-tight-radius-schedule}.  Since
\(R(\tau)\ge R_n\) on \(I_n\), the choice of \(R_n\) gives
\[
  \int_{I_n}J_4(\tau)\,d\tau
  =
  \int_{I_n}
  \left|\int P(y,\tau)V(y,\tau)\cdot\nabla\Phi(\tau,y)\,dy\right|d\tau
  \le 2^{-n}.
\]
If the retained pressure-tail scheduling were unavailable, the preceding lemma
would assign the branch to a pressure/exterior local alternative before this
windowwise retained branch is reached.
\par
Finally,
\[
  J_5(\tau)
  \le
  C_\phi \|b(\tau)\|_\infty R(\tau)^{-1}\int_{A_{R(\tau)}}|V|^2
  \le
  C_\phi \|b(\tau)\|_\infty \varepsilon_n^{2/3}.
\]
Hence
\[
  \int_{I_n}J_5(\tau)\,d\tau
  \le
  C_\phi \varepsilon_n^{2/3}\int_{I_n}\|b(\tau)\|_\infty\,d\tau
  \le
  C_\phi 2^{-2n}A_n^{-1}(A_n-1)
  \le C_\phi 2^{-2n}.
\]
The four resulting series are summable in \(n\), proving the claim.
\end{proof}

\begin{corollary}[Only transition and scale terms remain on the windowwise schedule]
\label{paper7:cor:windowwise-only-transition-and-scale-remain}
Under the hypotheses of
\cref{paper7:lem:windowwise-shell-pressure-translation}, the canonical
windowwise moving family
\[
  R(\tau)
  =
  R_n+(\tau-(\tau_0+n))(R_{n+1}-R_n),
  \qquad \tau\in I_n,
\]
has summable moving carrier components \(J_2,J_3,J_4,J_5\).  Thus the only
unresolved moving-carrier terms on that genuine moving family are the
transition term \(J_1\), the positive scale-shell term \(J_7\), and the
negative scale-drift term \(J_6\).
\end{corollary}

\begin{proof}
This is exactly the conclusion of
\cref{paper7:lem:windowwise-shell-pressure-translation}, together with the fact
that the moving carrier has exactly seven components, of which
\cref{paper7:lem:windowwise-shell-pressure-translation} already controls
\(J_2,J_3,J_4,J_5\).
\end{proof}

\begin{definition}[Windowwise transition/leakage alternative statement]
\label{paper7:def:windowwise-transition-routing}
The pre-carrier validated compact-cost regime has windowwise
transition/leakage alternative if, for every branch \(\omega\) in that regime,
the canonical windowwise moving family of
\cref{paper7:lem:windowwise-shell-pressure-translation} has the following
property: if
\[
  J_1\notin L^1(\tau_1,\infty)
\]
on every later tail \([\tau_1,\infty)\), then \(\omega\) enters the
exterior/leakage alternative of
\cref{paper7:def:stratified-moving-carrier-routing}.
\end{definition}

\begin{definition}[Windowwise positive-scale alternative statement]
\label{paper7:def:windowwise-positive-scale-routing}
The pre-carrier validated compact-cost regime has windowwise positive-scale
alternative if, for every branch \(\omega\) in that regime, the canonical
windowwise moving family of
\cref{paper7:lem:windowwise-shell-pressure-translation} has the following
property: if
\[
  J_7\notin L^1(\tau_1,\infty)
\]
on every later tail \([\tau_1,\infty)\), then \(\omega\) enters the
positive-scale alternative of
\cref{paper7:def:stratified-moving-carrier-routing}, unless it has already
entered the exterior/leakage, modulation-driven recentering, or negative
scale-drift alternatives of the same stratified moving-localized-energy criterion.
\end{definition}

\subsubsection{Proof of the windowwise alternative statements}

We now prove that the two alternative properties of
\cref{paper7:def:windowwise-transition-routing,paper7:def:windowwise-positive-scale-routing}
hold for the canonical windowwise moving family of
\cref{paper7:lem:windowwise-shell-pressure-translation}.  The proof does not
convert arbitrary nonintegrability into a false persistent unit-window lower
bound.  Instead it uses the local compactness interpretation of a failed carrier
component: the nonintegrable component defines normalized carrier blocks; each
block is then labelled by the ordered first-failure stratification already used
in \cref{paper7:lem:no-unrecorded-cost-escape}.

\begin{definition}[Normalized carrier blocks for a failed component]
\label{paper7:def:normalized-carrier-blocks}
Let \(g\in L^1_{\rm loc}([\tau_0,\infty))\) be nonnegative and assume
\[
  \int_T^\infty g(\tau)\,d\tau=\infty
  \qquad\text{for every }T\ge\tau_0 .
\]
A normalized \(g\)-carrier block sequence on a tail \([\tau_1,\infty)\) is a
sequence of intervals \([s_m,t_m]\), with
\(\tau_1\le s_m<t_m\), \(s_m\to\infty\), \(t_m\le s_{m+1}\), and
\[
  \int_{s_m}^{t_m}g(\tau)\,d\tau=1 .
\]
The associated carrier probability measures are
\[
  d\mu_m(\tau):=g(\tau)\mathbf 1_{[s_m,t_m]}(\tau)\,d\tau .
\]
For \(g=J_i\), \(i\in\{1,7\}\), a \(J_i\)-occupation state means the terminal
local occupation state obtained as follows.  Fix any compact retained
renormalized cylinder and record on it the repaired-gauge variables
\((V,P,a,b)\), the moving radius \(R\), the cutoff \(\Phi\), and the component
density \(J_i\), using the weak local topology supplied by
\cref{paper1:prop:compactness} and the pressure reconstruction for
\((V,P)\), the compactified repaired-gauge topology for \((a,b,R,\Phi)\), and
the one-point compactification of \([0,\infty)\) for the scalar component
density.  Push
\(\mu_m\) forward by this local-data map and pass, on every compact cylinder
and then diagonally, to a weak-\(*\) limit of probability measures on this
metrizable local state space.  Any such limit is a \(J_i\)-occupation state.  Its
label row is the first closed stratum met by its support in the ordered
validation procedure of \cref{paper7:lem:no-unrecorded-cost-escape}, with the
carrier row refined by the active component \(J_i\).  If several named
diagnostics are met at the same row, they are read as the declared label
cluster for that row.
\end{definition}

\begin{lemma}[Closed local carrier diagnostics]
\label{paper7:lem:closed-carrier-diagnostics}
On the local state space of
\cref{paper7:def:normalized-carrier-blocks}, each row of the ordered
validation procedure in \cref{paper7:lem:no-unrecorded-cost-escape} is
represented by a closed diagnostic set, after grouping diagnostics that are
declared to belong to the same row.  The component-refined carrier diagnostic
for \(J_i\), \(i\in\{1,7\}\), is also closed on normalized \(J_i\)-carrier
states.
\end{lemma}

\begin{proof}
Fix a compact retained cylinder.  The solution variables are recorded in the
weak local compactness topology of \cref{paper1:prop:compactness}; pressure is
recorded modulo spatial constants in the corresponding weak
\(L^{3/2}_{\rm loc}\) topology; and the auxiliary scale, center, radius, and
cutoff data are recorded in the compactified repaired-gauge topology, with the
boundary faces labelled by loss of scale admissibility, modulation
integrability, exterior tightness, multicore escape, or local \(H^1\) control.
Thus leaving a retained compact row is convergence to one of the prescribed
boundary faces, hence a closed condition.
\par
The quantitative rows are closed for the standard lower-semicontinuity
reasons used throughout the retained compactness arguments: \(L^3\), pressure
\(L^{3/2}\), \(L^2_tH^1_x\), and local energy quantities are lower
semicontinuous under the recorded weak/strong local convergence, while
failure of a required finite bound is the closed face at infinity in the
compactified coordinate.  Critical-annulus, local \(H^1\)-loss,
modulation-loss, cost-compatibility, and carrier-comparison failures are
therefore closed diagnostic sets in this topology.  Equal-scale, separated,
same-point cascade, gauge-degenerate, and scale-drift alternatives are closed
because they are defined by limits of the recorded scale-center parameters in
the finite parameter partition.
\par
Finally, the component-refined carrier row is the closed face consisting of
nonzero \(J_i\,d\tau\)-occupation states.  Each approximating block has total
occupation mass
\[
  \mu_m(\mathfrak X)=\int_{s_m}^{t_m}J_i(\tau)\,d\tau=1,
\]
where \(\mathfrak X\) denotes the local occupation-state space.  Weak-\(*\)
limits of probability measures have the same total mass.  Hence the row
``\(J_i\) is not negligible/summable on the carrier tail'' is closed on the
normalized occupation-measure compactification.  This is a statement about the
occupation measure; it does not choose pointwise times and does not assert a
unit-window lower bound.
\end{proof}

\begin{lemma}[Compactness and labels for normalized occupation states]
\label{paper7:lem:carrier-state-compactness-labels}
Let \(i\in\{1,7\}\).  Every normalized \(J_i\)-carrier block sequence for a
pre-carrier branch has a subsequential \(J_i\)-occupation state in the sense of
\cref{paper7:def:normalized-carrier-blocks}.  The state has a unique first
label row in the finite ordered list of
\cref{paper7:lem:no-unrecorded-cost-escape}.  If no non-carrier label occurs
before carrier validation, then the first carrier label is the
component-refined \(J_i\)-moving-annulus summability failure.
\end{lemma}

\begin{proof}
On every compact retained cylinder, the pre-carrier hypotheses give the
represented suitable estimates, pressure reconstruction, local \(L^3\) bounds,
and local \(L^2_tH^1_x\) bounds used in
\cref{paper1:prop:compactness}.  The repaired-gauge modulation data are locally
integrable by \cref{paper7:lem:local-modulation-integrability}, and the moving
radius is locally absolutely continuous.  Hence the local-data maps take values
in a metrizable product of weakly compact bounded sets after restricting to a
subsequence and a compact cylinder.  Banach--Alaoglu, the pressure compactness
part of \cref{paper1:prop:compactness}, and a diagonal extraction over compact
cylinders give weak-\(*\) convergence of the pushed-forward probability
measures.  This is precisely the occupation state in
\cref{paper7:def:normalized-carrier-blocks}.
\par
The ordered validation list is finite and its strata are evaluated in order.
By \cref{paper7:lem:closed-carrier-diagnostics}, each row is a closed
diagnostic set on the local occupation-measure compactification.  The first row
whose closed diagnostic meets the occupation-state support is therefore unique;
ties inside that row are exactly the declared label cluster.  If a
non-carrier diagnostic appears before carrier validation, that row is the
label row.  If all non-carrier diagnostics pass on the support, the carrier
row is reached.  Because the measures are normalized by
\(\int J_i\,d\tau=1\), the limiting occupation state has total carrier mass
one.  Thus the active carrier row cannot be the statement that \(J_i\) is
summable and negligible on the carrier tail.  The component-refined failure of
\(J_i\)-moving-annulus summability is then the first carrier label.  This is a
compactness label, not a pointwise lower bound for \(J_i\).
\end{proof}

\begin{lemma}[Occupation-state extraction from nonintegrability]
\label{paper7:lem:nonintegrable-component-certificate}
Let \(g\in L^1_{\rm loc}([\tau_0,\infty))\) be nonnegative and assume
\[
  \int_T^\infty g(\tau)\,d\tau=\infty
  \qquad\text{for every }T\ge\tau_0 .
\]
Then every tail \([\tau_1,\infty)\) contains a normalized \(g\)-carrier block
sequence in the sense of \cref{paper7:def:normalized-carrier-blocks}.  After
passing to a subsequence, exactly one of the following two density regimes
occurs:
\begin{enumerate}[label=\textup{(C\arabic*)}]
\item \emph{Persistent unit-window carrier mass.}  There are \(\gamma>0\) and
unit intervals \(U_m=[T_m,T_m+1]\), \(T_m\to\infty\), such that
\[
  \int_{U_m}g(\tau)\,d\tau\ge\gamma .
\]
\item \emph{Diffuse carrier mass.}  For the normalized block subsequence
retained in this case,
\[
  \sup_{T\in[s_m,t_m]}
  \int_{[T,T+1]\cap[s_m,t_m]}g(\tau)\,d\tau\longrightarrow0 .
\]
The probability measures \(d\mu_m=g\mathbf 1_{[s_m,t_m]}d\tau\) are then the
diffuse local occupation states of \(g\).
\end{enumerate}
If, in addition, \(q(\tau):=|a(\tau)|+\|b(\tau)\|_\infty\) satisfies the
uniform unit-window bound
\[
  Q:=\sup_{T\ge\tau_0}\int_T^{T+1}q(\tau)\,d\tau<\infty .
\]
then in case \textup{(C1)} either there are selected times
\(\tau_m\in U_m\) such that
\[
  g(\tau_m)\ge \gamma/2,\qquad q(\tau_m)\le 4Q/\gamma .
\]
after relabelling, or the normalized \(g\)-carrier mass on \(U_m\) is carried
by the high-modulation sets
\[
  H_m:=\{\tau\in U_m:q(\tau)>4Q/\gamma\},
  \qquad
  \int_{H_m}g(\tau)\,d\tau\ge\gamma/2 .
\]
The latter is a modulation occupation state; it is not converted into a shell or
scale lower bound.
\end{lemma}

\begin{proof}
\emph{Block construction.}  Fix a tail \([\tau_1,\infty)\).  Choose
\(s_1\ge\tau_1\).  Since every tail integral of \(g\) is infinite, there is a
finite \(t_1>s_1\) with \(\int_{s_1}^{t_1}g=1\); for instance take the first
level time of the absolutely continuous primitive
\(\int_{s_1}^T g\).  Having chosen \(t_m\), choose \(s_{m+1}\ge t_m+1\) and
then \(t_{m+1}>s_{m+1}\) with \(\int_{s_{m+1}}^{t_{m+1}}g=1\).  This gives
normalized carrier blocks and \(s_m\to\infty\).
\par
\emph{Persistent versus diffuse.}  If case \textup{(C1)} fails, then for each
j there is \(m(j)\) such that every interval of the form
\([T,T+1]\cap[s_m,t_m]\) inside every retained block with index
\(m\ge m(j)\) has \(g\)-mass at most \(1/j\).  A diagonal subsequence gives
the displayed diffuse condition.  Conversely, if
\textup{(C1)} holds then the block sequence has a persistent unit-window
carrier by definition.  The two alternatives are therefore exhaustive after
subsequence extraction.
\par
\emph{Bounded modulation on persistent windows.}  In case \textup{(C1)}, set
\[
  Q:=\sup_{T\ge\tau_0}\int_T^{T+1}q(\tau)\,d\tau .
\]
If \(Q=0\), then \(q=0\) for a.e. time in every \(U_m\).  Since
\(\int_{U_m}g\ge\gamma\) and \(|U_m|=1\), the set
\(\{\tau\in U_m:q(\tau)=0,\ g(\tau)\ge\gamma/2\}\) has positive measure after
modifying \(g,q\) only on null sets; otherwise \(g<\gamma/2\) a.e. on
\(\{q=0\}\cap U_m\) and \(\int_{U_m}g\le\gamma/2<\gamma\).  Choose
\(\tau_m\) to be a Lebesgue time in that set.  This gives the low-modulation
alternative.
Assume now that \(Q>0\).  Set
\[
  H_m:=\{\tau\in U_m:q(\tau)>4Q/\gamma\},
  \qquad E_m:=U_m\setminus H_m .
\]
Chebyshev gives \(|H_m|\le\gamma/4\).  Since
\(\int_{U_m}g\ge\gamma\), either \(\int_{E_m}g\ge\gamma/2\) or
\(\int_{H_m}g\ge\gamma/2\) along a subsequence.  In the first case
\(|E_m|\le1\) gives a Lebesgue time \(\tau_m\in E_m\) with
\(g(\tau_m)\ge\gamma/2\) and \(q(\tau_m)\le4Q/\gamma\).  In the second case
the displayed high-modulation carrier mass holds.  This proves the final
dichotomy and deliberately leaves the second case as a occupation state rather
than replacing it by a pointwise shell or scale estimate.
\end{proof}

\begin{lemma}[Transition carrier label is exterior/leakage]
\label{paper7:lem:transition-carrier-label-is-exterior}
\label{paper7:lem:transition-carrier-state-stratification}
Let \(\omega\) be a branch in the pre-carrier validated compact-cost regime,
and let \(R(\tau)\) be the canonical windowwise moving family of
\cref{paper7:lem:windowwise-shell-pressure-translation}.  If
\[
  J_1\notin L^1(\tau_1,\infty)
\]
on every later tail, then every \(J_1\)-occupation state of
\cref{paper7:def:normalized-carrier-blocks} satisfies the exterior/leakage
alternative of \cref{paper7:def:stratified-moving-carrier-routing}.  More
precisely, either an earlier non-carrier exterior/leakage label occurs, or the
component-refined \(J_1\)-moving-annulus summability failure is the
transition-leakage carrier label, which belongs to the same exterior/leakage
cluster.
\end{lemma}

\begin{proof}
By \cref{paper7:lem:carrier-state-compactness-labels}, every normalized
\(J_1\)-block has a subsequential labelled occupation state.  Since \(\omega\) is
in the pre-carrier regime, the analytic, representation, retained finite
critical-mass, compact-tightness, local \(H^1\), modulation, and cost-compatibility rows in
\cref{paper7:def:precarrier-validated-compact-channel} have already passed on
the retained branch.  By
\cref{paper7:lem:closed-carrier-diagnostics}, these rows are closed diagnostic
sets in the occupation-measure compactification.  Hence if one fails in the carrier
limit, the failure is exactly a loss of retained compactness, exterior
concentration/multicore escape, or local windowed \(H^1\)-loss by the ordered
classification in \cref{paper7:lem:no-unrecorded-cost-escape}.  Those are
exterior/leakage labels.
\par
It remains to identify the first carrier label when all earlier rows remain
retained.  The component
\[
  J_1(\tau)=\left|\int |V|^2\partial_\tau\Phi(\tau,\cdot)\right|
\]
is the time-transition part of the moving cutoff.  It contains no pressure,
convection, translation, or scale-drift factor.  The exact identity
\[
  \partial_\tau\Phi(\tau,y)
  =
  -\frac{R'(\tau)}{R(\tau)^2}\,y\cdot\nabla\phi(y/R(\tau))
\]
shows that \(J_1\) measures only occupation of the moving boundary of the
retained compact window.  Therefore the component-refined carrier row supplied
by \cref{paper7:lem:carrier-state-compactness-labels} is the
\emph{transition-leakage carrier label}: the retained moving boundary carries
non-summable transition occupation.  This is not a new compact-cost stratum.
In the finite compactness list it is assigned to the same exterior/leakage
cluster as critical-tightness failure, exterior concentration, and local
windowed \(H^1\)-loss, because the only local interpretation of non-summable
boundary transition in a retained compact single-core state is that mass is
crossing the retained boundary or the compact regularity needed to ignore that
crossing has failed.  Thus every possible first label satisfies the
exterior/leakage alternative.
\end{proof}

\begin{theorem}[Windowwise transition/leakage alternative by occupation states]
\label{paper7:thm:windowwise-transition-routing-certificate}
Let \(\omega\) be a branch in the pre-carrier validated compact-cost regime,
and let \(R(\tau)\) be the canonical windowwise moving family of
\cref{paper7:lem:windowwise-shell-pressure-translation}.  If
\[
  J_1\notin L^1(\tau_1,\infty)
\]
on every later tail, then \(\omega\) enters the exterior/leakage alternative
of \cref{paper7:def:stratified-moving-carrier-routing}.
\end{theorem}

\begin{proof}
Apply \cref{paper7:lem:nonintegrable-component-certificate} with
\(g=J_1\).  It gives normalized \(J_1\)-carrier blocks on every terminal tail,
with either persistent unit-window mass or diffuse carrier mass.  In the
persistent case, the supplementary modulation selection gives bounded
selected-time representatives when the low-modulation subcase occurs; in the
high-modulation subcase and in the diffuse case one keeps the normalized
occupation state instead of imposing an artificial pointwise lower bound.
\par
All three outputs are \(J_1\)-occupation states in the sense of
\cref{paper7:def:normalized-carrier-blocks}.  By
\cref{paper7:lem:transition-carrier-state-stratification}, each such state
either has an earlier exterior/leakage label or has the transition-leakage
carrier label identified there with the exterior/leakage cluster.  This is
exactly the exterior/leakage alternative allowed in
\cref{paper7:def:stratified-moving-carrier-routing}.
\end{proof}

\begin{corollary}[Windowwise transition/leakage alternative holds]
\label{paper7:cor:windowwise-transition-routing-holds}
The pre-carrier validated compact-cost regime has the windowwise
transition/leakage alternative property in the sense of
\cref{paper7:def:windowwise-transition-routing}.
\end{corollary}

\begin{proof}
This is the definition of
\cref{paper7:def:windowwise-transition-routing} applied with
\cref{paper7:thm:windowwise-transition-routing-certificate}.
\end{proof}

\begin{lemma}[Retained positive scale-work realizes a scale-rigid core]
\label{paper7:lem:positive-scale-carrier-core-realization}
Let \(\omega\) be a branch in the pre-carrier validated compact-cost regime,
and let \(R(\tau)\) be the canonical windowwise moving family.  Let
\(\mu_m\) be any normalized \(J_7\)-carrier block sequence and let \(\mathfrak
s\) be a subsequential \(J_7\)-occupation state.  Assume that \(\mathfrak s\) has
no earlier exterior/leakage label, no modulation-driven recentering label, and
no negative scale-drift label.  Then, after admissible scale-center
reselection, the represented sequence carried by \(\mathfrak s\) is produced
by Step~3 or Step~4 of \cref{paper3:thm:same-point-reduction} or by
\cref{paper3:thm:separated-point-reduction}, and satisfies the four clauses
(R1)--(R4) of \cref{paper3:def:scale-rigid-branch}.
\end{lemma}

\begin{proof}
\emph{Step 1: what the \(J_7\)-carrier records.}  By definition,
\[
  J_7(\tau)=
  a_+(\tau)\int |V(y,\tau)|^2
  \bigl(-y\cdot\nabla\Phi(\tau,y)\bigr)\,dy .
\]
Thus each normalized block satisfies
\[
  \int_{s_m}^{t_m}
  a_+(\tau)\int |V|^2(-y\cdot\nabla\Phi)\,dy\,d\tau=1 .
\]
This is positive scale-generator work carried by the retained moving window.
It is a scale-center compactness datum, not a consequence of an annular
\(L^2\)-to-CKN conversion.
\par
\emph{Step 2: compact-core realization.}  By
\cref{paper7:lem:carrier-state-compactness-labels}, \(\mathfrak s\) is a
labelled local occupation state.  The absence of an exterior/leakage label means
that the state remains in the retained compact single-core window: critical
tightness has not failed, no exterior/multicore loss has occurred, and local
windowed \(H^1\) control is retained.  Consequently the compactness,
pressure-reconstruction, and suitable-energy estimates inherited from the
pre-carrier regime give a fixed compact selected cylinder \(Q_{R_*}\) on
which the represented sequence has the uniform \(A+E+C+D\) bound required in
(R1).  This is only the usual compact-cylinder suitable bound; it is not an
\(L^\infty\) or regularity assertion.
\par
The positive inner concentration needed for (R2) is inherited from the retained
finite critical-mass row itself, not inferred from the \(J_7\)-density.  This is
the compact-window active mass supplied by
\cref{paper6a:lem:terminal-local-critical-mass-routing}.  Since the present
occupation state has no exterior/leakage label, the corresponding critical mass
remains inside the retained compact single-core region.  Cover a
slightly smaller retained cylinder by finitely many CKN cylinders.  If each
member of this finite cover had arbitrarily small velocity-pressure CKN
quantity along the represented sequence, then the retained positive local
critical mass would either vanish on the compact core or exit through the
exterior face; the first contradicts the retained mass row and the second is
exactly the exterior/leakage label excluded above.  Thus, after passing to a
subsequence and reselecting the scale-center within this finite cover, there
are \(Q_{r_*}\Subset Q_{R_*}\) and \(\eta_0>0\) with positive CKN mass.  This
is (R2).
\par
If the modulation were unbounded on the selected compact cylinder, the
component-refined carrier label would be the modulation-driven recentering
branch.  That label is excluded by hypothesis, so the repaired-gauge
modulation remains bounded on the selected cylinder; this is (R3).
\par
\emph{Step 3: scale-rigid relative selection.}  The normalized \(J_7\)-mass is
nonzero positive scale-generator work.  If this work could be absorbed by a
summable scale correction, the \(J_7\)-moving-annulus row would pass on a
terminal carrier tail, contradicting the normalized carrier label.  If the
work instead caused scale collapse or negative scale drift, the first label
would be the negative-drift branch, excluded by hypothesis.
\par
Apply the finite parameter partition
\cref{paper3:lem:parameter-partition} to the retained compact core selected in
Step~2 and to the inherited active ledger.  A separated physical-center
relation is exactly the input of
\cref{paper3:thm:separated-point-reduction}.  A same-point comparable-scale
relation is reduced by Step~3 of
\cref{paper3:thm:same-point-reduction}, and a strict same-point scale
separation is reduced by Step~4 of that theorem.  The nonretained outputs of
this partition are exterior/leakage, modulation-driven recentering, negative
drift, or perturbative absorption; the first three are excluded by hypothesis
and perturbative absorption would make the \(J_7\)-row summable.  The only
retained compact scale-center label left is therefore the scale-rigid residual
alternative produced by those reductions: the relative scale-center selection
is rigid in the sense of (R4), either through the gauge-degenerate selection or
through the innermost same-point cascade after outer profiles have been
absorbed.  Thus the production requirement in
\cref{paper3:def:scale-rigid-branch} and clauses (R1)--(R4) all hold.
\end{proof}

\begin{theorem}[Positive scale-work occupation states enter the scale-rigid branch]
\label{paper7:thm:positive-scale-certificate-enters-scale-rigid}
Let \(\omega\) be a branch in the pre-carrier validated compact-cost regime,
and let \(R(\tau)\) be the canonical windowwise moving family.  Let
\(\mu_m\) be any normalized \(J_7\)-carrier block sequence, and let
\(\mathfrak s\) be any subsequential \(J_7\)-occupation state generated by
\(\mu_m\).  If \(\mathfrak s\) is not labelled by exterior/leakage, by
modulation-driven recentering, or by negative scale drift, then the represented
branch carried by \(\mathfrak s\) is excluded from the local Type II class by
\cref{paper3:prop:scale-rigidity-discharged}.
\end{theorem}

\begin{proof}
\Cref{paper7:lem:positive-scale-carrier-core-realization} realizes the
\(J_7\)-occupation state \(\mathfrak s\) as a scale-rigid local branch satisfying
\cref{paper3:def:scale-rigid-branch}.  Applying
\cref{paper3:prop:scale-rigidity-discharged} excludes that branch.
\end{proof}

\begin{lemma}[Positive-scale occupation-state stratification]
\label{paper7:lem:positive-scale-work-stratification}
Let \(\omega\) be a branch in the pre-carrier validated compact-cost regime,
and let \(R(\tau)\) be the canonical windowwise moving family of
\cref{paper7:lem:windowwise-shell-pressure-translation}, with cutoff
\(\Phi(\tau,y)=\phi(y/R(\tau))\) and \(J_7\) defined as in
\cref{paper7:def:componentwise-annular-finite-error}.  If
\[
  J_7\notin L^1(\tau_1,\infty)
\]
on every later tail, then every \(J_7\)-occupation state is labelled either by
exterior/leakage, by modulation-driven recentering, by the Section~6 negative
scale-drift branch, or by the scale-rigid residual branch excluded in
\cref{paper7:thm:positive-scale-certificate-enters-scale-rigid}.
\end{lemma}

\begin{proof}
Occupation-state extraction is
\cref{paper7:lem:nonintegrable-component-certificate} with \(g=J_7\).
The ordered first-failure map of \cref{paper7:lem:no-unrecorded-cost-escape}
first removes exterior loss, critical-tightness failure, local \(H^1\)-loss,
and modulation-driven recentering.  The scale decomposition then separates
the favorable negative scale-drift obstruction, already assigned to Section~6,
from retained compact states carrying positive scale work.  The retained
compact states are realized as scale-rigid states by
\cref{paper7:lem:positive-scale-carrier-core-realization} and excluded by
\cref{paper7:thm:positive-scale-certificate-enters-scale-rigid}.  These
alternatives exhaust the \(J_7\)-carrier because no other moving component is
being tested here and \(J_2,J_3,J_4,J_5\) are already summable on the
canonical windowwise family by
\cref{paper7:lem:windowwise-shell-pressure-translation}.
\end{proof}

\begin{theorem}[Windowwise positive-scale alternative by occupation states]
\label{paper7:thm:windowwise-positive-scale-routing-certificate}
Let \(\omega\) be a branch in the pre-carrier validated compact-cost regime,
and let \(R(\tau)\) be the canonical windowwise moving family of
\cref{paper7:lem:windowwise-shell-pressure-translation}, with cutoff
\(\Phi(\tau,y)=\phi(y/R(\tau))\) and \(J_7\) defined as in
\cref{paper7:def:componentwise-annular-finite-error}.  If
\[
  J_7\notin L^1(\tau_1,\infty)
\]
on every later tail, then \(\omega\) enters the positive-scale alternative of
\cref{paper7:def:stratified-moving-carrier-routing}, unless it
has already entered the exterior/leakage, modulation-driven recentering, or
negative-drift alternatives.
\end{theorem}

\begin{proof}
Apply \cref{paper7:lem:positive-scale-work-stratification}.  Exterior/leakage
carrier labels already satisfy alternative \textup{(iii)} of
\cref{paper7:def:stratified-moving-carrier-routing}, modulation-driven
recentering satisfies alternative \textup{(iv)}, and negative drift satisfies
alternative \textup{(vi)}.  In the retained positive-scale case,
\cref{paper7:thm:positive-scale-certificate-enters-scale-rigid} places the
carrier in the scale-rigid residual branch and then applies
\cref{paper3:prop:scale-rigidity-discharged}.  The persistent unit-window
subcase gives the pointwise selected-time form of alternative \textup{(v)}
whenever bounded-modulation points are selected; the high-modulation and
diffuse subcases give the occupation-state form of the same alternative.  Hence
there is no unhandled diffuse \(J_7\) tail.
\end{proof}

\begin{corollary}[Windowwise positive-scale alternative holds]
\label{paper7:cor:windowwise-positive-scale-routing-holds}
The pre-carrier validated compact-cost regime has the windowwise
positive-scale alternative property in the sense of
\cref{paper7:def:windowwise-positive-scale-routing}.
\end{corollary}

\begin{proof}
This is exactly the alternatives allowed by
\cref{paper7:def:windowwise-positive-scale-routing}, supplied by
\cref{paper7:thm:windowwise-positive-scale-routing-certificate}.
\end{proof}

\begin{remark}[Status of the windowwise alternative definitions]
\label{paper7:rem:windowwise-routing-discharged}
\Cref{paper7:cor:windowwise-positive-scale-routing-holds,paper7:cor:windowwise-transition-routing-holds}
together prove unconditionally that the pre-carrier validated compact-cost
regime has the alternative properties of
\cref{paper7:def:windowwise-transition-routing,paper7:def:windowwise-positive-scale-routing}.
The two definitions are therefore proved properties in the present manuscript;
their proof is reduced to
the occupation-measure compactness and extraction argument
\cref{paper7:lem:closed-carrier-diagnostics,paper7:lem:carrier-state-compactness-labels,paper7:lem:nonintegrable-component-certificate},
the transition carrier label in
\cref{paper7:lem:transition-carrier-state-stratification}, the positive
scale-work compact-core realization
\cref{paper7:lem:positive-scale-carrier-core-realization}, the scale-rigid
exclusion argument
\cref{paper7:thm:positive-scale-certificate-enters-scale-rigid}, and the
canonical windowwise schedule of
\cref{paper7:lem:windowwise-tight-radius-schedule}.  The proof does not use
the invalid implication from nonintegrability to persistent unit-window lower
bounds; diffuse tails are retained as normalized occupation states and classified by
the local ordered stratification.
\end{remark}

\begin{theorem}[Canonical windowwise family reduces the alternatives to transition, positive scale, and negative drift]
\label{paper7:thm:windowwise-routing-gives-stratified-routing}
Assume every branch in the pre-carrier validated compact-cost regime uses a
validated compact cost that is energy-controlled by the fixed identity cost.
Assume repaired-gauge suitable local energy gives the moving carrier
subidentity in the sense of
\cref{paper7:lem:suitable-energy-gives-carrier-subidentities}.  Assume also
that the regime has the windowwise transition/leakage and windowwise
positive-scale alternative statements of
\cref{paper7:def:windowwise-transition-routing,paper7:def:windowwise-positive-scale-routing},
and that failure of \(J_6\in L^1(\tau_1,\infty)\) on every later tail for the
canonical windowwise moving family forces entry into the negative scale-drift
obstruction of \cref{paper6a:def:negative-drift}.  Then the regime has
stratified moving-cutoff alternative in the sense of
\cref{paper7:def:stratified-moving-carrier-routing}.
\end{theorem}

\begin{proof}
Let \(\omega\) be a branch in the pre-carrier validated compact-cost regime,
and let \(R(\tau)\) be the canonical windowwise moving family of
\cref{paper7:lem:windowwise-shell-pressure-translation}.  Write
\(J_1,\ldots,J_7\) for its moving carrier components.
\par
By \cref{paper7:lem:windowwise-shell-pressure-translation},
\[
  J_2,J_3,J_4,J_5\in L^1(\tau_0,\infty).
\]
If there exists \(\tau_1\ge\tau_0\) such that
\[
  J_1,J_6,J_7\in L^1(\tau_1,\infty),
\]
then all seven components are integrable on \([\tau_1,\infty)\).  Therefore
\cref{paper7:lem:successful-moving-test-family} gives the moving-cutoff
localized-energy closure criterion, so \(\omega\) satisfies alternative \textup{(ii)} of
\cref{paper7:def:stratified-moving-carrier-routing}.
\par
Assume now that no later tail has \(J_1,J_6,J_7\) all integrable.  Since
\(J_2,J_3,J_4,J_5\) are already integrable on the full tail, at least one of
\(J_1,J_6,J_7\) fails on arbitrarily late tails.  Fix such a component.
Failure on arbitrarily late tails implies failure on every later tail: if the
component were integrable on some tail, it would be integrable on every
subsequent tail, contradicting arbitrarily late failure.  If \(J_1\) fails on
every later tail, the windowwise
transition/leakage alternative statement gives alternative
\textup{(iii)} of
\cref{paper7:def:stratified-moving-carrier-routing}.  If \(J_6\) fails on
every later tail, the negative scale-drift routing statement gives alternative
\textup{(vi)}.  If \(J_7\) fails on every later tail, the windowwise
positive-scale alternative statement gives alternative \textup{(v)}, unless the
branch has already entered alternative \textup{(iii)}, \textup{(iv)}, or
\textup{(vi)}.
Therefore one of the alternatives in
\cref{paper7:def:stratified-moving-carrier-routing} always holds.
\end{proof}

\begin{corollary}[Adjacent closure reduced to windowwise transition and positive-scale alternative]
\label{paper7:cor:adjacent-closure-reduced-to-windowwise-routing}
Assume the non-carrier adjacent strata listed in
\cref{paper7:thm:noncarrier-exits-route} are closed by their corresponding
local compactness or interface arguments in the ambient Type II class.  Assume
also that the exterior concentration, local windowed \(H^1\)-loss,
modulation-driven recentering branches, scale-rigid residual branches, and the
negative scale-drift obstruction are closed by their assigned local arguments
in the ambient class.  Assume every branch in the pre-carrier validated
compact-cost regime uses a validated
compact cost that is energy-controlled by the fixed identity cost, that
repaired-gauge suitable local energy gives the moving carrier subidentity, and
that the windowwise transition/leakage and windowwise positive-scale alternative
statements of
\cref{paper7:def:windowwise-transition-routing,paper7:def:windowwise-positive-scale-routing}
hold, together with the negative-drift routing statement used in
\cref{paper7:thm:windowwise-routing-gives-stratified-routing}.  Then the
compact-cost regime has adjacent-stratum closure in the sense of
\cref{paper7:def:compact-cost-adjacent-closure}.
\end{corollary}

\begin{proof}
Let \(\omega\) be a branch with active supercritical scaling and validated
compact Type II cost divergence.  By
\cref{paper7:thm:classwise-retained-compact-coverage}, either \(\omega\)
remains in the retained compact stratum or it enters one of the named adjacent
strata.  The retained compact stratum is excluded by
\cref{paper7:cor:retained-stratum-cost-exclusion}.  If \(\omega\) enters a
non-carrier adjacent stratum before the pre-carrier regime, the first
closure condition closes it by \cref{paper7:thm:noncarrier-exits-route}.  It remains
to consider a branch in the pre-carrier validated compact-cost regime.
\Cref{paper7:thm:windowwise-routing-gives-stratified-routing} gives one of
the alternatives in \cref{paper7:def:stratified-moving-carrier-routing}.
Alternatives \textup{(i)} and \textup{(ii)} are excluded by the fixed and
moving carrier closure theorems
\cref{paper7:thm:fixed-carrier-closure,paper7:thm:moving-carrier-closure}.
The remaining alternatives are precisely the exterior concentration, local
windowed \(H^1\)-loss, modulation-driven recentering, scale-rigid residual,
positive-scale, or negative scale-drift branches listed in the hypotheses of
this corollary, hence are closed by their assigned local arguments.  Therefore
every compact-cost branch is either excluded on the retained stratum or closed
in a named adjacent branch, which is
\cref{paper7:def:compact-cost-adjacent-closure}.
\end{proof}

\begin{lemma}[Ambient non-carrier exits close by ordered local routing]
\label{paper7:lem:ambient-noncarrier-exits-close}
Let \(\omega\) be an admissible physical local Type II branch which has passed
the covered-entry theorem and then reaches the compact-cost validation layer.
If \(\omega\) leaves the retained compact-cost stratum through a non-carrier
first failure listed in \cref{paper7:thm:noncarrier-exits-route}, then the
branch is either class-incompatible, closed by a local closure theorem, or
assigned to a named local state-space row whose closure is supplied by the
ambient carrier/terminal assembly.
\end{lemma}

\begin{proof}
The ordered validation in
\cref{paper7:lem:no-unrecorded-cost-escape,paper7:thm:noncarrier-exits-route}
lists the possible non-carrier first failures.  We check them in that order.
Exhaustion failures are routed by
\cref{paper6a:thm:typeII-analytic-data-exhaustive}; the only alternatives not
retained there are class-incompatible alternatives or named local exits.
Representation, pressure, and modulation failures are the ordered failures in
\cref{cor:c2-ordered-representation,cor:c2r-ordered-classification} and the
absolutely continuous repaired-gauge consequences of
\cref{paper6a:thm:ac-repaired-gauge-consequences}; when the repaired chart is
available, pressure and modulation are consequences, and when it is not
available the failure is the representation alternative.  On a physical branch
which has already passed covered entry, representation failures cannot remain
as unclassified analytic gaps: if a nondegenerate represented chart is not
available, the local state-decomposition places the branch in a
multibubble/cascade, gauge-degenerate, scale-degenerate, finite-cost
scale-collapse, compact scale-collapse, or scale-rigid alternative.  These rows
are closed by the local closure theorems
\cref{paper6:thm:multibubble,paper6:thm:scale-cost,paper5:thm:compact-scale-collapse-rigidity-discharged,paper3:prop:scale-rigidity-discharged}.
\par
Critical-mass failures are used here only as local terminal-window diagnostics.
Failure of lower retained mass means that the candidate was not a retained
active core, while failure of a finite local upper \(L^3\) bound on a compact
selected cylinder produces an additional compact active label; this is routed by
\cref{paper6a:lem:terminal-local-critical-mass-routing}.  Once the retained
local finite critical-mass window is reached, the
compact tests are closed by
\cref{paper6a:thm:retained-compact-tests-closed,paper6a:cor:compact-test-closure-discharged}.
Exterior tightness and leakage are routed by
\cref{paper8:thm:exterior-stratification,paper8:thm:exterior-concentration-routed}.
Local windowed \(H^1\)-loss is the
rough-core row; by \cref{paper6:thm:rough-core} it forces failure of compact
CKN control on enclosing compact cylinders and hence returns to the
multibubble/cascade package \cref{paper6:thm:multibubble}.  Multibubble,
multicore, cascade, and gauge-degenerate rows are closed directly by
\cref{paper6:thm:multibubble}.  Finite-cost scale collapse is eliminated by
\cref{paper6:thm:scale-cost}.  Scale-rigid residual and retained compact
scale-collapse rows are eliminated by
\cref{paper3:prop:scale-rigidity-discharged,paper5:thm:compact-scale-collapse-rigidity-discharged}.
Finally, cost well-posedness and cost-comparison failures are precisely the
ordered cost rows of
\cref{def:c4-localized-cost,def:c4-cost-comparison,lem:c4-identity-cost,thm:c4-cost-comparison};
for the canonical identity cost the comparison is a theorem, and any
non-identity failure is recorded as its own cost-compatibility exit.
The thin-drift, autonomous-modulation, and compactness-package first failures
which can arise inside the scale-collapse stratification are the first failures
recorded by \cref{paper5:lem:retained-scale-collapse-first-failure}.  Thin
drift is routed by \cref{paper6a:thm:thin-negative-drift-closed}; compactness
package failure is routed by
\cref{paper5:thm:selected-compactness-failure-routed}; modulation failure
remains a non-retained scale-collapse exit.  The scale-drift face is recorded
by \cref{paper6a:cor:scale-drift-alternatives}, and the thin face is closed as
an independent terminal row by \cref{paper6a:thm:thin-negative-drift-closed}.
The retained compact scale-collapse state is excluded by
\cref{paper5:thm:compact-scale-collapse-rigidity-discharged}.
Thus every non-carrier first failure is either class-incompatible, closed by a
local closure theorem, or recorded as a named carrier/scale/cost row for the
ambient assembly.  No step uses a global \(L^3\) bound, a global profile
theorem, or vanishing localized dissipation outside the selected local branch.
\end{proof}

\begin{lemma}[Carrier applicability on physical compact-cost branches]
\label{paper7:lem:ambient-carrier-applicability}
Every physical local Type II branch in the pre-carrier validated compact-cost
regime either exits through the ordered cost-compatibility row, or else
satisfies the carrier applicability hypotheses used in
\cref{paper7:thm:windowwise-routing-gives-stratified-routing} on the retained
canonical channel: the validated compact cost is the fixed identity cost and is
energy-controlled with constant \(1\), and suitable repaired-gauge local energy
supplies the fixed and moving carrier subidentities.
\end{lemma}

\begin{proof}
The pre-carrier regime already includes the repaired-gauge representation,
local pressure reconstruction, local cost well-posedness, and compatibility of
the validated compact cost.  The ordered cost row is deliberately read before
carrier validation.  If the selected non-identity diagnostic has only the lower
comparison used for cost-divergence transfer, but lacks the upper
energy-control comparison needed by the carrier inequality, the branch is not
silently promoted to an energy-controlled carrier branch; it is recorded as the
cost-compatibility or carrier-comparison exit handled by
\cref{paper7:lem:ambient-noncarrier-exits-close}.  On the retained carrier
branch the canonical cost channel is selected:
\[
  \mathfrak C_{\rm II}^{NS}=\tilde{\mathfrak D}_{R_0}.
\]
By \cref{lem:c4-identity-cost} this is a valid localized cost, and by
\cref{paper7:def:energy-controlled-validated-cost} it is energy-controlled
with constant \(1\).  This uses the identity equality only; no unproved
upper comparison for a non-identity cost is invoked.
\par
The branch comes from a suitable physical solution and has an absolutely
continuous repaired-gauge chart on each compact renormalized window.  Hence
\cref{paper6a:lem:ac-pullback-local-energy} gives the renormalized local energy
inequality, and
\cref{paper7:lem:suitable-energy-gives-carrier-subidentities} converts it into
the fixed and moving carrier subidentities.  These are local compact-window
statements; no whole-space energy or global compactness estimate is used.
\end{proof}

\begin{lemma}[Negative drift is an assigned local branch]
\label{paper7:lem:ambient-negative-drift-branch}
On every physical local Type II branch in the canonical windowwise
pre-carrier regime, failure of \(J_6\in L^1(\tau_1,\infty)\) on every later
tail is exactly the negative scale-drift branch of
\cref{paper6a:def:negative-drift}.  This branch is routed by the local
scale-drift alternatives and the local scale-collapse packages; it is not
an additional ambient hypothesis.
\end{lemma}

\begin{proof}
The \(J_6\) component is the negative-scale contribution in the moving
localized-energy identity.  By
\cref{paper6a:def:negative-drift,paper6a:cor:scale-drift-alternatives}, its
nonintegrability on every later tail is precisely the scale-collapse drift
obstruction.  From this point the branch is not closed by feeding the same
adjacent-closure statement back into itself; it is passed to the local
scale-collapse stratification.
\par
Before applying scale-collapse rigidity, verify the local prerequisites by the
same first-failure order.  The pre-carrier retained branch supplies the
repaired-gauge representation, active supercritical scaling, retained
compact-window local critical mass, critical tightness, local \(H^1\) control,
and suitable local energy on compact cylinders.  If any of these inputs fails,
or if the selected core is not nonvanishing in the retained local window, then
the first failure is a critical-mass, exterior/leakage, rough-core,
representation, or compactness exit already closed by
\cref{paper7:lem:ambient-noncarrier-exits-close}.  If the scale path is not a
genuine retained collapse path, the branch is a finite-cost, scale-degenerate,
or bounded-window/scale-rigid exit, again a named local alternative.  Thus the
only \(J_6\)-failure case remaining after ordered routing is a represented
branch with a nonvanishing selected core and genuine retained scale collapse.
\par
Apply
\cref{paper6:thm:scale-stratification,paper5:thm:scale-collapse-alternatives}
to that remaining branch.  The finite scale-cost and finite absolute-cost
alternatives are excluded by \cref{paper6:thm:scale-cost}.  Thin drift is
routed by \cref{paper6a:thm:thin-negative-drift-closed}: persistent windows
return to the thick-window ledger, while diffuse windows either enter another
named local row or a closed retained row.  Autonomous-modulation failure,
compactness-package first failure, exterior/separated-core loss, rough-core
loss, and multicore or same-point cascade are exactly the remaining
non-retained first failures listed in
\cref{paper5:lem:retained-scale-collapse-first-failure}; by
\cref{paper5:thm:scale-collapse-first-failure-closure} they are assigned to
named local state-space rows before the retained compact scale-collapse state
is formed.  The only retained scale-collapse state left after those exits is the
retained compact scale-collapse state, and it is excluded by
\cref{paper5:thm:compact-scale-collapse-rigidity-discharged}; if the retained
state becomes scale-rigid, it is excluded by
\cref{paper3:prop:scale-rigidity-discharged}.  Thus \(J_6\)-failure is a
routed local scale-collapse branch and not a hidden input.
\end{proof}

\begin{theorem}[Ambient adjacent closure for physical Type II branches]
\label{paper7:thm:ambient-adjacent-closure-branchwise}
Every physical local Type II branch which reaches the compact-cost validation
layer satisfies adjacent-stratum closure for the compact-cost regime in the
sense of \cref{paper7:def:compact-cost-adjacent-closure}.  Equivalently, the
hypotheses of
\cref{paper7:cor:adjacent-closure-reduced-to-windowwise-routing} are theorems
on the physical local Type II state-space branch.
\end{theorem}

\begin{proof}
Non-carrier adjacent exits are routed by
\cref{paper7:lem:ambient-noncarrier-exits-close}; the local closure rows close
there, while carrier, scale, and cost rows are handled by the remaining
branchwise checks below.  The exterior,
local \(H^1\)-loss, modulation-driven recentering, scale-rigid residual, and
negative scale-drift branches appearing in the stratified moving-carrier
alternative are the named local branches routed in
\cref{paper8:thm:exterior-stratification,paper6:thm:rough-core,paper6a:thm:ac-repaired-gauge-consequences,paper3:prop:scale-rigidity-discharged,paper7:lem:ambient-negative-drift-branch}.
Carrier applicability and energy-control of the canonical cost are supplied by
\cref{paper7:lem:ambient-carrier-applicability}.  The windowwise transition
and positive-scale alternative statements are proved on the canonical
windowwise moving family by
\cref{paper7:cor:windowwise-transition-routing-holds,paper7:cor:windowwise-positive-scale-routing-holds}.
The negative-drift
alternative required by
\cref{paper7:thm:windowwise-routing-gives-stratified-routing} is
\cref{paper7:lem:ambient-negative-drift-branch}.
\par
All hypotheses of
\cref{paper7:cor:adjacent-closure-reduced-to-windowwise-routing} are therefore
verified branchwise.  Applying that corollary gives compact-cost
adjacent-stratum closure.
\end{proof}

\begin{corollary}[Ambient adjacent wrapper discharged]
\label{paper7:cor:ambient-adjacent-wrapper-discharged}
The ambient adjacent-stratum closure and carrier-applicability obligation for
the canonical retained compact-cost branch is discharged by local state-space
stratification and the canonical windowwise carrier package.  Nonretained first
failures are not removed by this corollary alone; they are assigned to the named
exit rows recorded in the final assembly.
\end{corollary}

\begin{proof}
This is \cref{paper7:thm:ambient-adjacent-closure-branchwise}.
\end{proof}

\begin{remark}[Cost divergence on the retained compact stratum]
The compact Type II cost argument gives infinite localized renormalization cost
under the compact PDE hypotheses.  Passing from that cost divergence to
exclusion from the Type II class uses the corrected localized energy inequality
constructed in
\cref{paper7:lem:monotonicity-carrier-gives-cost-to-exclusion-data}.  On the
retained compact stratum, the carrier exists by
\cref{paper7:thm:carrier-validation-retained-stratum}; failure of the carrier
test is not a retained compact branch but an adjacent stratum.
\end{remark}

\clearpage
\section{Terminal Profile Decoupling and Exterior Regularity}
\label{sec:terminal-profile-decoupling-exterior}

This section records the terminal local estimates used to remove multibubble
configurations, small residual components, and components carried at a positive
physical distance from the selected Type II core.  The statements are
local implications under the analytic inputs explicitly stated below: terminal
state-space completeness on compact windows, repaired-gauge representation,
local Caccioppoli estimates, exterior regularity, and the local compact
scale-collapse routing theorem
\cref{paper5:thm:compact-scale-collapse-rigidity-discharged}.

\subsection{Renormalized frames and active local states}

A renormalized frame consists of a center \(x_c(t)\), a scale
\(\lambda(t)>0\), and renormalized time
\[
  \frac{d\tau}{dt}=\lambda(t)^{-2}.
\]
It represents the physical solution by
\[
  V(y,\tau)=\lambda(t)u(x_c(t)+\lambda(t)y,t).
\]
When the repaired gauge is imposed, the represented velocity satisfies
\[
  \partial_\tau V+(V\cdot\nabla)V+\nabla P
  =
  \nu\Delta V
  +a(\tau)(V+y\cdot\nabla V)
  +b(\tau)\cdot\nabla V,
  \qquad \nabla\cdot V=0,
\]
where
\[
  a(\tau)=\partial_\tau\log\lambda(\tau),
\]
and \(b\) is the translation modulation coefficient determined by the
centering condition.  In a fixed frame, every other local state either remains
at bounded relative position and comparable scale, escapes to spatial infinity
in the frame variables, appears at a much larger same-point scale as a
perturbative ambient field in the innermost frame, or lies at a smaller scale
and therefore requires selecting that smaller active frame.

\begin{definition}[Active local terminal state data]
\label{paper8:def:active-critical-profile-data}
Let \(t_n\uparrow T^*\).  Active local terminal state data consist of local
state-space classes \(\mathcal A_j\), scales \(\lambda_j^n\to0\), centers
\(x_j^n\to x_j^*\), compact terminal windows \(I_j\), and representatives
\(W_j^n\) in the corresponding frames.  Each active class has a detecting
compact cylinder \(B_{R_j}\times I_j\) and a local mass floor
\[
  \limsup_{n\to\infty}
  \|W_j^n\|_{L^\infty(I_j;L^3(B_{R_j}))}
  \ge \varepsilon_{\rm loc}>0.
\]
After grouping comparable parameters based at the same physical point, distinct
classes are separated in the local state-space sense: they have different
physical points, strictly separated scales, or have already been assigned to an
adjacent local alternative.  On every compact terminal cylinder and for every
finite retained set \(J\), the local field is written
\[
  V_n=U_n^J+S_n^J,
\]
where \(U_n^J\) is the grouped retained local state and \(S_n^J\) is the
nonretained residual.  Positive local mass in \(S_n^J\) is not a new
unlabelled object: by terminal state-space completeness it is assigned either
to a named adjacent local alternative or to the retained comparable class, in
which case it must be included in \(U_n^J\).  No global profile decomposition,
Brézis--Lieb identity, or heat-flow remainder smallness is part of this data.
The corresponding local active mass floor is stated and proved in
\cref{paper8:thm:active-profile-mass-floor}.
\end{definition}

\begin{lemma}[Finite active local state count on a terminal slice]
\label{paper8:lem:finite-active-count}
Fix a compact terminal cylinder \(B_R\times I\), a selected time sequence
\(\tau_n\in I\), and assume
\[
  \sup_n\|V_n\|_{L^\infty(I;L^3(B_R))}\le M_R<\infty .
\]
Assume every active local state detected on the slice \(\tau_n\) inside
\(B_R\) has local mass at least \(\varepsilon_{\rm loc}\), and that, after
grouping comparable retained classes, the detecting balls on that slice have
overlap multiplicity at most \(N_R\).  Then the number of active local states
detected on that terminal slice is finite and satisfies
\[
  \#\mathcal J_{\rm act}(R,I)
  \le
  \left\lfloor
  \frac{N_R M_R^3}{\varepsilon_{\rm loc}^3}
  \right\rfloor .
\]
\end{lemma}

\begin{proof}
For each active state choose a detecting ball on the selected time slice,
shrunk if necessary so that the family has overlap multiplicity at most \(N_R\).
At that common detecting time,
the local mass floor gives
\[
  \int_{B_j}|V_n(y,\tau_n)|^3\,dy
  \ge \varepsilon_{\rm loc}^3+o_n(1)
\]
for the ball \(B_j\) associated with the \(j\)-th state.  Summing over any
finite active subfamily and using the overlap bound,
\[
  (\#J)\varepsilon_{\rm loc}^3
  \le
  N_R\sup_n\sup_{\tau\in I}\int_{B_R}|V_n(y,\tau)|^3\,dy
  \le N_RM_R^3 .
\]
Thus every finite active subfamily has cardinality at most
\(N_RM_R^3/\varepsilon_{\rm loc}^3\).  If there were more active states than
the displayed integer bound, a finite subfamily with one additional element
would contradict this estimate.
\end{proof}

\subsection{Same-point and separated-point reductions}

\begin{lemma}[Comparable same-point states form a compound state]
\label{paper8:lem:compound-same-point}
Let a finite set \(\mathcal C\) of active local states have the same physical
center \(x_j^*=x_*\).  Suppose that for a representative center and scale
\((x_*^n,\lambda_*^n)\),
\[
  \frac{\lambda_j^n}{\lambda_*^n}\to\rho_j\in(0,\infty),
  \qquad
  \frac{x_j^n-x_*^n}{\lambda_*^n}\to z_j\in\mathbb R^3 .
\]
Let \(W_{j,n}\) denote the local representative of the \(j\)-th state in its
own frame on compact cylinders, and assume
\[
  W_{j,n}\to W_j
  \quad\text{strongly in }L^3_{\rm loc}
\]
on the relevant terminal slice, or in \(L^\infty_{\rm loc}L^3_{\rm loc}\) on a
compact terminal window.  Write
\[
  \rho_j^n:=\frac{\lambda_j^n}{\lambda_*^n},
  \qquad
  z_j^n:=\frac{x_j^n-x_*^n}{\lambda_*^n}.
\]
Then, in the frame \((x_*^n,\lambda_*^n)\),
\[
  (\rho_j^n)^{-1}
  W_{j,n}\!\left(
  \frac{y-z_j^n}{\rho_j^n}
  \right)
  \to
  \rho_j^{-1}W_j\left(\frac{y-z_j}{\rho_j}\right)
\]
strongly in the same local topology on every compact subcylinder on which the
right-hand side is defined.  Hence the comparable class has compound local
representative
\[
  \Phi_{\mathcal C}(y)
  :=
  \sum_{j\in\mathcal C}
  \rho_j^{-1}W_j\left(\frac{y-z_j}{\rho_j}\right).
\]
If the compound representative carries the local mass floor
\[
  \|\Phi_{\mathcal C}\|_{L^3(B_R)}\ge\varepsilon_{\rm loc}
\]
on a detecting compact cylinder, it is one active local state in the selected
frame.  If no compact cylinder carries the local mass floor, it belongs to the
perturbative residual class by \cref{paper8:cor:small-profiles-perturbative}.
\end{lemma}

\begin{proof}
Set
\[
  \rho_j^n=\frac{\lambda_j^n}{\lambda_*^n},
  \qquad
  z_j^n=\frac{x_j^n-x_*^n}{\lambda_*^n}.
\]
The transformed local representative is
\[
  T_nW_{j,n}(y)
  =
  (\rho_j^n)^{-1}W_{j,n}\left(\frac{y-z_j^n}{\rho_j^n}\right),
\]
and the proposed limit is
\[
  TW_j(y)
  =
  \rho_j^{-1}W_j\left(\frac{y-z_j}{\rho_j}\right).
\]
Here \(T_{\rho,z}F(y):=\rho^{-1}F((y-z)/\rho)\).  Thus
\(T_nW_{j,n}=T_{\rho_j^n,z_j^n}W_{j,n}\) and
\(TW_j=T_{\rho_j,z_j}W_j\).
Fix a compact ball \(B_R\) in the representative frame.  For all large \(n\),
the preimage of \(B_R\) under \(y\mapsto (y-z_j^n)/\rho_j^n\) is contained in a
fixed compact ball \(B_{R'}\) in the \(j\)-th frame.  The local strong
convergence gives
\[
  \|T_{\rho_j^n,z_j^n}(W_{j,n}-W_j)\|_{L^3(B_R)}
  \le C\|W_{j,n}-W_j\|_{L^3(B_{R'})}\to0.
\]
It remains to move the fixed local limit \(W_j\).  Translation and dilation are
strongly continuous on \(L^3(B_{R'})\) after restriction to \(B_R\): prove this
first for \(C_c^\infty(B_{R'})\) by uniform convergence on compact sets, then
approximate \(W_j\) in \(L^3(B_{R'})\).  This proves the local convergence.
The same argument applied uniformly in \(\tau\) gives the
\(L^\infty_{\rm loc}L^3_{\rm loc}\) version on compact terminal windows.
Summing over the finite class gives \(\Phi_{\mathcal C}\).  The final
alternative is the local state-space criterion: a compound state with the
detection threshold is retained as active, whereas a component with no
compact-cylinder detection threshold is perturbative by the local residual
theorem.
\end{proof}

\begin{lemma}[Innermost-frame rescaling of strict outer states]
\label{paper8:lem:same-point-rescaled-form}
Suppose all active local states have the same physical center \(x_*\) and split
into distinct scale classes.  If \(\lambda_1^n\) is the smallest active scale,
then every outer scale \(\lambda_j^n\gg\lambda_1^n\) satisfies
\[
  \mu_j^n:=\lambda_1^n/\lambda_j^n\to0,
\]
and in the innermost frame the \(j\)-th outer state has the local form
\[
  W_j^n(y)
  =
  \mu_j^n
  \widetilde W_{j,n}\left(
  \frac{x_*^n-x_j^n}{\lambda_j^n}+\mu_j^n y\right),
\]
where \(\widetilde W_{j,n}\) is the representative of the \(j\)-th state in
its own outer frame.
\end{lemma}

\begin{proof}
The \(j\)-th local representative in physical coordinates is
\[
  u_j^n(x)
  =
  (\lambda_j^n)^{-1}\widetilde W_{j,n}\!\left(\frac{x-x_j^n}{\lambda_j^n}\right)
\]
on the compact part of its outer terminal frame.  Pull this back into the
innermost frame \(x=x_*^n+\lambda_1^ny\), with the
critical rescaling \(W_j^n(y)=\lambda_1^n u_j^n(x_*^n+\lambda_1^n y)\).
Substituting,
\[
  W_j^n(y)
  =
  \frac{\lambda_1^n}{\lambda_j^n}
  \widetilde W_{j,n}\!\left(\frac{x_*^n+\lambda_1^ny-x_j^n}{\lambda_j^n}\right)
  =
  \mu_j^n\widetilde W_{j,n}\!\left(
  \frac{x_*^n-x_j^n}{\lambda_j^n}+\mu_j^n y\right),
\]
which is the displayed form.  The convergence \(\mu_j^n\to0\) is the
hypothesis \(\lambda_j^n\gg\lambda_1^n\).
\end{proof}

\begin{lemma}[Local \(L^3\)-vanishing of strict outer states in the innermost frame]
\label{paper8:lem:same-point-outer-vanish}
Under the hypotheses of \cref{paper8:lem:same-point-rescaled-form}, the outer
states satisfy the following retained-frame alternative.  If
\[
  \xi_j^n:=\frac{x_*^n-x_j^n}{\lambda_j^n}
\]
stays in a compact subset of the outer frame and
\(\widetilde W_{j,n}\to W_j\) strongly in \(L^3_{\rm loc}\), then, for every
\(R<\infty\),
\[
  \|W_j^n\|_{L^3(B_R)}\to0\quad (n\to\infty).
\]
If \(\xi_j^n\) leaves every compact subset of the outer frame, the contribution
is assigned to the escaping-offset/exterior row of the local state-space
decomposition rather than to a retained strict outer state.
\end{lemma}

\begin{proof}
Assume first that \(\xi_j^n\) stays in a fixed compact set \(K\).  Choose a
compact \(K'\) containing \(K+\mu_j^nB_R\) for all large \(n\).  By the formula
in \cref{paper8:lem:same-point-rescaled-form},
\[
  \|W_j^n\|_{L^3(B_R)}^3
  =
  \int_{\xi_j^n+\mu_j^nB_R}|\widetilde W_{j,n}(z)|^3\,dz .
\]
The strong \(L^3(K')\) convergence reduces this to the same integral with
\(W_j\).  Since \(|W_j|^3\in L^1(K')\) and the sets
\(\xi_j^n+\mu_j^nB_R\) have measure tending to zero, absolute continuity of
the integral gives convergence to zero.  If \(\xi_j^n\) escapes every compact
set, then the outer state is not represented on a fixed compact outer
cylinder; this is exactly the escaping-offset/exterior boundary face in the
local state-space compactification.
\end{proof}

\begin{theorem}[Perturbative estimates for strict outer states]
\label{paper8:thm:same-point-scale-separation}
In the same setting as \cref{paper8:lem:same-point-rescaled-form}, the
strict outer states are perturbative in the innermost retained frame:
\begin{enumerate}[label=\textup{(\roman*)}]
\item finite-offset regular outer states, namely those for which
\((x_*^n-x_j^n)/\lambda_j^n\) is bounded, decompose on each fixed inner ball
into a constant drift plus an \(O((\mu_j^n)^2)\) shear remainder;
\item the constant drifts produced by item (i) are absorbed into the
translation-modulation coefficient of the innermost frame;
\item the remaining pressure, cutoff, and nonlinear interactions generated by
the outer states are perturbative in the compact scale-collapse topology
in which the local perturbation criterion
\cref{paper8:lem:local-perturbation} and the pressure-tail decoupling
\cref{paper8:lem:pressure-kernel-tail} apply.
\end{enumerate}
The escaping-offset alternative
\((x_*^n-x_j^n)/\lambda_j^n\to\infty\) is treated by
\cref{paper8:lem:same-point-outer-vanish}.
\end{theorem}

\begin{proof}
In the escaping-offset case, \cref{paper8:lem:same-point-outer-vanish} gives
local \(L^3\)-vanishing on every innermost compact cylinder, and the local
pressure-tail lemma gives the pressure and residual convergence.

It remains to consider bounded offsets.  If the outer state is not regular at
the observed offset, then some smaller offset cylinder has CKN quantity above
the threshold of \cref{paper8:thm:terminal-ckn-regularity}.  That cylinder is a
positive local state-space label at a larger same physical point scale, so the
branch is assigned to the same-point cascade/secondary-core alternative before
entering the retained innermost class.  Thus, in the retained case, every
bounded offset cylinder for the outer state is CKN-subcritical.  The CKN
criterion gives uniform \(C^1\) bounds on a smaller offset cylinder, after
subtracting the pressure mean.

On that retained regular cylinder, Taylor expansion at the offset gives a
spatially constant leading term plus a shear remainder of size \(O(\mu_j^n)\)
in \(C^1\) on every fixed inner ball.  The Navier--Stokes scaling contributes
the additional factor \(\mu_j^n\) in the velocity amplitude, so the nonconstant
velocity remainder is \(O((\mu_j^n)^2)\) locally.  The spatially constant term
is absorbed into the translation modulation coefficient.  The remaining
velocity, pressure, cutoff, and modulation terms vanish in the compact-window
topologies by
\cref{paper8:lem:local-perturbation,paper8:lem:pressure-kernel-tail}.  This is
exactly the asserted perturbative reduction.
\end{proof}

\begin{lemma}[Separated physical centers are locally invisible]
\label{paper8:lem:separated-centers-vanish}
Let \(p\ne q\) be two active physical concentration points,
\[
  x_p^*\ne x_q^*,
  \qquad
  d_{pq}:=|x_p^*-x_q^*|>0.
\]
In the \(p\)-frame define the \(q\)-state contribution
\[
  W_{q\to p}^n(y)
  :=
  \lambda_p^n(\lambda_q^n)^{-1}
  \phi^q\left(
  \frac{x_p^n+\lambda_p^ny-x_q^n}{\lambda_q^n}
  \right).
\]
If \(\phi^q\in L^3(\mathbb R^3)\), then for every \(R<\infty\),
\[
  \|W_{q\to p}^n\|_{L^3(B_R)}\to0.
\]
If the corresponding representative pressure \(\Pi^q\in L^{3/2}(\mathbb R^3)\), then
the pure pressure contribution
\[
  P_{q\to p}^n(y)
  =
  (\lambda_p^n)^2(\lambda_q^n)^{-2}
  \Pi^q\left(
  \frac{x_p^n+\lambda_p^ny-x_q^n}{\lambda_q^n}
  \right)
\]
satisfies
\[
  \|P_{q\to p}^n\|_{L^{3/2}(B_R)}\to0 .
\]
\end{lemma}

\begin{proof}
By the change of variables \(x=x_p^n+\lambda_p^ny\),
\[
  \|W_{q\to p}^n\|_{L^3(B_R)}^3
  =
  \int_{B(x_p^n,\lambda_p^nR)}
  (\lambda_q^n)^{-3}
  \left|\phi^q\left(\frac{x-x_q^n}{\lambda_q^n}\right)\right|^3\,dx .
\]
Set \(z=(x-x_q^n)/\lambda_q^n\).  Then
\[
  \|W_{q\to p}^n\|_{L^3(B_R)}^3
  =
  \int_{B(c_n,\rho_n)}|\phi^q(z)|^3\,dz,
\]
where
\[
  c_n=\frac{x_p^n-x_q^n}{\lambda_q^n},
  \qquad
  \rho_n=\frac{\lambda_p^nR}{\lambda_q^n}.
\]
Since \(x_p^n\to x_p^*\), \(x_q^n\to x_q^*\), and
\(\lambda_q^n\to0\),
\[
  |c_n|\sim d_{pq}/\lambda_q^n\to\infty.
\]
Moreover
\[
  \frac{\rho_n}{|c_n|}
  =
  \frac{\lambda_p^nR}{|x_p^n-x_q^n|}
  \to0.
\]
Thus the balls \(B(c_n,\rho_n)\) escape to spatial infinity: for every
\(A<\infty\), they are contained in \(\{|z|>A\}\) for all sufficiently large
\(n\).  Since \(\phi^q\in L^3\),
\[
  \int_{|z|>A}|\phi^q(z)|^3\,dz\to0
  \qquad(A\to\infty),
\]
and hence \(\|W_{q\to p}^n\|_{L^3(B_R)}\to0\).

For the pressure term the same change of variables gives
\[
  \|P_{q\to p}^n\|_{L^{3/2}(B_R)}^{3/2}
  =
  \int_{B(c_n,\rho_n)}|\Pi^q(z)|^{3/2}\,dz.
\]
The same escaping-ball argument and \(\Pi^q\in L^{3/2}\) imply the stated
pressure convergence.
\end{proof}

\begin{definition}[Multipoint decoupling conditions]
\label{paper8:def:multipoint-decoupling}
For each active point and selected frame, the required multipoint decoupling
conditions are:
\begin{enumerate}[label=\textup{(\roman*)}]
\item local velocity decoupling in the selected frame;
\item pressure decoupling modulo time-dependent constants in
\(L^1_\tau H^{-1}_y\) on compact cylinders;
\item localization errors from spatial cutoffs vanish in
\(L^1_\tau H^{-1}_y\);
\item repaired-gauge modulation coefficients remain bounded after exterior
ambient terms are absorbed;
\item each active point has positive finite critical mass in its own frame;
\item after states centered at other physical points are treated as
perturbative exterior terms, the selected frame is tested against
\cref{paper5:def:retained-compact-scale-collapse-state}: either the first
failed retained-state input is routed by
\cref{paper5:lem:retained-scale-collapse-first-failure} to a named local exit,
or all retained compact scale-collapse inputs hold.
\end{enumerate}
\end{definition}

\begin{theorem}[Reduction of active multibubbles]
\label{paper8:thm:multibubble-frame-reduction}
Assume the Type II local terminal state data, the active local mass floor, the
same-point nonlinear decoupling conditions, the multipoint decoupling conditions of
\cref{paper8:def:multipoint-decoupling}.  Then no active multibubble Type II
configuration can occur.
\end{theorem}

\begin{proof}
By \cref{paper8:lem:finite-active-count}, there are only finitely many active
local states on each selected terminal slice.  Partition them by physical
concentration point \(x_\alpha^*\), and
within each physical point partition them by comparable scale class.

If a scale class contains several comparable local states,
\cref{paper8:lem:compound-same-point} combines them into one compound state.
If the compound state has no detecting compact cylinder, it is perturbative by
\cref{paper8:cor:small-profiles-perturbative}; if it is active, it is a single
active local state for that frame.

If a physical point contains strict same-point scale separation,
\cref{paper8:thm:same-point-scale-separation} and the same-point nonlinear
decoupling conditions reduce the innermost active scale to a single
repaired-gauge solution with perturbative outer interactions.  The compact
scale-collapse routing theorem
\cref{paper5:thm:compact-scale-collapse-rigidity-discharged} then rules out
the corresponding retained same-point cascade; if one of its retained
hypotheses fails,
\cref{paper5:lem:retained-scale-collapse-first-failure} assigns the branch to a
named local alternative and it is not a retained active multibubble
configuration.

It remains to consider separated physical points.  Choose any active point
\(x_\alpha^*\) and place the selected frame at its active scale and center.
By \cref{paper8:lem:separated-centers-vanish}, every state centered at a
different physical point is locally \(L^3\)-invisible in this frame.  By the
multipoint decoupling conditions, its pressure, cutoff, and modulation
interactions are also perturbative in the compact scale-collapse topology.
Therefore the selected frame contains a single active repaired-gauge solution
whose scale tends to zero.  If any retained compact scale-collapse input fails,
the multipoint decoupling condition and
\cref{paper5:lem:retained-scale-collapse-first-failure} route the branch to a
named local exit, so it is not a retained active multibubble configuration.  If
no retained input fails, the selected frame is a retained compact
scale-collapse state in the sense of
\cref{paper5:def:retained-compact-scale-collapse-state}.  Applying
\cref{paper5:thm:compact-scale-collapse-rigidity-discharged} excludes that
selected active point.  Since the chosen active point was arbitrary, no active
multibubble configuration can occur.
\end{proof}

\subsection{Local perturbation and pressure decoupling}

\begin{lemma}[Local perturbation criterion]
\label{paper8:lem:local-perturbation}
Let \(I\subset\mathbb R\) be compact and \(Q_R=B_R\times I\).  Let \(U_n\) be
a repaired-gauge solution on \(Q_{2R}\) satisfying
\[
  \|U_n\|_{L^\infty_\tau L^3_y(B_{2R})}
  +
  \|U_n\|_{L^2_\tau H^1_y(B_{2R})}
  \le C_R.
\]
Let \(S_n\) be a divergence-free perturbation on \(Q_{2R}\) with
\[
  S_n\to0
  \quad\text{in }L^\infty_\tau L^3_y(B_{2R}),
\]
\[
  U_n\otimes S_n+S_n\otimes U_n+S_n\otimes S_n
  \to0
  \quad\text{in }L^1_\tau L^2_y(B_R),
\]
\[
  \mathcal R_n
  :=
  \partial_\tau S_n-\nu\Delta S_n
  +a_n(S_n+y\cdot\nabla S_n)
  +b_n\cdot\nabla S_n
  \to0
  \quad\text{in }L^1_\tau H^{-1}_y(B_R),
\]
where \(|a_n|+|b_n|\le M_{ab}\), and suppose the pressure cross-term satisfies
\[
  \nabla\mathcal P_{{\rm cross},n}\to0
  \quad\text{in }L^1_\tau H^{-1}_y(B_R).
\]
Then replacing \(U_n+S_n\) by \(U_n\) changes the repaired-gauge
Navier--Stokes equation on \(Q_R\) by an error tending to zero in
\[
  L^1_\tau H^{-1}_y(B_R).
\]
\end{lemma}

\begin{proof}
The linear part is precisely \(\mathcal R_n\), which tends to zero by
hypothesis.  The nonlinear difference is
\[
  (U_n+S_n)\cdot\nabla(U_n+S_n)-U_n\cdot\nabla U_n
  =
  \operatorname{div}F_n,
  \qquad
  F_n:=U_n\otimes S_n+S_n\otimes U_n+S_n\otimes S_n,
\]
using \(\nabla\cdot U_n=\nabla\cdot S_n=0\) distributionally.  By the assumed
\(L^1_\tau L^2_y\) stress convergence, this divergence tends to zero in
\[
  L^1_\tau H^{-1}_y(B_R),
\]
because for \(\varphi\in H^1_0(B_R;\mathbb R^3)\),
\[
\begin{aligned}
  \left|
  \int_{B_R}
  (U_n\otimes S_n+S_n\otimes U_n+S_n\otimes S_n):\nabla\varphi\,dy
  \right|
  &\le
  \|F_n\|_{L^2(B_R)} \\
  &\quad{}\times \|\nabla\varphi\|_{L^2(B_R)}.
\end{aligned}
\]
The pressure contribution vanishes by the assumed cross-pressure convergence.
Therefore the total perturbative error tends to zero in
\[
  L^1_\tau H^{-1}_y(B_R).
\]
\end{proof}

\begin{lemma}[Static same-point outer profiles vanish locally]
\label{paper8:lem:same-point-outer-static}
Let
\[
  W_n(y)=\mu_n\phi(\xi_n+\mu_n y),
  \qquad
  \mu_n\to0,
  \qquad
  \phi\in L^3(\mathbb R^3).
\]
Then, for every \(R<\infty\),
\[
  \|W_n\|_{L^3(B_R)}\to0.
\]
\end{lemma}

\begin{proof}
Changing variables \(z=\xi_n+\mu_n y\),
\[
  \|W_n\|_{L^3(B_R)}^3
  =
  \int_{B(\xi_n,\mu_nR)}|\phi(z)|^3\,dz.
\]
The measure of \(B(\xi_n,\mu_nR)\) tends to zero.  Since
\(|\phi|^3\in L^1(\mathbb R^3)\), absolute continuity of the integral gives
\[
  \int_{B(\xi_n,\mu_nR)}|\phi(z)|^3\,dz\to0,
\]
uniformly in \(\xi_n\).  Hence \(\|W_n\|_{L^3(B_R)}\to0\).
\end{proof}

\begin{lemma}[Pressure decoupling from kernel tails]
\label{paper8:lem:pressure-kernel-tail}
Let \(\Gamma_{ij}\) be the Calder\'on--Zygmund kernel for
\[
  P=\mathcal R_i\mathcal R_j(F_{ij}).
\]
Assume
\[
  F_n\to0\quad\text{in }L^1_\tau L^2_y(B_A)
  \quad\text{for every fixed }A,
\]
and
\[
  \sup_n\|F_n\|_{L^\infty_\tau L^{3/2}_y(\mathbb R^3)}<\infty.
\]
Then the pressure generated by \(F_n\) satisfies, for every \(R<\infty\),
\[
  P_n-c_{n,R}\to0
  \quad\text{in }L^1_\tau L^2_y(B_R)
\]
for suitable constants \(c_{n,R}(\tau)\), and consequently
\[
  \nabla P_n\to0
  \quad\text{in }L^1_\tau H^{-1}_y(B_R).
\]
\end{lemma}

\begin{proof}
Fix \(A>2R\) and write
\[
  F_n^{\rm loc}=F_n{\bf 1}_{B_A},
  \qquad
  F_n^{\rm far}=F_n{\bf 1}_{\mathbb R^3\setminus B_A}.
\]
For the local part, Calder\'on--Zygmund boundedness gives
\[
  \|P_n^{\rm loc}\|_{L^1_\tau L^2_y(B_A)}
  \le
  C_A\|F_n\|_{L^1_\tau L^2_y(B_A)}
  \to0.
\]
For the far part define
\[
  c_{n,R}(\tau)=
  \int_{|z|>A}\Gamma_{ij}(-z)F_{n,ij}(z,\tau)\,dz,
\]
first for compactly supported approximations and then by passage to the
\(L^{3/2}\) limit.  For \(y\in B_R\) and \(|z|>A>2R\),
\[
  |\Gamma_{ij}(y-z)-\Gamma_{ij}(-z)|
  \le C_R|z|^{-4}.
\]
Hence
\[
  |P_n^{\rm far}(y,\tau)-c_{n,R}(\tau)|
  \le
  C_R\int_{|z|>A}|z|^{-4}|F_n(z,\tau)|\,dz.
\]
By Hölder's inequality with exponents \(3/2\) and \(3\),
\[
  \int_{|z|>A}|z|^{-4}|F_n|\,dz
  \le
  \|F_n(\tau)\|_{L^{3/2}}
  \left(\int_{|z|>A}|z|^{-12}\,dz\right)^{1/3}
  \le
  CA^{-3}\|F_n(\tau)\|_{L^{3/2}}.
\]
Therefore
\[
  \|P_n^{\rm far}-c_{n,R}\|_{L^1_\tau L^\infty_y(B_R)}
  \le
  C_{R,I}A^{-3}
  \sup_n\|F_n\|_{L^\infty_\tau L^{3/2}_y}.
\]
Since \(B_R\) has finite measure, this also controls the far part in
\(L^1_\tau L^2_y(B_R)\).  Taking \(n\to\infty\) in the local term and then
\(A\to\infty\) in the far term gives
\[
  P_n-c_{n,R}\to0
  \quad\text{in }L^1_\tau L^2_y(B_R).
\]
For \(\varphi\in H^1_0(B_R;\mathbb R^3)\),
\[
  |\langle\nabla(P_n-c_{n,R}),\varphi\rangle|
  \le
  \left|\int_{B_R}(P_n-c_{n,R})\nabla\cdot\varphi\,dy\right|
  \le
  \|P_n-c_{n,R}\|_{L^2(B_R)}
  \|\nabla\varphi\|_{L^2(B_R)}.
\]
Thus \(\nabla P_n\to0\) in \(L^1_\tau H^{-1}_y(B_R)\).
\end{proof}

\begin{definition}[Terminal nonlinear local-state decoupling]
\label{paper8:def:terminal-nonlinear-decoupling}
On every compact terminal-frame cylinder \(B_R\times I\), let \(U_n\) be the
retained active local-state evolution and let \(S_n\) be the nonretained local
residual.  Terminal nonlinear local-state decoupling means:
\begin{enumerate}[label=\textup{(\roman*)}]
\item \(\nabla\cdot S_n=0\) distributionally;
\item \(S_n\to0\) in \(L^\infty_\tau L^3_y(B_{2R})\);
\item
\[
  U_n\otimes S_n+S_n\otimes U_n+S_n\otimes S_n
  \to0
  \quad\text{in }L^1_\tau L^2_y(B_R);
\]
\item
\[
  \partial_\tau S_n-\nu\Delta S_n
  +a_n(S_n+y\cdot\nabla S_n)
  +b_n\cdot\nabla S_n
  \to0
  \quad\text{in }L^1_\tau H^{-1}_y(B_R);
\]
\item pressure sources involving at least one factor of \(S_n\) satisfy
\cref{paper8:lem:pressure-kernel-tail}, or directly satisfy
\[
  \nabla\mathcal P_{{\rm cross},n}\to0
  \quad\text{in }L^1_\tau H^{-1}_y(B_R);
\]
\item cutoff errors introduced in isolating the selected frame vanish in
\(L^1_\tau H^{-1}_y(B_R)\);
\item the effective modulation coefficients remain bounded;
\item after subtracting \(S_n\), the retained solution has positive finite
critical mass, repaired-gauge representation, pressure reconstruction,
Caccioppoli regularity, compactness, tightness, and the compact
scale-collapse limiting inputs.
\end{enumerate}
\end{definition}

\begin{theorem}[Same-point and multipoint decoupling from terminal decoupling]
\label{paper8:thm:terminal-decoupling-implies-local-decoupling}
If terminal nonlinear local-state decoupling holds in the sense of
\cref{paper8:def:terminal-nonlinear-decoupling}, then both same-point
scale-decoupling and separated-point frame-decoupling hold.
\end{theorem}

\begin{proof}
For same-point active local states, group comparable scales by
\cref{paper8:lem:compound-same-point}.  For strict scale separation choose the
innermost active scale as the selected terminal frame.  By
\cref{paper8:lem:same-point-outer-static}, every outer scale that remains in the
selected frame is locally invisible in velocity, unless it was already grouped
as comparable or routed to an adjacent state.  Terminal nonlinear local-state
decoupling supplies the
nonlinear-stress, linear-residual, pressure-source, localization, modulation,
and compact scale-collapse admissibility conditions on every compact cylinder.
\Cref{paper8:lem:local-perturbation} then shows that the outer components and
residual alter the selected equation only by an error tending to zero in
\(L^1_\tau H^{-1}_y(B_R)\).  Thus the same-point configuration reduces to one
retained compact scale-collapse solution with perturbative interactions.

For separated physical points, choose an active physical point \(x_p^*\) and
its selected frame.  For every local state centered at \(x_q^*\ne x_p^*\),
\cref{paper8:lem:separated-centers-vanish} gives local velocity convergence
to zero in the \(p\)-frame.  If a pure profile pressure is represented in
\(L^2\), the same escaping-ball calculation gives local \(L^2\) pressure
convergence and therefore \(H^{-1}\) convergence of the gradient.  With only
\(L^{3/2}\)-pressure, the mixed pressure sources are controlled by
\cref{paper8:lem:pressure-kernel-tail} through the \(L^2\)-strength hypotheses
included in terminal nonlinear local-state decoupling.  The same decoupling
statement supplies localization-error convergence, bounded modulation, and
compact scale-collapse admissibility.  Applying
\cref{paper8:lem:local-perturbation} shows that separated physical-point
states change the selected equation only by an error tending to zero in
\(L^1_\tau H^{-1}_{\rm loc}\).  This proves separated-point decoupling.
\end{proof}

\subsection{Terminal local-state decoupling}

\begin{definition}[Terminal active local frame]
\label{paper8:def:terminal-active-frame}
Partition active local states by physical point \(x_j^n\to x_\alpha^*\).  At a
fixed physical point, group comparable scales.  For a comparable-scale class
\(\mathcal C\), a representative frame
\((x_{\mathcal C}^n,\lambda_{\mathcal C}^n)\) has compound local representative
\[
  \Phi_{\mathcal C}(y)=
  \sum_{j\in\mathcal C}
  \rho_j^{-1}W_j\left(\frac{y-z_j}{\rho_j}\right),
\]
where \(W_j\) is the local representative of the \(j\)-th state in its own
frame, transported as in \cref{paper8:lem:compound-same-point}.
The class \(\mathcal C\) is terminal at \(x_\alpha^*\) if no active class at
the same physical point has strictly smaller scale:
\[
  \lambda_{\mathcal D}^n/\lambda_{\mathcal C}^n\not\to0
  \quad\text{for every active class }\mathcal D.
\]
The terminal frame is
\[
  y=\frac{x-x_{\mathcal C}^n}{\lambda_{\mathcal C}^n},
  \qquad
  \tau=\int^t(\lambda_{\mathcal C}(s))^{-2}\,ds,
\]
and
\[
  V_n(y,\tau)=
  \lambda_{\mathcal C}^n
  u(x_{\mathcal C}^n+\lambda_{\mathcal C}^ny,
  t_n+(\lambda_{\mathcal C}^n)^2\tau).
\]
The retained solution \(U_n\) is the repaired-gauge nonlinear evolution of
\(\mathcal C\) in this frame, and \(S_n:=V_n-U_n\) is the nonretained
remainder field.
\end{definition}

\begin{theorem}[Terminal windowed state-space completeness]
\label{paper8:thm:terminal-windowed-profile-completeness}
For a sequence \(f_n\) and a frame \((x_n,\lambda_n)\), define
\[
  \mathfrak m_R(f_n;x_n,\lambda_n)
  :=
  \limsup_{n\to\infty}
  \left\|
  \lambda_nf_n(x_n+\lambda_n\cdot)
  \right\|_{L^3(B_R)}
\]
and
\[
  \mathfrak M_{R,I}(f_n;x_n,\lambda_n)
  :=
  \limsup_{n\to\infty}
  \sup_{\tau\in I}
  \left\|
  \lambda_nf_n(x_n+\lambda_n\cdot,t_n+\lambda_n^2\tau)
  \right\|_{L^3(B_R)}.
\]
If \(\mathfrak M_{R,I}(f_n;x_n,\lambda_n)>0\) for a sequence left in the
remainder, then after passing to a subsequence the local state-space
stratification assigns that residual mass either to a named adjacent local
alternative or to the retained comparable class.  In the latter case the
component must have been grouped into the retained terminal class before
selecting the frame.  Hence no positive nonretained residual mass remains on a
retained terminal branch.
\end{theorem}

\begin{proof}
This is the compact-window form of the local residual exhaustion theorem
\cref{paper6a:thm:local-residual-exhaustion}.  If
\(\mathfrak M_{R,I}(f_n;x_n,\lambda_n)>0\), then a subsequence and a smaller
compact terminal cylinder carry positive local \(L^3\) mass.  The local
residual-exhaustion theorem then assigns the resulting positive local defect to
the first closed local row.  An adjacent label is handled by its local closure
theorem and is not part of the retained terminal branch.  A retained
comparable label must already have been grouped into the selected terminal
class by the terminal frame selection rule.  Thus positive nonretained
residual mass cannot persist on a retained terminal branch.
\end{proof}

\begin{remark}[Status of \cref{paper8:thm:terminal-windowed-profile-completeness}]
\label{paper8:rem:terminal-completeness-replacement}
This theorem is deliberately local.  It is a compact-window state-space
exhaustion statement, not a terminal global profile theorem.  Its proof uses the
local compactness, closed-label, and residual-exhaustion lemmas in the terminal
state-space discharge package; it does not use a whole-space profile
decomposition or a heat-flow remainder estimate.
\end{remark}

\begin{lemma}[Terminal local velocity decoupling]
\label{paper8:lem:terminal-velocity-decoupling}
Let \(\mathcal C\) be a terminal active class and \(S_n\) be the nonretained
remainder field in the \(\mathcal C\)-frame.  Then for every compact cylinder
\(B_R\times I\),
\[
  S_n\to0
  \quad\text{in }L^\infty_\tau L^3_y(B_R).
\]
\end{lemma}

\begin{proof}
Assume otherwise.  Then there exist \(R<\infty\), a compact interval \(I\),
\(\delta>0\), a subsequence, and times \(\tau_n\in I\) such that
\[
  \|S_n(\tau_n)\|_{L^3(B_R)}\ge\delta.
\]
In physical variables, after the retained terminal class has been subtracted,
the remainder carries at least \(\delta\) critical \(L^3\)-mass inside a ball
of radius \(O(\lambda_{\mathcal C}^n)\), centered at
\[
  x_{\mathcal C}^n+O(\lambda_{\mathcal C}^n),
\]
at time
\[
  t_n+(\lambda_{\mathcal C}^n)^2\tau_n.
\]
By terminal windowed state-space completeness, this nonvanishing mass has a
first local state-space label.  If the label is an adjacent alternative, the
branch has left the retained terminal class and is handled by the corresponding
local closure theorem, contradicting that it remains in \(S_n\) on the retained
branch.  If the label is the retained comparable class, then it was already
grouped into \(\mathcal C\) before the terminal frame was selected, so it cannot
remain in \(S_n\).  This is again a contradiction.

Thus no such \(\delta\) exists.  The same argument covers same-point
larger-scale and separated-point components: any nonzero mass seen in the
terminal frame is either a named adjacent local state or part of the retained
comparable class, contradicting the terminal partition.
\end{proof}

\begin{lemma}[Uniform terminal \(H^1\) bounds]
\label{paper8:lem:terminal-h1-bounds}
Assume the repaired-gauge representation and the local Caccioppoli estimates
hold on every compact terminal cylinder \(B_{2R}\times I\).  Then
\[
  \sup_n
  \|V_n\|_{L^\infty_\tau L^3_y(B_{2R})}
  +
  \sup_n
  \|V_n\|_{L^2_\tau H^1_y(B_{2R})}
  \le C_R,
\]
\[
  \sup_n
  \|U_n\|_{L^\infty_\tau L^3_y(B_{2R})}
  +
  \sup_n
  \|U_n\|_{L^2_\tau H^1_y(B_{2R})}
  \le C_R,
\]
and, since \(S_n=V_n-U_n\),
\[
  \sup_n
  \|S_n\|_{L^\infty_\tau L^3_y(B_{2R})}
  +
  \sup_n
  \|S_n\|_{L^2_\tau H^1_y(B_{2R})}
  \le C_R.
\]
In particular, by Sobolev embedding on \(B_{2R}\),
\[
  U_n,S_n
  \quad\text{are bounded in}\quad
  L^2_\tau L^6_y(B_{2R}).
\]
\end{lemma}

\begin{proof}
The first estimate is the repaired-gauge Caccioppoli estimate for \(V_n\) on
the compact terminal cylinder.  The retained solution \(U_n\) is itself a
repaired-gauge Navier--Stokes solution with positive finite critical mass and
bounded modulation on the same compact cylinder, so the same estimate applies
to \(U_n\).  Since \(S_n=V_n-U_n\), the displayed bound for \(S_n\) follows by
the triangle inequality in the stated norms.  The \(L^2_\tau L^6_y\) bound
follows from the Sobolev embedding \(H^1(B_{2R})\hookrightarrow L^6(B_{2R})\).
\end{proof}

\begin{lemma}[Nonlinear stress decoupling in terminal frames]
\label{paper8:lem:terminal-stress-decoupling}
Let
\[
  F_n:=U_n\otimes S_n+S_n\otimes U_n+S_n\otimes S_n.
\]
Then, for every compact \(B_R\times I\),
\[
  F_n\to0
  \quad\text{in }L^1_\tau L^2_y(B_R).
\]
\end{lemma}

\begin{proof}
By \cref{paper8:lem:terminal-velocity-decoupling},
\[
  \|S_n\|_{L^\infty_\tau L^3_y(B_R)}\to0.
\]
Using Hölder's inequality in space and time,
\[
  \|U_n\otimes S_n\|_{L^1_\tau L^2_y(B_R)}
  \le
  \|S_n\|_{L^\infty_\tau L^3_y(B_R)}
  \|U_n\|_{L^1_\tau L^6_y(B_R)}
\]
and
\[
  \|U_n\|_{L^1_\tau L^6_y(B_R)}
  \le
  |I|^{1/2}\|U_n\|_{L^2_\tau L^6_y(B_R)}
  \le C_{R,I}.
\]
Thus \(U_n\otimes S_n\to0\) in \(L^1_\tau L^2_y(B_R)\), and the same estimate
gives \(S_n\otimes U_n\to0\).  Finally,
\[
  \|S_n\otimes S_n\|_{L^1_\tau L^2_y(B_R)}
  \le
  \|S_n\|_{L^\infty_\tau L^3_y(B_R)}
  \|S_n\|_{L^1_\tau L^6_y(B_R)}
  \to0,
\]
because \(\|S_n\|_{L^1_\tau L^6_y(B_R)}\le C_{R,I}\).  Summing the three
terms proves the local convergence.
\end{proof}

\begin{lemma}[Quantitative nested-ball harmonic tail]
\label{paper8:lem:nested-ball-harmonic-tail}
Fix radii \(0<r<R<R_1<R_2\) and \(1<p<\infty\).  Let
\(\theta\in C_c^\infty(B_{R_2}\setminus \overline{B_{R_1}})\), extend
\(\theta F\) by zero to \(\mathbb R^3\), and set
\[
  P^{\rm ann}=\mathcal R_i\mathcal R_j(\theta F_{ij}),
  \qquad
  H^{\rm ann}=P^{\rm ann}-(P^{\rm ann})_{B_{R_1}}
  \quad\hbox{on }B_{R_1}.
\]
Then \(H^{\rm ann}\) is harmonic and mean-zero on \(B_{R_1}\).  Moreover
there is a constant \(C_p(r,R,R_1,R_2,\theta)>0\), depending only on \(p\), the
nested balls and the cutoff, such that
\begin{enumerate}[label=\textup{(\roman*)}]
\item \emph{Interior derivative bound.} For every harmonic
mean-zero \(H\) on \(B_{R_1}\) and every \(R<R_1\),
\[
  \sup_{B_R}\bigl(|H|+R_1|\nabla H|\bigr)
  \le C_p(R,R_1)\,\|H\|_{L^p(B_{R_1})}.
\]
\item \emph{Annular harmonic-source bound.} The harmonic part generated by the
source in the fixed annulus satisfies
\[
  \|H^{\rm ann}\|_{L^p(B_R)}
  \le
  C_p(r,R,R_1,R_2,\theta)\,
  \|F\|_{L^p(B_{R_2}\setminus B_{R_1})}.
\]
\item \emph{Quantitative converse for annular tails.} If
\(\|H^{\rm ann}\|_{L^p(B_r)}\ge\eta>0\), then
\[
  \|F\|_{L^p(B_{R_2}\setminus B_{R_1})}
  \ge C_p(r,R,R_1,R_2,\theta)^{-1}\eta .
\]
\end{enumerate}
This lemma controls only the annular harmonic contribution generated by
sources in the displayed fixed annulus.  Harmonic pressure coming from outside
\(B_{R_2}\), or from a failure of the pressure reconstruction itself, is not
silently bounded by this lemma and is instead read as a pressure/exterior
state-space alternative.
\end{lemma}

\begin{proof}
Item (i) is the standard interior derivative estimate for harmonic
functions: by the mean-value property and Cauchy estimates on
\(B_{(R+R_1)/2}\), \(|H|+R_1|\nabla H|\) on \(B_R\) is controlled by the
\(L^p\) mass of \(H\) on \(B_{R_1}\) with constant depending only on
\(p\) and \(R/R_1\).

For item (ii), the source \(\theta F\) is supported outside \(B_{R_1}\), so
\(-\Delta P^{\rm ann}=0\) on \(B_{R_1}\).  Hence \(H^{\rm ann}\) is harmonic
there.  Calder\'on--Zygmund boundedness on \(\mathbb R^3\) gives
\[
  \|P^{\rm ann}\|_{L^p(\mathbb R^3)}
  \le C_p\|\theta F\|_{L^p(\mathbb R^3)}
  \le C_p(\theta)\|F\|_{L^p(B_{R_2}\setminus B_{R_1})},
\]
and subtracting the mean on \(B_{R_1}\) changes the \(L^p(B_R)\) norm by a
geometric constant.

For item (iii), this is the contrapositive of item (ii): if the
annular source were arbitrarily small, then item (ii) would force
\(\|H^{\rm ann}\|_{L^p(B_R)}\), and hence
\(\|H^{\rm ann}\|_{L^p(B_r)}\), to be arbitrarily small.
\end{proof}

\begin{lemma}[Harmonic cross pressure routes to pressure/exterior alternative]
\label{paper8:lem:harmonic-cross-pressure-routing}
Fix nested balls \(B_R\Subset B_{R_1}\Subset B_{R_2}\), \(1<p<\infty\), and
\(1\le q<\infty\).  Let
\[
  P_{{\rm cross},n}
  =
  P_{{\rm cross},n}^{\rm loc}+H_{n,R_1}+c_{n,R_1}(\tau)
\]
be the local pressure reconstruction on \(B_{R_1}\), where
\[
  -\Delta P_{{\rm cross},n}^{\rm loc}
  =
  \partial_i\partial_j\bigl(\zeta(F_n)_{ij}\bigr),
  \qquad
  \zeta\equiv1\text{ on }B_{R_1},\quad
  \operatorname{supp}\zeta\Subset B_{R_2},
\]
and \(H_{n,R_1}\) is harmonic with zero spatial mean on \(B_{R_1}\).  Suppose
\[
  F_n\to0\quad\text{in }L^q_\tau L^p_y(B_{R_2}).
\]
If \(H_{n,R_1}\) does not converge to zero in \(L^q_\tau L^p_y(B_R)\), then,
after passing to a subsequence, the branch enters the ordered
pressure/exterior-leakage alternative.  Consequently, on any retained terminal
branch, \(H_{n,R_1}\to0\) in \(L^q_\tau L^p_y(B_R)\).
\end{lemma}

\begin{proof}
The local part tends to zero by Calder\'on--Zygmund estimates:
\[
  \|P_{{\rm cross},n}^{\rm loc}\|_{L^q_\tau L^p_y(B_{R_1})}
  \le C\|F_n\|_{L^q_\tau L^p_y(B_{R_2})}\to0 .
\]
Hence any non-vanishing pressure oscillation modulo functions of time is
carried by the normalized harmonic part.  If
\[
  \limsup_{n\to\infty}\|H_{n,R_1}\|_{L^q_\tau L^p_y(B_R)}>0,
\]
then choose a finite smooth partition of unity
\(\{\theta_\ell\}_\ell\subset C_c^\infty(B_{R_2}\setminus\overline{B_{R_1}})\)
for the recorded annular part of the pressure atlas.  For each \(\ell\), the
corresponding annular harmonic contribution
\[
  H_{n,\ell}^{\rm ann}
  :=
  \mathcal R_i\mathcal R_j(\theta_\ell(F_n)_{ij})
  -\fint_{B_{R_1}}\mathcal R_i\mathcal R_j(\theta_\ell(F_n)_{ij})\,dy
\]
vanishes in \(L^q_\tau L^p_y(B_R)\) by
\cref{paper8:lem:nested-ball-harmonic-tail}, because
\[
  \|F_n\|_{L^q_\tau L^p_y(B_{R_2}\setminus B_{R_1})}\to0 .
\]
Summing over the finite partition shows that the whole harmonic contribution
generated by sources recorded in the fixed annulus
\(B_{R_2}\setminus B_{R_1}\) vanishes.

By construction of the local pressure atlas, after the local
Calder\'on--Zygmund piece and the recorded annular pieces have been removed, the
remaining mean-zero harmonic part has only two possible origins: either the
pressure gauge/reconstruction fails on the selected window, or it is a
far-field pressure tail generated outside the fixed ball \(B_{R_2}\).  The
former is one of the ordered pressure rows.  In the latter case,
\cref{paper8:thm:exterior-stratification} routes any exterior carrier to an
exterior or separated-core concentration alternative; if no such exterior
carrier exists, \cref{paper8:lem:exterior-pressure-vanishes} gives convergence
of the far-field pressure modulo functions of time on \(B_R\), contradicting
the positive lower bound for \(H_{n,R_1}\).  Therefore every non-vanishing
harmonic cross pressure is a named pressure/exterior-leakage alternative.  A
retained terminal branch is the complementary case in the ordered validation,
so the harmonic part vanishes there.
\end{proof}

\begin{lemma}[Terminal pressure decoupling]
\label{paper8:lem:terminal-pressure-decoupling}
Fix nested balls \(B_R\Subset B_{R_1}\Subset B_{R_2}\) and let
\(\zeta\in C_c^\infty(B_{R_2})\) be identically one on \(B_{R_1}\).  Let
\[
  P_{{\rm cross},n}^{\rm loc}
  :=
  \mathcal R_i\mathcal R_j\bigl(\zeta (F_n)_{ij}\bigr)
\]
be the local Calder\'on--Zygmund pressure generated by the cross stress.  Suppose
the terminal pressure reconstruction is used on \(B_{R_2}\), so that on
\(B_{R_1}\)
\[
  P_{{\rm cross},n}
  =
  P_{{\rm cross},n}^{\rm loc}+H_{n,R_1}+c_{n,R_1}(\tau),
\]
where \(H_{n,R_1}(\cdot,\tau)\) is harmonic and normalized by
\(\fint_{B_{R_1}}H_{n,R_1}(\cdot,\tau)\,dy=0\).  On a retained terminal branch,
\[
  P_{{\rm cross},n}-c_{n,R}
  \to0
  \quad\text{in }L^1_\tau L^2_y(B_R),
\]
and
\[
  \nabla P_{{\rm cross},n}
  \to0
  \quad\text{in }L^1_\tau H^{-1}_y(B_R).
\]
\end{lemma}

\begin{proof}
\Cref{paper8:lem:terminal-stress-decoupling} gives
\[
  F_n\to0
  \quad\text{in }L^1_\tau L^2_y(B_{R_2})
\]
on every compact terminal cylinder.  Calder\'on--Zygmund boundedness gives
\[
  \|P_{{\rm cross},n}^{\rm loc}\|_{L^1_\tau L^2_y(B_{R_1})}
  \le
  C_{R_1,R_2}
  \|F_n\|_{L^1_\tau L^2_y(B_{R_2})}
  \to0 .
\]
For the harmonic part, apply
\cref{paper8:lem:harmonic-cross-pressure-routing} with \((q,p)=(1,2)\).  A non-vanishing harmonic
cross pressure is a pressure/exterior-leakage alternative and therefore is not
present on a retained terminal branch.  Hence
\(H_{n,R_1}\to0\) in \(L^1_\tau L^2_y(B_R)\).  Adding the local part proves
\(P_{{\rm cross},n}-c_{n,R}\to0\) in
\(L^1_\tau L^2_y(B_R)\).

Finally, for \(\varphi\in H^1_0(B_R;\mathbb R^3)\),
\[
  |\langle\nabla(P_{{\rm cross},n}-c_{n,R}),\varphi\rangle|
  \le
  \|P_{{\rm cross},n}-c_{n,R}\|_{L^2(B_R)}
  \|\nabla\varphi\|_{L^2(B_R)}.
\]
Integrating in time gives
\(\nabla P_{{\rm cross},n}\to0\) in \(L^1_\tau H^{-1}_y(B_R)\).
\end{proof}

\begin{lemma}[Linear residual decoupling]
\label{paper8:lem:terminal-linear-residual}
Let \(V_n=U_n+S_n\) solve
\[
  \partial_\tau V_n+(V_n\cdot\nabla)V_n+\nabla P_n
  =
  \nu\Delta V_n
  +a_n(V_n+y\cdot\nabla V_n)
  +b_n\cdot\nabla V_n.
\]
Suppose \(U_n\) satisfies, in the same repaired gauge,
\[
  \partial_\tau U_n+(U_n\cdot\nabla)U_n+\nabla P_{U,n}
  =
  \nu\Delta U_n
  +a_n(U_n+y\cdot\nabla U_n)
  +b_n\cdot\nabla U_n+e_n,
\]
with
\[
  e_n\to0
  \quad\text{in }L^1_\tau H^{-1}_y(B_R)
\]
on compact terminal cylinders.  Then
\[
  \partial_\tau S_n-\nu\Delta S_n
  -a_n(S_n+y\cdot\nabla S_n)
  -b_n\cdot\nabla S_n
  \to0
  \quad\text{in }L^1_\tau H^{-1}_y(B_R).
\]
Equivalently, the same conclusion holds with the opposite sign convention for
the modulation terms.
\end{lemma}

\begin{proof}
Subtracting the \(U_n\)-equation from the \(V_n\)-equation gives
\[
  \partial_\tau S_n-\nu\Delta S_n
  =
  -\operatorname{div}F_n
  -\nabla P_{{\rm cross},n}
  -e_n
  +a_n(S_n+y\cdot\nabla S_n)
  +b_n\cdot\nabla S_n.
\]
Equivalently,
\[
  \partial_\tau S_n-\nu\Delta S_n
  -a_n(S_n+y\cdot\nabla S_n)
  -b_n\cdot\nabla S_n
  =
  -\operatorname{div}F_n-\nabla P_{{\rm cross},n}-e_n.
\]
By \cref{paper8:lem:terminal-stress-decoupling},
\[
  F_n\to0\quad\text{in }L^1_\tau L^2_y(B_R),
\]
hence
\[
  \operatorname{div}F_n\to0
  \quad\text{in }L^1_\tau H^{-1}_y(B_R).
\]
By \cref{paper8:lem:terminal-pressure-decoupling},
\[
  \nabla P_{{\rm cross},n}\to0
  \quad\text{in }L^1_\tau H^{-1}_y(B_R).
\]
Together with \(e_n\to0\), this proves the displayed residual convergence.

For the opposite sign convention, the two residuals differ by
\[
  2a_n(S_n+y\cdot\nabla S_n)+2b_n\cdot\nabla S_n.
\]
By \cref{paper8:lem:terminal-velocity-decoupling} and finite measure,
\[
  S_n\to0\quad\text{in }L^\infty_\tau L^2_y(B_R).
\]
For \(\varphi\in H^1_0(B_R)\),
\[
  |\langle y\cdot\nabla S_n,\varphi\rangle|
  =
  \left|\int_{B_R}S_n(3\varphi+y\cdot\nabla\varphi)\,dy\right|
  \le
  C_R\|S_n\|_{L^2(B_R)}\|\varphi\|_{H^1(B_R)},
\]
and
\[
  |\langle b_n\cdot\nabla S_n,\varphi\rangle|
  \le
  |b_n|\,\|S_n\|_{L^2(B_R)}\|\nabla\varphi\|_{L^2(B_R)}.
\]
The term \(S_n\) itself is controlled by the same \(L^2\)-convergence.  Since
\(a_n\) and \(b_n\) are bounded, the sign-convention difference tends to zero
in \(L^1_\tau H^{-1}_y(B_R)\).
\end{proof}

\begin{lemma}[Terminal modulation compatibility]
\label{paper8:lem:terminal-modulation-compatibility}
Same-point larger-scale profiles in a terminal frame have the form
\[
  W_j^n(y,\tau)
  =
  \mu_j^n
  U^j(\theta_j^n(\tau),\xi_j^n+\mu_j^ny),
  \qquad \mu_j^n\to0.
\]
When the profile is regular at the observed offset, its leading term is a
spatially constant ambient velocity.  Absorbing this constant into the
translation coefficient replaces \(b_n\) by \(b_n-b_{{\rm amb},n}\), and the
remaining terms vanish on compact terminal cylinders.  If the offset escapes,
the entire profile is locally invisible and no absorption is needed.  In
particular,
\[
  \sup_n\sup_{\tau\in I}\bigl(|a_n(\tau)|+|b_n(\tau)|\bigr)<\infty .
\]
\end{lemma}

\begin{proof}
The local invisibility of the nonconstant remainder follows from
\cref{paper8:lem:terminal-velocity-decoupling} and the \(H^{-1}\) product
estimates in \cref{paper8:lem:terminal-linear-residual}.  Separated
physical-point profiles are locally invisible by
\cref{paper8:lem:separated-centers-vanish}; their pressures and cutoff
commutators vanish by \cref{paper8:lem:terminal-pressure-decoupling} and
\cref{paper8:lem:terminal-linear-residual}.  Therefore only the spatially
constant ambient velocity changes the translation coefficient, and after this
absorption the effective modulation coefficients remain uniformly bounded.
\end{proof}

\begin{lemma}[Terminal compact scale-collapse admissibility]
\label{paper8:lem:terminal-scale-collapse-admissibility}
The retained terminal solution \(U_n\) has positive finite critical mass,
pressure reconstruction, bounded modulation, local windowed \(H^1\) control,
critical tightness on the selected windows, and the compactness input used by
the compact scale-collapse argument.
\end{lemma}

\begin{proof}
The terminal class is active, so after the perturbative residual and
nonretained local states have been removed,
\[
  \inf_{\tau\in I}\|U_n(\tau)\|_{L^3(B_R)}
  \ge c_{\rm act}>0
\]
on the selected compact windows, after passing to the subsequences used in the
local state-space selection.  The upper bound follows from the retained local
compact-cylinder bound:
\[
  \sup_{\tau\in I}\|U_n(\tau)\|_{L^3(B_R)}\le C(M).
\]

The terminal state-space completeness also gives the required tightness.  If
\(U_n\) were not uniformly \(L^3\)-tight on the selected windows, then there
would be \(\varepsilon>0\), radii \(R_k\to\infty\), and times \(\tau_n\) such
that
\[
  \int_{|y|>R_k}|U_n(y,\tau_n)|^3\,dy\ge\varepsilon.
\]
Applying the local state-space stratification to this escaping sequence
produces either an exterior-regular component, which is removed by the exterior
regularity argument below, or an additional active local component at a
different physical point or larger same-point scale.  Such a component belongs
to \(S_n\) and is handled by terminal decoupling.  Hence no non-tight retained
terminal solution remains.

The repaired-gauge representation gives pressure reconstruction and bounded
modulation.  The local Caccioppoli estimate gives local windowed \(H^1\)
control, and the Aubin--Lions compactness layer gives the compactness input on
the selected windows.
\end{proof}

\begin{theorem}[Terminal nonlinear local-state decoupling theorem]
\label{paper8:thm:terminal-nonlinear-decoupling}
Assume local compact critical bounds, terminal windowed state-space
completeness, local perturbative residual stability, the exterior annulus
alternative, repaired-gauge representation, and local Caccioppoli regularity.
Then
terminal nonlinear local-state decoupling holds in the sense of
\cref{paper8:def:terminal-nonlinear-decoupling}.
\end{theorem}

\begin{proof}
Let \(\mathcal C\) be any terminal active compound class, let \(U_n\) be the
retained repaired-gauge nonlinear solution in its terminal frame, and set
\[
  S_n=V_n-U_n.
\]
\Cref{paper8:lem:terminal-velocity-decoupling} proves
\[
  S_n\to0
  \quad\text{in }L^\infty_\tau L^3_y(B_{2R})
\]
on every compact terminal cylinder.  Divergence-free compatibility follows
from the divergence-free repaired-gauge equation and the local terminal
decomposition.

\Cref{paper8:lem:terminal-stress-decoupling} proves
\[
  U_n\otimes S_n+S_n\otimes U_n+S_n\otimes S_n
  \to0
  \quad\text{in }L^1_\tau L^2_y(B_R).
\]
The local compact critical bounds give the \(L^\infty_\tau L^3_y\) control
required in the pressure lemma below on the cylinders under consideration.
\Cref{paper8:lem:terminal-pressure-decoupling} proves
\[
  \nabla P_{{\rm cross},n}\to0
  \quad\text{in }L^1_\tau H^{-1}_y(B_R)
\]
after subtracting spatially constant pressure terms.  \Cref{paper8:lem:terminal-linear-residual}
proves
\[
  \partial_\tau S_n-\nu\Delta S_n
  +a_n(S_n+y\cdot\nabla S_n)
  +b_n\cdot\nabla S_n
  \to0
  \quad\text{in }L^1_\tau H^{-1}_y(B_R).
\]
\Cref{paper8:lem:terminal-modulation-compatibility} gives modulation
compatibility, and \cref{paper8:lem:terminal-scale-collapse-admissibility}
gives compact scale-collapse admissibility of the retained terminal solution.
These are the items in \cref{paper8:def:terminal-nonlinear-decoupling}.
\end{proof}

\begin{corollary}[Multibubble exclusion from terminal local decoupling]
\label{paper8:cor:multibubble-exclusion-terminal}
Under the Type II local terminal data, terminal windowed state-space
completeness, local perturbative residual stability, the exterior annulus
alternative, repaired-gauge representation, and local Caccioppoli regularity,
no active multibubble Type II configuration occurs.
\end{corollary}

\begin{proof}
\Cref{paper8:thm:terminal-nonlinear-decoupling} gives terminal nonlinear
decoupling.  \Cref{paper8:thm:terminal-decoupling-implies-local-decoupling}
then supplies the same-point and separated-point decoupling hypotheses used in
\cref{paper8:thm:multibubble-frame-reduction}.  The active local mass floor
follows from the retained positive local critical-mass row and local
perturbative residual stability.  Applying
\cref{paper8:thm:multibubble-frame-reduction} gives the claimed exclusion of
active multibubble configurations.
\end{proof}

\subsection{Small local residual components}

For a compact terminal cylinder \(Q_{R,I}:=B_R\times I\), define the local
residual seminorm
\[
  \|W\|_{\mathcal X(Q_{R,I})}
  :=
  \|W\|_{L^\infty(I;L^3(B_R))}
  +
  \|W\|_{L^2(I;H^1(B_R))}.
\]
This is only a compact-window norm.  It is not a whole-space heat-flow or
global mild-solution norm.

\begin{lemma}[Local residual product estimates]
\label{paper8:lem:kato-estimates}
Let \(Q_{2R,I}=B_{2R}\times I\) be compact.  Suppose
\[
  S_n\to0\quad\text{in }L^\infty(I;L^3(B_{2R}))
\]
and
\[
  \sup_n\bigl(\|U_n\|_{\mathcal X(Q_{2R,I})}
  +\|S_n\|_{\mathcal X(Q_{2R,I})}\bigr)<\infty .
\]
Then
\[
  U_n\otimes S_n+S_n\otimes U_n+S_n\otimes S_n
  \to0
  \quad\text{in }L^1(I;L^2(B_R)).
\]
If the pressures associated with these cross stresses are reconstructed by the
local pressure decomposition, then their gradients converge to zero in
\[
  L^1(I;H^{-1}(B_R)).
\]
\end{lemma}

\begin{proof}
Sobolev embedding on \(B_{2R}\) gives uniform
\[
  U_n,S_n\quad\text{in }L^2(I;L^6(B_{2R})).
\]
Hölder's inequality in space and time yields
\[
  \|U_n\otimes S_n\|_{L^1(I;L^2(B_R))}
  \le
  \|S_n\|_{L^\infty(I;L^3(B_R))}
  \|U_n\|_{L^1(I;L^6(B_R))}\to0,
\]
because
\[
  \|U_n\|_{L^1(I;L^6(B_R))}
  \le |I|^{1/2}\|U_n\|_{L^2(I;L^6(B_R))}.
\]
The same estimate treats \(S_n\otimes U_n\), and
\[
  \|S_n\otimes S_n\|_{L^1(I;L^2(B_R))}
  \le
  \|S_n\|_{L^\infty(I;L^3(B_R))}
  \|S_n\|_{L^1(I;L^6(B_R))}\to0.
\]
For the pressure statement, apply the local pressure reconstruction to the
cross source.  The preceding \(L^1L^2\) convergence gives local convergence of
the local pressure part by Calder\'on--Zygmund estimates.  Any non-vanishing
harmonic remainder is routed by
\cref{paper8:lem:harmonic-cross-pressure-routing} with \((q,p)=(1,2)\) to the pressure/exterior
alternative; on a retained branch this alternative is absent.  Therefore the
cross-pressure gradients vanish in \(L^1H^{-1}(B_R)\).
\end{proof}

\begin{theorem}[Small local residual components are perturbative]
\label{paper8:thm:small-data-critical}
Let \(S_n\) be a nonretained terminal residual on every compact terminal
cylinder \(Q_{2R,I}\).  Assume
\[
  S_n\to0\quad\text{in }L^\infty(I;L^3(B_{2R}))
\]
and that the local compact-window bounds, local pressure reconstruction, and
bounded modulation estimates hold on \(Q_{2R,I}\).  Then \(S_n\) is
perturbative in the terminal decoupling topology of
\cref{paper8:def:terminal-nonlinear-decoupling}.
\end{theorem}

\begin{proof}
\Cref{paper8:lem:kato-estimates} gives the vanishing nonlinear-stress and
cross-pressure terms.  The repaired-gauge equation for \(V_n=U_n+S_n\) and the
equation for the retained component \(U_n\), subtracted on \(Q_{R,I}\), give the
linear residual identity of
\cref{paper8:lem:terminal-linear-residual}.  The stress and pressure terms just
proved vanish, and bounded modulation sends the remaining
\(a_n(S_n+y\cdot\nabla S_n)+b_n\cdot\nabla S_n\) terms to zero in
\(L^1(I;H^{-1}(B_R))\).  Hence all perturbative clauses in
\cref{paper8:def:terminal-nonlinear-decoupling} hold on \(Q_{R,I}\).
\end{proof}

\begin{corollary}[Small local components are not active]
\label{paper8:cor:small-profiles-perturbative}
If a nonretained component has vanishing local \(L^\infty L^3\) seminorm on
every compact terminal cylinder and satisfies the local compact-window bounds,
then it cannot define an active terminal state.
\end{corollary}

\begin{proof}
By \cref{paper8:thm:small-data-critical}, such a component contributes only
vanishing stress, pressure, modulation, and residual errors in the terminal
decoupling topology.  Active terminal states are precisely retained local
state-space labels with positive nonperturbative local mass.  Therefore a
component that is perturbative on every compact terminal cylinder is not active.
\end{proof}

\begin{theorem}[Positive local mass floor for active terminal components]
\label{paper8:thm:active-profile-mass-floor}
There exists a local threshold \(\varepsilon_{\rm loc}>0\), depending only on
the Caffarelli--Kohn--Nirenberg threshold and the finite covering constants of
the selected compact cylinders, such that every retained active terminal
component has, after passing to a subsequence, a compact terminal cylinder
\(Q_{R,I}\) with
\[
  \liminf_{n\to\infty}
  \int_I\int_{B_R}|W_n|^3\,dy\,d\tau
  \ge \varepsilon_{\rm loc}^3
\]
and
\[
  \limsup_{n\to\infty}
  \|W_n\|_{L^\infty(I;L^3(B_R))}
  \ge \varepsilon_{\rm loc}.
\]
\end{theorem}

\begin{proof}
Let \(\varepsilon_{\rm CKN}\) be the threshold in
\cref{paper8:thm:terminal-ckn-regularity}.  Fix a finite nested covering of
each retained compact terminal cylinder by CKN cylinders whose smaller
cylinders still cover the retained core.  Choose
\(\eta_{\rm det}>0\) to be \(\varepsilon_{\rm CKN}\) divided by a sufficiently
large constant depending only on this cover.

If every CKN cylinder in the cover had velocity-pressure quantity below
\(\eta_{\rm det}\), then the CKN criterion would give a uniformly smooth
bounded representative on the retained smaller cylinder.  Such a bounded
selected-window representative is not an active Type II terminal component; it
is routed to the regular/bounded-window alternative.  Hence an active retained
component has a subcylinder on which the terminal CKN quantity is at least
\(\eta_{\rm det}\).

If the pressure part of that CKN quantity carries at least half of
\(\eta_{\rm det}\) while the velocity part is arbitrarily small on all smaller
subcylinders, then the local pressure reconstruction and harmonic-tail
alternative, made explicit in
\cref{paper8:lem:harmonic-cross-pressure-routing} with \((q,p)=(3/2,3/2)\), assign the branch to the
pressure/exterior-leakage stratum.  That stratum is adjacent and is not a
retained active terminal component.  Therefore on a retained active component
the velocity part carries a fixed fraction of the detected CKN quantity on some
smaller cylinder:
\[
  \rho^{-2}\int_{I_\rho}\int_{B_\rho}|W_n|^3\,dy\,d\tau
  \ge c\,\eta_{\rm det}
\]
along a subsequence.  Since \(|I_\rho|\simeq\rho^2\), there is a time in
\(I_\rho\) with
\[
  \int_{B_\rho}|W_n(y,\tau)|^3\,dy\ge c\,\eta_{\rm det}.
\]
After rescaling this subcylinder to a fixed compact terminal cylinder, this is
the asserted spacetime lower bound and the asserted
\(L^\infty_\tau L^3_y\) lower bound with
\(\varepsilon_{\rm loc}:=(c\,\eta_{\rm det})^{1/3}\).

If no such velocity lower bound existed on any compact terminal cylinder, then
the component would have vanishing local \(L^\infty L^3\) seminorm after
subsequence extraction and would be perturbative by
\cref{paper8:cor:small-profiles-perturbative}, contradicting retained
activity.
\end{proof}

\subsection{Exterior regularity}

\begin{definition}[Exterior annuli and terminal CKN quantity]
\label{paper8:def:exterior-ckn-quantity}
Let \((x_*,T^*)\) be the selected Type II core point.  For
\(0<r<R<\infty\), set
\[
  A_{r,R}:=\overline{B_R(x_*)}\setminus B_r(x_*).
\]
For \(x_0\in\mathbb R^3\) and \(\rho>0\), write
\[
  Q_\rho^*(x_0):=
  B_\rho(x_0)\times(T^*-\rho^2,T^*).
\]
For a suitable weak solution \((u,p)\), define the terminal local
velocity-pressure quantity
\[
  \mathcal C_*(u,p;x_0,\rho)
  :=
  \rho^{-2}
  \int_{T^*-\rho^2}^{T^*}\int_{B_\rho(x_0)} |u|^3\,dx\,dt
  +
  \rho^{-2}
  \int_{T^*-\rho^2}^{T^*}\int_{B_\rho(x_0)}
  |p-(p)_{B_\rho(x_0)}(t)|^{3/2}\,dx\,dt .
\]
The pressure is measured modulo functions of time.
\end{definition}

\begin{theorem}[Terminal Caffarelli--Kohn--Nirenberg regularity]
\label{paper8:thm:terminal-ckn-regularity}
There exists \(\varepsilon_{\rm CKN}>0\) with the following property.  Let
\((u,p)\) be suitable on \(Q_\rho^*(x_0)\).  If
\[
  \mathcal C_*(u,p;x_0,\rho)\le\varepsilon_{\rm CKN},
\]
then for each integer \(k\ge0\) there are \(0<c_k<1\) and \(C_k<\infty\), independent of
\((u,p),x_0,\rho\), such that \(u\) is smooth in
\[
  B_{c_k\rho}(x_0)\times(T^*-c_k\rho^2,T^*)
\]
and, after subtracting a function of time from \(p\),
\[
  \sup_{t\in(T^*-c_k\rho^2,T^*)}
  \left(
  \rho^{k+1}\|\nabla^k u(t)\|_{L^\infty(B_{c_k\rho}(x_0))}
  +
  \rho^{k+2}\|\nabla^k p(t)\|_{L^\infty(B_{c_k\rho}(x_0))}
  \right)
  \le C_k .
\]
\end{theorem}

\begin{proof}
This is the standard endpoint form of the
Caffarelli--Kohn--Nirenberg epsilon-regularity criterion
\cite{CaffarelliKohnNirenberg1982,Lin1998}, after translating the terminal
time to \(0\) and rescaling \(Q_\rho^*(x_0)\) to a unit parabolic cylinder.
The pressure normalization by spatial averages is the usual normalization in
the suitable weak-solution criterion.  Interior parabolic bootstrapping gives
the stated higher derivative bounds on smaller cylinders.
\end{proof}

Throughout the exterior regularity subsection we fix
\[
  \varepsilon_{\rm ext}:=\frac12\varepsilon_{\rm CKN}.
\]

\begin{definition}[No exterior concentration on a fixed annulus]
\label{paper8:def:no-exterior-concentration-stratum}
Let \(0<r<R<\infty\).  A suitable Type II branch has no exterior concentration on \(A_{r,R}\) if
there exists \(\rho_*>0\) such that whenever
\[
  x_0\in A_{r,R},\qquad 0<\rho\le\rho_*,
  \qquad Q_\rho^*(x_0)\subset
  (B_{2R}(x_*)\setminus \overline{B_{r/2}(x_*)})\times(0,T^*),
\]
one has
\[
  \mathcal C_*(u,p;x_0,\rho)\le\varepsilon_{\rm ext}.
\]
\end{definition}

\begin{definition}[Exterior concentration branch]
\label{paper8:def:exterior-concentration-branch}
An exterior concentration branch on \(A_{r,R}\) is a sequence
\[
  x_n\in A_{r,R},\qquad \rho_n\downarrow0,
\]
such that
\[
  Q_{\rho_n}^*(x_n)\subset
  (B_{2R}(x_*)\setminus \overline{B_{r/2}(x_*)})\times(0,T^*)
\]
and
\[
  \mathcal C_*(u,p;x_n,\rho_n)>\varepsilon_{\rm ext}.
\]
The associated exterior rescalings are
\[
  U_n(y,s):=\rho_n u(x_n+\rho_n y,T^*+\rho_n^2s),
  \qquad
  P_n(y,s):=\rho_n^2 p(x_n+\rho_n y,T^*+\rho_n^2s),
\]
for \(s<0\).
\end{definition}

\begin{lemma}[Exterior test dichotomy]
\label{paper8:lem:exterior-test-dichotomy}
Fix \(0<r<R<\infty\).  Then either the branch has no exterior concentration
on \(A_{r,R}\), or it admits an exterior concentration branch on \(A_{r,R}\).
\end{lemma}

\begin{proof}
If the no-concentration condition fails, then for every \(m\) there are
\[
  x_m\in A_{r,R},\qquad 0<\rho_m\le m^{-1},
\]
with \(Q_{\rho_m}^*(x_m)\) contained in the enlarged annular cylinder and
\[
  \mathcal C_*(u,p;x_m,\rho_m)>\varepsilon_{\rm ext}.
\]
After discarding repetitions and passing to a subsequence, \(\rho_m\to0\).
This is precisely an exterior concentration branch.  The converse is immediate
from the definitions.
\end{proof}

\begin{theorem}[Exterior concentration stratification]
\label{paper8:thm:exterior-stratification}
Fix \(0<r<R<\infty\).  A suitable Type II branch satisfies exactly one of the
following alternatives on \(A_{r,R}\):
\begin{enumerate}[label=\textup{(\roman*)}]
\item it has no exterior concentration on \(A_{r,R}\);
\item it admits an exterior concentration branch on \(A_{r,R}\).
\end{enumerate}
The second alternative is the complementary concentration branch.  According
to the relative position and scale of the exterior concentration sequence, it
belongs to one of the exterior-profile, separated-core, multicore, noncompact
exterior, or scale-collapse alternatives.
\end{theorem}

\begin{proof}
The dichotomy is \cref{paper8:lem:exterior-test-dichotomy}.  If an exterior
concentration branch exists, then the annular terminal cylinders carry a
positive Caffarelli--Kohn--Nirenberg quantity at scales \(\rho_n\downarrow0\)
and centers a positive physical distance from the selected core.  Let
\(\lambda_n\) denote the selected core scale along the same terminal sequence.
After passing to a subsequence, the relative scale-center configuration falls
into one of the standard local concentration-compactness alternatives.  If the
exterior frame is orthogonal to the selected core frame, it gives an exterior
profile.  If only finitely many separated active centers remain at comparable
terminal scales, it gives a separated-core or multicore configuration.  If
compactness fails in the exterior frame, one obtains the noncompact exterior
alternative.  If the exterior frame develops a further unresolved scale in its
own renormalized variables, one obtains the scale-collapse alternative.  These
alternatives exhaust the subsequential configurations at positive physical
distance from the selected core.  Thus the failure of the annular regularity
test is represented by a local concentration alternative rather than by an
additional regularity hypothesis.
\end{proof}

\begin{lemma}[Exterior CKN cover]
\label{paper8:lem:exterior-ckn-cover}
Assume there is no exterior concentration on \(A_{r,R}\).  Fix an integer
\(k\ge0\).  If
\[
  A_{r',R'}\Subset A_{r,R},
\]
then there are points \(x_j\in A_{r,R}\), radii \(\rho_j>0\), and a terminal
time length \(\delta>0\) such that
\[
  A_{r',R'}\times(T^*-\delta,T^*)
  \subset
  \bigcup_j
  B_{c_k\rho_j}(x_j)\times(T^*-c_k\rho_j^2,T^*)
\]
and
\[
  \mathcal C_*(u,p;x_j,\rho_j)\le\varepsilon_{\rm ext}
\]
for every \(j\).  Here \(c_k\) is the constant from
\cref{paper8:thm:terminal-ckn-regularity}.
\end{lemma}

\begin{proof}
Since \(A_{r',R'}\Subset A_{r,R}\), there is \(d>0\) such that
\[
  B_{4d}(x)\subset B_{2R}(x_*)\setminus \overline{B_{r/2}(x_*)}
  \qquad\text{for all }x\in A_{r',R'}.
\]
Let \(\rho_*\) be supplied by the no-concentration condition and set
\(\rho_0:=\min\{d,\rho_*\}\).  By compactness, finitely many balls
\[
  B_{c_k\rho_0/2}(x_j),\qquad x_j\in A_{r,R},
\]
cover \(A_{r',R'}\).  The terminal cylinders \(Q_{\rho_0}^*(x_j)\) are
contained in the enlarged annular cylinder, so the no-concentration condition
gives the displayed smallness.  Taking
\(\delta=c_k\rho_0^2/2\) gives the spacetime cover.
\end{proof}

\begin{lemma}[Exterior pressure localization]
\label{paper8:lem:exterior-pressure-localization}
Let \(A_{r',R'}\Subset A_{r,R}\), and assume that \(u\) is bounded on a
slightly larger annular terminal cylinder.  Then on
\[
  A_{r',R'}\times(T^*-\delta,T^*)
\]
the pressure has a decomposition
\[
  p=p_{\rm loc}+p_{\rm harm}+c(t),
\]
where \(c(t)\) is a function of time, \(p_{\rm harm}\) is harmonic in space on
the smaller annulus, and for every \(q<\infty\) and every integer \(k\ge0\)
\[
  \|p_{\rm loc}\|_{L^\infty_tL^q_x(A_{r',R'})}
  +
  \|p_{\rm harm}\|_{L^\infty_tC^k_x(A_{r',R'})}
  \le C_{r',R',q,k}.
\]
\end{lemma}

\begin{proof}
Choose \(\eta\in C_c^\infty(B_{2R}(x_*)\setminus \overline{B_{r/2}(x_*)})\)
with \(\eta\equiv1\) on a slightly larger annulus containing \(A_{r',R'}\).
Define \(p_{\rm loc}\) by
\[
  -\Delta p_{\rm loc}
  =
  \partial_i\partial_j(\eta u_i u_j)
\]
using the Newtonian potential, with the usual spatial normalization.  Since
\(u\) is locally bounded on the support of \(\eta\), Calder\'on--Zygmund
estimates give \(p_{\rm loc}\in L^\infty_tL^q_x\) for every \(q<\infty\) on
the terminal interval.  The difference
\[
  p-p_{\rm loc}
\]
is harmonic in space on the smaller annulus, modulo a function of time.  The
standard interior estimates for harmonic functions give the displayed
\(C^k\) bounds after fixing the time-dependent additive constant.
\end{proof}

\begin{theorem}[Exterior regularity on no-concentration annuli]
\label{paper8:thm:exterior-regularity-no-concentration}
Let \(u,p\) be a suitable weak solution on
\(\mathbb R^3\times[0,T^*)\), and let \((x_*,T^*)\) be the selected Type II
core point.  Fix \(0<r'<R'<\infty\).  If there is no exterior concentration on
some annulus \(A_{r,R}\) with
\[
  A_{r',R'}\Subset A_{r,R},
\]
then for every integer \(k\ge0\) there are \(\delta>0\), a function
\(c(t)\), and \(C_{r',R',k}<\infty\) such that
\[
  \sup_{t\in(T^*-\delta,T^*)}
  \left(
  \|u(t)\|_{C^k(A_{r',R'})}
  +
  \|p(t)-c(t)\|_{C^k(A_{r',R'})}
  \right)
  \le C_{r',R',k}.
\]
If the no-concentration hypothesis fails, the complementary exterior
concentration alternative holds in the sense of
\cref{paper8:def:exterior-concentration-branch}.
\end{theorem}

\begin{proof}
The no-concentration condition and \cref{paper8:lem:exterior-ckn-cover} give a
finite cover of \(A_{r',R'}\times(T^*-\delta,T^*)\) by terminal cylinders on
which the Caffarelli--Kohn--Nirenberg quantity is below
\(\varepsilon_{\rm CKN}\).  Applying
\cref{paper8:thm:terminal-ckn-regularity} on each cylinder gives local
\(C^k\) bounds for \(u\) and for \(p\) modulo functions of time on smaller
cylinders.  Since the cover is finite, taking the maximum over the cover gives
the stated local bound.
\par
The pressure normalizations on overlapping balls differ only by functions of
time.  Equivalently, one may use
\cref{paper8:lem:exterior-pressure-localization} after the \(k=0\) velocity
bound has been obtained and then apply standard interior parabolic estimates on
nested annuli.  This gives a single time-dependent normalization \(c(t)\) on
\(A_{r',R'}\).  If the no-concentration hypothesis fails, the exterior test
dichotomy gives an exterior concentration branch.
\end{proof}

\begin{lemma}[Divergence-free core and exterior decomposition]
\label{paper8:lem:core-exterior-decomposition}
Fix \(0<r_0<R_0<\infty\), and assume that there is no exterior concentration
on an annulus containing
\[
  \overline{B_{2r_0}(x_*)}\setminus B_{r_0}(x_*).
\]
Let \(\chi\in C_c^\infty(B_{2r_0}(x_*))\) satisfy
\[
  \chi\equiv1\quad\text{on }B_{r_0}(x_*).
\]
Then there are divergence-free fields \(u_{\rm core}\) and \(u_{\rm ext}\)
such that
\[
  u=u_{\rm core}+u_{\rm ext},\qquad
  u_{\rm core}=u\ \text{on }B_{r_0}(x_*),\qquad
  u_{\rm ext}=0\ \text{on }B_{r_0}(x_*),
\]
and all derivatives of \(u_{\rm ext}\) are uniformly bounded on compact
subannuli of the no-exterior-concentration region up to \(T^*\).
\end{lemma}

\begin{proof}
Set
\[
  g=\nabla\chi\cdot u.
\]
The support of \(g\) lies in the annulus
\[
  A_{r_0}=B_{2r_0}(x_*)\setminus B_{r_0}(x_*),
\]
where \(u\) is smooth by
\cref{paper8:thm:exterior-regularity-no-concentration}.  Since
\[
  \int_{A_{r_0}}g\,dx
  =
  \int_{\mathbb R^3}\nabla\chi\cdot u\,dx
  =
  -\int_{\mathbb R^3}\chi\,\nabla\cdot u\,dx
  =0,
\]
the Bogovskii operator on \(A_{r_0}\) gives a vector field
\[
  w=\mathcal B_{A_{r_0}}g
\]
supported in \(A_{r_0}\), smooth up to \(T^*\), and satisfying
\[
  \nabla\cdot w=g.
\]
Define
\[
  u_{\rm core}:=\chi u-w,
  \qquad
  u_{\rm ext}:=u-u_{\rm core}.
\]
Then
\[
  \nabla\cdot u_{\rm core}
  =
  \nabla\chi\cdot u+\chi\nabla\cdot u-\nabla\cdot w=0,
\]
and therefore \(\nabla\cdot u_{\rm ext}=0\).  Since \(\chi=1\) and \(w=0\)
on \(B_{r_0}(x_*)\), \(u_{\rm core}=u\) and \(u_{\rm ext}=0\) there.  The
uniform exterior bounds follow from
\cref{paper8:thm:exterior-regularity-no-concentration} on compact subannuli
and the standard boundedness of the Bogovskii construction on the fixed
annulus.
\end{proof}

\begin{lemma}[Exterior fields vanish in compact core frames]
\label{paper8:lem:exterior-velocity-vanishes}
Let \((x_n,\lambda_n)\) be a core frame with
\[
  x_n\to x_*,
  \qquad
  \lambda_n\to0.
\]
For a compact frame ball \(B_R\), define
\[
  V_{\rm ext}^n(y,\tau)
  :=
  \lambda_nu_{\rm ext}(x_n+\lambda_ny,t_n+\lambda_n^2\tau).
\]
Then, for every compact \(B_R\times I\),
\[
  V_{\rm ext}^n\equiv0
\]
on \(B_R\times I\) for all sufficiently large \(n\).  Consequently
\[
  V_{\rm ext}^n\to0
  \quad\text{in every local }L^q_\tau W^{k,q}_y(B_R)\text{ norm}.
\]
\end{lemma}

\begin{proof}
Since \(x_n\to x_*\) and \(\lambda_n\to0\), for fixed \(R\),
\[
  x_n+\lambda_nB_R\subset B_{r_0}(x_*)
\]
for all sufficiently large \(n\).  On \(B_{r_0}(x_*)\),
\[
  u_{\rm ext}=0.
\]
Therefore \(V_{\rm ext}^n=0\) on \(B_R\times I\) for all large \(n\), which
implies convergence in every stated local norm.
\end{proof}

\begin{lemma}[Exterior pressure vanishes modulo constants]
\label{paper8:lem:exterior-pressure-vanishes}
Let \(p_{\rm ext}\) be any pressure contribution generated by source terms
supported in
\[
  \{|x-x_*|\ge r_0\}.
\]
In a core frame set
\[
  P_{\rm ext}^n(y,\tau)
  :=
  \lambda_n^2p_{\rm ext}(x_n+\lambda_ny,t_n+\lambda_n^2\tau).
\]
Then, for every compact \(B_R\times I\), there are constants \(c_{n,R}(\tau)\)
such that
\[
  P_{\rm ext}^n-c_{n,R}
  \to0
  \quad\text{in }L^\infty_\tau C^1_y(B_R),
\]
and hence
\[
  \nabla_yP_{\rm ext}^n\to0
  \quad\text{in }L^1_\tau H^{-1}_y(B_R).
\]
\end{lemma}

\begin{proof}
The exterior pressure source is at positive physical distance from
\[
  x_n+\lambda_nB_R
\]
for all sufficiently large \(n\).  Therefore \(p_{\rm ext}\) is harmonic,
smooth, and uniformly \(C^2\)-bounded on the physical neighborhood swept out
by the compact frame cylinder.  Take
\[
  c_{n,R}(\tau)
  :=
  \lambda_n^2p_{\rm ext}(x_n,t_n+\lambda_n^2\tau).
\]
By the mean value theorem,
\[
  \|P_{\rm ext}^n-c_{n,R}\|_{C^1_y(B_R)}
  \le
  C_R\lambda_n^3\sup|\nabla_xp_{\rm ext}|
  +
  C_R\lambda_n^4\sup|\nabla_x^2p_{\rm ext}|.
\]
The right-hand side tends to zero because the physical derivatives are
uniformly bounded and \(\lambda_n\to0\).  This proves the asserted
convergence.
\end{proof}

\begin{lemma}[Positive-distance exterior states do not remain in compact core frames]
\label{paper8:lem:positive-distance-exterior}
Let \((x_n,\lambda_n)\) be a core frame with \(x_n\to x_*\) and
\(\lambda_n\to0\).  Let \((z_n,\mu_n)\) be an exterior local state satisfying
\[
  \liminf_{n\to\infty}|z_n-x_*|\ge r_*>0.
\]
Assume that the ordered terminal state-space test has already retained the
compact core at \((x_*,\lambda_n)\).  Then the positive-distance exterior state
does not leave an additional retained velocity component in any compact core
frame.  More precisely, one of the following ordered alternatives occurs:
\begin{enumerate}[label=\textup{(\roman*)}]
\item the state lies on a no-exterior-concentration annulus and is regular and
removable by \cref{paper8:thm:exterior-regularity-no-concentration};
\item the state contains a positive local CKN carrier, in which case it is an
exterior, separated-core, multicore, cascade, or rebased local branch rather
than part of the selected compact core;
\item its velocity contribution is zero in the retained compact-core topology,
and its pressure contribution vanishes modulo functions of time by
\cref{paper8:lem:positive-distance-sources-core-perturbative}, unless the
far-field harmonic tail is routed to the ordered pressure row by
\cref{paper8:lem:far-field-harmonic-tail-routes-to-dyadic-carrier}.
\end{enumerate}
No whole-space \(L^3\) profile decomposition is used in this decoupling.
\end{lemma}

\begin{proof}
For fixed \(R<\infty\), the physical frame ball
\[
  x_n+\lambda_nB_R
\]
lies in \(B_{r_*/4}(x_*)\) for all sufficiently large \(n\), while the profile
state center \(z_n\) lies outside \(B_{r_*/2}(x_*)\).  Thus every local source
supported in the positive-distance state is separated from the compact core
frame.
\par
Run the ordered local CKN test on the exterior state.  If the test is below
threshold on an annulus containing the state, then the no-concentration
regularity theorem gives item \textup{(i)}, and
\cref{paper8:thm:exterior-regular-removal} removes the regular contribution
from compact core frames.  If the test finds a positive CKN carrier, that
carrier is recorded before retention as an exterior-profile, separated-core,
multicore, cascade, or rebased local branch; this is item \textup{(ii)}.
\par
It remains to treat the nonretained part after these tests.  Its velocity
contribution is supported a positive distance from the physical core and
therefore vanishes identically on every compact core cylinder for all large
\(n\), as in \cref{paper8:lem:exterior-velocity-vanishes}.  The pressure
generated by fixed positive-distance sources is perturbative modulo functions
of time by \cref{paper8:lem:positive-distance-sources-core-perturbative}.  If
the only nonvanishing pressure contribution is a far-field harmonic tail, then
\cref{paper8:lem:far-field-harmonic-tail-routes-to-dyadic-carrier} routes that
tail to a dyadic exterior carrier, a regular dyadic tail, or the ordered
pressure row.  Hence no positive-distance exterior state is silently retained
inside the selected compact core.
\end{proof}

\begin{theorem}[Exterior regular components decouple from compact core frames]
\label{paper8:thm:exterior-regular-removal}
On the no-exterior-concentration stratum for the fixed annuli separating an
exterior component from the selected Type II core point, that component is a
regular exterior component and vanishes in compact Type II core frames,
including its pressure contribution modulo constants.
\end{theorem}

\begin{proof}
Apply the divergence-free decomposition of
\cref{paper8:lem:core-exterior-decomposition}.  The exterior component is
identically zero in every compact core frame for large \(n\) by
\cref{paper8:lem:exterior-velocity-vanishes}.  Its pressure contribution is a
locally scaled smooth harmonic field and vanishes modulo constants by
\cref{paper8:lem:exterior-pressure-vanishes}.  The cutoff and Bogovskii
correction terms are supported in the fixed exterior annulus \(A_{r_0}\), so
they also vanish in compact core frames by the same argument.

Thus a regular exterior component decouples from the Type II core without
changing the local repaired-gauge equation in any compact terminal frame, up
to errors converging to zero in the perturbative topology used in
\cref{paper8:lem:local-perturbation} and
\cref{paper8:def:terminal-nonlinear-decoupling}.
\end{proof}

\begin{corollary}[Exterior annulus alternative for the core decomposition]
\label{paper8:cor:exterior-annulus-alternative}
Fix a compact annulus \(A_{r,R}\) around the selected core.  Then exactly one
of the following alternatives holds for the Type II branch on this annulus:
\begin{enumerate}[label=\textup{(\roman*)}]
\item the branch has no exterior concentration on \(A_{r,R}\), and every
compact subannulus of \(A_{r,R}\) is covered by
\cref{paper8:thm:exterior-regularity-no-concentration}; consequently exterior
regular contributions from this annulus vanish in compact core frames
by \cref{paper8:thm:exterior-regular-removal};
\item the branch admits an exterior concentration branch on \(A_{r,R}\), and
therefore belongs to one of the adjacent exterior, separated-core, multicore,
noncompact exterior, or scale-collapse alternatives described in
\cref{paper8:thm:exterior-stratification}.
\end{enumerate}
\end{corollary}

\begin{proof}
This is the exterior concentration stratification of
\cref{paper8:thm:exterior-stratification}, combined in the first alternative
with the no-concentration regularity theorem and the exterior-removal theorem.
\end{proof}

\begin{lemma}[Exterior CKN mass split]
\label{paper8:lem:exterior-ckn-carrier-or-pressure-route}
Let \(x_n\in A_{r,R}\) and \(\rho_n\downarrow0\) form an exterior
concentration branch on \(A_{r,R}\).  After passing to a subsequence, one of
the following alternatives holds:
\begin{enumerate}[label=\textup{(\roman*)}]
\item \emph{velocity carrier:}
\[
  \rho_n^{-2}
  \int_{T^*-\rho_n^2}^{T^*}\int_{B_{\rho_n}(x_n)}
  |u|^3\,dx\,dt
  \ge \frac{\varepsilon_{\rm ext}}{2};
\]
\item \emph{pressure carrier:}
\[
  \rho_n^{-2}
  \int_{T^*-\rho_n^2}^{T^*}\int_{B_{\rho_n}(x_n)}
  |p-(p)_{B_{\rho_n}(x_n)}(t)|^{3/2}\,dx\,dt
  \ge \frac{\varepsilon_{\rm ext}}{2}.
\]
\end{enumerate}
\end{lemma}

\begin{proof}
By definition of exterior concentration,
\[
  \mathcal C_*(u,p;x_n,\rho_n)>\varepsilon_{\rm ext}.
\]
The displayed quantity is the sum of the velocity and pressure terms.  Hence at
least one term is at least \(\varepsilon_{\rm ext}/2\) for infinitely many
\(n\), and a subsequence gives one of the two alternatives.  This lemma is only
the local velocity/pressure split and does not itself close the pressure-only
case.
\end{proof}

\begin{lemma}[Positive-distance sources are perturbative in compact core frames]
\label{paper8:lem:positive-distance-sources-core-perturbative}
Let \((x_n,\lambda_n,t_n)\) be a compact core frame with
\[
  x_n\to x_*,
  \qquad
  \lambda_n\downarrow0,
  \qquad
  t_n\to T^*.
\]
Fix \(d>0\), and let \(K\Subset\{x:|x-x_*|\ge d\}\).  Suppose
\(F_n\) is supported in \(K\) for a.e. terminal time and satisfies the local
bound
\[
  \sup_n\|F_n\|_{L^1(T^*-\delta,T^*;L^1(K))}<\infty .
\]
Let \(q_n\) be the Newtonian pressure generated by \(F_n\),
\[
  q_n(\cdot,t)=\partial_i\partial_j(-\Delta)^{-1}F_{n,ij}(\cdot,t),
\]
restricted to a neighborhood of \(x_*\).  For every compact
\(B_M\times I\) in the core variables, set
\[
  Q_n(y,\tau)
  =
  \lambda_n^2 q_n(x_n+\lambda_n y,t_n+\lambda_n^2\tau),
  \qquad
  c_n(\tau)=\lambda_n^2 q_n(x_n,t_n+\lambda_n^2\tau).
\]
Then
\[
  Q_n-c_n\to0
  \quad\text{in }L^1_\tau C^1_y(B_M).
\]
If \(w_n\) is a velocity contribution supported in \(K\), then
\[
  \lambda_n w_n(x_n+\lambda_n y,t_n+\lambda_n^2\tau)\equiv0
  \quad\text{on }B_M\times I
\]
for all sufficiently large \(n\).  Thus positive-distance velocity sources and
their pressure tails are perturbative in every compact core frame, modulo
functions of time.
\end{lemma}

\begin{proof}
For fixed \(M\), the physical balls \(x_n+\lambda_nB_M\) lie in
\(B_{d/4}(x_*)\) for all sufficiently large \(n\), whereas \(K\subset
\{|x-x_*|\ge d\}\).  The velocity assertion follows immediately from the
support separation.
\par
On \(B_{d/4}(x_*)\) the kernel
\(\partial_i\partial_j(-\Delta)^{-1}\) applied to sources supported in \(K\) is
smooth.  For \(|\alpha|\le2\),
\[
  \|\nabla_x^\alpha q_n(t)\|_{L^\infty(B_{d/4}(x_*))}
  \le C_{\alpha,d,K}\|F_n(t)\|_{L^1(K)}
\]
for a.e. \(t\).  By the mean value theorem,
\[
  \|Q_n(\tau)-c_n(\tau)\|_{C^1_y(B_M)}
  \le
  C_M\lambda_n^3\|\nabla_x q_n(t_n+\lambda_n^2\tau)\|_{L^\infty(B_{d/4})}
  +C_M\lambda_n^4
  \|\nabla_x^2q_n(t_n+\lambda_n^2\tau)\|_{L^\infty(B_{d/4})}.
\]
Integrating in \(\tau\in I\) and using \(dt=\lambda_n^2\,d\tau\), the
right-hand side is bounded by
\[
  C_{M,I,d,K}(\lambda_n+\lambda_n^2)
  \|F_n\|_{L^1(T^*-\delta,T^*;L^1(K))},
\]
which tends to zero.  This gives the claimed convergence.  The proof is local
on the fixed set \(K\) and uses no whole-space critical norm.
\end{proof}

\begin{lemma}[Exterior local states decouple from the selected compact core]
\label{paper8:lem:exterior-profile-decouples-from-core}
Let an exterior concentration branch on \(A_{r,R}\) be in the exterior-profile
subcase of \cref{paper8:thm:exterior-stratification}.  Then the exterior
profile is not a terminal obstruction for the selected compact core.  More
precisely, in every compact core frame
\((x_n,\lambda_n,t_n)\) with \(x_n\to x_*\) and \(\lambda_n\downarrow0\), the
exterior local state has zero compact-core velocity contribution in the retained
terminal state-space topology.  Its pressure contribution either vanishes
modulo functions of time on compact core cylinders, or it is routed by the
far-field harmonic-tail theorem
\cref{paper8:lem:far-field-harmonic-tail-routes-to-dyadic-carrier} to a dyadic
exterior carrier or to the ordered pressure row.  If the exterior state itself
is selected as a possible singular branch, it is rebased at its own center and
tested by the same local Type II state-space decomposition; it is no longer an
exterior exit for the original selected core.
\end{lemma}

\begin{proof}
The exterior-profile subcase is read locally: its CKN witnesses remain at
positive physical distance from \(x_*\), or escape every compact set in the
selected core variables.  Suppose that a positive velocity contribution from
this exterior state were visible on some compact core cylinder.  Then
\cref{paper8:thm:terminal-windowed-profile-completeness} would assign that
positive compact-core mass to a first local state-space label.  If the label is
retained comparable mass, it must already have been grouped into the selected
core; if it is not retained comparable mass, it is an adjacent compact,
separated, cascade, rough-core, or pressure/exterior row.  Hence no positive
exterior velocity state remains in the nonretained exterior part of a retained
compact core branch.
\par
For pressure, decompose the exterior pressure seen on a compact core cylinder
into the local Calder\'on--Zygmund part generated by sources in a fixed annulus,
the far-field harmonic tail, and a function of time.  Fixed positive-distance
annular sources are perturbative by
\cref{paper8:lem:positive-distance-sources-core-perturbative}.  If the
remaining far-field harmonic tail does not vanish modulo functions of time,
then
\cref{paper8:lem:far-field-harmonic-tail-routes-to-dyadic-carrier} produces a
dyadic exterior stress carrier or the ordered pressure row.  Thus pressure is
also either perturbative for the selected core or routed before retention.
\par
Thus the exterior profile is perturbative for the selected core.  If one
chooses the exterior profile as the new active branch, the center and scale are
changed and the analysis restarts from the covered-entry theorem
\cref{paper6:cor:covered-class-entry-discharged}.  This rebasing does not
leave an independent exterior terminal row attached to the original core.
\end{proof}

\begin{lemma}[Separated exterior cores are multicore or rebased local branches]
\label{paper8:lem:separated-exterior-core-is-multicore}
Suppose an exterior concentration branch on \(A_{r,R}\) carries a velocity
carrier and that, after passing to a subsequence, the exterior carrier remains
as a positive active core at terminal scales comparable to another selected
active scale.  Then either the two active centers are grouped as a separated
multicore branch in the local state space, or the exterior center is rebased as
its own physical local Type II germ.  In the first case the branch is excluded
by \cref{paper6:thm:multibubble}; in the second case the row is no longer an
exterior row for the original selected core.
\end{lemma}

\begin{proof}
The exterior centers satisfy \(|x_n-x_*|\ge r>0\), while the selected compact
core scales tend to zero.  Hence the relative scale-center distance between
the original core and the exterior carrier tends to infinity whenever the
scales are comparable, and the separated-point reduction
\cref{paper3:thm:separated-point-reduction} applies.  Its output is exactly a
separated-core or multicore label unless the exterior center is selected as the
new base point.  The separated-core and multicore labels are part of the
multibubble/cascade branch excluded by \cref{paper6:thm:multibubble}.  If the
exterior center is chosen as the base point, the local CKN lower bound supplies
a physical local concentration sequence there and
\cref{paper6:cor:covered-class-entry-discharged} starts the local Type II
analysis for that new germ.  This is a rebasing, not an unresolved exterior
mechanism for the original core.
\end{proof}

\begin{lemma}[Exterior finite-scale packing is an ordered local row]
\label{paper8:lem:exterior-finite-scale-packing}
Fix a compact annulus \(A_{r,R}\).  In any fixed dyadic scale band of exterior
carrier cylinders contained in \(A_{r,R}\times(T^*-\delta,T^*)\), a Vitali
selection produces a pairwise disjoint subfamily with uniformly bounded overlap
after enlargement.  Therefore an exterior velocity carrier in that scale band
has an ordered local interpretation:
\begin{enumerate}[label=\textup{(\roman*)}]
\item a finite selected subfamily gives a separated-core or multicore label;
\item an unbounded number of comparable exterior carriers, or a sequence of
scale bands with no finite selected representative, is by definition a
noncompact exterior or cascade boundary face before the retained compact core is
declared.
\end{enumerate}
No uniform bound on the number of exterior carriers is asserted or used.
\end{lemma}

\begin{proof}
The geometric Vitali covering lemma applies on the fixed annulus for each
dyadic radius band.  It gives a disjoint subfamily whose fixed enlargement
covers all carrier balls in that band, with overlap bounded by an absolute
constant depending only on the enlargement factor.  The ordered local
state-space construction selects active carriers from this finite-overlap
family.  If only finitely many selected carriers remain, their positive CKN
lower bounds and separated scale-center relations give the separated-core or
multicore row.  If no finite selected list captures the active exterior
carriers, then the sequence has left every finite scale-center chart; this is
the noncompact exterior boundary face, or the same-point scale-cascade face
when the centers converge but the radii separate strictly.  The argument is
geometric and local; it does not invoke a global \(L^3\) mass bound to bound
the cardinality of the family.
\end{proof}

\begin{lemma}[Exterior scale-center compactification]
\label{paper8:lem:exterior-scale-center-compactification}
Let \((z_n,\sigma_n)\) be exterior active carrier data in an exterior frame,
with \(\sigma_n>0\), and suppose the corresponding terminal cylinders carry a
fixed positive CKN lower bound.  After passing to a subsequence, exactly one of
the following ordered alternatives occurs:
\begin{enumerate}[label=\textup{(\roman*)}]
\item the normalized centers \(z_n\) remain in a compact subset of the exterior
frame; this is a compact exterior-profile, separated-core, or multicore row;
\item \(|z_n|\to\infty\) and \(\sigma_n/|z_n|\to0\); this is a dyadic point
carrier in annuli \(\{|y|\simeq |z_n|\}\);
\item \(|z_n|\to\infty\) and
\[
  0<\liminf_{n\to\infty}\sigma_n/|z_n|
  \le
  \limsup_{n\to\infty}\sigma_n/|z_n|<\infty;
\]
this is a dyadic annular carrier;
\item the relative scale leaves every compact subset at a fixed exterior
point, or after rebasing at the exterior point; this is the same-point cascade
or exterior scale-collapse row.
\end{enumerate}
These alternatives are disjoint after the ordered first-failure convention and
exhaust the exterior scale-center configurations.  In particular, the phrase
``standard boundary faces'' in the exterior routing below means precisely the
four cases listed here.
\end{lemma}

\begin{proof}
If \(z_n\) is bounded, compactness of closed balls in the exterior frame gives a
subsequential limit, yielding item \textup{(i)}.  If \(|z_n|\to\infty\), pass to
a subsequence so that \(\sigma_n/|z_n|\) converges in the compactified interval
\([0,\infty]\).  Limit \(0\) gives item \textup{(ii)}.  A limit in
\((0,\infty)\) gives item \textup{(iii)}.  Limit \(+\infty\), or failure of
compactness after recentering at a bounded exterior point, means that the
active radius contains a further unresolved scale.  The ordered local
state-space compactification then records either strict same-point scale
separation or a genuine scale-collapse row, giving item \textup{(iv)}.  These
possibilities exhaust the one-point compactification of the center variable
together with the compactified relative-scale variable.
\end{proof}

\begin{lemma}[Dyadic exterior occupation selection]
\label{paper8:lem:dyadic-exterior-occupation-selection}
Work in an exterior frame and let
\[
  \mathcal A_k:=B_{2^{k+1}}\setminus B_{2^k},
  \qquad k\ge0 .
\]
For \(K\ge0\), define the far-dyadic CKN level
\[
  \Gamma_{n,K}
  :=
  \sup_{\substack{k\ge K,\ y\in\mathcal A_k\\
  0<r\le 2^{k-4}\\
  Q_r^*(y)\ {\rm admissible}}}
  \mathcal C_*(U_n,P_n;y,r),
\]
where the terminal cylinders are understood after pulling back to the physical
variables.  Then the ordered test gives one of the following alternatives after
passing to a subsequence:
\begin{enumerate}[label=\textup{(\roman*)}]
\item a dyadic CKN carrier exists:
\[
  \mathcal C_*(U_n,P_n;y_{n},r_n)\ge\varepsilon_{\rm ext}
\]
for a sequence \(|y_n|\to\infty\);
\item the far-dyadic CKN level is eventually below threshold:
\[
  \limsup_{K\to\infty}\limsup_{n\to\infty}\Gamma_{n,K}
  \le \varepsilon_{\rm ext},
\]
and, because \(\varepsilon_{\rm ext}<\varepsilon_{\rm CKN}\), every
sufficiently far dyadic annulus is regular on its interior by the terminal CKN
criterion.
\end{enumerate}
Thus a noncompact exterior branch can remain a terminal obstruction only in the
dyadic CKN-carrier case.
\end{lemma}

\begin{proof}
If
\[
  \limsup_{K\to\infty}\limsup_{n\to\infty}\Gamma_{n,K}>
  \varepsilon_{\rm ext},
\]
then for arbitrarily large \(K\) there are \(n\), \(k\ge K\), \(y\in\mathcal
A_k\), and an admissible \(r\le2^{k-4}\) with
\(\mathcal C_*(U_n,P_n;y,r)>\varepsilon_{\rm ext}\).  A diagonal subsequence
gives item \textup{(i)} with \(|y_n|\to\infty\).
\par
Otherwise the displayed limsup is at most \(\varepsilon_{\rm ext}\).  Fix
\(\eta\) with \(\varepsilon_{\rm ext}<\eta<\varepsilon_{\rm CKN}\).  After
discarding a finite initial set of annuli and passing to the selected tail in
\(n\), every admissible subcylinder in every far dyadic annulus has CKN
quantity at most \(\eta\), hence below \(\varepsilon_{\rm CKN}\).  A
bounded-overlap cover of each smaller annulus
\[
  B_{2^{k+1}-2^{k-5}}\setminus B_{2^k+2^{k-5}}
\]
by such cylinders and
\cref{paper8:thm:terminal-ckn-regularity} give local smoothness on the dyadic
annulus.  Since the cover is local on each annulus and the constants are
scale-invariant after rescaling that annulus to unit size, no global tail norm
is used.  Therefore the diffuse no-carrier case is regular on every far dyadic
annulus and cannot be a terminal noncompact exterior obstruction.
\end{proof}

\begin{lemma}[Exterior same-point cascades route to the cascade package]
\label{paper8:lem:exterior-same-point-cascade-routing}
Let exterior carriers on \(A_{r,R}\) have centers converging to the same
physical point \(z\ne x_*\) and have strictly separated scales.  Then, after
rebasing at \(z\), the branch is a same-point cascade in the sense of
\cref{paper3:thm:same-point-reduction}.  It is therefore closed by the
multibubble/cascade package \cref{paper6:thm:multibubble}, unless the ordered
rebased analysis exits earlier through one of the non-exterior first-failure
rows.
\end{lemma}

\begin{proof}
Because \(z\ne x_*\), the carriers are at positive physical distance from the
original core and are perturbative for that core by
\cref{paper8:lem:positive-distance-sources-core-perturbative}.  Recenter at
\(z\) and use the exterior carrier scales as the active scales.  Strict
separation of the radii is precisely the same-point scale-cascade boundary face
in the relative scale-center compactification.  The same-point cascade
reduction \cref{paper3:thm:same-point-reduction} then assigns the rebased
branch to the cascade branch, which is excluded by
\cref{paper6:thm:multibubble}.  If the rebased ordered tests fail before the
cascade label is reached, the failure is one of the ordinary representation,
modulation, active-core, pressure, rough-core, carrier, or cost rows; none is
an exterior row for the original selected core.
\end{proof}

\begin{lemma}[Noncompact exterior branches admit dyadic active carriers]
\label{paper8:lem:noncompact-exterior-dyadic-selection}
Let the exterior stratification of \cref{paper8:thm:exterior-stratification}
produce the noncompact exterior alternative in an exterior frame.  Then, after
passing to a subsequence, exactly one of the following occurs:
\begin{enumerate}[label=\textup{(\roman*)}]
\item the noncompact exterior occupation is diffuse below the CKN threshold on
all sufficiently far dyadic annuli, hence regular there by
\cref{paper8:lem:dyadic-exterior-occupation-selection};
\item there are dyadic annuli and terminal cylinders in the exterior frame with
CKN quantity bounded below by \(\varepsilon_{\rm ext}\), hence an exterior
velocity or pressure carrier by
\cref{paper8:lem:exterior-ckn-carrier-or-pressure-route};
\item the noncompact behavior is only a far-field harmonic pressure tail, which
is routed to item \textup{(i)}, item \textup{(ii)}, or the ordered pressure row
by \cref{paper8:lem:far-field-harmonic-tail-routes-to-dyadic-carrier}.
\end{enumerate}
\end{lemma}

\begin{proof}
The noncompact alternative is represented by active data leaving every compact
subset of the exterior scale-center chart.  Apply
\cref{paper8:lem:exterior-scale-center-compactification}.  The compact-center,
same-point cascade, and scale-collapse faces are not the noncompact exterior
face and are routed by their corresponding lemmas below.  The remaining
noncompact face is a dyadic occupation state in annuli \(\mathcal A_k\),
\(k\to\infty\).
\par
Apply the quantitative dyadic selection lemma
\cref{paper8:lem:dyadic-exterior-occupation-selection}.  If the occupation is
below the CKN threshold on all sufficiently far dyadic annuli, item
\textup{(i)} holds.  Otherwise a dyadic CKN carrier exists.  The velocity/pressure
split
\cref{paper8:lem:exterior-ckn-carrier-or-pressure-route} gives item
\textup{(ii)} unless the carrier is purely pressure.  A purely harmonic
far-field pressure tail is then routed by
\cref{paper8:lem:far-field-harmonic-tail-routes-to-dyadic-carrier}, giving item
\textup{(iii)}.  All tests are dyadic and local.
\end{proof}

\begin{lemma}[Noncompact exterior carriers route to structural rows]
\label{paper8:lem:noncompact-exterior-structural-route}
\label{paper8:lem:noncompact-exterior-routes-to-structural-row}
Every dyadic active carrier produced by
\cref{paper8:lem:noncompact-exterior-dyadic-selection} is routed, after passing
to a subsequence, to one of the following rows:
\begin{enumerate}[label=\textup{(\roman*)}]
\item a regular exterior-removal row, if all smaller dyadic tests are
no-concentration rows;
\item a separated-core or multicore row, closed by
\cref{paper6:thm:multibubble};
\item a same-point cascade row, closed by
\cref{paper8:lem:exterior-same-point-cascade-routing};
\item an exterior scale-collapse row, rebased and tested by
\cref{paper8:lem:exterior-scale-collapse-rebased};
\item a pressure/exterior row which is converted to a velocity carrier or to a
non-exterior first-failure row by
\cref{paper8:lem:exterior-pressure-only-not-terminal}.
\end{enumerate}
Thus noncompact exterior concentration is not an independent terminal
alternative.
\end{lemma}

\begin{proof}
Use the dyadic carrier selected in the preceding lemma.  If adjacent dyadic
tests become no-concentration tests after shrinking, then
\cref{paper8:thm:exterior-regular-removal} gives regular exterior removal.
Otherwise
\cref{paper8:lem:exterior-scale-center-compactification} gives the ordered
boundary face of the carrier.  Positive centers separated from the original core
at comparable scales give the separated-core or multicore row by
\cref{paper8:lem:separated-exterior-core-is-multicore}.  Centers converging to
one exterior point with strict scale separation give the same-point cascade row
by \cref{paper8:lem:exterior-same-point-cascade-routing}.  If the dyadic
carrier develops its own unresolved scale collapse, it is the exterior
scale-collapse row and is handled by
\cref{paper8:lem:exterior-scale-collapse-rebased}.  Finally, if the diagnostic
is pressure-only, the pressure-only routing lemma applies.  These cases are the
ordered local boundary faces of the exterior scale-center chart and exhaust the
noncompact exterior alternative.
\end{proof}

\begin{lemma}[Exterior scale collapse is an ordinary rebased scale-collapse row]
\label{paper8:lem:exterior-scale-collapse-rebased}
Let an exterior concentration branch develop scale collapse in its own exterior
variables.  Recenter at the exterior active carrier and use its scale as the
local active scale.  Then the rebased branch is a represented local
scale-collapse branch in the sense of
\cref{paper6:thm:scale-stratification}, unless the ordered rebased tests fail
first through one of the non-exterior rows: representation/pressure failure,
autonomous-modulation failure or modulation-driven recentering, loss of
retained active core or pressure-only critical mass, rough-core loss,
carrier/cost validation failure, or multicore/cascade.  In the represented
case, the finite-cost, retained compact scale-collapse, and scale-rigid
subcases are closed by
\[
  \cref{paper6:thm:scale-cost,paper5:thm:compact-scale-collapse-rigidity-discharged,paper3:prop:scale-rigidity-discharged}.
\]
\end{lemma}

\begin{proof}
The exterior carrier supplies a positive local CKN mass floor on a compact
cylinder in the rebased variables.  Run the covered-entry and repaired-gauge
selection tests at that carrier.  If a row fails before a represented
scale-collapse branch is obtained, the first-failure ledger records exactly the
listed non-exterior local row; this is the same ordered local test used in
\cref{paper6:cor:covered-class-entry-discharged} and
\cref{paper5:lem:retained-scale-collapse-first-failure}.
\par
If the represented scale-collapse branch is obtained, apply
\cref{paper6:thm:scale-stratification}.  Finite scale cost and finite absolute
scale cost are excluded by \cref{paper6:thm:scale-cost}.  Thin negative drift
is already closed as an independent terminal row by
\cref{paper6a:thm:thin-negative-drift-closed}.  Selected compactness-package
failure is routed by
\cref{paper5:thm:selected-compactness-failure-routed}.  The retained compact
scale-collapse state is excluded by
\cref{paper5:thm:compact-scale-collapse-rigidity-discharged}, and the
scale-rigid residual state is excluded by
\cref{paper3:prop:scale-rigidity-discharged}.  Autonomous modulation and the
other listed first failures remain ordinary non-exterior rows in the final
ledger; the exterior label itself has disappeared after rebasing.
\end{proof}

\begin{lemma}[Fixed annular stress is regular or has a local CKN carrier]
\label{paper8:lem:fixed-annular-stress-ckn-dichotomy}
Let \(\mathcal A\Subset\mathbb R^3\setminus B_2\) be a fixed annulus in a
local frame and let \(I\Subset(-\infty,0]\) be a compact terminal time
interval.  Suppose a subsequence has a positive annular stress level
\[
  \limsup_{n\to\infty}
  \|U_n\otimes U_n\|_{L^q_\tau(I;L^{3/2}_y(\mathcal A))}>0
\]
for some \(1<q<\infty\).  Then the ordered local CKN test on
\(\mathcal A\times I\) gives one of the following alternatives:
\begin{enumerate}[label=\textup{(\roman*)}]
\item a local CKN carrier occurs on a subcylinder contained in a
slightly larger annulus;
\item all sufficiently small admissible subcylinders in the annulus are below
\(\varepsilon_{\rm CKN}\), hence the annular stress is a regular
compact-window annular source.  It is absorbed into the retained smooth
annular part of the local pressure atlas, unless the harmonic-tail test sends
its pressure effect to the ordered pressure row.
\end{enumerate}
Thus a positive fixed-annulus stress level is not by itself promoted to a
singular carrier; it is tested locally first.
\end{lemma}

\begin{proof}
Apply the terminal CKN test to all admissible subcylinders whose doubled
cylinders remain in a fixed slightly larger annulus containing \(\mathcal A\).
If one of these cylinders has CKN quantity at least
\(\varepsilon_{\rm ext}\), then item \textup{(i)} holds.  If no such cylinder exists, every admissible subcylinder is
below \(\varepsilon_{\rm ext}<\varepsilon_{\rm CKN}\).  The terminal CKN
regularity theorem and a finite-overlap cover of compact subannuli give local
smoothness on \(\mathcal A\times I\), so the stress lower bound is a regular
annular source rather than a singular exterior carrier.  Since \(\mathcal A\)
is at bounded normalized distance from the core in the same local frame, this
regular source is part of the compact-window pressure atlas; if it is not
absorbed there, the first failed atlas condition is exactly the ordered
pressure row.  No cardinality bound or global \(L^3\) tail estimate is used.
\end{proof}

\begin{lemma}[Far-field harmonic tails route to dyadic carriers]
\label{paper8:lem:far-field-harmonic-tail-routes-to-dyadic-carrier}
Let \(B_1\Subset B_2\) be nested balls in a local frame and let
\[
  H_n(\tau)=H_n^{\rm far}(\tau)
\]
be the mean-zero harmonic pressure on \(B_2\) generated, in the local pressure
atlas, by the quadratic stress \(F_n=U_n\otimes U_n\) outside \(B_4\).  Let
\(\{\theta_k\}_{k\ge2}\) be a smooth dyadic partition of unity on
\(\mathbb R^3\setminus B_4\), with
\[
  \operatorname{supp}\theta_k
  \subset B_{2^{k+1}}\setminus B_{2^{k-1}} .
\]
Assume the pressure reconstruction and gauge rows have not failed.  If, for
some \(1<q<\infty\),
\[
  \limsup_{n\to\infty}
  \|H_n\|_{L^q_\tau L^{3/2}_y(B_1)}>0,
\]
then, after passing to a subsequence, either
\begin{enumerate}[label=\textup{(\roman*)}]
\item a dyadic exterior velocity-stress occupation state is present:
\[
  \limsup_{n\to\infty}
  \sum_{k\ge K_n}2^{-2k}
  \|U_n\otimes U_n\|_{L^q_\tau L^{3/2}_y(\mathcal A_k)}
  >0
\]
for some \(K_n\to\infty\); or
\item a fixed dyadic annulus carries a positive annular stress lower bound,
which is then tested by
\cref{paper8:lem:fixed-annular-stress-ckn-dichotomy} and is either a regular
annular source or a local CKN carrier.
\end{enumerate}
In item \textup{(i)}, the dyadic occupation state is then tested by
\cref{paper8:lem:dyadic-exterior-occupation-selection}: it either produces a
dyadic CKN carrier or is regular on the far dyadic annuli.  If the dyadic tail
is regular but its harmonic pressure contribution remains nonzero on the
interior cylinder, the branch is assigned to the ordered pressure row rather
than to exterior concentration.  Consequently a nonzero far-field harmonic
pressure tail is never silently discarded; it is a dyadic exterior carrier, a
regular dyadic tail with vanishing harmonic effect, or the ordered pressure row.
\end{lemma}

\begin{proof}
Write
\[
  H_n=\sum_{k\ge2}H_{n,k}+h_n,
  \qquad
  H_{n,k}
  =
  \mathcal R_i\mathcal R_j(\theta_kF_{n,ij})
  -\fint_{B_2}\mathcal R_i\mathcal R_j(\theta_kF_{n,ij})\,dy,
\]
where \(h_n\) is the pressure-gauge/reconstruction remainder.  By assumption
the gauge and reconstruction rows have not failed, so \(h_n\) is absent modulo
functions of time in the retained pressure atlas.
\par
For \(y\in B_2\) and \(z\in\operatorname{supp}\theta_k\), the second derivative
of the Newtonian kernel is \(O(2^{-3k})\).  Hence
\[
  \|H_{n,k}(\tau)\|_{L^{3/2}(B_1)}
  \le
  C2^{-3k}\|F_n(\tau)\|_{L^1(\mathcal A_k)}
  \le
  C2^{-2k}\|F_n(\tau)\|_{L^{3/2}(\mathcal A_k)},
\]
because \(|\mathcal A_k|^{1/3}\simeq2^k\).  Summing in \(k\) and taking the
\(L^q_\tau\) norm gives
\[
  \|H_n\|_{L^q_\tau L^{3/2}(B_1)}
  \le
  C
  \sum_{k\ge2}2^{-2k}
  \|F_n\|_{L^q_\tau L^{3/2}(\mathcal A_k)} .
\]
If the left-hand side has positive limsup, the weighted dyadic stress sum has
positive limsup.  Either a fixed finite set of dyadic annuli carries a positive
portion of that limsup, or the positive portion escapes to \(k\to\infty\).  The
first case gives item \textup{(ii)} and is routed by
\cref{paper8:lem:fixed-annular-stress-ckn-dichotomy}; a fixed annular stress
lower bound is not identified with a singular carrier until the local CKN test
has selected one.  The escaping case is precisely the dyadic occupation state
in item \textup{(i)}.  Applying
\cref{paper8:lem:dyadic-exterior-occupation-selection} to that occupation
state gives a dyadic CKN carrier or a regular dyadic tail.  In the regular-tail
case, either the corresponding harmonic contribution vanishes modulo functions
  of time on \(B_1\), or the nonzero harmonic contribution is by definition the
ordered pressure row.  No whole-space pressure or velocity norm is used.
\end{proof}

\begin{lemma}[Compact local-state factor with failure faces]
\label{paper6a:lem:compact-local-state-factor}
Fix a model cylinder \(Q\) and choose a countable exhaustion
\(Q^{(\ell)}\Subset Q\) by compact subcylinders, where \(B^{(\ell)}\) denotes
the spatial ball of \(Q^{(\ell)}\).  Choose also countable dense families of
smooth test fields \(\psi_j,\Psi_j,\varphi_j\) on those subcylinders.
Consider represented suitable branches on \(Q\).  For each compact
subcylinder record the coordinates
\[
  \int_{Q^{(\ell)}} V\cdot \psi_j,
  \qquad
  \int_{Q^{(\ell)}} \nabla V: \Psi_j,
  \qquad
  \int_{Q^{(\ell)}} \bigl(P-(P)_{B^{(\ell)}}(t)\bigr)\varphi_j,
\]
together with the local-energy, pressure-atlas, gauge, compact-upper,
rough-core, modulation, and cost diagnostic coordinates, each compactified by
adding the value \(\infty\) or a named failure face.

Then the closure of all represented local states satisfying the retained
finite bounds is compact metrizable.  Any limit point either:
\begin{enumerate}[label=\textup{(\roman*)}]
\item is again a represented local suitable state with the inherited pressure
    atlas and local energy inequality on smaller compact cylinders; or
\item lies on one of the explicitly named failure faces.
\end{enumerate}
Moreover, the occupation coordinates used in the pressure, shell, and
reselection compactifications are continuous coordinate projections in this
compactification.
\end{lemma}

\begin{proof}
On each retained chart the compact-window estimates furnish uniform weak-*\
\(L_t^\infty L_x^2\), weak \(L_t^2H_x^1\), and weak \(L^{3/2}\) pressure
bounds on the model cylinder.  Banach--Alaoglu therefore gives compactness of
the recorded linear coordinates, and Aubin--Lions gives strong local velocity
compactness on every subcylinder for which the time-derivative/modulation row
remains finite.

More explicitly, the compactified state space is defined as the closed subset
of the product
\[
  \prod_j K_j
\]
where each \(K_j\) is one of the following compact factors: a weakly compact
closed ball for a retained finite linear coordinate, a one-point compactified
interval \([0,\infty]\) for a nonnegative diagnostic size, or a finite
discrete set of named failure faces.  The occupation coordinates are included
as coordinate projections in this product topology.  Thus their positive
superlevel sets are closed, and weak-* convergence of normalized occupation
measures is exactly weak-* convergence against these continuous coordinate
projections.  The distributional equation and local energy inequality are
imposed as closed conditions on the retained finite-coordinate part; if any
condition fails, the state lies on the corresponding failure face rather than
inside the retained suitable-state component.

Passing to a diagonal subsequence over the countable test
families yields a compact metrizable closure.  Lower semicontinuity of the
local energy inequality and the stability of the local pressure atlas preserve
the represented local suitable structure on smaller compact cylinders.  If one
of the compactness inputs is not finite, then by definition the corresponding
coordinate lands on its named failure face.  This is exactly the stated
alternative.  The pressure, shell, and reselection occupation variables are
among the product coordinates just recorded, so their projections are
continuous in the compactified topology.
\end{proof}

\begin{lemma}[Normalized defect mass: atomic or diffuse visible row]
\label{paper6a:lem:normalized-defect-atomic-or-diffuse}
Let \(\mu_n\) be normalized nonnegative defect measures on a compactified local
state space \(X\), and suppose
\[
  \mu_n \rightharpoonup \mu,
  \qquad
  \mu(X)=1.
\]
Assume the compactification includes a countable separating family of
nonnegative local occupation coordinates
\[
  \mathcal O_j:X\to[0,\infty],
\]
generated by compactly supported nonnegative test functions \(\phi_j\) which
separate nonzero local \(L^1_+\)-mass on compact subcylinders.  These
coordinates are assumed to be continuous coordinate projections in the product
compactification, and for every \(j\) and \(\eta>0\) the superlevel set
\[
  \{x\in X:\mathcal O_j(x)\ge\eta\}
\]
is closed.  Then exactly one of the following occurs:
\begin{enumerate}[label=\textup{(\roman*)}]
\item there exist \(j\) and \(\eta>0\) such that
\[
  \mu\bigl(\{x\in X:\mathcal O_j(x)\ge\eta\}\bigr)>0.
\]
In this case a support point with positive compact local activity is realized
by descendant states;
\item for every \(j\) and every \(\eta>0\),
\[
  \mu\bigl(\{x\in X:\mathcal O_j(x)\ge\eta\}\bigr)=0.
\]
Then the defect is genuinely diffuse.  The nonzero probability mass \(\mu\)
itself is a visible diffuse occupation-state coordinate, and it must be routed
as a named diffuse row rather than replaced by a nonzero support profile.
\end{enumerate}
\end{lemma}

\begin{proof}
If alternative \textup{(i)} holds, the set
\(\{x\in X:\mathcal O_j(x)\ge\eta\}\) is closed and has positive \(\mu\)-mass,
so it meets \(\operatorname{supp}\mu\).  Any support point in that set carries
positive compact local activity, and the relevant support-realization theorem
turns it into an actual descendant realization.

If alternative \textup{(i)} fails, then alternative \textup{(ii)} holds by
definition.  Because the coordinates are countable and separate compact local
mass, no support point can carry positive compact local activity at a positive
level.  However \(\mu(X)=1\), so the limiting defect is still nonzero.  Since
the compactification records the normalized defect mass itself and not only its
support points, this nonzero diffuse measure remains a visible state-space
coordinate.  The two alternatives are manifestly exclusive and exhaustive.
\end{proof}

\begin{lemma}[Residual selected-time support transfer]
\label{paper6a:lem:residual-selected-time-support-transfer}
Let
\[
  V_n=U_n+S_n
\]
be an admissible retained terminal decomposition on a compact terminal cylinder
\(Q_{2R,I}=B_{2R}\times I\).  Fix selected times \(\tau_n\in I\) and a
compact ball \(B_\rho\Subset B_R\).  Let
\[
  \mu_n^{\rm res}:=|S_n(\tau_n)|^3\,dy\big|_{B_\rho}
\]
and assume \(\mu_n^{\rm res}(B_\rho)>0\) for all \(n\).  Recenter the
selected times by \cref{paper6a:lem:interior-time-selection} so that they are
represented on a fixed model interval.  For each point carrying residual mass,
record its compactified relative center-scale data, the rebased local
occupation state on a fixed model cylinder, and the named failure faces of the
compact local-state factor.  Then, after normalizing by
\(\mu_n^{\rm res}(B_\rho)\) and passing to a subsequence, these pushforwards
converge weakly-* to a probability measure \(\Pi_{\rm res}\) on a compact
metrizable residual state space
\[
  \mathfrak X_{\rm res}
  :=
  \overline{\mathcal C}_{\rm res}
  \times \mathfrak O_{\rm res}
  \times \mathfrak F_{\rm res} .
\]
The local occupation coordinates of \(\mathfrak X_{\rm res}\) are continuous
coordinate projections, so the dichotomy of
\cref{paper6a:lem:normalized-defect-atomic-or-diffuse} applies to
\(\Pi_{\rm res}\).  Every point of \(\operatorname{supp}\Pi_{\rm res}\) is
realized by an actual sequence of selected-time residual descendants and their
corresponding rebased local states from the original represented branch.
\end{lemma}

\begin{proof}
Use \cref{paper6a:lem:interior-time-selection} to place the selected time at
the origin of a fixed model interval in the local validation family.  For each
\(n\), send a point of \(B_\rho\) carrying residual mass to its compactified
relative center-scale data together with the associated rebased local state on
that fixed model cylinder and the named failure faces of
\cref{paper6a:lem:compact-local-state-factor}.  The compactified relative
scale-center chart is compact, and the local-state factor with failure faces is
compact metrizable by \cref{paper6a:lem:compact-local-state-factor}.  Hence
their product \(\mathfrak X_{\rm res}\) is compact metrizable.

Set
\[
  \widehat\mu_n^{\rm res}
  :=
  \frac{\mu_n^{\rm res}}{\mu_n^{\rm res}(B_\rho)}
\]
and push \(\widehat\mu_n^{\rm res}\) forward by the residual state map to
\(\mathfrak X_{\rm res}\).  By compactness, after passing to a subsequence,
these pushforwards converge weakly-* to a probability measure
\(\Pi_{\rm res}\).

Let \(\mathfrak r\in\operatorname{supp}\Pi_{\rm res}\).  Choose a nested
neighborhood basis \((U_m)\) of \(\mathfrak r\) with diameters tending to
zero.  Since \(\Pi_{\rm res}(U_m)>0\), weak-* convergence implies that the
approximating pushforwards assign positive mass to \(U_m\) for all large
\(n\).  For each \(m\), select an actual residual descendant whose rebased
state lies in \(U_m\).  These genuine rebased local states converge to
\(\mathfrak r\), proving the support-realization statement.

The occupation coordinates of \(\mathfrak X_{\rm res}\) come from the product
coordinates of the compact local-state factor, hence are continuous coordinate
projections with closed positive superlevel sets.  Therefore
\cref{paper6a:lem:normalized-defect-atomic-or-diffuse} applies directly to the
probability measure \(\Pi_{\rm res}\).
\end{proof}

\begin{lemma}[Pressure-defect support points are realized by descendant rebasing]
\label{paper8:lem:pressure-defect-support-realization}
Fix once and for all a canonical dyadic family of descendant cylinders inside the
unit cylinder.  For each normalized pressure defect \(\pi_j\) arising from the
nested parent cylinder \(Q_{\rho_j}(z_j)\), record for every dyadic descendant:
its compactified relative scale-center coordinates inside the parent, the
rebased velocity/pressure tuple on a fixed model cylinder, the named failure
faces for pressure atlas, pressure gauge, pressure reconstruction,
compact-upper, local-energy, and rough-core breakdown, and the harmonic jets
of the rebased harmonic remainder on a slightly smaller model cylinder.  By
\cref{paper6a:lem:compact-local-state-factor}, the local-state factor with
these named failure faces is compact metrizable, so the closure of this family
in the product topology is a compact metrizable
pressure-defect state space
\[
  \mathfrak X_{\rm pres}
  :=
  \overline{\mathcal C}_{\rm pres}
  \times \mathfrak O_{\rm pres}
  \times \mathfrak H_{\rm pres},
\]
where \(\overline{\mathcal C}_{\rm pres}\) is the compactified relative
scale-center chart of the descendants, \(\mathfrak O_{\rm pres}\) is the local
occupation-state space obtained from the compact-cylinder suitable bounds,
local energy structure, rough-core data, and the local pressure atlas on the
rebased model cylinder, together with the corresponding named failure faces, and
\(\mathfrak H_{\rm pres}\) is the compact harmonic-jet factor supplied by
interior harmonic estimates.
The local occupation coordinates in \(\mathfrak O_{\rm pres}\) are continuous
coordinate projections of this product compactification; hence each positive
superlevel set of a pressure-defect occupation coordinate is closed in
\(\mathfrak X_{\rm pres}\).

After passing to a subsequence, the pushforwards of \(\pi_j\) to
\(\mathfrak X_{\rm pres}\) converge weakly-* to a probability measure
\(\Pi_{\rm pres}\).  The limiting probability measure itself is retained as
part of the compactification data, so later diffuse arguments may use either
its support points or the diffuse measure itself via
\cref{paper6a:lem:normalized-defect-atomic-or-diffuse}.  Every point
\(\mathfrak p\in\operatorname{supp}\Pi_{\rm pres}\) is realized by an actual
sequence of descendant cylinders
\[
  Q_{\sigma_j\rho_j}(w_j)\subset Q_{\rho_j}(z_j)
\]
and the corresponding rebased suitable solutions from the original branch.
Moreover, on every realizing sequence the decomposition
\[
  P_j=P_j^{\rm loc}+H_j
\]
transports to the rebased cylinders: the rebased \(P_j^{\rm loc}\) is the local
Calder\'on--Zygmund pressure generated by the rebased truncated stress, and the
rebased \(H_j\) is harmonic on a slightly smaller model cylinder.  Hence every
realized support profile inherits the local Calder\'on--Zygmund/harmonic
splitting.
\end{lemma}

\begin{proof}
On each parent cylinder, first normalize to the unit cylinder and then assign to
every point its full itinerary of canonical dyadic descendants.  For a
descendant of relative radius \(\sigma\), translate and dilate by
\(\sigma\rho_j\) and subtract the appropriate pressure mean.  Translation,
dilation, and subtraction of time-dependent means preserve the suitable-solution
equation, the local energy inequality, and the local pressure decomposition on a
fixed interior model cylinder.  The rebased velocity variables therefore lie in
the compact local-state factor furnished by
\cref{paper6a:lem:compact-local-state-factor}, while the rebased harmonic
remainders have compact jet coordinates by interior harmonic estimates.
Together with the compactified relative scale-center chart, this gives the
compact metrizable state space
\(\mathfrak X_{\rm pres}\).

Push \(\pi_j\) forward by this measurable state map.  Since
\(\mathfrak X_{\rm pres}\) is compact metrizable, Prokhorov and Banach--Alaoglu
give a weak-* convergent subsequence \(\Pi_{\rm pres}^{(j)}\to\Pi_{\rm pres}\).
Now let \(\mathfrak p\in\operatorname{supp}\Pi_{\rm pres}\).  Choose a nested
neighborhood basis \((U_m)_{m\ge1}\) of \(\mathfrak p\) with diameters tending
to zero.  Since \(\Pi_{\rm pres}(U_m)>0\), weak-* convergence implies that for
each \(m\) and all sufficiently large \(j\) one has
\(\Pi_{\rm pres}^{(j)}(U_m)>0\).  Pick an index \(j_m\uparrow\infty\) and an
actual descendant cylinder/state from the original sequence whose rebased state
lies in \(U_m\).  By construction, these rebased states converge to
\(\mathfrak p\), so \(\mathfrak p\) is realized by a subsequential rebasing of
the original suitable solutions.

For each realizing descendant cylinder choose a cutoff equal to one on the
smaller rebased model cylinder and supported inside the parent cylinder.  The
usual local pressure atlas gives
\[
  P_j^{\rm rebased}
  =
  \mathcal R_i\mathcal R_j\bigl(\zeta(V_j^{\rm rebased}\otimes
  V_j^{\rm rebased})_{ij}\bigr)+H_j^{\rm rebased}
\]
modulo functions of time, with \(H_j^{\rm rebased}\) harmonic on a slightly
smaller cylinder.  This is exactly the transported
Calder\'on--Zygmund/harmonic splitting claimed above.
\end{proof}

In the pressure-defect compactification we use the same dyadic activity ladder
as in the shell compactification: a nonzero compact local stress or velocity
source is recorded at the first dyadic level \(\eta_m>0\) and first basis test
function for which the corresponding local occupation coordinate is positive.
Thus positive compact source is routed at an all-positive local level, not only
at the original retained threshold.

\begin{theorem}[Pressure-defect compactification and support transfer]
\label{paper8:thm:pressure-defect-compactification-support-transfer}
Let
\[
  Q_{\rho_j}(z_j)
\]
be a nested pressure-only selection with normalized pressure defects
\[
  \pi_j
  :=
  \frac{|P_j-(P_j)_{B_{\rho_j}}(t)|^{3/2}\,dx\,dt}
  {\iint_{Q_{\rho_j}(z_j)}
  |P_j-(P_j)_{B_{\rho_j}}(t)|^{3/2}\,dx\,dt}.
\]
After passing to a subsequence, the induced defect states of
\cref{paper8:lem:pressure-defect-support-realization} converge to a probability
measure \(\Pi_{\rm pres}\) on the compact metrizable pressure-defect state
space \(\mathfrak X_{\rm pres}\).  Every point in
\(\operatorname{supp}\Pi_{\rm pres}\) is realized by a subsequential rebasing of
the original suitable solutions, and every such realized support profile
satisfies exactly one of the following alternatives:
\begin{enumerate}[label=\textup{(\roman*)}]
\item a nonzero local Calder\'on--Zygmund source generated by velocity stress;
\item a non-affine harmonic tail;
\item an affine or pure-gauge harmonic pressure profile;
\item a pressure-atlas, pressure-gauge, pressure-reconstruction,
compact-upper, local-energy, or rough-core failure face.
\end{enumerate}
\end{theorem}

\begin{proof}
Apply \cref{paper8:lem:pressure-defect-support-realization}.  This gives the
compact metrizable state space, the weak-* limit measure
\(\Pi_{\rm pres}\), and an explicit realizing sequence of descendant cylinders
for every support point.

Fix \(\mathfrak p\in\operatorname{supp}\Pi_{\rm pres}\) and let
\((\widetilde V_m,\widetilde P_m)\) be a realizing rebased suitable sequence on
a fixed model cylinder.  If \(\mathfrak p\) lies on a compactification face at
which admissible pressure reconstruction, gauge normalization, the pressure
atlas, compact-cylinder upper control, local energy, or rough-core control
breaks down, then by the ordered state-space test we are in alternative
\textup{(iv)}.
Hence we may assume that the realizing sequence remains in the represented part
of the state space.  The support-realization lemma then supplies the rebased
splitting
\[
  \widetilde P_m=\widetilde P_m^{\rm loc}+\widetilde H_m
\]
modulo functions of time, with \(\widetilde P_m^{\rm loc}\) the local
Calder\'on--Zygmund pressure of the rebased stress and \(\widetilde H_m\)
harmonic on a slightly smaller model cylinder.  Passing to the limit in the
occupation-state and harmonic-jet coordinates yields the same splitting for the
realized support profile.

If the limiting Calder\'on--Zygmund component is nonzero, then some compact
occupation coordinate of the local velocity-stress source is positive.  By the
dyadic activity ladder fixed above, this source is recorded at the first
positive local level \(\eta_m>0\), so alternative \textup{(i)} occurs at an
all-positive local level rather than only at the original retained threshold.
If the
Calder\'on--Zygmund component vanishes, the realized profile is harmonic on the
interior cylinder.  A nonzero non-affine harmonic profile is exactly
alternative \textup{(ii)}, while a pure-gauge harmonic profile or an affine
harmonic profile after admissible accelerating rebasement is
alternative \textup{(iii)}.  Thus one and only one of the represented
support-point alternatives \textup{(i)}--\textup{(iv)} occurs.
\end{proof}

\begin{lemma}[Nested pressure-only tests are well-founded]
\label{paper8:lem:nested-pressure-only-well-founded}
Let
\[
   Q_{\rho_j}(z_j)
\]
be a nested sequence of cylinders selected by the exterior pressure-only test.
Assume that each cylinder carries a pressure lower bound at the pressure-only
threshold and that no velocity carrier has been selected at any previous stage.
Then, after passing to a subsequence, one of the following named rows occurs:
\begin{enumerate}[label=\textup{(\roman*)}]
\item a velocity-stress carrier row;
\item a far-field harmonic pressure-tail row;
\item a pressure-gauge or pressure-reconstruction failure row;
\item a same-point pressure cascade or scale-collapse row;
\item a diffuse pressure-only defect row.
\end{enumerate}
\end{lemma}

\begin{proof}
On each selected cylinder, decompose the pressure into the local
Calder\'on--Zygmund part and the harmonic part:
\[
   P_j=P_j^{\mathrm{loc}}+H_j.
\]
If \(P_j^{\mathrm{loc}}\) carries a fixed fraction of the pressure lower bound, then
the local pressure is generated by the local velocity stress \(V_j\otimes V_j\).
The fixed-annular stress dichotomy gives either a velocity carrier on a
slightly enlarged cylinder or a stress concentration at a smaller nested scale.
In the first case we are in alternative \textup{(i)}.  In the second case the
nested scale test continues.

If \(H_j\) carries a fixed fraction of the pressure lower bound, then the
harmonic interior estimates and the far-field harmonic-tail lemma give either a
far-field harmonic pressure-tail row or a pressure-gauge/reconstruction
failure.  These are alternatives \textup{(ii)} and \textup{(iii)}.

It remains to consider the possibility that the pressure-only test continues
indefinitely through smaller nested cylinders without producing a velocity
carrier, harmonic tail, or gauge failure.  Compare the selected centers and
scales.  If the centers remain at the same physical point and the scales form a
strict nested chain, the branch is a same-point pressure cascade.  If the
relative centers or relative scales fail to have a compact limit, the branch is
a scale-collapse row.  This gives alternative \textup{(iv)}.

Finally, normalize the pressure-defect measures on the nested cylinders and
pass to the compactified pressure-defect state space.  Apply
\cref{paper6a:lem:normalized-defect-atomic-or-diffuse}.  If the atomic/activity
alternative occurs, choose a realized support point with positive compact local
activity.  Then
\cref{paper8:thm:pressure-defect-compactification-support-transfer} routes that
point to one of the previously listed rows \textup{(i)}--\textup{(iv)}.  If the
genuinely diffuse alternative occurs, the nonzero diffuse pressure-defect
measure itself is the diffuse pressure-only row, giving alternative
\textup{(v)}.  The alternatives exhaust all infinite nested pressure-only
selections.
\end{proof}

\begin{lemma}[Diffuse pressure-only defects are not terminal]
\label{paper8:lem:diffuse-pressure-only-not-terminal}
A diffuse pressure-only defect produced by
\cref{paper8:lem:nested-pressure-only-well-founded} cannot be a terminal
retained branch.
\end{lemma}

\begin{proof}
Let \(\pi\) be the normalized diffuse pressure-defect measure.  By
\cref{paper6a:lem:normalized-defect-atomic-or-diffuse} to the normalized
pressure-defect measures.

If the atomic/activity alternative occurs, choose a realized support profile
with positive compact local activity furnished by
\cref{paper8:thm:pressure-defect-compactification-support-transfer}.  If that
support profile has nonzero local Calder\'on--Zygmund source, then some compact
test in the pressure-defect occupation chart detects a positive velocity-stress
level.  By the dyadic activity ladder, this is recorded at a positive local
level \(\eta_m>0\).  The fixed-annular stress dichotomy and the local CKN test
then produce either a regular annular source or a velocity carrier at that
positive level.  The regular source is absorbed by the local pressure atlas;
the carrier case contradicts pressure-only terminality.

If the atomic support profile is harmonic and remains non-affine after
quotienting by time-dependent constants and every admissible local accelerating
rebasement, then the far-field harmonic-tail row occurs.

If the realized atomic harmonic profile is spatially constant, it is removed by
the ordinary pressure gauge.  If it is spatially affine with nonzero linear
part, then \cref{paper8:lem:affine-harmonic-pressure-absorption} either absorbs
that affine forcing into the repaired accelerating-center modulation on a
smaller local cylinder, reducing the harmonic remainder to a time-gauge term,
or else the branch enters the affine-pressure/modulation failure row.  Thus no
terminal pressure-only branch can be supported by an untreated affine harmonic
profile.

If the genuinely diffuse alternative occurs, then the pressure-defect measure
itself is a visible diffuse pressure occupation coordinate.  Apply
\Cref{paper6a:thm:local-residual-exhaustion} to this visible pressure
coordinate.  The output is either a named
pressure/gauge/rough-core/local-energy row or a retained comparable class.  The
retained comparable class is impossible because pressure-only descendants are
not retained velocity classes.

Hence only named rows remain, and no diffuse pressure-only defect is terminal.
\end{proof}

\begin{lemma}[Affine harmonic pressure absorption by local accelerating rebasement]
\label{paper8:lem:affine-harmonic-pressure-absorption}
Let \((u,p)\) be a suitable weak solution on a local cylinder
\(Q:=B_R(x_0)\times I\), and suppose that on a smaller cylinder
\(Q':=B_{R'}(x_0)\times I'\Subset Q\) the harmonic pressure component has the
form
\[
H(x,t)=a(t)\cdot x+c(t)
\]
with \(a,c\in L^1_{\rm loc}(I')\).  Let \(B\in W^{2,1}_{\rm loc}(I';\RR^3)\).
On every compact subinterval, \(\dot B\) is bounded and absolutely continuous.
Assume
\[
\ddot B(t)=-a(t)
\qquad\text{a.e. on }I'.
\]
Define on
\[
\widetilde Q':=\{(y,t):(y+B(t),t)\in Q'\}
\]
the rebased fields
\[
v(y,t):=u(y+B(t),t)-\dot B(t),
\qquad
q(y,t):=p(y+B(t),t)+\ddot B(t)\cdot y.
\]
Then:
\begin{enumerate}[label=\textup{(\roman*)}]
\item \((v,q)\) is again a suitable weak solution on every compact
\(\widetilde Q''\Subset\widetilde Q'\);
\item the transformed harmonic component is spatially constant,
\[
\widetilde H(y,t)=a(t)\cdot B(t)+c(t),
\]
hence removable by the ordinary pressure gauge modulo functions of time;
\item all conclusions are local on compact subcylinders.
\end{enumerate}
If the repaired gauge/modulation package on the branch does not admit this
accelerating rebasement, then the first failed interface is the
affine-pressure/modulation failure row.
\end{lemma}

\begin{proof}
Set \(x=y+B(t)\).  Since \(\dot B\) is spatially constant,
\(\nabla_y\cdot v=\nabla_x\cdot u=0\).  Also
\[
\partial_t v
=\partial_t u(x,t)+\dot B(t)\cdot\nabla_x u(x,t)-\ddot B(t),
\]
while
\[
(v\cdot\nabla_y)v
=(u-\dot B)\cdot\nabla_x u
=u\cdot\nabla_x u-\dot B\cdot\nabla_x u.
\]
Therefore
\[
\partial_t v+(v\cdot\nabla)v-\Delta v
=
\partial_t u+(u\cdot\nabla)u-\Delta u-\ddot B,
\]
and, because
\[
\nabla_y q(y,t)=\nabla_x p(x,t)+\ddot B(t),
\]
the Navier--Stokes equation for \((u,p)\) implies
\[
\partial_t v+(v\cdot\nabla)v-\Delta v+\nabla q=0
\]
distributionally on \(\widetilde Q'\).

Decompose \(p=p_{\rm loc}+H\) on \(Q'\), with
\(H(x,t)=a(t)\cdot x+c(t)\).  Then
\[
q(y,t)
=p_{\rm loc}(y+B(t),t)+a(t)\cdot(y+B(t))+c(t)+\ddot B(t)\cdot y.
\]
Since \(\ddot B=-a\), the \(y\)-affine part cancels and
\[
q(y,t)=p_{\rm loc}(y+B(t),t)+a(t)\cdot B(t)+c(t).
\]
Hence the transformed harmonic remainder is spatially constant.  It is therefore
removable by subtracting a function of time.

It remains to justify suitability.  The transformed test
\[
  \psi(x,t):=\phi(x-B(t),t)
\]
is compactly supported after shrinking the cylinder, but since
\(B\in W^{2,1}_{\rm loc}\), it is in general only absolutely continuous in time.
This is admissible by the same AC-test approximation used for repaired-gauge
charts.

For completeness, choose smooth functions \(B^\varepsilon\to B\) in
\(W^{2,1}\) on the compact time interval of \(\widetilde Q''\), with
\(\dot B^\varepsilon\to\dot B\) uniformly and
\(\ddot B^\varepsilon\to\ddot B\) in \(L^1\).  Define
\[
  \psi^\varepsilon(x,t):=\phi(x-B^\varepsilon(t),t).
\]
Then \(\psi^\varepsilon\in C_c^\infty(Q')\) after shrinking once.  Apply the
local energy inequality for \((u,p)\) to \(\psi^\varepsilon\), perform the
change of variables \(x=y+B^\varepsilon(t)\), and rewrite the result in terms
of
\[
  v^\varepsilon(y,t):=u(y+B^\varepsilon(t),t)-\dot B^\varepsilon(t),
  \qquad
  q^\varepsilon(y,t):=p(y+B^\varepsilon(t),t)+\ddot B^\varepsilon(t)\cdot y.
\]
The transport terms generated by the moving center cancel against the constant
velocity shift and the affine pressure correction.  Passing
\(\varepsilon\to0\) uses the local finite-energy bounds, the uniform convergence
of \(\dot B^\varepsilon\), and the \(L^1\)-convergence of
\(\ddot B^\varepsilon\).  The limit is the local energy inequality for
\((v,q)\) on \(\widetilde Q''\).
\end{proof}

\begin{lemma}[Exterior pressure-only concentration is not terminal]
\label{paper8:lem:exterior-pressure-only-not-terminal}
Let an exterior concentration branch on \(A_{r,R}\) satisfy the pressure-carrier
alternative of
\cref{paper8:lem:exterior-ckn-carrier-or-pressure-route} and suppose that no
velocity carrier occurs on any smaller exterior or separated-core cylinder
after passing to subsequences.  Then the pressure carrier is not an independent
terminal exterior row.  More precisely, after local pressure decomposition on
nested cylinders, one of the following occurs:
\begin{enumerate}[label=\textup{(\roman*)}]
\item the local Calder\'on--Zygmund pressure part forces a velocity-stress
carrier on a slightly larger exterior cylinder;
\item the far-field harmonic pressure tail produces a dyadic exterior
occupation state, which is regular or gives a dyadic CKN carrier by
\cref{paper8:lem:dyadic-exterior-occupation-selection};
\item the pressure reconstruction or pressure gauge fails, which is the ordered
non-exterior pressure row;
\item the pressure-only selection continues through nested cylinders and is
well-foundedly routed by
\cref{paper8:lem:nested-pressure-only-well-founded}, with its diffuse case
excluded by \cref{paper8:lem:diffuse-pressure-only-not-terminal}.
\end{enumerate}
\end{lemma}

\begin{proof}
Work on a nested exterior cylinder and decompose the pressure into the local
Calder\'on--Zygmund part generated by \(\zeta u\otimes u\), the harmonic part,
and a function of time.  If the local Calder\'on--Zygmund part carries the
pressure lower bound, Calder\'on--Zygmund boundedness and the quantitative
converse in \cref{paper8:lem:nested-ball-harmonic-tail} force a nonzero
velocity-stress carrier on a slightly larger exterior cylinder.  Apply the
local split \cref{paper8:lem:exterior-ckn-carrier-or-pressure-route} on that
slightly larger cylinder.  If the resulting carrier is velocity-supported, this
is item \textup{(i)}.
\par
Thus any persistent pressure lower bound must lie in the harmonic remainder.
Apply \cref{paper8:lem:harmonic-cross-pressure-routing} on nested cylinders.
A nonvanishing harmonic remainder is either pressure-reconstruction/gauge
failure, which is item \textup{(iii)}, or it is generated by a far-field tail.
For the far-field tail, use
\cref{paper8:lem:far-field-harmonic-tail-routes-to-dyadic-carrier}.  A fixed
annular stress level is first tested by
\cref{paper8:lem:fixed-annular-stress-ckn-dichotomy}; a selected local CKN
carrier is then split by
\cref{paper8:lem:exterior-ckn-carrier-or-pressure-route}.  Its
velocity-supported subcase gives item \textup{(i)}, while its
pressure-supported subcase is handled by
\cref{paper8:lem:nested-pressure-only-well-founded}.  A regular annular source
is removable and cannot be an independent pressure-only exterior exit.  An
escaping dyadic stress occupation state is tested by
\cref{paper8:lem:dyadic-exterior-occupation-selection}; it is either regular on
far dyadic annuli or gives a dyadic CKN carrier, which is item \textup{(ii)}.
If the local pressure-only selection attempts to continue to smaller nested
cylinders, apply \cref{paper8:lem:nested-pressure-only-well-founded}.
Alternatives \textup{(i)}--\textup{(iv)} there are the already named
velocity-stress, harmonic tail, pressure-gauge, and scale-cascade rows.
Alternative \textup{(v)} is excluded by
\cref{paper8:lem:diffuse-pressure-only-not-terminal}.  Hence an exterior
pressure-only concentration cannot be terminal.
\end{proof}

\begin{lemma}[Locality audit for exterior routing]
\label{paper8:lem:exterior-routing-locality-audit}
The exterior routing lemmas from
\cref{paper8:lem:exterior-ckn-carrier-or-pressure-route} through
\cref{paper8:lem:exterior-pressure-only-not-terminal} use only local annular
CKN tests, finite-overlap covering, local pressure decomposition modulo
functions of time, and the ordered local state-space stratification.  They do
not use a whole-space \(L^3(\mathbb R^3)\) bound, a global critical profile
decomposition, global Kato smallness, or an assumption that localized
dissipation vanishes.
\end{lemma}

\begin{proof}
The CKN split is a local identity on the selected cylinder.  Positive-distance
decoupling uses kernel smoothness on sets separated by a fixed physical
distance, local \(L^1\) bounds for the recorded source, and the ordered
state-space rule that any positive compact-core mass must be retained or routed
before the exterior row is declared.  The exterior and separated-core lemmas use
the local scale-center compactification and the already stated separated-point
and same-point reductions.  The packing lemma uses only the Vitali
finite-overlap property on a fixed annulus and explicitly avoids any cardinality
bound from a global critical norm.  Noncompact exterior routing is a countable
dyadic application of the same local annular test.  A positive fixed-annulus
stress lower bound is not promoted directly to a singular carrier; it first
passes through the local CKN dichotomy of
\cref{paper8:lem:fixed-annular-stress-ckn-dichotomy}.  Exterior scale collapse
is handled by rebasing the active local CKN carrier and then applying the local
scale-collapse first-failure ledger.  Pressure-only routing uses the local
Calder\'on--Zygmund pressure atlas, the nested-ball harmonic-tail estimate, and
the harmonic pressure routing lemma.  None of these steps imports a global
\(L^3\) estimate, a whole-space profile theorem, global Kato perturbation, or
unproved vanishing of localized dissipation.
\end{proof}

\begin{theorem}[Exterior concentration routes to closed or non-exterior rows]
\label{paper8:thm:exterior-concentration-routed}
Fix \(0<r<R<\infty\).  If a Type II branch admits an exterior concentration
branch on \(A_{r,R}\), then that branch is not an independent terminal exterior
exit.  After passing to subsequences and applying the ordered local tests, one
of the following holds:
\begin{enumerate}[label=\textup{(\roman*)}]
\item the exterior contribution is regular or perturbative for the selected
compact core, by \cref{paper8:thm:exterior-regular-removal} or
\cref{paper8:lem:exterior-profile-decouples-from-core};
\item the branch enters the separated-core, multicore, or same-point cascade
rows, which are closed by \cref{paper6:thm:multibubble};
\item the branch enters an exterior scale-collapse row, which after rebasing is
handled by \cref{paper8:lem:exterior-scale-collapse-rebased};
\item the branch enters the noncompact exterior row, which routes to the
preceding alternatives by
\cref{paper8:lem:noncompact-exterior-routes-to-structural-row};
\item the branch enters a pressure-only exterior row, which is converted to an
exterior velocity/dyadic carrier, absorbed as a regular annular or dyadic
source, or routed to an ordered non-exterior pressure row by
\cref{paper8:lem:exterior-pressure-only-not-terminal}.
\end{enumerate}
Consequently, once the non-exterior rows in
\cref{paper6:def:remaining-named-exits} are handled, exterior concentration and
noncompact exterior concentration contribute no additional obstruction.
\end{theorem}

\begin{proof}
Start with the exterior concentration branch and apply the stratification
\cref{paper8:thm:exterior-stratification}.  In the exterior-profile subcase,
\cref{paper8:lem:exterior-profile-decouples-from-core} gives perturbative
decoupling from the selected compact core.  In the separated-core or multicore
subcase, \cref{paper8:lem:separated-exterior-core-is-multicore} places the
branch in the multibubble/cascade package, which is excluded by
\cref{paper6:thm:multibubble}.  If exterior carriers accumulate at one
positive-distance point with strict scale separation, then
\cref{paper8:lem:exterior-same-point-cascade-routing} gives the same cascade
closure.
\par
If the stratification gives noncompact exterior behavior, first apply
\cref{paper8:lem:noncompact-exterior-dyadic-selection}.  The no-concentration
dyadic case is regular and removable by the exterior regularity/removal
theorems.  The dyadic carrier case is routed by
\cref{paper8:lem:noncompact-exterior-routes-to-structural-row} to regular
removal, separated/multicore, cascade, exterior scale-collapse, or
pressure-only routing.  The first three are already covered.  The exterior
scale-collapse subcase is rebased and tested by
\cref{paper8:lem:exterior-scale-collapse-rebased}: finite-cost and retained
compact or scale-rigid scale-collapse subcases are closed by the cited
scale-collapse theorems, while any first failure is one of the ordinary
non-exterior rows in the final ledger.
\par
It remains only to treat the possibility that the CKN lower bound is
pressure-only.  The velocity-pressure split
\cref{paper8:lem:exterior-ckn-carrier-or-pressure-route} either supplies a
velocity carrier, returning to the preceding cases, or gives a pressure
carrier.  By \cref{paper8:lem:exterior-pressure-only-not-terminal}, a
pressure-only carrier either produces a velocity/dyadic carrier, is absorbed
as a regular annular or dyadic source, exposes the ordered pressure row, or
contradicts the pressure lower bound after all local and harmonic tails have
been removed.
Thus every exterior concentration subtype has been routed to a closed
structural row, regular removal, rebased scale-collapse, or one of the
remaining non-exterior first-failure rows.
\par
The proof is local by
\cref{paper8:lem:exterior-routing-locality-audit}.  Hence exterior
concentration is discharged as an independent named exit without adding a
global critical-norm, global profile, global Kato, or hidden dissipation
assumption.
\end{proof}

\begin{corollary}[Exterior concentration exit discharged]
\label{paper8:cor:exterior-concentration-exit-discharged}
Exterior concentration and noncompact exterior concentration are not remaining
independent named exits in the local Type II assembly.  They are either regular
and removable, closed by the structural multibubble/cascade or scale-collapse
packages after rebasing, or assigned to the remaining non-exterior rows of
\cref{paper6:def:remaining-named-exits}.
\end{corollary}

\begin{proof}
This is \cref{paper8:thm:exterior-concentration-routed}.
\end{proof}

\clearpage
\section{Corrected Monotonicity and Scale-Drift Cost Criteria}
\label{sec:corrected-monotonicity-scale-drift}

This section records the local analytic inputs for the localized
cost-divergence exclusion argument on renormalized Type II branches.  The
statements below are local criteria, and the later windowwise closure theorem
verifies the criteria on the canonical retained route.
The statements are written entirely in the variables of the repaired gauge:
\[
  \partial_\tau V+(V\cdot\nabla)V+\nabla P
  =
  \nu\Delta V+a(\tau)(V+y\cdot\nabla V)+b(\tau)\cdot\nabla V,
  \qquad \nabla\cdot V=0.
\]
All pressure terms are understood modulo functions of time.

\subsection{Data for the Cost Criterion}

\begin{definition}[Cost criterion data]\label{paper6a:def:cost-hypotheses}
A renormalized local Type II branch satisfies the cost criterion data if the
following local analytic facts hold.
\begin{enumerate}[label=\textup{(\roman*)}]
\item the branch lies in the local Type II concentration alternative;
\item the physical suitable-solution local energy inequality gives finite
localized energy on every compact selected window used by the branch, with the
renormalized local energy inequality available after the repaired-gauge
pullback; no whole-space energy bound is part of this data;
\item the branch is written in the repaired-gauge variables
\((V,P,a,b)\) above, including the scale, center, pressure normalization, and
modulation data;
\item the localized cost is identified with
\[
  \tilde{\mathfrak D}_{R_0}(\tau)
  =
  \nu\int_{\mathbb R^3}|\nabla V(y,\tau)|^2\phi_{R_0}(y)\,dy
  +a_+(\tau)\int_{\mathbb R^3}|V(y,\tau)|^2\phi_{R_0}(y)\,dy,
\]
or with an explicitly stated tail-comparable nonnegative replacement;
\item a localized energy identity or suitable-solution local energy inequality
is available for the cutoff used in the cost;
\item a corrected localized monotonicity inequality, in the sense of
\cref{paper6a:def:corrected-monotonicity-input}, is available.
\end{enumerate}
\end{definition}

\begin{definition}[Cost-divergence exclusion conclusion]
\label{paper6a:def:cost-divergence-exclusion-conclusion}
A class of renormalized local Type II branches has the cost-divergence
exclusion conclusion if every branch in the class satisfying
\cref{paper6a:def:cost-hypotheses} is excluded from the local Type II class
whenever the corresponding nonnegative localized cost diverges:
\[
  \int_{\tau_0}^{\infty}\mathcal D_{R_0}(\tau)\,d\tau=\infty.
\]
This is a conclusion verified below on the canonical retained route, not an
active input.
\end{definition}

\begin{lemma}[Concentration and localized energy]\label{paper6a:lem:scale-energy-data}
If a branch lies in the local Type II concentration alternative and satisfies
the physical suitable-solution local energy inequality on the compact cylinders
selected by the branch, then it satisfies
\cref{paper6a:def:cost-hypotheses}\textup{(i)} and
\textup{(ii)}.
\end{lemma}

\begin{proof}
The local Type II concentration alternative supplies the selected concentration
branch.  Suitability gives the local energy inequality on each compact physical
cylinder, and \cref{paper6a:lem:ac-pullback-local-energy} transports it to the
renormalized compact windows whenever the repaired-gauge chart is available.
These are precisely items \textup{(i)} and \textup{(ii)}; no global energy
quantity is used.
\end{proof}

\begin{lemma}[Cost identification]\label{paper6a:lem:cost-identification}
Assume the localized cost is the Navier--Stokes cost
\(\tilde{\mathfrak D}_{R_0}\).  Then divergence of
\[
  \int_{\tau_0}^{\infty}\tilde{\mathfrak D}_{R_0}(\tau)\,d\tau
\]
is the cost divergence in \cref{paper6a:def:cost-divergence-exclusion-conclusion}.
\end{lemma}

\begin{proof}
The active cost is assumed to be equal to
\(\tilde{\mathfrak D}_{R_0}\).  Thus the cost appearing in the local
Navier--Stokes estimates and the cost appearing in the local exclusion
conclusion are the same nonnegative functional.  The divergence condition is
therefore identical in the two formulations.
\end{proof}

\begin{lemma}[Parabolic representation class]\label{paper6a:lem:parabolic-class}
If the repaired-gauge representation theorem applies to a Type II branch, then
the branch has the parabolic transport-diffusion structure required in
\cref{paper6a:def:cost-hypotheses}.
\end{lemma}

\begin{proof}
The repaired-gauge theorem supplies the renormalized orbit \((V,P,a,b)\), the
scale and center, the pressure reconstruction modulo functions of time, and the
modulation coefficients.  The displayed equation is a parabolic
Navier--Stokes equation with transport, dilation, and translation drift terms.
Consequently the branch has the parabolic transport-diffusion structure required
by the cost criterion data.
\end{proof}

\begin{definition}[Corrected localized monotonicity input]
\label{paper6a:def:corrected-monotonicity-input}
A renormalized branch has corrected localized monotonicity if there are a
localized renormalized energy \(\mathcal E_{R_0}\), a nonnegative cost
\(\mathcal D_{R_0}\), a constant \(c_0>0\), and a nonnegative remainder
\(\mathcal R_{R_0}\) such that, on the tail of the branch,
\[
  \frac{d}{d\tau}\mathcal E_{R_0}(\tau)
  +c_0\mathcal D_{R_0}(\tau)
  +\mathcal R_{R_0}(\tau)\le0
\]
in distributions.  If the localized energy identity has an integrable error
\(B_{R_0}\), the corrected energy is
\[
  \mathcal E_{R_0}^{\mathrm{corr}}(\tau)
  =
  \mathcal E_{R_0}(\tau)
  +\int_\tau^\infty B_{R_0}(s)\,ds,
\]
and the same inequality is required for
\(\mathcal E_{R_0}^{\mathrm{corr}}\).  Multiplication of a nonnegative cost by
\(c_0>0\) does not change whether its tail integral diverges.
\end{definition}

\begin{lemma}[Energy and Caccioppoli bounds alone are not monotonicity]
\label{paper6a:lem:caccioppoli-not-monotonicity}
The global physical energy inequality and compact-cylinder Caccioppoli
estimates do not, by themselves, imply the corrected localized monotonicity
input of \cref{paper6a:def:corrected-monotonicity-input}.
\end{lemma}

\begin{proof}
The global physical energy inequality controls total physical energy and
physical dissipation, but it does not identify a localized renormalized energy
at the selected Type II scale whose derivative is monotone after pressure flux,
cutoff flux, transport, and modulation terms are included.  The Caccioppoli
estimate gives local spacetime \(H^1\) control from the local energy inequality.
It is an a priori estimate, not a corrected monotonicity formula.  Therefore an
additional absorption or tail-correction argument is needed.
\end{proof}

\begin{proposition}[Cost-divergence exclusion from the recorded conclusion]\label{paper6a:prop:cost-divergence-exclusion}
Let a renormalized Type II branch satisfy
\cref{paper6a:def:cost-hypotheses} and
let its ambient class have the cost-divergence exclusion conclusion of
\cref{paper6a:def:cost-divergence-exclusion-conclusion}.  If its nonnegative localized
cost has divergent tail integral, then the branch is excluded from the local
Type II class.
\end{proposition}

\begin{proof}
\Cref{paper6a:lem:scale-energy-data} gives the local Type II concentration and
bounded-energy inputs.  \Cref{paper6a:lem:cost-identification} identifies the
localized Navier--Stokes cost with the cost used in the exclusion criterion.
\Cref{paper6a:lem:parabolic-class} gives the parabolic transport-diffusion class.
The corrected monotonicity input is one item of
\cref{paper6a:def:cost-hypotheses}.  Once the cost diverges,
the conclusion recorded in \cref{paper6a:def:cost-divergence-exclusion-conclusion} gives
the stated exclusion.
\end{proof}

\subsection{Fixed-Cutoff Monotonicity}

Let \(\phi=\phi_{R_0}\in C_c^\infty(\mathbb R^3)\), \(0\le\phi\le1\), be the
radial nonincreasing cutoff used in the localized cost.  Define
\[
  \mathcal E_\phi(\tau)=\frac12\int_{\mathbb R^3}|V(y,\tau)|^2\phi(y)\,dy,
\]
\[
  A_\phi(\tau)=\int_{\mathbb R^3}|\nabla V(y,\tau)|^2\phi(y)\,dy,\qquad
  M_\phi(\tau)=\int_{\mathbb R^3}|V(y,\tau)|^2\phi(y)\,dy,
\]
and
\[
  \tilde{\mathfrak D}_{R_0}(\tau)=\nu A_\phi(\tau)+a_+(\tau)M_\phi(\tau),
  \qquad a_+(\tau)=\max\{a(\tau),0\}.
\]
The localized energy identity used below is
\begin{align}
  \frac{d}{d\tau}\mathcal E_\phi
  ={}&-\nu A_\phi
  +\frac{\nu}{2}\int |V|^2\Delta\phi
  +\frac12\int |V|^2V\cdot\nabla\phi
  +\int P\,V\cdot\nabla\phi \notag\\
  &-\frac{a}{2}M_\phi
  -\frac{a}{2}\int |V|^2y\cdot\nabla\phi
  -\frac12 b\cdot\int |V|^2\nabla\phi .
  \label{paper6a:eq:fixed-local-energy}
\end{align}

\begin{definition}[Fixed-cutoff error density]\label{paper6a:def:fixed-error-density}
Set \(a_-(\tau)=\max\{-a(\tau),0\}\) and define
\begin{align*}
  B_{R_0}(\tau)
  :={}&
  \frac{\nu}{2}\left|\int |V|^2\Delta\phi\right|
  +\frac12\left|\int |V|^2V\cdot\nabla\phi\right|
  +\left|\int P\,V\cdot\nabla\phi\right|\\
  &+\frac12|b(\tau)|\left|\int |V|^2\nabla\phi\right|
  +\frac12a_-(\tau)M_\phi(\tau)
  +\frac12a_+(\tau)\int |V|^2\bigl(-y\cdot\nabla\phi\bigr).
\end{align*}
The fixed-cutoff finite-error condition is
\[
  \int_{\tau_0}^{\infty}B_{R_0}(\tau)\,d\tau<\infty .
\]
\end{definition}

\begin{remark}[Pressure normalization]
The pressure term in \(B_{R_0}\) is independent of the representative of the
pressure modulo functions of time.  Replacing \(P\) by \(P+c(\tau)\) changes
the pressure flux by
\[
  c(\tau)\int V\cdot\nabla\phi
  =
  -c(\tau)\int \phi\,\nabla\cdot V
  =0.
\]
\end{remark}

\begin{theorem}[Fixed-cutoff corrected monotonicity]
\label{paper6a:thm:fixed-corrected-monotonicity}
Assume \eqref{paper6a:eq:fixed-local-energy} holds on
\((\tau_0,\infty)\) and that \(B_{R_0}\in L^1(\tau_0,\infty)\).  Define
\[
  \mathcal E_\phi^{\mathrm{corr}}(\tau)
  =
  \mathcal E_\phi(\tau)+\int_\tau^\infty B_{R_0}(s)\,ds .
\]
Then, in distributions on \((\tau_0,\infty)\),
\[
  \frac{d}{d\tau}\mathcal E_\phi^{\mathrm{corr}}(\tau)
  +\frac12\tilde{\mathfrak D}_{R_0}(\tau)\le0 .
\]
\end{theorem}

\begin{proof}
Add
\[
  \frac12\tilde{\mathfrak D}_{R_0}
  =
  \frac{\nu}{2}A_\phi+\frac12a_+M_\phi
\]
to both sides of \eqref{paper6a:eq:fixed-local-energy}.  The gradient term
becomes
\[
  -\nu A_\phi+\frac{\nu}{2}A_\phi=-\frac{\nu}{2}A_\phi\le0.
\]
The core scale term satisfies
\[
  -\frac{a}{2}M_\phi+\frac12a_+M_\phi
  =
  \frac12a_-M_\phi
  \le B_{R_0}(\tau).
\]
Since \(\phi\) is radial nonincreasing, one has
\[
  y\cdot\nabla\phi(y)\le0
  \qquad\text{for all }y\in\mathbb R^3.
\]
Hence
\[
  -\frac{a}{2}\int |V|^2y\cdot\nabla\phi
  \le
  \frac12a_+(\tau)\int |V|^2\bigl(-y\cdot\nabla\phi\bigr),
  \]
because the \(a_-\)-part contributes a nonpositive term.  Every remaining term
in \eqref{paper6a:eq:fixed-local-energy} is bounded above by its corresponding
absolute-value contribution in \(B_{R_0}\).  Therefore
\[
  \frac{d}{d\tau}\mathcal E_\phi(\tau)
  +\frac12\tilde{\mathfrak D}_{R_0}(\tau)
  \le B_{R_0}(\tau)
\]
in distributions.  Since \(B_{R_0}\in L^1(\tau_0,\infty)\), the correction
tail is finite for every \(\tau\ge\tau_0\), absolutely continuous on compact
subintervals, and
\[
  \frac{d}{d\tau}\int_\tau^\infty B_{R_0}(s)\,ds=-B_{R_0}(\tau).
\]
Adding this identity gives the claimed inequality.
\end{proof}

\begin{corollary}[Fixed-cutoff monotonicity input]\label{paper6a:cor:fixed-input}
If the repaired-gauge representation, the cost identification, the localized
energy identity, and the fixed-cutoff finite-error condition all hold, then
\cref{paper6a:def:corrected-monotonicity-input} holds with
\(\mathcal D_{R_0}=\tilde{\mathfrak D}_{R_0}\) and \(c_0=1/2\).
\end{corollary}

\begin{proof}
The repaired-gauge representation supplies \((V,P,a,b)\).  The cost
identification supplies \(\tilde{\mathfrak D}_{R_0}\).  The localized energy
identity and finite-error condition allow
\cref{paper6a:thm:fixed-corrected-monotonicity} to be applied.
\end{proof}

\begin{remark}[Finite-tail subconditions]\label{paper6a:rem:fixed-tail-subconditions}
The fixed-cutoff finite-error condition is equivalent to the conjunction of tail
integrability for the viscous cutoff term, the convective flux, the pressure
flux, the center-drift cutoff term, the negative-scale term, and the
positive scale-shell term:
\[
\begin{gathered}
  \left|\int |V|^2\Delta\phi\right|,\qquad
  \left|\int |V|^2V\cdot\nabla\phi\right|,\qquad
  \left|\int P\,V\cdot\nabla\phi\right|,\\
  |b|\left|\int |V|^2\nabla\phi\right|,\qquad
  a_-M_\phi,\qquad
  a_+\int |V|^2\bigl(-y\cdot\nabla\phi\bigr).
\end{gathered}
\]
For radial nonincreasing cutoffs, only the positive part \(a_+\) contributes to
the scale-shell error.  The negative part is favorable and does not belong to
the residual finite-error burden.
Uniform compact-window bounds do not by themselves give this global tail
integrability on \([\tau_0,\infty)\).
\end{remark}

\begin{definition}[Fixed-cutoff cost hypotheses]\label{paper6a:def:direct-fixed-data}
A renormalized Type II branch satisfies the fixed-cutoff cost hypotheses if it
has the local Type II concentration alternative, bounded physical energy,
repaired-gauge representation, cost identification with
\(\tilde{\mathfrak D}_{R_0}\), divergent nonnegative localized cost, localized
energy identity, and the finite-tail bound
\(\int_{\tau_0}^{\infty}B_{R_0}<\infty\).
\end{definition}

\begin{theorem}[Fixed-cutoff cost criterion]\label{paper6a:thm:fixed-cost-criterion}
Every branch satisfying the fixed-cutoff cost hypotheses is covered by
\cref{paper6a:prop:cost-divergence-exclusion}.
\end{theorem}

\begin{proof}
By \cref{paper6a:cor:fixed-input}, the finite-tail condition supplies the
corrected monotonicity input.  The remaining inputs are precisely the other
items in \cref{paper6a:def:cost-hypotheses}, together with cost divergence.
The local exclusion step is later proved on the canonical windowwise
schedule in \cref{paper6a:thm:cost-divergence-exclusion-discharged}.
\end{proof}

\begin{definition}[Compact cost-divergence replacement]\label{paper6a:def:compact-cost-replacement}
A compact cost-divergence replacement is available when the compact Type II hypotheses
of the preceding sections imply
\[
  \int_{\tau_0}^{\infty}\tilde{\mathfrak D}_{R_0}(\tau)\,d\tau=\infty,
\]
and this divergence is the cost divergence used in
\cref{paper6a:prop:cost-divergence-exclusion}.  In particular, it may be
supplied by the combination of classification of local alternatives, uniform
critical tightness, local windowed \(H^1\) control, and the cost identification
above whenever those are the active hypotheses in the compact Type II theorem.
\end{definition}

\begin{theorem}[Compact cost-divergence criterion]\label{paper6a:thm:compact-cost-criterion}
If a renormalized branch has all fixed-cutoff cost hypotheses except that cost
divergence is supplied through the compact cost-divergence replacement of
\cref{paper6a:def:compact-cost-replacement}, then
\cref{paper6a:prop:cost-divergence-exclusion} applies.
\end{theorem}

\begin{proof}
The compact cost-divergence replacement supplies the missing divergent cost.
Substituting this into the fixed-cutoff hypotheses gives the hypotheses of
\cref{paper6a:thm:fixed-cost-criterion}.
\end{proof}

\subsection{Moving-Cutoff Monotonicity}

Let \(\phi\in C_c^\infty(\mathbb R^3)\) be radial and nonincreasing,
\(0\le\phi\le1\), \(\phi=1\) on \(B_1\), and \(\phi=0\) outside \(B_2\).  For
\(R(\tau)\ge R_0\), locally absolutely continuous and nondecreasing, set
\[
  \Phi(\tau,y)=\phi(y/R(\tau)),\qquad
  A_{R(\tau)}=\{R(\tau)\le |y|\le2R(\tau)\},
\]
\[
  \mathcal E_R(\tau)=\frac12\int |V(y,\tau)|^2\Phi(\tau,y)\,dy,
\]
and
\[
  \tilde{\mathfrak D}_{R(\tau)}(\tau)
  =
  \nu\int|\nabla V|^2\Phi
  +a_+(\tau)\int|V|^2\Phi .
\]

\begin{definition}[Moving-cutoff error density]\label{paper6a:def:moving-error-density}
Define
\begin{align*}
  B_{\mathrm{mov}}(\tau)
  :={}&
  \frac12\left|\int |V|^2\partial_\tau\Phi\right|
  +\frac{\nu}{2}\left|\int |V|^2\Delta\Phi\right|
  +\frac12\left|\int |V|^2V\cdot\nabla\Phi\right|
  +\left|\int P\,V\cdot\nabla\Phi\right|\\
  &+\frac12|b(\tau)|\left|\int |V|^2\nabla\Phi\right|
  +\frac12a_-(\tau)\int |V|^2\Phi
  +\frac12a_+(\tau)\int |V|^2\bigl(-y\cdot\nabla\Phi\bigr).
\end{align*}
The moving finite-error condition is
\[
  \int_{\tau_0}^{\infty}B_{\mathrm{mov}}(\tau)\,d\tau<\infty .
\]
\end{definition}

\begin{theorem}[Moving-cutoff corrected monotonicity]
\label{paper6a:thm:moving-corrected-monotonicity}
Assume the localized energy identity is valid with the time-dependent cutoff
\(\Phi\), and assume \(B_{\mathrm{mov}}\in L^1(\tau_0,\infty)\).  Then
\[
  \mathcal E_R^{\mathrm{corr}}(\tau)
  =
  \mathcal E_R(\tau)+\int_\tau^\infty B_{\mathrm{mov}}(s)\,ds
\]
satisfies
\[
  \frac{d}{d\tau}\mathcal E_R^{\mathrm{corr}}(\tau)
  +\frac12\tilde{\mathfrak D}_{R(\tau)}(\tau)\le0
\]
in distributions.
\end{theorem}

\begin{proof}
Use the localized energy identity with the time-dependent test function
\(\Phi(\tau,y)\).  Compared with the fixed-cutoff identity, the only new term
is
\[
  \frac12\int |V|^2\partial_\tau\Phi,
\]
which comes from differentiating the cutoff inside \(\mathcal E_R\).  All
other terms are the fixed-cutoff terms with \(\phi\) replaced by \(\Phi\).  Add
\(\frac12\tilde{\mathfrak D}_{R(\tau)}\) to both sides.  The gradient term
leaves
\[
  -\frac{\nu}{2}\int|\nabla V|^2\Phi\le0,
\]
and the core scale term contributes
\[
  \frac12a_-\int |V|^2\Phi .
\]
Since \(\phi\) is radial nonincreasing, one has
\[
  y\cdot\nabla\Phi(\tau,y)\le0
  \qquad\text{for all }(\tau,y).
\]
Therefore
\[
  -\frac{a}{2}\int |V|^2y\cdot\nabla\Phi
  \le
  \frac12a_+(\tau)\int |V|^2\bigl(-y\cdot\nabla\Phi\bigr),
\]
because the \(a_-\)-part is nonpositive.  Each remaining term is bounded above
by the corresponding absolute-value term in \(B_{\mathrm{mov}}\).  Hence
\[
  \frac{d}{d\tau}\mathcal E_R(\tau)
  +\frac12\tilde{\mathfrak D}_{R(\tau)}(\tau)
  \le B_{\mathrm{mov}}(\tau).
\]
The moving finite-error condition makes the correction tail finite and
absolutely continuous on compact subintervals.  Differentiating that tail
subtracts \(B_{\mathrm{mov}}\), giving the stated inequality.
\end{proof}

\begin{definition}[Summable moving-annulus errors]\label{paper6a:def:summable-moving-errors}
The moving-annulus summability condition is the existence of a nondecreasing
locally absolutely continuous \(R(\tau)\to\infty\) such that each of the seven
quantities
\[
\begin{gathered}
  \left|\int |V|^2\partial_\tau\Phi\right|,\quad
  \left|\int |V|^2\Delta\Phi\right|,\quad
  \left|\int |V|^2V\cdot\nabla\Phi\right|,\quad
  \left|\int P\,V\cdot\nabla\Phi\right|,\\
  |b|\left|\int |V|^2\nabla\Phi\right|,\quad
  a_-\int |V|^2\Phi,\quad
  a_+\int |V|^2\bigl(-y\cdot\nabla\Phi\bigr)
\end{gathered}
\]
is integrable on \([\tau_0,\infty)\).
\end{definition}

\begin{lemma}[Summable annuli imply moving finite error]
\label{paper6a:lem:summable-moving-errors}
If \cref{paper6a:def:summable-moving-errors} holds, then
\(B_{\mathrm{mov}}\in L^1(\tau_0,\infty)\).
\end{lemma}

\begin{proof}
The density \(B_{\mathrm{mov}}\) is a finite positive linear combination of
the seven quantities listed in
\cref{paper6a:def:summable-moving-errors}.  If each is integrable, then their
linear combination is integrable.
\end{proof}

\begin{definition}[Summable tightness schedule]\label{paper6a:def:summable-tightness}
A summable tightness schedule is a nondecreasing sequence \(R_n\to\infty\) such
that, on unit windows \(I_n=[\tau_0+n,\tau_0+n+1]\),
\[
  \sum_{n=0}^{\infty}
  \sup_{\tau\in I_n}
  \|V(\tau)\|_{L^3(A_{R_n})}^3<\infty,
\]
and the same schedule is compatible with the pressure-tail normalization used
for \(P\).  This is stronger than uniform critical tightness, which permits
large radii on each window but does not by itself give a monotone radius
schedule with a summable series.
\end{definition}

\begin{definition}[Moving-cost compatibility]\label{paper6a:def:moving-cost-compatibility}
The moving localized cost is compatible with the cost-divergence criterion if
\[
  \tilde{\mathfrak D}_{R(\tau)}(\tau)
  =
  \nu\int|\nabla V|^2\phi_{R(\tau)}
  +a_+(\tau)\int|V|^2\phi_{R(\tau)}
\]
is the localized cost appearing in
\cref{paper6a:def:cost-divergence-exclusion-conclusion}, or is compared to that cost by a
positive constant on the tail so that cost divergence is unchanged.
\end{definition}

\begin{theorem}[Moving-cutoff replacement]\label{paper6a:thm:moving-replacement}
Assume the repaired-gauge representation, localized energy identity,
moving-cost compatibility, and summable moving-annulus errors.  Then the moving
cutoff supplies the corrected monotonicity input needed in
\cref{paper6a:prop:cost-divergence-exclusion}.
\end{theorem}

\begin{proof}
\Cref{paper6a:lem:summable-moving-errors} gives
\(B_{\mathrm{mov}}\in L^1(\tau_0,\infty)\).  Applying
\cref{paper6a:thm:moving-corrected-monotonicity} gives corrected monotonicity
with cost \(\tilde{\mathfrak D}_{R(\tau)}\).  The moving-cost compatibility
input identifies this cost, or a tail-comparable version of it, with the
localized cost in the recorded cost-divergence conclusion.
\end{proof}

\subsection{Negative Scale Drift}

Set
\[
  M_R(\tau)=\int |V(y,\tau)|^2\Phi(\tau,y)\,dy .
\]
The negative-scale contribution in the moving cutoff argument is
\[
  \int_{\tau_0}^{\infty}a_-(\tau)M_R(\tau)\,d\tau.
\]

\begin{definition}[Negative scale-drift alternatives]\label{paper6a:def:negative-drift}
We say that the branch has finite negative drift if
\[
  \int_{\tau_0}^{\infty}a_-(\tau)\,d\tau<\infty.
\]
It has bounded moving \(L^2\) core mass if
\[
  \sup_{\tau\ge\tau_0}M_R(\tau)<\infty.
\]
It has a moving \(L^2\) core floor if there are \(m_0>0\) and
\(\tau_1\ge\tau_0\) such that
\[
  M_R(\tau)\ge m_0
  \qquad\text{for a.e. }\tau\ge\tau_1.
\]
The scale-collapse drift obstruction is the divergent alternative
\[
  \int_{\tau_0}^{\infty}a_-(\tau)M_R(\tau)\,d\tau=\infty .
\]
\end{definition}

\begin{lemma}[Finite negative drift controls the scale term]
\label{paper6a:lem:finite-negative-drift}
If the branch has bounded moving \(L^2\) core mass and finite negative drift,
then
\[
  \int_{\tau_0}^{\infty}a_-(\tau)M_R(\tau)\,d\tau<\infty .
\]
\end{lemma}

\begin{proof}
Let \(M_*=\sup_{\tau\ge\tau_0}M_R(\tau)<\infty\).  Since \(a_-\ge0\),
\[
  \int_{\tau_0}^{\infty}a_-(\tau)M_R(\tau)\,d\tau
  \le
  M_*\int_{\tau_0}^{\infty}a_-(\tau)\,d\tau<\infty.
\]
\end{proof}

\begin{theorem}[Finite-or-infinite scale-negative alternative]
\label{paper6a:thm:finite-infinite-scale-negative}
For every renormalized orbit for which \(a_-\) and \(M_R\) are measurable and
nonnegative, precisely one of the following alternatives holds:
\[
  \int_{\tau_0}^{\infty}a_-(\tau)M_R(\tau)\,d\tau<\infty,
  \qquad
  \int_{\tau_0}^{\infty}a_-(\tau)M_R(\tau)\,d\tau=\infty.
\]
\end{theorem}

\begin{proof}
The integral
\[
  I_R=\int_{\tau_0}^{\infty}a_-(\tau)M_R(\tau)\,d\tau
\]
is an extended nonnegative number.  Hence either \(I_R<\infty\) or
\(I_R=\infty\), and the two cases are mutually exclusive.
\end{proof}

\begin{lemma}[Logarithmic scale identity]\label{paper6a:lem:log-scale}
For the repaired-gauge convention
\[
  \tau_t=\lambda(t)^{-2},\qquad a(\tau)=\lambda(t)\lambda_t(t),
\]
one has
\[
  \int_{\tau_1}^{\tau_2}a(\tau)\,d\tau
  =
  \log\lambda(\tau_2)-\log\lambda(\tau_1)
\]
whenever \(\tau_0\le\tau_1<\tau_2<\infty\).
\end{lemma}

\begin{proof}
The chain rule gives
\[
  \frac{d}{d\tau}\log\lambda
  =
  \frac{\lambda_t}{\lambda}\frac{dt}{d\tau}
  =
  \frac{\lambda_t}{\lambda}\lambda^2
  =
  \lambda\lambda_t=a.
\]
Integrating over \([\tau_1,\tau_2]\) proves the identity.
\end{proof}

\begin{theorem}[Collapse with core floor forces scale-drift obstruction]
\label{paper6a:thm:collapse-core-floor}
Assume
\[
  \lambda(\tau)\to0\qquad\text{as }\tau\to\infty,
\]
and assume the moving \(L^2\) core floor of
\cref{paper6a:def:negative-drift}.  Then
\[
  \int_{\tau_0}^{\infty}a_-(\tau)M_R(\tau)\,d\tau=\infty .
\]
\end{theorem}

\begin{proof}
By \cref{paper6a:lem:log-scale},
\[
  \int_{\tau_1}^{\tau_2}a(\tau)\,d\tau
  =
  \log\lambda(\tau_2)-\log\lambda(\tau_1).
\]
Since \(\lambda(\tau)\to0\), the right-hand side tends to \(-\infty\) as
\(\tau_2\to\infty\).  If
\(\int_{\tau_1}^{\infty}a_-(\tau)\,d\tau<\infty\), then
\[
  \int_{\tau_1}^{\tau_2}a(\tau)\,d\tau
  =
  \int_{\tau_1}^{\tau_2}a_+(\tau)\,d\tau
  -
  \int_{\tau_1}^{\tau_2}a_-(\tau)\,d\tau
  \ge
  -\int_{\tau_1}^{\infty}a_-(\tau)\,d\tau>-\infty,
\]
which contradicts the preceding limit.  Thus
\(\int_{\tau_1}^{\infty}a_-(\tau)\,d\tau=\infty\).

By the moving \(L^2\) core floor, \(M_R(\tau)\ge m_0\) for a.e.
\(\tau\ge\tau_1\).  Therefore
\[
  \int_{\tau_0}^{\infty}a_-(\tau)M_R(\tau)\,d\tau
  \ge
  m_0\int_{\tau_1}^{\infty}a_-(\tau)\,d\tau
  =
  \infty.
\]
\end{proof}

\begin{corollary}[Scale-drift alternatives]\label{paper6a:cor:scale-drift-alternatives}
In the moving-cutoff setup, the negative-scale term has the following
exhaustive alternatives.  If bounded moving \(L^2\) core mass and finite negative
drift hold, then the negative-scale term is integrable.  If the weighted
negative-drift integral is infinite, the branch lies in the scale-collapse
drift obstruction.  If genuine collapse and a moving \(L^2\) core floor also
hold, then the obstruction is forced.
\end{corollary}

\begin{proof}
The first case is \cref{paper6a:lem:finite-negative-drift}.  The exhaustive
split is \cref{paper6a:thm:finite-infinite-scale-negative}.  The final
assertion is \cref{paper6a:thm:collapse-core-floor}.
\end{proof}

\subsection{Scale-Drift Cost Criteria}

\begin{definition}[Scale-collapse drift cost]\label{paper6a:def:scale-collapse-cost}
For a renormalized branch define
\[
  \mathfrak C_{\mathrm{sc}}(\tau)
  =
  a_-(\tau)M_R(\tau),
  \qquad
  M_R(\tau)=\int |V(y,\tau)|^2\phi_{R(\tau)}(y)\,dy,
\]
and
\[
  \mathcal C_{\mathrm{sc}}[\tau_0,\infty)
  =
  \int_{\tau_0}^{\infty}\mathfrak C_{\mathrm{sc}}(\tau)\,d\tau .
\]
Thus the scale-collapse drift obstruction is precisely
\(\mathcal C_{\mathrm{sc}}[\tau_0,\infty)=\infty\).  The density is
nonnegative and local in the renormalized core.
\end{definition}

\begin{definition}[Scale-collapse cost compatibility]\label{paper6a:def:scale-collapse-cost-compatibility}
The scale-collapse drift cost is compatible with the cost-divergence criterion
when:
\begin{enumerate}[label=\textup{(\roman*)}]
\item \(\mathfrak C_{\mathrm{sc}}\) is measurable, nonnegative, and locally
integrable on finite \(\tau\)-intervals for renormalized branches;
\item divergence of \(\mathcal C_{\mathrm{sc}}[\tau_0,\infty)\) is one of the
divergent-cost cases covered by the conclusion recorded in
\cref{paper6a:def:cost-divergence-exclusion-conclusion};
\item the resulting conclusion is the same local Type II exclusion conclusion
used for the fixed localized cost.
\end{enumerate}
\end{definition}

\begin{theorem}[Scale-collapse drift cost criterion]
\label{paper6a:thm:scale-collapse-drift-cost}
If the scale-collapse drift obstruction holds and the cost
\(\mathfrak C_{\mathrm{sc}}\) is compatible with the cost-divergence criterion, then the divergent-cost hypothesis
of \cref{paper6a:prop:cost-divergence-exclusion} is satisfied for the
scale-collapse branch.
\end{theorem}

\begin{proof}
By \cref{paper6a:def:scale-collapse-cost}, the obstruction is precisely
\[
  \mathcal C_{\mathrm{sc}}[\tau_0,\infty)=\infty .
\]
By \cref{paper6a:def:scale-collapse-cost-compatibility}, this divergence is a
nonnegative cost-divergence hypothesis with the same local conclusion.
\end{proof}

\begin{corollary}[Scale-collapse exclusion by drift cost]
\label{paper6a:cor:scale-collapse-drift-exclusion}
Assume a renormalized local Type II branch has scale-collapse drift
obstruction, compatible scale-collapse drift cost, and the remaining cost
criterion data.  If
\cref{paper6a:def:cost-divergence-exclusion-conclusion} holds, then the branch is excluded
from the local Type II class under consideration.
\end{corollary}

\begin{proof}
\Cref{paper6a:thm:scale-collapse-drift-cost} supplies the divergent
localized cost.  Applying
\cref{paper6a:prop:cost-divergence-exclusion} gives the exclusion.
\end{proof}

\begin{definition}[Absolute scale-drift cost]\label{paper6a:def:absolute-scale-cost}
Define
\[
  \mathfrak C_{\mathrm{abs}}(\tau)
  =
  \nu\int|\nabla V|^2\phi_{R(\tau)}
  +|a(\tau)|M_R(\tau).
\]
The absolute scale-drift cost is compatible with the cost-divergence criterion
if divergence of
\(\int_{\tau_0}^{\infty}\mathfrak C_{\mathrm{abs}}(\tau)\,d\tau\) is included
as a Type II cost-divergence condition with the same local conclusion.
\end{definition}

\begin{lemma}[Scale-collapse drift forces absolute-cost divergence]
\label{paper6a:lem:scale-drift-absolute}
If the scale-collapse drift obstruction holds, then
\[
  \int_{\tau_0}^{\infty}\mathfrak C_{\mathrm{abs}}(\tau)\,d\tau=\infty .
\]
\end{lemma}

\begin{proof}
Since
\[
  \mathfrak C_{\mathrm{abs}}(\tau)
  =
  \nu\int|\nabla V|^2\phi_{R(\tau)}
  +|a(\tau)|M_R(\tau)
  \ge a_-(\tau)M_R(\tau),
\]
divergence of \(\int a_-M_R\) forces divergence of
\(\int\mathfrak C_{\mathrm{abs}}\).
\end{proof}

\begin{corollary}[Absolute scale-cost criterion]\label{paper6a:cor:absolute-scale-cost-criterion}
If the scale-collapse drift obstruction holds and the absolute scale-drift
cost is compatible with the cost-divergence criterion, then the divergent-cost hypothesis of
\cref{paper6a:prop:cost-divergence-exclusion} is satisfied.
\end{corollary}

\begin{proof}
\Cref{paper6a:lem:scale-drift-absolute} gives divergence of the absolute
scale-drift cost, and compatibility makes this divergence a Type II
cost-divergence condition.
\end{proof}

\begin{theorem}[Moving cutoff with scale-collapse alternative]
\label{paper6a:thm:moving-cutoff-scale-collapse-alternative}
Assume the moving-cutoff setup and a renormalized local Type II branch.
\begin{enumerate}[label=\textup{(\roman*)}]
\item If the negative-scale term is integrable and the remaining moving-annulus
errors are summable, then the moving-cutoff corrected monotonicity estimate is
available.
\item If the scale-collapse drift obstruction holds and either the
scale-collapse drift cost or the absolute scale-drift cost is compatible with
the cost-divergence criterion, then the divergent-cost hypothesis is satisfied
directly.
\item If, in addition, the remaining local data of
\cref{paper6a:def:cost-hypotheses} and
\cref{paper6a:def:cost-divergence-exclusion-conclusion} hold, then the branch is excluded.
\end{enumerate}
\end{theorem}

\begin{proof}
The first branch is \cref{paper6a:thm:moving-replacement}.  The second branch
is \cref{paper6a:thm:scale-collapse-drift-cost} or
\cref{paper6a:cor:absolute-scale-cost-criterion}, depending on which compatible
cost is used.  The third branch is
\cref{paper6a:prop:cost-divergence-exclusion}.
\end{proof}

\subsection{Cost-Divergence Exclusion on the Canonical Windowwise Schedule}
\label{paper6a:subsec:cost-divergence-canonical-schedule}

The exclusion conclusion recorded in
\cref{paper6a:def:cost-divergence-exclusion-conclusion} is, given the corrected
monotonicity input that is already part of
\cref{paper6a:def:cost-hypotheses}, an arithmetic consequence of the
nonnegativity of the localized renormalized energy together with the
windowwise alternative argument of
\cref{paper7:cor:windowwise-positive-scale-routing-holds,paper7:cor:windowwise-transition-routing-holds}.
We make this implication explicit here so that no cost-divergence input remains
in the assembly.

\begin{lemma}[Nonnegativity of the corrected localized energy]
\label{paper6a:lem:corrected-energy-nonnegative}
Let \(\phi=\phi_{R_0}\in C_c^\infty(\mathbb R^3)\) be a radial nonnegative
cutoff, and let \(R(\tau)\) be either fixed at \(R_0\) or the canonical
windowwise interpolated radius of
\cref{paper7:lem:windowwise-shell-pressure-translation}.  Then
\(\mathcal E_\phi(\tau)=\tfrac12\int|V(y,\tau)|^2\phi(y)\,dy\ge 0\) and
\(\mathcal E_R(\tau)=\tfrac12\int|V(y,\tau)|^2\Phi(\tau,y)\,dy\ge 0\) for
every \(\tau\), and the correction tails
\(\int_\tau^\infty B_{R_0}(s)\,ds\) and
\(\int_\tau^\infty B_{\mathrm{mov}}(s)\,ds\) of
\cref{paper6a:def:fixed-error-density,paper6a:def:moving-error-density} are
nonnegative.  Consequently the corrected energies
\(\mathcal E_\phi^{\mathrm{corr}}\) and
\(\mathcal E_R^{\mathrm{corr}}\) of
\cref{paper6a:thm:fixed-corrected-monotonicity,paper6a:thm:moving-corrected-monotonicity}
are nonnegative.
\end{lemma}

\begin{proof}
Both \(\mathcal E_\phi\) and \(\mathcal E_R\) are integrals of \(|V|^2\)
against a nonnegative weight, hence nonnegative.  Inspecting the seven
summands of \(B_{R_0}\) (\cref{paper6a:def:fixed-error-density}) and the
seven summands of \(B_{\mathrm{mov}}\)
(\cref{paper6a:def:moving-error-density}): five are absolute values of
integrals and are therefore nonnegative; the remaining two are
\(\tfrac12 a_-(\tau)\!\int|V|^2\Phi\) and
\(\tfrac12 a_+(\tau)\!\int|V|^2(-y\cdot\nabla\Phi)\) (and likewise with
\(\Phi\) replaced by \(\phi\)).  In the first, \(a_-\ge 0\) and the
integrand \(|V|^2\Phi\ge 0\), so the term is nonnegative.  In the second,
\(a_+\ge 0\), and \(-y\cdot\nabla\Phi\ge 0\) because
\(\Phi(\tau,y)=\phi(y/R(\tau))\) with \(\phi\) radial nonincreasing forces
\(\nabla\Phi(\tau,y)\) to point opposite to \(y\); the term is nonnegative.
Hence \(B_{R_0},B_{\mathrm{mov}}\ge 0\) pointwise, and their tail integrals
are nonnegative.  Adding a nonnegative tail to a nonnegative quantity
preserves nonnegativity, so
\(\mathcal E_\phi^{\mathrm{corr}}\) and \(\mathcal E_R^{\mathrm{corr}}\) are
nonnegative.
\end{proof}

\begin{lemma}[Monotonicity versus cost divergence]
\label{paper6a:lem:corrected-monotonicity-versus-divergent-cost}
Let \(\mathcal E^{\mathrm{corr}}\colon[\tau_0,\infty)\to[0,\infty)\) be
locally integrable with finite initial value
\(\mathcal E^{\mathrm{corr}}(\tau_0)<\infty\), let
\(\mathcal D\colon[\tau_0,\infty)\to[0,\infty)\) be locally integrable, and
let \(c_0>0\).  If
\[
  \frac{d}{d\tau}\mathcal E^{\mathrm{corr}}(\tau)+c_0\mathcal D(\tau)\le 0
  \qquad\text{in distributions on }(\tau_0,\infty),
\]
then
\[
  \int_{\tau_0}^{\infty}\mathcal D(\tau)\,d\tau
  \le \frac{1}{c_0}\,\mathcal E^{\mathrm{corr}}(\tau_0)<\infty .
\]
In particular, the simultaneous validity of corrected monotonicity in this
sense and of \(\int_{\tau_0}^{\infty}\mathcal D=\infty\) is contradictory.
The same conclusion holds if the inequality is strengthened to the form
\((d/d\tau)\mathcal E^{\mathrm{corr}}+c_0\mathcal D+\mathcal R\le 0\) with
\(\mathcal R\ge 0\), as in
\cref{paper6a:def:corrected-monotonicity-input}.
\end{lemma}

\begin{proof}
The distributional inequality
\((\mathcal E^{\mathrm{corr}})'+c_0\mathcal D\le 0\) on
\((\tau_0,\infty)\), with \(\mathcal D\ge 0\) and
\(\mathcal E^{\mathrm{corr}}\ge 0\) locally integrable, implies that
\(\mathcal E^{\mathrm{corr}}\) admits a non-increasing right-continuous
representative.  Testing the inequality against a smooth approximation of
\(\mathbf 1_{[\tau_0,\tau]}\) and passing to the limit using
\(\mathcal D\ge 0\) and \(\mathcal E^{\mathrm{corr}}\ge 0\) (so the
boundary terms have a definite sign) yields
\[
  c_0\int_{\tau_0}^{\tau}\mathcal D(\sigma)\,d\sigma
  \le \mathcal E^{\mathrm{corr}}(\tau_0)-\mathcal E^{\mathrm{corr}}(\tau)
  \le \mathcal E^{\mathrm{corr}}(\tau_0).
\]
Letting \(\tau\to\infty\) and using monotone convergence
(\(\mathcal D\ge 0\)) yields the displayed bound.  If
\(\int\mathcal D=\infty\) the left-hand side is infinite while the
right-hand side is finite, a contradiction.
\end{proof}

\begin{theorem}[Cost-divergence/corrected-monotonicity dichotomy on the canonical schedule]
\label{paper6a:thm:cost-divergence-routing-dichotomy}
Let \(\omega\) be a renormalized local Type II branch in the pre-carrier
validated compact-cost regime of
\cref{paper7:def:precarrier-validated-compact-channel}.  Equip \(\omega\)
with the canonical windowwise interpolated cutoff
\(\Phi(\tau,y)=\phi(y/R(\tau))\) of
\cref{paper7:lem:windowwise-shell-pressure-translation}, and let
\(\widetilde{\mathfrak D}_{R(\tau)}\) be the moving localized cost of
\cref{paper6a:def:moving-cost-compatibility}.  Assume in addition the
localized energy identity for the moving cutoff \(\Phi(\tau,\cdot)\) of
\cref{paper6a:def:cost-hypotheses}\textup{(v)} (supplied, e.g., by
\cref{paper7:lem:suitable-energy-gives-carrier-subidentities} on the
regime).  If
\[
  \int_{\tau_0}^{\infty}\widetilde{\mathfrak D}_{R(\tau)}(\tau)\,d\tau=\infty,
\]
then the following sharpened dichotomy holds: at least one of the three
components
\[
  J_1,\qquad J_6,\qquad J_7
\]
fails to be in \(L^1(\tau_1,\infty)\) for every \(\tau_1\ge\tau_0\).  More
precisely, exactly one of the following two mutually exclusive alternatives
holds:
\begin{enumerate}[label=\textup{(D\arabic*)}]
\item at least one of \(J_1\), \(J_6\), \(J_7\) fails to be in
\(L^1(\tau_1,\infty)\) for every \(\tau_1\ge\tau_0\);
\item every \(J_k\), \(k\in\{1,\dots,7\}\), belongs to
\(L^1(\tau_0,\infty)\).
\end{enumerate}
Moreover, case \textup{(D2)} is incompatible with cost divergence, so
divergence of the cost forces case \textup{(D1)}, i.e.~failure of one of
\(J_1\), \(J_6\), \(J_7\).
\end{theorem}

\begin{proof}
The moving error density of \cref{paper6a:def:moving-error-density} satisfies
the identity
\[
  B_{\mathrm{mov}}(\tau)
  =\tfrac12 J_1(\tau)+\tfrac{\nu}{2}J_2(\tau)+\tfrac12 J_3(\tau)
  +J_4(\tau)+\tfrac12 J_5(\tau)+\tfrac12 J_6(\tau)+\tfrac12 J_7(\tau),
\]
with \(J_k\) as in \cref{paper7:def:componentwise-annular-finite-error},
because the seven nonnegative terms in \(B_{\mathrm{mov}}\) coincide with
\(J_1,\dots,J_7\) up to the listed multiplicative constants, all of which
are strictly positive (recall \(\nu>0\) is fixed throughout).  In
particular \(B_{\mathrm{mov}}\in L^1(\tau_0,\infty)\) if and only if every
\(J_k\in L^1(\tau_0,\infty)\).

By \cref{paper7:lem:windowwise-shell-pressure-translation}, on the canonical
windowwise schedule one has automatically
\[
  J_2,\,J_3,\,J_4,\,J_5\in L^1(\tau_0,\infty).
\]
Therefore the only \(J_k\) that can fail to be in \(L^1\) on a tail are
\(J_1\), \(J_6\), \(J_7\).  We obtain the following dichotomy: either
at least one of \(J_1,J_6,J_7\) fails to be in \(L^1(\tau_1,\infty)\) for
every \(\tau_1\ge\tau_0\) (case \textup{(D1)}), or each of \(J_1,J_6,J_7\) is
integrable on \([\tau_0,\infty)\) and hence (combined with the automatic
integrability of \(J_2,J_3,J_4,J_5\)) every \(J_k\) is integrable on
\([\tau_0,\infty)\) (case \textup{(D2)}).

In case \textup{(D2)}, \cref{paper6a:lem:summable-moving-errors} gives
\(B_{\mathrm{mov}}\in L^1(\tau_0,\infty)\), and the localized energy identity
of \cref{paper6a:def:cost-hypotheses}\textup{(v)} together with
\cref{paper6a:thm:moving-corrected-monotonicity} produces
\[
  \frac{d}{d\tau}\mathcal E_R^{\mathrm{corr}}(\tau)
  +\tfrac12\,\widetilde{\mathfrak D}_{R(\tau)}(\tau)\le 0
  \qquad\text{in distributions on }(\tau_0,\infty),
\]
with \(c_0=\tfrac12\) in the sense of
\cref{paper6a:def:corrected-monotonicity-input}.  We verify the finite
initial value locally.  By
\cref{paper7:lem:precarrier-finite-local-energy}, after a harmless terminal
restriction the initial time is a compact-cylinder local-energy time for the
finite ball supporting \(\Phi(\tau_0,\cdot)\).  Hence
\[
  \mathcal E_R(\tau_0)
  =\tfrac12\!\int|V(\tau_0)|^2\Phi(\tau_0,\cdot)\,dy<\infty .
\]
Combined with the case-\textup{(D2)} integrability
\(\|B_{\mathrm{mov}}\|_{L^1(\tau_0,\infty)}<\infty\),
\[
  \mathcal E_R^{\mathrm{corr}}(\tau_0)
  =\mathcal E_R(\tau_0)+\int_{\tau_0}^{\infty}\!\!B_{\mathrm{mov}}(s)\,ds
  <\infty.
\]
By \cref{paper6a:lem:corrected-energy-nonnegative}, \(\mathcal
E_R^{\mathrm{corr}}\ge 0\), so
\cref{paper6a:lem:corrected-monotonicity-versus-divergent-cost} yields
\[
  \int_{\tau_0}^{\infty}\widetilde{\mathfrak D}_{R(\tau)}(\tau)\,d\tau
  \le 2\,\mathcal E_R^{\mathrm{corr}}(\tau_0)<\infty,
\]
contradicting the cost divergence.  Hence case \textup{(D2)} cannot occur,
and case \textup{(D1)} must hold; that is, at least one of \(J_1,J_6,J_7\)
fails to be in \(L^1\) on every later tail.
\end{proof}

\begin{theorem}[Cost-divergence exclusion on the canonical windowwise schedule]
\label{paper6a:thm:cost-divergence-exclusion-discharged}
Fix the canonical windowwise moving family of
\cref{paper7:lem:windowwise-shell-pressure-translation} on the pre-carrier
validated compact-cost regime of
\cref{paper7:def:precarrier-validated-compact-channel}, and assume the
localized energy identity for the moving cutoff
(\cref{paper7:lem:suitable-energy-gives-carrier-subidentities}).  Assume
further that
\textup{(R1)} failure of \(J_6\in L^1(\tau_1,\infty)\) on every later tail
\([\tau_1,\infty)\) places the branch into the negative scale-drift
alternative \textup{(vi)} of
\cref{paper7:def:stratified-moving-carrier-routing}, in the sense of
\cref{paper7:thm:windowwise-routing-gives-stratified-routing} (this is the
same negative-drift routing statement used in that theorem,
realized through \cref{paper6a:def:negative-drift,paper6a:cor:scale-drift-alternatives});
and \textup{(R2)} the named adjacent strata of
\cref{paper7:thm:noncarrier-exits-route} together with the
exterior/translation/positive-scale/negative-drift alternatives
\textup{(iii)}--\textup{(vi)} of
\cref{paper7:def:stratified-moving-carrier-routing} are closed by their
assigned local arguments in the ambient Type II class.
Then for every renormalized local Type II branch \(\omega\) in that regime
satisfying \cref{paper6a:def:cost-hypotheses}, divergence of the localized
cost
\[
  \int_{\tau_0}^{\infty}\widetilde{\mathfrak D}_{R(\tau)}(\tau)\,d\tau=\infty
\]
excludes \(\omega\) from the local Type II class.  Equivalently, the
negative-drift alternative \textup{(R1)} together with the adjacent-stratum
closures \textup{(R2)} imply
\cref{paper6a:def:cost-divergence-exclusion-conclusion} as a theorem on the canonical
windowwise schedule.
\end{theorem}

\begin{proof}
Let \(\omega\) be a candidate branch.  By
\cref{paper6a:thm:cost-divergence-routing-dichotomy}, case \textup{(D2)} is
impossible, so case \textup{(D1)} occurs: at least one of
\(J_1\), \(J_6\), \(J_7\) fails to be in \(L^1(\tau_1,\infty)\) for every
\(\tau_1\ge\tau_0\).

Apply the alternative statement assigned to that index:
\begin{itemize}
\item if \(J_1\) fails, the windowwise transition/leakage alternative of
\cref{paper7:cor:windowwise-transition-routing-holds} (proved
unconditionally on the regime via
\cref{paper7:thm:windowwise-transition-routing-certificate}) places
\(\omega\) into the exterior/leakage alternative \textup{(iii)} of
\cref{paper7:def:stratified-moving-carrier-routing}, entering the named
critical-tightness failure stratum of
\cref{paper7:thm:noncarrier-exits-route};
\item if \(J_6\) fails, hypothesis \textup{(R1)} places \(\omega\) into the
negative scale-drift alternative \textup{(vi)};
\item if \(J_7\) fails, the windowwise positive-scale alternative of
\cref{paper7:cor:windowwise-positive-scale-routing-holds} (proved
unconditionally via the argument
\cref{paper7:thm:positive-scale-certificate-enters-scale-rigid} and
\cref{paper3:prop:scale-rigidity-discharged}) places \(\omega\) into the
positive-scale alternative \textup{(v)}.
\end{itemize}
By hypothesis \textup{(R2)}, each of these named alternatives is closed by
an already proved argument in the ambient Type II class.  Therefore
\(\omega\) is excluded.  This proves the conclusion of
\cref{paper6a:def:cost-divergence-exclusion-conclusion} on the canonical windowwise
schedule under premises \textup{(R1)} and \textup{(R2)}; the next corollary
verifies these premises on the physical local Type II branch.
\end{proof}

\begin{corollary}[Cost-divergence exclusion holds under the same hypothesis set as adjacent closure]
\label{paper6a:cor:cost-divergence-exclusion-unconditional}
The hypotheses \textup{(R1)}--\textup{(R2)} of
\cref{paper6a:thm:cost-divergence-exclusion-discharged} on the canonical
windowwise schedule coincide, after combining with the unconditional
windowwise alternatives
\cref{paper7:cor:windowwise-positive-scale-routing-holds,paper7:cor:windowwise-transition-routing-holds},
with the verified negative-drift routing statement for \(J_6\) and the
adjacent-stratum closure of
\cref{paper7:cor:adjacent-closure-reduced-to-windowwise-routing}.  In
particular, on any ambient Type II class on which the conclusion of
\cref{paper7:cor:adjacent-closure-reduced-to-windowwise-routing} (compact-cost
adjacent-stratum closure) is established,
\cref{paper6a:def:cost-divergence-exclusion-conclusion} is a theorem and may be removed
from the hypothesis lists of
\cref{paper6a:def:abstract-final-data,paper6a:def:expanded-terminal-data}.
For the physical local Type II class considered in the final assembly, this
adjacent closure is established by
\cref{paper7:cor:ambient-adjacent-wrapper-discharged}.
\end{corollary}

\begin{proof}
The dependency order is fixed in
\cref{paper6a:lem:cost-divergence-closure-order}: windowwise routing and
ambient adjacent closure are verified before the cost-divergence theorem is
used.
The windowwise transition/leakage alternative
\cref{paper7:cor:windowwise-transition-routing-holds} and windowwise
positive-scale alternative
\cref{paper7:cor:windowwise-positive-scale-routing-holds} are unconditional
on the regime.  Combined with the negative-drift alternative for \(J_6\)
verified on physical local Type II branches by
\cref{paper7:lem:ambient-negative-drift-branch}, and closed as an independent
terminal row by \cref{paper6a:thm:thin-negative-drift-closed},
the alternative component \textup{(R1)} of
\cref{paper6a:thm:cost-divergence-exclusion-discharged} is satisfied.
The adjacent-stratum closure component \textup{(R2)} is exactly the
closure statement of
\cref{paper7:cor:adjacent-closure-reduced-to-windowwise-routing}; in the
physical local Type II class it is supplied by
\cref{paper7:cor:ambient-adjacent-wrapper-discharged}.  Therefore
\cref{paper6a:thm:cost-divergence-exclusion-discharged} produces exclusion
without any further hypothesis, and
\cref{paper6a:def:cost-divergence-exclusion-conclusion} is a theorem in that ambient
class.
\end{proof}

\begin{remark}[Status of the cost-divergence conclusion]
\label{paper6a:rem:cost-divergence-conclusion-status}
The proof has three independent ingredients:
\begin{enumerate}[label=\textup{(\roman*)}]
\item the trivial arithmetic consequence
\cref{paper6a:lem:corrected-monotonicity-versus-divergent-cost} that
corrected localized monotonicity together with finite initial corrected
energy and \(\mathcal D\ge 0\) forces \(\int\mathcal D<\infty\);
\item the moving error-density identity
\(B_{\mathrm{mov}}=\sum_k c_k J_k\) of
\cref{paper6a:def:moving-error-density,paper7:def:componentwise-annular-finite-error}
(positive linear combination), combined with the canonical schedule
integrability \(J_2,J_3,J_4,J_5\in L^1(\tau_0,\infty)\) of
\cref{paper7:lem:windowwise-shell-pressure-translation};
\item the windowwise alternative for \(J_1\) and \(J_7\)
(\cref{paper7:cor:windowwise-transition-routing-holds,paper7:cor:windowwise-positive-scale-routing-holds},
proved unconditionally on the regime) plus the negative-drift alternative for
\(J_6\) imported from
\cref{paper7:thm:windowwise-routing-gives-stratified-routing}.
\end{enumerate}
Thus \cref{paper6a:def:cost-divergence-exclusion-conclusion} is not an independent
analytic input.  Its content is the conjunction of the
trivial monotonicity-vs-divergence arithmetic and the windowwise alternative
argument that is already used to close adjacent strata in
\cref{paper7:cor:adjacent-closure-reduced-to-windowwise-routing}.
\end{remark}

\begin{lemma}[Acyclic order for cost-divergence closure]
\label{paper6a:lem:cost-divergence-closure-order}
In the physical local Type II assembly, the cost-divergence closure is used only
after the local carrier and adjacent-routing rows have been verified.  The
dependency order is:
\begin{enumerate}[label=\textup{(\roman*)}]
\item the suitable local energy inequality and repaired-gauge pullback supply the
local carrier subidentities on compact windows;
\item the canonical windowwise moving family supplies the unconditional
transition and positive-scale routing alternatives
\cref{paper7:cor:windowwise-transition-routing-holds,paper7:cor:windowwise-positive-scale-routing-holds};
\item nonintegrability of the negative scale-drift component \(J_6\) is routed
by the local negative-drift branch
\cref{paper7:lem:ambient-negative-drift-branch};
\item the preceding rows give ambient adjacent-stratum closure through
\cref{paper7:cor:ambient-adjacent-wrapper-discharged};
\item only then is
\cref{paper6a:thm:cost-divergence-exclusion-discharged} applied to remove a
retained branch with divergent validated localized cost;
\item terminal retained-state elimination is applied after the adjacent and cost
rows have been routed or closed.
\end{enumerate}
No statement in \textup{(i)}--\textup{(iv)} invokes
\cref{paper6a:thm:cost-divergence-exclusion-discharged},
\cref{paper6a:cor:cost-divergence-exclusion-unconditional}, or the final
retained Type II closure theorem.
\end{lemma}

\begin{proof}
Item \textup{(i)} is
\cref{paper6a:lem:ac-pullback-local-energy,paper7:lem:suitable-energy-gives-carrier-subidentities}.
Item \textup{(ii)} is exactly the pair of windowwise routing corollaries cited
there.  Item \textup{(iii)} is
\cref{paper7:lem:ambient-negative-drift-branch}, whose proof routes \(J_6\)
failure through the local scale-drift, scale-collapse, compact
scale-collapse, and scale-rigid packages, not through cost-divergence closure.
Item \textup{(iv)} is
\cref{paper7:cor:ambient-adjacent-wrapper-discharged}; its proof cites the
windowwise routing and negative-drift branch but not the cost-divergence theorem.
The proof of
\cref{paper6a:thm:cost-divergence-exclusion-discharged} then uses those already
verified rows as inputs, which is item \textup{(v)}.  Item \textup{(vi)} is an
ordering convention for the later final assembly: terminal retained-state
elimination is allowed to use the fact that the named adjacent and cost rows
have already been routed or closed, but none of the preceding rows uses terminal
retained-state elimination as an input.
\end{proof}

\subsection{Absolutely Continuous Repaired-Gauge Charts and Renormalized Equations}
\label{subsec:ac-repaired-gauge-renormalized-equations}

The following statements isolate the parts of the repaired-gauge
representation that follow from an absolutely continuous concentration chart
and an absolutely continuous repaired-gauge path.  They do not prove extraction
of the chart from arbitrary data and do not prove solvability of the gauge
constraints.

\begin{definition}[Absolutely continuous concentration chart]
\label{paper6a:def:ac-initial-concentration-chart}
Let \(u,p\) solve the Navier--Stokes equations on
\(\mathbb R^3\times(t_0,T^*)\),
\[
  \partial_tu+(u\cdot\nabla)u+\nabla p=\nu\Delta u,
  \qquad
  \nabla\cdot u=0.
\]
An absolutely continuous concentration chart consists of
absolutely continuous functions
\[
  x_c:(t_0,T^*)\to\mathbb R^3,\qquad
  \lambda:(t_0,T^*)\to(0,\infty),
\]
and a strictly increasing absolutely continuous time \(\rho\) such that
\[
  \rho_t=\lambda^{-2}\quad\text{a.e.},
\]
\[
  y=\frac{x-x_c(t)}{\lambda(t)},\qquad
  u(x,t)=\lambda(t)^{-1}U(y,\rho(t)),
\]
and
\[
  p(x,t)=\lambda(t)^{-2}Q(y,\rho(t))+c(t)
\]
for some scalar function \(c(t)\).  The chart is assumed to have the local
integrability required for the distributional chain rule on compact sets in
the \(y\)-variables.  Define
\[
  A(\rho(t)):=\lambda(t)\lambda_t(t),
  \qquad
  B(\rho(t)):=\lambda(t)x_c'(t),
\]
or equivalently
\[
  A=\partial_\rho\log\lambda,\qquad
  B=\lambda^{-1}\partial_\rho x_c ,
\]
where \(\lambda\) and \(x_c\) are regarded as functions of \(\rho\).
\end{definition}

\begin{lemma}[Initial renormalized equation]
\label{paper6a:lem:initial-renormalized-equation}
Assume Definition~\ref{paper6a:def:ac-initial-concentration-chart}.  Then \(U,Q\) satisfy
\[
  \partial_\rho U+(U\cdot\nabla)U+\nabla Q
  =
  \nu\Delta U+A(\rho)(U+y\cdot\nabla U)+B(\rho)\cdot\nabla U,
  \qquad
  \nabla\cdot U=0
\]
distributionally on compact \(y\)-sets.
\end{lemma}

\begin{proof}
For smooth functions the computation is pointwise.  The distributional case is
obtained by testing against compactly supported functions and using the
assumed chart chain rule.  The spatial derivatives are
\[
  \nabla_xu=\lambda^{-2}\nabla_yU,\qquad
  \Delta_xu=\lambda^{-3}\Delta_yU,
  \qquad
  (u\cdot\nabla_x)u=\lambda^{-3}(U\cdot\nabla_y)U.
\]
Since
\[
  y_t=-\frac{x_c'}{\lambda}-\frac{\lambda_t}{\lambda}y,
  \qquad
  \rho_t=\lambda^{-2},
\]
one has
\[
  \partial_tu
  =
  \lambda^{-3}\partial_\rho U
  -\lambda^{-2}\lambda_t(U+y\cdot\nabla U)
  -\lambda^{-2}x_c'\cdot\nabla U.
\]
Also
\[
  \nabla_xp=\lambda^{-3}\nabla_yQ.
\]
Substitution into the physical Navier--Stokes equation and multiplication by
\(\lambda^3\) gives
\[
  \partial_\rho U
  -\lambda\lambda_t(U+y\cdot\nabla U)
  -\lambda x_c'\cdot\nabla U
  +(U\cdot\nabla)U+\nabla Q
  =
  \nu\Delta U.
\]
Moving the two chart-velocity terms to the right gives the displayed equation.
Finally,
\[
  \nabla_x\cdot u=\lambda^{-2}\nabla_y\cdot U,
\]
so \(\nabla_y\cdot U=0\).
\end{proof}

\begin{definition}[Absolutely continuous repaired-gauge path]
\label{paper6a:def:ac-repaired-gauge-path}
An absolutely continuous repaired-gauge path for the initial chart consists, on
every compact final time interval, of absolutely continuous functions
\[
  \mu(\tau)>0,\qquad q(\tau)\in\mathbb R^3,\qquad \rho=\rho(\tau),
\]
such that
\[
  \rho_\tau=\mu(\tau)^2\quad\text{a.e.},
\]
\[
  V(Y,\tau)=\mu(\tau)U(\mu(\tau)Y+q(\tau),\rho(\tau)),
\]
and
\[
  P(Y,\tau)=\mu(\tau)^2Q(\mu(\tau)Y+q(\tau),\rho(\tau)).
\]
The scale and centering constraints defining the repaired gauge are assumed to
hold along this path.  This hypothesis assumes the path has already been
selected; it does not prove solvability of the gauge constraints.
\end{definition}

\begin{lemma}[Physical chart after gauge repair]
\label{paper6a:lem:repaired-physical-chart}
Assume Definitions~\ref{paper6a:def:ac-initial-concentration-chart}
and~\ref{paper6a:def:ac-repaired-gauge-path}.
Let \(t=t(\rho(\tau))\), and define
\[
  \Lambda(\tau):=\lambda(t(\rho(\tau)))\mu(\tau),
  \qquad
  X(\tau):=x_c(t(\rho(\tau)))+\lambda(t(\rho(\tau)))q(\tau).
\]
Then, on every compact final time interval, \(\Lambda\) and \(X\) are
absolutely continuous, \(\Lambda>0\), and
\[
  \frac{d\tau}{dt}=\Lambda(\tau)^{-2}.
\]
Moreover,
\[
  x=X(\tau)+\Lambda(\tau)Y
\]
is the final physical chart and
\[
  u(x,t)=\Lambda(\tau)^{-1}V(Y,\tau),
  \qquad
  p(x,t)=\Lambda(\tau)^{-2}P(Y,\tau)+c(t).
\]
\end{lemma}

\begin{proof}
The intermediate time \(\rho\) is strictly increasing because
\(\rho_t=\lambda^{-2}>0\) a.e.; hence it has an absolutely continuous inverse
on compact subintervals of its range.  The compositions
\(\lambda(t(\rho(\tau)))\) and \(x_c(t(\rho(\tau)))\) are therefore
absolutely continuous on compact final windows.  Products and sums with the
absolutely continuous functions \(\mu,q\) are absolutely continuous, so
\(\Lambda\) and \(X\) are absolutely continuous.

The identity \(\rho_\tau=\mu^2\) implies
\[
  \frac{d\tau}{d\rho}=\mu^{-2}.
\]
Since \(d\rho/dt=\lambda^{-2}\),
\[
  \frac{d\tau}{dt}
  =
  \frac{d\tau}{d\rho}\frac{d\rho}{dt}
  =
  \mu^{-2}\lambda^{-2}
  =
  (\lambda\mu)^{-2}
  =
  \Lambda^{-2}.
\]
If \(x=X+\Lambda Y\), then
\[
  x=x_c+\lambda q+\lambda\mu Y,
\]
and hence the intermediate coordinate is
\[
  y=\frac{x-x_c}{\lambda}=q+\mu Y.
\]
Therefore
\[
  u(x,t)
  =
  \lambda^{-1}U(q+\mu Y,\rho)
  =
  \lambda^{-1}\mu^{-1}V(Y,\tau)
  =
  \Lambda^{-1}V(Y,\tau),
\]
and similarly
\[
  p(x,t)
  =
  \lambda^{-2}Q(q+\mu Y,\rho)+c(t)
  =
  \lambda^{-2}\mu^{-2}P(Y,\tau)+c(t)
  =
  \Lambda^{-2}P(Y,\tau)+c(t).
\]
\end{proof}

\begin{lemma}[Modulation coefficients in repaired gauge]
\label{paper6a:lem:repaired-modulation-coefficients}
Define
\[
  a(\tau):=\partial_\tau\log\Lambda(\tau),
  \qquad
  b(\tau):=\Lambda(\tau)^{-1}X_\tau(\tau),
\]
where derivatives are understood a.e. on compact final windows.  Under the
hypotheses of Lemma~\ref{paper6a:lem:repaired-physical-chart},
\[
  a(\tau)=\mu(\tau)^2A(\rho(\tau))+\frac{\mu_\tau(\tau)}{\mu(\tau)}
\]
and
\[
  b(\tau)=\mu(\tau)B(\rho(\tau))
  +\mu(\tau)A(\rho(\tau))q(\tau)
  +\frac{q_\tau(\tau)}{\mu(\tau)}.
\]
In particular \(a,b\in L^1_{\rm loc}\) on every compact final time interval.
\end{lemma}

\begin{proof}
Since \(\Lambda=\lambda(\rho(\tau))\mu(\tau)\),
\[
  \partial_\tau\log\Lambda
  =
  \rho_\tau\partial_\rho\log\lambda+\partial_\tau\log\mu
  =
  \mu^2A+\frac{\mu_\tau}{\mu}.
\]
This proves the formula for \(a\).  Next
\[
  X=x_c(\rho(\tau))+\lambda(\rho(\tau))q(\tau).
\]
Using
\[
  \partial_\rho x_c=\lambda B,\qquad
  \partial_\rho\lambda=\lambda A,\qquad
  \rho_\tau=\mu^2,
\]
we get
\[
  X_\tau
  =
  \mu^2\lambda B+\mu^2\lambda Aq+\lambda q_\tau.
\]
Division by \(\Lambda=\lambda\mu\) gives
\[
  \Lambda^{-1}X_\tau
  =
  \mu B+\mu Aq+\frac{q_\tau}{\mu}.
\]

Local integrability follows from absolute continuity and positivity of
\(\mu\).  On every compact final interval, \(\mu\) is bounded above and below
away from zero, \(q\) is bounded, and \(\mu_\tau,q_\tau\in L^1\).  The
\(A\)-terms are locally integrable because
\[
  \int |A(\rho(\tau))|\mu(\tau)^2\,d\tau
  =
  \int |A(\rho)|\,d\rho,
\]
and the extra factors \(\mu\), \(\mu^{-1}\), and \(q\) are bounded on the same
compact window.  Likewise, since \(\mu\) is bounded away from zero,
\[
  \int |B(\rho(\tau))|\mu(\tau)\,d\tau
  \le
  C\int |B(\rho(\tau))|\mu(\tau)^2\,d\tau
  =
  C\int |B(\rho)|\,d\rho,
\]
so the \(B\)-term in \(b\) is locally integrable.
\end{proof}

\begin{lemma}[Renormalized equation in repaired gauge]
\label{paper6a:lem:repaired-gauge-equation}
The final variables \(V,P\) satisfy
\[
  \partial_\tau V+(V\cdot\nabla)V+\nabla P
  =
  \nu\Delta V+a(\tau)(V+Y\cdot\nabla V)+b(\tau)\cdot\nabla V,
  \qquad
  \nabla\cdot V=0
\]
distributionally on compact \(Y\)-sets.
\end{lemma}

\begin{proof}
It suffices to compute in the smooth case and then pass to distributions by
testing against compactly supported functions.  Set
\[
  z=\mu Y+q.
\]
Since \(V(Y,\tau)=\mu U(z,\rho(\tau))\) and
\(\rho_\tau=\mu^2\), the spatial derivatives are
\[
  \nabla_YV=\mu^2\nabla_zU,\qquad
  \Delta_YV=\mu^3\Delta_zU,
\]
\[
  (V\cdot\nabla_Y)V=\mu^3(U\cdot\nabla_z)U,
  \qquad
  \nabla_YP=\mu^3\nabla_zQ.
\]
Differentiation in \(\tau\) gives
\[
\begin{aligned}
  V_\tau
  &=
  \mu_\tau U
  +\mu\mu^2U_\rho
  +\mu(\mu_\tau Y+q_\tau)\cdot\nabla_zU .
\end{aligned}
\]
Using the initial renormalized equation,
\[
  U_\rho
  =
  \nu\Delta_zU-(U\cdot\nabla_z)U-\nabla_zQ
  +A(U+z\cdot\nabla_zU)+B\cdot\nabla_zU,
\]
we obtain
\[
\begin{aligned}
&V_\tau+(V\cdot\nabla_Y)V+\nabla_YP-\nu\Delta_YV \\
&\quad =
\mu_\tau U
{}+\mu^3A(U+z\cdot\nabla_zU)
{}+\mu^3B\cdot\nabla_zU
{}+\mu(\mu_\tau Y+q_\tau)\cdot\nabla_zU .
\end{aligned}
\]
Now
\[
  U=\mu^{-1}V,
  \qquad
  \nabla_zU=\mu^{-2}\nabla_YV,
  \qquad
  z=\mu Y+q.
\]
Therefore
\[
\begin{aligned}
\mu_\tau U+\mu^3A U
&=
\left(\frac{\mu_\tau}{\mu}+\mu^2A\right)V,\\
\mu^3A\,z\cdot\nabla_zU
{}+\mu(\mu_\tau Y+q_\tau)\cdot\nabla_zU
&=
\left(\mu^2A+\frac{\mu_\tau}{\mu}\right)Y\cdot\nabla_YV
{}+\left(\mu Aq+\frac{q_\tau}{\mu}\right)\cdot\nabla_YV,\\
\mu^3B\cdot\nabla_zU
&=
\mu B\cdot\nabla_YV .
\end{aligned}
\]
Combining these identities, the coefficient multiplying
\[
  V+Y\cdot\nabla_YV
\]
is
\[
  \frac{\mu_\tau}{\mu}+\mu^2A,
\]
and the coefficient multiplying \(\nabla_YV\) is
\[
  \mu B+\mu Aq+\frac{q_\tau}{\mu}.
\]
By Lemma~\ref{paper6a:lem:repaired-modulation-coefficients} these coefficients are
\(a\) and \(b\).  Hence the displayed repaired-gauge equation holds in the
distributional sense on compact \(Y\)-sets.  Incompressibility follows from
\[
  \nabla_Y\cdot V=\mu^2\nabla_z\cdot U=0.
\]
\end{proof}

\begin{lemma}[Pressure reconstruction under the charts]
\label{paper6a:lem:pressure-reconstruction-under-charts}
Assume the physical pressure satisfies
\[
  -\Delta_xp=\partial_i\partial_j(u_iu_j)
\]
in distributions.  Then the intermediate pressure satisfies
\[
  -\Delta_yQ=\partial_i\partial_j(U_iU_j)
\]
modulo functions of \(\rho\), and the final pressure satisfies
\[
  -\Delta_YP=\partial_i\partial_j(V_iV_j)
\]
modulo functions of \(\tau\).
\end{lemma}

\begin{proof}
The additive function \(c(t)\) in
\[
  p=\lambda^{-2}Q+c(t)
\]
has zero spatial derivatives.  Since
\[
  \Delta_xp=\lambda^{-4}\Delta_yQ,
  \qquad
  \partial_{x_i}\partial_{x_j}(u_iu_j)
  =
  \lambda^{-4}\partial_{y_i}\partial_{y_j}(U_iU_j),
\]
the intermediate pressure equation follows.  For the repaired variables,
\[
  P(Y)=\mu^2Q(\mu Y+q),
  \qquad
  V(Y)=\mu U(\mu Y+q).
\]
Therefore
\[
  -\Delta_YP
  =
  \mu^4(-\Delta_zQ)(z)
  =
  \mu^4\partial_{z_i}\partial_{z_j}(U_iU_j)(z)
  =
  \partial_{Y_i}\partial_{Y_j}(V_iV_j)(Y).
\]
Adding functions of time to \(Q\) or \(P\) does not change the identity.
\end{proof}

\begin{lemma}[Scaling of local \(L^r\) norms]
\label{paper6a:lem:critical-lp-scaling}
For every \(1\le r<\infty\) for which the norms are finite,
\[
  \|V(\tau)\|_{L^r_Y}^r
  =
  \Lambda(\tau)^{r-3}\|u(t(\tau))\|_{L^r_x}^r.
\]
In particular,
\[
  \|V(\tau)\|_{L^3_Y}=\|u(t(\tau))\|_{L^3_x}.
\]
Thus membership in \(L^3(\mathbb R^3)\), finiteness, infiniteness, and positivity
are invariant under the final chart.
\end{lemma}

\begin{proof}
Using
\[
  u(X+\Lambda Y,t)=\Lambda^{-1}V(Y,\tau),
  \qquad
  dx=\Lambda^3\,dY,
\]
we obtain
\[
  \|u(t)\|_{L^r_x}^r
  =
  \int |\Lambda^{-1}V(Y,\tau)|^r\Lambda^3\,dY
  =
  \Lambda^{3-r}\|V(\tau)\|_{L^r_Y}^r.
\]
Rearranging gives the stated formula.  The case \(r=3\) is scale invariant.
\end{proof}

\begin{lemma}[Regularity of scale, center, and modulation]
\label{paper6a:lem:ac-chart-regularity}
If \(\lambda,x_c,\mu,q\) are absolutely continuous and \(\mu>0\), then on
compact final time intervals the final scale \(\Lambda\), center \(X\), and
coefficients \(a,b\) are absolutely continuous or locally integrable as
specified above: \(\Lambda,X\in W^{1,1}_{\rm loc}\) and
\[
  a,b\in L^1_{\rm loc}.
\]
\end{lemma}

\begin{proof}
This follows from
Lemmas~\ref{paper6a:lem:repaired-physical-chart}
and~\ref{paper6a:lem:repaired-modulation-coefficients}.
Positivity of \(\mu\) and compactness of the time window give a positive lower
bound for \(\mu\), so the divisions by \(\mu\) in the formulas for \(a,b\) are
legitimate.
\end{proof}

\begin{theorem}[Consequences of absolutely continuous repaired-gauge charts]
\label{paper6a:thm:ac-repaired-gauge-consequences}
Assume an absolutely continuous concentration chart and an absolutely
continuous repaired-gauge path in the sense of
Definitions~\ref{paper6a:def:ac-initial-concentration-chart}
and~\ref{paper6a:def:ac-repaired-gauge-path}.  Then the
following analytic conclusions hold:
\begin{enumerate}[label=\textup{(\roman*)}]
\item the initial and repaired renormalized Navier--Stokes equations hold
distributionally on compact renormalized sets;
\item the final physical chart satisfies
\[
  u(x,t)=\Lambda(\tau)^{-1}V(Y,\tau),
  \qquad
  p(x,t)=\Lambda(\tau)^{-2}P(Y,\tau)+c(t),
  \qquad
  \frac{d\tau}{dt}=\Lambda^{-2};
\]
\item pressure reconstruction holds modulo functions of time;
\item the final modulation coefficients are the functions \(a,b\) in
Lemma~\ref{paper6a:lem:repaired-modulation-coefficients};
\item the critical \(L^3\) norm is invariant under the final chart;
\item the final scale, center, and modulation coefficients have the local
absolute-continuity and integrability stated in
Lemma~\ref{paper6a:lem:ac-chart-regularity}.
\end{enumerate}
\end{theorem}

\begin{proof}
Lemma~\ref{paper6a:lem:initial-renormalized-equation} gives the initial
renormalized equation.
Lemmas~\ref{paper6a:lem:repaired-physical-chart}
and~\ref{paper6a:lem:repaired-gauge-equation}
give the final chart and the final repaired-gauge equation with the stated
coefficients.  Lemma~\ref{paper6a:lem:pressure-reconstruction-under-charts} gives
pressure reconstruction modulo functions of time.
Lemma~\ref{paper6a:lem:critical-lp-scaling} gives critical norm
invariance, and Lemma~\ref{paper6a:lem:ac-chart-regularity} gives the
regularity inheritance.
\end{proof}

\begin{corollary}[Representation inputs supplied by the repaired gauge]
\label{paper6a:cor:representation-inputs-repaired-gauge}
Once the absolutely continuous concentration chart and absolutely continuous
repaired-gauge path are available, pressure reconstruction, modulation
coefficient realization, local \(W^{1,1}\) regularity of the chart, and
critical \(L^3\)-invariance are analytic consequences rather than separate
inputs.  The representation inputs that must be supplied before applying this
corollary are:
\begin{enumerate}[label=\textup{(\roman*)}]
\item extraction of an absolutely continuous concentration chart from the
selected Type II branch;
\item solvability of the repaired scale and centering equations on the
admissible profile class;
\item selection of the repaired scale and center as an absolutely continuous
path satisfying
\(\rho_\tau=\mu^2\);
\item terminal admissibility of the final scale and time, including
\(\tau\to\infty\) and preservation of the Type II core scale.
\end{enumerate}
\end{corollary}

\begin{proof}
This is Theorem~\ref{paper6a:thm:ac-repaired-gauge-consequences}.  The
enumerated items are the complementary analytic inputs that must be
available before that theorem can be applied: one must produce the initial concentration chart,
solve the repaired gauge equations, select the resulting scale and center
absolutely continuously,
and verify that the resulting final chart is terminally admissible.
\end{proof}

\subsection{Local Caccioppoli Estimate for Absolutely Continuous Suitable Branches}
\label{subsec:ac-suitable-local-caccioppoli}

\begin{definition}[Suitable repaired-gauge Type II branch]
\label{paper6a:def:suitable-repaired-gauge-typeII-branch}
A suitable repaired-gauge Type II branch is a repaired-gauge Type II branch for
which:
\begin{enumerate}[label=\textup{(\roman*)}]
\item an absolutely continuous concentration chart and an absolutely continuous
repaired-gauge path are available;
\item the conclusions of
Theorem~\ref{paper6a:thm:ac-repaired-gauge-consequences} hold, including
pressure reconstruction and final modulation coefficients;
\item the physical branch \((u,p)\) is a local suitable weak solution on the
terminal interval:
\[
  u\in L^\infty_tL^2_{\rm loc}\cap L^2_tH^1_{\rm loc},
  \qquad
  p\in L^{3/2}_{\rm loc},
  \qquad
  \nabla\cdot u=0,
\]
the Navier--Stokes equations hold distributionally, and the physical local
energy inequality holds for every nonnegative smooth compactly supported test
function.
\end{enumerate}
\end{definition}

\begin{lemma}[Local energy inequality under absolutely continuous Type II charts]
\label{paper6a:lem:ac-pullback-local-energy}
Let \((u,p)\) be physically suitable on a compact physical cylinder, and let
\((V,P)\) be obtained from an absolutely continuous final chart
\[
  u(x,t)=\Lambda(s)^{-1}V(Y,s),
  \qquad
  p(x,t)=\Lambda(s)^{-2}P(Y,s)+c(t),
  \qquad
  Y=\frac{x-X(s)}{\Lambda(s)},
\]
with \(s=\tau\), \(dt=\Lambda(s)^2\,ds\), \(\Lambda>0\), and
\[
  X,\Lambda\in W^{1,1}_{\rm loc}.
\]
Then \((V,P)\) satisfies the renormalized local energy inequality on compact
\((Y,s)\)-cylinders:
\[
\begin{aligned}
&\int_{\mathbb R^3} \frac{|V(s_2)|^2}{2}\phi(s_2)\,dY
+\nu\int_{s_1}^{s_2}\int_{\mathbb R^3} |\nabla V|^2\phi\,dY\,ds \\
&\le
\int_{\mathbb R^3} \frac{|V(s_1)|^2}{2}\phi(s_1)\,dY
+\int_{s_1}^{s_2}\int_{\mathbb R^3}
\frac{|V|^2}{2}(\partial_s\phi+\nu\Delta\phi)\,dY\,ds \\
&\quad
+\int_{s_1}^{s_2}\int_{\mathbb R^3}
\left(\frac{|V|^2}{2}+P\right)V\cdot\nabla\phi\,dY\,ds \\
&\quad
+\int_{s_1}^{s_2}\int_{\mathbb R^3}
 a(s)\frac{|V|^2}{2}(-\phi-Y\cdot\nabla\phi)\,dY\,ds\\
&\quad
-\int_{s_1}^{s_2}\int_{\mathbb R^3}
 \frac{|V|^2}{2}b(s)\cdot\nabla\phi\,dY\,ds,
\end{aligned}
\]
for every nonnegative
\(\phi\in C^\infty([s_1,s_2]\times\mathbb R^3)\) that is compactly supported
in the spatial variable.  Here \(a,b\) are the final modulation coefficients.
\end{lemma}

\begin{proof}
First suppose \(X,\Lambda\) are \(C^1\).  Use the physical test function
\[
  \psi(x,t)=\Lambda(s)^{-1}
  \phi\left(\frac{x-X(s)}{\Lambda(s)},s\right),
  \qquad t=t(s),
\]
where \(dt=\Lambda^2ds\).  The spatial support of \(\phi\) is compact, and
\(\Lambda\) is bounded above and below on the selected compact time interval,
so \(\psi\) is an admissible compactly supported test function in the physical
local energy inequality.

Let
\[
  Y=\frac{x-X(s)}{\Lambda(s)}.
\]
At fixed \(x\),
\[
  \frac{dY}{ds}
  =
  -\frac{X_s}{\Lambda}-\frac{\Lambda_s}{\Lambda}Y .
\]
Consequently,
\[
  \nabla_x\psi=\Lambda^{-2}\nabla_Y\phi,\qquad
  \Delta_x\psi=\Lambda^{-3}\Delta_Y\phi,
\]
and, since \(ds/dt=\Lambda^{-2}\),
\[
  \partial_t\psi
  =
  \Lambda^{-3}
  \left[
  \partial_s\phi
  -\frac{\Lambda_s}{\Lambda}(\phi+Y\cdot\nabla\phi)
  -\frac{X_s}{\Lambda}\cdot\nabla\phi
  \right].
\]
Changing variables
\[
  x=X(s)+\Lambda(s)Y,\qquad dt=\Lambda^2ds,
\]
and using
\[
  u=\Lambda^{-1}V,\qquad p=\Lambda^{-2}P+c(t),
\]
each term in the physical local energy inequality has the following
renormalized form.  The endpoint and dissipation terms transform as
\[
  \int \frac{|u(t_i)|^2}{2}\psi(t_i)\,dx
  =
  \int \frac{|V(s_i)|^2}{2}\phi(s_i)\,dY,
\]
and
\[
  \nu\int_{t_1}^{t_2}\int |\nabla_xu|^2\psi\,dx\,dt
  =
  \nu\int_{s_1}^{s_2}\int |\nabla_YV|^2\phi\,dY\,ds.
\]
The test-function derivative term becomes
\[
\begin{aligned}
&\int_{t_1}^{t_2}\int \frac{|u|^2}{2}
(\partial_t\psi+\nu\Delta_x\psi)\,dx\,dt\\
&\quad =
\int_{s_1}^{s_2}\int \frac{|V|^2}{2}
\left[
\partial_s\phi+\nu\Delta_Y\phi
-a(\phi+Y\cdot\nabla\phi)-b\cdot\nabla\phi
\right]\,dY\,ds,
\end{aligned}
\]
where
\[
  a=\frac{\Lambda_s}{\Lambda},
  \qquad
  b=\frac{X_s}{\Lambda}.
\]
The transport and pressure term becomes
\[
\begin{aligned}
&\int_{t_1}^{t_2}\int
\left(\frac{|u|^2}{2}+p\right)u\cdot\nabla_x\psi\,dx\,dt\\
&\quad =
\int_{s_1}^{s_2}\int_{\mathbb R^3}
\left(\frac{|V|^2}{2}+P\right)V\cdot\nabla_Y\phi\,dY\,ds
{}+\int_{s_1}^{s_2}\Lambda(s)^2c(t(s))
\left(\int_{\mathbb R^3} V\cdot\nabla_Y\phi\,dY\right)\,ds .
\end{aligned}
\]
The pressure constant drops out
because
\[
  \int_{\mathbb R^3} V\cdot\nabla_Y\phi\,dY
  =
  -\int_{\mathbb R^3} \phi\,\nabla_Y\cdot V\,dY=0.
\]
Combining the transformed terms gives the displayed renormalized local energy
inequality.

For \(X,\Lambda\in W^{1,1}_{\rm loc}\), the same identity follows at the level
of test functions.  On the selected compact time interval, \(\Lambda\) has
positive lower and finite upper bounds.  The pulled-back test
\[
  \psi(x,t)=\Lambda(s)^{-1}\phi((x-X(s))/\Lambda(s),s)
\]
is compactly supported, nonnegative, bounded, has bounded spatial derivatives
on the compact cylinder, and has time derivative in \(L^1_tL^\infty_x\)
because \(X',\Lambda'\in L^1\) and \(\phi\) is smooth compactly supported.
The physical local energy inequality, initially stated for \(C_c^\infty\)
tests, extends to this class by standard mollification of the test function:
the suitability classes give the local integrability needed for every term in
the inequality, and the mollified tests keep support inside a slightly larger
compact cylinder.  Applying the extended inequality to \(\psi\) and changing
variables gives the displayed renormalized inequality.  The only derivatives
of the chart that appear are the \(L^1\)-functions \(X',\Lambda'\), hence the
modulation terms are locally integrable against local \(|V|^2\).
\end{proof}

\begin{lemma}[Pressure and modulation for suitable repaired-gauge branches]
\label{paper6a:lem:suitable-gauge-pressure-modulation}
For a suitable repaired-gauge branch, the pressure identity
\[
  -\Delta P=\partial_i\partial_j(V_iV_j)
\]
holds modulo functions of \(\tau\), and the final coefficients \(a,b\) in the
renormalized equation are measurable and locally integrable on compact
renormalized windows.
\end{lemma}

\begin{proof}
This is the pressure and modulation part of
Theorem~\ref{paper6a:thm:ac-repaired-gauge-consequences}.  Pressure pullback
comes from the physical pressure equation and the chart/gauge scaling.  The
final coefficients are forced by the absolutely continuous scale-translation
gauge transform:
\[
  a=\mu^2A+\frac{\mu_\tau}{\mu},
  \qquad
  b=\mu B+\mu Aq+\frac{q_\tau}{\mu}.
\]
Since the chart and gauge parameters are absolutely continuous on compact
windows, these coefficients are locally integrable.
\end{proof}

\begin{theorem}[Local Caccioppoli estimate for repaired gauges]
\label{paper6a:thm:ac-local-caccioppoli}
Every suitable repaired-gauge Type II branch satisfies the local renormalized
Caccioppoli estimate on each compact repaired-gauge cylinder.  The constants
may depend on the cylinder, the cutoff, and local upper and lower bounds for
the chart.  No local Caffarelli--Kohn--Nirenberg smallness, global windowed
\(H^1\), critical \(L^3\)-norm, tightness, bounded modulation, finite energy at
infinity, or whole-space compactness conclusion is asserted.
\end{theorem}

\begin{proof}
By the suitable repaired-gauge hypotheses and
Theorem~\ref{paper6a:thm:ac-repaired-gauge-consequences}, the branch has an
absolutely continuous final chart, an admissible absolutely continuous
repaired-gauge path, pressure reconstruction, and final modulation
coefficients.  Lemma~\ref{paper6a:lem:ac-pullback-local-energy} pulls the physical
local energy inequality to the repaired variables.
Lemma~\ref{paper6a:lem:suitable-gauge-pressure-modulation} supplies the pressure and
coefficient identities needed to interpret the renormalized inequality in the
same PDE class used by the local windowed estimates.

Choose a standard compactly supported Caccioppoli cutoff
\[
  \phi(Y,\tau)=\eta(\tau)\zeta(Y)^2.
\]
Let \(I=(\tau_1,\tau_2)\), let \(I'\Subset I\), and let
\(0<r<R\).  Take \(0\le\eta\le1\), \(\eta\in C_c^\infty(I)\), with
\(\eta\equiv1\) on \(I'\), and take \(0\le\zeta\le1\),
\(\zeta\in C_c^\infty(B_R)\), with \(\zeta\equiv1\) on \(B_r\).  Applying
Lemma~\ref{paper6a:lem:ac-pullback-local-energy} on \((\tau_1,\sigma)\) and
then taking the essential supremum over \(\sigma\in I'\) controls the endpoint
term.  Applying the same inequality on the full interval \(I\) controls the
dissipation term.  Adding the two estimates gives
\[
\begin{aligned}
&\operatorname*{ess\,sup}_{\tau\in I'}
  \int_{B_r}\frac{|V(\tau)|^2}{2}\,dY
  +\nu\int_{I'}\int_{B_r}|\nabla V|^2\,dY\,d\tau \\
&\le C
\int_I\int_{B_R}\frac{|V|^2}{2}
   \left(|\eta'|\zeta^2+\nu\eta|\Delta(\zeta^2)|\right)\,dY\,d\tau \\
&\quad
+C\int_I\int_{B_R}
   \left(\frac{|V|^2}{2}+|P|\right)|V|\,\eta|\nabla(\zeta^2)|\,dY\,d\tau\\
&\quad
+C\int_I |a(\tau)|
  \int_{B_R}\frac{|V|^2}{2}\eta
    \left(\zeta^2+|Y|\,|\nabla(\zeta^2)|\right)\,dY\,d\tau\\
&\quad
+C\int_I |b(\tau)|
  \int_{B_R}\frac{|V|^2}{2}\eta|\nabla(\zeta^2)|\,dY\,d\tau .
\end{aligned}
\]
Here \(C\) is an absolute constant.  Since the cutoff derivatives are bounded
by constants depending only on \(I,I',r,R\), this is the local Caccioppoli
estimate on the compact renormalized cylinder.

No global estimate is used in this step.  On the compact cylinder under
consideration, the absolutely continuous chart has
\[
  0<\Lambda_-\leq \Lambda(\tau)\leq \Lambda_+<\infty.
\]
Pulling back physical suitability therefore gives, for every compact spatial
set \(K\),
\[
  V\in L^\infty_\tau L^2_Y(K)\cap L^2_\tau H^1_Y(K),
  \qquad
  P\in L^{3/2}_{\tau,Y}(K).
\]
Local Sobolev interpolation on the finite cylinder gives
\[
  V\in L^{10/3}_{\tau,Y}(K),
\]
hence \(V\in L^3_{\tau,Y}(K)\).  Thus \(PV\cdot\nabla\phi\) is integrable by
Hölder's inequality.  The representation identities give
\[
  a,b\in L^1_{\rm loc}(d\tau),
\]
and
\[
  \tau\mapsto\int_K |V(Y,\tau)|^2\,dY
\]
is locally essentially bounded; therefore the scale and translation modulation
terms in the pulled-back local energy inequality are integrable on the same
compact cylinder.  This proves the local Caccioppoli estimate, but not a
global windowed \(H^1\) or critical \(L^3\)-norm bound.
\end{proof}

\begin{corollary}[No separate Caccioppoli obstruction for suitable absolutely continuous branches]
\label{paper6a:cor:no-independent-caccioppoli-obstruction-ac-suitable}
On suitable repaired-gauge branches, failure of local Caccioppoli regularity,
failure of pressure reconstruction, and failure of local gauge regularity are
not separate obstructions to the local \(H^1\) analysis.  The additional inputs
used when invoking the uniform windowed \(H^1\) estimates are the bounded
critical \(L^3\) and modulation bounds on the retained window; if either input
fails, the ordered local state-space test records the corresponding active-core,
rough-core, modulation, or scale-drift row before the uniform estimate is
invoked.
\end{corollary}

\begin{proof}
Theorem~\ref{paper6a:thm:ac-local-caccioppoli} gives local Caccioppoli
regularity.  Theorem~\ref{paper6a:thm:ac-repaired-gauge-consequences} gives
pressure reconstruction and the final modulation coefficients from the
final chart and absolutely continuous gauge path.  Thus the listed
failures cannot occur separately on suitable repaired-gauge branches.  The
bounded critical norm and bounded modulation used by the uniform windowed
\(H^1\) estimate are therefore retained-window inputs or routed first failures,
not independent Caccioppoli obstructions.
\end{proof}

\subsection{Caccioppoli Regularity Criterion}

\begin{definition}[Suitable physical hypotheses]\label{paper6a:def:suitable-physical-hypotheses}
Let \(u,p\) be a local suitable weak solution of the three-dimensional
Navier--Stokes equations on \(\mathbb R^3\times(0,T^*)\).  Thus
\[
  u\in L^\infty_tL^2_{\rm loc}\cap L^2_tH^1_{\rm loc},
  \qquad
  p\in L^{3/2}_{\rm loc},
  \qquad \nabla\cdot u=0,
\]
the equations hold distributionally, and for every nonnegative
\(\psi\in C_c^\infty(\mathbb R^3\times(0,T^*))\) and a.e.
\(0<t_1<t_2<T^*\),
\[
\begin{aligned}
&\int_{\mathbb R^3}\frac{|u(t_2)|^2}{2}\psi(t_2)\,dx
+\nu\int_{t_1}^{t_2}\int_{\mathbb R^3}|\nabla u|^2\psi\,dx\,dt \\
&\le
\int_{\mathbb R^3}\frac{|u(t_1)|^2}{2}\psi(t_1)\,dx
+\int_{t_1}^{t_2}\int_{\mathbb R^3}
\frac{|u|^2}{2}(\partial_t\psi+\nu\Delta\psi)\,dx\,dt\\
&\quad
+\int_{t_1}^{t_2}\int_{\mathbb R^3}
\left(\frac{|u|^2}{2}+p\right)u\cdot\nabla\psi\,dx\,dt .
\end{aligned}
\]
\end{definition}

\begin{lemma}[Renormalized local energy inequality]\label{paper6a:lem:renormalized-lei}
Assume the suitable physical hypotheses of
\cref{paper6a:def:suitable-physical-hypotheses}.  Let
\[
  u(x,t)=\lambda(t)^{-1}
  V\left(\frac{x-x_c(t)}{\lambda(t)},\tau(t)\right),
  \qquad
  \frac{d\tau}{dt}=\lambda(t)^{-2},
\]
with
\[
  P(y,\tau)=\lambda(t)^2p(x_c(t)+\lambda(t)y,t),
\]
where \(\lambda>0\), \(x_c\), and \(\tau\) are \(C^1\) on compact windows.
Assume that \(V,P\) satisfy the repaired-gauge equation
\[
  \partial_\tau V+(V\cdot\nabla)V+\nabla P
  =
  \nu\Delta V+a(\tau)(V+y\cdot\nabla V)+b(\tau)\cdot\nabla V,
  \qquad \nabla\cdot V=0 .
\]
Then, for every nonnegative
\(\phi\in C_c^\infty(\mathbb R^3\times(\tau_0,\infty))\) and a.e.
\(\tau_1<\tau_2\) in the corresponding renormalized time interval,
\[
\begin{aligned}
&\int_{\mathbb R^3}\frac{|V(\tau_2)|^2}{2}\phi(\tau_2)\,dy
+\nu\int_{\tau_1}^{\tau_2}\int_{\mathbb R^3}|\nabla V|^2\phi\,dy\,d\tau \\
&\le
\int_{\mathbb R^3}\frac{|V(\tau_1)|^2}{2}\phi(\tau_1)\,dy
+\int_{\tau_1}^{\tau_2}\int_{\mathbb R^3}
\frac{|V|^2}{2}(\partial_\tau\phi+\nu\Delta\phi)\,dy\,d\tau\\
&\quad
+\int_{\tau_1}^{\tau_2}\int_{\mathbb R^3}
\left(\frac{|V|^2}{2}+P\right)V\cdot\nabla\phi\,dy\,d\tau\\
&\quad
+\int_{\tau_1}^{\tau_2}\int_{\mathbb R^3}
a(\tau)\frac{|V|^2}{2}\bigl(-\phi-y\cdot\nabla\phi\bigr)\,dy\,d\tau\\
&\quad
-\int_{\tau_1}^{\tau_2}\int_{\mathbb R^3}
\frac{|V|^2}{2}\,b(\tau)\cdot\nabla\phi\,dy\,d\tau .
\end{aligned}
\]
\end{lemma}

\begin{proof}
For smooth solutions, multiply the repaired-gauge equation by \(V\phi\) and
integrate by parts.  Put \(e=|V|^2/2\).  Since \(\nabla\cdot V=0\),
\[
  V\cdot(V\cdot\nabla V)=V\cdot\nabla e=\nabla\cdot(eV),
  \qquad
  V\cdot\nabla P=\nabla\cdot(PV),
\]
and
\[
  V\cdot(y\cdot\nabla V)=y\cdot\nabla e,
  \qquad
  V\cdot(b\cdot\nabla V)=b\cdot\nabla e .
\]
The scale contribution is
\[
  \int a(2e+y\cdot\nabla e)\phi\,dy
  =
  \int ae(-\phi-y\cdot\nabla\phi)\,dy,
\]
because
\[
  \int y\cdot\nabla e\,\phi\,dy
  =
  -\int e\,\nabla\cdot(y\phi)\,dy
  =
  -\int e(3\phi+y\cdot\nabla\phi)\,dy .
\]
The translation contribution is
\[
  \int b\cdot\nabla e\,\phi\,dy
  =
  -\int e\,b\cdot\nabla\phi\,dy,
\]
and the viscous term gives
\[
  \int \nu\Delta V\cdot V\phi\,dy
  =
  -\nu\int|\nabla V|^2\phi\,dy
  +\nu\int e\Delta\phi\,dy .
\]
Combining these identities gives the displayed equality in the smooth case.

For suitable weak solutions, apply
\cref{paper6a:lem:ac-pullback-local-energy} with
\[
  \Lambda=\lambda,\qquad X=x_c,\qquad s=\tau.
\]
Its physical test function is
\[
  \psi(x,t)=
  \lambda(t)^{-1}
  \phi\left(\frac{x-x_c(t)}{\lambda(t)},\tau(t)\right),
\]
which is exactly the critically scaled pullback needed for the endpoint and
dissipation terms.  Since the present \(C^1\) case is a special case of the
absolutely continuous pullback lemma, the same inequality follows directly.
\end{proof}

\begin{lemma}[Caccioppoli estimate from the renormalized local energy inequality]
\label{paper6a:lem:caccioppoli-from-renormalized-lei}
Assume \cref{paper6a:lem:renormalized-lei} on a compact renormalized cylinder
\(B_{2R}\times I\), integrated modulation on \(I\), and local bounds
\[
  V\in L^\infty(I;L^2(B_{2R}))\cap L^3(I\times B_{2R}),
  \qquad
  P\in L^{3/2}(I\times B_{2R}).
\]
Then the compact-cylinder renormalized Caccioppoli estimate used in
\cref{paper6a:lem:windowed-h1-from-integrated-modulation} is justified.
\end{lemma}

\begin{proof}
Choose a standard nonnegative cutoff
\[
  \phi(y,\tau)=\eta(\tau)\zeta(y)^2,
\]
where \(\zeta\equiv1\) on the inner ball and \(\zeta\) is supported in the
outer ball.  Insert this test function in
\cref{paper6a:lem:renormalized-lei}.  The left-hand side contains
\[
  \nu\int |\nabla V|^2\eta\zeta^2 .
\]
Every term on the right-hand side is finite under the stated hypotheses.  The term
\[
  \int |V|^2|\partial_\tau\phi+\Delta\phi|
\]
is controlled by \(V\in L^\infty(I;L^2(B_{2R}))\) and finite measure.  The
convective term
\[
  \int |V|^3|\nabla\phi|
\]
is controlled by the local spacetime \(L^3\) bound.  For the pressure term, choose
a time-dependent constant \(c(\tau)\).  Since
\[
  \int c(\tau)V\cdot\nabla\phi\,dy
  =
  -\int c(\tau)\phi\,\nabla\cdot V\,dy=0,
\]
Hölder's inequality controls
\[
  \int |P-c(\tau)|\,|V|\,|\nabla\phi|
\]
by \(P-c(\tau)\in L^{3/2}(I\times B_{2R})\) and
\(V\in L^3(I\times B_{2R})\).  The
modulation terms are bounded by
\[
  \int_I(|a(\tau)|+|b(\tau)|)\,d\tau
\]
times the local \(L^\infty_I L^2\)-mass.  These estimates give the
renormalized Caccioppoli estimate on compact cylinders.
\end{proof}

\begin{proposition}[Caccioppoli regularity criterion]\label{paper6a:prop:caccioppoli-regularity-criterion}
Assume the suitable physical hypotheses, the repaired-gauge representation, pressure
reconstruction modulo functions of time, and \(C^1\) scale and center functions
on compact renormalized windows.  Then the compact-cylinder Caccioppoli
regularity hypothesis used in the local windowed \(H^1\) estimate holds.
\end{proposition}

\begin{proof}
The repaired-gauge representation writes the suitable weak solution in the
variables \(V,P,a,b\).  \Cref{paper6a:lem:renormalized-lei} gives the
renormalized local energy inequality, and
\cref{paper6a:lem:caccioppoli-from-renormalized-lei} shows that this inequality
justifies the renormalized Caccioppoli estimate on every compact cylinder.
\end{proof}

\begin{corollary}[Caccioppoli estimates give windowed control under boundedness hypotheses]
\label{paper6a:cor:caccioppoli-windowed-control}
Assume the suitable physical hypotheses, the repaired-gauge representation,
pressure reconstruction, a bounded compact spacetime \(L^3\) amount on compact
cylinders, the local \(L^\infty_\tau L^2_y\) bound from the suitable local
energy package, and bounded integrated modulation on unit windows.  Then local
windowed \(H^1\) control follows from
\cref{paper6a:lem:windowed-h1-from-integrated-modulation}.  Thus failure of the
Caccioppoli regularity hypothesis is not an independent rough-core alternative
under these hypotheses.
\end{corollary}

\begin{proof}
\Cref{paper6a:prop:caccioppoli-regularity-criterion} supplies the Caccioppoli
regularity hypothesis.  The bounded compact spacetime \(L^3\) amount, the
local \(L^\infty_\tau L^2_y\) control, pressure reconstruction, and the
windowwise integrated modulation bound are the remaining local inputs used in
\cref{paper6a:lem:windowed-h1-from-integrated-modulation}.  Therefore the
local windowed \(H^1\) estimate holds.
\end{proof}

\subsection{Terminal Local State-Space Compactness}

\begin{definition}[Admissible retained terminal decomposition and completeness conclusion]
\label{paper6a:def:critical-ns-profile-decomposition-local}
Let \(Q_{2R,I}:=B_{2R}\times I\) be a compact terminal-frame cylinder and let
\[
  V_n=U_n+S_n
\]
be the local decomposition into the retained terminal component \(U_n\) and the
nonretained residual \(S_n\).  The critical-profile interface is expressed
through the following local notions.
\begin{enumerate}[label=\textup{(\roman*)}]
\item The decomposition is an admissible retained terminal decomposition if, on
every \(Q_{2R,I}\), the represented fields satisfy the repaired-gauge
equation, local energy inequality, pressure reconstruction modulo functions of
time, bounded modulation, and the compact-cylinder estimates
\[
  \sup_n\Big(
  \|V_n\|_{L^\infty(I;L^3(B_{2R}))}
  +\|V_n\|_{L^\infty(I;L^2(B_{2R}))}
  +\|V_n\|_{L^2(I;H^1(B_{2R}))}
  +\|P_n-c_{n,2R}\|_{L^{3/2}(Q_{2R,I})}
  \Big)<\infty ,
\]
and the same local velocity bounds hold for \(U_n\), hence also for
\(S_n=V_n-U_n\) after enlarging the constant.  The represented time
derivatives \(\partial_sV_n\) are bounded in
\(L^1(I;H^{-m}(B_{2R}))\) for some \(m\) large enough for the Aubin--Lions
compactness argument.  No time compactness of \(U_n\) or \(S_n\) is assumed
unless a componentwise time-derivative row is explicitly present; otherwise
those components are recorded by weak local velocity coordinates,
selected-time residual measures, and mixed-interaction stress occupation
coordinates.
\item A positive residual local label is assigned as follows.  If, after
passing to a subsequence,
\[
  \limsup_{n\to\infty}\|S_n\|_{L^\infty(I;L^3(B_R))}>0,
\]
then a selected subwindow and subcylinder carry an augmented local occupation
state for the tuple \((V_n,P_n,U_n,S_n,a_n,b_n)\), together with the local
residual concentration measures \(|S_n|^3\,dy\) on selected time slices, whose
residual coordinate or residual concentration measure is nonzero.  The first
closed label of this augmented state is read in the ordered local validation
list of \(\cref{paper7:lem:no-unrecorded-cost-escape}\), with the compact local
Type II row refined by \(\cref{paper6:thm:state-decomposition}\).
\item The local terminal completeness conclusion is the statement that, if all
named adjacent alternatives are closed by their local arguments and comparable
retained classes have been grouped before selecting the terminal frame, then no
nonretained residual with a positive residual local label remains on the
retained terminal branch.  Equivalently,
\[
  S_n\to0
  \quad\text{in }L^\infty(I;L^3(B_R))
\]
on every compact terminal cylinder of the retained branch.
\end{enumerate}
No global \(L^3\) profile decomposition, global heat-flow norm, or global
Kato estimate is part of this statement.
\end{definition}

\begin{remark}[Terminal \(L^\infty(I;L^3)\) coordinates are interface inputs]
\label{paper6a:rem:terminal-l3-interface}
The \(L^\infty(I;L^3)\) bounds in
\cref{paper6a:def:critical-ns-profile-decomposition-local}\textup{(i)} are
part of the local terminal decomposition interface inherited from the
admissible perturbative family of \cref{paper3:def:admissible-terminal}.  They
are not produced by the spacetime compactness package of
\cref{paper5:def:selected-scale-collapse-compactness-package}.  When a
physical branch fails this pointwise critical control on a compact terminal
cylinder, the branch is routed first by
\cref{paper6a:lem:terminal-local-critical-mass-routing} or by an earlier
active-family failure, rather than being treated inside the perturbative
terminal decomposition layer.
\end{remark}

\subsubsection*{Local proof of terminal state-space completeness}
\label{paper6a:subsec:discharge-critical-ns-profile-decomposition}

\begin{lemma}[Interior-time selection on a compact terminal window]
\label{paper6a:lem:interior-time-selection}
Let \(I=[\tau_-,\tau_+]\) be a compact time-interval in the terminal-frame
coordinates, and let \(\tau_n\in I\) be a sequence of selected times for which
\(\|S_n(\tau_n)\|_{L^3(B_R)}\ge\delta>0\).  Then, after passing to a
subsequence, exactly one of the following two admissible selections holds:
\begin{enumerate}[label=\textup{(\roman*)}]
\item there is a compact subinterval \(J\Subset I\) such that
\(\tau_n\in J\) for all large \(n\);
\item the selected times approach an endpoint of \(I\), and the terminal
window is replaced by the admissible recentered windows
\[
  I_n^*:=\tau_n+J_0,\qquad J_0=[-\sigma,\sigma],
\]
with fixed \(\sigma>0\), written in coordinates centered at \(\tau_n\).  In
these recentered coordinates the selected time is \(0\in J_0^\circ\) and
\[
  \|S_n(0)\|_{L^3(B_R)}\ge\delta .
\]
\end{enumerate}
No lower bound is transferred from one time to another.
\end{lemma}

\begin{proof}
After passing to a subsequence, \(\tau_n\to\tau^*\in I\).  If
\(\tau^*\in(\tau_-,\tau_+)\), choose \(J\Subset I\) with
\(\tau^*\in J^\circ\).  Then \(\tau_n\in J\) for all large \(n\), giving
\textup{(i)}.

If \(\tau^*\in\partial I\), we do not move the lower bound to a nearby time.
Instead we use the admissibility of the terminal validation family under finite
covariant time recentering: selected compact windows are local objects, and
composing the terminal-frame chart with the finite translation
\(\tau\mapsto\tau-\tau_n\) produces another member of the same local validation
family on every fixed model interval \(J_0=[-\sigma,\sigma]\).  The
compact-cylinder estimates are then the estimates of
\cref{paper6a:def:critical-ns-profile-decomposition-local}\textup{(i)} on these
recentered selected windows, with constants depending on the model cylinder
and not on \(n\).  In the recentered coordinates the original selected time
\(\tau_n\) is exactly \(0\), so the lower bound remains
\(\|S_n(0)\|_{L^3(B_R)}\ge\delta\).  This proves \textup{(ii)} and uses no
continuity in time of the \(L^3\) norm.
\end{proof}

\begin{lemma}[Local compactness of terminal occupation states]
\label{paper6a:lem:l3-profile-extraction-cited}
Assume the estimates in
\Cref{paper6a:def:critical-ns-profile-decomposition-local}\textup{(i)}
on \(Q_{2R,I}\).  Then every sequence of selected subwindows and
subcylinders has a subsequential local occupation state.

More precisely, after passing to a subsequence, the full represented branch
\(V_n\) converges weakly in \(L^2(I;H^1(B_R))\) and strongly in
\(L^2(I;L^q(B_R))\) for every \(q<6\).  Its pressure gradients converge
distributionally modulo functions of time.

The components \(U_n\) and \(S_n\) have the same strong compactness only on
subbranches for which their componentwise time-derivative rows are available.
In the general terminal decomposition, \(U_n\) and \(S_n\) are retained through
their weak \(L^2H^1\) bounds, selected-time measures
\[
   |U_n(\tau_n)|^3\,dy,\qquad |S_n(\tau_n)|^3\,dy,
\]
and mixed-stress occupation coordinates.
\end{lemma}

\begin{proof}
The Aubin--Lions argument applies directly to \(V_n\), because the represented
Navier--Stokes equation gives the required time-derivative bound.  The local
velocity bounds and the pressure bound give compactness of \(V_n\) on smaller
cylinders and pressure convergence modulo functions of time.

For \(U_n\) and \(S_n\), the same strong spacetime compactness is available
only when their componentwise time-derivative rows are part of the retained
terminal decomposition.  Otherwise the augmented state space records their
weak velocity bounds, their selected-time concentration measures, and their
mixed-stress/pressure diagnostics.  These are precisely the coordinates used
in the local residual exhaustion theorem.
\end{proof}

\begin{lemma}[Closed local labels for terminal residuals]
\label{paper6a:lem:critical-mass-decoupling}
In the augmented local topology for
\[
  (V_n,P_n,U_n,S_n,a_n,b_n)
\]
supplied by \cref{paper6a:lem:l3-profile-extraction-cited}, the ordered
validation rows in \cref{paper7:lem:no-unrecorded-cost-escape}, with the compact
local Type II row refined by \cref{paper6:thm:state-decomposition}, are closed
under subsequential limits of selected compact terminal windows, after shrinking
the tested cylinders inside the fixed compact cylinder when necessary.
\end{lemma}

\begin{proof}
The topology records the represented velocities \(V_n\) weakly in \(L^2H^1\)
and strongly in \(L^2L^q_{\rm loc}\) for every \(q<6\).  The component fields
\(U_n\) and \(S_n\) are recorded weakly in their local velocity topologies;
they are recorded strongly only on subbranches where the corresponding
componentwise time-derivative rows are selected.  In all cases the augmented
topology also records the selected-time residual measures
\(|S_n(\tau_n)|^3\,dy\) weakly-\(*\) as finite Radon measures on compact balls
and the mixed stress coordinates
\(U_n\otimes S_n+S_n\otimes U_n+S_n\otimes S_n\) as occupation coordinates.  It
records the pressures weakly in \(L^{3/2}_{\rm loc}\) modulo functions of time,
and the modulation, scale, center, cutoff, and cost variables in the
compactified repaired-gauge coordinates used in the ordered validation
procedure.  The compactification adds boundary faces exactly for loss of
representation, scale admissibility, bounded modulation, pressure
reconstruction, local \(H^1\) control, critical tightness, componentwise
time-compactness when needed, and cost compatibility.

Positive local velocity-pressure concentration is closed on smaller cylinders.
For spacetime rows, the velocity part follows from strong local \(L^3\)
convergence after interpolating the strong \(L^2L^q\) convergence with the
uniform \(L^\infty L^3\) bound.  For the pressure part, the weak
\(L^{3/2}\) bound alone gives only lower semicontinuity and would allow
positive pressure mass to disappear in the weak limit.  We therefore
upgrade pressure convergence to strong local convergence on the retained branch
as follows.  After subtracting the reconstructed pressure of the limit branch,
write the pressure difference on \(B_{R_1}\), modulo functions of time, as
\[
  \widetilde P_n
  =
  P_n-P-c_n
  =
  \widetilde P_n^{\rm loc}+\widetilde H_n+c_{n,R_1}(\tau),
  \qquad
  \widetilde P_n^{\rm loc}
  =
  \mathcal R_i\mathcal R_j\bigl(\zeta((V_n\otimes V_n)-(V\otimes V))_{ij}\bigr),
\]
with \(\zeta\equiv1\) on \(B_{R_1}\) and
\(\operatorname{supp}\zeta\Subset B_{R_2}\).  Since \(V_n\to V\) strongly in
\(L^2_\tau L^3_y(B_{R_2})\), while \(V_n\) and \(V\) are uniformly bounded in
\(L^\infty_\tau L^3_y(B_{R_2})\), interpolation gives strong convergence in
\(L^3_\tau L^3_y(B_{R_2})\).  Hence
\[
  V_n\otimes V_n\to V\otimes V
  \quad\text{strongly in }L^{3/2}_\tau L^{3/2}_y(B_{R_2}).
\]
Calder\'on--Zygmund continuity therefore gives strong convergence of
\(\widetilde P_n^{\rm loc}\) to zero in
\(L^{3/2}_\tau L^{3/2}_y(B_R)\).  The harmonic remainder
\(\widetilde H_n\) is not controlled by a weak lower-semicontinuity argument:
it is routed
with \((q,p)=(3/2,3/2)\) by the pressure/exterior alternative in
\cref{paper8:lem:harmonic-cross-pressure-routing}.  Thus on the retained
branch the harmonic part vanishes locally modulo functions of time, while off
that branch the nonzero harmonic part is a named pressure or exterior label.
Consequently positive pressure mass cannot be hidden by weak convergence on
the retained branch; pressure rows are strongly closed in the local
\(L^{3/2}_\tau L^{3/2}_y\) topology after shrinking cylinders.  For
selected-time rows, the velocity part is recorded by the weak-\(*\) limits of
the local measures \(|S_n(\tau_n)|^3dy\), so a positive lower bound persists
as mass of the limiting measure.  Hence a lower CKN threshold cannot disappear
after passing to a subsequence and shrinking the cylinder.

Multibubble and multicore splitting are finite collections of such positive
CKN lower bounds with separated centers or separated scales; the separation
relations are closed in the scale-center compactification.  Same-point cascade
labels are closed because they are the boundary faces
\(\lambda_i^n/\lambda_j^n\to0\) or \(\infty\) at a common physical center.
Comparable same-point labels are closed by convergence of the relative
scale-center parameters and are grouped by
\cref{paper8:lem:compound-same-point}.

Exterior leakage and critical-tightness failure are closed because they are
defined by a positive lower bound for local \(L^3\) or CKN mass outside every
retained compact set.  If that mass returns to a fixed compact set along a
subsequence, the state is relabelled by the corresponding compact concentration
row; otherwise it remains on the exterior/tightness boundary face.  Loss of
local windowed \(H^1\) control is the face at infinity for the nonnegative
\(L^2H^1\) seminorm on a compact cylinder.

Finite-cost scale collapse, negative-drift, and compact cost-divergence labels
are determined by nonnegative localized cost functionals.  Their sublevel or
positive-carrier conditions are lower semicontinuous in the recorded
occupation-state topology.  The scale-rigid residual label is closed by the
relative scale-center convergence and by the compact realization of the
scale-rigid branch in \cref{paper3:def:scale-rigid-branch}.  Representation,
modulation, pressure, carrier, and cost-compatibility failures are not hidden
analytic assumptions: they are the explicitly added boundary faces of the
ordered validation compactification.  Thus every first row of the ordered
local validation is a closed diagnostic set.
\end{proof}

\begin{lemma}[Residual measure routing for retained decompositions]
\label{paper6a:lem:local-no-cancellation-retained-decomposition}
Let
\[
  V_n=U_n+S_n
\]
be an admissible retained terminal decomposition on a compact terminal cylinder
\(Q_{2R,I}=B_{2R}\times I\).  Assume the compact-window bounds, pressure
reconstruction, and augmented occupation-state coordinates of
\Cref{paper6a:def:critical-ns-profile-decomposition-local}.  Suppose that for
selected times \(\tau_n\in I\) and a compact ball \(B_\rho\Subset B_R\),
\[
  |S_n(\tau_n)|^3\,dy\big|_{B_\rho}
  \stackrel{*}{\rightharpoonup}\mu_S,
  \qquad
  \mu_S(B_\rho)>0 .
\]
Then, after passing to a subsequence and shrinking \(B_\rho\) if necessary,
one of the following occurs:
\begin{enumerate}[label=\textup{(\roman*)}]
\item the residual measure is supported on a separated-center,
    separated-scale, noncompact, or cascade face; this is a named adjacent
    row;

\item the residual measure is supported on a same-point comparable retained
    class; then the class is grouped into \(U_n\) by
    \Cref{paper6a:lem:retained-frame-maximality};

\item the mixed-stress occupation coordinate
\[
  U_n\otimes S_n+S_n\otimes U_n+S_n\otimes S_n
\]
has positive mass; this is a named pressure, rough-core, compactness,
multibubble/cascade, exterior, or cost/carrier row;

\item the residual measure is genuinely diffuse in the augmented local
  occupation-state coordinates; then
  \Cref{paper6a:thm:local-residual-exhaustion} routes it to a named adjacent
  row or to a retained comparable class, and the retained comparable class is
  grouped into \(U_n\);

\item none of the above occurs, in which case
\[
  S_n\to0
  \quad\text{in the terminal residual topology recorded by the augmented
  state space.}
\]
\end{enumerate}
\end{lemma}

\begin{proof}
The augmented terminal state records the residual measures
\[
  |S_n(\tau_n)|^3\,dy
\]
and the mixed-stress occupation coordinates
\[
  U_n\otimes S_n+S_n\otimes U_n+S_n\otimes S_n .
\]
If the mixed-stress coordinate has positive mass, then the pressure/stress
diagnostics in the augmented state record one of the named pressure,
rough-core, compactness, multibubble/cascade, exterior, or cost/carrier rows.
This is \textup{(iii)}.

Assume therefore that the mixed-stress coordinate is not positive.  Normalize
the positive residual measure on \(B_\rho\) and apply
\Cref{paper6a:lem:residual-selected-time-support-transfer}.  This furnishes a
probability measure \(\Pi_{\rm res}\) on a compact residual local-state space
to which \Cref{paper6a:lem:normalized-defect-atomic-or-diffuse} applies, and
every point of \(\operatorname{supp}\Pi_{\rm res}\) is realized by an actual
sequence of selected-time residual descendants from the original represented
branch.

If the atomic/activity alternative occurs, the support-realization principle
of \Cref{paper6a:lem:residual-selected-time-support-transfer} realizes a
descendant state with positive compact local activity.  The relative
scale-center partition places that descendant in a separated-center,
separated-scale, noncompact, cascade, or same-point comparable class.  The
separated, noncompact, and cascade cases are named adjacent rows, giving
\textup{(i)}.  The same-point comparable case is grouped into \(U_n\) by
\Cref{paper6a:lem:retained-frame-maximality}, giving \textup{(ii)}.

If the normalized residual measure is genuinely diffuse, then the diffuse
probability measure \(\Pi_{\rm res}\) itself is a visible residual coordinate
of the augmented state space.  Applying
\Cref{paper6a:thm:local-residual-exhaustion} to that visible diffuse residual
coordinate yields either a named adjacent local alternative or a retained
comparable class.  The retained comparable class is grouped into \(U_n\) by
\Cref{paper6a:lem:retained-frame-maximality}.  Thus the diffuse case gives
\textup{(iv)}.

If none of \textup{(i)}--\textup{(iv)} occurs, then all residual occupation
coordinates and all mixed-stress coordinates vanish in the augmented topology.
This is exactly perturbativity of \(S_n\) in the terminal residual topology
recorded by the augmented state space, giving \textup{(v)}.
\end{proof}

\begin{lemma}[Residual concentration is visible in the augmented state]
\label{paper6a:lem:residual-concentration-visible}
Let \(V_n=U_n+S_n\) on a compact terminal cylinder, with the local
\(L^\infty L^3\) bounds of
\cref{paper6a:def:critical-ns-profile-decomposition-local}\textup{(i)}.  Suppose
that, for selected times \(\tau_n\) and a compact ball \(B_\rho\),
\[
  |S_n(\tau_n)|^3dy\big|_{B_\rho}
  \stackrel{*}{\rightharpoonup}\mu_S,
  \qquad
  \mu_S(B_\rho)>0 .
\]
Then, after shrinking \(B_\rho\) if necessary, the residual is assigned to a
named adjacent row, grouped into the retained comparable class, recorded as a
visible diffuse residual coordinate routed by
\cref{paper6a:thm:local-residual-exhaustion} to a named adjacent local
alternative or the retained comparable class, or is perturbative in the
terminal residual topology.  In particular,
residual \(L^3\) concentration cannot be invisible in the augmented local state
space by cancellation between \(V_n\) and \(U_n\).
\end{lemma}

\begin{proof}
Apply
\Cref{paper6a:lem:local-no-cancellation-retained-decomposition}.  The lemma
routes the positive residual measure either to a named adjacent row, to a
same-point comparable retained class grouped into \(U_n\), to a positive
mixed-stress row, to a visible diffuse residual coordinate routed by
\cref{paper6a:thm:local-residual-exhaustion} to a named adjacent local
alternative or a retained comparable class, or to perturbativity in the
terminal residual topology.  Thus residual mass is visible in the augmented
local state and cannot be hidden by cancellation between \(V_n\) and \(U_n\).
\end{proof}

\begin{lemma}[Maximality of the selected retained frame]
\label{paper6a:lem:retained-frame-maximality}
The terminal selection rule for the retained component
\(U_n\) on a compact terminal cylinder \(Q_{2R,I}\) produces a finite
collection of comparable retained classes
\(\{U_n^\alpha\}_{\alpha\in A}\) such that:
\begin{enumerate}[label=\textup{(\roman*)}]
\item \emph{Pairwise scale-center separation.} For
\(\alpha\ne\beta\in A\), the relative scale-center parameters
\((\lambda_n^\alpha/\lambda_n^\beta,(x_n^\alpha-x_n^\beta)/\lambda_n^\beta)\)
either tend to \(0\) or to \(\infty\) (separated scales) or stay
bounded with the centers escaping each other in scale-comparable units
(separated centers).
\item \emph{Maximality.} If a candidate retained class
\((\tilde\lambda_n,\tilde x_n)\) at the same physical center and with
scale parameter comparable to some \(\lambda_n^\alpha\) survives the
selection rule (i.e.\ carries a positive active local mass floor in
the sense of \cref{paper8:thm:active-profile-mass-floor}), then it is already
asymptotically equivalent to \((\lambda_n^\alpha,x_n^\alpha)\) and is
absorbed into \(U_n^\alpha\).
\end{enumerate}
\end{lemma}

\begin{proof}
The terminal selection rule is local on the fixed compact cylinder
\(Q_{2R,I}\).  It extracts only classes which carry the active local mass floor
of \cref{paper8:thm:active-profile-mass-floor} on a selected compact
subcylinder.  Classes whose witnesses leave every compact subcylinder, or form
an infinite same-point scale chain, are not retained classes in this branch:
they are read first as exterior, separated-scale, or cascade boundary faces in
the ordered local validation.  After those adjacent faces are removed, the
remaining witnesses lie in a compact scale-center chart.  A finite-overlap
cover of that chart and the local bound
\(\sup_n\|V_n\|_{L^\infty(I;L^3(B_{2R}))}<\infty\) allow only finitely many
disjoint active-floor witnesses.  Hence the retained list on the compact
terminal branch is finite.

For two retained classes in this finite list, either their relative
scale-center parameters leave every compact subset of the relative
scale-center chart, giving separated scales or separated centers, or those
parameters remain in a compact subset.  In the latter case the two
concentrations are comparable in the same local observer coordinates and are
grouped before the terminal frame is fixed.  This proves the separation
alternative in \textup{(i)} for distinct retained classes.

For maximality (ii), suppose toward contradiction that
\((\tilde\lambda_n,\tilde x_n)\) is comparable to
\((\lambda_n^\alpha,x_n^\alpha)\) (i.e.\ both
\(\tilde\lambda_n/\lambda_n^\alpha\to c\in(0,\infty)\) and
\(|\tilde x_n-x_n^\alpha|/\lambda_n^\alpha\to0\)) but not asymptotically
equivalent.  Then the rescaled profile
\(\tilde\lambda_n V_n(x_n^\alpha+\tilde\lambda_n y,\cdot)\) carries an active
local mass floor in a relative scale-center chart compactly comparable to the
chart of \(U_n^\alpha\).  Thus the two candidates are not separated by any
ordered local boundary face.  By \cref{paper8:lem:compound-same-point}, all
same-point comparable concentrations in such a compact relative chart are
assembled into one compound retained class before terminal selection is
closed.  The candidate \((\tilde\lambda_n,\tilde x_n)\) must therefore already
be absorbed into \(U_n^\alpha\), contradicting the assumption that it is a
surviving unabsorbed comparable class.
\end{proof}

\begin{lemma}[Terminal local critical-mass failures are adjacent labels]
\label{paper6a:lem:terminal-local-critical-mass-routing}
Let a physical local Type II branch have passed covered entry and let a
terminal active frame be selected from its local state-space decomposition.  On
each fixed compact terminal cylinder, exactly one of the following occurs after
passing to a subsequence:
\begin{enumerate}[label=\textup{(\roman*)}]
\item the retained terminal component has positive finite local critical
  spacetime \(L^3\) mass on the cylinder;
\item the retained spacetime activity floor fails, so the candidate is not a retained
  terminal component;
\item the local critical-mass upper control fails on a compact spacetime
  subcylinder, in which case the ordered local state space records an
  additional compact active label: a critical-mass adjacent label at the same
  retained frame, a comparable finite-mass retained class which must be
  grouped into the retained component, a multibubble/multicore label, a
  cascade label, a rough-core label, or an exterior/separated label.
\end{enumerate}
In the retained terminal branch only \textup{(i)} remains.
\end{lemma}

\begin{proof}
The lower bound in \textup{(i)} is the active local spacetime mass floor of
\cref{paper8:thm:active-profile-mass-floor}.  If this floor is absent, the
candidate frame is not retained by the terminal selection rule, giving
\textup{(ii)}.
\par
Assume the lower floor is present and the finite local upper bound fails on a
compact spacetime subcylinder \(B_R\times J\Subset B_{2R}\times I\).  Then
\[
  \int_J\int_{B_R}|V_n(y,\tau)|^3\,dy\,d\tau\to\infty .
\]
Cover \(B_R\times J\) by finitely many smaller parabolic cylinders.  On one
member of this finite cover, after passing to a subsequence and renormalizing
at the corresponding local scale if necessary, the branch carries a further
positive compact spacetime \(L^3\) diagnostic.  The ordered local state-space
compactification records the first label of that diagnostic.  If the
diagnostic stays at the same comparable retained frame and the compact upper
control still fails, the first label is the critical-mass adjacent face
itself; by the retained-state definition this is not a retained terminal
branch.  If the diagnostic is comparable and finite-mass, then
\cref{paper6a:lem:retained-frame-maximality} says it was already grouped into
the retained component.  If it is not comparable, it is by definition one of
the adjacent compact active labels: multibubble/multicore splitting, cascade,
rough core, or exterior/separated concentration.  These are precisely the
labels in \textup{(iii)}.
\par
On a retained terminal branch all comparable labels have already been grouped,
and all non-comparable labels are adjacent local alternatives routed by the
state-space closure package.  Hence only the positive finite local critical
mass alternative remains.  This is a compact-cylinder statement; it does not
require \(V(\tau)\in L^3(\mathbb R^3)\) or a whole-space critical norm bound.
\end{proof}

\begin{theorem}[Adjacent terminal labels are locally routed]
\label{paper6a:thm:terminal-adjacent-label-routing}
For every admissible retained terminal decomposition, every first label in the
ordered augmented local state space is either the retained comparable class or
one of the named local alternatives.  The adjacent labels are closed and routed
out of the retained terminal component as follows:
representation, modulation, pressure, critical-mass, carrier,
cost-compatibility, exterior tightness, local \(H^1\)-loss, multibubble,
multicore, cascade, finite-cost scale-collapse, and scale-rigid residual labels
are assigned to their named local alternatives and do not remain on the
retained terminal branch.  This theorem is a routing theorem only: it does not
invoke the final state-elimination proposition and does not claim, inside the
terminal-completeness argument, that the adjacent alternatives have already
been globally excluded.
\end{theorem}

\begin{proof}
\Cref{paper6a:lem:critical-mass-decoupling} proves closedness of the rows in
the augmented local topology.  Representation and modulation failures are the
boundary faces of the repaired-gauge compactification and are exactly the
representation alternatives.  Pressure rows are routed by
\cref{paper8:lem:harmonic-cross-pressure-routing} and by the pressure
reconstruction stability \cref{paper1:lem:pressure-stability}.  Critical-mass
rows are routed by
\cref{paper6a:lem:terminal-local-critical-mass-routing}: failure of the
retained spacetime activity floor
is not retained, and an upper local \(L^3\)-failure produces a comparable or
adjacent active label.  Exterior tightness and exterior leakage are routed by
\cref{paper8:thm:exterior-stratification}.  Local \(H^1\)-loss is exactly the
rough-core row supplied by \cref{paper6:thm:rough-core}.  Multibubble,
multicore, and cascade are primitive structural labels in the ordered active
selection record: they arise from failure of unique retained-core selection,
strict nested same-point descendants, or separated comparable descendants
before any terminal retained class is fixed.  Gauge-degenerate labels are
already part of the repaired-gauge representation failure record.  Finite-cost
scale-collapse is the finite-cost scale-collapse row of
\cref{paper6:thm:scale-stratification}, and scale-rigid residual labels are
recorded as the scale-rigid row governed by
\cref{paper3:prop:scale-rigidity-discharged}.  Carrier and compact-cost
compatibility rows are the carrier/cost adjacent rows of the ordered
compact-cost validation.
\par
The point needed here is only separation from the retained terminal component.
Each listed row is a boundary face of the ordered local state space, not the
retained comparable class.  Its later exclusion is supplied in the assembly
layer by the local closure theorems and, for carrier/cost rows, by the
ambient adjacent wrapper.  The terminal-completeness proof therefore uses this
theorem only to say that a positive residual label cannot stay in the
nonretained remainder of a retained terminal branch.  The only first label not
routed away is the retained comparable class, which is grouped into the
selected terminal component before the terminal frame is fixed.
\end{proof}

\begin{theorem}[Residual local mass gives an admissible augmented state]
\label{paper6a:thm:residual-mass-admissible-state}
Assume the local compact-window estimates of
\cref{paper6a:def:critical-ns-profile-decomposition-local}\textup{(i)}.  If
\[
  \limsup_{n\to\infty}\|S_n\|_{L^\infty(I;L^3(B_R))}>0,
\]
then, after passing to a subsequence, there are selected times, still denoted
\(\tau_n\) in the selected coordinates, a compact selected or recentered
subcylinder \(Q_{\rho,J}\), and an augmented local occupation state
\[
  \mathfrak s
  =
  \lim_{n\to\infty}
  (V_n,P_n,U_n,S_n,a_n,b_n)\big|_{Q_{\rho,J}}
\]
whose residual coordinate or selected-time residual \(L^3\) concentration
measure is nonzero.  The Navier--Stokes admissibility conditions in the ordered
local validation are imposed on the full represented branch
\((V_n,P_n,a_n,b_n)\), not on \(S_n\) as an independent solution.  Consequently
the first nonretained residual row is either a named adjacent local alternative
or the retained comparable local class.
\end{theorem}

\begin{proof}
Choose \(\delta>0\), times \(\tau_n\in I\), and a subsequence such that
\[
  \|S_n(\tau_n)\|_{L^3(B_R)}\ge\delta .
\]
By \cref{paper6a:lem:interior-time-selection}, after passing to a further
subsequence, either the original selected times lie in a fixed interior
subinterval \(J\Subset I\), or we replace the terminal window by the
admissible recentered windows supplied by that lemma and write the selected
time as \(0\in J_0^\circ\).  In both cases we keep the same selected time
rather than transferring the lower bound to a different time; below \(J\)
denotes the fixed model interval in the selected coordinates and \(\tau_n\)
denotes the selected interior time in those coordinates.
Cover \(B_R\) by finitely many smaller balls compactly contained in \(B_R\).
One ball, after passing to a subsequence, carries residual \(L^3\) mass bounded
below by a fixed multiple of \(\delta\).  The compact subinterval
\(J\) supplied by \cref{paper6a:lem:interior-time-selection}
contains the selected times with positive distance from the boundary of the
model window; we use this \(J\) without further time movement.  The estimates in
\cref{paper6a:def:critical-ns-profile-decomposition-local}\textup{(i)} and
\cref{paper6a:lem:l3-profile-extraction-cited} give a subsequential limit of
the augmented tuple on \(Q_{\rho,J}\).  In addition, the measures
\[
  |S_n(\tau_n)|^3\,dy\big|_{B_\rho}
\]
have uniformly bounded total mass on compact balls and therefore converge
weakly-\(*\), after passing to a subsequence, to a finite nonzero Radon measure.
The lower mass bound on the chosen ball passes to this measure.  Thus the
augmented state has a nonzero residual coordinate or, in the concentration
case, a nonzero residual \(L^3\) concentration measure.
\Cref{paper6a:lem:residual-concentration-visible} shows that this residual
diagnostic is visible either as a named adjacent local label of the full
represented branch, as a positive mixed-stress row, as a visible diffuse
residual coordinate in the augmented state space, or as a positive comparable
retained label.  Hence it cannot be hidden by cancellation with \(U_n\).

The PDE structure used by the state-space decomposition is carried by
\((V_n,P_n,a_n,b_n)\): it satisfies the repaired-gauge Navier--Stokes equation,
the local energy inequality, pressure reconstruction modulo functions of time,
bounded modulation, and the compact-cylinder bounds.  The retained component
\(U_n\) and residual coordinate \(S_n\) are bookkeeping variables in the
augmented state space.  Thus no step applies the Navier--Stokes decomposition
theorem to \(S_n\) alone.

Evaluate the ordered local validation list of
\cref{paper7:lem:no-unrecorded-cost-escape} on the augmented state.  If a
representation, pressure, modulation, exterior-tightness, local \(H^1\),
finite-cost, scale-rigid, or carrier/cost row fails first, that row is the
named adjacent alternative.  If the compact local Type II row is reached, apply
\cref{paper6a:thm:terminal-adjacent-label-routing} to the first nonretained row
of the full represented branch \((V_n,P_n)\) on the selected compact window.
If that first row is not the retained comparable class, then it is a named
adjacent local alternative.  If the visible residual diagnostic is diffuse
rather than atomic, we apply the same first-row extraction to that diffuse
residual coordinate in the augmented state space; closedness of the ordered
rows again yields a named adjacent local alternative or the retained
comparable class.  The remaining possibility is therefore that the residual
diagnostic is carried by the retained compact class.  That class is comparable
to the selected terminal frame; otherwise the separated-center or
separated-scale decoupling lemmas would put it into one of the preceding
adjacent rows.  Hence the first nonretained residual row is either a named
adjacent local alternative or the retained comparable class.
\end{proof}

\begin{lemma}[Positive residual mass produces a local state-space label]
\label{paper6a:lem:heat-flow-kato-smallness}
Assume the local compact-window estimates of
\cref{paper6a:def:critical-ns-profile-decomposition-local}\textup{(i)}.  If
\[
  \limsup_{n\to\infty}\|S_n\|_{L^\infty(I;L^3(B_R))}>0,
\]
then, after passing to a subsequence and selecting times and subcylinders, the
residual produces one of the first closed labels listed in
\cref{paper6a:def:critical-ns-profile-decomposition-local}\textup{(ii)}.
\end{lemma}

\begin{proof}
\Cref{paper6a:thm:residual-mass-admissible-state} constructs the augmented
occupation state and verifies that the Navier--Stokes admissibility belongs to
the full represented branch \((V_n,P_n)\).  The ordered first row of that state
is closed by \cref{paper6a:lem:critical-mass-decoupling}.  Therefore, after
passing to a subsequence, the positive residual mass has a first closed local
label in the sense of
\cref{paper6a:def:critical-ns-profile-decomposition-local}\textup{(ii)}.
\end{proof}

\begin{lemma}[Residual mass cannot remain in the retained terminal branch]
\label{paper6a:lem:nonlinear-profile-evolution-cited}
Assume the local terminal frame has been chosen after grouping comparable
retained classes.  Then a nonretained residual satisfying the hypotheses of
\cref{paper6a:lem:heat-flow-kato-smallness} cannot remain on the retained
terminal branch.
\end{lemma}

\begin{proof}
If the residual has positive local mass, then
\cref{paper6a:lem:heat-flow-kato-smallness} assigns it to a first closed local
label.  By \cref{paper6a:thm:terminal-adjacent-label-routing}, an adjacent
label is a named local alternative and is not part of the retained branch.  If
the label is the retained compact class, then the corresponding component is
comparable to the selected retained terminal frame; by the terminal selection
rule together with \cref{paper6a:lem:retained-frame-maximality},
comparable retained classes based at the same physical point were already
grouped into \(U_n\).  It therefore cannot remain in \(S_n\).  Both cases
contradict the
assumption that positive residual mass persists in the nonretained component of
the retained terminal branch.
\end{proof}

\begin{lemma}[Subsequential exhaustion of local residual mass]
\label{paper6a:lem:profile-exhaustion-by-diagonal-extraction}
Under the hypotheses of
\cref{paper6a:lem:nonlinear-profile-evolution-cited}, every compact terminal
cylinder satisfies
\[
  S_n\to0
  \quad\text{in }L^\infty(I;L^3(B_R)).
\]
\end{lemma}

\begin{proof}
If the convergence failed, then a subsequence would have positive residual
mass in the sense of \cref{paper6a:lem:heat-flow-kato-smallness}.  The preceding
lemma shows that such mass cannot remain in the retained terminal branch after
local state-space closure and grouping of comparable retained classes.  This
contradiction proves the convergence.
\end{proof}

\begin{theorem}[Visible local residual exhaustion theorem]
\label{paper6a:thm:local-residual-exhaustion}
Assume the local compact-window estimates, local pressure atlas, repaired-gauge
chart on the retained branch, and local energy inequality on a fixed compact
terminal cylinder.  Assume moreover that the positive residual diagnostic is
visible in the augmented local occupation-state coordinates used in
\cref{paper6a:lem:critical-mass-decoupling}; this includes, in particular,
positive local \(L^3\) residual mass and any pressure/cost/compactness defect
already encoded there as a closed diagnostic row.  Then after passing to a
subsequence and, if necessary, shrinking to a selected compact subcylinder, one
of the explicitly named local rows occurs.  More precisely, the first closed
row is either a named adjacent local alternative, or the retained comparable
class, in which case the corresponding component must already have been grouped
into the retained terminal class before the frame was fixed.
\end{theorem}

\begin{proof}
For positive residual \(L^3\) mass,
\cref{paper6a:thm:residual-mass-admissible-state} produces an augmented local
occupation state for the tuple
\[
  (V_n,P_n,U_n,S_n,a_n,b_n)
\]
on a selected compact subcylinder.  The no-cancellation step is supplied before
the first-row extraction by
\cref{paper6a:lem:local-no-cancellation-retained-decomposition}: a positive
selected-time residual measure produces a named adjacent row, a same-point
comparable retained class, a positive mixed-stress row, or a visible diffuse
residual coordinate; if none of these occurs, the residual is perturbative in
the terminal residual topology and no positive residual measure remains.
For any other residual diagnostic covered by the statement, visibility in the
same augmented coordinates is part of the hypothesis.  By
\cref{paper6a:lem:critical-mass-decoupling}, each row of the ordered local
validation is a closed diagnostic set in that augmented topology.  Therefore a
positive residual defect has a first closed row after subsequence extraction.

If the first row is representation, pressure, modulation, rough-core, exterior,
multibubble, cascade, carrier, or cost incompatibility, then
\cref{paper6a:thm:terminal-adjacent-label-routing} identifies it as a named
adjacent local alternative.  If the first row is positive residual \(L^3\)
mass, then \cref{paper6a:lem:residual-concentration-visible} routes it through
the no-cancellation alternatives: a named adjacent row, a positive
mixed-stress row, a visible diffuse residual coordinate, or a retained
comparable class.  By
\cref{paper6a:lem:critical-mass-decoupling}, the diffuse residual coordinate
also has a first closed row in the ordered validation.  By
\cref{paper6a:lem:retained-frame-maximality}, every such retained comparable
class was already grouped into the selected retained component before the
terminal frame was fixed.

Consequently no positive nonretained residual defect can remain unlabelled on
the retained terminal branch: it is exhausted by a named local row or absorbed
into the retained comparable class.  This is precisely the required local
exhaustion statement.
\end{proof}

\begin{theorem}[Local terminal completeness from state-space stratification]
\label{paper6a:thm:critical-ns-profile-decomposition-discharged}
Every admissible retained terminal decomposition in the sense of
\cref{paper6a:def:critical-ns-profile-decomposition-local}\textup{(i)} satisfies the
local terminal completeness conclusion of
\cref{paper6a:def:critical-ns-profile-decomposition-local}\textup{(iii)}, after the
terminal frame is selected by grouping comparable retained classes.
\end{theorem}

\begin{proof}
\Cref{paper6a:lem:l3-profile-extraction-cited} gives the local compactness
needed to pass to occupation states.  The actual exhaustion step for positive
local residual \(L^3\) mass, and more generally for every visible residual
diagnostic recorded in the augmented state space, is
\cref{paper6a:thm:local-residual-exhaustion}: such a defect is assigned, after
subsequence extraction, to a named adjacent row or to the retained comparable
class.  The latter possibility is removed by
\cref{paper6a:lem:retained-frame-maximality}, and the former is incompatible
with persistence of nonretained residual mass on the retained branch by
\cref{paper6a:lem:nonlinear-profile-evolution-cited}.  Hence
\cref{paper6a:lem:profile-exhaustion-by-diagonal-extraction} gives \(S_n\to0\)
on compact terminal cylinders.  This is the local terminal completeness
conclusion named in
\cref{paper6a:def:critical-ns-profile-decomposition-local}\textup{(iii)}.
\end{proof}

\begin{remark}[Status of the terminal completeness statement]
\label{paper6a:rem:profile-decomposition-status}
\Cref{paper6a:thm:critical-ns-profile-decomposition-discharged} is local.  It
uses only compact-cylinder estimates, Aubin--Lions compactness, pressure
reconstruction modulo functions of time, and the local state-space
stratification.  It does not use a global critical profile decomposition, a
global heat-flow smallness norm, or a global Kato perturbation theorem.
\end{remark}

\begin{definition}[Local compact terminal windows]
\label{paper6a:def:bounded-critical-terminal-sequences}
The local compact terminal-window hypothesis is the assertion that every
terminal active sequence used in the terminal analysis is represented, on each
compact terminal-frame cylinder \(Q_{2R,I}\), by fields satisfying the local
compact-cylinder bounds in
\cref{paper6a:def:critical-ns-profile-decomposition-local}\textup{(i)}.  Equivalently,
the analysis is carried out on the local state variables
\[
  (V_n,U_n,S_n,P_n,a_n,b_n)\big|_{Q_{2R,I}},
\]
not on an arbitrary global sequence of initial data.
Here the time interval \(I\) is always taken from the fixed local validation
family of model intervals, up to the admissible finite covariant recenterings
of \cref{paper6a:lem:interior-time-selection}.  In particular the terminal
compactness package is local on fixed-size model windows rather than on
arbitrary long global intervals.
\end{definition}

\begin{theorem}[Local perturbative residual stability]
\label{paper6a:thm:kato-small-data-critical-stability}
Let \(Q_{2R,I}\) be a compact terminal cylinder.  Assume
\[
  S_n\to0\quad\text{in }L^\infty(I;L^3(B_{2R})),
\]
and assume \(U_n\) and \(S_n\) are uniformly bounded in
\[
  L^2(I;H^1(B_{2R}))\cap L^\infty(I;L^3(B_{2R})).
\]
Then all velocity interactions involving at least one factor \(S_n\) vanish on
the smaller cylinder:
\[
  U_n\otimes S_n+S_n\otimes U_n+S_n\otimes S_n
  \to0
  \quad\text{in }L^1(I;L^2(B_R)).
\]
If the corresponding cross pressures are reconstructed by the local
Calder\'on--Zygmund decomposition, then their gradients vanish in
\[
  L^1(I;H^{-1}(B_R)).
\]
\end{theorem}

\begin{proof}
By Sobolev on \(B_{2R}\), the \(L^2H^1\) bounds give uniform
\(L^2(I;L^6(B_{2R}))\) bounds.  Hölder's inequality gives
\[
  \|U_n\otimes S_n\|_{L^1(I;L^2(B_R))}
  \le
  \|S_n\|_{L^\infty(I;L^3(B_R))}
  \|U_n\|_{L^1(I;L^6(B_R))}\to0,
\]
and the same estimate treats \(S_n\otimes U_n\).  The term
\(S_n\otimes S_n\) is identical, using the uniform \(L^1L^6\) bound for
\(S_n\).  The pressure statement follows from the local pressure-kernel
estimate applied to the cross stress just shown to vanish in \(L^1L^2\), with
the pressure normalized modulo functions of time.  Concretely, the
local Calder\'on--Zygmund piece of the cross pressure vanishes in
\(L^1_\tau L^2_y(B_R)\) by Calder\'on--Zygmund continuity applied to the
vanishing cross stress, and the harmonic remainder is routed by
\cref{paper8:lem:harmonic-cross-pressure-routing} with \((q,p)=(1,2)\) (using the
quantitative nested-ball harmonic-tail estimate of
\cref{paper8:lem:nested-ball-harmonic-tail}): on the retained terminal
branch any non-vanishing harmonic remainder triggers a named
pressure/exterior-leakage alternative, which is excluded.  The two
contributions sum to zero in \(L^1_\tau H^{-1}_y(B_R)\).
\end{proof}

\begin{definition}[Local perturbative stability]
\label{paper6a:def:small-data-critical-stability}
Local perturbative stability means that every residual component which is
small in \(L^\infty(I;L^3(B_R))\) on compact terminal cylinders contributes
only vanishing velocity-stress, pressure, and linear-residual terms in the
local topologies used in
\cref{paper8:def:terminal-nonlinear-decoupling}.
\end{definition}

\begin{corollary}[Local residual stability gives perturbative stability]
\label{paper6a:cor:kato-gives-small-data-stability}
\Cref{paper6a:thm:kato-small-data-critical-stability} implies the
local perturbative stability condition of
\cref{paper6a:def:small-data-critical-stability}.
\end{corollary}

\begin{proof}
This is exactly the conclusion of
\cref{paper6a:thm:kato-small-data-critical-stability} on each compact terminal
cylinder, combined with the local pressure and residual estimates already
included in terminal nonlinear decoupling.
\end{proof}

\subsubsection*{Discharge of the retained compact-test closure}
\label{paper6a:subsec:discharge-retained-compact-tests}

\begin{lemma}[Logarithmic shell support points are realized after shell rebasing]
\label{paper9:lem:logarithmic-tail-shell-realization}
Fix compact model intervals \(J_{\rm sh}\Subset J_*\Subset\mathbb R\) with
nonempty interior and set
\[
  Q_{\rm sh}:=B_1\times J_{\rm sh}.
\]
Let \(J_n\) be selected terminal windows modeled on \(J_*\).  For each occupied
spacetime shell block
\[
  A_{n,k}\times J_n
  =
  \{2^kR_n<|y|\le 2^{k+1}R_n\}\times J_n
\]
cover the block by finitely overlapping translated shell cylinders
\[
  Q_{\rm sh}(y_{n,k,\alpha},\tau_{n,k,\alpha})
  :=
  B_1(y_{n,k,\alpha})\times (\tau_{n,k,\alpha}+J_{\rm sh})
  \subset \mathbb R^3\times J_n,
\]
with \(y_{n,k,\alpha}\in A_{n,k}\).  Record the logarithmic shell radius
\[
  \rho_{n,k}:=\log(2^kR_n),
\]
the compactified shell-location data of the chosen cylinder inside the shell
block, the rebased local velocity/pressure state on the fixed model cylinder
\(Q_{\rm sh}\) obtained by spatial translation and admissible finite time
recentering only, and the boundary faces where representation, pressure-atlas,
gauge, local-energy, or rough-core admissibility fails.  The
local-state factor is compact by
\cref{paper6a:lem:compact-local-state-factor}, so the closure of this family is
a compact metrizable shell-state space
\[
  \mathfrak X_{\rm shell}
  :=
  (\mathbb N\cup\{\infty\})
  \times \overline{\mathcal C}_{\rm shell}
  \times \mathfrak O_{\rm shell}
  \times \mathfrak F_{\rm shell}.
\]
The shell occupation coordinates in \(\mathfrak O_{\rm shell}\) are continuous
coordinate projections of the product compactification.  Consequently every
positive superlevel set of a shell occupation coordinate is closed in
\(\mathfrak X_{\rm shell}\).
After passing to a subsequence, the normalized spacetime shell measures induce a weak-*
probability limit on \(\mathfrak X_{\rm shell}\).  The limiting shell measure
itself is retained as part of the shell compactification, so later diffuse
arguments may apply
\cref{paper6a:lem:normalized-defect-atomic-or-diffuse} directly to it.  Every support point is
realized by an actual sequence of occupied shell blocks and corresponding
translated shell-cylinder extractions from the original represented branch.

If the support point avoids the failure faces \(\mathfrak F_{\rm shell}\), then
the realizing shell sequence satisfies the represented local equation, pressure
atlas modulo functions of time, repaired-gauge structure, local energy
inequality, and compact-cylinder bounds on the fixed model cylinder
\(Q_{\rm sh}\);
hence the ordered local state-space validation applies there.  If the support
point lies on a face of \(\mathfrak F_{\rm shell}\), that face itself is the
corresponding named local row.
\end{lemma}

\begin{proof}
For each occupied shell block, choose a finitely overlapping family of
translated shell cylinders
\(Q_{\rm sh}(y_{n,k,\alpha},\tau_{n,k,\alpha})\subset \mathbb R^3\times J_n\).
Translate each chosen cylinder by \(-y_{n,k,\alpha}\) and use the admissible
finite covariant recentering of \cref{paper6a:lem:interior-time-selection} to
identify its time window with the fixed model interval \(J_{\rm sh}\).  No
dilation by the shell radius is performed in the local Navier--Stokes state;
the shell radius enters only through the logarithmic coordinate \(\rho_{n,k}\)
and the compactified shell-location data.  Spatial translation, finite time
recentering, and subtraction of the appropriate pressure mean preserve the
represented local equation, divergence-free condition, local energy inequality,
and local pressure reconstruction on the fixed model cylinder \(Q_{\rm sh}\).
Thus every rebased shell cylinder either remains in a uniform local compactness
class on \(Q_{\rm sh}\) or hits one of the explicit failure faces recorded in
\(\mathfrak F_{\rm shell}\).  The logarithmic shell index belongs to the
one-point compactification of \(\mathbb N\), the shell-location chart is
compactified exactly as in the local state-space package, and the rebased
local-state factor is compact by
\cref{paper6a:lem:compact-local-state-factor}.  This gives the compact metrizable space
\(\mathfrak X_{\rm shell}\).

Push forward the normalized spacetime shell measures to \(\mathfrak X_{\rm shell}\) and
extract a weak-* limit.  If \(\mathfrak z\) lies in the support of that limit,
choose a nested neighborhood basis \((U_m)\) of \(\mathfrak z\) with diameters
tending to zero.  Positive limit mass of \(U_m\) implies positive mass of
\(U_m\) for the approximating shell measures along a subsequence, so for each
\(m\) we may select an actual occupied shell block and a rebased
model-cylinder state lying in \(U_m\).  These genuine rebased shell states converge to
\(\mathfrak z\), proving realization.  The last paragraph is built into the
definition of the failure faces and the represented part of the shell-state
space.
\end{proof}

\begin{theorem}[Logarithmic tail support-profile transfer]
\label{paper9:thm:logarithmic-tail-support-transfer}
Let
\[
  M_n
  :=
  \int_{J_n}\int_{|y|>R_n}|V(y,\tau)|^3\,dy\,d\tau,
\]
and
\[
  \nu_n
  :=
  M_n^{-1}|V(y,\tau)|^3\mathbf 1_{\{\tau\in J_n,\ |y|>R_n\}}\,dy\,d\tau
\]
be the normalized spacetime tail measures associated to a failure sequence of
windowed critical tightness, where the selected terminal windows \(J_n\) are
all modeled on one fixed compact interval \(J_*\Subset\mathbb R\),
\(M_n\ge\varepsilon_0>0\), and \(R_n\to\infty\).  After passing to a
subsequence, the logarithmic shell compactification of the spacetime tail
carries a probability limit state.  Every point in the support of that limit is
realized by a translated shell-cylinder extraction of the original represented
branch, with no shell-radius Navier--Stokes dilation, and each realized
support profile has a first closed local row in the ordered local state-space
validation.
\end{theorem}

\begin{proof}
Apply \cref{paper9:lem:logarithmic-tail-shell-realization}.  This gives the
compact metrizable shell-state space, a probability limit state, and an actual
translated shell-cylinder realizing sequence for every support point.

Let \(\Pi_{\rm shell}\) denote the weak-* limit shell state furnished by that
lemma.

Let \(\mathfrak z\in\operatorname{supp}\Pi_{\rm shell}\) and let
\((W_m,\widetilde P_m)\) be the corresponding realizing rebased shell sequence on a fixed
model cylinder \(Q_{\rm sh}\), obtained only by spatial translation and finite
time recentering.  If \(\mathfrak z\) lies on one of the failure faces
of \(\mathfrak F_{\rm shell}\), then that face is already the first closed local
row.  Otherwise \((W_m,\widetilde P_m)\) is a genuine represented local sequence with the
pressure atlas, repaired gauge, local energy inequality, and compact-cylinder
bounds on the fixed model cylinder.  The ordered local rows are closed in this
topology by
\cref{paper6a:lem:critical-mass-decoupling,paper6a:thm:terminal-adjacent-label-routing},
so the realized support profile has a first closed local row.  This proves the
support-transfer statement.
\end{proof}

From this point onward, ``rebased unit-shell'' means a local state on the
fixed model cylinder \(Q_{\rm sh}\) obtained from an occupied shell block only
by spatial translation and admissible finite time recentering.  The shell
radius is recorded as a logarithmic occupation-state coordinate and is not used
as a Navier--Stokes scaling parameter.

\begin{theorem}[Minimal mesoscopic extraction for diffuse critical tails]
\label{paper9:thm:minimal-mesoscopic-extraction}
Assume the support profiles furnished by
\cref{paper9:thm:logarithmic-tail-support-transfer} carry no positive compact
retained local velocity activity at any predeclared dyadic ladder level
\(\eta_m\).  Then, after passing to a
subsequence, there exists either a realized support profile or a genuinely
diffuse shell measure whose first closed local row is minimal among the
atomic/activity and diffuse critical-tail alternatives.  This first row is one
of the named local rows below the compact retained single-core branch.
\end{theorem}

\begin{proof}
The ordered local validation has finitely many rows.  Apply
\cref{paper6a:lem:normalized-defect-atomic-or-diffuse} to the shell measures
furnished by \cref{paper9:thm:logarithmic-tail-support-transfer}.

If the atomic/activity alternative occurs, choose among the realized support
profiles with positive compact local activity one whose first closed local row
is minimal in that finite ordered list.  By the assumption that no support
profile carries positive compact retained local velocity activity, the carrier
rows of \cref{paper9:lem:critical-tightness-failure-row} do not occur for this
minimal profile.  Therefore its first row is one of the remaining named local
rows: pressure/gauge/reconstruction failure, rough-core or local-energy
failure, affine/pure-gauge behavior, hidden-scale/scale-collapse behavior, or
the diffuse coherent critical-tail row.

If the genuinely diffuse alternative occurs, then the nonzero shell measure
itself is the required diffuse critical-tail state, and its first closed row is
the diffuse critical-tail row.  This exhausts the minimal mesoscopic
alternatives.
\end{proof}

From this point onward, fix once and for all a dyadic activity ladder
\[
\eta_m:=2^{-m}\eta_{\rm ref},\qquad m\in\mathbb N,
\]
for some \(\eta_{\rm ref}>0\).  Fix also a countable family of nonnegative
functions \(\{\phi_j\}_{j\ge1}\subset C_c^\infty(Q_{\rm sh})\) with the following
separation property: for every compact subcylinder \(Q'\Subset Q_{\rm sh}\)
and every nonzero \(f\in L^1_+(Q')\), there exists \(j\) with
\(\operatorname{supp}\phi_j\Subset Q'\) and
\[
  \int_{Q_{\rm sh}}\phi_j f>0.
\]
The local occupation-state coordinates are generated by the integrals
\[
  f\mapsto \int \phi_j f.
\]
Positive local \(L^3\) mass is recorded at the first basis pair \((j,m)\) for which
\(\int \phi_j |W|^3\ge \eta_m\).

\begin{theorem}[Young-variance compactness dichotomy for realized diffuse tails]
\label{paper9:thm:young-variance-strong-compactness}
Let \(W_n\) be a realized diffuse critical-tail extraction on a rebased
unit-shell cylinder whose earlier representation, pressure, gauge, and
rough-core rows have passed.  Then exactly one of the following occurs:
\begin{enumerate}[label=\textup{(\roman*)}]
\item one of the remaining hypotheses needed for the selected compactness
package---namely the local modulation bound, the retained inner spacetime
\(L^3\) floor, or the selected upper local
bound, or the selected upper local \(L^3\) rows---fails on the rebased
unit-shell window, and the first such failure is already a named lower local
row;
\item all hypotheses of the selected compactness package hold on the rebased
unit-shell window, and
\cref{paper5:lem:selected-compactness-aubin-lions} gives strong local
\(L^3\)-compactness on the rebased unit-shell windows.
\end{enumerate}
\end{theorem}

\begin{proof}
Passing the earlier representation, pressure, gauge, and rough-core rows gives
the represented equation, local pressure reconstruction modulo functions of
time, and local compact-cylinder \(H^1\) control on the rebased unit-shell
window.  What remains to verify for the selected compactness package of
\cref{paper5:def:selected-scale-collapse-compactness-package} are precisely the
local modulation bound, the retained inner spacetime \(L^3\) floor, and the upper
local spacetime \(L^3\) rows.

If the modulation bound fails, the ordered local validation records the
autonomous-modulation or modulation-driven recentering row, which is one of the
named lower rows in
\cref{paper5:lem:selected-compactness-defect-label}.  If the retained active
spacetime \(L^3\) floor is lost, then the branch enters the failure-of-retained-spacetime-activity / pressure-only
critical-mass row, again named in
\cref{paper5:lem:selected-compactness-defect-label}.  If the upper local
\(L^3\) rows fail, then
\cref{paper6a:lem:terminal-local-critical-mass-routing} supplies the named lower
local label.  These are exactly the missing hypotheses noted in
alternative \textup{(i)}.

If none of those failures occurs, then all clauses of the selected compactness
package hold on the rebased unit-shell window.  We may therefore apply
\cref{paper5:lem:selected-compactness-aubin-lions}, which yields strong local
\(L^3\)-compactness on the rebased unit-shell windows.  This is
alternative \textup{(ii)}.
\end{proof}

\begin{lemma}[Diffuse shell states enter the augmented residual coordinates]
\label{paper9:lem:diffuse-shell-visible-residual}
Let \((W_n,\widetilde P_n,a_n,b_n)\) be a realized shell-state sequence from
\cref{paper9:thm:logarithmic-tail-support-transfer} on a fixed unit-shell model
cylinder \(Q_{\rm sh}\).  Assume that:
\begin{enumerate}[label=\textup{(\roman*)}]
\item the branch lies away from the pressure/gauge/local-energy failure faces,
so the shell realization package supplies the local pressure splitting,
repaired gauge coordinates, and local energy inequality on compact
subcylinders of \(Q_{\rm sh}\);
\item case \textup{(ii)} of \cref{paper9:thm:young-variance-strong-compactness}
holds on compact subcylinders of \(Q_{\rm sh}\), so that after passing to a
subsequence
\[
  W_n\to W_\infty
  \quad\text{strongly in }L^3_{\rm loc}(Q_{\rm sh});
\]
\item \(W_\infty\) carries positive local \(L^3\) mass on some compact
subcylinder \(Q'\Subset Q_{\rm sh}\):
\[
\int_{Q'}|W_\infty|^3\,dy\,ds>0.
\]
\end{enumerate}
Then there exists a basis test function \(\phi_j\) from the fixed countable
local occupation chart on \(Q_{\rm sh}\) such that
\[
\liminf_{n\to\infty}\int_{Q_{\rm sh}}\phi_j\,|W_n|^3\,dy\,ds>0.
\]
Consequently, after setting
\[
  U_n:=W_\infty,
  \qquad
  S_n:=W_n-W_\infty,
\]
the augmented tuple
\[
  (W_n,\widetilde P_n,U_n,S_n,a_n,b_n)
\]
belongs to the local occupation-state framework of
\cref{paper6a:lem:l3-profile-extraction-cited,paper6a:lem:critical-mass-decoupling}.
Hence every positive nonretained local mass/pressure defect that survives after
subtracting the strong local limit is visible either as a residual coordinate in
that augmented state space or as an already named pressure/gauge/rough-core
failure row.  In particular, the shell branch legitimately enters the augmented
residual coordinates of \cref{paper6a:thm:local-residual-exhaustion}.
\end{lemma}

\begin{proof}
Choose a nonnegative \(\phi\in C_c^\infty(Q_{\rm sh})\) supported in \(Q'\) with
\[
\int_{Q_{\rm sh}}\phi|W_\infty|^3\,dy\,ds>0.
\]
Strong local \(L^3\)-convergence gives
\[
\int \phi|W_n|^3\,dy\,ds\to \int \phi|W_\infty|^3\,dy\,ds>0.
\]
Because the augmented local occupation-state topology is built from a fixed
countable family of nonnegative compact test functions on \(Q_{\rm sh}\), choose
\(\phi_j\) in that family so that \(\phi_j\) is supported in \(Q'\) and
\(\int \phi_j|W_\infty|^3>0\).  Then
\[
\liminf_{n\to\infty}\int \phi_j|W_n|^3\,dy\,ds>0,
\]
which is the required visible occupation coordinate.

By \cref{paper9:thm:logarithmic-tail-support-transfer}, the realized shell
sequence is a represented local sequence on the fixed model cylinder and hence
satisfies the repaired-gauge equation, the local energy inequality, pressure
reconstruction modulo functions of time, and the compact-cylinder bounds on
smaller compact windows.  Since \(U_n\equiv W_\infty\) and
\(S_n=W_n-W_\infty\to0\) strongly in local \(L^3\), the tuple
\((W_n,\widetilde P_n,U_n,S_n,a_n,b_n)\) belongs to the augmented local
occupation-state framework on every smaller compact subcylinder.

For the pressure, let \(\widetilde P_\infty^{\rm loc}\) be the local
Calder\'on--Zygmund pressure generated by \(W_\infty\) on a slightly larger
compact cylinder.  By \cref{paper5:lem:selected-compactness-aubin-lions}, either
\[
  \widetilde P_n-c_{n,R_1}(s)-\widetilde P_\infty^{\rm loc}
  \to0
  \quad\text{strongly in }L^{3/2}_{\rm loc}
\]
modulo functions of time, or the branch has already entered the named
pressure/exterior row.  Thus every surviving local pressure defect is visible in
the same augmented local state space.  Consequently, if a positive nonretained
local diagnostic remains after subtracting the strong local limit, then on a
smaller compact subcylinder it is either a nonzero residual coordinate for
\(S_n\) in that augmented state space, or one of the already named
pressure/gauge/rough-core failure faces.  This is exactly the interface needed
to apply \cref{paper6a:thm:local-residual-exhaustion} on the diffuse shell
branch.
\end{proof}

\begin{lemma}[Strong diffuse-tail limits create a named row]
\label{paper9:lem:strong-diffuse-tail-limit-row}
Let \((W_n,\widetilde P_n)\) be a realized diffuse critical-tail extraction on a
rebased unit-shell model cylinder \(Q_{\rm sh}\), and suppose case \textup{(ii)} of
\cref{paper9:thm:young-variance-strong-compactness} holds on compact
subcylinders of \(Q_{\rm sh}\).  Let
\[
  W_n\to W_\infty
  \quad\text{strongly in }L^3_{\rm loc}(Q_{\rm sh}).
\]
Assume that the pressure/gauge/local-energy structure required to test the local
CKN, bounded-window, and occupation-state rows has not already failed.  If
\(W_\infty\not\equiv0\), then after shrinking to a smaller compact shell
subwindow one
of the following occurs:
\begin{enumerate}[label=\textup{(\roman*)}]
\item \(W_\infty\in L^\infty\) on that compact subwindow, so
\cref{paper6:lem:bounded-window-local-regularization,paper6:thm:bounded-selected-window-physical-routing}
route the branch to a named bounded-window/local row;
\item \(W_\infty\) has a singular point on that compact subwindow.  Then the
contrapositive of
\cref{paper8:thm:terminal-ckn-regularity} yields a compact subcylinder on which
the local CKN quantity of \((W_\infty,\widetilde P_\infty)\) is positive.  Strong
local convergence of the velocity and local pressure transfers a positive amount
of this CKN quantity to \((W_n,\widetilde P_n)\), and
\cref{paper6a:lem:terminal-local-critical-mass-routing} assigns a compact active,
pressure, rough-core, or other named adjacent local row;
\item \(W_\infty\) is regular on the chosen compact subwindow but carries
positive local \(L^3\) mass that is only visible below the original retained
activity level.  Then for some basis test function \(\phi_j\) from the fixed
countable shell occupation chart and some dyadic level \(\eta_m\) from the
predeclared activity ladder,
\[
\int_{Q_{\rm sh}}\phi_j|W_\infty|^3\,dy\,ds\ge \eta_m>0,
\]
and strong local \(L^3\)-convergence yields
\[
\liminf_{n\to\infty}\int_{Q_{\rm sh}}\phi_j|W_n|^3\,dy\,ds\ge \eta_m/2.
\]
Thus the shell support-transfer bookkeeping records a retained carrier at the
predeclared lower activity level \(\eta_m\);
\item the pressure/gauge/local-energy structure needed for the preceding tests
fails on the chosen compact subwindow, in which case the first failed test is an
already named local row.
\end{enumerate}
\end{lemma}

\begin{proof}
Because the realized diffuse support profile is nonzero, there is a compact
subcylinder \(Q'\Subset Q_{\rm sh}\) on which \(W_\infty\) has positive \(L^3\)
mass.  If during passage to the limit the pressure atlas, repaired gauge, or
local-energy structure breaks on a compact shell window, then the first failed
row is already one of the named local failure rows.  This is
alternative \textup{(iv)}.

Assume henceforth that the local suitable structure survives on compact shell
windows.  If \(W_\infty\) is bounded on a slightly smaller compact subwindow,
then
\cref{paper6:lem:bounded-window-local-regularization} applies directly and gives
alternative \textup{(i)} unless a named remaining-activity row has already been
detected, in which case we are done anyway.

Assume next that \(W_\infty\) is not bounded on the chosen compact shell
subwindow.  Because a local suitable weak solution has open regular set, there is
then a singular point on that subwindow.  If every sufficiently small compact
subcylinder around that point had CKN quantity below the threshold of
\cref{paper8:thm:terminal-ckn-regularity}, that theorem would imply local
boundedness of \(W_\infty\), contradiction.  Hence there is a compact
subcylinder on which the local CKN quantity of the limit is positive.  By the
strong local \(L^3\) convergence and the pressure compactness conclusion of
\cref{paper5:lem:selected-compactness-aubin-lions}, the same positive CKN
quantity transfers to \((W_n,\widetilde P_n)\) for all large \(n\), unless
pressure compactness fails first; such a pressure failure is
alternative \textup{(iv)}.  In the non-failure case,
\cref{paper6a:lem:terminal-local-critical-mass-routing} reads this positive local
diagnostic as a named compact active/pressure/rough-core/adjacent row.  This is
alternative \textup{(ii)}.

It remains to treat the case where \(W_\infty\) is regular on a compact shell
subwindow and hence locally bounded after shrinking, but the bookkeeping of the
original shell branch has not yet registered activity at the previously retained
level.  Choose a nonnegative basis test function \(\phi_j\) from the fixed
countable shell occupation chart, supported in such a subwindow, with
\(\int \phi_j|W_\infty|^3>0\).  By the dyadic activity ladder there exists
\(m\) such that
\[
0<\eta_m\le \int \phi_j|W_\infty|^3\,dy\,ds.
\]
Strong local \(L^3\)-convergence then gives
\[
\int \phi_j|W_n|^3\,dy\,ds\to \int \phi_j|W_\infty|^3\,dy\,ds,
\]
so \(\int \phi_j|W_n|^3\ge \eta_m/2\) for all large \(n\).  Thus the shell
support-transfer mechanism records a retained local carrier at the predeclared
activity level \(\eta_m\).  This is alternative \textup{(iii)}.

Taking the first closed row in the ordered shell/local state space makes the
four outcomes exclusive.  Therefore one of them occurs.
\end{proof}

\begin{theorem}[Coherent critical-tail closure]
\label{paper9:thm:coherent-critical-tail-closure}
A nonzero realized diffuse critical-tail state cannot remain terminal on the
retained compact single-core branch.  More precisely, every such state either
produces a lower named local row immediately, or else after strong local
compactness its local limit enters the bounded-window/local-routing package and
any surviving shell remainder is exhausted by the augmented residual state space.
In particular no realized diffuse critical-tail state is terminal.
\end{theorem}

\begin{proof}
Let \(W_n\) be a realized diffuse critical-tail state.  By
\cref{paper9:thm:young-variance-strong-compactness}, either the selected
compactness package fails and we immediately obtain a lower named local row, or
else \(W_n\) is strongly compact in local \(L^3\) on rebased unit-shell
windows.  Choose a compact unit-shell subwindow on which the strong local limit
is nonzero.  By \cref{paper9:lem:strong-diffuse-tail-limit-row}, after shrinking
that subwindow if necessary, either the limit is locally bounded and the
bounded-window row closes the branch, or it has a singular point and the local
critical-mass routing lemma places the branch in one of the already named local
rows, or a smaller positive dyadic activity level is visible and the shell
support-transfer bookkeeping records a retained carrier at that predeclared
level, or the pressure/gauge/local-energy interface fails and gives the
corresponding named failure row.

If, after subtracting the strong local limit on such a compact subwindow, a
positive nonretained shell remainder still survives, then
\cref{paper9:lem:diffuse-shell-visible-residual} places that remainder inside the
augmented residual coordinates of
\cref{paper6a:thm:local-residual-exhaustion}.  That theorem yields either a
named adjacent local row or a retained comparable class.  A retained comparable
class cannot occur for a diffuse tail, because comparable retained compact
carriers were already grouped before the shell extraction and the diffuse shell
selection was made outside every retained compact carrier.  Hence only named
lower rows remain.

Therefore no realized diffuse critical-tail state is terminal on the retained
compact single-core branch.
\end{proof}

\begin{definition}[Windowed critical tightness]
\label{paper9:def:windowed-critical-tightness}
A selected terminal window modeled on a compact interval
\(J_*\Subset\mathbb R\) means a window in the local terminal validation family
obtained from \(J_*\) by the admissible finite covariant recenterings of
\cref{paper6a:lem:interior-time-selection}.  A retained represented branch is
windowed critically tight on the terminal tail if for every
\(\varepsilon>0\) and every compact model interval \(J_*\Subset\mathbb R\)
there exists \(R_{\varepsilon,J_*}<\infty\) such that for every selected
terminal window \(J\) modeled on \(J_*\),
\[
  \int_J\int_{|y|>R_{\varepsilon,J_*}}|V(y,\tau)|^3\,dy\,d\tau
  <\varepsilon .
\]
Failure of windowed critical tightness means that there exist
\(\varepsilon_0>0\), a compact model interval \(J_*\Subset\mathbb R\),
selected terminal windows \(J_n\) modeled on \(J_*\), and radii
\(R_n\to\infty\) such that
\[
  \int_{J_n}\int_{|y|>R_n}|V(y,\tau)|^3\,dy\,d\tau
  \ge \varepsilon_0 .
\]
\end{definition}

\begin{remark}[Pointwise time-slice \(L^3\) persistence is not used in the unconditional routing]
\label{paper9:lem:timeslice-l3-persistence}
The unconditional terminal routing does not convert a single-time \(L^3\) tail
lower bound into a spacetime row.  Instead it works directly with the
windowed spacetime tail defect of
\cref{paper9:def:windowed-critical-tightness}, whose shell occupation
measures are available without an additional time-persistence argument.
\end{remark}

\begin{lemma}[Windowed diffuse critical tail becomes a spacetime row]
\label{paper9:lem:diffuse-timeslice-tail-spacetime-row}
Let \(V\) be a retained compact single-core branch, and suppose there exist
\(\varepsilon_0>0\), a compact model interval \(J_*\Subset\mathbb R\),
selected terminal windows \(J_n\) modeled on \(J_*\), and radii
\(R_n\to\infty\) such that
\[
  \int_{J_n}\int_{|y|>R_n}|V(y,\tau)|^3\,dy\,d\tau\ge\varepsilon_0 .
\]
Assume no rebased unit shell cylinder carries a positive compact spacetime
\(L^3\) activity level that is already known to realize a retained carrier.
Then exactly one of the following occurs:
\begin{enumerate}[label=\textup{(\roman*)}]
\item during shell extraction the represented local structure, pressure atlas,
  gauge chart, local energy, rough-core, or compact-upper package fails on the
  rebased shell cylinders, producing the corresponding named failure row;
\item after passing to a subsequence, the normalized spacetime tail measures
  define a nonzero diffuse shell occupation coordinate in the
  logarithmic shell compactification;
\item a parabolic shell support profile with positive compact spacetime
  activity is
  realized, reducing to the carrier case of
  \Cref{paper9:lem:critical-tightness-failure-row}.
\end{enumerate}
\end{lemma}

\begin{proof}
Decompose the spacetime tail into dyadic shell blocks
\[
  A_{n,k}\times J_n
  =
  \{2^kR_n<|y|\le 2^{k+1}R_n\}\times J_n .
\]
Cover each occupied block by finitely overlapping unit parabolic cylinders in
rebased shell coordinates.  If the represented local structure or one of the
pressure-atlas, gauge, local-energy, rough-core, or compact-upper inputs fails
on those rebased shell cylinders, then the first failed test is exactly the
corresponding named row.  This is alternative \textup{(i)}.

Assume these failure faces do not occur.  If some rebased unit shell cylinder
carries a positive compact spacetime \(L^3\) level, then
\cref{paper9:lem:logarithmic-tail-shell-realization,paper9:thm:logarithmic-tail-support-transfer}
produce a realized shell support profile with positive compact activity on a
fixed model cylinder.  This is alternative \textup{(iii)}.

If no unit shell cylinder carries such a positive compact spacetime level,
normalize the spacetime tail measures by their total mass and pass to the
logarithmic shell compactification.  Apply
\cref{paper6a:lem:normalized-defect-atomic-or-diffuse}.  If the
atomic/activity alternative occurs, a realized support point with positive
compact local activity appears on a fixed shell cylinder, returning us to
alternative \textup{(iii)}.  If the genuinely diffuse alternative occurs,
then the nonzero limiting shell measure itself is a diffuse shell occupation
coordinate.  This is alternative \textup{(ii)}.
\end{proof}

\begin{lemma}[Failure of windowed critical tightness is a named local row]
\label{paper9:lem:critical-tightness-failure-row}
\label{paper6a:lem:retained-tightness-routing}
Let \(V\) be a retained compact single-core branch in the renormalized
variables, with scale function \(\lambda(\tau)\).  Suppose windowed critical
tightness fails in the sense of
\Cref{paper9:def:windowed-critical-tightness}: there exist
\(\varepsilon_0>0\), a compact model interval \(J_*\Subset\mathbb R\),
selected terminal windows \(J_n\) modeled on \(J_*\), and radii
\(R_n\to\infty\) such that
\[
  \int_{J_n}\int_{|y|>R_n}|V(y,\tau)|^3\,dy\,d\tau
  \ge \varepsilon_0 .
\]
Then, after passing to a subsequence, one of the following named rows occurs:
\begin{enumerate}[label=\textup{(\roman*)}]
\item a same-point larger-scale cascade row;
\item a separated physical exterior concentration row;
\item a noncompact escaping exterior row;
\item a pressure/gauge/reconstruction failure row;
\item a rough-core, local-energy, or compact-upper failure row;
\item a diffuse critical-tail row.
\end{enumerate}
In particular, failure of critical tightness cannot remain inside the retained
compact single-core terminal stratum.
\end{lemma}

\begin{proof}
Decompose the spacetime tail
\[
  \{(y,\tau):\tau\in J_n,\ |y|>R_n\}
\]
into dyadic spacetime shell blocks
\[
  A_{n,k}\times J_n
  =
  \{2^kR_n<|y|\le 2^{k+1}R_n\}\times J_n .
\]
Cover each occupied shell by finitely overlapping translated shell cylinders
modeled on \(Q_{\rm sh}\).  If the represented local structure, pressure atlas,
or repaired-gauge package fails during this extraction, we are in item
\textup{(iv)}.  If the local-energy, rough-core, or compact-upper package
fails, we are in item \textup{(v)}.

Assume these failure faces do not occur.  Suppose first that some rebased unit
shell cylinder carries a positive compact spacetime \(L^3\) mass level.  Then
\cref{paper9:lem:logarithmic-tail-shell-realization,paper9:thm:logarithmic-tail-support-transfer}
furnish a realized shell support profile coming from translated shell
cylinders with centers \((y_n,\tau_n)\) in the original renormalized
variables and \(|y_n|\simeq R_n\).  Set
\[
  d_n:=\lambda(\tau_n)|y_n| .
\]
Passing to a subsequence, exactly one of the following occurs.  If
\(d_n\to0\), the support profile stays at the same physical point while moving
to a larger renormalized radius, giving item \textup{(i)}.  If
\(d_n\to d\in(0,\infty)\), the profile lies at positive physical distance from
the selected core, giving item \textup{(ii)}.  If \(d_n\to\infty\), the profile
escapes every compact physical neighborhood in the selected chart, giving item
\textup{(iii)}.

The only remaining possibility is that no positive local carrier exists at any
fixed level.  Normalize the spacetime tail measures by their total mass and
pass to the logarithmic shell compactification.  Apply
\cref{paper6a:lem:normalized-defect-atomic-or-diffuse}.  If the
atomic/activity alternative occurs, then
\cref{paper9:thm:logarithmic-tail-support-transfer} realizes a shell support
point with positive compact spacetime activity on a translated shell cylinder
modeled on \(Q_{\rm sh}\),
returning us to the carrier cases already handled above and therefore to one
of the items \textup{(i)}--\textup{(iii)}.  If the genuinely diffuse
alternative occurs, then the nonzero limiting shell measure itself is a
genuinely diffuse shell occupation coordinate, yielding item \textup{(vi)}.
This exhausts all possibilities.
\end{proof}

\begin{lemma}[Diffuse critical-tail row is closed]
\label{paper9:lem:diffuse-critical-tail-row-closed}
A diffuse critical-tail defect produced by
\cref{paper9:lem:critical-tightness-failure-row} cannot remain as a terminal
retained compact single-core branch.  It either produces a lower named row
already present in the state-space ledger, or it contradicts the minimal
mesoscopic/critical-tail closure theorem.
\end{lemma}

\begin{proof}
Let \(\nu\) be the normalized limiting tail-defect measure produced in
\cref{paper9:lem:critical-tightness-failure-row}.  By construction, \(\nu\) is
nonzero.  Apply \cref{paper9:lem:logarithmic-tail-shell-realization} to the
logarithmic shell compactification.  If the shell compactification hits a
pressure, gauge, local-energy, or rough-core failure face first, that face is
already the corresponding named local row.

Apply \cref{paper6a:lem:normalized-defect-atomic-or-diffuse} to the
logarithmic shell defect measures.  If the atomic/activity alternative occurs,
then \cref{paper9:thm:logarithmic-tail-support-transfer} realizes a support
point with positive compact local activity on a fixed shell model cylinder.
This is a retained local carrier in the original branch, contrary to the
diffuse case.

If the genuinely diffuse alternative occurs, then the nonzero diffuse shell
measure itself is a visible diffuse occupation-state coordinate.  If the
selected compactness package fails before the strong-shell stage, the first
failed test is already a named lower row.  Otherwise the same augmented-shell
mechanism used in \cref{paper9:lem:diffuse-shell-visible-residual} places this
diffuse coordinate inside the residual coordinates of
\cref{paper6a:thm:local-residual-exhaustion}.  That theorem yields either a
named local row or a retained comparable class.  The retained comparable class
is impossible because comparable retained compact carriers were grouped before
the shell extraction.

Thus a diffuse critical-tail defect cannot be terminal in the retained compact
single-core stratum.
\end{proof}

\begin{lemma}[Windowed Caccioppoli with integrated modulation]
\label{paper6a:lem:windowed-h1-from-integrated-modulation}
Let \((V,P,a,b)\) be a suitable repaired-gauge branch on a terminal tail.
Fix \(R>0\) and a unit window \(I_n=[\tau_0+n,\tau_0+n+1]\), and set
\[
  I_n^*:=[\tau_0+n-1,\tau_0+n+2],
  \qquad
  A_{n,R}:=1+\int_{I_n^*}\bigl(|a(\tau)|+\|b(\tau)\|_\infty\bigr)\,d\tau .
\]
Assume on \(I_n^*\times B_{4R}\) the local pressure reconstruction, the
compact-cylinder critical spacetime bound
\[
  \int_{I_n^*}\int_{B_{4R}}|V|^3\,dy\,d\tau\le M_{3,4R},
\]
the local energy bound
\[
  \sup_{\tau\in I_n^*}\int_{B_{4R}}|V(y,\tau)|^2\,dy\le E_{2,4R},
\]
and the local Caccioppoli estimate of
\cref{paper6a:thm:ac-local-caccioppoli}.  Then
\[
  \int_{I_n}\int_{B_R}|\nabla V|^2\,dy\,d\tau
  \le
  C_R\bigl(M_{3,4R}+E_{2,4R}(1+A_{n,R})\bigr),
\]
where \(C_R\) depends only on the fixed compact-cylinder cutoffs and \(\nu\).
\end{lemma}

\begin{proof}
Apply the local Caccioppoli estimate on \(I_n^*\times B_{2R}\), with a cutoff
equal to one on \(I_n\times B_R\).  The endpoint and cutoff-error terms are
bounded by the local energy control:
\[
  \int_{I_n^*}\int_{B_{2R}}|V|^2\,dy\,d\tau\le C_R E_{2,4R}.
\]
The cubic convection term is controlled directly by the spacetime \(L^3\)
bound:
\[
  \int_{I_n^*}\int_{B_{2R}}|V|^3\,dy\,d\tau\le M_{3,4R}.
\]
The pressure term is controlled, after subtracting spatial means, by the local
pressure reconstruction and Calder\'on--Zygmund on the larger ball:
\[
  \|P-c_{4R}\|_{L^{3/2}(I_n^*\times B_{2R})}
  \le C_R\|V\otimes V\|_{L^{3/2}(I_n^*\times B_{4R})}
  \le C_R M_{3,4R}^{2/3} .
\]
Hence H\"older yields
\[
  \int_{I_n^*}\int_{B_{2R}}|P-c_{4R}|\,|V|\,dy\,d\tau
  \le C_R M_{3,4R} .
\]
The modulation terms in the local energy inequality have the form
\[
  \int_{I_n^*}\bigl(|a(\tau)|+\|b(\tau)\|_\infty\bigr)
  \int_{B_{2R}}|V(\tau)|^2\,dy\,d\tau,
\]
up to cutoff constants; the local \(L^\infty_\tau L^2_y\) bound controls the
inner integral by \(E_{2,4R}\).  Therefore
\[
  \int_{I_n^*}\bigl(|a(\tau)|+\|b(\tau)\|_\infty\bigr)
  \int_{B_{2R}}|V(\tau)|^2\,dy\,d\tau
  \le E_{2,4R}A_{n,R} .
\]
Combining these estimates gives the displayed bound.  No pointwise-in-time
critical \(L^3\) control or pointwise-in-time modulation bound is used.
\end{proof}

\begin{lemma}[Unbounded modulation windows are adjacent exits]
\label{paper6a:lem:modulation-window-bound-or-exit}
Let a physical branch be in the retained compact single-core stratum after the
ordered local validation tests.  Fix \(R>0\).  Then either
\[
  \sup_n
  \int_{\tau_0+n-1}^{\tau_0+n+2}
  \bigl(|a(\tau)|+\|b(\tau)\|_\infty\bigr)\,d\tau<\infty,
\]
or the branch leaves the retained compact single-core stratum through one of
the named modulation, scale-drift, modulation-driven recentering, exterior, or
gauge-degenerate local alternatives.
\end{lemma}

\begin{proof}
If the displayed window integrals are unbounded, choose windows \(I_{n_j}^*\)
on which the integral tends to infinity.  The coefficients \(a,b\) are the
absolutely continuous scale and translation velocities of the repaired-gauge
chart, as identified in
\cref{paper6a:lem:repaired-modulation-coefficients}.  Large accumulated
negative scale motion is the negative scale-drift branch of
\cref{paper6a:def:negative-drift}; large positive scale-shell work is the
positive-scale carrier branch routed by
\cref{paper7:cor:windowwise-positive-scale-routing-holds}; and large
translation work either is absorbed by the admissible recentering rule or
produces the modulation-driven recentering/exterior label in
\cref{paper7:def:stratified-moving-carrier-routing}.  If the accumulated
modulation cannot be represented by the absolutely continuous chart variables,
the first failed row is the repaired-gauge or gauge-degenerate alternative.
These are all named adjacent local alternatives, not retained compact
single-core states.  Hence on the retained compact single-core stratum the
window integrals are uniformly bounded.
\end{proof}

\begin{lemma}[Local windowed \(H^1\)-loss routes to the rough-core alternative]
\label{paper6a:lem:retained-h1-routing}
Assume the retained compact single-core stratum, with a suitable repaired-gauge
representation, local pressure reconstruction, compact-cylinder critical
spacetime bounds, local \(L^\infty_\tau L^2_y\) energy control, and the local
Caccioppoli estimate on the compact cylinders used in
\cref{paper6a:lem:windowed-h1-from-integrated-modulation}.  These are the
retained-stratum inputs supplied by the ordered tests in
\cref{paper7:lem:no-unrecorded-cost-escape}\textup{(i),(ii),(iv),(vii)} and by
the suitable local energy package.  Assume also that the unbounded-modulation
window exit of \cref{paper6a:lem:modulation-window-bound-or-exit} has not
occurred.
Then, for every \(R>0\), the local windowed \(H^1\) bound
\[
  \sup_{T\ge\tau_0+1}
  \int_T^{T+1}\int_{B_R}|\nabla V|^2\,dy\,d\tau
  \le C\!\left(R,\nu,M_{3,4R},E_{2,4R},\sup_n A_{n,R}\right)
\]
holds, so the local windowed \(H^1\)-control condition of
\cref{def:c5-compact-tests} is automatic on the retained stratum.
Equivalently, failure of local windowed \(H^1\) control is incompatible
with the retained-stratum input package and is routed to the rough-core
alternative \textup{(iii)} of \cref{paper6:thm:state-decomposition},
which by \cref{paper6:thm:rough-core} forces failure of compact CKN
control on enclosing cylinders and therefore returns to the
multibubble/cascade alternative \textup{(ii)}, already excluded from the
retained branch.
\end{lemma}

\begin{proof}
By \cref{paper6a:lem:modulation-window-bound-or-exit}, absence of the
modulation-window exit gives, for every fixed \(R\),
\[
  \sup_n A_{n,R}<\infty .
\]
The compact-cylinder critical spacetime bound and local pressure reconstruction are
retained-stratum inputs, together with the local
\(L^\infty_\tau L^2_y\) energy control.  Applying
\cref{paper6a:lem:windowed-h1-from-integrated-modulation} on each unit window
therefore gives a bound independent of \(n\).  This proves the displayed
windowed \(H^1\) estimate.  The standard mollification and lower
semicontinuity argument passes the estimate from smooth approximants to
suitable repaired-gauge branches.

If, contrapositively, the windowed \(H^1\) bound fails for some \(R>0\)
along a terminal sequence, then by
\cref{paper6a:lem:windowed-h1-from-integrated-modulation} and
\cref{paper6a:lem:modulation-window-bound-or-exit} at least one retained input
must fail along that sequence: local pressure reconstruction, compact-cylinder
critical spacetime control, the local \(L^\infty_\tau L^2_y\) energy control,
local Caccioppoli regularity, or the uniform integrated modulation-window
bound.  By the ordered tests
\cref{paper7:lem:no-unrecorded-cost-escape} this earlier-test failure
routes the branch out of the retained stratum into the corresponding
representation, repaired-gauge, critical-mass, local-regularity, or
modulation/scale-drift alternative.  Independently, in the local state-space
terminology of
\cref{paper6:thm:state-decomposition} the same failure is recorded as
the rough-core alternative \textup{(iii)}, and by
\cref{paper6:thm:rough-core} that alternative is equivalent to failure
of compact CKN control on enclosing cylinders, hence to active
multibubble or cascade behaviour, alternative \textup{(ii)} of
\cref{paper6:thm:state-decomposition}.  Both routings exit the retained
single-core stratum.  Therefore on the retained stratum
\(H^1\)-loss is impossible.
\end{proof}

\begin{theorem}[Retained compact-test closure]
\label{paper6a:thm:retained-compact-tests-closed}
On the retained compact single-core stratum, both compact tests of
\cref{def:c5-compact-tests} pass: critical tightness holds and the local
windowed \(H^1\) control holds.  Equivalently, the entries
\textup{(v),(vi)} of \cref{paper7:lem:no-unrecorded-cost-escape} are not
genuine new assumptions on the retained stratum; their failure is routed
to one of the named non-retained alternatives.
\end{theorem}

\begin{proof}
Windowed critical tightness on selected terminal windows holds by
\cref{paper9:lem:critical-tightness-failure-row}: failure routes the branch
to a named local row.  The carrier rows are handled by the same-point cascade,
separated exterior, or noncompact exterior closures.  The pressure/gauge and
the rough-core, local-energy, and compact-upper rows are closed by the
corresponding first-failure theorems.  The
diffuse tail case is closed by
\cref{paper9:lem:diffuse-critical-tail-row-closed}.  Local
windowed \(H^1\) control holds by
\cref{paper6a:lem:retained-h1-routing}: failure routes the branch to the
rough-core alternative, which by \cref{paper6:thm:rough-core} is
absorbed into the multibubble/cascade alternative, also non-retained.
Hence both compact tests are theorems, not assumptions, on the retained
stratum.
\end{proof}

\begin{corollary}[Discharge of the retained compact-test closure row]
\label{paper6a:cor:compact-test-closure-discharged}
The retained compact-test closure obligation listed as item~3 of the
local closure ledger is discharged on every retained compact single-core branch
produced by the ordered local validation procedure of
\cref{paper7:lem:no-unrecorded-cost-escape}.  No additional global
tightness or global windowed \(H^1\) hypothesis is required.
\end{corollary}

\begin{proof}
\Cref{paper6a:thm:retained-compact-tests-closed} reduces both compact
tests to local consequences of inputs already present on the retained
stratum or to named non-retained alternatives.  No global tightness or
\(H^1\) bound is invoked.
\end{proof}

\subsection{Closure of the thin negative-drift exit}
\label{paper6a:subsec:thin-negative-drift-closure}

The remaining point in the scale-collapse drift branch is the case in which
the weighted negative drift
\[
  g_{\rm sc}(\tau):=a_-(\tau)M_R(\tau),
  \qquad
  M_R(\tau)=\int |V(y,\tau)|^2\Phi(\tau,y)\,dy,
\]
has infinite tail integral but does not produce fixed-length windows with a
uniform positive average.  The following argument closes this row as an
independent terminal alternative.  It does not turn an infinite integral into a
pointwise lower bound; instead it passes to normalized occupation blocks and
uses the local state-space labels.

\begin{definition}[Thin negative-drift occupation blocks]
\label{paper6a:def:thin-negative-drift-occupation-blocks}
Let a represented branch satisfy the scale-collapse drift obstruction
\[
  \int_{\tau_0}^{\infty} g_{\rm sc}(\tau)\,d\tau=\infty .
\]
A sequence of intervals \([s_m,t_m]\), \(s_m\to\infty\), is a normalized
negative-drift occupation block sequence if
\[
  \int_{s_m}^{t_m}g_{\rm sc}(\tau)\,d\tau=1
  \qquad\text{for every }m.
\]
It is diffuse if
\[
  \sup_{T\in[s_m,t_m]}
  \int_{[T,T+1]\cap[s_m,t_m]} g_{\rm sc}(\tau)\,d\tau
  \longrightarrow0 .
\]
The associated occupation measures are
\[
  d\mu_m(\tau)=g_{\rm sc}(\tau)\mathbf 1_{[s_m,t_m]}(\tau)\,d\tau .
\]
\end{definition}

\begin{definition}[Canonical thin-drift terminal schedule]
\label{paper6a:def:canonical-thin-drift-schedule}
In the canonical windowwise pre-carrier regime, after a harmless terminal
restriction, the terminal unit-window schedule is
\[
  K_k=[\tau_\sharp+k,\tau_\sharp+k+1],
  \qquad k=0,1,2,\ldots .
\]
These are the canonical cells on which the windowwise moving family and the
ordered local validation rows are tested.
\end{definition}

\begin{lemma}[Canonical terminal windows are admissible or already labelled]
\label{paper6a:lem:terminal-unit-drift-window-admissibility}
Let a physical local Type II branch be in the canonical windowwise
pre-carrier regime, and let \(K_k\) be the canonical cells of
\cref{paper6a:def:canonical-thin-drift-schedule}.  For every sequence
\(K_{k_m}\) with \(k_m\to\infty\), after passing to a subsequence, exactly one
of the following alternatives holds:
\begin{enumerate}[label=\textup{(A\arabic*)}]
\item a non-thin first-failure row occurs on these windows: autonomous
      modulation failure, modulation-driven recentering, compactness-package
      first failure routed by
      \cref{paper5:thm:selected-compactness-failure-routed}, exterior or
      separated-core concentration, rough-core loss, multicore or same-point
      cascade, carrier/cost validation failure outside the canonical retained
      channel, or loss of the retained active core;
\item after the finite covariant recentering \(s=\tau-(\tau_\sharp+k_m)\), the
      canonical windows \(K_{k_m}\) are admissible selected terminal windows
      for the ordered local validation: the repaired-gauge representation,
      local pressure
      reconstruction, local suitable-energy bounds, moving \(L^2\) core floor,
      retained local critical-mass row, compact tests, and modulation
      bookkeeping hold on every fixed compact cylinder, with constants
      independent of \(m\).
\end{enumerate}
\end{lemma}

\begin{proof}
The canonical windowwise moving family is the schedule on which the carrier,
compactness, repaired-gauge, pressure, and modulation rows are validated.  The
finite translation \(\tau\mapsto\tau-(\tau_\sharp+k_m)\) does not transfer any
lower bound from one physical time to another; it only writes the canonical
cell \(K_{k_m}\) as the model interval \([0,1]\) in its terminal chart, as in
the finite recentering step of
\cref{paper6a:lem:interior-time-selection}.  Thus every row in the ordered
validation list is tested on the same canonical cells.
\par
Run the finite ordered validation list on \(K_{k_m}\).  If any required
local row fails before retention, the first failed row is recorded by
\cref{paper7:lem:no-unrecorded-cost-escape,paper5:lem:retained-scale-collapse-first-failure}
and belongs to the list in \textup{(A1)} after excluding the thin-drift row
itself.  If no such first failure occurs, all retained local inputs listed in
\textup{(A2)} hold on the recentered windows.  The moving \(L^2\) floor is
part of this retained local state data on the same canonical windows; if that
floor fails, the ordered validation records the first such failure before any
thin-drift conclusion.  Thus only the stated alternatives remain.
\end{proof}

\begin{lemma}[Thin drift gives diffuse occupation blocks]
\label{paper6a:lem:thin-drift-diffuse-blocks}
Assume \(g_{\rm sc}\ge0\), \(g_{\rm sc}\in L^1_{\rm loc}([\tau_0,\infty))\),
and
\[
  \int_T^\infty g_{\rm sc}(\tau)\,d\tau=\infty
  \qquad\text{for every }T\ge\tau_0 .
\]
Then every terminal tail contains a normalized negative-drift occupation block
sequence.  After passing to a subsequence, exactly one of the following holds:
\begin{enumerate}[label=\textup{(D\arabic*)}]
\item there are \(\gamma>0\) and unit windows \(U_m=[\sigma_m,\sigma_m+1]\),
      \(\sigma_m\to\infty\), such that
      \[
        \int_{U_m}g_{\rm sc}(\tau)\,d\tau\ge\gamma ;
      \]
\item the retained block sequence is diffuse in the sense of
      \cref{paper6a:def:thin-negative-drift-occupation-blocks}.
\end{enumerate}
If the branch is in the thin-drift alternative of
\cref{paper5:thm:scale-collapse-stratification}, then case
\textup{(D1)} either produces one of the non-thin first-failure rows of
\cref{paper6a:lem:terminal-unit-drift-window-admissibility}\,\textup{(A1)}, or
is impossible on the retained selected terminal tail.  Hence, after non-thin
first failures are removed, only \textup{(D2)} remains.
\end{lemma}

\begin{proof}
Apply \cref{paper7:lem:nonintegrable-component-certificate} with
\(g=g_{\rm sc}\).  Its block construction gives disjoint terminal intervals
\([s_m,t_m]\) with \(\int_{s_m}^{t_m}g_{\rm sc}=1\).  The same lemma gives the
two alternatives: persistent unit-window carrier mass or diffuse carrier mass.
These are exactly \textup{(D1)} and \textup{(D2)} above.
\par
If \textup{(D1)} holds, compare \(U_m\) with the canonical schedule
\(K_k=[\tau_\sharp+k,\tau_\sharp+k+1]\).  Each unit interval \(U_m\) meets at
most three canonical cells.  Hence, after passing to a subsequence, there are
canonical cells \(K_{k_m}\) with \(k_m\to\infty\) and
\[
  \int_{K_{k_m}}g_{\rm sc}(\tau)\,d\tau\ge\gamma/3 .
\]
Apply \cref{paper6a:lem:terminal-unit-drift-window-admissibility} to these
canonical cells.  If alternative \textup{(A1)} occurs, the branch has already
left the thin terminal row through a non-thin first failure.  Otherwise the
canonical cells are admissible selected terminal windows.  On those selected
windows,
\[
  \int_{K_{k_m}}a_-(\tau)M_R(\tau)\,d\tau
  \ge \gamma/3 .
\]
This is the negation of the thin-drift alternative in
\cref{paper5:thm:scale-collapse-stratification}, which was precisely the
absence of a uniform positive lower bound for fixed-length averaged weighted
negative drift on the selected terminal tail.  Thus a retained thin-drift
branch with no non-thin first failure can only produce diffuse occupation
blocks.
\end{proof}

\begin{lemma}[Small weighted negative drift is perturbative]
\label{paper6a:lem:small-negative-drift-perturbative}
Let \(I\subset[\tau_0,\infty)\) be an interval of length at most one.  Assume
that the branch has a moving \(L^2\) core floor
\[
  M_R(\tau)\ge m_0>0
  \qquad\text{for a.e. }\tau\in I,
\]
and that on a compact set \(K^+\Subset\mathbb R^3\), containing a unit
neighborhood of \(K\), the retained local energy bound
\[
  \operatorname*{ess\,sup}_{\tau\in I}\|V(\tau)\|_{L^2(K^+)}\le L
\]
holds.  If
\[
  \int_I g_{\rm sc}(\tau)\,d\tau\le\varepsilon ,
\]
then the negative scale correction
\[
  F_-(\tau):=a_-(\tau)\bigl(V(\tau)+y\cdot\nabla V(\tau)\bigr)
\]
satisfies
\[
  \|F_-\|_{L^1(I;H^{-1}(K))}
  \le C_K\,L\,m_0^{-1}\varepsilon ,
\]
where \(C_K\) depends only on \(K\) and the fixed enlargement \(K^+\).
Moreover, for every \(\psi\in C_c^\infty(K)\),
\[
  \int_I
  \left|\left\langle F_-(\tau),V(\tau)\psi\right\rangle\right|\,d\tau
  \le C_{K,\psi}L^2m_0^{-1}\varepsilon .
\]
\end{lemma}

\begin{proof}
The core floor gives
\[
  \int_I a_-(\tau)\,d\tau
  \le m_0^{-1}\int_I a_-(\tau)M_R(\tau)\,d\tau
  \le m_0^{-1}\varepsilon .
\]
Let \(\varphi\in C_c^\infty(K;\mathbb R^3)\) with
\(\|\varphi\|_{H^1_0(K)}\le1\).  By integration by parts in \(y\),
\[
  \int_K (y\cdot\nabla V)\cdot\varphi\,dy
  =
  -\int_K V\cdot\bigl(y\cdot\nabla\varphi+3\varphi\bigr)\,dy .
\]
Hence, for a.e. \(\tau\in I\),
\[
  |\langle V+y\cdot\nabla V,\varphi\rangle|
  \le C_K\|V(\tau)\|_{L^2(K^+)}\|\varphi\|_{H^1_0(K)}
  \le C_KL .
\]
Taking the supremum over \(\varphi\) and integrating against \(a_-(\tau)\)
gives
\[
  \|F_-\|_{L^1(I;H^{-1}(K))}
  \le C_KL\int_Ia_-(\tau)\,d\tau
  \le C_KLm_0^{-1}\varepsilon .
\]
For the local energy pairing, integrate the scale-generator term by parts:
\[
  \int_K (y\cdot\nabla V)\cdot V\,\psi\,dy
  =
  \frac12\int_K y\cdot\nabla(|V|^2)\psi\,dy
  =
  -\frac12\int_K |V|^2(3\psi+y\cdot\nabla\psi)\,dy .
\]
Therefore
\[
  \left|\left\langle V+y\cdot\nabla V,V\psi\right\rangle\right|
  \le C_{K,\psi}\|V(\tau)\|_{L^2(K^+)}^2
  \le C_{K,\psi}L^2 .
\]
Integrating against \(a_-\) and using
\(\int_Ia_-\le m_0^{-1}\varepsilon\) gives the second estimate.  No bound on
\(\nabla V\) and no global norm is used.
\end{proof}

\begin{lemma}[Active core persistence on admissible drift windows]
\label{paper6a:lem:admissible-drift-window-active-core}
Let \(I_m=[\sigma_m,\sigma_m+1]\) be admissible selected terminal windows in the sense of
\cref{paper6a:lem:terminal-unit-drift-window-admissibility}\,\textup{(A2)}.
Then either the retained spacetime activity floor fails (equivalently, the
loss-of-retained-active-core row occurs), or, after passing to a
subsequence, there are a compact cylinder
\(Q_{r_*}=B_{r_*}\times I_{r_*}\Subset B_R\times(0,1)\) and \(\eta_*>0\) such
that the represented sequence on \(I_m\) carries the retained active local mass
floor
\[
  \liminf_{m\to\infty}
  \int_{I_{r_*}}\int_{B_{r_*}}|V_m|^3\,dy\,ds
  \ge\eta_*^3,
\]
where \(V_m(y,s)=V(y,\sigma_m+s)\) and \(P_m(y,s)=P(y,\sigma_m+s)\).
Consequently the minimal active-family extraction and the same-point/separated
reductions are applicable on these windows; if pressure-only CKN detection is
needed instead of velocity detection, the branch is routed to the
pressure/exterior row rather than retained here.
\end{lemma}

\begin{proof}
The retained local critical-mass row in
\cref{paper6a:lem:terminal-unit-drift-window-admissibility}\,\textup{(A2)} is
interpreted by \cref{paper6a:lem:terminal-local-critical-mass-routing}.  If the
retained spacetime activity floor fails, the selected candidate is not a retained active
component; this is exactly the loss-of-retained-active-core row.
\par
If the retained spacetime activity floor does not fail, then
\cref{paper8:thm:active-profile-mass-floor} supplies a positive active local
spacetime velocity floor on a compact subcylinder of the selected terminal window.  A
finite cover of the retained compact core and the terminal maximality rule
\cref{paper6a:lem:retained-frame-maximality} allow one fixed compact
subcylinder \(Q_{r_*}\) and one threshold \(\eta_*>0\) after passing to a
subsequence.  If the witnesses escaped every such fixed subcylinder, the first
label would be exterior/separated concentration, cascade, or compactness
failure, all of which are non-thin rows already listed in
\cref{paper6a:lem:terminal-unit-drift-window-admissibility}\,\textup{(A1)}.
If the active CKN diagnostic were carried only by pressure while the velocity
floor above failed, then
\cref{paper8:lem:harmonic-cross-pressure-routing,paper8:lem:nested-ball-harmonic-tail}
routes the branch to the pressure/exterior row, not to the retained thin-drift
row.  Thus the displayed active spacetime velocity floor holds on the retained
branch.
\par
The hypotheses of the minimal active-family extraction and of the
same-point/separated reductions are exactly the repaired-gauge representation,
local pressure reconstruction, compact-cylinder suitable bounds, and retained
active local mass floor now available on \(Q_{r_*}\).  Hence those reductions
may be applied by the terminal active-frame machinery.
\end{proof}

\begin{lemma}[Vanishing negative scale forcing has no terminal label]
\label{paper6a:lem:vanishing-negative-drift-forcing-absorbed}
Let \(I_m=[\sigma_m,\sigma_m+1]\) be admissible selected terminal windows and
assume that on every compact \(K\Subset\mathbb R^3\)
\[
  a_-(\sigma_m+\cdot)
  \bigl(V_m+y\cdot\nabla V_m\bigr)
  \longrightarrow0
  \quad\text{in }L^1(0,1;H^{-1}(K)).
\]
Assume also the corresponding local energy pairings with compact cutoffs tend
to zero.  Then the negative scale-drift term contributes no nonzero local
defect measure and no independent terminal state-space label.  Any retained
limit or occupation state obtained from the windows is the same local
Navier--Stokes state that would be obtained after moving this term to a
vanishing \(L^1H^{-1}_{\rm loc}\) forcing; the ordered labels are therefore the
ordinary active-core, compactness, exterior, multicore/cascade, scale-rigid,
or bounded-window labels.
\end{lemma}

\begin{proof}
Let \(\varphi\in C_c^\infty((0,1)\times K;\mathbb R^3)\).  The assumed
\(L^1H^{-1}\) convergence gives
\[
  \int_0^1
  \left\langle
  a_-(\sigma_m+s)(V_m+y\cdot\nabla V_m),\varphi(s)
  \right\rangle ds
  \longrightarrow0 .
\]
Hence the negative scale term disappears in the distributional limit of the
represented equation.  The compact-cutoff energy pairings tend to zero by
assumption, so the local energy inequality also passes to the same limit with
no additional defect measure.  The pressure reconstruction is unchanged: since
\(\nabla\cdot V_m=0\), the field \(V_m+y\cdot\nabla V_m\) is divergence-free
in the spatial variables, so this lower-order velocity forcing contributes no
new pressure source.
\par
Thus, in the augmented local compactness topology, the extra coordinate
recording the negative scale forcing converges to the zero face.  Closedness of
the ordered local diagnostics
\cref{paper7:lem:closed-carrier-diagnostics,paper6a:lem:critical-mass-decoupling}
then implies that this vanishing coordinate cannot be the first nonzero label
of a terminal state.  The only possible first labels are the ordinary local
state-space labels listed in the statement.
\par
This is a zero-face statement for the augmented local state space, not a
small-data statement for a velocity remainder.  The proof uses only
distributional closedness, compact-cutoff local-energy closedness, and the
closedness of the already defined local diagnostics.  No global Kato theory and
no small \(L^\infty L^3\) velocity remainder is invoked.
\end{proof}

\begin{lemma}[Diffuse thin drift routes to another local row]
\label{paper6a:lem:diffuse-thin-drift-routes}
Let a physical local Type II branch be in the canonical windowwise
pre-carrier regime and suppose that, after the finite-cost and finite
absolute-cost alternatives have been removed, it satisfies the scale-collapse
drift obstruction and the thin-drift alternative.  Let
\([s_m,t_m]\) be the diffuse blocks supplied by
\cref{paper6a:lem:thin-drift-diffuse-blocks}.  Then exactly one of the
following occurs after passing to a subsequence:
\begin{enumerate}[label=\textup{(T\arabic*)}]
\item one of the non-thin local first-failure rows occurs: autonomous
      modulation failure, modulation-driven recentering, compactness-package
      first failure routed by
      \cref{paper5:thm:selected-compactness-failure-routed}, exterior or
      separated-core concentration, rough-core loss, multicore or same-point
      cascade, carrier/cost validation failure outside the canonical retained
      channel, or loss of the retained active core;
\item a closed retained row occurs: compact single-core, scale-rigid residual
      behavior, retained compact scale-collapse, or a nonzero bounded
      selected-window limit;
\item the negative scale-drift term is perturbative on the canonical selected
      compact unit windows and no additional terminal thin-drift state remains.
\end{enumerate}
\end{lemma}

\begin{proof}
Because the blocks are diffuse,
\[
  \varepsilon_m:=
  \sup_{T\in[s_m,t_m]}
  \int_{[T,T+1]\cap[s_m,t_m]}g_{\rm sc}(\tau)\,d\tau
  \longrightarrow0 .
\]
The block lengths tend to infinity.  Indeed, if a subsequence of blocks had
length bounded by \(L\), each such block would be covered by at most
\(\lceil L\rceil+2\) unit intervals, and its total \(g_{\rm sc}\)-mass would be
bounded by \((\lceil L\rceil+2)\varepsilon_m\to0\), contradicting the
normalization \(\int_{s_m}^{t_m}g_{\rm sc}=1\).  Hence, after discarding a
finite prefix, each block contains a canonical unit cell
\[
  K_{k_m}=[\tau_\sharp+k_m,\tau_\sharp+k_m+1]\subset[s_m,t_m],
  \qquad k_m\to\infty .
\]
The same estimate holds for every canonical cell contained in the block.  In
particular, since \(K_{k_m}\subset[s_m,t_m]\), diffuseness gives
\[
  \int_{K_{k_m}}g_{\rm sc}\le\varepsilon_m .
\]
\par
Apply
\cref{paper6a:lem:terminal-unit-drift-window-admissibility} to the canonical
cells \(K_{k_m}\).  If alternative \textup{(A1)} occurs, the branch is in
\textup{(T1)}.  Hence assume that alternative \textup{(A2)} holds: after
recentring, \(K_{k_m}\) are admissible selected terminal windows with the retained
local representation, pressure, energy, compactness, modulation, and moving
core-floor data.  Write \(I_m:=K_{k_m}\) for these canonical selected windows.
\par
On these admissible windows, the moving \(L^2\) core floor and the retained
local \(L^\infty_tL^2_x\) bounds are part of
\cref{paper6a:lem:terminal-unit-drift-window-admissibility}\,\textup{(A2)}.
Hence
\cref{paper6a:lem:small-negative-drift-perturbative} implies
\[
  \|a_-(V+y\cdot\nabla V)\|_{L^1(I_m;H^{-1}(K))}
  \longrightarrow0
\]
for every compact \(K\Subset\mathbb R^3\), and its compact-cutoff local energy
pairings also vanish.  By
\cref{paper6a:lem:vanishing-negative-drift-forcing-absorbed}, the diffuse
negative drift is a vanishing local forcing on the selected unit windows; it is
not a new compact state-space object and cannot be a terminal first label.
\par
It remains to classify the retained compact sequence on these same admissible
windows.  Apply
\cref{paper6a:lem:admissible-drift-window-active-core}.  If the retained
spacetime activity floor is lost, the branch is in \textup{(T1)}.  Otherwise a
fixed compact subcylinder carries a positive retained local spacetime mass
floor, so the
structural reductions are applicable.
Apply the minimal active-family extraction and the same-point/separated-point
reductions given by
\cref{paper3:thm:minimal-bad-extraction,paper3:thm:same-point-reduction,paper3:thm:separated-point-reduction}.
If those reductions produce another active core, a strict same-point cascade,
a separated physical point, exterior pressure persistence, rough-core loss, or
loss of compact estimates, then the branch is in \textup{(T1)}.  If they
produce the compact single-core row, the retained compact scale-collapse row,
or the scale-rigid residual row, those rows are the closed alternatives in
\textup{(T2)}, closed respectively by
\cref{paper6:thm:single-core,paper5:thm:compact-scale-collapse-rigidity-discharged,paper3:prop:scale-rigidity-discharged}.
In the scale-rigid case, if the compact limit has no further CKN
subconcentration, \cref{paper3:lem:scale-rigid-no-subconcentration-bounded}
gives a nonzero bounded selected-window limit; if it has a further
subconcentration, the same lemma routes it to the active structural rows
already listed in \textup{(T1)}.  Finally,
\cref{paper3:lem:scale-rigid-bounded-limit-exit} excludes the bounded-limit
case from the local Type II class.
\par
Thus, after all ordered first failures and closed retained rows have been
accounted for, the only remaining effect of thin negative drift on the
selected unit windows is the vanishing \(H^{-1}_{\rm loc}\) perturbative term
proved above.  By
\cref{paper6a:lem:vanishing-negative-drift-forcing-absorbed}, such a vanishing
forcing carries no independent terminal state-space label.  This proves
\textup{(T3)} and the asserted trichotomy.
\end{proof}

\begin{theorem}[Thin negative drift is not a terminal exit]
\label{paper6a:thm:thin-negative-drift-closed}
On every physical local Type II branch in the canonical windowwise
pre-carrier regime, the thin negative-drift alternative is not an independent
remaining terminal exit.  More precisely, if the branch enters the
scale-collapse drift obstruction and does not have persistent fixed-length
negative-drift windows, then the branch either enters one of the other named
local exit rows, or it enters an already closed retained row.  Consequently
the row ``thin drift or negative scale-drift without thick autonomous
windows'' may be removed from
\cref{paper6:def:remaining-named-exits}.
\end{theorem}

\begin{proof}
The scale-collapse drift obstruction gives the infinite tail integral of
\(g_{\rm sc}\).  By
\cref{paper6a:lem:thin-drift-diffuse-blocks}, every terminal tail has either
persistent unit-window weighted negative drift or diffuse normalized drift
blocks.  In the first alternative, finite overlap with the canonical schedule
of \cref{paper6a:def:canonical-thin-drift-schedule} gives canonical cells
carrying a uniform positive weighted negative-drift lower bound, unless one of
the non-thin first-failure rows in
\cref{paper6a:lem:terminal-unit-drift-window-admissibility}\textup{(A1)} has
already occurred.  If such a first failure occurs, the branch has left through
another named row.  Otherwise those canonical cells are admissible selected
terminal windows, and the branch is in the thick-window ledger of
\cref{paper5:thm:scale-collapse-stratification}: selected-window compactness
can fail, autonomous modulation convergence can fail, or the retained thick
compact scale-collapse state is reached.  The first two are other named rows,
and the retained thick state is closed by
\cref{paper5:thm:compact-scale-collapse-rigidity-discharged}.  No part of this
persistent-window branch validates arbitrary intervals or uses the diffuse
thin-drift closure itself.
\par
On a genuine thin-drift tail, only the diffuse block alternative remains.
\Cref{paper6a:lem:diffuse-thin-drift-routes} then gives three possibilities.
The first is one of the other named local exit rows.  The second is a closed
retained row, including the bounded selected-window exit excluded by
\cref{paper3:lem:scale-rigid-bounded-limit-exit}.  The third says that the
negative drift is a vanishing local \(H^{-1}\) perturbation on the selected
compact windows and therefore does not define a terminal state-space row after
the retained perturbative remainder has been formed.
\par
Thus no local Type II branch can remain whose first and only unresolved label
is thin negative drift.  The argument used only finite-window local estimates,
occupation blocks, the retained spacetime activity floor, and the ordered local
state-space partition; it did not use a global \(L^3\) bound, a whole-space
profile decomposition, or a hidden conversion of diffuse drift into thick
windows.
\end{proof}

\subsection{Closure of the Type II analytic-data exhaustiveness row}
\label{paper6a:subsec:discharge-typeII-analytic-data}

The exhaustion row behind \cref{thm:c1-typeII-branch-exhaustion} is the entry
condition for every admissible local Type II branch.  It is closed by an
ordered local extraction package: each component of
\cref{def:c1-typeII-analytic-data} is either
supplied by an existing local lemma on the admissible class
\(\mathcal U_{\mathrm{II}}^{NS}\), or its first failure is routed to a named
alternative that is class-incompatible or already covered by the local
state-space stratification.  No global critical profile theorem and no
global \(L^3\) bound are invoked.

\begin{lemma}[Concentration extraction routes]
\label{paper6a:lem:typeII-concentration-extraction-routes}
Let \(\omega=(u,p,T^*)\in\mathcal U_{\mathrm{II}}^{NS}\).  Exactly one of the
following holds:
\begin{enumerate}[label=\textup{(\alph*)}]
\item the local concentration dichotomy
\cref{paper1:thm:paper0-dichotomy} produces \(x_0\in\Sigma(T^*)\),
\(\eta>0\), and \(r_n\downarrow0\) with
\(C((x_0,T^*),r_n)+D((x_0,T^*),r_n)\ge\eta\) for all \(n\), and the
concentration point lies in the local Type II alternative;
\item the singular set \(\Sigma(T^*)\) is empty;
\item \(\Sigma(T^*)\ne\emptyset\), and every concentration point of
\(\Sigma(T^*)\) lies in the local Type I alternative.
\end{enumerate}
Cases \textup{(b)} and \textup{(c)} are class-incompatible with
\cref{def:c1-admissible-typeII-class}.  Hence on
\(\mathcal U_{\mathrm{II}}^{NS}\) only case \textup{(a)} occurs, and the
first ordered failure of \cref{def:c1-exhaustion-failures}
(``concentration extraction fails'') does not arise on the admissible class.
\end{lemma}

\begin{proof}
If \(\Sigma(T^*)=\emptyset\), case \textup{(b)} holds.  If
\(\Sigma(T^*)\ne\emptyset\), apply the local concentration dichotomy at
singular points.  If some singular point is not in the local Type I
alternative, choose that point as \(x_0\); the Caffarelli--Kohn--Nirenberg
failure at \(x_0\) supplies \(\eta>0\) and \(r_n\downarrow0\) with
\[
  C((x_0,T^*),r_n)+D((x_0,T^*),r_n)\ge\eta,
\]
and case \textup{(a)} holds.  If no such point exists, every concentration
point lies in the local Type I alternative and case \textup{(c)} holds.  The
three cases are therefore disjoint and exhaustive.

In case \textup{(b)} there is no terminal singular concentration point, so the
branch is the local continuation/no-terminal-concentration alternative rather
than a local Type II branch.  This is excluded by clause \textup{(iv)} of
\cref{def:c1-admissible-typeII-class}; no global continuation estimate is used.
In case \textup{(c)} every concentration point lies in the local Type I
alternative of \cref{paper1:thm:paper0-dichotomy}, which is exactly the local
Type I scale-invariant criterion.  This contradicts clause \textup{(iii)} of
\cref{def:c1-admissible-typeII-class}.  In case \textup{(a)} the dichotomy
supplies the positive CKN concentration scales and the exclusion of the Type I
alternative at \((x_0,T^*)\); this is the content of the first analytic datum of
\cref{def:c1-typeII-analytic-data}.
\end{proof}

\begin{lemma}[Supercritical scale alternative routes]
\label{paper6a:lem:typeII-scale-alternative-routes}
Let \(\omega\in\mathcal U_{\mathrm{II}}^{NS}\) and let \((x_0,T^*)\) be a
concentration point produced by
\cref{paper6a:lem:typeII-concentration-extraction-routes}\,\textup{(a)}.
Then the supercritical scale alternative of
\cref{def:c1-typeII-analytic-data}\,\textup{(ii)} holds at \((x_0,T^*)\).
\end{lemma}

\begin{proof}
The dichotomy of \cref{paper1:thm:paper0-dichotomy} at \((x_0,T^*)\) is
exclusive: either a local Type I scale-invariant bound holds, or no such
local Type I bound holds.  The first option contradicts clause \textup{(iii)}
of \cref{def:c1-admissible-typeII-class}, which excludes branches in the
local Type I class.  Hence the second option holds at every concentration
point produced by the previous lemma; this is exactly the supercritical
Type II scale alternative.  The second ordered failure of
\cref{def:c1-exhaustion-failures} therefore does not arise on
\(\mathcal U_{\mathrm{II}}^{NS}\).
\end{proof}

\begin{lemma}[Local profile data routes]
\label{paper6a:lem:typeII-profile-data-routes}
Let \(\omega\in\mathcal U_{\mathrm{II}}^{NS}\) and let \((x_0,T^*)\) be a
concentration point in the local Type II alternative as produced by
\cref{paper6a:lem:typeII-concentration-extraction-routes,paper6a:lem:typeII-scale-alternative-routes}.
Run the local scale-translation selection of
\cref{paper1:prop:paper0-typeII-reduction} at the positive concentration
scales and test whether it produces a nondegenerate compact single core.  Then
exactly one of the following holds:
\begin{enumerate}[label=\textup{(\alph*)}]
\item the localized scale and center are nondegenerate and no represented
selected-window subsequence has a nonzero bounded selected-window limit on
any compact cylinder; \cref{paper1:prop:representation-input} supplies
represented variables \((V_n,P_n,a_n,b_n)\) on compact cylinders satisfying
the represented Navier--Stokes equation
\eqref{paper1:eq:represented-ns}; the represented selected windows form a
local Type II sequence in the sense of \cref{paper1:def:typeII}, providing
the local profile data of
\cref{def:c1-typeII-analytic-data}\,\textup{(iii)};
\item the localized nondegeneracy fails, and the branch falls into the
multibubble, cascade, or gauge-degenerate alternative of
\cref{paper3:def:local-multibubble}\,\textup{(ii)}--\textup{(v)};
\item nondegeneracy holds but the represented selected windows admit a
nonzero bounded selected-window limit on some compact cylinder; the branch
then enters the bounded-selected-window alternative of
\cref{paper1:def:bounded-window-limit} and leaves the strict selected-window
branch by \cref{paper3:lem:scale-rigid-bounded-limit-exit}, after which the
physical bounded-window event is routed by
\cref{paper6:thm:bounded-selected-window-physical-routing}.
\end{enumerate}
In cases \textup{(b)} and \textup{(c)} the branch is non-retained: case
\textup{(b)} is treated by the multibubble/cascade machinery of
\cref{sec:multibubble-cascade-exclusion}, and case \textup{(c)} is the
bounded selected-window exit of the scale-rigid stratum.  Neither case is
a compact single-core branch.
\end{lemma}

\begin{proof}
The local scale-translation selection at \((x_0,T^*)\) attempts to extract
a single compact core at the positive CKN concentration scales
\(r_n\downarrow0\).  Either the localized scale and center are
nondegenerate (cases \textup{(a)} and \textup{(c)}), or they are
degenerate (case \textup{(b)}); under nondegeneracy, either no represented
selected-window subsequence has a nonzero bounded selected-window limit
(case \textup{(a)}), or some such subsequence does (case \textup{(c)}).
The three cases are exhaustive.

In case \textup{(a)},
\cref{paper1:prop:paper0-typeII-reduction} applies and rewrites the
shrinking cylinders \(Q_{r_n}(x_0,T^*)\) as fixed selected windows.  The
represented sequence has positive local CKN concentration on a fixed
compact cylinder; combined with the defining case-\textup{(a)} absence of
nonzero bounded selected-window limits, it is a local Type II sequence in the
sense of
\cref{paper1:def:typeII}.  \Cref{paper1:prop:representation-input} then
supplies the represented variables \((V_n,P_n,a_n,b_n)\) on compact
cylinders satisfying \eqref{paper1:eq:represented-ns}; this is the local
profile datum of \cref{def:c1-typeII-analytic-data}\,\textup{(iii)}.

In case \textup{(b)}, failure of localized nondegeneracy is, by the second
sentence of \cref{paper1:prop:representation-input}, the multibubble,
cascade, or gauge-degenerate alternative.  The corresponding stratum is
named in \cref{paper3:def:local-multibubble}\,\textup{(ii)}--\textup{(v)}
and is non-retained: it is handled by the multibubble/cascade machinery
and never enters the retained compact single-core stratum produced by
\cref{paper7:lem:no-unrecorded-cost-escape}.

In case \textup{(c)}, the represented selected-window sequence has a
nonzero bounded selected-window limit on some compact cylinder, placing
the branch in \cref{paper1:def:bounded-window-limit};
\cref{paper3:lem:scale-rigid-bounded-limit-exit} then routes the branch
out of Type II via the scale-rigid bounded-window exit, which is also
non-retained.

Hence the third ordered failure of \cref{def:c1-exhaustion-failures}
(``local profile data are insufficient'') is routed to the
multibubble/cascade/gauge-degenerate non-retained alternative or to the
scale-rigid bounded-window exit; on the retained compact single-core
stratum it does not arise.
\end{proof}

\begin{theorem}[Type II analytic-data routing on the admissible class]
\label{paper6a:thm:typeII-analytic-data-exhaustive}
The Type II analytic data of \cref{def:c1-typeII-analytic-data} are
exhaustive by ordered local routing on
\(\mathcal U_{\mathrm{II}}^{NS}\): every
\(\omega\in\mathcal U_{\mathrm{II}}^{NS}\) either has the concentration
profile, the supercritical scale alternative, and the local profile data, or
its first applicable failure in
\cref{def:c1-exhaustion-failures}\,\textup{(i)}--\textup{(iv)} is matched to a
named alternative that is class-incompatible with
\cref{def:c1-admissible-typeII-class} or to a non-retained state-space branch
already covered by the local stratification.  In particular, after those named
exits are removed, every retained admissible branch has the three Type II
analytic data.
\end{theorem}

\begin{proof}
Process the three analytic data of \cref{def:c1-typeII-analytic-data} in the
order specified by \cref{def:c1-exhaustion-failures}.

\emph{Datum (i): concentration extraction.}
\Cref{paper6a:lem:typeII-concentration-extraction-routes} shows that on
\(\mathcal U_{\mathrm{II}}^{NS}\) the only realised case is the dichotomy
case \textup{(a)} producing \((x_0,T^*)\) with positive CKN concentration in
the local Type II alternative; the other two cases are class-incompatible.

\emph{Datum (ii): supercritical scale alternative.}
\Cref{paper6a:lem:typeII-scale-alternative-routes} then forces the
supercritical Type II scale alternative at \((x_0,T^*)\); the Type I
alternative is class-incompatible by clause \textup{(iii)} of
\cref{def:c1-admissible-typeII-class}.

\emph{Datum (iii): local profile data.}
\Cref{paper6a:lem:typeII-profile-data-routes} either supplies the local
profile data via the local scale-translation selection and the representation
input \cref{paper1:prop:representation-input}, or routes the branch to the
multibubble/cascade/gauge-degenerate non-retained alternative
\cref{paper3:def:local-multibubble}\,\textup{(ii)}--\textup{(v)}, or to the
bounded selected-window exit
\cref{paper3:lem:scale-rigid-bounded-limit-exit}.

\emph{Residual datum (iv).}
The fourth ordered failure of \cref{def:c1-exhaustion-failures}, ``the
admissible Type II branch lies outside the concentration, scale, and
profile alternatives under consideration'', is the catch-all for a branch that
is not one of the local Navier--Stokes Type II alternatives after the first
three tests have been applied.  Clause \textup{(iv)} of
\cref{def:c1-admissible-typeII-class} excludes continuation, Type I, boundary,
and non-Navier--Stokes-domain alternatives.  Any remaining local branch outside
the retained profile alternatives is one of the named non-retained state-space
exits produced in the preceding steps.

In every case the branch either has all three analytic data or its first
ordered failure is recorded as a class-incompatible alternative or as a
non-retained state-space branch already covered by the local stratification.
This is the ordered local-routing exhaustiveness statement associated with
\cref{cor:c1-ordered-exhaustion}.  Removing the named exits leaves only the
first alternative of that corollary, namely the full Type II analytic data.
\end{proof}

\begin{corollary}[Type II analytic-data exhaustiveness discharged]
\label{paper6a:cor:typeII-analytic-data-exhaustiveness-discharged}
The ordered-routing exhaustion input used by
\cref{thm:c1-typeII-branch-exhaustion} is a theorem on
\(\mathcal U_{\mathrm{II}}^{NS}\): the analytic-data entry row is
discharged by ordered local extraction.  On retained admissible branches the
three Type II analytic data hold; non-retained failures are assigned to named
local alternatives before the downstream compact analysis is invoked.  No
global critical profile theorem and no global \(L^3\) bound are invoked.
\end{corollary}

\begin{proof}
\Cref{paper6a:thm:typeII-analytic-data-exhaustive} establishes ordered local
routing of all failures of the Type II analytic data on
\(\mathcal U_{\mathrm{II}}^{NS}\).  By
\cref{cor:c1-ordered-exhaustion}, after the named routed exits are removed, the
remaining retained branch is in the first alternative and therefore has the
Type II analytic data required by \cref{thm:c1-typeII-branch-exhaustion}.
Each step uses only the local concentration dichotomy
\cref{paper1:thm:paper0-dichotomy}, the local scale-translation selection
\cref{paper1:prop:paper0-typeII-reduction}, the local representation input
\cref{paper1:prop:representation-input}, the local Type II definition
\cref{paper1:def:typeII}, the admissibility clauses of
\cref{def:c1-admissible-typeII-class}, and the multibubble/cascade
stratum \cref{paper3:def:local-multibubble}, together with the bounded-window
exit \cref{paper3:lem:scale-rigid-bounded-limit-exit}.
\end{proof}

\begin{definition}[Terminal local compactness input hypotheses]
\label{paper6a:def:terminal-critical-profile-input-hypotheses}
The terminal local compactness \emph{input} hypotheses consist of the local
compact terminal-window hypothesis
\cref{paper6a:def:bounded-critical-terminal-sequences} and local
perturbative stability
\cref{paper6a:def:small-data-critical-stability}, together with item
\textup{(i)} of
\cref{paper6a:def:critical-ns-profile-decomposition-local} (the admissibility
estimates on the represented branch).  These hypotheses do \emph{not}
include the local terminal completeness conclusion
\textup{(iii)} of
\cref{paper6a:def:critical-ns-profile-decomposition-local}.
\end{definition}

\begin{definition}[Terminal local compactness package]
\label{paper6a:def:terminal-critical-profile-hypotheses}
The terminal local compactness package is the input hypotheses
\cref{paper6a:def:terminal-critical-profile-input-hypotheses} together with the
local terminal completeness theorem
\cref{paper6a:thm:critical-ns-profile-decomposition-discharged}.  This is a
record of what has been proved and is not an additional assumption.
\end{definition}

\begin{theorem}[Terminal critical local compactness]
\label{paper6a:thm:terminal-critical-local-compactness}
Assume the terminal local compactness input hypotheses of
\cref{paper6a:def:terminal-critical-profile-input-hypotheses}.  Assume also
that the ordered adjacent local labels are routed as in
\cref{paper6a:thm:terminal-adjacent-label-routing} and that comparable retained
classes are grouped before the terminal frame is fixed.  Then, on every
retained terminal branch and every compact terminal cylinder, all nonretained
residual components vanish in the local perturbative topology of
\cref{paper8:def:terminal-nonlinear-decoupling}.
\end{theorem}

\begin{proof}
The local compact terminal-window hypothesis gives the compact-cylinder
estimates.  Applying
\cref{paper6a:thm:critical-ns-profile-decomposition-discharged} under the
stated routing and grouping gives
\[
  S_n\to0\quad\text{in }L^\infty(I;L^3(B_R))
\]
on each compact terminal cylinder.  Local perturbative stability then converts
this velocity convergence into the nonlinear-stress, pressure, and residual
convergences required in terminal nonlinear decoupling.
\end{proof}

\begin{corollary}[Terminal completeness from local compactness hypotheses]
\label{paper6a:cor:profile-completeness-from-critical-hypotheses}
The terminal local completeness statement in the final argument may be
replaced by the terminal local compactness package of
\cref{paper6a:def:terminal-critical-profile-hypotheses}.
\end{corollary}

\begin{proof}
\Cref{paper6a:thm:terminal-critical-local-compactness} gives the asserted
terminal local completeness.
\end{proof}

\begin{remark}[Boundary of the terminal compactness hypothesis]
\label{paper6a:rem:terminal-compactness-boundary}
The preceding compactness conclusion applies only to terminal windows on which
the repaired-gauge representation, pressure reconstruction, local Caccioppoli
control, bounded modulation, and local state-space labels are available.  If
one of these inputs is absent, the branch is assigned to the corresponding
representation, pressure, rough-core, modulation, or state-space alternative;
it is not treated by importing a global profile theorem.
\end{remark}

\begin{definition}[Terminal sequences sampled from the renormalized orbit]
\label{paper6a:def:terminal-sequences-sampled-from-orbit}
A terminal active sequence is sampled from the renormalized orbit if it is of
the form
\[
  u_{0,n}=\Lambda_{\lambda_n,x_n}V(\tau_n),
  \qquad
  \tau_n\to\infty,\qquad \lambda_n>0,\qquad x_n\in\mathbb R^3,
\]
or by the equivalent inverse critical change of variables, where
\[
  (\Lambda_{\lambda,x_0}f)(x)
  =
  \lambda^{-1}f\left(\frac{x-x_0}{\lambda}\right).
\]
\end{definition}

\begin{lemma}[Critical symmetry preserves the \(L^3\) norm]
\label{paper6a:lem:critical-symmetry-preserves-l3}
For every \(f\in L^3(\mathbb R^3)\), \(\lambda>0\), and
\(x_0\in\mathbb R^3\),
\[
  \|\Lambda_{\lambda,x_0}f\|_{L^3}=\|f\|_{L^3}.
\]
\end{lemma}

\begin{proof}
By the change of variables \(x=x_0+\lambda y\),
\[
  \|\Lambda_{\lambda,x_0}f\|_{L^3}^3
  =
  \int_{\mathbb R^3}\lambda^{-3}
  \left|f\left(\frac{x-x_0}{\lambda}\right)\right|^3\,dx
  =
  \int_{\mathbb R^3}|f(y)|^3\,dy .
\]
Taking cube roots gives the claim.
\end{proof}

\begin{theorem}[Bounded terminal sequences from local critical bounds]
\label{paper6a:thm:bounded-terminal-sequences-from-l3}
Assume the renormalized orbit satisfies the local compact-window bounds
\[
  0<\inf_{\tau\in I}\|V(\tau)\|_{L^3(B_R)}
  \le
  \sup_{\tau\in I}\|V(\tau)\|_{L^3(B_{2R})}
  <\infty,
\]
on the terminal interval under consideration, together with the local
Caccioppoli and pressure bounds on \(Q_{2R,I}\).  If every terminal active
window is sampled from the renormalized orbit, then every terminal active
window satisfies the local compact terminal-window hypothesis.
\end{theorem}

\begin{proof}
The local \(L^3\) upper bound is one component of the compact-window
hypothesis.  Local \(L^\infty L^2\), \(L^2H^1\), and pressure bounds are supplied
by the local Caccioppoli estimate and pressure reconstruction on \(Q_{2R,I}\).
The repaired-gauge equation gives the required time-derivative bound in
\(H^{-m}\).  The lower bound records that the retained window is nontrivial.
\end{proof}

\begin{theorem}[Terminal compactness alternatives]
\label{paper6a:thm:terminal-boundedness-alternatives}
For a candidate terminal active sequence in the Type II analysis, the following
alternatives exhaust the possible local compactness outcomes, after passing to a
subsequence when necessary:
\begin{enumerate}[label=\textup{(\roman*)}]
\item the sequence is not sampled from the renormalized orbit;
\item the repaired-gauge representation, pressure reconstruction, or local
Caccioppoli estimate is unavailable on the selected compact cylinder;
\item the local \(L^3\) upper bound fails along a terminal sampling subsequence;
\item the retained local \(L^3\) lower bound vanishes along a terminal sampling
subsequence;
\item after excluding the preceding alternatives, the terminal active sequence
satisfies the local compact terminal-window hypothesis.
\end{enumerate}
\end{theorem}

\begin{proof}
If the terminal sequence is not sampled from the renormalized orbit, then
\textup{(i)} holds.  If one of the local PDE inputs is missing, then
\textup{(ii)} holds and the branch is assigned to the corresponding local
alternative.  If the local critical upper bound fails, then \textup{(iii)}
holds after passing to a subsequence.  If the retained lower bound vanishes,
then \textup{(iv)} holds.  Otherwise the local hypotheses of
\cref{paper6a:thm:bounded-terminal-sequences-from-l3} hold, and
\textup{(v)} follows.
\end{proof}

\begin{definition}[Terminal coordinate-frame construction]
\label{paper6a:def:terminal-coordinate-frame-construction}
The terminal coordinate-frame construction means that every terminal active
frame used in the analysis is obtained from the repaired-gauge orbit as
follows.  Choose sampling times \(\tau_n\to\infty\) and critical symmetry
parameters \((\lambda_n,x_n)\); form the local terminal windows by applying the
critical coordinate-frame map to \(V\) on compact time intervals around
\(\tau_n\); group comparable retained local classes based at the same physical
point before selecting the terminal active frame.  No external data are
introduced.
\end{definition}

\begin{theorem}[Terminal coordinate-frame construction from local state selection]
\label{paper6a:thm:terminal-coordinate-frame-construction-profile-selection}
Assume terminal active frames are selected by first partitioning local active
states by physical concentration point, then grouping comparable retained
states at the same point into compound local states, and finally choosing a
minimal-scale retained class at that point.  Then the terminal coordinate-frame
construction of \cref{paper6a:def:terminal-coordinate-frame-construction} holds
for the selected terminal frames.
\end{theorem}

\begin{proof}
By the stated selection rule, each terminal active frame is attached to a
compound active class after all comparable retained states at the same physical
point have been grouped.  The scale and center of the terminal frame are the
scale and center of this compound class.  The terminal data are then obtained
by applying the Navier--Stokes critical change of variables to the local
renormalized branch along the chosen sampling sequence \(\tau_n\to\infty\).
Thus the construction uses only the repaired-gauge orbit and the
Navier--Stokes critical symmetries, and it satisfies all clauses of
\cref{paper6a:def:terminal-coordinate-frame-construction}.
\end{proof}

\begin{theorem}[Terminal coordinate frames are sampled from the orbit]
\label{paper6a:thm:terminal-coordinate-frames-sampled}
Assume the terminal coordinate-frame construction.  Then every terminal active
sequence belongs to the class described in
\cref{paper6a:def:terminal-sequences-sampled-from-orbit}.
\end{theorem}

\begin{proof}
By definition of the terminal coordinate-frame construction, each terminal
active frame is obtained by choosing times \(\tau_n\), selecting scale and
center parameters \((\lambda_n,x_n)\), and applying the Navier--Stokes critical
change of variables to the renormalized solution \(V(\tau_n)\).  Hence the
associated terminal data have the form
\[
  u_{0,n}=\Lambda_{\lambda_n,x_n}V(\tau_n),
\]
or the equivalent inverse formulation, which is the orbit-sampling
condition.
\end{proof}

\begin{theorem}[Repaired-gauge representation gives terminal orbit sampling]
\label{paper6a:thm:representation-gives-terminal-orbit-sampling}
Assume a repaired-gauge representation of the Type II branch and the terminal
coordinate-frame construction.  Then every terminal active sequence is sampled
from the renormalized orbit.
\end{theorem}

\begin{proof}
The repaired-gauge representation supplies the renormalized orbit \(V(\tau)\)
and the critical coordinate maps between physical variables and renormalized
variables.  The terminal coordinate-frame construction supplies the terminal
sampling times and scale-center parameters.  Combining these hypotheses gives
\[
  u_{0,n}=\Lambda_{\lambda_n,x_n}V(\tau_n),
\]
which is the assertion of
\cref{paper6a:def:terminal-sequences-sampled-from-orbit}.
\end{proof}

\begin{corollary}[Failure of terminal orbit sampling]
\label{paper6a:cor:failure-terminal-orbit-sampling}
If a sequence reaches the terminal local-state stage but is not sampled from the
renormalized orbit, then either the repaired-gauge representation is not
available for that sequence or the terminal coordinate-frame construction has
not been applied as specified.
\end{corollary}

\begin{proof}
This is the contrapositive of
\cref{paper6a:thm:representation-gives-terminal-orbit-sampling}.
\end{proof}

\subsection{Terminal Retained-State Assembly}

\begin{definition}[Abstract terminal elimination data]\label{paper6a:def:abstract-final-data}
The abstract terminal elimination data for this section are:
\begin{enumerate}[label=\textup{(\roman*)}]
\item every renormalized local Type II branch enters one of the local
alternatives used in \cref{paper6:thm:state-decomposition};
\item compact single-core vanishing-dissipation branches are excluded by
\cref{paper6:thm:single-core};
\item multibubble, cascade, and gauge-degenerate branches are excluded by
\cref{paper6:thm:multibubble}, and scale-rigid terminal branches are excluded by
\cref{paper3:prop:scale-rigidity-discharged};
\item finite-cost scale-collapse branches are excluded by
\cref{paper6:thm:scale-cost}, while compact retained single-core branches with
divergent localized cost are covered by
\cref{thm:c3-local-retained-mass-cost-divergence} and the cost-divergence
theorem
\cref{paper6a:thm:cost-divergence-exclusion-discharged,paper6a:cor:cost-divergence-exclusion-unconditional};
scale-collapse branches with divergent negative scale-drift cost are covered by
\cref{paper6a:thm:scale-collapse-drift-cost} or
\cref{paper6a:cor:absolute-scale-cost-criterion} and then excluded by the same
cost-divergence theorem;
\item failure of the windowed critical-tightness test on selected terminal
windows is excluded, equivalently that test holds for every renormalized
branch in the class;
\item failure of local windowed \(H^1\) control is excluded, equivalently the required local
windowed \(H^1\) control holds for every renormalized branch in the class;
\end{enumerate}
\end{definition}

\begin{theorem}[Abstract terminal Type II exclusion]
\label{paper6a:thm:abstract-final-typeII}
Under the abstract terminal elimination data of
\cref{paper6a:def:abstract-final-data}, no renormalized local Type II branch in
the stated class remains.
\end{theorem}

\begin{proof}
Let a renormalized local Type II branch be given.  By item \textup{(i)}, it
lies in one of the alternatives in \cref{paper6:thm:state-decomposition}.
Item \textup{(ii)} excludes the vanishing-dissipation compact single-core
subcase, while item \textup{(iv)} excludes the retained compact single-core
subcase with divergent localized cost.  Item \textup{(iii)} excludes
multibubble, cascade, gauge-degenerate, and scale-rigid alternatives.  Item
\textup{(iv)} also excludes finite-cost scale collapse and the divergent
negative-drift scale-collapse branch whenever the corresponding cost is one of
the costs covered by the cost-divergence theorem
\cref{paper6a:thm:cost-divergence-exclusion-discharged} and its status
corollary \cref{paper6a:cor:cost-divergence-exclusion-unconditional}.  Items
\textup{(v)} and
\textup{(vi)} exclude failure of the windowed critical-tightness test and failure of
local windowed \(H^1\) control.  These cases exhaust the alternatives, so no
branch remains.
\end{proof}

\begin{definition}[Expanded terminal data]\label{paper6a:def:expanded-terminal-data}
The following analytic data imply the abstract terminal elimination data:
\begin{enumerate}[label=\textup{(\roman*)}]
\item the local alternative decomposition for renormalized Type II branches;
\item the compact single-core vanishing-dissipation, multibubble/cascade,
finite-cost scale-collapse, compact-cost-divergence, and scale-rigidity
exclusions from the preceding sections;
\item a critical-tightness estimate giving the windowed critical-tightness
condition on selected terminal windows;
\item an exhaustive terminal compactness partition and a local finite
critical-mass statement for retained active terminal strata;
\item local perturbative residual stability and terminal state-space
completeness on compact windows;
\item removal of physically exterior regular components;
\item the repaired-gauge representation theorem;
\item the local Caccioppoli/windowed regularity input;
\item the local compact scale-collapse routing theorem
\cref{paper5:thm:compact-scale-collapse-rigidity-discharged}.
\end{enumerate}
\end{definition}

\begin{lemma}[Terminal data imply critical tightness]
\label{paper6a:lem:terminal-critical-tightness}
The expanded terminal data imply the windowed critical-tightness condition of
\cref{def:c5-compact-tests}.
\end{lemma}

\begin{proof}
The repaired-gauge representation turns each renormalized orbit into the
terminal windows used in the local state-space analysis.  The terminal
compactness partition, finite critical-mass statement, local perturbative
stability, and terminal state-space completeness give the required terminal
completeness.  Exterior regular components are removed by the exterior-removal
datum in \cref{paper6a:def:expanded-terminal-data}.
Together with the multibubble and scale-collapse exclusions, this leaves only
possible loss of windowed critical tightness on selected terminal windows; the
critical-tightness estimate excludes that loss.
\end{proof}

\begin{lemma}[Terminal data imply windowed \(H^1\) control]
\label{paper6a:lem:terminal-windowed-h1}
The expanded terminal data imply local windowed \(H^1\) control.
\end{lemma}

\begin{proof}
The windowed \(H^1\) estimate requires a repaired-gauge branch, local pressure and
modulation control, compact-cylinder critical spacetime bounds, the local
\(L^\infty_\tau L^2_y\) energy control, and the Caccioppoli/windowed regularity
input.  These are among the expanded terminal
data.  Hence the required local windowed \(H^1\) control holds, and the
failure of local windowed \(H^1\) control is unavailable.
\end{proof}

\begin{lemma}[Expanded terminal data imply abstract terminal data]
\label{paper6a:lem:expanded-to-abstract}
The expanded terminal data imply the abstract terminal elimination data of
\cref{paper6a:def:abstract-final-data}.
\end{lemma}

\begin{proof}
Items \textup{(i)} and \textup{(ii)} of
\cref{paper6a:def:expanded-terminal-data} give the local alternative
decomposition and the compact vanishing-dissipation, multibubble,
scale-collapse, compact-cost-divergence, and scale-rigidity exclusions.  The
cost-divergence criterion used in the divergent retained-core subcase is
supplied by the cost-divergence theorem
\cref{paper6a:thm:cost-divergence-exclusion-discharged} and its status
corollary \cref{paper6a:cor:cost-divergence-exclusion-unconditional}.  By
\cref{paper6a:lem:terminal-critical-tightness} they also give critical
tightness, and by
\cref{paper6a:lem:terminal-windowed-h1} they give the local windowed \(H^1\)
control.  These are the data listed in
\cref{paper6a:def:abstract-final-data}.
\end{proof}

\begin{theorem}[Expanded terminal Type II exclusion]
\label{paper6a:thm:expanded-terminal-typeII}
If the expanded terminal data hold, then no renormalized local Type II
branch in the stated class remains.
\end{theorem}

\begin{proof}
\Cref{paper6a:lem:expanded-to-abstract} gives the abstract terminal elimination
data.  Applying \cref{paper6a:thm:abstract-final-typeII} gives the
exclusion.
\end{proof}

\begin{remark}[Branchwise reading]
Compact single-core branches are split by
\cref{paper6:lem:compact-core-dissipation-dichotomy}: vanishing-dissipation
cores are removed by the compact criterion, and divergent retained-local-cost
cores are removed by
\cref{thm:c3-local-retained-mass-cost-divergence,paper6a:cor:cost-divergence-exclusion-unconditional}.
Retained scale-collapse branches are removed by the local compact
scale-collapse routing theorem
\cref{paper5:thm:compact-scale-collapse-rigidity-discharged}, or by
\cref{paper6a:thm:scale-collapse-drift-cost} when the scale-collapse drift cost
is compatible with the cost-divergence theorem; multibubble and non-tight cases
are removed by the terminal profile and critical-tightness inputs; failures of
local windowed \(H^1\) control are removed by the local Caccioppoli estimate;
and every divergent compatible localized cost on the canonical windowwise
schedule is handled by
\cref{paper6a:thm:cost-divergence-exclusion-discharged}.
\end{remark}

\begin{theorem}[Expanded terminal data hold branchwise]
\label{paper6a:thm:expanded-terminal-data-branchwise}
Let \(\omega\) be a physical local Type II branch after
\cref{paper6:cor:covered-class-entry-discharged}.  Then either \(\omega\)
exits through a named local or ambient alternative assigned to its specified
closure row in the acyclic local state-space order, or the nine items of
\cref{paper6a:def:expanded-terminal-data} hold on a retained terminal tail.
\end{theorem}

\begin{proof}
The proof is an ordered verification of the nine terminal inputs.
\par
\emph{Local decomposition and structural exclusions.}
Covered entry supplies a positive local Type II concentration sequence or an
earlier named exit.  Such an exit is not silently discarded; it is assigned to
the closure row specified by
\cref{paper7:cor:ambient-adjacent-wrapper-discharged,paper6a:lem:cost-divergence-closure-order}.
In the retained case
\cref{paper6:thm:state-decomposition} gives the local alternative
decomposition.  The compact single-core alternative is split by
\cref{paper6:lem:compact-core-dissipation-dichotomy}: the
vanishing-dissipation subcase is closed by
\cref{paper6:thm:single-core}, and the divergent retained-local-cost subcase is
closed, after the carrier and adjacent rows have been verified, by
\cref{thm:c3-local-retained-mass-cost-divergence,paper6a:cor:cost-divergence-exclusion-unconditional}.
The multibubble/cascade, finite-cost scale-collapse, retained compact
scale-collapse, and scale-rigid alternatives are closed by
\cref{paper6:thm:multibubble,paper6:thm:scale-cost,paper5:thm:compact-scale-collapse-rigidity-discharged,paper3:prop:scale-rigidity-discharged}.
\par
\emph{Critical tightness and compact terminal completeness.}
On the retained compact stratum, critical tightness and local windowed
\(H^1\)-control are theorems by
\cref{paper6a:thm:retained-compact-tests-closed}.  If either test fails before
the retained stratum is reached, the failure is an exterior, rough-core, or
multibubble/cascade state-space exit, routed by
\cref{paper8:thm:exterior-stratification,paper6:thm:rough-core,paper6:thm:multibubble}.
The finite critical-mass input used here is local on compact terminal windows:
\cref{paper6a:lem:terminal-local-critical-mass-routing} proves that failure of
positive finite local \(L^3\) mass is either non-retention or an adjacent active
label.  The proof does not require the whole-space critical norm
\(\|V(\tau)\|_{L^3(\mathbb R^3)}\).  With this local finite-mass input, the
terminal compactness partition is supplied by
\cref{paper6a:thm:terminal-boundedness-alternatives}, and terminal completeness
on the retained branch is supplied by
\cref{paper6a:thm:terminal-critical-local-compactness}.  In the retained
compact-cost path, the compactness actually used by the cost-divergence
argument is supplied from the local compact tests by
\cref{lem:c5-retained-data-supply-aubin-lions}.  The
adjacent labels in the partition are routed by
\cref{paper6a:thm:terminal-adjacent-label-routing}.  That theorem only removes
adjacent labels from the retained terminal component; the later exclusion of
carrier and cost adjacent rows is supplied in the assembly layer by the ambient
adjacent wrapper \cref{paper7:cor:ambient-adjacent-wrapper-discharged}.
\par
\emph{Residual stability, representation, exterior removal, and local
regularity.}
Terminal state-space completeness is
\cref{paper8:thm:terminal-windowed-profile-completeness}, proved locally from
the discharged critical profile decomposition
\cref{paper6a:thm:critical-ns-profile-decomposition-discharged}.  Local
perturbative residual stability is
\cref{paper6a:cor:kato-gives-small-data-stability}.  Physically exterior
regular components are removed by
\cref{paper8:thm:exterior-stratification}.  The repaired-gauge representation
is the ordered consequence of
\cref{cor:c2-ordered-representation,cor:c2r-ordered-classification} together
with the absolutely continuous chart consequences
\cref{paper6a:thm:ac-repaired-gauge-consequences}; if the chart or gauge row
fails, that failure is the representation alternative and not a retained
terminal branch.  The local Caccioppoli/windowed regularity input is supplied
by
\cref{paper6a:thm:ac-local-caccioppoli,paper6a:cor:no-independent-caccioppoli-obstruction-ac-suitable}
on suitable repaired-gauge branches, with retained windowed control supplied
by \cref{paper6a:thm:retained-compact-tests-closed}.  Finally, compact
scale-collapse routing is exactly
\cref{paper5:thm:compact-scale-collapse-rigidity-discharged}.
\par
Thus each item in \cref{paper6a:def:expanded-terminal-data} is either verified
on the retained terminal tail or its first failure is one of the named local or
ambient alternatives assigned to a closure row before terminal retained-state
elimination is invoked.  The dependency order is the one recorded in
\cref{paper6a:lem:cost-divergence-closure-order}: carrier and adjacent rows are
routed before cost-divergence closure, and terminal retained-state elimination
is last.  The verification uses only local compact windows, ordered state-space
labels, suitable local energy, and compact terminal estimates; it does not
import a global profile theorem, global Kato smallness, or a whole-space
\(L^3\) bound.
\end{proof}

\begin{corollary}[Expanded terminal wrapper discharged]
\label{paper6a:cor:expanded-terminal-wrapper-discharged}
The expanded terminal data are verified on the retained terminal path: every
physical local Type II branch either satisfies those data on a retained
terminal tail or exits through a named local or ambient
alternative assigned to the closure row specified by the acyclic local
state-space order.  This corollary verifies the retained terminal inputs; it
does not by itself exclude every named exit row.
\end{corollary}

\begin{proof}
This is \cref{paper6a:thm:expanded-terminal-data-branchwise}.
\end{proof}

\clearpage
\section{Assembly of the Local Type II Exclusion}\label{sec:typeII-assembly}
\subsection{Local Entry}

Let \((u,p)\) be a suitable weak solution of the three-dimensional
Navier--Stokes equations
\[
  \partial_tu+(u\cdot\nabla)u+\nabla p=\nu\Delta u,
  \qquad \nabla\cdot u=0,
\]
on \(\mathbb R^3\times(T-\delta,T)\).  For \(z_0=(x_0,T)\) and \(r>0\), set
\[
  Q_r(z_0)=B_r(x_0)\times(T-r^2,T),
\]
\[
  C(z_0,r)=r^{-2}\int_{Q_r(z_0)}|u|^3\,dx\,dt,
\]
and
\[
  D(z_0,r)=r^{-2}\int_{T-r^2}^{T}\int_{B_r(x_0)}
  |p(x,t)-(p)_{B_r(x_0)}(t)|^{3/2}\,dx\,dt.
\]
The pressure is always taken modulo functions of time, and the spatial mean
normalization is the one used in the Caffarelli--Kohn--Nirenberg criterion.

\begin{definition}[Singular set and local Type II point]
Let \(\Sigma(T)\) be the set of points \(x_0\) for which \(u\) is not locally
bounded in any backward cylinder \(Q_r(x_0,T)\).  A point \(x_0\in\Sigma(T)\)
is in the local Type II alternative if it carries a positive local
Caffarelli--Kohn--Nirenberg concentration sequence and no local Type I
scale-invariant bound holds at \((x_0,T)\).
\end{definition}

\begin{definition}[Positive local Type II concentration sequence]
\label{paper6:def:local-typeII-sequence}
A sequence \(r_n\downarrow0\) is a positive local Type II concentration sequence
at \((x_0,T)\) if there exists \(\eta>0\) such that
\[
  C((x_0,T),r_n)+D((x_0,T),r_n)\ge \eta
  \qquad\text{for all }n,
\]
and \((x_0,T)\) is in the local Type II alternative.  The associated parabolic
rescalings are
\[
  u_n(y,s)=r_nu(x_0+r_ny,T+r_n^2s),
  \qquad
  p_n(y,s)=r_n^2p(x_0+r_ny,T+r_n^2s).
\]
\end{definition}

\begin{theorem}[Local concentration and Type I/Type II dichotomy]
\label{paper6:thm:paper0}
If \(\Sigma(T)\ne\emptyset\), then for every \(x_0\in\Sigma(T)\) there are
\(\eta>0\) and \(r_n\downarrow0\) such that
\[
  C((x_0,T),r_n)+D((x_0,T),r_n)\ge\eta
  \qquad\text{for all }n.
\]
At the same point exactly one of the following alternatives holds: there is a
local Type I scale-invariant bound, or \((x_0,T)\) is in the local Type II
alternative.
\end{theorem}

\begin{proof}
This is the pointwise CKN concentration dichotomy of
\cref{paper1:thm:paper0-dichotomy}, restated in the notation used in the final
assembly.
\end{proof}

\subsection{Local Theorems Used in the Assembly}

The local results needed for the final assembly are restated in the following form.

\begin{theorem}[Repaired-gauge representation]
\label{paper6:thm:representation}
Let a positive local Type II concentration sequence contain a nondegenerate
selected core.  Then, after passing to selected windows, there are local
scale-center variables \(\lambda(\tau)>0\), \(x_c(\tau)\), renormalized fields
\[
  V(y,\tau)=\lambda(\tau)u(x_c(\tau)+\lambda(\tau)y,t(\tau)),
  \qquad
  P(y,\tau)=\lambda(\tau)^2p(x_c(\tau)+\lambda(\tau)y,t(\tau))+\pi(\tau),
\]
and modulation coefficients \(a,b\) such that
\[
  \partial_\tau V+(V\cdot\nabla)V+\nabla P
  =
  \nu\Delta V+a(\tau)(V+y\cdot\nabla V)+b(\tau)\cdot\nabla V,
  \qquad \nabla\cdot V=0.
\]
On each compact renormalized window one has the local suitable energy
inequality, local pressure reconstruction modulo functions of time, stability
under subsequences, and the compact-window bounds proved in \cref{sec:repaired-gauge-construction}.  Failure of the nondegenerate scale-center selection is part of the
multibubble, cascade, or scale-rigidity alternatives.
\end{theorem}

\begin{proof}
\Cref{paper2:thm:C} is precisely the compact-window representation theorem.
\end{proof}

\begin{theorem}[Compact single-core exclusion]
\label{paper6:thm:single-core}
No compact single-core vanishing-dissipation branch can be a local Type II
sequence.  More explicitly, let a represented selected-window sequence have:
\begin{enumerate}[label=\textup{(\roman*)}]
\item retained positive local velocity concentration on a fixed compact cylinder;
\item uniformly bounded compact-cylinder suitable estimates \(A+E+C+D\);
\item nondegenerate local scale-center representation;
\item bounded modulation coefficients;
\item local pressure reconstruction modulo spatial means;
\item vanishing localized dissipation on some retained compact cylinder.
\end{enumerate}
Then the sequence is incompatible with the local Type II condition.
\end{theorem}

\begin{proof}
\Cref{paper1:thm:single-core-criterion} gives the stated exclusion.
\end{proof}

\begin{theorem}[Multibubble and cascade exclusion]
\label{paper6:thm:multibubble}
In the local terminal class where terminal admissibility is supplied by
\cref{paper3:prop:terminal-admissibility-active-frames} and scale-rigidity is discharged by
\cref{paper3:prop:scale-rigidity-discharged}, no local multibubble, cascade, or
gauge-degenerate Type II branch can occur in the class produced by the local
active-core extraction.  Equivalently, every failure of the single
nondegenerate compact-core alternative is excluded by the local
multibubble/cascade analysis.
\end{theorem}

\begin{proof}
\Cref{paper3:thm:multi}, together with its removal corollary
\cref{paper3:cor:multi-removal}, gives the stated exclusion.
\end{proof}

\begin{theorem}[Rough-core reduction]
\label{paper6:thm:rough-core}
On any compact renormalized window satisfying the local energy inequality and
bounded modulation hypotheses, finite outer velocity-pressure CKN control
\(\mathcal C+\mathcal D\) implies finite inner windowed \(H^1\) control.  Hence
failure of local windowed \(H^1\) control on a retained inner cylinder forces
failure of the compact-cylinder \(\mathcal C+\mathcal D\) bounds on every
relevant enclosing cylinder.
\end{theorem}

\begin{proof}
\Cref{paper4:cor:rough-core} gives the stated reduction.
\end{proof}

\begin{theorem}[Finite-cost scale-collapse exclusion]
\label{paper6:thm:scale-cost}
No renormalized branch with finite scale-collapsing cost, and hence no branch
with finite absolute scale cost, can have genuine scale collapse together with a
nonvanishing selected core.
\end{theorem}

\begin{proof}
\Cref{paper5:thm:cost-exclusion,paper5:cor:absolute-cost-exclusion} give the stated exclusion.
\end{proof}

\begin{theorem}[Scale-collapse decomposition of alternatives]
\label{paper6:thm:scale-stratification}
Let a represented Type II branch have genuine scale collapse and a nonvanishing
selected core.  After passing to selected terminal windows, exactly one of the
alternatives isolated in \cref{sec:scale-collapse-cost} occurs:
\begin{enumerate}[label=\textup{(\roman*)}]
\item finite scale-collapsing cost;
\item finite absolute scale cost;
\item thin-drift alternative, routed by
\cref{paper6a:thm:thin-negative-drift-closed};
\item autonomous-modulation alternative;
\item compactness-package first failure, routed by
\cref{paper5:thm:selected-compactness-failure-routed};
\item a nonzero autonomous reduced limit in the scale-rigid state.
\end{enumerate}
In the local Type II class of alternatives used by the multibubble/cascade and
scale-collapse analyses, alternatives
\textup{(iii)}--\textup{(vi)} do not constitute additional compact single-core
mechanisms: thin drift is routed to another local row or a closed retained row,
the modulation alternative lies outside the retained compact scale-collapse
branch, the compactness-package failure is routed to a more specific named
local row, and the remaining retained scale-collapse state is removed by
\cref{paper5:thm:compact-scale-collapse-rigidity-discharged}.
\end{theorem}

\begin{proof}
The assertion follows by combining the scale-collapse decomposition of
alternatives with the local compact scale-collapse routing theorem
\cref{paper5:thm:compact-scale-collapse-rigidity-discharged}.
\end{proof}

\begin{lemma}[Compact retained core dissipation dichotomy]
\label{paper6:lem:compact-core-dissipation-dichotomy}
Let a represented retained compact core carry positive local Type II
concentration on a fixed compact cylinder and satisfy the retained
compact-cylinder bounds.  Then, after passing to selected terminal windows, one
of the following ordered alternatives holds:
\begin{enumerate}[label=\textup{(\roman*)}]
\item there is a retained compact cylinder on which the selected windows have
vanishing localized dissipation, hence the branch is a compact single-core
vanishing-dissipation branch in the sense of
\cref{paper1:def:single-core-branch};
\item an exterior, rough-core, modulation, scale, carrier, or cost adjacent
label occurs before the branch is a retained compact single-core state;
\item on a retained core cylinder carrying the positive local velocity concentration, the
localized dissipation has divergent tail integral, hence the canonical
localized cost diverges for every cutoff which is identically one on that core.
\end{enumerate}
\end{lemma}

\begin{proof}
Fix a retained cylinder \(Q_\rho\) on which the positive velocity core is
recorded.  If
\[
  \int_{\tau_0}^{\infty}\int_{B_\rho}|\nabla V|^2\,dy\,d\tau<\infty,
\]
then choose retained selected windows \(I_{n_j}\), with disjoint enlarged
windows after passing to a subsequence, such that
\[
  \int_{I_{n_j}}\int_{B_\rho}|\nabla V|^2\,dy\,d\tau\to0 .
\]
The retained compact-cylinder bounds, pressure reconstruction, admissible
modulation control on the selected windows, and positive velocity core are
exactly the remaining clauses of \cref{paper1:def:single-core-branch}; if any
retained compact-cylinder, pressure, exterior, rough-core, modulation, scale,
carrier, or cost input needed for that definition fails, the ordered
first-failure rule assigns item \textup{(ii)}.  In particular, failure of the
integrated modulation control is the modulation-window exit of
\cref{paper6a:lem:modulation-window-bound-or-exit}.  If no such exit occurs,
the selected subsequence gives item \textup{(i)}.
\par
If the displayed dissipation integral is infinite, then for any cutoff
\(\phi_{R_0}\equiv1\) on \(B_\rho\) the canonical localized cost satisfies
\[
  \int_{\tau_0}^{\infty}\tilde{\mathfrak D}_{R_0}(\tau)\,d\tau
  \ge
  \nu\int_{\tau_0}^{\infty}\int_{B_\rho}|\nabla V|^2\,dy\,d\tau
  =\infty,
\]
unless one of the ordered local tests needed to define the retained compact
core has failed first.  Such first failures are precisely the exterior,
rough-core, modulation, scale, or carrier/cost adjacent labels of the local
state-space validation, giving item \textup{(ii)}.  If no first failure occurs,
the displayed divergence is item \textup{(iii)}.
\end{proof}

\subsection{Assembly of Alternatives}

\begin{lemma}[Bounded selected-window limits are locally regular modulo named remaining activity]
\label{paper6:lem:bounded-window-local-regularization}
\label{paper10:lem:bounded-selected-window-local-regularity}
Let \((V_n,P_n)\) be represented suitable weak solutions on \(Q_{2R}\), with
the local pressure reconstruction and compact-cylinder suitable bounds supplied
by the repaired-gauge state-space construction.  Assume
\[
  V_n\to V_*
  \qquad\text{strongly in }L^3(Q_{2R}),
  \qquad
  V_*\in L^\infty(Q_{2R}).
\]
Then for every compact \(K\Subset Q_R\), exactly one of the following occurs
after passing to a subsequence:
\begin{enumerate}[label=\textup{(\roman*)}]
\item \(V_n\) is uniformly bounded on a finite subcover of \(K\), hence the
corresponding physical branch is regular on the physical image of \(K\);
\item a positive CKN amount survives outside that regular subcover and is
captured by one of the named ordered rows: a smaller active core, an exterior
core, a pressure-only carrier, a multibubble/cascade branch, a
gauge-degenerate branch, or a rough-core/local-energy failure row.
\end{enumerate}
\end{lemma}

\begin{proof}
Fix \(z_0=(y_0,s_0)\in K\).  Since \(V_*\in L^\infty(Q_{2R})\), for sufficiently
small \(\rho=\rho(z_0)>0\) with \(Q_{2\rho}(z_0)\Subset Q_{2R}\),
\[
  \rho^{-2}\int_{Q_{2\rho}(z_0)} |V_*|^3
\]
is below a fixed fraction of the represented-gauge CKN threshold.  Strong
\(L^3\)-convergence gives the same smallness for \(V_n\) on
\(Q_{2\rho}(z_0)\) for all large \(n\).

By \cref{paper3:lem:bounded-window-vmo-pressure}, the bounded selected-window
limit has uniformly vanishing velocity--pressure CKN quantity at sufficiently
small scales on compact subcylinders.  The strong convergence of \(V_n\), the
strong convergence of the local Calder\'on--Zygmund pressure pieces, and
harmonic pressure decay therefore
transfer this smallness to \((V_n,P_n)\) for all large \(n\), unless the
pressure-atlas or harmonic-pressure bounds fail.  Such a failure is one of the
named pressure rows in alternative \textup{(ii)}.  In the non-failure case, the
represented-gauge CKN regularity lemma gives uniform boundedness of \(V_n\) on
a smaller cylinder around \(z_0\).

A finite subcover of \(K\) gives alternative \textup{(i)} unless, during the
finite covering process, a positive amount of CKN mass is found outside the
regularized subcover.  That remaining positive CKN mass is, by the ordered
state-space construction, a smaller active core, exterior core, pressure-only
carrier, multibubble/cascade, gauge-degenerate, or local-energy/rough-core
failure.  These are exactly the named alternatives in \textup{(ii)}.
\end{proof}

\begin{definition}[Model carrier versus original-window carrier]
A model carrier is a positive CKN lower bound on a fixed cylinder after
rebasing a selected descendant to the model cylinder.  An original-window
carrier is a positive CKN lower bound on a fixed subcylinder of the original
bounded selected window, independent of the generation of the reselection tree.

The diffuse reselection-tree case means that model carriers exist at every
generation, but no subsequence yields an original-window carrier.
\end{definition}

\begin{lemma}[Reselection-tree support points are realized by descendant states]
\label{paper10:lem:reselection-tree-support-realization}
Consider an infinite bounded-window reselection tree whose selected descendant
cylinders all satisfy the fixed lower CKN threshold.  Normalize the activity
mass on each generation and record for every selected descendant its
compactified relative scale-center coordinates inside the original bounded
window, together with the rebased local occupation state on a fixed model
cylinder and the usual failure faces for representation, pressure atlas, gauge
chart, compact-upper, local-energy, and rough-core breakdown.  After passing
to a subsequence in the generation index,
these normalized tree measures converge weakly-* on a compact metrizable
reselection state space.  The limiting tree measure itself is retained as part
of the compactification, so later diffuse arguments may apply
\cref{paper6a:lem:normalized-defect-atomic-or-diffuse} directly to it.
The tree occupation coordinates are continuous coordinate projections in this
product compactification, and their positive superlevel sets are closed in the
reselection state space.

Every support point of the limit state is realized by an actual sequence of
descendant cylinders and corresponding rebased local states from the original
reselection tree.  If such a realized support profile carries positive compact
local activity on some fixed model cylinder, then it gives a model carrier.  It
gives an original-window carrier only if the relative scale-center coordinates
of the realizing descendants remain in a compact chart of the original bounded
window.
\end{lemma}

\begin{proof}
The compactified relative scale-center chart is compact by construction, and the
rebased local-state factor is compact by
\cref{paper6a:lem:compact-local-state-factor} once one records failure faces
for loss of represented structure, compact-upper control, local energy,
pressure atlas, gauge chart, and rough-core breakdown.  Hence the normalized
tree measures are tight and possess weak-* limits.

Let \(\Theta\) be one such limit and let
\(\mathfrak s\in\operatorname{supp}\Theta\).  Choose nested neighborhoods
\((U_m)\) of \(\mathfrak s\) with diameters tending to zero.  Since
\(\Theta(U_m)>0\), for every \(m\) and all sufficiently large generations there
is an actual selected descendant cylinder whose rebased state lies in \(U_m\).
Selecting one such descendant for each \(m\) produces a sequence of genuine
rebased local states converging to \(\mathfrak s\), proving realization.

If the realized support profile has positive compact local activity on a fixed
model cylinder \(Q_\rho\), then some neighborhood of \(\mathfrak s\) is defined
by a quantitative lower bound for the CKN quantity on \(Q_\rho\).  The support
property then forces actual descendants from the original tree to satisfy the
same lower bound for all large \(m\), which is exactly a model carrier.  If the
relative scale-center coordinates of those descendants remain in a compact chart
of the original bounded window, undoing the rebasing sends them to physical
subcylinders of the original bounded window and produces an original-window
carrier.
\end{proof}

\begin{lemma}[Diffuse reselection-tree row is closed]
\label{paper10:lem:diffuse-reselection-tree-row-closed}
Let an infinite bounded-window reselection tree be produced after a bounded
selected-window region has been regularized.  Assume that the tree carries the
fixed CKN lower threshold at each selected generation, but that no subsequence
of descendants produces an original-window carrier in the original
reselection branch.

Then the diffuse reselection-tree alternative is not terminal.  More precisely,
after passing to the compactified reselection-tree state space of
\Cref{paper10:lem:reselection-tree-support-realization}, one of the following
named rows occurs:
\begin{enumerate}[label=\textup{(\roman*)}]
\item a representation, pressure atlas, gauge chart, compact-upper,
  local-energy, or rough-core failure face;
\item a same-point cascade, scale-collapse, or noncompact exterior row;
\item an original-window carrier, contradicting the diffuse/no-original-window-carrier case;
\item a visible positive residual defect in the augmented local occupation-state
  coordinates of \Cref{paper6a:thm:local-residual-exhaustion}, hence a
  named adjacent local row or a retained comparable class already grouped
  before the bounded-window reselection was declared.
\end{enumerate}
Consequently a diffuse bounded-window reselection tree cannot remain as a
terminal physical Type II branch.
\end{lemma}

\begin{proof}
Apply \Cref{paper10:lem:reselection-tree-support-realization}.  The normalized
tree measures have a weak-* limit on a compact metrizable reselection state
space, and every support point is realized by actual descendant cylinders from
the original reselection tree.

If a support point lies on a failure face for representation, pressure atlas,
gauge chart, compact upper estimates, local energy, or rough-core control, then
the corresponding named row occurs.  This is alternative \textup{(i)}.

Apply \Cref{paper6a:lem:normalized-defect-atomic-or-diffuse} to the normalized
reselection-tree measures.  If the atomic/activity alternative occurs, then the
support-realization lemma gives a model carrier.  If the relative scale-center
coordinates of the corresponding descendants stay in a compact chart, that model
carrier yields an original-window carrier, contradicting the diffuse case.
If those relative coordinates do not stay in a compact chart, then the ordered
relative scale-center partition gives the same-point cascade,
scale-collapse, or noncompact exterior row.  This is alternatives
\textup{(ii)} and \textup{(iii)}.

If the genuinely diffuse alternative occurs, then the normalized tree measure
itself is a visible diffuse occupation coordinate.  Applying
\Cref{paper6a:thm:local-residual-exhaustion} gives a named adjacent row or a
retained comparable class.  The retained comparable class is impossible because
comparable retained classes are grouped before bounded-window reselection is
declared.  Hence only named adjacent rows remain.

Thus the diffuse reselection-tree alternative is not terminal.
\end{proof}

\begin{lemma}[Bounded-window reselection is well-founded]
\label{paper10:lem:bounded-window-reselection-well-founded}
Let a physical local Type II germ produce a bounded selected-window limit on a
selected compact cylinder.  Remove the regular portion of that bounded window
using the local CKN criterion, and rerun the ordered local selection on any
remaining positive CKN concentration in the physical image of a smaller
terminal neighborhood.

Then exactly one of the following alternatives occurs:
\begin{enumerate}[label=\textup{(\roman*)}]
\item there is a physical neighborhood \(Q_{\rho_0}(x_0,T)\) and a scale
\(r_0>0\) such that every admissible subcylinder
\(Q_r(z)\subset Q_{\rho_0}(x_0,T)\), \(0<r<r_0\), satisfies
\[
   C(z,r)+D(z,r)<\varepsilon_{\rm CKN};
\]
hence the germ is locally regular there;
\item after finitely many reselections, the branch enters one of the named
state-space rows: active core, pressure-only row, separated exterior row,
multibubble row, same-point cascade row, gauge-degenerate row, rough-core row,
or retained terminal compact branch;
\item the reselection process does not terminate finitely.  Then the selected
centers and scales generate one of the infinite rows: same-point scale cascade,
scale-collapse row, noncompact exterior row, or diffuse reselection-tree row.
\end{enumerate}
\end{lemma}

\begin{proof}
Each reselection produces a selected cylinder \(Q_{r_j}(z_j)\) with CKN
quantity bounded from below by the fixed activity threshold.  If no such
cylinder exists below some physical scale, the finite-cover form of the CKN
criterion gives item \textup{(i)}.

Assume therefore that selected cylinders continue to exist.  Compare each new
selected cylinder with the previous selected core through the local complexity
triple
\[
  \Xi_j
  :=
  \bigl(\mathfrak c_j,\mathfrak s_j,\eta_*\bigr),
\]
where \(\mathfrak c_j\) records whether the new center is same-point,
physically separated at comparable scale, or escapes every compact relative
chart, and \(\mathfrak s_j\) records whether the new scale is strictly smaller,
comparable, strictly larger in the renormalized variables, or has no compact
relative limit.  This is a finite ordered partition of the relative
scale-center chart.

If infinitely many reselections fall in the comparable separated-center class,
then a finite-cover selection on the corresponding comparable cylinders gives a
separated exterior or multibubble row.  If infinitely many reselections stay at
the same physical point with strictly decreasing physical scales, then the
selected cylinders form an infinite nested chain and we are in the same-point
scale-cascade row.  If infinitely many reselections stay at the same physical
point while the renormalized radius strictly increases, then the branch is the
same-point larger-scale cascade row.  If the relative centers or relative
scales leave every compact subset of the relative scale-center chart, then by
definition the branch is in the scale-collapse or noncompact exterior row.

Finally, if the CKN mass persists but no subsequence produces an
original-window carrier, normalize the remaining CKN mass on the reselection
tree at the first scale on which each descendant cylinder is selected.  Apply
\cref{paper10:lem:reselection-tree-support-realization} and then
\cref{paper6a:lem:normalized-defect-atomic-or-diffuse}.  If the
atomic/activity alternative occurs, then one obtains a model carrier.  If its
relative scale-center coordinates stay in a compact chart, this produces an
original-window carrier, contradiction.  If not, the branch is in the
same-point cascade, scale-collapse, or noncompact exterior row.  If the
genuinely diffuse alternative occurs, then by definition the branch is in the
diffuse reselection-tree row.  Thus infinite reselection is routed to one of
the named infinite rows.
\end{proof}

\begin{theorem}[Bounded selected-window branch is physically routed]
\label{paper6:thm:bounded-selected-window-physical-routing}
\label{paper6:thm:bounded-window-physical-closure}
Let \((u,p,x_0,T)\) be a physical local Type II singular germ, and suppose that
one selected-window extraction has a nonzero bounded selected-window limit in
the sense of \cref{paper1:def:bounded-window-limit}.  Then the bounded-window
branch is not a terminal unrecorded physical Type II branch.  More precisely,
after applying
\cref{paper6:lem:bounded-window-local-regularization}, exactly one of the
following occurs:
\begin{enumerate}[label=\textup{(\roman*)}]
\item the selected window is regular/removable and the ordered local
concentration algorithm continues on any surviving smaller, exterior, pressure,
multicore, cascade, gauge-degenerate, or rough-core row;
\item no surviving CKN concentration remains in the physical neighborhood
under consideration, in which case the local Type I/regular alternative holds
there, contradicting the assumption that the germ is a physical local Type II
singular germ;
\item the surviving concentration is one of the named rows in the local
state-space ledger and is closed by the corresponding row theorem.
\end{enumerate}
Consequently bounded selected-window limits are closed physical routing rows,
not terminal residual Type II branches.
\end{theorem}

\begin{proof}
Apply \cref{paper6:lem:bounded-window-local-regularization} on an exhausting
finite family of compact subcylinders inside the selected window.  If the
bounded-window region regularizes and no CKN concentration remains, the
physical image has a local boundedness/Type I neighborhood, contrary to the
physical Type II assumption.  If CKN concentration remains, the ordered
state-space tests are not restarted ad infinitum without classification:
apply \cref{paper10:lem:bounded-window-reselection-well-founded} to any
remaining positive CKN concentration.  Alternative \textup{(i)} gives physical
local regularity.  Alternative \textup{(ii)} enters one of the finite named
rows.  Alternative \textup{(iii)} enters one of the infinite named rows.  All
rows are closed by the terminal ledger.  Hence the bounded selected-window
branch cannot remain as a physical Type II germ.
\end{proof}

\begin{theorem}[Ordered local state-space exhaustion]
\label{paper6:thm:ordered-local-state-space-exhaustion}
\label{paper6:thm:ordered-state-space-exhaustion}
Let \(r_n\downarrow0\) be a positive local Type II concentration sequence in the
sense of \cref{paper6:def:local-typeII-sequence}.  Starting from the associated
rescalings, run the following ordered local tests:
\[
\begin{array}{l}
\text{C1 admissibility}
\to
\text{representation/gauge}
\to
\text{single-core versus multibubble},
\\[2pt]
\text{pressure-only versus velocity carrier}
\to
\text{compact upper estimates}
\to
\text{rough-core/}H^1\text{ control},
\\[2pt]
\text{scale-collapse/scale-rigid}
\to
\text{canonical cost validation}
\to
\text{terminal retained branch}.
\end{array}
\]
After passing to a diagonal subsequence, exactly one first successful retained
state or one first failed local row occurs.  Every first failed row is one of
the named local rows later closed in the ledger, and every retained state is
one of the following: a represented compact single-core branch, a
multibubble/cascade/gauge-degenerate branch, loss of local windowed \(H^1\)
control on a retained compact cylinder, a finite-cost scale-collapse branch
with nonvanishing selected core, or a retained compact scale-collapse or
scale-rigid terminal state.
\end{theorem}

\begin{proof}
The C1 admissibility test is the pointwise local Type I/Type II dichotomy of
\cref{paper6:thm:paper0}; on a physical local Type II germ its class-side
failures are recorded by
\cref{paper6:lem:no-extra-outside-class-row,paper6:cor:covered-class-entry-discharged}.
Given admissibility, the ordered repaired-gauge representation tests of
\cref{paper6:thm:representation,cor:c2-ordered-representation,cor:c2r-ordered-classification}
either produce a represented local branch on selected windows or record the
first representation/gauge failure.
\par
On a represented branch, the active-core extraction is tested first.  Failure
of uniqueness, same-point splitting, separated cores, or unrepaired scale
degeneracy is the multibubble/cascade side of
\cref{paper6:thm:multibubble}.  If a single retained core is selected, then the
ordered pressure atlas and carrier tests decide whether one is on a retained
velocity core, a pressure-only row, or an exterior/separated carrier row; this
is exactly the local routing package of
\cref{paper6:lem:pressure-only-critical-mass-routing,paper8:cor:exterior-concentration-exit-discharged}.
The compact-window upper estimates are then supplied by
\cref{paper5:prop:selected-window-suitable-estimates}, while the retained
spacetime mass floor is supplied by \cref{paper5:thm:core-retention}
whenever the loss-of-active-core row has not already been recorded.
\par
Before applying any repaired gauge chart, the strong gauge-chart entry test
\cref{paper2:lem:strong-gauge-chart-entry-test} is run.  Failure of this test
is the repaired-gauge degeneracy row.  If a nonzero bounded selected-window
limit appears, then \cref{paper3:lem:bounded-window-vmo-pressure} gives the
local velocity--pressure CKN smallness of the bounded limit, and
\cref{paper10:lem:bounded-window-reselection-well-founded} shows that
infinite reselection is not left open: it is routed to the same-point cascade,
scale-collapse, noncompact exterior, or diffuse reselection-tree rows.
\par
If a nonzero bounded selected-window limit appears at this stage, the branch is
not retained in the strict selected-window Type II row by
\cref{paper3:lem:scale-rigid-bounded-limit-exit}; in the physical assembly that
event is closed by \cref{paper6:thm:bounded-selected-window-physical-routing}, and in
the local branchwise assembly it is returned to the ordered active-window
decomposition rather than declared terminal.  If no bounded selected-window
limit appears, the rough-core test \cref{paper6:thm:rough-core} decides whether
local \(H^1\) control is retained or the branch returns to the multibubble or
compact-cylinder-failure side.
\par
When critical tightness fails, apply
\cref{paper9:lem:critical-tightness-failure-row}.  This produces one of the
named rows: same-point larger-scale cascade, separated exterior concentration,
noncompact exterior, pressure/gauge/reconstruction failure,
rough-core/local-energy/compact-upper failure, or diffuse critical-tail.  When
pressure-only exterior selection nests indefinitely, apply
\cref{paper8:lem:nested-pressure-only-well-founded}; this produces one of the
named rows: same-point pressure cascade, separated exterior concentration,
scale-collapse, or diffuse pressure-only.  The later closure of the diffuse
critical-tail and diffuse pressure-only rows is deferred to the remaining-exit
closure layer and is not used in the routing step here.
\par
The remaining represented retained branch is tested by the scale package of
\cref{paper6:thm:scale-stratification}: either one reaches finite-cost
scale-collapse, thin drift, autonomous modulation, compactness-package first
failure, or the retained compact scale-collapse/scale-rigid states.  Finally,
the canonical cost-validation layer is available on every retained compact
branch by \cref{paper6:lem:canonical-cost-channel-available}, and its first-failure map
is the ordered list in \cref{paper7:lem:no-unrecorded-cost-escape}.  Because
the list of tests is finite and each stage either records its first failure or
passes to the next named layer after subsequence extraction, the diagonal
subsequence construction terminates in exactly one first failed row or one
retained terminal state.  The retained terminal states are exactly the five
retained alternatives listed in the statement.
\end{proof}

\begin{theorem}[Local Type II decomposition into local alternatives]
\label{paper6:thm:state-decomposition}
Every positive local Type II concentration sequence in \cref{paper6:def:local-typeII-sequence}, after passing to subsequences and selected windows, enters one of the following alternatives:
\begin{enumerate}[label=\textup{(\roman*)}]
\item a represented compact-core branch carrying a single retained active core,
which by
\cref{paper6:lem:compact-core-dissipation-dichotomy} either has
vanishing-dissipation selected windows, exits through a named adjacent local
row before retention, or enters the retained local cost-divergence branch;
\item a multibubble, cascade, or gauge-degenerate branch;
\item loss of local windowed \(H^1\) control on a retained compact cylinder;
\item a finite-cost scale-collapse branch with nonvanishing selected core;
\item a retained compact scale-collapse or scale-rigid terminal state.
\end{enumerate}
\end{theorem}

\begin{proof}
This is the list of retained terminal outcomes produced by the ordered
exhaustion theorem \cref{paper6:thm:ordered-local-state-space-exhaustion}.  The
single-core represented branch is refined by
\cref{paper6:lem:compact-core-dissipation-dichotomy}; all failures of unique
retained-core selection belong to the multibubble/cascade side of
\cref{paper6:thm:multibubble}; rough-core loss is returned by
\cref{paper6:thm:rough-core} to the compact-cylinder failure side; and genuine
scale collapse is divided by \cref{paper6:thm:scale-stratification} into
finite-cost scale collapse or one of the retained compact
scale-collapse/scale-rigid terminal states.  No other retained terminal state
survives the ordered first-failure algorithm.
\end{proof}

\begin{proposition}[Elimination of local Type II states]
\label{paper6:prop:state-elimination}
No positive local Type II concentration sequence can remain in any retained
state among alternatives \textup{(i)}--\textup{(v)} of
\cref{paper6:thm:state-decomposition}.  If the ordered validation inside
alternative \textup{(i)} exits through a named adjacent local row before
retention, that branch is not a retained local state and is closed in the
ambient adjacent-routing layer.
\end{proposition}

\begin{proof}
Alternative \textup{(i)} is excluded by the compact single-core criterion,
\cref{paper6:thm:single-core}, in the vanishing-dissipation subcase of
\cref{paper6:lem:compact-core-dissipation-dichotomy}.  The adjacent-label
subcase of that lemma is not a retained compact single-core state: the ordered
validation has stopped earlier at a named local row.  In the divergent
localized-dissipation subcase, the canonical localized cost diverges on the
retained core and the cost-divergence closure
\cref{paper6a:thm:cost-divergence-exclusion-discharged,paper6a:cor:cost-divergence-exclusion-unconditional}
excludes the branch after the windowwise adjacent alternatives have been
routed.

Alternative \textup{(ii)} is excluded by the local multibubble and cascade
theorem, \cref{paper6:thm:multibubble}.

Alternative \textup{(iii)} is not an independent terminal state.  By
\cref{paper6:thm:rough-core}, rough-core loss of local \(H^1\) control forces
failure of the compact-cylinder \(\mathcal C+\mathcal D\) bounds on enclosing
compact cylinders.  Such failure yields an active compact-cylinder
concentration branch and therefore belongs to the multibubble/cascade
alternative, already excluded by \cref{paper6:thm:multibubble}.

Alternative \textup{(iv)} is excluded by the finite-cost scale-collapse theorem,
\cref{paper6:thm:scale-cost}.

Alternative \textup{(v)} is the retained compact scale-collapse or scale-rigid
terminal state.  Retained compact scale-collapse states are excluded by
\cref{paper5:thm:compact-scale-collapse-rigidity-discharged}.  Scale-rigid
terminal states are excluded by
\cref{paper3:prop:scale-rigidity-discharged}: the state either has a nonzero
bounded selected-window limit, in which case the strict selected-window branch
stops and the bounded-window event is routed by
\cref{paper6:thm:bounded-selected-window-physical-routing} in the physical assembly, or
it reduces to the compact single-core branch, which is excluded by
\cref{paper6:thm:single-core}.  Hence it cannot remain as a retained local Type II
branch.
\end{proof}

\subsection{Physical Entry into the Covered Local Type II Path}

\begin{definition}[Physical local Type II singular germ]
\label{paper6:def:physical-local-typeII-germ}
Let \((u,p)\) be a suitable weak solution on
\(\mathbb R^3\times(T-\delta,T)\).  A tuple
\((u,p,x_0,T)\) is a physical local Type II singular germ if:
\begin{enumerate}[label=\textup{(\roman*)}]
\item \(x_0\in\Sigma(T)\);
\item the local Type I scale-invariant bound of
\cref{paper1:thm:paper0-dichotomy} fails at \((x_0,T)\);
\item the germ is an interior Navier--Stokes germ, meaning that all sufficiently
small backward cylinders \(Q_r(x_0,T)\) are contained in the Navier--Stokes
domain under consideration.
\end{enumerate}
\end{definition}

\begin{definition}[Covered local Type II entry]
\label{paper6:def:covered-local-typeII-entry}
A physical local Type II singular germ has covered local Type II entry if one
of the following two alternatives occurs:
\begin{enumerate}[label=\textup{(\roman*)}]
\item one of its Caffarelli--Kohn--Nirenberg concentration sequences, after
passing to a subsequence and applying the selected-window reductions used in
this paper, enters one of the alternatives in
\cref{paper6:thm:state-decomposition};
\item before the selected-window state decomposition is reached, the ordered
local tests assign the germ to a class-incompatible alternative or to a named
non-retained local state-space exit.
\end{enumerate}
\end{definition}

\begin{lemma}[Physical Type II germs are admissible local branches]
\label{paper6:lem:physical-typeII-germ-admissible}
Let \((u,p,x_0,T)\) be a physical local Type II singular germ.  After translating
space and time and restricting to a smaller terminal interval if necessary, the
terminal germ defines an element of
\(\mathcal U_{\mathrm{II}}^{NS}\) in the sense of
\cref{def:c1-admissible-typeII-class}.
\end{lemma}

\begin{proof}
Choose \(0<\delta'<\delta\) so that the cylinders under consideration remain
interior.  Set
\[
  \tilde u(y,s)=u(y+x_0,T-\delta'+s),\qquad
  \tilde p(y,s)=p(y+x_0,T-\delta'+s),
  \qquad 0<s<\delta',
\]
and \(T^*=\delta'\).  This space-time translation and restriction preserve the
local Navier--Stokes equations, suitability, and the pressure normalization
modulo functions of time.  The terminal time \(T^*\) is finite by construction.
Since \(x_0\in\Sigma(T)\), the terminal germ is not a local continuation
branch.  By the definition of physical local Type II germ, the local Type I
scale-invariant bound fails at \((x_0,T)\), so the germ is not in the local
Type I class.  The interior-domain clause excludes boundary and
non-Navier--Stokes-domain alternatives.  These are exactly the admissibility
clauses in \cref{def:c1-admissible-typeII-class}.
\end{proof}

\begin{lemma}[Physical Type II germs carry local Type II concentration sequences]
\label{paper6:lem:physical-typeII-germ-concentration-sequence}
Let \((u,p,x_0,T)\) be a physical local Type II singular germ.  Then there are
\(\eta>0\) and \(r_n\downarrow0\) such that
\[
  C((x_0,T),r_n)+D((x_0,T),r_n)\ge\eta
  \qquad\text{for all }n,
\]
and \(r_n\) is a positive local Type II concentration sequence in the sense of
\cref{paper6:def:local-typeII-sequence}.
\end{lemma}

\begin{proof}
Since \(x_0\in\Sigma(T)\), \cref{paper1:thm:paper0-dichotomy} gives
\(\eta>0\) and \(r_n\downarrow0\) such that the displayed CKN lower bound
holds.  The same dichotomy has exactly two pointwise alternatives: the local
Type I scale-invariant bound holds, or the point lies in the local Type II
alternative.  The first alternative is excluded by
\cref{paper6:def:physical-local-typeII-germ}.  Hence \((x_0,T)\) lies in the
local Type II alternative, and the displayed concentration sequence is exactly
a positive local Type II concentration sequence as defined in
\cref{paper6:def:local-typeII-sequence}.
\end{proof}

\begin{lemma}[No extra outside-class row for physical Type II germs]
\label{paper6:lem:no-extra-outside-class-row}
Let \((u,p,x_0,T)\) be a physical local Type II singular germ and let
\(\omega\) be the admissible terminal branch supplied by
\cref{paper6:lem:physical-typeII-germ-admissible}.  Any failure of the C1
analytic data for \(\omega\) is either class-incompatible or one of the named
non-retained local exits in
\cref{paper6a:thm:typeII-analytic-data-exhaustive}.  In particular, there is no
additional outside-covered-class alternative for such a physical Type II germ.
\end{lemma}

\begin{proof}
By \cref{paper6:lem:physical-typeII-germ-admissible}, \(\omega\in
\mathcal U_{\mathrm{II}}^{NS}\).  The ordered C1 discharge
\cref{paper6a:thm:typeII-analytic-data-exhaustive} applies on this class.  Its
conclusion is precisely that the first failure of
\cref{def:c1-exhaustion-failures} is assigned either to a class-incompatible
alternative excluded by \cref{def:c1-admissible-typeII-class}, or to a named
non-retained state-space branch already covered by the local stratification.
Thus a physical Type II germ has no residual ``outside the covered class'' row
after the ordered local tests are run.
\end{proof}

\begin{theorem}[Covered-class entry for physical Type II germs]
\label{paper6:thm:physical-typeII-covered-entry}
Every physical local Type II singular germ has covered local Type II entry in
the sense of \cref{paper6:def:covered-local-typeII-entry}.  More explicitly,
after passing to a CKN concentration sequence and to the selected windows used
in this paper, the germ enters one of the alternatives in
\cref{paper6:thm:state-decomposition} through the ordered exhaustion route of
\cref{paper6:thm:ordered-local-state-space-exhaustion}, unless it has already exited through one
of the class-incompatible or named non-retained local alternatives of
\cref{paper6a:thm:typeII-analytic-data-exhaustive}.
\end{theorem}

\begin{proof}
Let \((u,p,x_0,T)\) be a physical local Type II singular germ.  By
\cref{paper6:lem:physical-typeII-germ-admissible}, the germ defines an
admissible terminal branch.  Run the ordered C1 analytic-data tests on that
branch.  If a first failure occurs, then
\cref{paper6:lem:no-extra-outside-class-row} routes it to a class-incompatible
alternative or to a named non-retained local state-space exit.  This is covered
entry in the sense of
\cref{paper6:def:covered-local-typeII-entry}\,\textup{(ii)}.

It remains to consider the retained case, when no C1 failure occurs.  By
\cref{paper6:lem:physical-typeII-germ-concentration-sequence}, the germ carries
a positive local Type II concentration sequence.  The ordered exhaustion theorem
\cref{paper6:thm:ordered-local-state-space-exhaustion} is formulated exactly
for such sequences, and its retained outcomes are listed in
\cref{paper6:thm:state-decomposition}.  After passing to subsequences and
selected windows, the sequence therefore enters one of the five alternatives
listed there.  This is covered entry in the sense of
\cref{paper6:def:covered-local-typeII-entry}\,\textup{(i)}.  Thus every case is
a covered local state-space outcome.  The proof uses only
pointwise CKN concentration, local first-failure routing, and selected-window
state-space decomposition; it does not use a global \(L^3(\mathbb R^3)\) bound,
a global profile decomposition, global Kato smallness, or vanishing localized
dissipation beyond what is produced inside the selected local branch.
\end{proof}

\begin{corollary}[Discharge of covered-class entry]
\label{paper6:cor:covered-class-entry-discharged}
The covered-class entry obligation for the local Type II state-space analysis is
discharged:
every physical local Type II singular germ either enters the local
state-decomposition path of \cref{paper6:thm:state-decomposition}, or exits
earlier through a class-incompatible or named non-retained local alternative.
\end{corollary}

\begin{proof}
This is \cref{paper6:thm:physical-typeII-covered-entry}.
\end{proof}

\begin{remark}[Scope of the covered-entry discharge]
\label{paper6:rem:covered-entry-scope}
\Cref{paper6:cor:covered-class-entry-discharged} proves only the entry
statement: a physical local Type II germ is not lost before the local
state-space analysis begins.  It does not by itself supply the ambient
adjacent-stratum closure data or the expanded terminal data used in
the final assembly.  Nor does it assert that every named non-retained exit has
already been eliminated; it asserts that any such exit is recorded in the local
state-space taxonomy rather than being an untracked outside-class branch.  The
closure of those named exits is supplied by the local closure theorems and by
the ambient/terminal wrappers
\cref{paper7:cor:ambient-adjacent-wrapper-discharged,paper6a:cor:expanded-terminal-wrapper-discharged}.
The entry argument is also independent of
\cref{paper6:thm:local-typeII-exclusion}; it uses the pointwise CKN dichotomy,
the local Type I/Type II split, and the ordered C1 routing theorem, so the
discharge is not circular.
\end{remark}

\subsection{Closure of the remaining named exit rows}
\label{paper6:subsec:remaining-exit-closures}

The three closure theorems in this subsection convert the rows previously
listed in \cref{paper6:def:remaining-named-exits} into closed adjacent rows,
i.e.\ rows that are routed by an existing local closure or wrapper into a
retained-state alternative that has already been excluded.  The arguments are
local on compact terminal cylinders.  They use only the ordered local
state-space tests, the modulation-window dichotomy, the local pressure
reconstruction with harmonic remainder, and the existing wrappers
\cref{paper7:cor:ambient-adjacent-wrapper-discharged,paper6a:cor:expanded-terminal-wrapper-discharged}.
No global \(L^3\) bound, no global profile decomposition, no global Kato
smallness, no hidden vanishing localized dissipation, and no noncanonical
schedule are introduced.

\begin{lemma}[Bounded modulation defects are local carrier labels]
\label{paper6:lem:bounded-modulation-defect-routing}
Let \((a_n,b_n)\) be the repaired-gauge modulation coefficients on a fixed
compact terminal window \([0,T]\) for a represented branch with a nonvanishing
selected core \(V_n\), and suppose
\[
  \sup_n\int_0^T\bigl(|a_n(\tau)|+\|b_n(\tau)\|\bigr)\,d\tau<\infty .
\]
After passing to a subsequence, the ordered local test gives one of the
following outcomes:
\begin{enumerate}[label=\textup{(\roman*)}]
\item writing
  \[
    \bar a_n=T^{-1}\int_0^T a_n(\tau)\,d\tau,\qquad
    \bar b_n=T^{-1}\int_0^T b_n(\tau)\,d\tau,
  \]
  there are constants \(a_\infty\in\mathbb R\) and
  \(b_\infty\in\mathbb R^3\) such that, after passing to the same subsequence,
  \((\bar a_n,\bar b_n)\to(a_\infty,b_\infty)\) and either
  \[
    a_n\to a_\infty\quad\hbox{in }L^1(0,T),\qquad
    b_n\to b_\infty\quad\hbox{in }L^1(0,T;\mathbb R^3);
  \]
  or the nonconstant part of the modulation is invisible in the terminal local
  residual topology:
  \[
    (a_n-\bar a_n)(V_n+y\cdot\nabla V_n)
    +(b_n-\bar b_n)\cdot\nabla V_n
    \to0
  \]
  in \(L^1(0,T;H^{-1}(K))\) for every compact \(K\);
\item a signed component of the normalized modulation defect is visible in the
  local residual topology and defines a local occupation state.  The ordered
  state-space test records one of the adjacent labels:
  modulation-driven recentering, positive scale-shell work, negative
  scale-drift, loss of the selected active core, or repaired-gauge degeneracy.
\end{enumerate}
The conclusion is local on the fixed terminal window and uses no compactness
stronger than finite-measure weak-\(*\) compactness of the modulation defect
measures.
\end{lemma}

\begin{proof}
Write \(m_n=T^{-1}\int_0^T(a_n,b_n)\,d\tau\in\mathbb R^4\).  The \(L^1\) bound
implies that, after passing to a subsequence, \(m_n\to m_\infty=(a_\infty,
b_\infty)\).  Set
\[
  \alpha_n(\tau)=a_n(\tau)-\bar a_n,\qquad
  \beta_n(\tau)=b_n(\tau)-\bar b_n,
  \qquad m_n=(\bar a_n,\bar b_n).
\]
If
\[
  \|\alpha_n\|_{L^1(0,T)}+\|\beta_n\|_{L^1(0,T;\mathbb R^3)}\to0,
\]
then item \textup{(i)} holds.
\par
Assume instead that the displayed defect norm has a positive liminf.  Define
the local modulation-defect residual
\[
  \mathcal F_n^{\rm mod}
  :=
  \alpha_n(V_n+y\cdot\nabla V_n)+\beta_n\cdot\nabla V_n .
\]
If \(\mathcal F_n^{\rm mod}\to0\) in \(L^1(0,T;H^{-1}(K))\) for every compact
\(K\), then the represented local equation has the same terminal residual
limit as the sequence of constant-coefficient local equations with
coefficients \(m_n\).  Since \(m_n\to m_\infty\) in the finite-dimensional
coefficient space, this is the constant-coefficient autonomous channel after
the retained local compactness test is run.  Item \textup{(i)} holds in its
perturbative form.
\par
It remains to consider the case in which the defect residual is not
perturbative on some compact \(K\).  Passing to a further subsequence and
applying the ordered component convention, one of
\[
  \int_0^T(\alpha_n)_+\,d\tau,\qquad
  \int_0^T(\alpha_n)_-\,d\tau,\qquad
  \int_0^T|\beta_n|\,d\tau
\]
has a subsequence whose normalized signed residual contribution is nonzero in
the local terminal residual topology.  Normalize that visible component by its
mass.  The resulting probability measures are bounded in \(\mathcal M([0,T])\),
and hence have a weak-\(*\) subsequential limit.  Together with the nonzero
residual pairing, this is a visible local modulation-defect occupation state on
the selected window.
\par
The ordered local state-space test evaluates visible signed defect residuals
before declaring a retained autonomous window.  A visible \(\beta_n\)-component
is the modulation-driven recentering/moving-translation label of
\cref{paper7:def:stratified-moving-carrier-routing}.  A visible
\((\alpha_n)_+\)-component is the positive scale-shell label, and a visible
\((\alpha_n)_-\)-component is the negative scale-drift label.  If the selected
core on which the modulation acts has lost its active mass floor, the first
failed row is the retained active-core row.  If the defect cannot be represented
in the repaired chart, the first failed row is repaired-gauge degeneracy.  These
alternatives are ordered before autonomous retention, and they exhaust the
nonperturbative signed components of the modulation defect.  Hence item
\textup{(ii)} holds.
\end{proof}

\begin{theorem}[Autonomous-modulation failure is a closed adjacent row]
\label{paper6:thm:autonomous-modulation-failure-closed}
Let \((u,p,x_0,T)\) be a physical local Type II singular germ in the sense
of \cref{paper6:def:physical-local-typeII-germ}, and assume the germ has
admissible covered entry into the local state-space path.  If a represented
selected scale-collapse branch with nonvanishing selected core enters the
autonomous-modulation alternative
\cref{paper5:thm:scale-collapse-stratification}\textup{(v)} on a compact
terminal cylinder, then the branch enters one of the following closed
adjacent rows: the modulation-window-unbounded row routed by
\cref{paper6a:lem:modulation-window-bound-or-exit}; the modulation-driven
recentering or moving-carrier row routed by
\cref{paper7:def:stratified-moving-carrier-routing} and discharged in the
ambient adjacent wrapper
\cref{paper7:cor:ambient-adjacent-wrapper-discharged}; the negative
scale-drift row closed by
\cref{paper6a:thm:thin-negative-drift-closed}; the positive scale-shell row
routed by \cref{paper7:cor:windowwise-positive-scale-routing-holds}; the
loss-of-selected-active-core row closed by
\cref{paper6:thm:active-core-loss-closed}; or the gauge-degenerate /
repaired-gauge row already excluded from the retained branch by the ordered
representation tests.  Equivalently, the
autonomous-modulation alternative is not an independent remaining terminal
exit.
\end{theorem}

\begin{proof}
Fix a represented selected scale-collapse branch on a compact terminal
cylinder \(B_R\times[\tau_n,\tau_n+T]\) in the autonomous regime, and let
\(a_n(\tau)=a(\tau_n+\tau)\) and \(b_n(\tau)=b(\tau_n+\tau)\) on
\([0,T]\) be the absolutely continuous repaired-gauge modulation
coefficients identified in
\cref{paper6a:lem:repaired-modulation-coefficients}.  Cover the selected
window by finitely many unit windows of the type used in
\cref{paper6a:lem:modulation-window-bound-or-exit}, with the same
repaired-gauge chart variables as that lemma; the construction of the
chart variables and their absolute continuity are supplied for any
represented branch by
\cref{paper6a:lem:repaired-modulation-coefficients} and the ordered
representation tests
\cref{cor:c2-ordered-representation,cor:c2r-ordered-classification}.
Exactly one of the following two cases occurs.

\smallskip\noindent\emph{Case A (window-integrals unbounded).}  On at least
one unit window, the supremum
\(\sup_n\int(|a_n|+\|b_n\|_\infty)\,d\tau\) is infinite.  By the same
stratification of accumulated chart velocities used in the proof of
\cref{paper6a:lem:modulation-window-bound-or-exit}, namely the splitting
into large accumulated negative scale motion, large positive scale-shell
work, large translation work, and unrepresentable accumulated modulation,
the branch is recorded as one of the named modulation, scale-drift,
modulation-driven recentering, exterior, or gauge-degenerate local
alternatives.  The same stratification applies on the scale-collapse
stratum because it depends only on the absolutely continuous chart
variables of
\cref{paper6a:lem:repaired-modulation-coefficients} and on the
carrier-routing definition
\cref{paper7:def:stratified-moving-carrier-routing}, neither of which uses
the single-core hypothesis.  All five rows are closed
adjacent rows: the negative scale-drift row by
\cref{paper6a:thm:thin-negative-drift-closed}; the positive scale-shell
component by
\cref{paper7:cor:windowwise-positive-scale-routing-holds}; modulation-driven
recentering by
\cref{paper7:def:stratified-moving-carrier-routing} together with the
ambient adjacent wrapper
\cref{paper7:cor:ambient-adjacent-wrapper-discharged}; the exterior row by
\cref{paper8:cor:exterior-concentration-exit-discharged}; and the
gauge-degenerate / representation row by the ordered representation tests
\cref{cor:c2-ordered-representation,cor:c2r-ordered-classification}.  This
is the conclusion claimed.

\smallskip\noindent\emph{Case B (window-integrals bounded).}  On every unit
window in the cover, the opposite alternative of
\cref{paper6a:lem:modulation-window-bound-or-exit} holds, so
\(\sup_n\int_0^T(|a_n|+\|b_n\|_\infty)\,d\tau<\infty\).  By the
local modulation-defect routing lemma
\cref{paper6:lem:bounded-modulation-defect-routing}, after passing to a
subsequence either the modulation coefficients have a constant \(L^1\) limit on
the window, the nonconstant defect is perturbative and the represented local
equation has the same terminal residual limit as the sequence of
constant-coefficient branches with the same window-average coefficients, or a
nonzero visible local modulation-defect occupation state is recorded.  In the
first two alternatives the finite-dimensional window averages converge, and the
retained local compactness test gives the corresponding constant-coefficient
autonomous local limit, which is excluded by the autonomous-modulation
alternative
\cref{paper5:thm:scale-collapse-stratification}\textup{(v)}.  In the remaining
nonperturbative alternative the defect is recorded, before autonomous retention, as
modulation-driven recentering, positive scale-shell work, negative scale-drift,
loss of selected active core, or repaired-gauge degeneracy.  These rows are
closed respectively by
\cref{paper7:cor:ambient-adjacent-wrapper-discharged},
\cref{paper7:cor:windowwise-positive-scale-routing-holds},
\cref{paper6a:thm:thin-negative-drift-closed},
\cref{paper6:thm:active-core-loss-closed}, and the ordered representation
tests
\cref{cor:c2-ordered-representation,cor:c2r-ordered-classification}.
Thus the bounded-modulation case also leaves the autonomous-modulation
alternative through a closed adjacent row.

The two cases exhaust the possibilities.  No part of the argument used a
whole-space critical norm \(\|V(\tau)\|_{L^3(\mathbb R^3)}\), a global
profile decomposition, a global Kato smallness condition, a noncanonical
schedule, or a hidden vanishing localized dissipation: only finite-window
\(L^1\) boundedness, weak-\(*\) compactness of normalized defect measures, and
the existing closure theorems for the named adjacent rows are used.
\end{proof}

\begin{lemma}[Canonical identity-cost channel is always available on retained compact branches]
\label{paper6:lem:canonical-identity-cost-available}
\label{paper6:lem:canonical-cost-channel-available}
\label{paper6:lem:canonical-cost-available}
Let \(\omega\) be a retained compact branch that has passed the representation,
pressure-atlas, compactness, and retained-core tests leading to the local
cost-validation layer.  Then the canonical identity-cost validation channel
coming from the suitable local energy inequality is available on that branch.
If this canonical validation fails, the first failed canonical test is one of
the already named local rows in
\cref{paper7:lem:no-unrecorded-cost-escape}.
\end{lemma}

\begin{proof}
On every retained compact branch, the represented equation, repaired-gauge
chart, and local suitable energy inequality are available by
\cref{paper6:thm:representation,paper6a:lem:renormalized-lei}.  Hence the
canonical identity-cost channel is defined using the identity cutoff and the
localized dissipation/carrier package attached to that suitable-energy
inequality.  The ordered validation tests are exactly those of
\cref{paper7:lem:no-unrecorded-cost-escape}.  If one fails before retained
canonical divergence is validated, that failure is by definition the first
failed canonical test and is one of the named local rows listed there.  Thus
the canonical channel is always available once the retained compact branch
reaches the cost-validation layer, and any failure is already part of the
existing first-failure ledger.
\end{proof}

\begin{theorem}[Noncanonical carrier and cost validation route to the canonical channel]
\label{paper6:thm:noncanonical-cost-routing-closed}
Let \(\omega\) be a renormalized Type II solution with active supercritical
scaling on a terminal tail.  Let a candidate continuation lie in a compact
cost-divergence regime through a noncanonical carrier or cost-validation
channel different from the canonical retained identity-cost channel.  Then
exactly one of the following occurs after running the ordered local
validation tests of \cref{paper7:lem:no-unrecorded-cost-escape}:
\begin{enumerate}[label=\textup{(\alph*)}]
\item the candidate is, after the comparable-cost compatibility supplied by
\cref{paper7:lem:comparable-cost-compatibility}, the canonical
identity-cost retained channel up to the terminal one-sided comparison of
\cref{def:c4-cost-comparison}, in which case the cost-divergence exclusion
on the retained stratum is supplied by
\cref{paper7:cor:retained-stratum-cost-exclusion}, hence closed;
\item the first failed test lies in the cost-side list
\cref{paper7:lem:no-unrecorded-cost-escape}\textup{(b)}.  These are the
cost-compatibility, carrier-subidentity, finite fixed-cutoff error,
summable moving-annulus error, carrier-comparison for the validated cost,
modulation-integrability, or compact cost-divergence failures; each such first
failure is a closed adjacent row, routed in the ambient adjacent wrapper
\cref{paper7:cor:ambient-adjacent-wrapper-discharged}.
\end{enumerate}
Consequently the row ``carrier or cost-validation failure outside the
canonical retained identity-cost channel'' is not an independent remaining
terminal exit.
\end{theorem}

\begin{proof}
By \cref{paper7:lem:no-unrecorded-cost-escape}, after running the ordered
local validation tests \textup{(i)}--\textup{(ix)} of that lemma, exactly
one of \textup{(a)} all tests pass with \(\omega\) on the retained
cost-divergence regime, or \textup{(b)} the first failed test is one of the
named local cost-side first failures listed there.

In case \textup{(a)}, the comparable-cost compatibility lemma
\cref{paper7:lem:comparable-cost-compatibility} identifies the validated
noncanonical compact Navier--Stokes cost on the retained stratum with the
canonical identity-cost up to the terminal one-sided comparison of
\cref{def:c4-cost-comparison}.  By
\cref{paper6:lem:canonical-cost-channel-available}, the canonical identity-cost
validation channel is available on the retained compact branch.  Therefore the retained cost-divergence
exclusion supplied by
\cref{paper7:cor:retained-stratum-cost-exclusion} applies, and the candidate
is excluded.

In case \textup{(b)}, the first failed test is by construction one of the
named cost-side adjacent rows enumerated in
\cref{paper7:lem:no-unrecorded-cost-escape}\textup{(b)}.  We show each
such row is closed.

\smallskip\noindent\emph{Class-side first failures.}  Ordered exhaustion
failure, representation failure, and repaired-gauge or scale-degeneracy
failure are the C1/C2 first failures already routed before the
state-decomposition path is reached.  On a physical Type II germ they are not
left as analytic gaps: the covered-entry theorem records them either as
class-incompatible alternatives or as named local state-space rows.  Thus they
are discharged by \cref{paper6:cor:covered-class-entry-discharged} together
with the ordered representation tests
\cref{cor:c2-ordered-representation,cor:c2r-ordered-classification}.

\smallskip\noindent\emph{State-decomposition first failures.}
Multibubble/multicore splitting is closed by \cref{paper6:thm:multibubble}
and \cref{paper6:thm:single-core}; cascade is closed by the finite-cost
scale-collapse theorem \cref{paper6:thm:scale-cost} together with
\cref{paper3:cor:multi-removal}; exterior-tightness failure is closed by
\cref{paper8:cor:exterior-concentration-exit-discharged}; loss of local
windowed \(H^1\) control is closed by \cref{paper6:thm:rough-core}, which
returns to the multibubble/cascade closed row; finite-cost scale-collapse
is closed by \cref{paper6:thm:scale-cost}; the scale-rigid residual is
closed by \cref{paper3:prop:scale-rigidity-discharged}; critical-mass
adjacent labels are closed by
\cref{paper6:thm:active-core-loss-closed}; and modulation-integrability
failure is closed by
\cref{paper6:thm:autonomous-modulation-failure-closed} (and by
\cref{paper6a:lem:modulation-window-bound-or-exit} in the ambient
single-core stratum).

\smallskip\noindent\emph{Cost-side first failures.}  These are the irreducible
noncanonical-cost rows.  The remaining items in
\cref{paper7:lem:no-unrecorded-cost-escape}\textup{(b)} are the cost-side
labels: cost-compatibility failure, carrier-subidentity failure, finite
fixed-cutoff error failure, summable moving-annulus error failure,
carrier-comparison failure for the validated cost, and failure of compact
cost divergence for the validated cost.  These rows are not closed by declaring
the Navier--Stokes solution impossible; they are closed by the ordered
canonical-channel rule: the final retained proof uses only the canonical
identity cost, or an explicitly comparable compact cost, and every failure of a
noncanonical diagnostic is routed back to the canonical local validation layer.

\emph{Cost-compatibility and carrier-comparison failure.}  By definition of
\cref{paper7:lem:canonical-cost-compatibility,paper7:lem:comparable-cost-compatibility},
cost-compatibility failure means that the validated cost used to record
divergence is neither the canonical identity cost nor a comparable cost in
the sense of \cref{def:c4-cost-comparison}.  Such a diagnostic is not used to
retain a Type II branch.  By
\cref{paper6:lem:canonical-cost-channel-available}, the canonical identity-cost
channel from the suitable local energy inequality is available on every
retained compact branch.  Hence the ordered local validation is rerun on that
canonical channel.  If the canonical channel passes, the retained stratum is closed by
\cref{paper7:cor:retained-stratum-cost-exclusion}.  If it fails, the first
failed canonical test is one of the ordered local rows already treated above.
Thus an incomparable noncanonical cost is only a bookkeeping failure of that
diagnostic and does not create a new terminal exit.

\emph{Carrier-subidentity, finite fixed-cutoff error, and summable
moving-annulus error failure.}  By
\cref{paper7:lem:suitable-energy-gives-carrier-subidentities}, suitable
local energy on the repaired-gauge branch supplies the moving carrier
subidentities and the corresponding fixed and moving carrier closure
packages of
\cref{paper7:def:fixed-carrier-closure-package,paper7:def:moving-carrier-closure-package},
with finite-error selection
\cref{paper7:def:finite-error-carrier-selection} and the carrier
completion of
\cref{paper7:cor:finite-error-selection-gives-carrier-completion,paper7:cor:annular-finite-error-selection-gives-carrier-completion}.
If one of these carrier rows fails before the retained carrier is validated,
the failure is not repaired by a global estimate.  The ordered componentwise
carrier decomposition records the first nonsummable local term.  Transition,
pressure, convection, and cutoff terms route to exterior/leakage, critical
tightness failure, or local \(H^1\)-loss; translation terms route to
modulation-driven recentering; positive scale terms route to the positive
scale-shell/scale-rigid row; and negative scale terms route to the thin-drift
row.  These are precisely the alternatives in
\cref{paper7:def:stratified-moving-carrier-routing}, closed by
\cref{paper7:cor:ambient-adjacent-wrapper-discharged},
\cref{paper7:cor:windowwise-positive-scale-routing-holds}, and
\cref{paper6a:thm:thin-negative-drift-closed}.  If the suitable-energy
carrier identity itself is unavailable, the first failed row is the already
routed representation/repaired-gauge or suitable-energy row, not a new
cost-validation exit.

\emph{Failure of compact cost divergence for the validated cost.}  This is
the negation of the entry hypothesis: if the validated cost does not
diverge on the compact terminal cylinder, the candidate is not in the
compact cost-divergence regime under consideration.  The final assembly then
continues through the ordinary retained-state decomposition: finite cost is
handled by \cref{paper6:thm:scale-cost}, while failure before finite-cost
classification is one of the ordered local rows already listed in
\cref{paper7:lem:no-unrecorded-cost-escape}.  Thus nondivergence of a
validated noncanonical cost cannot survive as an independent terminal exit.

\smallskip
The argument uses only the ordered finite-step local validation tests, the
comparable-cost compatibility on compact windows, the suitable
local energy package on the retained stratum, and the canonical-channel
fallback for diagnostics that are not explicitly comparable.  It introduces no whole-space
\(L^3\) bound, no global profile decomposition, and no noncanonical
schedule outside the comparison framework already established by
\cref{paper7:lem:comparable-cost-compatibility,paper7:lem:abstract-cost-exclusion}.
\end{proof}

\begin{lemma}[Pressure-only critical mass routes through local pressure atlas]
\label{paper6:lem:pressure-only-critical-mass-routing}
Let \(B_R\Subset B_{R_1}\Subset B_{R_2}\) be nested balls in a retained
terminal frame on a compact time interval \(I\), and fix \(1\le q<\infty\).
Suppose a pressure oscillation modulo functions of time persists in
\(L^q_\tau L^{3/2}_y(B_R)\), while the ordered terminal selection has not
retained a velocity active core for that candidate.  With the local pressure
decomposition
\[
  P_n=P_n^{\mathrm{loc}}+H_{n,R_1}+c_{n,R_1}(\tau)
  \qquad\text{on }B_{R_1},
\]
where \(P_n^{\mathrm{loc}}\) is generated by
\(\zeta(V_n\otimes V_n)\) on \(B_{R_2}\) and \(H_{n,R_1}\) is mean-zero
harmonic on \(B_{R_1}\), one of the following ordered alternatives occurs:
\begin{enumerate}[label=\textup{(\roman*)}]
\item the localized velocity source
  \(\zeta(V_n\otimes V_n)\) has a nonzero compact-window
  \(L^q_\tau L^{3/2}_y\) limsup; the local CKN test on the selected compact
  subcylinders then either absorbs the source as regular, or records a
  comparable retained class, a multibubble/multicore label, a cascade label, a
  rough-core label, or an exterior/separated label;
\item the localized velocity source tends to zero in
  \(L^q_\tau L^{3/2}_y(B_{R_2})\) for the fixed exponent \(q\), and any
  remaining pressure oscillation is carried by the harmonic part
  \(H_{n,R_1}\), which is routed by
  \cref{paper8:lem:harmonic-cross-pressure-routing} to the ordered
  pressure/exterior row;
\item the pressure reconstruction or gauge needed for the displayed
  decomposition fails, which is itself the ordered pressure/repaired-gauge row.
\end{enumerate}
Thus a pressure-only critical-mass signal is never closed by assuming
vanishing velocity mass; it is routed by the local source/harmonic dichotomy.
\end{lemma}

\begin{proof}
The pressure decomposition is one of the ordered local pressure-atlas tests.  If
that test fails, item \textup{(iii)} is recorded.  Assume it holds.
\par
If
\[
  \limsup_{n\to\infty}
  \|\zeta(V_n\otimes V_n)\|_{L^q_\tau L^{3/2}_y(B_{R_2})}>0,
\]
then the compact-window velocity stress is nonzero.  Since
\(\|V_n\otimes V_n\|_{L^{3/2}(B)}=\|V_n\|_{L^3(B)}^2\) on each spatial ball,
a finite cover of \(B_{R_2}\) and a subsequence select a compact subcylinder
with positive local velocity critical mass in the corresponding spacetime norm.
Run the local CKN test on that subcylinder.  If all smaller admissible
cylinders are below the CKN threshold, the source is regular on the compact
window and is absorbed into the regular pressure atlas.  Otherwise a local CKN
carrier is selected, and the ordered terminal critical-mass routing
\cref{paper6a:lem:terminal-local-critical-mass-routing} records either a
comparable retained class or one of the adjacent active labels listed in item
\textup{(i)}.  This uses only the compact-window stress norm, local CKN
testing, and finite covering.
\par
If the localized source tends to zero in
\(L^q_\tau L^{3/2}_y(B_{R_2})\), local Calder\'on--Zygmund estimates give
\[
  P_n^{\mathrm{loc}}\to0
  \quad\text{in }L^q_\tau L^{3/2}_y(B_{R_1})
\]
modulo functions of time.  Hence any persistent pressure oscillation on
\(B_R\) must lie in the harmonic remainder.  Applying
\cref{paper8:lem:harmonic-cross-pressure-routing} with \(p=3/2\) and the fixed
time exponent \(q\) routes that nonvanishing harmonic remainder to the
pressure/exterior row.  No global
pressure compactness or lower semicontinuity argument is used.
\end{proof}

\begin{theorem}[Loss of retained active core is a closed adjacent row]
\label{paper6:thm:active-core-loss-closed}
Let \((u,p,x_0,T)\) be a physical local Type II singular germ that has
covered entry into the local state-space path, and let a terminal active
frame be selected from its local state-space decomposition.  If on a fixed
compact terminal cylinder the candidate fails to retain an active core in
the sense of \cref{paper6a:lem:terminal-local-critical-mass-routing}, then
the candidate enters one of the following closed adjacent rows:
\begin{enumerate}[label=\textup{(\roman*)}]
\item the retained spacetime activity-floor failure, in which case the candidate is
not a retained terminal frame by the terminal selection rule and the
state-space partition records the next available retained frame; if no velocity
retained frame remains, the original CKN signal is pressure/exterior,
rough-core, or regular, and is routed by the pressure-only and exterior
closures;
\item the local \(L^3\)-upper-bound failure, which by
\cref{paper6a:lem:terminal-local-critical-mass-routing}\textup{(iii)} is
one of: a comparable finite-mass class absorbed into the retained frame by
\cref{paper6a:lem:retained-frame-maximality}, a multibubble or multicore
splitting closed by \cref{paper6:thm:multibubble} and
\cref{paper6:thm:single-core}, a cascade closed by the finite-cost
scale-collapse theorem \cref{paper6:thm:scale-cost} together with the
cascade absorption used in \cref{paper3:cor:multi-removal}, a rough-core
label closed by \cref{paper6:thm:rough-core}, or an exterior/separated
label closed by
\cref{paper8:cor:exterior-concentration-exit-discharged};
\item the pressure-only critical-mass row, in which case the local pressure
atlas routes the signal by
\cref{paper6:lem:pressure-only-critical-mass-routing}: the local
Calder\'on--Zygmund source is either regular on the compact window, produces a
velocity-supported terminal label, or the harmonic remainder is routed into the
pressure/exterior-leakage alternative, closed by
\cref{paper8:cor:exterior-concentration-exit-discharged}.
\end{enumerate}
Consequently the row ``loss of a retained active core, including
pressure-only critical-mass detection'' is not an independent remaining
terminal exit.
\end{theorem}

\begin{proof}
Apply
\cref{paper6a:lem:terminal-local-critical-mass-routing} on the fixed
compact terminal cylinder \(B_R\times I\).  After passing to a subsequence,
exactly one of its three alternatives \textup{(i)}--\textup{(iii)} holds.

If alternative \textup{(i)} of that lemma holds, the candidate has positive
finite local critical \(L^3\) mass on the cylinder; this is the retained
case and not the loss-of-core case under consideration.  If alternative
\textup{(ii)} holds, the lower active-mass floor of
\cref{paper8:thm:active-profile-mass-floor} fails, so the candidate frame
is not retained by the terminal selection rule.  The selection rule then moves
to the next available velocity frame.  If there is no such frame, the positive
CKN concentration of the physical germ cannot be a retained velocity core; it
is either pressure-only, exterior/separated, rough-core, or below the CKN
threshold on all smaller cylinders.  The first three are routed by
\cref{paper6:lem:pressure-only-critical-mass-routing},
\cref{paper8:cor:exterior-concentration-exit-discharged}, and
\cref{paper6:thm:rough-core}; the last contradicts singular CKN concentration
by \cref{paper8:thm:terminal-ckn-regularity}.  This is item \textup{(i)} of
the present conclusion.  If alternative \textup{(iii)} holds, the local
\(L^3\) upper bound fails on the compact cylinder.  By
\cref{paper6a:lem:terminal-local-critical-mass-routing}\textup{(iii)}, the
ordered local state space then records a comparable finite-mass class
absorbed into the retained component, a multibubble or multicore label, a
cascade label, a rough-core label, or an exterior/separated label.  Each
such label is closed by the theorems cited in item \textup{(ii)} of the
present conclusion: the rough-core label is closed by
\cref{paper6:thm:rough-core}, which forces failure of compact-cylinder CKN
control on every relevant enclosing cylinder and therefore returns to the
multibubble/cascade closed row; the exterior/separated label is closed by
\cref{paper8:cor:exterior-concentration-exit-discharged}.

It remains to handle the pressure-only subcase, in which the active core
fails because the critical mass detected on the cylinder is carried by the
pressure rather than by a velocity profile, i.e.\ the local velocity
\(L^3\)-mass remains below the active floor while a pressure oscillation
modulo a function of time persists on \(B_R\) along a subsequence.  Apply the
local pressure-only routing lemma
\cref{paper6:lem:pressure-only-critical-mass-routing} with the time exponent
in which the pressure oscillation persists.  If the localized
Calder\'on--Zygmund source has a nonzero compact-window stress limsup, the
finite-cover and CKN test in that lemma either absorbs the source as regular,
or records a velocity-supported terminal label; the latter is either absorbed
into a comparable retained class by
\cref{paper6a:lem:retained-frame-maximality} or routed to one of the closed
multibubble/cascade, rough-core, or exterior/separated rows.  If the localized
source tends to zero, the local pressure part vanishes by
Calder\'on--Zygmund estimates and any remaining pressure oscillation is harmonic;
\cref{paper8:lem:harmonic-cross-pressure-routing} routes it to the ordered
pressure/exterior-leakage alternative, closed by
\cref{paper8:cor:exterior-concentration-exit-discharged}.  If pressure
reconstruction fails, that is the ordered pressure/repaired-gauge row.  This is
item \textup{(iii)} of the present conclusion.

Each step uses only compact terminal windows, the local active-mass floor,
the ordered local state-space partition, and the local pressure
reconstruction with harmonic remainder.  No global \(L^3\) bound, no global
profile decomposition, no global Kato smallness, and no pressure
compactness by weak lower semicontinuity is used.
\end{proof}

\begin{lemma}[Acyclic terminal ledger]
\label{paper6:lem:acyclic-terminal-ledger}
The routing layer used before the remaining-exit closure subsection is
acyclic.  More precisely,
\cref{paper6a:thm:terminal-adjacent-label-routing,paper6a:thm:local-residual-exhaustion,paper6:thm:ordered-local-state-space-exhaustion,paper8:lem:nested-pressure-only-well-founded,paper9:lem:critical-tightness-failure-row,paper10:lem:bounded-window-reselection-well-founded}
produce only named local rows or the retained comparable class.  They do not
use the later closure theorems
\cref{paper8:lem:diffuse-pressure-only-not-terminal,paper9:lem:diffuse-critical-tail-row-closed,paper10:lem:diffuse-reselection-tree-row-closed}
to create those rows.

Consequently the terminal routing package first records a named row, and only
later, in the remaining-exit closure layer, invokes the theorem that closes
that row.
\end{lemma}

\begin{proof}
By construction,
\cref{paper8:lem:nested-pressure-only-well-founded,paper9:lem:critical-tightness-failure-row,paper10:lem:bounded-window-reselection-well-founded}
stop at the diffuse pressure-only, diffuse critical-tail, and diffuse
reselection-tree rows by name.  They do not appeal to the later theorems that
eliminate those rows.  Theorem
\cref{paper6a:thm:terminal-adjacent-label-routing} is a routing theorem only,
and \cref{paper6a:thm:local-residual-exhaustion} routes a visible residual
defect to a named adjacent row or to the retained comparable class.  The proof
of \cref{paper6:thm:ordered-local-state-space-exhaustion} uses these routing
outputs only as named ledger rows and explicitly defers their closure to the
remaining-exit closure layer.  This is exactly the asserted acyclic ordering.
\end{proof}

\begin{corollary}[All four remaining named exit rows are closed]
\label{paper6:cor:remaining-named-exits-closed}
The four rows listed in \cref{paper6:def:remaining-named-exits} are all
closed adjacent rows: the bounded-window row by
\cref{paper3:lem:bounded-window-vmo-pressure,paper10:lem:bounded-window-reselection-well-founded,paper6:thm:bounded-selected-window-physical-routing};
the autonomous-modulation row by
\cref{paper6:thm:autonomous-modulation-failure-closed}; the noncanonical
carrier/cost-validation row by
\cref{paper6:thm:noncanonical-cost-routing-closed}; and the loss-of-active-core
row, including pressure-only critical-mass detection, by
\cref{paper6:thm:active-core-loss-closed,paper8:lem:nested-pressure-only-well-founded,paper8:lem:diffuse-pressure-only-not-terminal,paper9:lem:critical-tightness-failure-row,paper9:lem:diffuse-critical-tail-row-closed}.
\end{corollary}

\begin{proof}
Direct application of the bounded-window, autonomous-modulation, noncanonical
cost-routing, active-core-loss, critical-tail, and nested pressure-only
closure theorems cited in the statement.
\end{proof}

\begin{definition}[Remaining named exit rows]
\label{paper6:def:remaining-named-exits}
The present retained-state assembly does not count a first-failure row as
closed merely because it has been named.  The ledger rows outside the retained
closure theorem before the closure subsection is applied are:
\begin{enumerate}[label=\textup{(\roman*)}]
\item a bounded selected-window limit producing a bounded local profile on a
fixed physical window, before the physical routing theorem is applied;
its closure uses
\cref{paper3:lem:bounded-window-vmo-pressure,paper10:lem:bounded-window-reselection-well-founded};
\item autonomous-modulation failure or modulation-driven recentering;
\item carrier or cost-validation failure outside the canonical retained
identity-cost channel, including noncanonical cost-comparison failures before
the canonical channel of \cref{paper6:lem:canonical-cost-channel-available} is
invoked;
\item loss of a retained active core, including pressure-only critical-mass
detection, unless the pressure/exterior routing theorem places it in one of the
closed structural rows listed below.  This row also absorbs critical-tightness
failure and diffuse critical-tail behavior through
\cref{paper9:lem:critical-tightness-failure-row,paper9:lem:diffuse-critical-tail-row-closed},
and nested exterior pressure-only behavior through
\cref{paper8:lem:nested-pressure-only-well-founded,paper8:lem:diffuse-pressure-only-not-terminal}.
\end{enumerate}
The closed rows are not in this list.  They are finite scale cost, finite
absolute scale cost, thin negative drift as an independent terminal row,
compactness failure on selected scale-collapse windows as an independent
terminal row,
compact single-core vanishing dissipation, multibubble, cascade,
gauge-degenerate behavior, scale-rigid residual behavior, the retained compact
scale-collapse state, exterior concentration and noncompact exterior
concentration as independent terminal rows, and the retained canonical
compact-cost branch.  Thin
negative drift is closed by
\cref{paper6a:thm:thin-negative-drift-closed}; selected compactness failure is
routed by \cref{paper5:thm:selected-compactness-failure-routed}; exterior
concentration is routed by
\cref{paper8:cor:exterior-concentration-exit-discharged}; the other closed
rows are
closed by the theorems cited in \cref{paper6:prop:state-elimination} and in the
ambient compact-cost package.

By \cref{paper6:lem:acyclic-terminal-ledger}, these rows are recorded in the
terminal routing layer only by name.  Their actual closure is invoked only in
the present subsection.

\medskip
The four rows \textup{(i)}, \textup{(ii)}, \textup{(iii)}, \textup{(iv)} listed above are
closed adjacent rows, by the closure theorems of
\cref{paper6:subsec:remaining-exit-closures}: the bounded-window row by
\cref{paper3:lem:bounded-window-vmo-pressure,paper10:lem:bounded-window-reselection-well-founded,paper6:thm:bounded-selected-window-physical-routing}; the autonomous-modulation row
by \cref{paper6:thm:autonomous-modulation-failure-closed}; the
noncanonical carrier/cost-validation row by
\cref{paper6:thm:noncanonical-cost-routing-closed}; and the
loss-of-active-core row, including pressure-only critical-mass
detection, by \cref{paper6:thm:active-core-loss-closed} together with the
critical-tail and nested pressure-only closures
\cref{paper9:lem:critical-tightness-failure-row,paper9:lem:diffuse-critical-tail-row-closed,paper8:lem:nested-pressure-only-well-founded,paper8:lem:diffuse-pressure-only-not-terminal}.  This is
summarised in \cref{paper6:cor:remaining-named-exits-closed}.  The list is
the controlling ledger for
\cref{paper6:thm:unconditional-no-typeII}; the unconditional upgrade is recorded in
\cref{paper6:thm:unconditional-no-typeII-discharged}.
\end{definition}

\begin{theorem}[Retained Type II closure with named-exit ledger]
\label{paper6:thm:unconditional-no-typeII}
Let \((u,p)\) be a suitable weak solution on
\(\mathbb R^3\times(T-\delta,T)\), and let
\((u,p,x_0,T)\) be a physical local Type II singular germ in the sense of
\cref{paper6:def:physical-local-typeII-germ}.  Then either the germ enters one
of the remaining named exit rows in
\cref{paper6:def:remaining-named-exits}, or its retained local branch is
impossible.  Consequently, together with
\cref{paper6:cor:remaining-named-exits-closed}, no physical local Type II
singular germ remains.
\end{theorem}

\begin{proof}
Let \((u,p,x_0,T)\) be a physical local Type II singular germ.  By
\cref{paper6:cor:covered-class-entry-discharged}, the germ either enters the
local state-decomposition path or exits earlier through a class-incompatible
or named non-retained local alternative.
\par
A class-incompatible exit contradicts the definition of physical local Type II
germ and the admissibility theorem
\cref{paper6:lem:physical-typeII-germ-admissible}.  If a named non-retained
exit occurs, then either it is one of the closed local rows cited in
\cref{paper6:prop:state-elimination} and
\cref{paper7:cor:ambient-adjacent-wrapper-discharged}, or it belongs to the
remaining named-exit ledger
\cref{paper6:def:remaining-named-exits}.  In the latter case the first
alternative in the statement holds.
\par
It remains to consider the case in which the branch reaches the retained
terminal path without entering any remaining named exit.  The expanded terminal
wrapper
\cref{paper6a:cor:expanded-terminal-wrapper-discharged} supplies the expanded
terminal data of \cref{paper6a:def:expanded-terminal-data}, and
\cref{paper6a:thm:expanded-terminal-typeII} excludes the branch.
\par
By \cref{paper6:lem:acyclic-terminal-ledger}, all uses of critical-tightness,
pressure-only, and bounded-window reselection inside the terminal routing layer
occur only at the row-naming stage.  Their closure is invoked only now through
\cref{paper6:cor:remaining-named-exits-closed}.
\par
Equivalently, starting from a physical local Type II germ, the ordered
state-space exhaustion produces one first retained row or one first failed row.
The strong gauge-chart entry test prevents the gauge implicit-function theorem
from being applied outside its \(C^1\) chart.  Bounded selected-window limits
are locally regular or well-foundedly reselected by
\cref{paper3:lem:bounded-window-vmo-pressure,paper10:lem:bounded-window-reselection-well-founded}.
Failure of critical tightness is routed by
\cref{paper9:lem:critical-tightness-failure-row}, and the diffuse tail case is
closed by \cref{paper9:lem:diffuse-critical-tail-row-closed}.  Exterior
pressure-only nesting is well-founded by
\cref{paper8:lem:nested-pressure-only-well-founded}, and diffuse pressure-only
defects are excluded by
\cref{paper8:lem:diffuse-pressure-only-not-terminal}.  All remaining rows are
the previously closed active, exterior, multibubble, cascade, rough-core,
pressure-gauge, cost, and terminal compact rows.  Hence no row can contain a
physical local Type II germ.
\end{proof}

\begin{theorem}[No physical local Type II germ]
\label{paper6:thm:unconditional-no-typeII-discharged}
Let \((u,p)\) be a suitable weak solution on
\(\mathbb R^3\times(T-\delta,T)\).  No physical local Type II singular germ
\((u,p,x_0,T)\) in the sense of
\cref{paper6:def:physical-local-typeII-germ} exists.
\end{theorem}

\begin{proof}
By \cref{paper6:thm:unconditional-no-typeII}, every physical local Type II
singular germ either enters one of the rows in
\cref{paper6:def:remaining-named-exits} or is excluded on the retained
branch.  By \cref{paper6:cor:remaining-named-exits-closed}, every row in
\cref{paper6:def:remaining-named-exits} is a closed adjacent row: it is routed
by its local closure theorem either to a closed retained-state alternative, to
a structural row already excluded, or to a regular/removable row.  Hence no
physical local Type II singular germ remains.
\end{proof}

\subsection{Assembly Theorem}

\begin{theorem}[Local Type II exclusion]
\label{paper6:thm:local-typeII-exclusion}
Let \((u,p)\) be a suitable weak solution on
\(\mathbb R^3\times(T-\delta,T)\).  No point \(x_0\in\Sigma(T)\) lies in the
local Type II alternative.
\end{theorem}

\begin{proof}
\Cref{paper6:thm:unconditional-no-typeII-discharged} excludes every physical
local Type II singular germ.  If a point \(x_0\in\Sigma(T)\) lay in the local
Type II alternative, then it would satisfy the three clauses of
\cref{paper6:def:physical-local-typeII-germ}: singularity gives
\(x_0\in\Sigma(T)\), the local Type I bound fails by definition of the Type II
alternative, and the theorem is stated for interior cylinders of
\(\mathbb R^3\times(T-\delta,T)\).  This is a physical local Type II singular
germ, contradicting
\cref{paper6:thm:unconditional-no-typeII-discharged}.
\end{proof}

\begin{corollary}[Final local dichotomy after Type I separation]
\label{paper6:cor:final-local-dichotomy}
Every singular point at time \(T\) belongs to the local Type I alternative.
Equivalently, once the Type I case is covered by the Type I part of the theory,
there is no remaining local Type II singular branch.
\end{corollary}

\begin{proof}
Let \(x_0\in\Sigma(T)\).  By \cref{paper6:thm:paper0}, \(x_0\) carries a
positive CKN concentration sequence and lies either in the local Type I
alternative or in the local Type II alternative.
\Cref{paper6:thm:local-typeII-exclusion} excludes the latter.
\end{proof}

\clearpage
\section{Conclusion}

The paper gives a local branchwise exclusion mechanism for Type II
Navier--Stokes concentration.  Beginning
with a positive Caffarelli--Kohn--Nirenberg concentration sequence, the argument
passes to repaired-gauge variables around a retained core and keeps the analysis
inside local PDE estimates: suitable energy inequalities, pressure
reconstruction, compactness on fixed cylinders, Caccioppoli estimates,
profile decoupling, and localized scale-cost identities.

The compact single-core mechanism is the central obstruction.  A retained core
with positive local velocity density and vanishing localized dissipation has
only two compact limits after passing to selected windows.  The zero limit
contradicts the persistence of the retained velocity core by strong \(L^3\)
convergence, while a nonzero spatially constant limit gives a bounded
selected-window profile and hence exits the Type II regime.  Pressure-only
harmonic modes are routed through the pressure/exterior rows and are not used
to close the compact single-core velocity row.  The rest
of the proof shows that the remaining alternatives to this compact
mechanism are also accounted for: rough cores are reduced to compact-cylinder
control, active multibubbles and cascades are separated by the profile
analysis, and finite-cost scale collapse is incompatible with a nonvanishing
retained core.

The scale-collapse analysis is intentionally explicit about its routing.  Finite
scale-collapsing cost, and hence finite absolute scale cost in the covered
setting, are excluded within the renormalized setting.  Thick autonomous
windows produce nontrivial autonomous limits only after the relevant local
compactness and modulation rows have passed.  The retained compact autonomous
branch is not closed by silently invoking a global Liouville theorem: it is
routed through
\cref{paper5:thm:compact-scale-collapse-rigidity-discharged} into the
scale-rigid discharge.  If a retained input fails, the failure is one of the
named modulation, compactness, exterior, rough-core, or multicore local rows;
the compactness, exterior, and thin-drift boundary faces are routed by their
local closure theorems, including
\cref{paper5:thm:selected-compactness-failure-routed},
\cref{paper8:cor:exterior-concentration-exit-discharged}, and
\cref{paper6a:thm:thin-negative-drift-closed}.

Thus every physical local Type II concentration branch is either eliminated on
the retained path or routed through one of the named local rows.  The remaining
named rows are closed in
\cref{paper6:cor:remaining-named-exits-closed}, and the final assembly
\cref{paper6:thm:unconditional-no-typeII-discharged} excludes physical local
Type II singular germs.  The conclusion remains local and branchwise: it does
not import a whole-space critical estimate, a global profile decomposition,
global Kato smallness, or hidden vanishing localized dissipation.


\clearpage
\appendix
\section{PDE Interface Definitions Used in the Local Type II Proof}
\label{sec:c1-c5-pde-translation}

This appendix records, in PDE notation, the local Type II classification
interfaces used earlier in the proof.  The statements are interface definitions
and local consequences for admissible Type II branches; they do not provide a
separate post-assembly proof path.  In the final assembly, their inputs are
verified on the retained local path or routed to closed named rows.

\subsection{Admissible Type II branches and exhaustion}

\begin{definition}[Admissible local Type II class]
\label{def:c1-admissible-typeII-class}
Let \(\mathcal U_{\mathrm{II}}^{NS}\) be the class of admissible terminal
branches \((u,p,T^*)\) such that:
\begin{enumerate}[label=\textup{(\roman*)}]
\item \((u,p)\) is a three-dimensional incompressible Navier--Stokes solution
on a terminal interval \([0,T^*)\) in the analytic class under consideration;
\item \(T^*<\infty\) is the terminal time of the branch;
\item the branch is not in the local Type I class associated with the
scale-invariant Type I criterion;
\item the branch is a Type II branch rather than a continuation, Type I,
boundary, or non-Navier--Stokes-domain alternative.
\end{enumerate}
Membership in \(\mathcal U_{\mathrm{II}}^{NS}\) specifies the class under
analysis; no nonemptiness assertion is made.
\end{definition}

\begin{definition}[Type II analytic data]
\label{def:c1-typeII-analytic-data}
A branch \(\omega=(u,p,T^*)\in\mathcal U_{\mathrm{II}}^{NS}\) has the Type II
analytic data if the following three analytic assertions hold:
\begin{enumerate}[label=\textup{(\roman*)}]
\item a concentration profile or blow-up germ modulo the Navier--Stokes
symmetry group;
\item a supercritical scale alternative, meaning that the branch enters the Type II
scale alternative rather than the Type I scale alternative;
\item local profile data in the Navier--Stokes profile class used
for the subsequent compactness and representation arguments.
\end{enumerate}
The analytic data are \emph{strictly exhaustive} if every admissible Type II
branch has these three properties.  They are \emph{exhaustive by ordered local
routing} if every branch which fails one of the three properties is assigned,
at its first failed C1 test, to one of the ordered failures in
\cref{def:c1-exhaustion-failures} and that failure is then routed to a
class-incompatible alternative or to a named non-retained local state-space
exit.  On a retained subbranch, ordered local-routing exhaustion reduces to
strict exhaustion.
\end{definition}

\begin{definition}[Ordered exhaustion failures]
\label{def:c1-exhaustion-failures}
For an admissible Type II branch, the ordered failures of the Type II analytic
data are:
\begin{enumerate}[label=\textup{(\roman*)}]
\item concentration extraction fails;
\item concentration extraction succeeds but the supercritical scale alternative
does not hold;
\item concentration extraction and the scale alternative hold but the local
profile data are insufficient for the subsequent compactness and representation
arguments;
\item the admissible Type II branch lies outside the concentration, scale, and
profile alternatives under consideration.
\end{enumerate}
The first applicable failure in this order is the exhaustion alternative.
\end{definition}

\begin{theorem}[Type II branch exhaustion]
\label{thm:c1-typeII-branch-exhaustion}
Assume the Type II analytic data of \cref{def:c1-typeII-analytic-data} are
exhaustive by ordered local routing on \(\mathcal U_{\mathrm{II}}^{NS}\), and
that each ordered failure is assigned to a class-incompatible alternative or to
a named non-retained local state-space exit.  Then every
\(\omega\in\mathcal U_{\mathrm{II}}^{NS}\) either has the concentration
profile, the supercritical scale alternative, and the local profile data, or is
already in one of those named exits.  In particular, every retained admissible
Type II branch has the three analytic data.
\end{theorem}

\begin{proof}
Fix \(\omega=(u,p,T^*)\in\mathcal U_{\mathrm{II}}^{NS}\).  By ordered local
routing, exactly one of two possibilities occurs.  If no ordered failure occurs,
then the concentration-extraction, scale, and profile components all hold, so
\(\omega\) has the three Type II analytic data.  If an ordered failure occurs,
the first such failure is, by hypothesis, assigned to a class-incompatible
alternative or to a named non-retained local state-space exit.  Therefore a
retained admissible branch cannot be in the failure case, and must have the
three analytic data.
\end{proof}

\begin{corollary}[Ordered exhaustion classification]
\label{cor:c1-ordered-exhaustion}
For each \(\omega\in\mathcal U_{\mathrm{II}}^{NS}\), exactly one ordered
alternative occurs: the Type II analytic data hold, or one of the four failures
in \cref{def:c1-exhaustion-failures} occurs.  The data are strictly exhaustive
on \(\mathcal U_{\mathrm{II}}^{NS}\) if and only if the first alternative holds
for every \(\omega\in\mathcal U_{\mathrm{II}}^{NS}\).  They are exhaustive by
ordered local routing if every failure alternative is assigned to a
class-incompatible alternative or to a named non-retained local state-space
exit; on the retained subbranch this is equivalent to the first alternative.
\end{corollary}

\begin{proof}
Test the four analytic requirements in the stated order.  If concentration
extraction, the scale alternative, and the profile condition all hold, the
branch has the Type II analytic data.  If a test fails, the first failed test gives the
corresponding ordered failure.  The alternatives are disjoint by the first
failure convention and exhaustive by construction.  The final equivalence is
the definition of strict exhaustion.  The ordered-routing statement follows by
recording the target assigned to each first failure.  Once class-incompatible
and named non-retained exits have been removed, a retained branch cannot lie in
any failure alternative, so it lies in the data alternative.
\end{proof}

\subsection{Repaired-gauge representation from Type II analytic data}

\begin{definition}[Raw Type II chart]
\label{def:c2-raw-chart}
A raw Type II chart consists of
\[
  (u,p,T^*,x_c,\lambda,\tau,V,P,\mathcal I),
  \qquad \mathcal I=(t_0,T^*),
\]
where \(u,p\) solve
\[
  \partial_tu+(u\cdot\nabla)u+\nabla p=\nu\Delta u,\qquad
  \nabla\cdot u=0,
\]
\(\lambda(t)>0\), \(x_c(t)\in\mathbb R^3\), \(x_c,\lambda\in AC_{\rm loc}\),
\(\lambda(t)\to0\) as \(t\uparrow T^*\), and
\[
  \frac{d\tau}{dt}=\lambda(t)^{-2},\qquad
  \tau(t)\to\infty\quad\text{as }t\uparrow T^*.
\]
The renormalized variables are
\[
  y=\frac{x-x_c(t)}{\lambda(t)},\qquad
  u(x,t)=\lambda(t)^{-1}V(y,\tau(t)),\qquad
  p(x,t)=\lambda(t)^{-2}P(y,\tau(t)),
\]
with enough regularity to justify the distributional chain rule on compact
\(y\)-sets.
\end{definition}

\begin{lemma}[Raw renormalized equation]
\label{lem:c2-raw-renormalized-equation}
For a raw Type II chart, define
\[
  a_{\rm raw}(\tau(t))=\lambda(t)\lambda_t(t),\qquad
  b_{\rm raw}(\tau(t))=\lambda(t)x_c'(t).
\]
Then \(V,P\) satisfy, distributionally on compact \(y\)-sets,
\[
  \partial_\tau V+(V\cdot\nabla)V+\nabla P
  =
  \nu\Delta V
  +a_{\rm raw}(\tau)(V+y\cdot\nabla V)
  +b_{\rm raw}(\tau)\cdot\nabla V,
  \qquad \nabla\cdot V=0 .
\]
\end{lemma}

\begin{proof}
Write
\[
u(x,t)=\lambda(t)^{-1}V(y,\tau(t)),\qquad
y=\frac{x-x_c(t)}{\lambda(t)},\qquad \tau_t=\lambda^{-2}.
\]
Then
\[
  \nabla_xu=\lambda^{-2}\nabla_yV,\qquad
  \Delta_xu=\lambda^{-3}\Delta_yV,\qquad
  (u\cdot\nabla_x)u=\lambda^{-3}(V\cdot\nabla_y)V.
\]
Since
\[
  y_t=-\frac{x_c'(t)}{\lambda(t)}
      -\frac{\lambda_t(t)}{\lambda(t)}y,
\]
one has
\[
  \partial_tu
  =
  \lambda^{-3}\partial_\tau V
  -\lambda_t\lambda^{-2}(V+y\cdot\nabla V)
  -\lambda^{-2}x_c'(t)\cdot\nabla V .
\]
Moreover \(\nabla_xp=\lambda^{-3}\nabla_yP\).  Substitution in
Navier--Stokes and multiplication by \(\lambda^3\) give
\[
  \partial_\tau V-\lambda\lambda_t(V+y\cdot\nabla V)
  -\lambda x_c'(t)\cdot\nabla V
  +(V\cdot\nabla)V+\nabla P
  =\nu\Delta V .
\]
Moving the modulation terms to the right-hand side gives the equation with
\(a_{\rm raw}=\lambda\lambda_t\) and \(b_{\rm raw}=\lambda x_c'\).  Finally,
\(\nabla_x\cdot u=\lambda^{-2}\nabla_y\cdot V\), so incompressibility of \(u\)
implies \(\nabla_y\cdot V=0\).
\end{proof}

\begin{definition}[Repaired-gauge realization]
\label{def:c2-repaired-gauge-realization}
Let \(0<p<3\), \(\Theta_0>0\), and let \(\psi_R\) be a smooth compactly
supported centering cutoff.  A raw Type II chart admits a repaired-gauge
realization if, after an admissible time-dependent scale and translation
correction, the profile satisfies
\[
  G_{\rm sc}(V)=\int_{\mathbb R^3}|y|^{-p}|V(y)|^3\,dy-\Theta_0=0,
\]
and
\[
  G_j(V)=\int_{\mathbb R^3}y_j|V(y)|^2\psi_R(y)\,dy=0,
  \qquad j=1,2,3,
\]
for a.e. renormalized time.  The corrected profile belongs to the admissible
gauge class \(\mathcal A_p\), has the regularity needed to differentiate these
functionals, and all time-dependent symmetry corrections have been absorbed
into final modulation coefficients \(a,b\).
\end{definition}

\begin{definition}[Repaired-gauge representation]
\label{def:c2-representation}
A Type II branch has a repaired-gauge representation if it has:
\begin{enumerate}[label=\textup{(\roman*)}]
\item a raw Type II chart;
\item a repaired-gauge realization;
\item pressure reconstruction modulo functions of time;
\item final modulation coefficients \(a,b\) appearing in the repaired-gauge
equation.
\end{enumerate}
The representation consists of the tuple
\[
  (V,P,a,b,G_{\rm sc},G_1,G_2,G_3,\tau_0).
\]
\end{definition}

\begin{theorem}[Representation implication]
\label{thm:c2-representation-implication}
Assume a Type II profile branch has a raw Type II chart, a repaired-gauge
realization, pressure reconstruction modulo functions of time, and final
modulation coefficients.  Then the branch has a repaired-gauge representation
in the sense of \cref{def:c2-representation}.  In particular,
\((V,P,a,b)\) satisfies the PDE class required by the compact Type II
cost-divergence criterion.
\end{theorem}

\begin{proof}
The raw chart gives \(V,P\) and, by \cref{lem:c2-raw-renormalized-equation},
the renormalized Navier--Stokes equation with raw modulation coefficients.  The
repaired-gauge realization places the profile on the scale and centering gauge
surface and absorbs the time-dependent symmetry correction into updated
coefficients.  The pressure reconstruction identifies the pressure associated
with \(V\), modulo the usual functions of time.  The modulation data identify
the final coefficients \(a,b\).  These are the components of
\cref{def:c2-representation}.
\end{proof}

\begin{definition}[Ordered representation failures]
\label{def:c2-representation-failures}
If a repaired-gauge representation is not obtained, the ordered failures are:
\begin{enumerate}[label=\textup{(\roman*)}]
\item no raw orbit is supplied by the concentration and profile data;
\item the raw orbit exists but the repaired gauge cannot be realized;
\item pressure reconstruction modulo functions of time is unavailable;
\item the final modulation coefficients are unavailable.
\end{enumerate}
The first applicable failure is the representation alternative.
\end{definition}

\begin{corollary}[Ordered representation classification]
\label{cor:c2-ordered-representation}
Every Type II profile branch in the admissible class has exactly one ordered
alternative: a repaired-gauge representation, or one of the four failures in
\cref{def:c2-representation-failures}.
\end{corollary}

\begin{proof}
Verify the four representation inputs in order.  If all are present,
\cref{thm:c2-representation-implication} gives the repaired-gauge
representation.  If not, the first missing input gives the corresponding
ordered failure.  This is exhaustive and disjoint by first-failure ordering.
\end{proof}

\subsection{Minimal chart and gauge realization}

\begin{definition}[Minimal representation data]
\label{def:c2r-minimal-representation-data}
The minimal data for a Type II branch satisfying the analytic data consist of:
\begin{enumerate}[label=\textup{(\roman*)}]
\item an absolutely continuous concentration chart
\[
  (u,p,T^*,x_c,\lambda,\rho,U,Q,\mathcal I),
\]
where \(\mathcal I=(t_0,T^*)\), \(\lambda(t)>0\), \(\lambda(t)\to0\),
\(\rho_t=\lambda^{-2}\), and
\[
  u(x,t)=\lambda(t)^{-1}U(y,\rho(t)),\qquad
  p(x,t)=\lambda(t)^{-2}Q(y,\rho(t))+c(t),\qquad
  y=\frac{x-x_c(t)}{\lambda(t)};
\]
\item an absolutely continuous repaired-gauge correction consisting of
\(\rho=\rho(\tau)\), \(\mu(\tau)>0\), and \(q(\tau)\in\mathbb R^3\), with
\[
  \rho_\tau=\mu(\tau)^2,
\]
\[
  V(Y,\tau)=\mu(\tau)U(\mu(\tau)Y+q(\tau),\rho(\tau)),\qquad
  P(Y,\tau)=\mu(\tau)^2Q(\mu(\tau)Y+q(\tau),\rho(\tau)),
\]
and
\[
  G_{\rm sc}(V(\tau))=0,\qquad
  G_j(V(\tau))=0,\quad j=1,2,3.
\]
The condition \(\rho_\tau=\mu^2\) preserves the viscosity coefficient after the
time-dependent scale correction.
\end{enumerate}
\end{definition}

\begin{lemma}[Chart data imply the raw orbit]
\label{lem:c2r-chart-raw-orbit}
Under the chart data in \cref{def:c2r-minimal-representation-data}, \(U,Q\)
satisfy
\[
  \partial_\rho U+(U\cdot\nabla)U+\nabla Q
  =
  \nu\Delta U+A(\rho)(U+y\cdot\nabla U)+B(\rho)\cdot\nabla U,
  \qquad \nabla\cdot U=0,
\]
where
\[
  A(\rho(t))=\lambda(t)\lambda_t(t),\qquad
  B(\rho(t))=\lambda(t)x_c'(t).
\]
\end{lemma}

\begin{proof}
The chart data are the raw orbit data with raw time \(\rho\).  The additive
pressure function \(c(t)\) has zero spatial gradient.  The chain-rule
computation is the one in \cref{lem:c2-raw-renormalized-equation}, with
\(\tau\) replaced by \(\rho\).  Thus spatial derivatives scale as
\(\nabla_xu=\lambda^{-2}\nabla_yU\), \(\Delta_xu=\lambda^{-3}\Delta_yU\), the
convection term scales as \(\lambda^{-3}(U\cdot\nabla)U\), and
\(\rho_t=\lambda^{-2}\).  The terms generated by \(\lambda_t\) and \(x_c'\)
give \(A=\lambda\lambda_t\) and \(B=\lambda x_c'\).  Incompressibility pulls
back as \(\nabla_x\cdot u=\lambda^{-2}\nabla_y\cdot U\).
\end{proof}

\begin{lemma}[Pressure pullback]
\label{lem:c2r-pressure-pullback}
Under the chart data,
\[
  -\Delta Q=\partial_i\partial_j(U_iU_j)
\]
in distributions, modulo functions of \(\rho\).  After any admissible
repaired-gauge correction, the transformed pressure satisfies
\[
  -\Delta P=\partial_i\partial_j(V_iV_j)
\]
in distributions, modulo functions of \(\tau\).
\end{lemma}

\begin{proof}
Taking divergence of the physical Navier--Stokes equation gives
\[
  -\Delta p=\partial_i\partial_j(u_iu_j),
\]
with pressure determined up to a function of time.  Pulling this identity back
through \(p=\lambda^{-2}Q+c(t)\), \(u=\lambda^{-1}U\), and
\(y=(x-x_c)/\lambda\) gives the pressure equation for \(Q\).  If
\[
  V(Y)=\mu U(\mu Y+q),\qquad P(Y)=\mu^2Q(\mu Y+q),
\]
then
\[
  -\Delta_YP
  =\mu^4(-\Delta Q)(\mu Y+q)
  =\mu^4\partial_i\partial_j(U_iU_j)(\mu Y+q)
  =\partial_i\partial_j(V_iV_j).
\]
Adding a function of \(\rho\) or \(\tau\) to the pressure does not affect the
identity.
\end{proof}

\begin{lemma}[Forced modulation coefficients]
\label{lem:c2r-forced-modulation}
Assume the minimal data of \cref{def:c2r-minimal-representation-data}.  Then
\((V,P)\) satisfies
\[
  \partial_\tau V+(V\cdot\nabla)V+\nabla P
  =
  \nu\Delta V
  +a(\tau)(V+Y\cdot\nabla V)
  +b(\tau)\cdot\nabla V,
\]
where
\[
  a(\tau)=\mu(\tau)^2A(\rho(\tau))+\frac{\mu_\tau(\tau)}{\mu(\tau)}
\]
and
\[
  b(\tau)=\mu(\tau)B(\rho(\tau))
      +\mu(\tau)A(\rho(\tau))q(\tau)
      +\frac{q_\tau(\tau)}{\mu(\tau)}.
\]
\end{lemma}

\begin{proof}
Let
\[
  z=\mu(\tau)Y+q(\tau),\qquad
  V(Y,\tau)=\mu(\tau)U(z,\rho(\tau)),\qquad
  \rho_\tau=\mu^2.
\]
The raw equation in \((z,\rho)\) is
\[
  U_\rho+(U\cdot\nabla_z)U+\nabla_zQ
  =
  \nu\Delta_zU+A(U+z\cdot\nabla_zU)+B\cdot\nabla_zU.
\]
The identities
\[
  \nabla_YV=\mu^2\nabla_zU,\qquad
  \Delta_YV=\mu^3\Delta_zU,\qquad
  (V\cdot\nabla_Y)V=\mu^3(U\cdot\nabla_z)U,\qquad
  \nabla_YP=\mu^3\nabla_zQ
\]
together with differentiation in \(\tau\) give
\[
  V_\tau+(V\cdot\nabla_Y)V+\nabla_YP-\nu\Delta_YV
  =
  \left(\frac{\mu_\tau}{\mu}+\mu^2A\right)(V+Y\cdot\nabla_YV)
  +
  \left(\mu B+\mu Aq+\frac{q_\tau}{\mu}\right)\cdot\nabla_YV.
\]
This is the displayed equation.  Since \(\mu,q,\rho\) are absolutely continuous
and \(A,B\) are the raw chart coefficients, \(a,b\) are measurable and locally
integrable on every finite renormalized window.
\end{proof}

\begin{theorem}[Minimal representation theorem]
\label{thm:c2r-minimal-representation}
If a Type II branch satisfying the analytic data has the minimal data of
\cref{def:c2r-minimal-representation-data}, then it has a repaired-gauge
representation in the sense of \cref{def:c2-representation}.
\end{theorem}

\begin{proof}
\Cref{lem:c2r-chart-raw-orbit} gives the raw orbit.  The gauge correction
places the profile on the repaired scale and centering surface, hence gives the
repaired-gauge realization.  \Cref{lem:c2r-pressure-pullback} gives pressure
reconstruction modulo functions of time.  \Cref{lem:c2r-forced-modulation}
gives the final modulation coefficients.  These are the four inputs in
\cref{def:c2-representation}.
\end{proof}

\begin{definition}[Minimal representation failures]
\label{def:c2r-minimal-failures}
Outside the admissible repaired-gauge representation class, the ordered failures
are:
\begin{enumerate}[label=\textup{(\roman*)}]
\item no absolutely continuous concentration chart is available;
\item the chart exists but the repaired scale and centering gauges cannot be
solved admissibly and absolutely continuously.
\end{enumerate}
Pressure-pullback and modulation-coefficient failures are not independent
alternatives when the chart and the absolutely continuous gauge solve exist.
\end{definition}

\begin{corollary}[Ordered minimal representation classification]
\label{cor:c2r-ordered-classification}
Every Type II branch satisfying the analytic data has exactly one ordered
alternative: a repaired-gauge representation, failure of the raw chart, or
failure of the repaired-gauge solve.
\end{corollary}

\begin{proof}
If the chart is absent, the first failure occurs.  If the chart exists but the
admissible absolutely continuous gauge solve is absent, the second failure
occurs.  If both are present, \cref{thm:c2r-minimal-representation} gives the
repaired-gauge representation.  The alternatives are disjoint by their order
and exhaustive for the representation row.
\end{proof}

\subsection{Local Retained Mass Cost Divergence}

\begin{lemma}[Spatially constant retained limits are bounded windows]
\label{lem:c3-spatially-constant-retained-limit-bounded}
Let \(0<r<R_*\), let \(J\Subset(-\theta,\theta)\), and let
\[
  V_j(y,s)=V(y,\sigma_j+s)
\]
be selected retained windows on \(B_{R_*}\times(-\theta,\theta)\).  Assume
\[
  V_j\to c(s)
  \quad\text{strongly in }L^3(B_{R_*}\times(-\theta,\theta)),
  \qquad
  \nabla_y c=0,
\]
and assume the retained time-derivative estimate
\[
  \sup_j\|\partial_s V_j\|_{L^1(J;H^{-m}(B_r))}<\infty
\]
for some \(m\) with \(C_c^\infty(B_r)\subset H^m_0(B_r)\).  Then \(c\) has a
representative in \(BV(J;\mathbb R^3)\), hence \(c\in L^\infty(J)\).  If
\(c\not\equiv0\) on \((-\theta,\theta)\), then after restricting to a smaller
selected parabolic subwindow and recentering its time coordinate, the sequence
has a nonzero bounded selected-window limit in the sense of
\cref{paper1:def:bounded-window-limit}.  Consequently the branch exits the
local Type II class by \cref{paper3:lem:scale-rigid-bounded-limit-exit}.
\end{lemma}

\begin{proof}
Choose \(\zeta\in C_c^\infty(B_r)\) with \(\int_{B_r}\zeta\,dy=1\).  For each
component \(k\), set
\[
  C_{j,k}(s)=\int_{B_r}V_{j,k}(y,s)\zeta(y)\,dy .
\]
The strong \(L^3\) convergence implies \(C_{j,k}\to c_k\) in \(L^1(J)\).  In
distributions on \(J\),
\[
  \frac{d}{ds}C_{j,k}(s)
  =
  \langle \partial_sV_{j,k}(s),\zeta\rangle ,
\]
and therefore
\[
  \operatorname{Var}_J(C_{j,k})
  \le
  \|\zeta\|_{H^m_0(B_r)}
  \|\partial_sV_j\|_{L^1(J;H^{-m}(B_r))}
  \le C_{J,r}.
\]
Lower semicontinuity of total variation under \(L^1\) convergence gives
\(\operatorname{Var}_J(c_k)\le C_{J,r}\).  Since \(c_k\in L^1(J)\), this proves
\(c_k\in BV(J)\).  In one dimension \(BV(J)\subset L^\infty(J)\), with the
bound depending only on \(\|c_k\|_{L^1(J)}\), \(\operatorname{Var}_J(c_k)\), and
\(|J|\).  Thus \(c\in L^\infty(J)\).

If \(c\not\equiv0\), choose a density point \(s_0\in(-\theta,\theta)\) of the
set where \(|c|>0\).  Since the preceding argument applies to every compact
subinterval \(J\Subset(-\theta,\theta)\), choose \(J_0\Subset(-\theta,\theta)\)
containing \(s_0\) and use the \(BV(J_0)\subset L^\infty(J_0)\) representative.
Then choose a radius \(\rho>0\) such that
\[
  B_\rho\times(s_0-\rho^2,s_0)\Subset B_r\times J_0,
  \qquad
  \int_{s_0-\rho^2}^{s_0}|c(s)|^3\,ds>0 .
\]
Replacing the selected centers by \(\sigma_j+s_0\) identifies this subwindow
with \(Q_\rho\).  This finite time recentering is admissible by the same local
selected-window covariance used in
\cref{paper6a:lem:interior-time-selection}: no lower bound is transferred to a
different physical time; the original subwindow is merely rewritten with its
own center as the origin.  The strong convergence remains valid there and the
limit is bounded and nonzero.  This is exactly a nonzero bounded
selected-window limit, so \cref{paper3:lem:scale-rigid-bounded-limit-exit}
applies.
\end{proof}

\begin{theorem}[Local retained-mass cost divergence]
\label{thm:c3-local-retained-mass-cost-divergence}
Let \((V,P,a,b)\) be a repaired-gauge local Type II branch on a terminal tail.
Assume there are \(R_*>0\), \(\theta>0\), \(\eta_*>0\), and an unbounded family
\(\mathcal S\subset[\tau_0,\infty)\) of retained selected active-window centers
such that every \(\sigma\in\mathcal S\) carries positive local velocity mass,
\[
  \int_{\sigma-\theta}^{\sigma+\theta}\int_{B_{R_*}}|V(y,\tau)|^3\,dy\,d\tau
  \ge \eta_* .
\]
Assume also the retained compact-window bounds, the retained
\(L^1H^{-m}\) time-derivative bounds on compact subwindows, the local
Aubin--Lions compactness on \(B_{R_*}\times(-\theta,\theta)\), local windowed
\(H^1\) control, and the scale-rigid routing of nonzero bounded
selected-window limits.  These are all local compact-window hypotheses; no
whole-space critical bound, global profile theorem, or a priori vanishing
dissipation is assumed.  Then finite total localized dissipation on the
retained core is impossible.  In particular, for the canonical identity cost on
a cutoff \(\phi_{R_0}\equiv1\) on \(B_{R_*}\),
\[
  \int_{\tau_0}^{\infty}\tilde{\mathfrak D}_{R_0}(\tau)\,d\tau=\infty
\]
unless the branch has already exited through a bounded selected-window,
rough-core, exterior, or adjacent state-space alternative.
\end{theorem}

\begin{proof}
Assume the localized cost is finite on the terminal tail and no listed exit
has occurred.  Since \(\phi_{R_0}\equiv1\) on \(B_{R_*}\), finite cost gives
finite localized dissipation on \(B_{R_*}\).  Choose retained selected
active-window centers \(\sigma_j\in\mathcal S\), \(\sigma_j\to\infty\), with
disjoint enlarged windows and such that
\[
  \int_{\sigma_j-\theta}^{\sigma_j+\theta}
  \int_{B_{R_*}}|\nabla V|^2\,dy\,d\tau\to0 .
\]
This selection is possible because the total dissipation on the tail is finite
and \(\mathcal S\) is unbounded; if no unbounded retained active-window family
with the stated mass floor remains, the lower active-mass row has failed and
the branch has already exited by the terminal critical-mass routing.
\par
Set \(V_j(y,s)=V(y,\sigma_j+s)\).  The retained compact-window bounds give,
after passing to a subsequence,
\[
  V_j\to V_\infty
  \quad\text{strongly in }L^3(B_{R_*}\times(-\theta,\theta)),
\]
by the local Aubin--Lions compactness in the retained state-space topology.
The vanishing dissipation gives
\(\nabla V_\infty=0\) distributionally on
\(B_{R_*}\times(-\theta,\theta)\).  Hence
\(V_\infty(y,s)=c(s)\) on the retained core cylinder.
\par
If \(c\not\equiv0\), \cref{lem:c3-spatially-constant-retained-limit-bounded}
turns the spatially constant compact limit into a nonzero bounded
selected-window limit on a smaller selected parabolic subwindow.  The
scale-rigid exit theorem \cref{paper3:lem:scale-rigid-bounded-limit-exit}
therefore removes the branch from the retained local Type II path.  If
\(c\equiv0\), strong convergence on the core cylinder contradicts the retained
local mass floor:
\[
  \int_{-\theta}^{\theta}\int_{B_{R_*}}|V_j(y,s)|^3\,dy\,ds\to0
  \quad\text{while}\quad
  \int_{-\theta}^{\theta}\int_{B_{R_*}}|V_j(y,s)|^3\,dy\,ds\ge\eta_* .
\]
Both alternatives are impossible on the retained branch.  Therefore the
localized cost cannot be finite.
\end{proof}

\subsection{Localized cost and the Type II scale-exclusion condition}

\begin{definition}[Localized Type II cost]
\label{def:c4-localized-cost}
For a repaired-gauge Type II branch, let
\[
  \tilde{\mathfrak D}_{R_0}(\tau)
  =
  \nu\int_{\mathbb R^3}|\nabla V(y,\tau)|^2\phi_{R_0}(y)\,dy
  +
  a_+(\tau)\int_{\mathbb R^3}|V(y,\tau)|^2\phi_{R_0}(y)\,dy,
\]
where \(a_+=\max(a,0)\), \(\phi_{R_0}\) is the compact Type II cutoff, and
\(\nu>0\).  A Type II scale-exclusion cost is a measurable nonnegative function
\(\mathfrak C_{\mathrm{II}}\) that is locally integrable on finite
\(\tau\)-intervals and whose divergence on \([\tau_0,\infty)\) gives the Type
II scale-exclusion condition.
\end{definition}

\begin{definition}[Cost comparison]
\label{def:c4-cost-comparison}
A cost comparison for a repaired-gauge Type II branch consists of
\(R_0,\phi_{R_0},\mathfrak C_{\mathrm{II}},c_0,\tau_1\), with \(c_0>0\) and
\(\tau_1\ge\tau_0\), such that
\[
  \mathfrak C_{\mathrm{II}}(\tau)
  \ge c_0\,\tilde{\mathfrak D}_{R_0}(\tau)
  \qquad\text{for a.e. }\tau\ge\tau_1.
\]
The canonical Navier--Stokes choice is
\[
  \mathfrak C_{\mathrm{II}}(\tau)=\tilde{\mathfrak D}_{R_0}(\tau),
  \qquad c_0=1,\qquad \tau_1=\tau_0.
\]
\end{definition}

\begin{lemma}[Identity-cost comparison]
\label{lem:c4-identity-cost}
If \(\tilde{\mathfrak D}_{R_0}\) is measurable, nonnegative, and locally
integrable on finite \(\tau\)-intervals, then the canonical choice
\(\mathfrak C_{\mathrm{II}}=\tilde{\mathfrak D}_{R_0}\) defines a valid Type II
scale-exclusion cost and satisfies the cost comparison with \(c_0=1\) and
\(\tau_1=\tau_0\).
\end{lemma}

\begin{proof}
The assumed measurability, nonnegativity, and local finite-interval
integrability give the cost properties in \cref{def:c4-localized-cost}.  The
identity
\[
  \mathfrak C_{\mathrm{II}}=\tilde{\mathfrak D}_{R_0}
\]
gives the comparison in \cref{def:c4-cost-comparison} with \(c_0=1\) and
\(\tau_1=\tau_0\).
\end{proof}

\begin{theorem}[Cost comparison theorem]
\label{thm:c4-cost-comparison}
Let \((V,P,a,b)\) be a repaired-gauge Type II branch.  Assume:
\begin{enumerate}[label=\textup{(\roman*)}]
\item \(\tilde{\mathfrak D}_{R_0}\) is measurable, nonnegative, and locally
integrable on finite \(\tau\)-intervals;
\item
\[
  \int_{\tau_0}^{\infty}
  \tilde{\mathfrak D}_{R_0}(\tau)\,d\tau=\infty;
\]
\item either the canonical cost
\(\mathfrak C_{\mathrm{II}}=\tilde{\mathfrak D}_{R_0}\) is used, or a cost
comparison as in \cref{def:c4-cost-comparison} is available.
\end{enumerate}
Then the Type II scale-exclusion condition associated with
\(\mathfrak C_{\mathrm{II}}\) holds.
\end{theorem}

\begin{proof}
In the canonical case, \cref{lem:c4-identity-cost} supplies the cost comparison.
In the general case, the comparison is assumed.  Hence there are
\(c_0>0\) and \(\tau_1\ge\tau_0\) such that
\[
  \mathfrak C_{\mathrm{II}}(\tau)
  \ge c_0\tilde{\mathfrak D}_{R_0}(\tau)
\]
for a.e. \(\tau\ge\tau_1\).  Since
\(\tilde{\mathfrak D}_{R_0}\ge0\) and is locally integrable on finite
intervals, divergence on \([\tau_0,\infty)\) implies
\[
  \int_{\tau_1}^{\infty}
  \tilde{\mathfrak D}_{R_0}(\tau)\,d\tau=\infty.
\]
Therefore
\[
  \int_{\tau_1}^{\infty}
  \mathfrak C_{\mathrm{II}}(\tau)\,d\tau
  \ge
  c_0\int_{\tau_1}^{\infty}
  \tilde{\mathfrak D}_{R_0}(\tau)\,d\tau
  =\infty.
\]
Since \(\mathfrak C_{\mathrm{II}}\) is nonnegative and locally integrable on
finite intervals,
\[
  \int_{\tau_0}^{\infty}
  \mathfrak C_{\mathrm{II}}(\tau)\,d\tau
  =
  \int_{\tau_0}^{\tau_1}\mathfrak C_{\mathrm{II}}(\tau)\,d\tau
  +
  \int_{\tau_1}^{\infty}\mathfrak C_{\mathrm{II}}(\tau)\,d\tau
  =\infty.
\]
By definition of the Type II scale-exclusion cost, this divergence is the Type II
scale-exclusion condition.
\end{proof}

\subsection{Local compact-cost classification}

\begin{definition}[Remaining compact cost-divergence tests]
\label{def:c5-compact-tests}
For a repaired-gauge Type II orbit, define the critical-tightness
condition by
\[
  \forall\varepsilon>0\ \forall J_*\Subset\mathbb R\ \exists R_{\varepsilon,J_*}:\quad
  \int_J\int_{|y|>R_{\varepsilon,J_*}}|V(y,\tau)|^3\,dy\,d\tau<\varepsilon
\]
for every selected terminal window \(J\Subset[\tau_0,\infty)\) modeled on
\(J_*\) in the sense of \cref{paper9:def:windowed-critical-tightness}.
Failure of critical tightness is the condition
\[
  \exists\varepsilon_0>0\ \exists J_*\Subset\mathbb R\ \exists J_n\Subset[\tau_0,\infty),\ R_n\to\infty:
  \quad
  \int_{J_n}\int_{|y|>R_n}|V(y,\tau)|^3\,dy\,d\tau\ge\varepsilon_0.
\]
Here each \(J_n\) is modeled on the same fixed compact interval \(J_*\).
Define local windowed \(H^1\) control by
\[
  \forall m\ge1:\quad
  \sup_{n\in\mathbb N}
  \int_{\tau_0+n}^{\tau_0+n+1}
  \|V(\tau)\|_{H^1(B_m)}^2\,d\tau<\infty.
\]
Failure of local windowed \(H^1\) control is the condition
\[
  \exists m\ge1:\quad
  \sup_{n\in\mathbb N}
  \int_{\tau_0+n}^{\tau_0+n+1}
  \|V(\tau)\|_{H^1(B_m)}^2\,d\tau=\infty.
\]
\end{definition}

\begin{definition}[Compact Type II data]
\label{def:c5-compact-typeII-data}
A Type II branch has compact Type II data if
it has:
\begin{enumerate}[label=\textup{(\roman*)}]
\item the Type II analytic data of \cref{def:c1-typeII-analytic-data};
\item the repaired-gauge representation of \cref{def:c2-representation};
\item retained positive finite local critical mass on every compact terminal
window, in the sense of \cref{paper6a:lem:terminal-local-critical-mass-routing};
\item the locally integrable nonnegative cost and cost comparison of
\cref{def:c4-localized-cost,def:c4-cost-comparison}.
\end{enumerate}
This condition does not include critical tightness or windowed \(H^1\) control;
those are the two remaining compact cost-divergence tests.
\end{definition}

\begin{definition}[Routed pre-compact-data exits]
\label{def:c5-precompact-routed-exits}
A routed pre-compact-data exit is the first failed ordered test before the
compact cost-divergence tests are reached.  Thus it is one of the routed C1
exhaustion failures in \cref{def:c1-exhaustion-failures}, one of the ordered
representation failures in \cref{def:c2-representation-failures} or
\cref{def:c2r-minimal-failures}, one of the local terminal critical-mass
failures routed by \cref{paper6a:lem:terminal-local-critical-mass-routing}, or a
failure of local cost well-posedness or cost comparison before
\cref{def:c4-localized-cost,def:c4-cost-comparison} is available.  The
first-failure convention makes these exits disjoint from the compact Type II
data of \cref{def:c5-compact-typeII-data}.
\end{definition}

\begin{lemma}[Pre-compact-data routing]
\label{lem:c5-precompact-data-routing}
Every admissible Type II branch has exactly one of the following ordered
outcomes:
\begin{enumerate}[label=\textup{(\roman*)}]
\item it has compact Type II data in the sense of
\cref{def:c5-compact-typeII-data};
\item it has a routed pre-compact-data exit in the sense of
\cref{def:c5-precompact-routed-exits}.
\end{enumerate}
\end{lemma}

\begin{proof}
Fix an admissible branch.  The ordered C1 routing hypothesis of
\cref{thm:c1-typeII-branch-exhaustion} is supplied by
\cref{paper6a:thm:typeII-analytic-data-exhaustive}.  Applying
\cref{thm:c1-typeII-branch-exhaustion}, if the concentration, scale, and local
profile data are not all present, the first failed C1 test is a routed
pre-compact-data exit.  Otherwise the first component of
\cref{def:c5-compact-typeII-data} holds.

Next run the ordered representation classification
\cref{cor:c2-ordered-representation}, equivalently the minimal representation
classification \cref{cor:c2r-ordered-classification} when the minimal chart
package is used.  A first representation, repaired-gauge, pressure, or
modulation failure is a routed pre-compact-data exit.  If no such failure
occurs, the branch has the repaired-gauge representation of
\cref{def:c2-representation}.

On that repaired-gauge branch, apply the local terminal critical-mass routing
\cref{paper6a:lem:terminal-local-critical-mass-routing}.  Failure of the
retained spacetime activity floor means the candidate was not a retained active
core; an upper local
\(L^3\)-failure produces a comparable retained class or an adjacent active
label.  These are routed pre-compact-data exits unless they are grouped into
the retained terminal component.  In the retained case the positive finite
local critical-mass row holds on every compact terminal window.

Finally test whether the localized cost is well posed and compatible with the
Navier--Stokes identity cost.  Availability of the local cost and comparison is
the fourth component of \cref{def:c5-compact-typeII-data}.  If it is absent,
the first absence is the local cost-well-posedness or cost-comparison failure
listed in \cref{def:c5-precompact-routed-exits}; if it is present, all four
components of compact Type II data hold.  The first-failure convention makes
the two outcomes disjoint and exhaustive.
\end{proof}

\begin{lemma}[Uniform retained active-window mass floor]
\label{lem:c5-uniform-retained-active-window-mass}
Assume an admissible Type II branch has compact Type II data in the sense of
\cref{def:c5-compact-typeII-data} and remains in the retained compact
cost-divergence stratum after the ordered validation tests of
\cref{paper7:lem:no-unrecorded-cost-escape}.  Then there exist
\[
  R_*>0,\qquad \theta>0,\qquad \eta_*>0,
\]
and an unbounded family \(\mathcal S\subset[\tau_0,\infty)\) of retained
selected active-window centers such that
\[
  \int_{\sigma-\theta}^{\sigma+\theta}\int_{B_{R_*}}
  |V(y,\tau)|^3\,dy\,d\tau\ge\eta_*
  \qquad\text{for every }\sigma\in\mathcal S .
\]
If such a fixed retained active-window family does not exist, the branch has
already left the retained compact stratum through a lower-mass, pressure-only,
exterior/separated, cascade, or comparable-class grouping row.
\end{lemma}

\begin{proof}
The retained finite critical-mass component of
\cref{def:c5-compact-typeII-data} is read through
\cref{paper6a:lem:terminal-local-critical-mass-routing}.  Thus the lower active
mass row does not fail on the retained branch; if it did, the candidate would
not be a retained terminal component.
\par
Apply the CKN detection argument in
\cref{paper8:thm:active-profile-mass-floor} to the retained active component on
the selected compact terminal cylinders.  Its proof gives a fixed local
threshold \(\eta_{\rm det}>0\).  If the velocity part of the detected CKN
quantity were not bounded below on a compact subcylinder, then either the
component would be perturbative by
\cref{paper8:cor:small-profiles-perturbative}, or the detected mass would be
pressure-only and would be routed by
\cref{paper8:lem:harmonic-cross-pressure-routing} to the pressure/exterior
row.  Neither row remains in the retained compact cost-divergence stratum.
Therefore each retained active witness carries a positive spacetime velocity
mass on one member of a finite compact covering of the selected terminal core.
\par
The covering is finite and the terminal selection rule groups all comparable
same-point retained classes by
\cref{paper6a:lem:retained-frame-maximality}.  If the active witnesses escaped
every fixed member of this finite compact chart along the tail, the first
failure would be an exterior/separated or cascade row in the ordered local
state space.  Since the branch remains retained, one fixed compact member of
the cover occurs along an unbounded family of selected windows.  Writing that
member as \(B_{R_*}\times(-\theta,\theta)\), and absorbing the scale-invariant
covering constants into
\(\eta_*:=c\,\eta_{\rm det}R_*^2>0\), gives the displayed lower bound for an
unbounded family \(\mathcal S\) of retained active-window centers.
\end{proof}

\begin{lemma}[Local upper estimates give retained Aubin--Lions compactness]
\label{lem:c5-retained-upper-estimates-aubin-lions}
Let \((V,P,a,b)\) be a retained repaired-gauge branch on a terminal tail.  Fix
a compact model interval \(I\Subset\mathbb R\), a radius \(R>0\), and selected
centers \(\sigma_j\to\infty\).  Assume that on the shifted windows
\[
  V_j(y,s)=V(y,\sigma_j+s),\qquad
  P_j(y,s)=P(y,\sigma_j+s),
  \qquad s\in I,
\]
the ordered retained rows supply:
\begin{enumerate}[label=\textup{(\roman*)}]
\item a uniform local critical upper bound
\[
  \sup_j\|V_j\|_{L^\infty(I;L^3(B_{2R}))}<\infty;
\]
\item local windowed \(H^1\) control on \(B_{2R}\);
\item local pressure reconstruction modulo functions of time;
\item the repaired-gauge equation with uniformly bounded selected-window
modulation integrals
\[
  \sup_j\int_I\bigl(|a(\sigma_j+s)|+|b(\sigma_j+s)|\bigr)\,ds<\infty .
\]
\end{enumerate}
Then, after subtracting spatial pressure means, the sequence satisfies
\[
  \sup_j\Big(
  \|V_j\|_{L^\infty(I;L^2(B_{2R}))}
  +\|V_j\|_{L^2(I;H^1(B_{2R}))}
  +\|P_j-c_{j,2R}\|_{L^{3/2}(B_{2R}\times I)}
  \Big)<\infty,
\]
and for some large \(m\),
\[
  \sup_j\|\partial_sV_j\|_{L^1(I;H^{-m}(B_R))}<\infty .
\]
Consequently, after passing to a subsequence,
\[
  V_j\to V_\infty
  \quad\text{strongly in }L^3(B_R\times I).
\]
\end{lemma}

\begin{proof}
The \(L^\infty_sL^2_y(B_{2R})\) bound follows from the local
\(L^\infty_sL^3_y(B_{2R})\) bound and finite measure:
\[
  \|V_j(s)\|_{L^2(B_{2R})}
  \le |B_{2R}|^{1/6}\|V_j(s)\|_{L^3(B_{2R})}.
\]
The \(L^2_sH^1_y(B_{2R})\) bound is the local windowed \(H^1\) control; a fixed
compact model interval \(I\) meets only finitely many unit windows, so the
uniform unit-window bound gives a uniform bound on \(I\).  The pressure bound
is the local Calder\'on--Zygmund pressure reconstruction, modulo functions of
time, applied to \(V_j\otimes V_j\); the finite length of \(I\) and the
uniform \(L^\infty_sL^3_y\) bound give the displayed
\(L^{3/2}_{s,y}\) estimate.
\par
It remains to bound the time derivative.  On \(B_R\) the repaired-gauge equation
with
\[
  a_j(s):=a(\sigma_j+s),\qquad b_j(s):=b(\sigma_j+s)
\]
gives
\[
  \partial_sV_j
  =
  \nu\Delta V_j-\nabla\cdot(V_j\otimes V_j)-\nabla P_j
  +a_j(V_j+y\cdot\nabla V_j)+b_j\cdot\nabla V_j .
\]
Choose \(m\) so large that
\[
  L^1(B_{2R})+L^{3/2}(B_{2R})+H^{-1}(B_{2R})
  \hookrightarrow H^{-m}(B_R).
\]
The diffusion term is bounded in \(L^1(I;H^{-m})\) by the \(L^2H^1\) bound.
The nonlinear term is bounded because
\[
  \|V_j\otimes V_j\|_{L^1(I;L^{3/2}(B_{2R}))}
  \le |I|\|V_j\|_{L^\infty(I;L^3(B_{2R}))}^2 .
\]
The pressure-gradient term is bounded by the pressure estimate.  For the
modulation terms, integrate by parts in the distributional pairing:
\[
  y\cdot\nabla V_j=\nabla\cdot(y\otimes V_j)-3V_j,\qquad
  b_j\cdot\nabla V_j=\nabla\cdot(b_j\otimes V_j).
\]
Thus their \(H^{-m}(B_R)\) norms are bounded by
\[
  C_R\bigl(|a_j(s)|+|b_j(s)|\bigr)\|V_j(s)\|_{L^2(B_{2R})}.
\]
The retained modulation row gives \(a_j,b_j\in L^1(I)\) with uniform selected
window bounds; if this row failed, the ordered validation would route the
branch to the modulation/scale-drift exit before the retained compact stratum
was declared.  Hence \(\partial_sV_j\) is bounded in \(L^1(I;H^{-m}(B_R))\).
\par
The Aubin--Lions--Simon theorem applied to
\[
  H^1(B_R)\Subset L^2(B_R)\hookrightarrow H^{-m}(B_R)
\]
gives strong convergence in \(L^2(B_R\times I)\), after passing to a
subsequence.  The same bounds give a uniform \(L^{10/3}(B_R\times I)\) bound by
the standard parabolic interpolation between \(L^\infty L^2\) and \(L^2H^1\).
Interpolating the strong \(L^2\) convergence with this uniform
\(L^{10/3}\)-bound gives strong convergence in \(L^p\) for every
\(2\le p<10/3\), in particular in \(L^3(B_R\times I)\).
\end{proof}

\begin{lemma}[Retained compact data supply the retained mass and compactness package]
\label{lem:c5-retained-data-supply-aubin-lions}
Assume an admissible Type II branch has compact Type II data in the sense of
\cref{def:c5-compact-typeII-data}, satisfies the two compact tests of
\cref{def:c5-compact-tests}, and remains in the retained compact
cost-divergence stratum after the ordered validation tests of
\cref{paper7:lem:no-unrecorded-cost-escape}.  Then, on every compact retained
terminal window after a harmless terminal restriction, the requirements of
\cref{thm:c3-local-retained-mass-cost-divergence} are satisfied.  More precisely:
\begin{enumerate}[label=\textup{(\roman*)}]
\item there are \(R_*>0\), \(\theta>0\), \(\eta_*>0\), and an unbounded retained
selected active-window family \(\mathcal S\) satisfying the positive spacetime
local mass floor of \cref{thm:c3-local-retained-mass-cost-divergence};
\item the local compact-cylinder bounds
\[
  V\in L^\infty(I;L^3(B_{2R}))\cap L^\infty(I;L^2(B_{2R}))
  \cap L^2(I;H^1(B_{2R})),
  \qquad
  P-c_{2R}\in L^{3/2}(B_{2R}\times I)
\]
hold for each fixed \(R\) and compact terminal interval \(I\);
\item on smaller balls, \(\partial_\tau V\) is bounded in
\(L^1(I;H^{-m}(B_R))\) for some large \(m\), hence the local Aubin--Lions
compactness used in
\cref{thm:c3-local-retained-mass-cost-divergence} is available;
\item any nonzero bounded selected-window limit is routed out of the retained
Type II branch by \cref{paper3:lem:scale-rigid-bounded-limit-exit}.
\end{enumerate}
If any item in this list fails before the retained stratum is reached, the first
failure is one of the named representation, pressure, modulation, critical-mass,
rough-core, exterior, or adjacent compact-cost exits of
\cref{paper7:lem:no-unrecorded-cost-escape}; it is not repaired by inserting a
global estimate.
\end{lemma}

\begin{proof}
The positive spacetime retained-window mass floor is
\cref{lem:c5-uniform-retained-active-window-mass}.  Its proof uses the active
CKN detection threshold, pressure-only routing, and the terminal maximality rule;
if the fixed unbounded active-window family does not exist, the branch has
already left the retained compact stratum through a named first-failure row.

The repaired-gauge representation, pressure reconstruction, and modulation data
are item \textup{(ii)} of \cref{def:c5-compact-typeII-data} together with the
passed representation rows of \cref{paper7:lem:no-unrecorded-cost-escape}.  The
two compact tests give critical tightness and local windowed \(H^1\) control.
Equivalently, by
\cref{paper6a:thm:retained-compact-tests-closed,paper6a:cor:compact-test-closure-discharged},
their failure is a named non-retained alternative.  Combining the retained local
\(L^3\) upper bound, local windowed \(H^1\) control, local pressure
reconstruction, and the retained modulation-window bound of
\cref{paper6a:lem:modulation-window-bound-or-exit},
\cref{lem:c5-retained-upper-estimates-aubin-lions} gives the displayed local
compact-cylinder bounds, the \(L^1H^{-m}\) time-derivative bound, and the
strong \(L^3\) Aubin--Lions compactness on selected windows.  This compactness
lemma uses only local upper bounds and local PDE estimates; it does not require
the pointwise-in-time lower bound in
\cref{paper6a:thm:bounded-terminal-sequences-from-l3}.

Finally, a nonzero bounded selected-window limit leaves the strict
selected-window retained branch by
\cref{paper3:lem:scale-rigid-bounded-limit-exit}; on a physical germ it is
then closed by \cref{paper6:thm:bounded-selected-window-physical-routing}.  Hence all
hypotheses used in \cref{thm:c3-local-retained-mass-cost-divergence} are local
consequences of the retained compact data and validation rows.
\end{proof}

\begin{lemma}[Compact positive case is excluded]
\label{lem:c5-compact-positive-excluded}
Assume an admissible Type II branch has compact Type II data and satisfies both
compact cost-divergence tests in \cref{def:c5-compact-tests}.  If the branch
remains in the retained compact cost-divergence stratum after the ordered
validation tests of \cref{paper7:lem:no-unrecorded-cost-escape}, then the
branch satisfies the Type II exclusion conclusion.
\end{lemma}

\begin{proof}
The compact Type II data give the supercritical scale alternative,
a repaired-gauge orbit \((V,P,a,b)\), retained positive local critical mass,
local compact-cylinder critical upper bounds, and the cost comparison.  By
\cref{lem:c5-retained-data-supply-aubin-lions}, the retained branch also has the
local compact-window, pressure, \(H^1\), time-derivative, and bounded-limit
routing package required by
\cref{thm:c3-local-retained-mass-cost-divergence}.
Thus
\[
  \int_{\tau_0}^{\infty}
  \tilde{\mathfrak D}_{R_0}(\tau)\,d\tau=\infty.
\]
\Cref{thm:c4-cost-comparison} converts this divergence into the Type II
scale-exclusion condition, hence the validated compact cost-divergence
conclusion.  Since the branch remains in the retained compact cost-divergence
stratum, \cref{paper7:cor:retained-stratum-cost-exclusion} gives the Type II
exclusion conclusion.
\end{proof}

\clearpage
\bibliographystyle{alpha}
\bibliography{type_II_regularity}

\end{document}
